\documentclass[11pt,reqno]{amsart}
\usepackage[T1]{fontenc}
\usepackage{lmodern,amsmath,amscd,amssymb,amsthm,mathtools,mathrsfs,microtype,booktabs,longtable}
\usepackage[margin=1.08in]{geometry}
\usepackage{xr-hyper}
\usepackage[hidelinks]{hyperref}
\numberwithin{equation}{section}
\newtheorem{theorem}{Theorem}[section]
\newtheorem{proposition}[theorem]{Proposition}
\newtheorem{lemma}[theorem]{Lemma}
\newtheorem{corollary}[theorem]{Corollary}
\theoremstyle{definition}\newtheorem{definition}[theorem]{Definition}
\newtheorem{example}[theorem]{Example}
\theoremstyle{remark}\newtheorem{remark}[theorem]{Remark}
\newcommand{\C}{\mathbb C}\newcommand{\Q}{\mathbb Q}
\newcommand{\Pj}{\mathbb P}\newcommand{\A}{\mathbb A}
\newcommand{\OO}{\mathcal O}\newcommand{\II}{\mathcal I}
\newcommand{\NN}{\mathcal N}\newcommand{\LL}{\mathcal L}
\newcommand{\ZZ}{\mathcal Z}\newcommand{\RR}{\mathcal R}
\DeclareMathOperator{\Sym}{Sym}\DeclareMathOperator{\Gr}{Gr}
\DeclareMathOperator{\Hom}{Hom}\DeclareMathOperator{\Tot}{Tot}
\DeclareMathOperator{\Fitt}{Fitt}\DeclareMathOperator{\Ann}{Ann}
\DeclareMathOperator{\rank}{rank}\DeclareMathOperator{\Proj}{Proj}
\DeclareMathOperator{\Spec}{Spec}\DeclareMathOperator{\PGL}{PGL}
\DeclareMathOperator{\Pic}{Pic}\DeclareMathOperator{\diag}{diag}
\DeclareMathOperator{\Tor}{Tor}\DeclareMathOperator{\im}{im}
\DeclareMathOperator{\coker}{coker}\DeclareMathOperator{\codim}{codim}
\DeclareMathOperator{\End}{End}
\DeclareMathOperator{\disc}{disc}\DeclareMathOperator{\Disc}{Disc}
\DeclareMathOperator{\Supp}{Supp}\DeclareMathOperator{\Hilb}{Hilb}
\DeclareMathOperator{\rad}{rad}\DeclareMathOperator{\Ass}{Ass}
\DeclareMathOperator{\gr}{gr}\DeclareMathOperator{\Gal}{Gal}
\DeclareMathOperator{\length}{length}
\allowdisplaybreaks[2]
\setlength{\emergencystretch}{3em}

\begin{document}
\title[Finite failure schemes and quadratic pencils]{Finite failure schemes and the reconstruction of quadratic pencils}
\author{Qian Qi}
\date{September 26, 2026}
\subjclass[2020]{14M12, 14B05, 14D20, 13C40}
\keywords{Quadratic pencil, finite failure algebra, intrinsic reconstruction, common divisor, conductor, normalization fibre}
\hypersetup{pdftitle={Finite failure schemes and the reconstruction of quadratic pencils, revision 170},pdfauthor={Qian Qi}}
\begin{abstract}
We reconstruct quadratic pencils from their sharp unmarked finite
failure algebras and analyze the associated multiplication graphs.
The complete first nonreduced Hilbert fibre has reduced components
$\Pj^4$ and $\mathbb F_2$ meeting along doubled lines, together with
a square-zero nilpotent ideal. Products give all fibres with smaller
corank-two Smith exponents at most two. An explicit Hilbert--Burch
chart, a fixed-target embedded comparison, and a connectedness
argument establish the full scheme statement. This master contains
the complete proofs of both companion papers and their supplementary
results. We also determine the collision of
two reduced incidences, separating its flat horizontal closure from
its vertical conic excess and identifying the global nilradical sheaf.
Universal horizontal models, all-arc double-contact presentations,
the nonsplit algebra extension, and arbitrary ramification are
proved in the additional sections.
The determinantal model for every polynomial contact order, its finite
jet and decorated chamber theorems, and the complete nonsymmetric
two-wall slice are included without removing predecessor proofs.
The all-order marked-contact chain, its normal Rees algebra,
punctual genus-correction rings, and fixed-degree equivariant gluing
are also included in full.
The normalization, full parameter fibre, global conductor, and root-label
identifications of every marked coefficient cover are included as well.
\end{abstract}
\maketitle
\section*{The paired manuscripts}
\label{sec:introduction-v153}
The sharp inverse is the main theorem of the first paper. The second
paper's new main results are
Theorems~\ref{thm:complete-fibre-v163} and
\ref{thm:product-fibres-v163}, based on
Proposition~\ref{prop:embedded-functors-v163} and
Lemma~\ref{lem:HB-family-v163}. This unified source is a preservation
master, not a third submission. Prior proof blocks are retained,
while the focused companion sources provide the primary reading order.
The effective-family application uses reconstruction; it does not
assert an internal boundary operation before it. The Ballico 1993
full-text comparison and an external independent audit of the whole
inverse proof remain uncompleted.

\subsection{Collision families}
Theorem~\ref{thm:collision-family-v164} identifies a full discriminant
pullback and its horizontal closure. Theorem~\ref{thm:ordered-resolution-v164}
gives its ordered normalization after a quadratic base change.
Theorem~\ref{thm:nilradical-sheaf-v164} identifies the global square-zero
sheaf, and Proposition~\ref{prop:genuine-collision-v164} supplies a
fixed-target pencil realization. The parameter-space semistable
model is kept distinct from the universal embedded curve.

% BEGIN PRESERVED PART 33-first-relations-v147.tex
\subsection{Universal horizontal specialization}
The new primary results are Theorems~\ref{thm:horizontal-universal-v166},
\ref{thm:all-arcs-v166}, \ref{thm:extension-class-v166}, and
\ref{thm:ramified-v166}. The complete companion papers supply their
focused reading order; this master retains every preceding proof.
\section{First relations, ungraded isomorphisms, and inverse systems}
\label{sec:first-relations-v147}

The passage from a finite local algebra to its first relation space is
independent of pencils.  We isolate it before undertaking the geometric
inverse.  In this section \(E\) denotes a space of degree-one
\emph{generators}; thus it is a cotangent space, rather than the dual
normal space used later.  Let \(e=\dim E\), \(d\ge2\), and
\(K\subset\Sym^d E\).  Set
\begin{equation}\label{eq:general-first-algebra-v147}
 A(K)=\Sym E/\bigl((K)+\mathfrak m^{d+1}\bigr),
 \qquad \mathfrak m=(E),\qquad \mathfrak n=\mathfrak m A(K).
\end{equation}
The grading in this presentation is not supplied to an isomorphism.

\begin{proposition}[The first-relation orbit principle]
\label{prop:first-relation-orbit-v147}
The abstract algebra \(A(K)\) recovers \(E\) and \(K\), up to
linear isomorphism, by
\begin{equation}\label{eq:intrinsic-first-relation-v147}
 E=\mathfrak n/\mathfrak n^2,\qquad
 K=\ker\bigl[\Sym^d(\mathfrak n/\mathfrak n^2)
                                    \longrightarrow\mathfrak n^d\bigr].
\end{equation}
Consequently \(A(K)\simeq A(K')\) as ungraded complex algebras
if and only if some \(a:E\xrightarrow{\sim}E'\) satisfies
\((\Sym^d a)K=K'\).

More precisely, after fixing the homogeneous inclusions of generators,
all isomorphisms with tangent map \(a\) are exactly the substitutions
\[
                 x\longmapsto a(x)+u(x),\qquad
                 u\in\Hom_{\C}(E,\mathfrak n'^2).
\]
These assertions about substitutions remain valid after scalar
extension to any commutative complex algebra.
\end{proposition}
\begin{proof}
The positive ideal is the unique maximal ideal, and
\(\mathfrak n^{d+1}=0\).  Multiplication of \(d\) lifts of
cotangent vectors is independent of their lifts: replacing any one
by an element of \(\mathfrak n^2\) changes the product in
\(\mathfrak n^{d+1}\).  There are no relations of smaller degree,
and the degree-\(d\) quotient is \(\Sym^d E/K\).  This proves
\eqref{eq:intrinsic-first-relation-v147}.  In particular every
isomorphism induces the stated linear equivalence.

Conversely substitute \(a(x)+u(x)\) for each generator.  A relation
of degree \(d\) changes, beyond its linear substitution, only in
degrees at least \(d+1\).  Those terms vanish.  The relations of
degree \(d+1\) also vanish.  The induced map on associated graded
algebras is the isomorphism induced by \(a\); a finite filtration
then makes the algebra map invertible.  Every map with this tangent
map has precisely such generator images.  These polynomial identities
hold over any commutative complex coefficient algebra.  The intrinsic
preservation of the augmentation under that scalar extension is
verified in the next proof, so no preservation of a chosen origin has
tacitly been added to the isomorphism problem.
\end{proof}

\begin{proposition}[The full automorphism group scheme]
\label{prop:first-relation-group-v147}
Write \(\ell=\dim_{\C}A(K)=\binom{e+d}{d}-\dim K\) and
\(H_K=\operatorname{Stab}_{\operatorname{GL}(E)}(K)\), with its
closed subgroup scheme structure.  The functor
\(D\mapsto\operatorname{Aut}_{D}(A(K)\otimes_{\C}D)\) is an
affine group scheme of finite type and fits into a split exact sequence
of affine group schemes over \(\C\):
\begin{equation}\label{eq:first-group-v147}
 1\longrightarrow U_K\longrightarrow\operatorname{Aut}(A(K))
       \longrightarrow H_K\longrightarrow1.
\end{equation}
The kernel is unipotent.  Choosing the homogeneous lifting of \(E\)
identifies its underlying scheme, but not generally its group law,
with the affine space associated to \(\Hom_{\C}(E,\mathfrak n^2)\).
Its dimension is \(e(\ell-1-e)\).  Intrinsically, \(U_K\) acts
simply transitively on the affine space of linear splittings of
\(\mathfrak n\twoheadrightarrow\mathfrak n/\mathfrak n^2\).
The splitting of \eqref{eq:first-group-v147} uses a choice of such a
lifting; it is not an additional intrinsic grading.
\end{proposition}
\begin{proof}
Choose a basis of the finite-dimensional vector space \(A=A(K)\).
Inside \(\operatorname{GL}(A)\), the equations
\(f(1)=1\) and \(f(xy)=f(x)f(y)\), imposed on basis elements,
are polynomial equations.  They represent the automorphism functor,
including its values on nonreduced complex algebras.

The augmentation can be recovered without mistaking the absolute
nilradical after scalar extension for the augmentation ideal.  If
\(M_z\) denotes multiplication by \(z\), the homogeneous filtration
gives
\begin{equation}\label{eq:trace-augmentation-v147}
                      \operatorname{tr}(M_z)=\ell\,\epsilon(z).
\end{equation}
Indeed multiplication by a positive-degree element is strictly
triangular, whereas multiplication by a scalar has that scalar on
every diagonal entry.  Identity~\eqref{eq:trace-augmentation-v147}
holds over every commutative complex algebra \(D\).  Algebra
automorphisms conjugate multiplication matrices, so they preserve
\(\epsilon=\ell^{-1}\operatorname{tr}(M_{(-)})\), its kernel
\(\mathfrak n_D\), and all its powers.  Thus the tangent map is
a morphism of group schemes into \(H_K\), not merely a map on
complex points.

Linear substitutions give a morphism \(H_K\to\operatorname{Aut}(A)\)
sectioning this tangent map.  By Proposition~\ref{prop:first-relation-orbit-v147},
its kernel has, functorially in \(D\), the unique parametrization
\(x\mapsto x+u(x)\) with \(u:E_D\to\mathfrak n_D^2\).
Reading generator images and making these substitutions are inverse
polynomial maps.  On a basis ordered by the positive filtration all
kernel matrices are unitriangular, proving unipotence.  Multiplication
is substitution, which need not be addition of corrections.  The
dimension follows from \(\dim\mathfrak n^2=\ell-1-e\).

Finally two splittings of \(\mathfrak n\twoheadrightarrow E\)
differ by exactly a map \(E\to\mathfrak n^2\).  Any splitting
generates \(A\), and its degree-\(d\) products impose exactly
\(K\); the preceding substitution argument therefore gives the
unique element of \(U_K\) relating any two splittings.  This is
the asserted coordinate-free torsor description.
\end{proof}

\subsection{The inverse-system description and its precise role}
Put \(E^{\vee}=\Hom_{\C}(E,\C)\).  Use the apolar pairing in
which dual divided monomials satisfy
\(\langle x^{\alpha},y^{\beta}/\beta!\rangle=\delta_{\alpha\beta}\).
Here the pairing means evaluation at the origin after applying the
constant-coefficient differential operator.  For
\(I=(K)+\mathfrak m^{d+1}\), its Macaulay inverse system is
\begin{equation}\label{eq:inverse-system-v147}
 I^{\perp}=\bigoplus_{j<d}\Sym^j E^{\vee}\ \oplus\ K^{\perp},
              \qquad K^{\perp}\subset\Sym^d E^{\vee}.
\end{equation}
To check this equality, the relations of degree \(d+1\) bound the
degree of every inverse-system element by \(d\).  The degree-\(d\)
relations pair only with its top component.  All lower components are
unrestricted.  Conversely that direct sum is stable under every
partial derivative: differentiating its top part lands in the full
lower-degree space.  It is therefore annihilated by \(I\), proving
the displayed equality.  Thus the orbit problem for \(K\) is
equivalently the contragredient orbit problem for \(K^{\perp}\).
This uses the classical inverse-system correspondence; it does not
assert a new canonical-grading theorem.

Elias and Rossi prove that Gorenstein local algebras with Hilbert
function \((1,e,e,1)\) are canonically graded
\cite[Theorem~3.3 and Corollaries~3.4--3.5]{EliasRossiShort}.
Their proof removes a lower-degree term of the inverse-system
generator by a nonlinear coordinate change.  In their later work,
compressed Gorenstein algebras of socle degree at most four are
canonically graded, with counterexamples beyond the stated range
\cite[Theorem~3.1 and Examples~3.3--3.5]{EliasRossiCompressed}.
In contrast, \eqref{eq:general-first-algebra-v147} is homogeneous by
construction, and its truncation makes all higher substitutions
invisible to its first relation.  We do not deduce that an arbitrary
filtered deformation with the same Hilbert function is graded.

For the pencil algebras, \(e=n^2\), \(d=n^2+2n-4\ge11\), and
\(\dim K=\binom N2\).  They are not Gorenstein: their socles contain
\(\mathfrak n^d\), whose dimension is
\(\binom{e+d-1}{d}-\binom N2>1\).  For example, the first binomial
is at least \(\binom{e+1}{2}\), while \(e=n^2>N\).  The issue
left by Proposition~\ref{prop:first-relation-orbit-v147} is accordingly
not the removal of a nonhomogeneous inverse-system term.  It is the
recovery of the matrix factors and defining coefficient data from the
\emph{unrestricted} \(\operatorname{GL}(E)\)-orbit of \(K\).
The next sections establish that recovery and its global descent.

% BEGIN PRESERVED PART 00-ruling-foundations-v145.tex
\section{Intrinsic tensor rulings}
\label{sec:ruling-foundations-v143}
The first step of the inverse problem is to recognize the two tensor
factors of a determinant cone without an ambient marking.  The following
scheme-theoretic ruling calculation and descent lemma supply precisely
this step.  They apply to the pencil argument independently of the web
and exterior-contraction constructions in the separate research archive.

\begin{lemma}[The relative scheme of maximal rank-one spaces]
\label{lem:rank-one-fano-scheme}
Let \(B\) be any complex scheme and \(E,F\) vector bundles of
rank \(n\ge2\).  In \(\Gr_B(n,E\otimes F)\), let
\(Z(E,F)\) be the zero scheme of the restriction of the quadratic
rank-one equations to the universal subbundle \(\mathcal T\).
Equivalently it is the zero scheme of
\[
 \Sym^2\mathcal T\longrightarrow
              \bigwedge^2E\otimes\bigwedge^2F.
\]
The natural morphism
\begin{equation}\label{eq:fano-rank-one-scheme}
 \Pj_B(E)\ \sqcup\ \Pj_B(F)\longrightarrow Z(E,F),
 \qquad \ell\mapsto\ell\otimes F,
 \quad m\mapsto E\otimes m,
\end{equation}
is an isomorphism.  It commutes with arbitrary change of complex
base, including change to a nonreduced base.
\end{lemma}
\begin{proof}
We first work over \(\C\).  Each displayed map to the
Grassmannian is a closed immersion.  Its composition with the
Pluecker embedding is the degree-\(n\) Veronese embedding on
\(\ell^{\otimes n}\otimes\det F\), respectively
\(\det E\otimes m^{\otimes n}\), followed by its natural linear
inclusion into \(\bigwedge^n(E\otimes F)\).  This linear inclusion
is injective: polarization of the nonzero pure-power map gives the
corresponding irreducible symmetric-power summand.  The rank-one
equations vanish identically on both universal families, so the
closed immersions factor through \(Z(E,F)\).

Every geometric point is in one of the two images.  Indeed the sum
of two rank-one tensors is rank at most one only if their first
factors or their second factors are proportional.  Starting with two
independent tensors, this observation forces a linear space consisting
of rank-one tensors to fix the same factor throughout.  A space of
dimension \(n\) is then exactly one of the displayed spaces.
The two images are disjoint: a space fixing both factors has dimension
one, not \(n\).

We verify that these pointwise statements omit no infinitesimal
structure.  At \(\ell\otimes F\), choose a generator of \(\ell\)
and a complement in \(E\).  A Grassmannian tangent vector is a map
\(\phi:F\to(E/\ell)\otimes F\).  For each coordinate of
\(E/\ell\), write its component as \(A:F\to F\).
Linearizing the restricted minors gives
\[
                       u\wedge A(u)=0\quad\text{for every }u\in F.
\]
Polarizing this diagonal identity gives
\[
               u\wedge A(v)+v\wedge A(u)=0\quad(u,v\in F).
\]
In characteristic zero the diagonal and polarized identities are
equivalent.  These are precisely all coefficients of the bilinear
linearization of the restricted quadratic map on \(\Sym^2F\),
so no infinitesimal equation is omitted by testing the diagonal.
Taking the basis vectors shows that \(A\) is diagonal; taking their
pairwise sums makes all its diagonal entries equal.  Conversely a
scalar \(A\) satisfies all the linearized equations.  Thus the
Zariski tangent space to \(Z(E,F)\) is precisely
\(\Hom(\ell,E/\ell)\), of dimension \(n-1\).
The other image has the symmetric calculation.

At every closed point the closed immersed projective space gives
local dimension at least \(n-1\), whereas the tangent calculation
gives embedding dimension exactly \(n-1\).  The local ring is
therefore regular.  A finite-type complex scheme whose closed-point
local rings are regular is reduced: a nonempty support of its
nilradical would contain a closed point.  In particular there is no
nilpotent thickening or embedded associated structure at either
image.  The disjoint union of the two closed immersions is itself a
closed immersion, surjective on geometric points.  Its ideal in the
now reduced target vanishes at every closed point, so it is zero.
This proves the scheme isomorphism over \(\C\).

For general \(B\), trivialize \(E,F\) on a common open cover.
The Grassmannian and the zero scheme of the displayed universal
quadratic map are, on such an open, the base change of the constant
construction just proved.  Thus \eqref{eq:fano-rank-one-scheme} is
an isomorphism there even if the base has nilpotents.  The maps are
natural under changes of frame and hence glue.  This argument is
an identification by equations and base change, not an inference of
relative flatness from a fibrewise tangent calculation.
\end{proof}


\begin{lemma}[Oriented relative Segre descent]
\label{lem:oriented-segre-descent}
Let \(B\) be a connected projective complex variety with
\(H^0(B,\OO_B)=\C\), and let \(V,V'\) have the same dimension
\(n\ge2\).  Let \(\mathcal U,\mathcal U'\) be vector bundles
of rank \(n\) on \(B\).  Suppose a linear bundle isomorphism
\[
 \alpha:V\otimes\mathcal U^*\xrightarrow{\sim}
                         V'\otimes\mathcal U'^*
\]
preserves the rank-one cones and the ruling whose parameter bundle
is \(\Pj(V)\times B\).  The induced isomorphism
\(\Pj(V)\times B\simeq\Pj(V')\times B\) is regular and is
induced by one constant projective isomorphism \([g]:\Pj(V)\to\Pj(V')\).
Locally \(\alpha(T)=gTh^{-1}\), with the same \([g]\) on every
chart.  The possible line-bundle twists of vector-bundle lifts do
not change the projective left coefficient lines of any homogeneous
functorial coefficient construction.
\end{lemma}
\begin{proof}
Lemma~\ref{lem:rank-one-fano-scheme}, applied to
\(E=V\otimes\OO_B\) and \(F=\mathcal U^*\), identifies the
closed relative Fano scheme with
\((\Pj(V)\times B)\sqcup\Pj_B(\mathcal U^*)\).
In particular its two components carry their displayed scheme structures,
with no additional nilpotent or embedded components.
The isomorphism \(\alpha\) induces a regular map of these relative
parameter schemes, and the orientation hypothesis preserves the
first component.

For completeness, an isomorphism of two trivial \(\Pj^{n-1}\)-bundles
over any base gives a regular projective-matrix-valued map, not
merely a fibrewise collection.  Its pullback of the relative
\(\OO(1)\) has the form \(\OO(1)\otimes p^*L\) for a line
bundle \(L\) on the base.  To see this, tensor by \(\OO(-1)\);
the resulting line bundle is trivial on every projective-space
fibre, and its pushforward and evaluation identify it with a
pullback from the base.  Pushing forward \(\OO(1)\) now gives
an isomorphism of the corresponding trivial rank-\(n\) bundles,
with that line-bundle twist.  On charts trivializing \(L\), it
is an invertible matrix \(g_i\) of regular functions.  On overlaps,
\(g_j=u_{ij}g_i\) for an invertible regular function \(u_{ij}\).
The classes \([g_i]\) thus glue to a regular morphism to the
projective-isomorphism scheme, identified after choosing fixed
bases with \(\PGL_n\).  This accounts explicitly for the possible
absence of a global linear lift.

The group \(\PGL_n\) is affine: it is the determinant-nonzero
principal open in the projective space of matrices.  A morphism
from \(B\) to an affine scheme factors through
\(\Spec H^0(B,\OO_B)=\Spec\C\), so the projective map is constant.
Choose a constant linear lift \(g\).  Each local \(g_i\) then
differs from \(g\) by a scalar unit.  After removing its two
projective-factor actions, a linear map preserving every
decomposable line is scalar: apply it to sums of two tensors with
one common factor.  Consequently the remaining action is a right
bundle map, and the scalar units may be absorbed in its local
lifts \(h\).  This proves the local tensor formula.

Finally a scalar multiplying a lift acts on a homogeneous functor
of degree \(a\) by its scalar to the power \(a\) (or \(-a\)
on its dual).  It cannot alter a projective coefficient line.
The conclusion applies to every homogeneous degree.  In particular,
no right change of frame can replace the single projective left map
by separate transformations on different coefficient spaces.
\end{proof}

% BEGIN PRESERVED PART 34-coefficient-descent-v147.tex
\section{Actual factor descent and moving coefficient spaces}
\label{sec:coefficient-descent-v147}

We now separate the linear-preserver input, the descent of vector
bundles, and the recovery of coefficients.  The resulting principle
allows several varying coefficient spaces, of arbitrary prescribed
ranks, rather than only the line associated with a pencil.

\begin{lemma}[Linear preservers of the determinant cone]
\label{lem:linear-preserver-v147}
Let \(U,V,U',V'\) be complex vector spaces of dimension \(n\ge2\).
An invertible linear map
\(\Phi:\Hom(U,V)\to\Hom(U',V')\) carrying the determinant
hypersurface onto the determinant hypersurface has one of the forms
\[
 T\longmapsto gTh^{-1},\qquad
 T\longmapsto gT^{\mathsf t}h^{-1},
\]
where the second formula is written after choosing identifications
for transposition.  The two cases respectively preserve and exchange
the two rank-one rulings.  In the first case the pair \((g,h)\)
is unique up to a common nonzero scalar.
\end{lemma}
\begin{proof}
At a matrix of rank \(r\), an invertible \(r\)-minor reduces the
local rank condition to the rank condition on a generic
\((n-r)\)-square matrix.  The singular locus of the rank-at-most
\(s\) variety is therefore the rank-at-most \(s-1\) variety for
\(1\le s<n\): on the rank-\(s\) stratum the Schur-complement
entries give independent equations, whereas on a lower stratum the
linear terms of its defining \((s+1)\)-minors vanish.  Starting
with the determinant cone and iterating reduced singular loci hence
recovers the rank-one cone.  For \(n=2\) that cone is already the
determinant cone.

If \(v\otimes u\) and \(w\otimes z\) have rank one, their sum
has rank at most one only if \(v,w\) or \(u,z\) are dependent.
It follows that each maximal linear space of rank-one tensors has
one factor fixed.  The two families are
\(\Pj(V)\) and \(\Pj(U^*)\).  The induced Segre automorphism
must preserve these families or exchange them.  In the preserving
case its restriction to a ruling gives projective linear maps on
each factor.  Remove linear lifts of both maps.  The remaining
linear transformation fixes each decomposable line.  In tensor bases
its eigenvalues at \(v_i\otimes u_j\) agree along each row and
each column by applying it to
\((v_i+v_k)\otimes u_j\) and \(v_i\otimes(u_j+u_l)\).
They are all equal, so the remaining map is scalar.  This proves
the first form and its uniqueness.  Composing with transposition
reduces the exchanging case to the first.

This is the square invertible case of the classical rank-one
preserver theorem of Marcus and Moyls
\cite[Theorem~1, pp.~1218--1219]{MarcusMoyls}; the argument above
specifies the reduction from the determinant hypersurface that is
needed here.
\end{proof}

\begin{lemma}[An actual tensor map removes the line twist]
\label{lem:actual-factor-descent-v147}
Let \(B\) be a reduced complex variety, \(V\) a complex vector
space with \(\dim V\ge2\), and \(\mathcal A_1,\mathcal A_2\)
vector bundles.  Suppose an actual bundle isomorphism
\[
 \Psi:V\otimes\mathcal A_1\xrightarrow{\sim}V\otimes\mathcal A_2
\]
preserves the rank-one cones and induces the identity on the
\(\Pj(V)\)-parameter space of rulings.  There is a unique bundle
isomorphism \(a:\mathcal A_1\to\mathcal A_2\) with
\(\Psi=1_V\otimes a\).

Equivalently, the intrinsic operation of taking the commutant of the
constant matrix algebra gives
\begin{equation}\label{eq:commutant-descent-v147}
 \mathcal Hom_{\End(V)\otimes\OO_B}
       (V\otimes\mathcal A_1,V\otimes\mathcal A_2)
             \simeq\mathcal Hom_{\OO_B}(\mathcal A_1,\mathcal A_2).
\end{equation}
Thus after a constant left projective map has been lifted linearly,
no residual line-bundle twist is left in the recovery of the right
factor from \(\Psi\).
\end{lemma}
\begin{proof}
For every constant line \(\C v\subset V\), the corresponding
ruling is fixed, so on every geometric fibre \(\Psi\) maps
\(v\otimes\mathcal A_1\) onto \(v\otimes\mathcal A_2\).
Choose a basis \(v_1,\ldots,v_n\) for the purpose of checking
this assertion as an equality of bundle maps.  The rulings for
\(v_i\) force all off-diagonal blocks of \(\Psi\) to vanish.
The rulings for \(v_i+v_j\) force its diagonal blocks to agree.
Since \(B\) is reduced, fibrewise zero morphisms of vector bundles
are zero.  Therefore \(\Psi\) commutes with every constant
matrix unit, hence with \(\End(V)\otimes\OO_B\).

For clarity, commutant recovery itself is canonical and does not
use these chosen matrix units as extra data.  Tensoring a bundle
map \(a\) with \(1_V\) gives the forward map in
\eqref{eq:commutant-descent-v147}.  Commutation with matrix units
shows locally that it is bijective, with inverse reading any diagonal
block.  This proves that it is an isomorphism of sheaves, independent
of the checking basis.  Its inverse applied to the global section
\(\Psi\) produces a global section \(a\), not projective local
sections requiring arbitrary scalar lifts.  Applying the same
operation to \(\Psi^{-1}\) proves that \(a\) is invertible.
Uniqueness follows from injectivity of \(a\mapsto1_V\otimes a\).

Projective-factor data alone would permit replacing a bundle by a
line twist.  The additional datum used here is the actual linear map
of the tensor bundles.  It is this datum, recovered from the dual of
the first nilradical quotient, that makes the commutant argument
applicable.
\end{proof}

\subsection{A general finite-order coefficient inverse}
Fix \(n\ge2\), \(q\ge1\), and integers
\(0\le r_\lambda\le m_\lambda=\dim\mathbb S_\lambda\C^n\)
for partitions \(\lambda\) of \(q\) of length at most \(n\).
Assume that for at least one partition
\begin{equation}\label{eq:strict-support-v147}
                         0<r_\lambda<m_\lambda.
\end{equation}
Let \(B\) be a smooth connected projective complex variety,
\(V\) a constant \(n\)-dimensional space, and \(\mathcal U\)
a rank-\(n\) bundle.  For each \(\lambda\), choose a subbundle
with locally free quotient
\[
 \mathcal L_\lambda\subset\mathbb S_\lambda V^*\otimes\OO_B,
                       \qquad\rank\mathcal L_\lambda=r_\lambda.
\]
Zero and full subbundles are permitted.  In the matrix bundle
\(\mathcal E=\Hom(\mathcal U,V\otimes\OO_B)\), the
multiplicity-one Cauchy decomposition supplies the subbundle
\begin{equation}\label{eq:general-moving-coeff-v147}
 \mathcal J_q=\bigoplus_\lambda
       \mathcal L_\lambda\otimes\mathbb S_\lambda\mathcal U
             \ \subset\ \Sym^q\mathcal E^*.
\end{equation}
Let \(\mathcal D_n\subset\Sym^n\mathcal E^*\) be the
determinant line, put \(D=n+q\), and define
\begin{equation}\label{eq:coefficient-thickening-v147}
 \mathcal K=\mathcal D_n\mathcal J_q,
 \quad \mathcal A=\Sym\mathcal E^*/\bigl((\mathcal K)
                                  +\mathfrak a^{D+1}\bigr),
 \quad Y=\Spec_B\mathcal A,
\end{equation}
where \(\mathfrak a\) is the positive homogeneous ideal.  The
notation \(D\) in this subsection is an integer, not a failure
locus.  The scheme \(Y\) is finite over \(B\), and need not be
zero-dimensional over \(\C\).

\begin{theorem}[Unmarked descent for moving coefficient spaces]
\label{thm:general-moving-coefficients-v147}
For fixed \(n,q,(r_\lambda)\) satisfying
\eqref{eq:strict-support-v147}, the abstract complex scheme \(Y\)
recovers its coefficient data as follows.  Two such schemes \(Y_1,Y_2\)
are isomorphic if and only if there are an isomorphism
\(\sigma:B_1\to B_2\), a constant linear isomorphism
\(g:V_1\to V_2\), and an actual bundle isomorphism
\(h:\mathcal U_1\to\sigma^*\mathcal U_2\) such that, for every
\(\lambda\),
\begin{equation}\label{eq:coefficient-equivalence-v147}
 \mathcal L_{1,\lambda}=
       (\mathbb S_\lambda g)^*\sigma^*\mathcal L_{2,\lambda}.
\end{equation}
Every unmarked isomorphism induces a normal-bundle map
\(T\mapsto gTh^{-1}\); for that map the pair \((g,h)\) is
unique up to a common nonzero constant scalar.  Conversely every
such triple lifts homogeneously to an isomorphism.  The scheme is
finite locally free over its reduced base of rank
\[
             \binom{n^2+D}{D}-\sum_\lambda r_\lambda m_\lambda.
\]
No nontriviality hypothesis on \(\Pj_B(\mathcal U^*)\) is required.
\end{theorem}
\begin{proof}
\emph{Recovering the intrinsic linear data.}
Multiplication by a determinant polynomial is an injective linear
map of polynomial spaces.  In local bundle frames its matrix is
constant over \(\C\) and its cokernel is free.  The Cauchy
summands and the \(\mathcal L_\lambda\) are subbundles, so
\(\mathcal K\) is a subbundle of rank
\(\sum r_\lambda m_\lambda\).  The graded pieces of
\(\mathcal A\) below \(D\) are the full symmetric powers;
its top piece is \(\Sym^D\mathcal E^*/\mathcal K\).
This proves local freeness and the rank formula.

The nilradical \(\mathfrak n\) of \(\OO_Y\) is its positive
ideal, since \(B\) is reduced.  Thus \(Y_{\mathrm{red}}=B\) and
\(\mathfrak n/\mathfrak n^2=\mathcal E^*\) intrinsically.
The multiplication kernel
\[
 \ker\bigl[\Sym^D(\mathfrak n/\mathfrak n^2)
                                    \longrightarrow\mathfrak n^D\bigr]
\]
is precisely \(\mathcal K\), by the proof of
Proposition~\ref{prop:first-relation-orbit-v147} in local frames.
An abstract isomorphism consequently gives the base isomorphism
\(\sigma\) and an actual linear normal-bundle isomorphism carrying
\(\mathcal K_1\) to \(\mathcal K_2\) contravariantly.  The
normal map is linear over the reduced base even when the original
scheme map does not preserve the chosen graph projection.  Indeed
the action on \(\mathfrak n/\mathfrak n^2\) is linear over
\(\OO_Y/\mathfrak n\).

\emph{Recovering and orienting the factors.}
Generate the homogeneous ideal \((\mathcal K)\) in the recovered
symmetric algebra; do not confuse its cone with the finite
truncation.  Its radical is the determinant ideal.  To verify this,
at an invertible matrix \(T\), the map
\(\mathbb S_\lambda T\) is invertible.  For a partition with
\(r_\lambda>0\), a nonzero covector in \(\mathcal L_\lambda\)
has a nonzero matrix coefficient there.  Thus the residual
coefficients cannot all vanish off the determinant.  All generators
vanish on the determinant.  The relative determinant hypersurface
is reduced over the smooth base; the equality of radicals follows
on affine charts from this equality of closed points.  Successive
reduced singular loci and Lemma~\ref{lem:rank-one-fano-scheme}
recover its two unordered tensor rulings with their relative scheme
structures.

Cancellation of the determinant line is intrinsic and recovers the
ideal generated by \(\mathcal J_q\), and hence its degree-\(q\)
piece.  In a Cauchy summand the two contraction supports of
\(\mathcal L_\lambda\otimes\mathbb S_\lambda\mathcal U\)
have ranks \((r_\lambda,m_\lambda)\).  Separate factor changes
preserve these ranks.  By Lemma~\ref{lem:linear-preserver-v147},
an exchange of the two matrix factors interchanges the ranks in
each summand.  At a partition satisfying
\eqref{eq:strict-support-v147}, the result
\((m_\lambda,r_\lambda)\) cannot equal
\((r_\lambda,m_\lambda)\) for the other datum.  Thus no fibre
permits an exchange.  Twisting local factor lifts by lines changes
no support rank.  The \(V\)-ruling is thereby oriented intrinsically,
including when both projective ruling bundles are trivial.

\emph{Descending the actual map.}
By Lemma~\ref{lem:oriented-segre-descent}, the induced morphism
\(B_1\to\operatorname{Isom}(\Pj(V_1),\Pj(V_2))\) is constant:
the target is affine and \(H^0(B_1,\OO_{B_1})=\C\).  Choose a
constant linear lift \(g\).  Remove it from the actual normal map.
Lemma~\ref{lem:actual-factor-descent-v147}, applied to
\(\mathcal A_1=\mathcal U_1^*\) and
\(\mathcal A_2=(\sigma^*\mathcal U_2)^*\), supplies the unique
global right factor.  Dualizing its inverse gives the actual
isomorphism \(h\).  This step uses more than an isomorphism of
projective bundles; it is exactly what excludes an undetected line
twist.  Replacing \(g\) by \(cg\) replaces \(h\) by \(ch\).
These are all ambiguities because the left lift is constant.

\emph{Recovering every coefficient subbundle.}
The oriented Cauchy decomposition is now functorial for the recovered
left-right map.  Project the residual degree-\(q\) module onto each
summand and contract against the full right dual factor.  The image
is precisely \(\mathcal L_\lambda\), as a subbundle of the
constant left module, not just a collection of fibrewise projective
subspaces.  The right map acts invertibly on the full Schur factor
and hence does not change this image.  Pullback by the normal map
therefore gives \eqref{eq:coefficient-equivalence-v147} for the same
constant \(g\) in all summands.  Scalar choices in the determinant
line do not change these subbundles.  This proves necessity.

Conversely \(T\mapsto gTh^{-1}\) carries the determinant line,
all coefficient subbundles, and every power of the positive ideal
to the corresponding objects.  It induces the required homogeneous
scheme isomorphism.  The stated uniqueness refers to the induced
normal map; it does not eliminate the higher-order automorphisms of
that scheme.
\end{proof}

\subsection{Identical fibres and graded bundles do not determine a family}
The global conclusion has content even when neither fibre isomorphism
classes nor the underlying graded vector bundles vary.  The following
example uses the general coefficient theorem, not a further spectral
readout of a reconstructed pencil.

\begin{example}[Two isotrivial thickenings separated by coefficient descent]
\label{ex:isotrivial-covers-v147}
Take \(B=\Pj^1\), \(V=\C^3\), \(\mathcal U=\OO_B^3\),
\(q=1\), and \(r_{(1)}=1\).  The two everywhere nonzero columns
\[
 (s^3,t^3,0)^{\mathsf t},\qquad
 (s^3+st^2,t^3,0)^{\mathsf t}
\]
define subbundles \(\mathcal L_0,\mathcal L_1\subset V^*\otimes\OO_B\),
both isomorphic to \(\OO_B(-3)\).  Let \(Y_0,Y_1\) be their
order-\(4\) thickenings \eqref{eq:coefficient-thickening-v147}.
They are finite flat of rank \(712\) over \(B\) and Zariski
locally trivial with the same Artin algebra as fibre.  Their
nilradical graded vector bundles are isomorphic in every degree,
although \(Y_0\) and \(Y_1\) are not isomorphic as abstract
complex schemes.
\end{example}
\begin{proof}
Neither column vanishes, since its second entry is nonzero when
\(t\ne0\), and its first is nonzero at \(t=0\).  The left
coefficient map to its image line \(\Pj^1\subset\Pj(V^*)\)
is, in the coordinate \(z=s/t\), respectively
\[
                        f_0(z)=z^3,\qquad f_1(z)=z^3+z.
\]
A local frame completing a generator of \(\mathcal L_i\) to a
basis of \(V^*\otimes\OO\) changes the coefficient line to a
fixed line.  This gives Zariski local product presentations of both
thickenings.  Alternatively, equality of the fibre isomorphism
classes follows from transitivity of \(\operatorname{GL}(V)\)
on nonzero coefficient lines.  The rank formula is
\(\binom{13}{4}-3=712\).

Let \(E=V^*\otimes\C^3\), of dimension nine.  For \(j<4\),
the degree-\(j\) nilradical quotient is \(\Sym^j E\otimes\OO_B\).
Multiplication by the fixed determinant embeds \(E\) as a fixed
nine-dimensional subspace of \(\Sym^4 E\).  Choosing a constant
linear complement gives, for the top quotient,
\[
 \mathfrak n_i^4\simeq
 \OO_B^{486}\oplus
           \bigl((V^*\otimes\OO_B)/\mathcal L_i\bigr)^{\oplus3}.
\]
Each column has only its first two entries nonzero, and these two
entries have no common zero.  The rank-one quotient of \(\OO_B^2\)
by \(\OO_B(-3)\) is consequently \(\OO_B(3)\), by determinants.
Thus the displayed top bundle is
\(\OO_B^{489}\oplus\OO_B(3)^{\oplus3}\) for both \(i\).
The assertions about every graded vector bundle follow.  They are
not assertions that the graded \emph{algebra} structures agree.

If an abstract isomorphism \(Y_0\simeq Y_1\) existed,
Theorem~\ref{thm:general-moving-coefficients-v147} would identify
their coefficient maps by one automorphism of \(B\) and one
constant projective transformation of \(V^*\).  Both maps have
as image the same projective line, so the latter transformation
would restrict to a projective automorphism of that line.  Hence
\(f_0\) and \(f_1\) would be equivalent by independent projective
changes in source and target.  This is impossible.  The map \(f_0\)
has two ramification points, each of index three.  The map \(f_1\)
has the two simple zeros of \(3z^2+1\), each of index two, and
infinity, of index three.  Ramification indices are invariant under
those projective changes.  Thus the unmarked algebra scheme
distinguishes these two globally different coefficient maps even
though its fibre algebras, conormal bundle, abstract first-relation
bundle, and all graded vector bundles agree.
\end{proof}

\subsection{Pencils and the category of their automorphisms}
For pencils, take \(q=2p\), \(p=N-2\), and
\(F(H)=\bigwedge^{N-2}\Sym^2 H\).  By
Lemma~\ref{lem:pencil-irreducibility} this is a single Schur module
of dimension \(M=\binom N2\), and
Lemma~\ref{lem:moving-coefficients-v146} identifies its coefficient
subbundle with
\(\mathcal L_{\mathcal R}=\bigwedge^p(\Sym^2V/\mathcal R)^*\).
Its rank is one, so the strict-support hypothesis holds for every
pencil.  Exterior duality recovers the subbundle \(\mathcal R\)
from that line.  Theorem~\ref{thm:general-moving-coefficients-v147}
therefore proves the unrestricted moving-pencil inverse, including
its actual source-bundle conclusion.  This is a simultaneous
reconstruction theorem for a map into a coefficient Grassmannian,
not just a classification of its individual fibres.

\begin{corollary}[Algebraic automorphism sequence for a pencil]
\label{cor:pencil-group-scheme-v147}
Let \(\mathfrak A_R^{[d]}\) be the unmarked local algebra of
Theorem~\ref{thm:artin-local-inverse-v146}, and let
\(G_R=\{g\in\operatorname{GL}(V):(\Sym^2g)R=R\}\).
There is a split exact sequence of complex affine algebraic groups
\begin{equation}\label{eq:pencil-algebraic-group-v147}
 1\longrightarrow U_R\longrightarrow
 \operatorname{Aut}(\mathfrak A_R^{[d]})
 \longrightarrow (G_R\times\operatorname{GL}(U))/\mathbb G_m
 \longrightarrow1,
\end{equation}
where \(\mathbb G_m\) is the diagonal scalar subgroup.  Its
intrinsic quotient is the closed stabilizer of \(K_R\) in the
general linear group of the cotangent space.  In that form the
sequence represents automorphisms after extension to every
commutative complex algebra.  Its splitting is homogeneous after
choosing generators, and \(U_R\) has the intrinsic description and
dimension in Proposition~\ref{prop:first-relation-group-v147}.
\end{corollary}
\begin{proof}
Proposition~\ref{prop:first-relation-group-v147} gives the represented
sequence with quotient \(H_{K_R}\).  The coefficient inverse and
Lemma~\ref{lem:linear-preserver-v147} identify its complex points
with the image of the left-right action; transposition is excluded
by the support ranks \((1,M)\).  Use the contragredient of the
normal action to obtain a homomorphism on the cotangent space.
Its kernel is exactly the diagonal scalar subgroup: a tensor product
acting identically has both factors scalar, with reciprocal actions
on the dual factors.  The image is a closed algebraic subgroup and
all the groups in question are smooth in characteristic zero
\cite[Lemma~39.8.2]{StacksGroups}.  Thus the stabilizer and this
image agree as closed subgroup schemes, not only on complex points.
The induced quotient identification is an isomorphism of algebraic
groups: after division by the kernel it is a bijective separable
homomorphism between smooth groups.  The splitting and kernel are
those already constructed by substitutions.
\end{proof}

For a positive-dimensional projective reduction,
Theorem~\ref{thm:automorphism-kernel-v146} is used as an exact
sequence of groups of complex scheme automorphisms, with its stated
filtered kernel.  We do not identify that sequence with a finite-type
affine group scheme or a moduli-stack equivalence.  Likewise, the
marked construction over nonreduced parameter bases later in the
paper retains its specified relative socle: it does not replace that
ideal by the absolute nilradical.  These distinctions preserve the
full unmarked geometric assertion while specifying its category.

% BEGIN PRESERVED PART 13a-criterion-v145.tex
\section{An intrinsic coefficient principle}
\label{sec:structural-coefficient-principle}

Extracting coefficient data from an abstract scheme and reconstructing
a defining relation space from those data are distinct operations.
The criterion below isolates the first.  Its hypotheses will be
verified directly for the multiplication schemes and then on curves.

\subsection{The structural criterion}

Let \(X\) be a finite-type complex scheme with a closed subvariety
\(B\) specified by a construction preserved by every abstract scheme
isomorphism.  Write \(\mathfrak b\subset\OO_X\) for its ideal and
\[
 \mathcal G=\bigoplus_{r\ge0}\mathfrak b^r/\mathfrak b^{r+1},
 \qquad \mathcal A=\Sym_{\OO_B}(\mathcal G_1),
 \qquad \mathcal I=\ker(\mathcal A\twoheadrightarrow\mathcal G).
\]
Assume the following conditions, for fixed integers \(n\ge2,m\ge1\).

\smallskip
\noindent\textup{(i)} \(B\) is connected and projective,
\(H^0(B,\OO_B)=\C\), and \(\mathcal G_1\) is locally free.
Its dual admits a presentation \(V\otimes\mathcal U^*\), where
\(\dim V=\rank\mathcal U=n\).

\smallskip
\noindent\textup{(ii)} The reduction of \(\Spec_B\mathcal G\)
in its degree-one ambient bundle is the generic determinant cone.
The bundle \(\Pj_B(\mathcal U^*)\) is not projectively trivial.
Thus precisely one of the two ruling parameter bundles of its rank-one
locus is projectively trivial.

\smallskip
\noindent\textup{(iii)} If \(\mathcal D\subset\mathcal A\)
is the ideal of that reduced cone and \(\mathcal L=\mathcal D_n\),
then \(\mathcal I=\mathcal D\mathcal J\), where
\(\mathcal J\) is generated in degree \(m\).  Under the Cauchy
decomposition in any local frame of \(\mathcal U\),
\begin{equation}\label{eq:structural-coefficient-hypothesis}
 \mathcal J_m=
 \bigoplus_{\lambda\in\Lambda}
          L_\lambda\otimes\mathbb S_\lambda\mathcal U
 \ \subset\
 \Sym^m(V^*\otimes\mathcal U),
 \qquad L_\lambda\subset\mathbb S_\lambda V^*,
\end{equation}
where \(\Lambda\) is a specified collection of partitions of \(m\)
of length at most \(n\), and the left subspaces are constant over
\(B\).  Zero subspaces are permitted.  The Cauchy inclusions, not
just their representation labels, are the coefficient realizations
\(\beta\otimes\xi\mapsto[T\mapsto\beta(\mathbb S_\lambda(T)\xi)]\).

These conditions are hypotheses on the normal cone and its residual
coefficient module.  They do not assume that an abstract isomorphism
extends to an ambient space, preserves a chosen frame, or already acts
on all the left subspaces through a common projective transformation.

\begin{theorem}[Intrinsic coefficient reconstruction]
\label{thm:structural-coefficient-reconstruction}
Suppose \(X,B\) satisfy \textup{(i)--(iii)}.  The abstract scheme
\(X\) determines the collection \((L_\lambda)_{\lambda\in\Lambda}\)
up to one common projective coordinate transformation on \(V\).
More precisely, an isomorphism between two such schemes induces a
constant projective isomorphism \([g]:\Pj(V)\to\Pj(V')\) such that
\[
                L_\lambda=\mathbb S_\lambda(g)^*L'_\lambda
                       \quad(\lambda\in\Lambda).
\]
The equality means equality of subspaces; scalar changes of a linear
lift of \([g]\) do not affect it.
\end{theorem}
\begin{proof}
An isomorphism transports \(B\), every power of \(\mathfrak b\),
\(\mathcal G_1\), and the surjection
\(\mathcal A\twoheadrightarrow\mathcal G\).  In particular it
transports \(\mathcal I\) in the intrinsic degree-one ambient
bundle.  Iterated reduced singular loci of the determinant cone give
its rank-one locus.  Lemma~\ref{lem:rank-one-fano-scheme} reconstructs
its two ruling parameter schemes, including their scheme structure.
Condition \textup{(ii)} distinguishes the trivial ruling intrinsically.
Lemma~\ref{lem:oriented-segre-descent} then gives one constant
projective map \([g]\).  After choosing a constant linear lift and
local frames on both bases, the ambient bundle map has the form
\(T\mapsto gTh^{-1}\).  These local choices are made only after
constructing the graded algebra, its degree-one part, and the ruling.

Locally \(\mathcal D=(d)\) with \(d=\det T\).  Its leading
coefficient is a unit, so it is a nonzerodivisor over every coefficient
ring.  Hence \((\mathcal I:\mathcal D)=\mathcal J\), and
\[
 \mathcal J_m\simeq\mathcal L^*\otimes\mathcal I_{n+m}.
\]
This recovers the residual module without choosing a scalar determinant
equation.  Under \(T\mapsto gTh^{-1}\), the Cauchy coefficient
map transports a tensor by
\(\mathbb S_\lambda(g)^*\otimes\mathbb S_\lambda(h^{-1})\).
Projection to the specified Cauchy summand, followed by contraction
against its whole right dual, therefore transports exactly
\(L'_\lambda\) to \(L_\lambda\).  The right map is invertible
and has no effect on that left image.  Both determinant-line transition
functions and changes of a lift of \([g]\) act by scalar units.
Since every partition has size \(m\), this ambiguity is common
in degree and disappears on subspaces.  This proves the assertion.
\end{proof}


% BEGIN PRESERVED PART 18-pencil-proof-v145.tex
\section{Quadratic relation spaces and intrinsic reconstruction of pencils}
\label{sec:natural-pencils}

The multiplication construction itself supplies a wider natural class
for the coefficient criterion.  In this section the number of quadratic
relations is allowed to vary independently of the embedding dimension.
For two relations the resulting theorem recovers every pencil of
quadrics, including singular pencils.  No extra block is added to the
multiplication map.

\subsection{Automatic coefficient extraction for quadratic multiplication}

Let \(\dim V=n\ge3\), \(W=\Sym^2V\), \(N=n(n+1)/2\), and
\(R\in\Gr(r,W)\), where \(1\le r<N-n\).  Put \(p=N-r>n\),
\(S=W/R\), and
\[
 A_R=\C\oplus V\oplus S,\qquad
 X_R=\Gr(n,V\oplus S).
\]
The only nonzero product in the maximal ideal is the quotient map
\(\gamma:\Sym^2V\to S\).  Write \(\mathcal K\) for the
tautological bundle on \(X_R\), and use the same definition
\eqref{eq:failure-intro} of \(\widehat D_R\), with \(A_R\) in
place of \(B_R\).  Thus \(\widehat D_R\) is the failure scheme
of generating \(A_R\) by a varying \(n\)-plane and the unit.
The notation in this section agrees with the web construction when
\((n,r)=(4,4)\).

Let \(\Lambda_{n,p}\) denote the partitions occurring in the
classical multiplicity-free decomposition
\begin{equation}\label{eq:general-exterior-decomposition}
 F(V):=\bigwedge^p\Sym^2V
       =\bigoplus_{\lambda\in\Lambda_{n,p}}\mathbb S_\lambda V.
\end{equation}
Every partition has size \(2p\); terms of length greater than \(n\)
are omitted.  Choose a nonzero volume form on \(S\) and let
\(\beta_R\in F(V)^*\) be its pullback by \(\bigwedge^p\gamma\).
Its component in \(\mathbb S_\lambda V^*\) is denoted by
\(\beta_{R,\lambda}\).  Only the lines of the nonzero components
will be used, so the volume choice is immaterial.

\begin{theorem}[Automatic intrinsic coefficient extraction]
\label{thm:automatic-quadratic-coefficients}
For every \(R\in\Gr(r,\Sym^2V)\) in the stated range, the
abstract scheme \(\widehat D_R\) determines
\[
       \bigl(\C\beta_{R,\lambda}\bigr)_{\lambda\in\Lambda_{n,p}}
\]
up to one common projective transformation of \(V\), with zero
components recorded as zero subspaces.  More precisely, the intrinsic
construction recovers the base \(B=\Gr(n,S)\), its degree-one
normal-cone bundle, and the residual ideal generated in degree \(2p\)
with coefficient module
\begin{equation}\label{eq:automatic-residual-module}
 \mathcal J_{2p}=
 \bigoplus_{\lambda\in\Lambda_{n,p}}
       \C\beta_{R,\lambda}\otimes\mathbb S_\lambda\mathcal U
 \ \subset\ \Sym^{2p}(V^*\otimes\mathcal U),
\end{equation}
where \(\mathcal U\) is the tautological rank-\(n\) bundle on
\(B\).  Consequently all geometric hypotheses of
Theorem~\ref{thm:structural-coefficient-reconstruction} hold for this
entire Grassmannian family, not just for smooth webs.
\end{theorem}
\begin{proof}
On a frame chart of \(\mathcal K\), write its inclusion as
\(\binom M{A_0}\), with \(M\) an \(n\times n\) matrix.  Since
\(A_R\) is cube-zero, columns of degree greater than two vanish.
After retaining the scalar, linear, and quadratic columns, multiplication
has the matrix
\begin{equation}\label{eq:general-multiplication-matrix}
 \begin{pmatrix}
  1&0&0\\
  0&M&0\\
  0&A_0&\gamma\Sym^2M
 \end{pmatrix}.
\end{equation}
A nonzero maximal minor must use the scalar column and all \(n\)
linear columns: these are the only columns meeting the first \(1+n\)
rows.  Expanding along those rows gives the identity of ideals
\begin{equation}\label{eq:general-fitting-factorization}
 I_{\widehat D_R}=(\det M)I_p(\gamma\Sym^2M).
\end{equation}
The same expansion works for every \(m\ge2\).  Where \(M\) is
invertible, the quadratic map is surjective.  Hence the reduction of
\(\widehat D_R\) is exactly the Schubert divisor where \(M\)
is singular, independent of \(R\) after identifying the vector
spaces \(S\).

Successive reduced singular loci of that divisor are the projection
rank loci.  Their last nonempty member is \(B=\Gr(n,S)\).
For clarity, the usual determinantal assertion here is local: at a
projection-rank-\(k\) point, an invertible \(k\)-minor splits
away, leaving a generic square \((n-k)\)-matrix and smooth
coordinates.  Every remaining entry can vary independently by
moving the kernel of the projection in \(V\) modulo its image.
Thus its reduced singular locus is the next rank locus.  Since
\(p>n\), the final base is nonempty, connected and projective,
with \(H^0(B,\OO_B)=\C\).

Near \(B\), a Grassmannian chart consists of a varying plane in
\(S\) and its graph \(T:\mathcal U\to V\).  Formula
\eqref{eq:general-fitting-factorization} is independent of the other
graph coordinates and homogeneous in \(T\).  It therefore gives
its actual normal-cone ideal, not merely an initial reduced support:
\begin{equation}\label{eq:general-normal-cone}
 C_B\widehat D_R\subset\Tot\Hom(\mathcal U,V),\qquad
 \mathcal I=(\det T)\mathcal J,\quad
 \mathcal J=I_p(\gamma\Sym^2T).
\end{equation}
There are no degree-one equations.  The ambient bundle in
\eqref{eq:general-normal-cone} is consequently the dual of the
intrinsically recovered degree-one conormal.  Its reduced ideal is
the determinant ideal \(\mathcal D\).  The determinant polynomial
is monic for a suitable monomial order, so cancellation gives
\((\mathcal I:\mathcal D)=\mathcal J\).  This is an intrinsic
ideal quotient; choosing a generator of the determinant line is not
part of the recovered datum.

The rank-one locus of the reduced cone determines the two ruling
parameter schemes by Lemma~\ref{lem:rank-one-fano-scheme}.  They are
\(\Pj(V)\times B\) and \(\Pj_B(\mathcal U^*)\).
Exactly the first is projectively trivial.  Indeed a Schubert line in
\(\Gr(n,S)\) restricts \(\mathcal U^*\) to
\(\OO(1)\oplus\OO^{n-1}\).  A trivial projective bundle on
that line would make this bundle a line-bundle twist of
\(\OO^n\); its determinant would then have degree divisible by
\(n\), rather than degree one.  Thus the trivial ruling is
intrinsically oriented.  Lemma~\ref{lem:oriented-segre-descent}
applies: an abstract isomorphism gives a morphism \(B\to\PGL(V)\),
constant because \(B\) is connected projective and \(\PGL(V)\)
is affine.  Local linear lifts and right changes of frame give
\(T\mapsto gTh^{-1}\) for this one fixed \([g]\).

It remains to prove \eqref{eq:automatic-residual-module}, rather
than assume it as a condition on the example.  The minors are the
image, for fixed \(\beta_R\), of the universal coefficient map
\begin{equation}\label{eq:general-coefficient-map}
 F(V)^*\otimes F(U)\longrightarrow
 \Sym^{2p}(V^*\otimes U),\qquad
 \beta\otimes\xi\longmapsto[T\mapsto\beta(F(T)\xi)].
\end{equation}
In characteristic zero, the polynomial functor \(F\) has the
multiplicity-free decomposition \eqref{eq:general-exterior-decomposition};
see \cite[Theorem~3.5]{ChengWang}.  The Cauchy decomposition of the
target contains each diagonal
\(\mathbb S_\lambda V^*\otimes\mathbb S_\lambda U\)
once, and no off-diagonal pair.  Equivariance kills every
off-diagonal pair.  On each diagonal pair the coefficient map is
nonzero: identify \(U\) with \(V\), evaluate at \(T=1\), and
choose \(\beta,\xi\) in that summand with
\(\beta(\xi)\ne0\).  Schur's lemma then identifies its image
with the specified Cauchy summand.  This is the actual matrix-coefficient
inclusion, not merely a dimension count.  Keeping \(\beta_R\)
fixed proves \eqref{eq:automatic-residual-module}.  All generators
have degree \(2p\), and at least one minor is nonzero since
\(\gamma\) is surjective.

Under \(T\mapsto gTh^{-1}\), the right map acts within the full
factor \(\mathbb S_\lambda U\).  Contracting against its dual
therefore leaves exactly the left coefficient line transformed by
\(\mathbb S_\lambda(g)^*\).  Determinant-line scalars and local
lift ambiguities do not change these lines.  This proves the theorem.
\end{proof}

The theorem does not assert that separately projectivized Schur
components always determine their relative scalars.  For webs this
is the recombination problem solved in the separate research archive.  For pencils that issue
disappears altogether.

\subsection{Every pencil, without a regularity hypothesis}

\begin{lemma}[The exterior square of the quadratic representation]
\label{lem:pencil-irreducibility}
For \(n\ge2\), \(\bigwedge^2\Sym^2V=\mathbb S_{(3,1)}V\)
is irreducible.  Consequently
\(\bigwedge^{N-2}\Sym^2V\) is irreducible as well.
\end{lemma}
\begin{proof}
The vector \(x_1^2\wedge x_1x_2\) is a highest-weight vector of
weight \((3,1)\).  By complete reducibility it generates a copy of
\(\mathbb S_{(3,1)}V\).  The hook formula gives
\[
 \dim\mathbb S_{(3,1)}V
   =\frac{n(n-1)(n+1)(n+2)}8
   =\binom{N}{2}.
\]
This exhausts the exterior square.  The last assertion follows by
exterior duality.  These are classical representation identities.
\end{proof}

\begin{theorem}[Intrinsic reconstruction of all quadratic pencils]
\label{thm:natural-pencil-reconstruction}
Let \(n\ge3\) and let \(R,R'\) be two-dimensional subspaces
of \(\Sym^2V\).  For the natural multiplication-failure schemes
of their cube-zero algebras of Hilbert function \((1,n,N-2)\),
\begin{equation}\label{eq:natural-pencil-torelli}
 \widehat D_R\simeq\widehat D_{R'}
 \quad\Longleftrightarrow\quad
 R'=gR\text{ for some }g\in\PGL(V).
\end{equation}
No pencil, base locus, or discriminant is required to be smooth or
regular.  The isomorphism on the left is abstract and carries no
supplied ambient embedding, base, or tensor-factor marking.  All the
reduced schemes are isomorphic to the same Schubert divisor.
\end{theorem}
\begin{proof}
Here \(p=N-2>n\).  By
Theorem~\ref{thm:automatic-quadratic-coefficients} and
Lemma~\ref{lem:pencil-irreducibility}, the normal cone has only one
left coefficient line to recover: the entire line \(\C\beta_R\).
The canonical duality
\begin{equation}\label{eq:pencil-exterior-duality}
 \bigwedge^{N-2}W^*\simeq
       \bigwedge^2W\otimes(\det W)^*
       =\bigwedge^2W\otimes(\det V)^{-(n+1)}
\end{equation}
sends that line to the Pluecker line of \(R\), with the same
projective transformation \(g\).  Indeed the pullback of a
quotient volume is the annihilator-volume dual of the kernel
\(R\).  The common determinant character is a scalar on
projective space.  A nonzero decomposable two-vector recovers its
plane as \(\{w\in W:w\wedge w_R=0\}\).  Thus an abstract
scheme isomorphism implies \(R'=gR\).

Conversely such a linear coordinate change induces the quotient
isomorphism \(S_R\to S_{R'}\), the algebra isomorphism
\(A_R\to A_{R'}\), and the corresponding Grassmannian isomorphism.
It intertwines multiplication and hence preserves its zeroth Fitting
ideal.  This proves the reverse implication.  The assertion about
reductions was proved in \eqref{eq:general-fitting-factorization}.
\end{proof}

\subsection{Discriminants, elementary divisors, and hyperelliptic moduli}

A regular pencil means one containing an invertible symmetric matrix;
it need not have a reduced discriminant.  The classical invariant of
such a pencil consists of its discriminant points together with their
Jordan-block partitions, usually called its Weierstrass--Segre data.
We use the root-labelled partitions, not a pairing of an unlabelled
list with an independently unlabelled divisor.  Their classification
is classical; compare \cite[Theorem~1.1, Corollary~2.1 and \S2]{FevolaMandelshtamSturmfels}.
The following consequence identifies precisely what the intrinsic
failure scheme remembers beyond the discriminant.

\begin{corollary}[Intrinsic recovery of elementary divisors]
\label{cor:intrinsic-segre-data}
For a regular pencil in any dimension \(n\ge3\), the abstract
scheme \(\widehat D_R\) determines all root-labelled elementary
divisors of a matrix pencil, up to projective reparametrization of its
parameter line.  Conversely these data determine its abstract failure
scheme.  In particular the invariant distinguishes different Jordan
partitions over a fixed multiple discriminant root.
\end{corollary}
\begin{proof}
The first implication follows by reconstructing \(R\), then
choosing an invertible member \(A\) and a complementary member
\(B\).  For completeness, the relevant classical sufficiency can
be seen directly.  Put \(C=A^{-1}B\); it is self-adjoint for
\(A\).  If \(hC=C'h\) identifies the Jordan forms for another
pair \((A',B')\), set
\[
 H=A^{-1}h^tA'h.
\]
Then \(H\) is invertible, \(A\)-self-adjoint, and commutes with
\(C\).  There is a polynomial \(k\) in \(H\) with
\(k^2=H\): on each nonzero eigenvalue choose a square root and
truncate its Taylor series to the size of the largest Jordan block,
then combine by the Chinese remainder theorem.  Thus \(k\) is
invertible, self-adjoint, and commutes with \(C\).
The map \(hk^{-1}\) is an isometry from \(A\) to \(A'\)
and intertwines \(C,C'\), so it simultaneously identifies
\(A,B\) with \(A',B'\).

Changing the basis of the pencil replaces \(C\) by a fractional
linear expression with invertible denominator.  At each eigenvalue
its derivative is nonzero; hence each Jordan-block size is preserved
while its root moves by the same projective transformation.  A
parameter change can avoid infinity for both finite spectra.
Root-labelled elementary divisors therefore suffice, and
\eqref{eq:natural-pencil-torelli} proves the converse for the
failure schemes.  This argument recalls the classical classification;
the new input is recovery from an unmarked abstract failure scheme.
\end{proof}

\begin{example}[Identical discriminant, distinct intrinsic failure schemes]
\label{ex:same-discriminant-distinct-failure}
Choose distinct nonzero \(\mu_3,\ldots,\mu_n\), and put
\[
 S_2=\begin{pmatrix}0&1\\1&0\end{pmatrix},\qquad
 J_2=\begin{pmatrix}0&1\\0&0\end{pmatrix}.
\]
Consider the symmetric pairs
\begin{align*}
 A_{\mathrm{s}}&=I_n,&
 B_{\mathrm{s}}&=\diag(0,0,\mu_3,\ldots,\mu_n),\\
 A_{\mathrm{j}}&=S_2\oplus I_{n-2},&
 B_{\mathrm{j}}&=S_2J_2\oplus\diag(\mu_3,\ldots,\mu_n).
\end{align*}
Their determinant binary forms agree up to a nonzero scalar:
\begin{equation}\label{eq:same-discriminant-pencils}
 \det(sA_{\mathrm{s}}+tB_{\mathrm{s}})
 =-\det(sA_{\mathrm{j}}+tB_{\mathrm{j}})
 =s^2\prod_{i=3}^n(s+\mu_i t).
\end{equation}
At the unique double root the ranks are respectively \(n-2\)
and \(n-1\), so no simultaneous projective equivalence exists.
The two root partitions are \((1,1)\) and \((2)\).
Their abstract failure schemes are therefore nonisomorphic although
both their reduced failure schemes and their discriminant divisors,
including multiplicities, agree.  This realizes inside the same
multiplication construction the classical loss of elementary-divisor
information under passage to the determinant.
\end{example}

\begin{corollary}[Hyperelliptic isomorphism classes]
\label{cor:hyperelliptic-failure-invariant}
Fix \(g\ge2\) and \(n=2g+2\).  On the locus of pencils with
simple discriminant, \(\widehat D_R\) gives an injective invariant
of the isomorphism classes of smooth hyperelliptic curves of genus
\(g\).  All these failure schemes have the same reduction; their
nonreduced structures distinguish a \((2g-1)\)-dimensional family
of curve classes.  The pencil invariant also distinguishes the two
regular boundary pencils in
Example~\ref{ex:same-discriminant-distinct-failure}, where their
singular discriminant double covers alone do not.
\end{corollary}
\begin{proof}
The classical curve associated with a simple-discriminant pencil is
the double cover of its parameter line branched at
\(\det(sA+tB)=0\); see \cite[\S4.2, equation (4.3)]{SmithQuadrics}.
Choose an invertible member.  Distinct eigenvalues of \(A^{-1}B\)
have mutually \(A\)-orthogonal eigenlines, each nonisotropic, so
simultaneous diagonalization identifies the pencil with
\(\langle\sum x_i^2,\sum\lambda_i x_i^2\rangle\).
The unordered branch divisor modulo \(\PGL_2\) therefore
determines the pencil, and every reduced divisor of degree \(2g+2\)
arises.  A smooth curve of genus at least two has a unique degree-two
map to a rational curve up to automorphism of that curve: two distinct
such maps would give a birational map to a curve of bidegree \((2,2)\)
in \(\Pj^1\times\Pj^1\), whose arithmetic genus is one.
Thus curve isomorphism is exactly branch-divisor equivalence.
Applying Theorem~\ref{thm:natural-pencil-reconstruction} gives the
asserted injectivity and well-definedness.

Ordered distinct branch points have dimension \(2g+2\).
Normalizing three of them removes the three-dimensional projective
action; taking the finite permutation quotient leaves dimension
\(2g-1\).  Equality of all reductions follows from
\eqref{eq:general-fitting-factorization}.  Finally the two binary
forms in \eqref{eq:same-discriminant-pencils} differ only by a
scalar, a square over \(\C\), and define isomorphic double covers,
whereas their failure schemes have already been separated.
\end{proof}

These are statements about complex isomorphism classes.  They do not
identify moduli stacks, assert a universal family over a coarse moduli
space, or eliminate the different automorphism groups of pencils and
curves.  Nor is the classical pencil--hyperelliptic correspondence a
new Torelli theorem.  The conclusion here is reconstruction of those
objects and of the elementary-divisor information on their boundary
from the abstract nonreduced multiplication-failure scheme, while all
reduced failure data are fixed.

% BEGIN PRESERVED PART 20-finite-neighbourhoods.tex
\section{The exact infinitesimal order of pencil reconstruction}
\label{sec:finite-neighbourhoods}

The all-pencil theorem admits a finite-order strengthening.  It
identifies the first infinitesimal neighbourhood that retains the
pencil, and shows that every preceding neighbourhood is independent
of it.  We use ideal-adic neighbourhoods of the deepest stratum;
these are not arc spaces or jet schemes.

Throughout this section let \(n\ge3\), \(W=\Sym^2V\), and
\begin{equation}\label{eq:pencil-detection-order}
 N=\binom{n+1}{2},\qquad p=N-2,\qquad
 d=n+2p=n^2+2n-4.
\end{equation}
For a pencil \(R\subset W\), put \(S_R=W/R\) and
\(B=\Gr(n,S_R)\subset\widehat D_R\).  This notation for the
socle Grassmannian is local to this section; the cube-zero algebra
is denoted by \(A_R\), as in Section~\ref{sec:natural-pencils}.
Write \(\mathfrak b\subset\OO_{\widehat D_R}\) for the
ideal of \(B\), and define
\begin{equation}\label{eq:finite-neighbourhood-definition}
 Z_R^{[k]}=(B,\OO_{\widehat D_R}/\mathfrak b^{k+1})
       \qquad(k\ge0).
\end{equation}
Thus \(Z_R^{[0]}=B\).  No embedding or identification of its
reduction is supplied with the abstract scheme \(Z_R^{[k]}\).

\begin{theorem}[Sharp finite-neighbourhood reconstruction]
\label{thm:sharp-finite-pencil}
For any two pencils \(R,R'\subset\Sym^2V\), the following
conditions are equivalent:
\begin{equation}\label{eq:finite-pencil-equivalences}
\begin{split}
 Z_R^{[d]}\simeq Z_{R'}^{[d]}
 &\quad\Longleftrightarrow\quad
 \widehat D_R\simeq\widehat D_{R'}\\
 &\quad\Longleftrightarrow\quad
 R'=gR\quad\text{for some }g\in\PGL(V).
\end{split}
\end{equation}
In contrast, for every \(0\le k<d\), all the schemes
\(Z_R^{[k]}\) are isomorphic.  Consequently \(d\) is the
smallest uniform order at which these neighbourhoods classify all
pencils up to projective equivalence.  Singular pencils are included.
\end{theorem}

\subsection{The first relation in the conormal algebra}

Let \(\mathcal U\) be tautological on \(B\).  The open set
of \(X_R\) where projection to \(S_R\) is injective is
canonically \(\Tot E\), where
\(E=\Hom(\mathcal U,V)\): a plane is the graph of a unique
map from its image in \(S_R\) to \(V\).  In particular this
is one open neighbourhood of all of \(B\), not merely a
collection of unrelated formal charts.  Let \(\mathfrak a\)
be its zero-section ideal.  The exact multiplication calculation
\eqref{eq:general-fitting-factorization} gives
\begin{equation}\label{eq:finite-homogeneous-ideal}
 \mathcal I_R=(\det T)I_p(\gamma_R\Sym^2T)
            \subset\Sym E^*.
\end{equation}
It is generated in degree \(d\), with a nonzero degree-\(d\)
piece \(K_R\).  Indeed each quadratic minor has degree \(2p\),
and an invertible \(T\) makes \(\gamma_R\Sym^2T\)
surjective.  There are no terms of lower degree, even on special
pencils.  Homogeneity is an equality of ideals in the actual graph
neighbourhood, not just an equality of initial supports.

\begin{lemma}[Intrinsic recovery from the first relation]
\label{lem:first-relation-cone}
The abstract scheme \(Z_R^{[d]}\) determines, as a graded
\(\OO_B\)-algebra, the complete normal-cone algebra
\(\gr_{\mathfrak b}\OO_{\widehat D_R}\).
\end{lemma}
\begin{proof}
The reduction of \(Z_R^{[d]}\) is \(B\), since its
zero-section ideal is nilpotent and \(B\) is reduced.  Hence
its nilradical \(\mathfrak n\) is intrinsic.  Put
\(L=\mathfrak n/\mathfrak n^2\).  Formula
\eqref{eq:finite-homogeneous-ideal} and \(d>1\) identify
\(L=E^*\).  Multiplication in the given finite scheme recovers
\begin{equation}\label{eq:intrinsic-first-relation}
 K_R=\ker\bigl[\Sym^d L\longrightarrow
                    \mathfrak n^d/\mathfrak n^{d+1}\bigr].
\end{equation}
The last denominator is zero, but retaining it displays the grading.
No splitting of the reduction map is needed to form this kernel.
Because \(\mathcal I_R\) is generated in this single degree,
\begin{equation}\label{eq:finite-cone-completion}
 \gr_{\mathfrak b}\OO_{\widehat D_R}
       =\Sym_{\OO_B}L/(K_R).
\end{equation}
Here \((K_R)\) denotes the homogeneous ideal generated by
\(K_R\), not its radical.  The right side is intrinsic to the
finite scheme, proving the assertion.
\end{proof}

\begin{proof}[Proof of Theorem~\ref{thm:sharp-finite-pencil}]
For \(k<d\), equations \eqref{eq:finite-homogeneous-ideal}
and \eqref{eq:finite-neighbourhood-definition} give
\[
 \OO_{Z_R^{[k]}}=\Sym E^*/\mathfrak a^{k+1}.
\]
A linear identification \(S_R\simeq S_{R'}\) identifies
\(B\), its tautological bundle, and \(E\).  It therefore
identifies every one of these lower-order neighbourhoods.

An abstract isomorphism at order \(d\) induces an isomorphism
of reductions, of conormal bundles, and, by
Lemma~\ref{lem:first-relation-cone}, of complete graded normal
cones.  We spell out why the reconstruction in
Theorem~\ref{thm:automatic-quadratic-coefficients} requires no
information outside this cone.  Its reduced support is the
square-matrix determinant cone; its iterated singular loci give
the rank-one cone.  The relative scheme of maximal linear spaces
in that cone has the two rulings
\(\Pj(V)\times B\) and \(\Pj_B(\mathcal U^*)\).
Exactly the first is projectively trivial, as is seen by restriction
to a Schubert line.  Thus it is intrinsically oriented.  The
induced projective transformation of the trivial factor is constant
on the connected projective base.  Local lifts have the form
\(T\mapsto gTh^{-1}\) for this one \([g]\).
The determinant ideal line is intrinsic, so its ideal quotient in
\eqref{eq:finite-homogeneous-ideal} recovers the residual ideal.
Finally Lemma~\ref{lem:pencil-irreducibility} identifies the entire
left coefficient line \(\C\beta_R\), and exterior duality
\eqref{eq:pencil-exterior-duality} recovers the Pluecker line of
\(R\).  These are operations on the recovered normal cone;
we do not assume that its isomorphism extends to \(X_R\) or
\(\widehat D_R\).  They imply \(R'=gR\).

Conversely such a projective transformation induces the algebra
and failure-scheme isomorphisms, carrying the socle Grassmannian
to the socle Grassmannian and preserving all its ideal powers.
This proves \eqref{eq:finite-pencil-equivalences}, also using
Theorem~\ref{thm:natural-pencil-reconstruction}.  For sharpness,
there are inequivalent pencils in every dimension under consideration:
\(\langle x_1^2,x_1x_2\rangle\) has a common linear factor,
whereas \(\langle x_1^2,x_2^2\rangle\) does not.  Their
neighbourhoods are nevertheless isomorphic at every order below
\(d\).  Thus no smaller order classifies all pencils.
\end{proof}

\subsection{A closed immersion of the full parameter scheme}

The finite-order result has a scheme-theoretic parameter version,
which records infinitesimal variation as well as closed points.
Here tensor coordinates are fixed.  This additional qualification
is important: we are not asserting an equivalence of unmarked
moduli stacks.

Fix an auxiliary \(n\)-dimensional vector space \(U\), and put
\[
 F(V)=\bigwedge^p\Sym^2V,\quad F(U)=\bigwedge^p\Sym^2U,
 \quad M=\dim F(U)=\binom N2,\quad H=\Hom(U,V).
\]
Choose a nonzero scalar representative \(\delta\) of the
determinant polynomial on \(H\); its scalar has no effect on
the subspaces below.  Let
\(K_R\subset\Sym^dH^*\) be the degree-\(d\) space of
\((\det T)I_p(\gamma_R\Sym^2T)\) in these fixed coordinates.

\begin{theorem}[Faithful finite-order parameterization]
\label{thm:finite-parameter-immersion}
The assignment
\begin{equation}\label{eq:finite-parameter-map}
 \Phi:\Gr(2,\Sym^2V)\longrightarrow\Gr(M,\Sym^dH^*),
       \qquad R\longmapsto K_R,
\end{equation}
is a closed immersion.  It is defined on the whole Grassmannian,
including singular pencils.  Its construction and its inverse on
the image commute with arbitrary base change over \(\C\).
In particular it separates families over nonreduced bases and has
injective differential at every pencil.
\end{theorem}

\begin{lemma}[Tensoring a line with a fixed space]
\label{lem:line-subspace-immersion}
For finite-dimensional vector spaces \(A,B_0\), with
\(b=\dim B_0>0\), the map
\(\Pj(A)\to\Gr(b,A\otimes B_0)\),
\(\ell\mapsto\ell\otimes B_0\), is a closed immersion.
\end{lemma}
\begin{proof}
Choose bases \(a_0,\ldots,a_t\) and a basis of \(B_0\).
Above the projective chart
\(\ell=\C(a_0+\sum x_i a_i)\), the subspace is the graph
of the block matrices \(x_i I_b\).  The corresponding
Grassmannian chart has arbitrary \(b\times b\) blocks.  The
image is cut out by their off-diagonal entries and the differences
between their diagonal entries; an inverse reads one diagonal
entry in each block.  These are scheme-theoretic equations over
any \(\C\)-algebra, not just pointwise conditions.  The
preimage of this Grassmannian chart is exactly the stated projective
chart, since the projection has determinant \(a_0^b\).
Such charts cover the image.  The map is proper because its source
is projective, so these chartwise closed immersions give a global
closed immersion.
\end{proof}

\begin{proof}[Proof of Theorem~\ref{thm:finite-parameter-immersion}]
The coefficient map \eqref{eq:general-coefficient-map} becomes
\[
 \mu:F(V)^*\otimes F(U)\longrightarrow\Sym^{2p}H^*.
\]
It is injective: by Lemma~\ref{lem:pencil-irreducibility} its
source is irreducible for \(\operatorname{GL}(V)\times
\operatorname{GL}(U)\), and evaluation at an isomorphism
\(U\simeq V\) shows that the equivariant map is nonzero.
Multiplication by the nonzero polynomial \(\delta\) is
injective.  Thus \(\nu=\delta\mu\) is a fixed linear
inclusion and
\begin{equation}\label{eq:finite-coefficient-factorization}
 K_R=\nu(\C\beta_R\otimes F(U)).
\end{equation}
In particular \(\dim K_R=M\), also on every singular pencil.
By exterior duality the quotient-volume map
\(R\mapsto\C\beta_R\) is the Pluecker closed immersion
\(\Gr(2,W)\hookrightarrow\Pj(F(V)^*)\).
Factor \(\Phi\) through this map, the closed immersion of
Lemma~\ref{lem:line-subspace-immersion}, and the Grassmannian
closed immersion induced by \(\nu\).  This proves the theorem.
All are morphisms of schemes defined by the displayed polynomial
and linear equations.  A closed immersion is a monomorphism and
remains one after arbitrary base change.  This proves the family
assertion, not merely separation of complex-valued points.
\end{proof}

\begin{corollary}[A flat family of finite normal algebras]
\label{cor:finite-flat-family}
Over \(\Gr(2,W)\), the algebras
\begin{equation}\label{eq:finite-flat-algebra}
 Q_R=\Sym H^*/\bigl((K_R)+\mathfrak m^{d+1}\bigr),
 \qquad \mathfrak m=\bigoplus_{j>0}\Sym^jH^*,
\end{equation}
form a finite locally free family of rank
\[
              \binom{n^2+d}{d}-\binom N2.
\]
In fixed coordinates this family determines every base-changed
family of pencils.  All its truncations below degree \(d\)
are constant families.
\end{corollary}
\begin{proof}
Theorem~\ref{thm:finite-parameter-immersion} makes the spaces
\(K_R\) a rank-\(M\) subbundle of the constant degree-\(d\)
bundle, with locally free quotient.  In degree \(j<d\) the
algebra is \(\Sym^jH^*\), and its degree-\(d\) piece is
\(\Sym^dH^*/K_R\); higher degrees vanish.  This proves
local freeness and base-change compatibility.  Summing the
ranks gives the displayed binomial identity.  The kernel of the
coordinate quotient in degree \(d\) recovers \(K_R\), and
\(\Phi\) recovers the pencil family.  This is a statement
about coordinate quotient algebras; it does not discard their
coordinate maps and then infer a family reconstruction theorem.
\end{proof}

\subsection{Classical moduli at a sharp order}

\begin{corollary}[Hyperelliptic classes in finite neighbourhoods]
\label{cor:finite-hyperelliptic}
Let \(g\ge2\), \(n=2g+2\), and
\(d=4g^2+12g+4\).  For pencils with simple discriminant,
\(Z_R^{[d]}\) and \(Z_{R'}^{[d]}\) are isomorphic if
and only if their associated hyperelliptic curves are isomorphic.
Every lower-order neighbourhood is independent of the curve.
Thus the same \((2g-1)\)-dimensional family of curve classes
is invisible to all these lower-order neighbourhoods and is
separated at order \(d\).
\end{corollary}
\begin{proof}
Combine Theorem~\ref{thm:sharp-finite-pencil} with
Corollary~\ref{cor:hyperelliptic-failure-invariant}, and substitute
\(n=2g+2\) in \eqref{eq:pencil-detection-order}.
\end{proof}

The order here concerns a neighbourhood of the entire projective
socle Grassmannian.  Its global ruling distinguishes the constant
\(V\)-factor from the tautological factor.  A single unmarked
Artinian normal slice does not supply this orientation; the
coordinate statement of Corollary~\ref{cor:finite-flat-family}
is deliberately different.  The result is therefore not a claim
that one unmarked local algebra replaces the projective family.
The matrix-pencil classification and the hyperelliptic
correspondence remain classical.  The assertion proved here is
that this classical information first appears at the specified
order in the intrinsic multiplication-failure neighbourhoods.

% BEGIN PRESERVED PART 32-local-inverse-v146.tex
\section{Coefficient-support orientation and a local algebraic inverse}
\label{sec:local-inverse-v146}

For pencils there is a second, entirely local way to distinguish the
two tensor rulings.  It uses the asymmetry of the residual coefficient
module rather than the topology of a projective bundle.  Consequently
the nontrivial-projectivization hypothesis in the preceding family
theorem can be removed, and a single unmarked Artin local algebra
already recovers the pencil.

\subsection{The coefficient module orients its own tensor factors}
If \(K\subset A\otimes B\) is a linear subspace, define its two
factor supports by
\[
 s_A(K)=\im(K\otimes B^*\longrightarrow A),\qquad
 s_B(K)=\im(K\otimes A^*\longrightarrow B).
\]
The maps here are contractions.  Their ranks are invariant under
separate changes of basis and are exchanged by interchanging the
factors.  In particular, if \(K=L\otimes B\), then
\(s_A(K)=L\) and \(s_B(K)=B\).  The support ranks are therefore
\((\dim L,\dim B)\), not the common dimension \(\dim K\).

\begin{lemma}[Asymmetric coefficient support]
\label{lem:coefficient-orientation-v146}
Suppose a homogeneous ideal in the matrix space \(\Hom(U,V)\), with
\(\dim U=\dim V=n\ge3\), has reduced support the determinant cone.
Suppose its intrinsic residual ideal after determinant cancellation
has, in one multiplicity-one Cauchy summand of degree \(q\), the
coefficient subspace
\[
 L\otimes\mathbb S_\lambda U
 \ \subset\ \mathbb S_\lambda V^*\otimes\mathbb S_\lambda U,
 \qquad 0<\dim L<\dim\mathbb S_\lambda U.
\]
A linear isomorphism carrying this ideal to another ideal with a
coefficient subspace of the same displayed type cannot interchange
the two matrix rulings.  The conclusion is unchanged by line twists
of the locally chosen tensor factors.  It also applies fibrewise to
subbundles of this type over a reduced base.
\end{lemma}
\begin{proof}
A linear isomorphism of determinant cones preserves their successive
rank loci and therefore the rank-one cone.  Its action on the two
ruling components either preserves them or interchanges them.  To
make the familiar two alternatives explicit, identify the rank-one
projectivization with \(\Pj(V)\times\Pj(U^*)\).  An automorphism
preserving its two families of maximal linear spaces acts separately
on the two projective factors.  Each factor map is projective linear,
by restricting the ambient linear map to a ruling.  Removing their
linear lifts leaves a map fixing every decomposable line; sums of
tensors with one factor in common show that this remaining map is
scalar.  The preserving alternative is thus a left-right matrix
map.  In the interchanging alternative, compose with transposition
to obtain the preserving alternative.

The determinant ideal is carried to the determinant ideal, so its
ideal quotient is carried to the residual ideal.  The Cauchy
decomposition is functorial for separate changes of matrix factors;
transposition exchanges the factors in its \(\lambda\)-summand.
Put \(m=\dim\mathbb S_\lambda U\), and \(r=\dim L<m\).
The displayed subspace has support ranks \((r,m)\).  Separate
factor maps preserve these ranks, while transposition changes them
to \((m,r)\).  An ideal of the prescribed type on the other side
has support ranks \((r',m)\) with \(r'<m\).  It cannot be the
transposed image.  This excludes interchanging the rulings.

Tensoring a factor by a line changes its Schur module by a line
power, not its contraction rank.  The argument therefore does not
require a specified linear lift of either projective factor.  For
bundle maps it applies to each geometric fibre; an interchange on
any fibre would contradict the same strict rank inequality there.
\end{proof}

For pencils, take
\(F(H)=\bigwedge^{N-2}\Sym^2H\).
Lemma~\ref{lem:pencil-irreducibility} makes this a single irreducible
Schur module, of dimension \(M=\binom N2>1\).  The residual module
is precisely
\(\mathcal L_{\mathcal R}\otimes F(\mathcal U)\), by
Lemma~\ref{lem:moving-coefficients-v146}.  Its support ranks are
\((1,M)\).  The relevant strict inequality holds for every pencil,
including singular pencils; it is not a generic-rank condition.

\begin{theorem}[Removal of the projective-bundle restriction]
\label{thm:unrestricted-moving-v146}
In Theorem~\ref{thm:moving-reconstruction-v146}, the hypothesis
that \(\Pj_B(\mathcal U^*)\) is nontrivial can be omitted.  The
coefficient-sub\-bundle reconstruction, the smallest uniform order,
and the automorphism exact sequence remain valid for every
rank-\(n\) source bundle on a smooth connected projective complex
base.  In particular the base is allowed to be a point.
\end{theorem}
\begin{proof}
The first-relation construction, the recovered homogeneous cone,
and cancellation of the determinant use neither projective
nontriviality nor positive dimension of the base.  They recover
both the unordered rulings and the residual coefficient module
before a ruling is chosen.  At this stage apply
Lemma~\ref{lem:coefficient-orientation-v146} with support ranks
\((1,M)\).  It excludes the transposed alternative and hence
orients the left ruling even when both ruling bundles are trivial.
This replaces exactly the nontrivial-projectivization step in the
proof of Theorem~\ref{thm:moving-reconstruction-v146}.

The left ruling is the trivial bundle \(\Pj(V)\times B\).
Projectivity and connectedness still make its projective-matrix
map constant.  Choose its constant linear lift; uniqueness of the
right bundle map on overlaps gives the same actual isomorphism
\(h\), and coefficient contraction recovers the entire pencil
subbundle.  The converse uses the same homogeneous map.  The
lower-order algebras and the two inequivalent constant pencils in
Corollary~\ref{cor:moving-sharpness-v146} are unchanged.  Finally,
the proof of Theorem~\ref{thm:automorphism-kernel-v146} uses only
the resulting classification of graded data, not the earlier way
of orienting the rulings, so it too applies without the restriction.
\end{proof}

\subsection{Every pencil from one Artin local algebra}
Choose \(n\)-dimensional vector spaces \(U,V\).  Let
\(T=(t_{ij})\) be the universal matrix, and put
\(\mathfrak m=(t_{ij})\).  For a pencil
\(R\subset\Sym^2V\), let \(\gamma_R\) be its quotient map and set
\begin{equation}\label{eq:local-algebra-v146}
 \mathfrak A_R^{[k]}=
 \C[t_{ij}]/\bigl((\det T)I_p(\gamma_R\Sym^2T)
                                      +\mathfrak m^{k+1}\bigr).
\end{equation}
An isomorphism below means an isomorphism of complex algebras,
without a grading, tangent coordinates, matrix decomposition or
chosen generators.

\begin{theorem}[Unmarked local reconstruction]
\label{thm:artin-local-inverse-v146}
For every \(n\ge3\) and every pair of quadratic pencils,
\begin{equation}\label{eq:artin-local-inverse-v146}
 \mathfrak A_R^{[d]}\simeq\mathfrak A_{R'}^{[d]}
 \quad\Longleftrightarrow\quad
 R'=(\Sym^2g)R\quad\text{for some }g\in\operatorname{GL}(V).
\end{equation}
The algebra has embedding dimension \(n^2\) and length
\(\binom{n^2+d}{d}-\binom N2\).  Every lower-order algebra
\(\mathfrak A_R^{[k]}\), \(k<d\), is independent of \(R\).
Thus \(d=n^2+2n-4\) is the smallest uniform reconstruction order
already among these single-point objects.
\end{theorem}
\begin{proof}
Theorem~\ref{thm:unrestricted-moving-v146} applies with
\(B=\Spec\C\) and the bundle \(U\).  To spell out what its
intrinsic construction uses in this case, the unique maximal ideal
\(\mathfrak n\) of the finite algebra determines its tangent
space and the multiplication kernel
\[
 K_R=\ker\left[\Sym^d(\mathfrak n/\mathfrak n^2)
                                      \longrightarrow\mathfrak n^d\right].
\]
This is precisely the degree-\(d\) ideal in the chosen presentation;
nonlinear changes of generators contribute only terms above degree
\(d\) to a degree-\(d\) relation.  Hence an ungraded algebra
isomorphism identifies the homogeneous cones generated by these
kernels.  Their determinant reductions recover the unordered
matrix factors.  The residual support ranks \((1,M)\) rule out
transposition, and the remaining left-right map carries the left
coefficient line to the left coefficient line.  Exterior duality
then gives \(R'=(\Sym^2g)R\).  Conversely such a map, together
with any chosen right identification, induces the algebra
isomorphism.  There are no relations below degree \(d\), so the
embedding dimension and the lower-order assertion are immediate.
The relation subspace at degree \(d\) has dimension \(M\), giving
the asserted length.  The inequivalent common-factor and
no-common-factor pencils prove uniform minimality as before.
\end{proof}

In particular the earlier Schubert-line construction is a sufficient
geometric realization, not a lower bound on the dimension of a
supporting reduced scheme.  The local inverse just proved retains no
curve at all.  Its additional ingredient is the coefficient-support
asymmetry, which was not used in the earlier geometric orientation
argument.  This is a stronger inverse statement, not a conclusion
from knowing the discriminant or the reduced determinant cone.

\subsection{The local automorphism kernel is explicit}
Put \(a=n^2\), \(\ell=\length\mathfrak A_R^{[d]}\), and
\(G_R=\{g\in\operatorname{GL}(V):(\Sym^2g)R=R\}\).
The symbol \(\mathfrak n^2\) in the next statement is the square
of the maximal ideal of the finite algebra.

\begin{proposition}[Tangent-identity automorphisms]
\label{prop:local-automorphisms-v146}
There is a split exact sequence of complex automorphism groups
\begin{equation}\label{eq:local-aut-sequence-v146}
 1\longrightarrow U_R\longrightarrow
 \operatorname{Aut}_{\C}(\mathfrak A_R^{[d]})
 \longrightarrow (G_R\times\operatorname{GL}(U))/\C^*
 \longrightarrow1,
\end{equation}
where the scalar subgroup is diagonal.  The kernel \(U_R\) is a
unipotent algebraic group whose underlying variety is an affine
space of dimension
\[
                            a(\ell-1-a).
\]
More explicitly, after choosing the \(a\) degree-one generators
\(x_1,\ldots,x_a\), its elements are exactly the substitutions
\(x_i\mapsto x_i+u_i\) with arbitrary \(u_i\in\mathfrak n^2\).
The affine-space parametrization is not an assertion that this
group law is additive.
\end{proposition}
\begin{proof}
The local inverse and the support-rank argument identify the linear
stabilizer of \(K_R\) with the indicated left-right quotient;
transposition is absent.  These linear maps lift homogeneously,
giving the quotient and splitting.  An automorphism in the kernel
is the identity on \(\mathfrak n/\mathfrak n^2\), and consequently
has exactly the displayed generator images.

Conversely choose arbitrary \(u_i\in\mathfrak n^2\).  Under this
substitution a homogeneous defining relation of degree \(d\)
changes only by terms in \(\mathfrak n^{d+1}=0\), so all defining
relations are respected.  The induced map is the identity on the
associated graded algebra, hence is invertible as a finite filtered
linear map and therefore as an algebra map.  There are no further
conditions on the \(u_i\).  Reading the generator images and
performing substitutions are polynomial mutually inverse maps
between the kernel and the vector space
\((\mathfrak n^2)^{\oplus a}\).  Since
\(\dim\mathfrak n^2=\ell-1-a\), this proves the variety and
dimension assertions.  On a basis ordered by degree every such
automorphism is upper unitriangular; the algebraic group is
therefore unipotent.  Its law includes higher substitution terms,
which accounts for its generally noncommutative structure.
\end{proof}

The positive-dimensional family statement remains useful after the
local inverse: it recovers an entire varying subbundle and the
actual source bundle up to one global left transformation, whereas
separate fibrewise local isomorphisms need not glue.  The geometric
automorphism statements here still do not identify unmarked moduli
stacks or all deformations over nonreduced parameter bases.

\section{The intrinsic polarization and a projective moduli completion}
\label{sec:polarization-v154}
Put $P=\Gr(2,\Sym^2V)$, $N=\binom{n+1}{2}$,
$M=\binom N2$, and $D=n^2+2n-4$. Fix temporarily an auxiliary
$n$-space $U$ and put $E=\Hom(U,V)^*$. Write $L=\OO_P(1)$ for
the Pluecker line and $\Gamma=\Phi(P)\subset\Gr(M,\Sym^D E)$ for
the closed image of Theorem~\ref{thm:finite-parameter-immersion}.
The group in this section is $\operatorname{SL}(V)$, acting by
congruence on pencils and by its induced action on $E$, with $U$
fixed. It is not the full group $\operatorname{GL}(E)$.

\begin{theorem}[Polarized completion of the effective failure moduli]
\label{thm:polarized-completion-v154}
There is an $\operatorname{SL}(V)$-equivariant isomorphism
\begin{equation}\label{eq:polarization-v154}
 \Phi^*\OO_{\Gr(M,\Sym^D E)}(1)\simeq L^{\otimes M}
                 \otimes C_U,
\end{equation}
where $C_U$ is a constant one-dimensional space on which
$\operatorname{SL}(V)$ acts trivially. Consequently the stable and
semistable loci of $\Gamma$, for the restricted Pluecker
linearization, are respectively the images of $P^s$ and $P^{ss}$.
The quotient defined by the complete restricted section ring is
canonically
\begin{equation}\label{eq:git-completion-v154}
 \overline{\mathcal M}^{\mathrm{fail}}_n
 =\Proj\!\left(\bigoplus_{q\geq0}
                  H^0(P,L^{\otimes q})^{\operatorname{SL}(V)}\right)
 \simeq\Gamma/\!/\operatorname{SL}(V).
\end{equation}
It is a normal projective variety. It contains the geometric quotient
of the simple-discriminant pencils as a dense open subvariety, and
has dimension $n-3$. The complement of that open is the image of the
semistable discriminant-zero locus.

The intrinsic inverse identifies the effective semistable failure
stack with $[P^{ss}/\PGL(V)]$ and gives a moduli morphism to
\eqref{eq:git-completion-v154}. On framed charts the failure-algebra
family extends over all $P$, remains finite locally free, and
commutes with arbitrary base change. These assertions do not require,
and do not assert, the existence of a universal ordinary algebra on
the coarse quotient or properness of the entire quotient stack.
\end{theorem}
\begin{proof}
Let $\mathcal K$ be the tautological relation bundle pulled back by
$\Phi$. The coefficient factorization
\eqref{eq:finite-coefficient-factorization} identifies it with
$\mathcal B\otimes F(U)$, where $\mathcal B$ is the quotient-volume
coefficient line. Exterior duality gives
$\mathcal B\simeq L^{-1}$ up to a fixed determinant character of
$\operatorname{GL}(V)$. That character is trivial on
$\operatorname{SL}(V)$. Taking the dual determinant gives
\eqref{eq:polarization-v154}, with the harmless constant factor
$(\det F(U))^*$ and the corresponding fixed determinant line of $V$.
A trivialization of the constant factor identifies the section rings;
a different trivialization rescales each graded piece and has no
effect on its projective spectrum. The section ring for $L^M$ is
the $M$th Veronese of the section ring for $L$. Taking invariants
commutes with this operation. The Hilbert--Mumford weights are
multiplied by $M>0$, so stability and semistability are unchanged.
This also explains why one uses the complete section ring on
$\Gamma$, rather than assuming surjectivity of unspecified ambient
invariant restrictions.

Finite generation and the projective good quotient are the reductive
GIT construction of \cite[Chapter~1]{MFK1994}. The complete section
ring of the normal projective variety $P$ is integrally closed: it
is the intersection of the valuation conditions for rational sections
of powers of $L$ at the prime divisors of $P$. Its invariant subring
is integrally closed as well. Indeed an element of its fraction
field integral over the invariants is integral over the original
ring, belongs to that ring by normality, and is invariant. This proves
normality of the projective quotient.

For completeness, the simple-discriminant open is nonempty and
stable. In a frame $(Q_0,Q_1)$ of a pencil, take the discriminant of
$\det(sQ_0+tQ_1)$. Under a frame change $h\in\operatorname{GL}_2$
it has weight $(\det h)^{n(n-1)}$, and it is invariant under
$\operatorname{SL}(V)$. Thus it is a nonzero invariant section of
$L^{n(n-1)}$. It is nonzero at a diagonal pencil with $n$ distinct
roots, so its nonvanishing locus is semistable and saturated.
A pencil there has a nonsingular member. Relative to that member,
the other member is self-adjoint with distinct eigenvalues;
orthogonal diagonalization identifies its congruence class with an
unordered configuration of $n$ distinct points of $\Pj^1$.
The latter configurations are stable for the usual binary-form
linearization: the root-multiplicity criterion is strict since
$1<n/2$ for $n\geq3$. Their quotient is geometric and separated.
It follows that the congruence orbits are closed in the saturated
nonvanishing locus. Their stabilizers are finite: the induced
projective transformation permutes a set of at least three points,
and the kernel consists of diagonal sign changes modulo scalars.
They are therefore stable. This gives a geometric dense open of
dimension $n-3$. The complement is the image of the invariant
section's zero locus by the good-quotient property. This reasoning
uses classical pencil classification and GIT, not a new classification
of binary forms; compare \cite{FevolaMandelshtamSturmfels,MFK1994}.

Finally apply Theorem~\ref{thm:pencil-stack-equivalence} and the
rigidification results of Section~\ref{sec:rigidification-v151}.
The Pluecker polarization is taken on the recovered pencil, not on
an arbitrary first-relation orbit. A sufficiently divisible power
of $L$, for example $L^n$, has trivial central character and descends
from $\operatorname{SL}(V)$ to $\PGL(V)$. A change of auxiliary frame
changes the constant factor in \eqref{eq:polarization-v154}, not the
resulting projective quotient. On the framed Grassmannian the extended
algebras are exactly those of Corollary~\ref{cor:finite-flat-family}.
Their equivariance yields the stated classifying morphisms. Removing
ineffective inertia does not turn these into a universal ordinary
algebra on the coarse space, and no such descent is used.
\end{proof}

\begin{remark}[What is compactified]\label{rem:git-scope-v154}
The completion is the classical polarized pencil quotient recovered
from the intrinsic failure invariant, with its polarization verified
by \eqref{eq:polarization-v154}. It is not a new construction of
classical GIT. Its role is to supply a proper coarse target for the
effective invariant and a specified class of semistable limiting
families. Strictly semistable points represent closed-orbit
representatives, not every orbit in their fibres. In particular Segre
data on the stack need not descend to an invariant separating all
points of a coarse-quotient fibre. For $n=3$ the quotient is a point
and its discriminant boundary can be empty. The full
$\operatorname{GL}(E)$ common limit of
Theorem~\ref{thm:universal-boundary-v152} is a different construction.
\end{remark}

\section{Spectral divisor fibres and the recovery of degeneration type}
\label{sec:spectral-divisor-v154}
The preceding completion concerns the native congruence action. To
retain degeneration type before passage to a coarse quotient, we
associate a tower of binary divisor incidences to a regular pencil.
This construction is intrinsic after the unmarked inverse and is
unchanged by a simultaneous change of the pencil and source frames.

Let $\Lambda=\Pj(R)$ be the pencil line and let
$\mathcal T_R$ be the cokernel of the universal symmetric map on
$\Lambda$. A regular pencil means that its determinant is not
identically zero. Thus $\mathcal T_R$ is a torsion sheaf. Define
$D_j(R)$ by
\begin{equation}\label{eq:spectral-Fitting-v154}
 \II_{D_j(R)}=\Fitt_j(\mathcal T_R),\qquad 0\leq j\leq n,
\end{equation}
where $D_n=\varnothing$. These are effective divisors on the smooth
curve $\Lambda$. Write $s_j=\deg D_j$. If $s_j>0$, let $G_j$ be
its binary equation and put
\[
 J_j=G_j H^0(\Lambda,\OO_\Lambda(1))
       \subset H^0(\Lambda,\OO_\Lambda(s_j+1)).
\]
Changing the scalar representative of $G_j$ does not change $J_j$.
Let $\mathcal F_j(R)$ be the whole scheme fibre at $J_j$ of the
degree-one divisor normalization $\nu_1$, with $r=2$ and $D=s_j+1$.
For $s_j=0$ put $\mathcal F_j=\varnothing$.

\begin{theorem}[A fibre-theoretic Segre invariant]\label{thm:spectral-fibres-v154}
There is a natural identification of finite schemes over $\Lambda$
\begin{equation}\label{eq:fibre-divisor-v154}
                        \mathcal F_j(R)\simeq D_j(R).
\end{equation}
At a spectral point $x$, write the positive elementary divisors of
$\mathcal T_R$ as
$\lambda_{x,1}\geq\cdots\geq\lambda_{x,r_x}>0$, and extend this
sequence by zeroes. Then
\begin{equation}\label{eq:fibre-length-v154}
 (\mathcal F_j(R))_x=\Spec\C[\epsilon]/(\epsilon^{b_{x,j}}),
 \qquad b_{x,j}=\sum_{i>j}\lambda_{x,i},
\end{equation}
with the empty-scheme convention when $b_{x,j}=0$. Consequently
\begin{equation}\label{eq:Segre-difference-v154}
                 \lambda_{x,j+1}=b_{x,j}-b_{x,j+1}.
\end{equation}
The tower, with its maps to the same intrinsic spectral line,
therefore determines the full Segre symbol. Conversely, the
elementary-divisor partitions together with their spectral support
on that line determine the tower. The tower retains the projective
configuration of the support as well as the Segre symbol.
It is an invariant of the unmarked finite failure algebra.

The construction is compatible with base change on every family
stratum where the Fitting ideals in \eqref{eq:spectral-Fitting-v154}
define relative effective Cartier divisors of fixed degrees. On such
a stratum $\mathcal F_j$ is finite locally free of degree $s_j$.
These hypotheses are part of the relative statement; no assertion of
flatness across a jump of $s_j$ is intended.
\end{theorem}
\begin{proof}
Over the discrete valuation ring $\OO_{\Lambda,x}$, Smith normal
form gives diagonal entries with exponents
$\lambda_{x,r_x},\ldots,\lambda_{x,1}$ and $n-r_x$ zero
exponents. The minors of size $n-j$ have greatest common divisor
of order $\sum_{i>j}\lambda_{x,i}$. This proves the formula for
$D_j$, including its nonreduced structure. It is the usual
Fitting-ideal description of elementary divisors, used here on the
intrinsic spectral line.

The residual two-space $H^0(\OO_\Lambda(1))$ is gcd-free. Hence
Theorem~\ref{thm:binary-fibres-v153} identifies the degree-one
normalization fibre at $J_j$ with the scheme of linear factors of
$G_j$. On an affine chart write a monic linear factor as $z-a$.
Monic division gives $G_j=(z-a)H+G_j(a)$, so the entire factor
scheme is $G_j(a)=0$, not its reduction. These descriptions agree
on overlaps of the projective line. They give
\eqref{eq:fibre-divisor-v154} and the local algebra
\eqref{eq:fibre-length-v154}. Taking successive differences proves
\eqref{eq:Segre-difference-v154}. The common map to $\Lambda$ is
essential: forgetting it separately in each member of the tower
would discard the alignment of the spectral supports.

Pencil frame changes act on $\Lambda$ by projective transformations;
source congruence changes give isomorphic cokernels and preserve all
Fitting ideals. The intrinsic inverse of
Theorem~\ref{thm:artin-local-inverse-v146} therefore defines this
invariant from the unmarked algebra. The equivalence with Segre data
uses the classical elementary-divisor convention of
\cite[\S5]{FevolaMandelshtamSturmfels}.

For a family with the stated Cartier and degree conditions, the same
monic division argument works over any base algebra, including one
with nilpotents. It represents the divisor's functor of points, so
it commutes with base change. A relative effective divisor of fixed
degree on a projective-line bundle is finite locally free of that
degree. Local choices of $G_j$ differ by units and glue. Formation
of Fitting ideals itself commutes with base change; requiring them
to remain flat Cartier divisors is precisely what permits the
additional finite-flat conclusion.
\end{proof}

\begin{example}[Equal discriminant, different fibre towers]
\label{ex:segre-fibres-v154}
At a root of multiplicity four, the partitions $(3,1)$ and $(2,2)$
have the same zeroth fibre $\C[\epsilon]/(\epsilon^4)$.
Their first fibres are respectively $\C$ and
$\C[\epsilon]/(\epsilon^2)$; all higher fibres are empty there.
Thus the discriminant alone does not detect the distinction, whereas
the divisor-fibre tower does. Both are realized by regular symmetric
pencils: for a block of size $l$, take the reversal matrix $H_l$ and
the symmetric pair $(H_l,H_lJ_l(\alpha))$, where $J_l(\alpha)$ is
a Jordan block. Direct sums realize every indicated partition.
Additional simple blocks may be appended. The example is a statement
on the regular-pencil stack, not a claim that distinct Segre strata
always remain distinct after strictly semistable GIT identification.
\end{example}

\begin{remark}[The two divisor constructions]\label{rem:two-divisors-v154}
The primitive first relation of an actual pencil has exact common
factor $(\det T)^{n-1}$, by Lemma~\ref{lem:primitive-pencil}.
Its common-factor degree does not jump when the pencil changes
Segre type. The divisors $D_j$ above instead lie on the recovered
spectral line and define auxiliary binary first-relation subspaces
$J_j$. Thus \eqref{eq:Segre-difference-v154} relates pencil data to
intrinsically induced divisor ramification, not to a nonexistent
change in the conductor of the original exact-gcd locus. The
identification supplies a moduli consequence of the all-binary-fibre
theorem while respecting this distinction.
\end{remark}

\part*{Part II. Divisor incidence and homological fibres}

% BEGIN PRESERVED PART 37-common-divisor-strata-v149.tex
\section{Common-divisor strata and intrinsic families}
\label{sec:common-divisor-strata}

The coefficient condition in Theorem~\ref{thm:recognition-dichotomy-v148}
is not imposed in this section.  We instead retain the common divisor
of a space of first relations.  This geometric condition gives a
larger smooth parameter space, includes relations without either
factor symmetry, and has a scheme-theoretic inverse.  The distinction
between a stratum and the condition on its geometric points is
essential when the parameter base is nonreduced.

Write \(S_j(E)=\Sym^j E\), where \(E\) is a complex vector space.
For \(r\ge2\), denote by
\[
 X_{h,r}(E)\subset\Gr(r,S_h(E))
\]
the locus of subspaces with greatest common divisor equal to \(1\),
up to a scalar.  Empty Grassmannians are omitted throughout.  When
\(h=0\) this notation can also be used for \(r=1\).

\begin{lemma}[Universal common divisor]\label{lem:universal-gcd}
Let \(D=g+h\).  The spaces of \(r\) degree-\(D\) forms whose exact
greatest common divisor has degree \(g\) form a locally closed
reduced subscheme \(Z_g\) of \(\Gr(r,S_D(E))\).  Multiplication
induces an isomorphism
\begin{equation}\label{eq:universal-gcd}
 \Pj(S_g(E))\times X_{h,r}(E)\ \xrightarrow{\ \sim\ }\ Z_g,
                 \qquad ([f],J)\longmapsto fJ.
\end{equation}
In particular, the common-divisor line and the residual subbundle of
a family with classifying map in \(Z_g\) are recovered uniquely and
commute with arbitrary complex base change.
\end{lemma}
\begin{proof}
For each \(a\), multiplication from
\(\Pj(S_a(E))\times\Gr(r,S_{D-a}(E))\) has closed image: its domain
is projective.  The union for \(a>g\) is closed.  Remove this union
from the image for \(a=g\), and give the resulting locally closed
set its reduced structure.  Its points are exactly those in the
statement.  The inverse image of this set under multiplication is
\(\Pj(S_g(E))\times X_{h,r}(E)\).  The same closed-image argument
with \(D=h\) shows that \(X_{h,r}(E)\) is open.

The restricted multiplication map is proper, by base change.  It
is injective on geometric points over every algebraically closed
extension of \(\C\), since a greatest common divisor is unique up
to scalar in a polynomial ring.  Its differential is injective as
well.  Indeed, a tangent vector in its kernel satisfies
\[
                   \dot f\, j+f\dot j\in fJ\quad(j\in J).
\]
Consequently \(f\mid\dot f\,j\) for every \(j\in J\).  For each
irreducible factor of \(f\), some element of \(J\) is not divisible
by that factor.  Comparison of its valuation gives
\(f\mid\dot f\).  Equality of degrees makes \(\dot f\) a scalar
multiple of \(f\); then \(\dot J=0\) in
\(\Hom(J,S_h(E)/J)\).

In characteristic zero the geometric-point injectivity and this
differential calculation imply that the map is universally
injective and unramified.  A proper, universally injective,
unramified morphism is a closed immersion
\cite[Lemma~41.7.2]{StacksClosedImmersion}.
Here its image is all of the reduced target.  The defining ideal
of a closed immersion with this support is zero in a reduced
scheme.  This proves the isomorphism, including its assertion for
nonreduced test schemes.  Equivalently, one can check properness
on the base change of the original projective multiplication map;
that base change has the same underlying open preimage.  Its source
is reduced because it is this open subscheme of the original
smooth product.
\end{proof}

\begin{remark}\label{rem:gcd-scheme-condition}
A map from a nonreduced scheme is required to factor through the
locally closed \emph{scheme} \(Z_g\), not merely to have geometric
points in its underlying set.  These are different requirements.
For example, over \(\C[\epsilon]/(\epsilon^2)\) the span of
\(x^2+\epsilon y^2,xy\) specializes to a pair with common divisor
\(x\), but has no factor \(x+\epsilon(ax+by)\).  Comparison of the
\(y^2\) coefficient in the first form proves this.  The stratum in
Lemma~\ref{lem:universal-gcd} excludes this transverse infinitesimal
family.  No use of a fibrewise radical can replace its scheme structure.
\end{remark}

Now let \(\dim V=\dim U=n\ge3\),
\(E=V^*\otimes U\), and \(\delta=\det T\).  Denote by
\(G\subset\operatorname{GL}(E)\) the stabilizer of the line
\(\C\delta\).  By Lemma~\ref{lem:linear-preserver-v147},
\begin{equation}\label{eq:determinant-group}
 G^\circ=(\operatorname{GL}(V)\times\operatorname{GL}(U))/\mathbb G_m,
                 \qquad G/G^\circ=\mathbf Z/2\mathbf Z.
\end{equation}
The second component exchanges the two tensor factors; the diagonal
scalars in the displayed product act trivially on the matrix space.
Fix \(s\ge1\) and \(D=sn+h\).  Inside \(Z_{sn}\), impose the condition
that the common-divisor line is the \(s\)-th power of a determinant
form after a linear change of the \(n^2\) variables.  Call the resulting
locally closed scheme \(Z^{\det}_{s,h,r}\).

\begin{theorem}[Intrinsic determinant-divisor stratum]
\label{thm:determinant-stratum}
There are canonical isomorphisms of schemes and algebraic stacks
\begin{equation}\label{eq:determinant-stratum}
 Z^{\det}_{s,h,r}\simeq
 \operatorname{GL}(E)\times^G X_{h,r}(E),\qquad
 [Z^{\det}_{s,h,r}/\operatorname{GL}(E)]\simeq[X_{h,r}(E)/G].
\end{equation}
They hold over arbitrary complex parameter schemes.  The stack is
smooth.  At a residual relation \(J\), its tangent complex, in degrees
\(-1,0\), is
\begin{equation}\label{eq:ambient-tangent-complex}
 [\,\operatorname{Lie}(G)\longrightarrow
                    \Hom(J,S_h(E)/J)\,],
             \qquad X\longmapsto(j\mapsto Xj\bmod J).
\end{equation}
No invariance of \(J\) under either full factor group is assumed.
The equivalence identifies stabilizers, infinitesimal families, and
specializations whose classifying maps remain in this stratum.
\end{theorem}
\begin{proof}
The power map \(\Pj(S_n(E))\to\Pj(S_{sn}(E))\),
\([f]\mapsto[f^s]\), is a closed immersion in characteristic zero.
It is proper and geometrically injective, and its differential has
kernel given by \(sf^{s-1}\dot f\in\C f^s\), hence by
\(\dot f\in\C f\).  The same closed-immersion criterion as above
applies.  The orbit of \([\delta]\) is the smooth locally closed
homogeneous space \(\operatorname{GL}(E)/G\).  The preimage of its
power orbit under the first projection in
\eqref{eq:universal-gcd} is therefore a locally closed subscheme.
It is exactly \(Z^{\det}_{s,h,r}\).

Pulling back along the \(G\)-torsor
\(\operatorname{GL}(E)\to\operatorname{GL}(E)/G\) trivializes the
common-divisor form, leaving the unrestricted residual relation
\(J\in X_{h,r}(E)\).  Descent gives the first isomorphism in
\eqref{eq:determinant-stratum}; descent of equivariant maps and
morphisms gives the second.  This proof uses morphisms of schemes,
not a comparison only on complex points.  It also supplies local
lifts of all isomorphisms over an fppf cover.

The space \(X_{h,r}(E)\) is open in a Grassmannian and \(G\) is
smooth.  The quotient is algebraic and smooth
\cite[Theorem~97.17.2]{StacksQuotient}; see also
\cite[Lemma~101.33.7]{StacksSmoothQuotient}.
For completeness, write a first-order residual subspace as the
graph of \(\epsilon\phi:J\to S_h(E)/J\).  An infinitesimal change
\(1+\epsilon X\) replaces \(\phi\) by \(\phi+X|_J\bmod J\).
This proves the tangent complex.  Its kernel records infinitesimal
automorphisms and its cokernel records first-order isomorphism
classes.  The same torsor description proves compatibility with
specialization.  No assertion is made about limits outside the
specified common-divisor stratum.
\end{proof}

The next step connects these parameter spaces with ungraded finite
algebras.  It is useful to state the family condition intrinsically,
since an absolute nilradical is unsuitable over a nonreduced base.
Put \(e=\dim E\) and
\(\ell=\binom{e+D}{D}-r\).  Consider finite locally free commutative
unital algebras \(\mathcal A\) of rank \(\ell\) on a complex scheme
\(T_0\), subject to the following conditions.  The map
\(\epsilon(a)=\ell^{-1}\operatorname{tr}(M_a)\) is multiplicative.
For \(\mathcal N=\ker\epsilon\), one has
\(\mathcal N^{D+1}=0\); all quotients of its power filtration are
locally free; \(\mathcal E=\mathcal N/\mathcal N^2\) has rank \(e\);
and multiplication gives isomorphisms
\(\Sym^j\mathcal E\simeq\mathcal N^j/\mathcal N^{j+1}\) for
\(j<D\).  In degree \(D\) its kernel has rank \(r\).  These
conditions specify the maximal-lower-Hilbert-function stratum; they
do not assert that arbitrary smoothings of a local algebra remain
in it.

\begin{proposition}[Relative first-relation algebra]
\label{prop:relative-first-relation}
Every such algebra is locally isomorphic, over its base, to
\[
 \mathcal A(\mathcal K)=
 \Sym\mathcal E/((\mathcal K)+\mathfrak m^{D+1}),\qquad
 \mathcal K=\ker(\Sym^D\mathcal E\longrightarrow\mathcal N^D).
\]
Its trace augmentation, \(\mathcal E\), and \(\mathcal K\) are
intrinsic and compatible with arbitrary base change in this
stratum.  For two such algebras, taking the linear part gives a
locally surjective map of isomorphism sheaves to the sheaf of
isomorphisms of \((\mathcal E,\mathcal K)\).  Its kernel is the
smooth substitution group whose underlying scheme in a local
presentation is
\(\Hom(\mathcal E,\mathcal N^2)\).

After quotienting morphisms by this tangent-identity kernel and
fppf stackifying, the algebras with first-relation map in
\(Z^{\det}_{s,h,r}\) have stack \([X_{h,r}(E)/G]\).
\end{proposition}
\begin{proof}
Locally lift a basis of \(\mathcal E\) to \(\mathcal N\).
Nilpotence and multiplication on the successive quotients make the
induced map from the truncated symmetric algebra surjective.
A relation with a nonzero term of least degree \(j<D\) would give
a nonzero kernel in degree \(j\) of the associated graded algebra,
which is impossible.  Its entire kernel therefore lies in degree
\(D\) and is precisely \(\mathcal K\).  Changing a lift by an
element of \(\mathcal N^2\) does not alter a product of \(D\)
generators, since the difference has degree at least \(D+1\).
This also proves the asserted uniqueness of the first relation.

In this presentation multiplication by a positive-degree element
is strictly filtration-increasing.  Its trace is zero over any
base ring.  Multiplication by a scalar has trace \(\ell\) times
that scalar, proving that the presentation's augmentation is the
specified intrinsic trace augmentation.  All filtration sequences
are locally split sequences of vector bundles, so their kernels
and quotients commute with base change.

An isomorphism must preserve trace and induces a linear map carrying
\(\mathcal K\) to \(\mathcal K'\).  Conversely such a map lifts
locally to an isomorphism by the homogeneous presentations.  All
other lifts are obtained by substituting generators plus arbitrary
elements of \(\mathcal N'^2\); the inverse is constructed degree by
degree in a finite filtration.  This gives the stated kernel and
its smoothness, without replacing its substitution law by addition.
After identifying morphisms with the same linear part, frames of
\(\mathcal E\) give the quotient
\([Z^{\det}_{s,h,r}/\operatorname{GL}(E)]\).
Theorem~\ref{thm:determinant-stratum} completes the proof.
\end{proof}

\begin{remark}\label{rem:algebra-inertia}
The original algebra stack is not this quotient stack: it has the
additional nonlinear inertia just computed.  The phrase
\emph{effective first-relation stack} will refer to the specified
quotient on morphisms followed by stackification.  The construction
retains the full automorphism group before taking that quotient.
It does not identify algebras with the same geometric fibres, nor
does it discard nilpotent parameters.  In particular, its relation
to a stratum of first relations is stronger than a bijection of
isomorphism classes over \(\C\).
\end{remark}

% BEGIN PRESERVED PART 40-canonical-divisor-boundary-v151.tex
\section{The common-divisor boundary and its normalization}
\label{sec:canonical-boundary-v151}

The exact-common-divisor stratum has a canonical completion which need
not be normal.  Its normalization records a divisor, and its failure to
be an isomorphism is controlled by two different phenomena: choices of
that divisor and overlap with the residual common factor.  We establish
these assertions for arbitrary spaces of forms.  No determinant, Cauchy,
or Pluecker equation is imposed in this section.

Let \(E\) have dimension \(e\geq2\), put \(S_j=\Sym^j E\) and
\(s_j=\dim S_j=\binom{e+j-1}{j}\), and fix
\[
             g,h\geq1,\qquad D=g+h,\qquad 2\leq r\leq s_h.
\]
Write \(\mathscr G=\Gr(r,S_D)\).  A \emph{degree-\(g\) divisor
incidence} over a complex scheme \(T_0\) consists of a line subbundle
\(\mathcal L\subset S_g\otimes\OO_{T_0}\) and a rank-\(r\)
subbundle \(\mathcal K\subset S_D\otimes\OO_{T_0}\) such that
\begin{equation}\label{eq:divisor-incidence-v151}
 \mathcal K\longrightarrow
 (S_D\otimes\OO_{T_0})/(\mathcal L\cdot S_h)=0.
\end{equation}
Subbundle always means that the quotient is locally free.  In
particular, the form generating \(\mathcal L\) is nonzero on every
geometric fibre.  Condition~\eqref{eq:divisor-incidence-v151} is the
vanishing of a bundle map, not a condition imposed only on closed points.

\begin{proposition}[Canonical divisor incidence]
\label{prop:incidence-v151}
The incidence functor is represented, with its full scheme structure, by
\[
 \mathscr X_g=\Pj(S_g)\times\Gr(r,S_h).
\]
Its forgetful morphism is multiplication
\begin{equation}\label{eq:normalization-map-v151}
 \nu_g:\mathscr X_g\longrightarrow\mathscr G,
                   \qquad ([f],J)\longmapsto fJ.
\end{equation}
This morphism is finite.  Its scheme-theoretic image \(Y_g\) is an
integral projective scheme.  If \(K\) is a geometric point of \(Y_g\)
and \(F_K=\gcd(K)\), the underlying set of \(\nu_g^{-1}(K)\) is
\begin{equation}\label{eq:divisor-fibre-set-v151}
       \{[f]: f\mid F_K,\ \deg f=g\},
\end{equation}
with residual space \(J=K/f\).  The incidence construction and its
scheme-theoretic fibres commute with arbitrary complex base change.
\end{proposition}
\begin{proof}
On \(\Pj(S_g)\), multiplication from the tautological line tensored
with \(S_h\) to \(S_D\) has constant rank \(s_h\).  Indeed,
multiplication by a nonzero polynomial over a field is injective.
A nonzero maximal minor locally splits this map, so its image and
cokernel are vector bundles, also after a nonreduced base change.
The closed incidence~\eqref{eq:divisor-incidence-v151} is therefore
the relative Grassmannian of rank-\(r\) subbundles of
\(\mathcal L\otimes S_h\).  Tensoring by \(\mathcal L^{-1}\)
identifies it canonically with the stated product.  In more elementary
terms, \(\mathcal K\) has a unique inverse image in
\(\mathcal L\otimes S_h\); its quotient there is locally free
because the quotient of \(S_D/\mathcal K\) onto
\(S_D/(\mathcal L\cdot S_h)\) is a surjection of vector bundles.
This proves the functorial assertion, including its scheme structure.

The source of \(\nu_g\) is projective.  Unique factorization in the
polynomial ring over an algebraically closed field gives
\eqref{eq:divisor-fibre-set-v151}.  A fixed polynomial has only finitely
many divisors of a fixed degree, up to scalar.  The morphism is thus
proper with finite fibres and is finite
\cite[Tag~02LS]{StacksFiniteV151}.
Define \(Y_g\) by the kernel of
\(\OO_{\mathscr G}\to(\nu_g)_*\OO_{\mathscr X_g}\), as in
\cite[Tag~01R5]{StacksImageV151}.  Its integrality follows from the
integrality of \(\mathscr X_g\): on an affine open meeting the image,
the coordinate ring of the image injects into the domain which is the
coordinate ring of its finite inverse image.  Reducedness is therefore
a consequence of this incidence construction, not an extra choice.
Representability of incidence gives the base-change assertion.
\end{proof}

\begin{remark}\label{rem:degree-divisor-v151}
The condition defining \(Y_g\) is the existence of a divisor of
\emph{degree \(g\)}.  In several variables it is not the condition
\(\deg\gcd(K)\geq g\): an irreducible polynomial of larger degree
need not have a degree-\(g\) factor.  Nor do we assert that formation
of the scheme-theoretic image, or normalization itself, commutes with
arbitrary nonflat base change.  The assertion for arbitrary bases in
Proposition~\ref{prop:incidence-v151} concerns the incidence functor
and the base change of its representing morphism.
\end{remark}

We next retain the nilpotents of a fibre.  This will distinguish a
ramified divisor choice from several distinct divisor choices.

\begin{lemma}[Fibre equations and the overlap space]
\label{lem:fibre-equations-v151}
Let \(K=fJ\), set \(H=\gcd(J)\), and put \(c=\gcd(f,H)\) and
\(k=\deg c\).  Then the tangent space of the scheme-theoretic fibre
at \(([f],J)\) is naturally, after the displayed representatives have
been chosen,
\begin{equation}\label{eq:overlap-kernel-v151}
 \ker(d\nu_g)_{([f],J)}\simeq S_k/\C c,
                     \qquad \dim\ker(d\nu_g)=s_k-1.
\end{equation}
More explicitly, choose a projective coefficient chart
\(f(u)=f+\sum u_\alpha m_\alpha\), with one coefficient normalized
to its value in \(f\).  Let \(M(f(u))\) be the matrix of
multiplication \(S_h\to S_D\).  Choose a set \(A_0\) of \(s_h\)
rows for which \(M_{A_0}(f)\) is invertible, and let \(B_0\) denote
the remaining rows.  For a matrix \(C\) whose columns form a basis
of \(K\), the full fibre on this chart is defined by
\begin{equation}\label{eq:fibre-schur-v151}
 C_{B_0}-M_{B_0}(f(u))M_{A_0}(f(u))^{-1}C_{A_0}=0,
             \qquad \det M_{A_0}(f(u))\ne0.
\end{equation}
In particular its completed local ring is the power-series ring in
the \(s_g-1\) variables \(u_\alpha\), modulo the entries in
\eqref{eq:fibre-schur-v151}; no radical is taken.
\end{lemma}
\begin{proof}
For each column of \(C\), its possible quotient by \(f(u)\) is forced
by the \(A_0\) rows to be \(M_{A_0}^{-1}C_{A_0}\).  The remaining
rows are exactly the condition that this quotient works.  These
quotients span a subbundle wherever the equations hold, by the
incidence argument above.  Thus elimination has introduced neither
additional solutions nor discarded nilpotents, proving the fibre
presentation over arbitrary coefficient rings in the chart.

A tangent vector is a pair \((\dot f\bmod\C f,\dot J)\).  It lies
in the fibre precisely when
\[
              \dot f\,j+f\dot j\in fJ\quad\hbox{for all }j\in J.
\]
It is necessary and sufficient that \(f\mid\dot f\,j\) for every
\(j\in J\); the class of \(\dot j\) modulo \(J\) is then forced.
Taking the minimum valuation over a basis of \(J\), for each
irreducible polynomial, makes this equivalent to \(f\mid\dot f H\).
Writing \(f=cf_1\), \(H=cH_1\), with \(\gcd(f_1,H_1)=1\), gives
\(\dot f=f_1w\) with \(w\in S_k\).  The scalar tangent direction
\(\dot f\in\C f\) is exactly \(w\in\C c\).  This proves
\eqref{eq:overlap-kernel-v151} and its dimension assertion.
\end{proof}

\begin{theorem}[Normalization and the exact normal locus]
\label{thm:boundary-normalization-v151}
The finite map
\[
 \nu_g:\Pj(S_g)\times\Gr(r,S_h)\longrightarrow Y_g
\]
is the normalization of the natural degree-\(g\) divisor locus.
Its dimension is
\begin{equation}\label{eq:boundary-dimension-v151}
                    \dim Y_g=s_g-1+r(s_h-r).
\end{equation}
For a geometric point \(K\), the scheme \(Y_g\) is normal at \(K\)
if and only if \(F_K\) has exactly one degree-\(g\) divisor \([f]\)
and
\begin{equation}\label{eq:normal-criterion-v151}
                         \gcd(f,F_K/f)=1.
\end{equation}
Normality at a point is also equivalent to smoothness there.
Thus the two precise obstructions are more than one divisor choice,
or overlap between the selected divisor and the residual common
factor.  The latter gives the fibre tangent dimension in
\eqref{eq:overlap-kernel-v151}, and
\eqref{eq:fibre-schur-v151} determines its higher nilpotents.
\end{theorem}
\begin{proof}
The locus of gcd-free \(r\)-planes in \(S_h\) is nonempty: a space
containing \(x_1^h,x_2^h\) is one example, and it can be extended
to every allowed rank \(r\).  It is open, since the loci possessing
a positive-degree divisor are the finitely many closed images of
projective multiplication maps.  It is therefore dense in the
irreducible Grassmannian.  Its product with \(\Pj(S_g)\) maps to
the locus where the exact gcd has degree \(g\).

This last locus is open in \(Y_g\): remove the finitely many closed
images for divisor degrees \(a>g\).  Call it \(U_g\).  Every fibre
over a geometric point of \(U_g\) has one point and, by
Lemma~\ref{lem:fibre-equations-v151} with \(H=1\), zero tangent
space.  A finite local algebra over an algebraically closed field
with zero cotangent space is the field itself, by Nakayama's lemma.
The fibre consequently has length one, not merely one point.

Here is the local argument which also settles the scheme structure.
At a closed point of \(U_g\), write \(A\) for the local ring of
\(Y_g\), and \(B\) for its finite inverse-image algebra.  The map
\(A\to B\) is injective, and \(B/\mathfrak m_A B=\C\), generated
by \(1\).  Nakayama's lemma, applied to the finite module \(B/A\),
gives \(B=A\).  The non-isomorphism locus of a finite morphism is
closed, so the same conclusion holds throughout \(U_g\).  Hence
\(\nu_g\) is birational.  The source is smooth and integral, thus
normal.  On affine opens, a finite birational extension \(A\subset B\)
with \(B\) integrally closed is the integral closure of \(A\) in
the common fraction field: every element integral over \(A\) is
integral over \(B\) and hence belongs to \(B\).  This proves the
normalization assertion; compare \cite[Tag~035E]{StacksNormalV151}.
The dimension follows from the product description.

If \(Y_g\) is normal at \(K\), its normalization is an isomorphism
there, so its fibre is one reduced point.  Formula
\eqref{eq:divisor-fibre-set-v151} gives uniqueness of \([f]\), and
\eqref{eq:overlap-kernel-v151} gives \(s_k-1=0\), equivalent to
\(k=0\) because \(e\geq2\).  Conversely uniqueness and \(k=0\)
again make the entire fibre a single reduced point.  The same finite
local Nakayama argument shows that the normalization is an isomorphism
near \(K\).  Its source is normal, so \(Y_g\) is normal at \(K\).
The source is smooth, so every normal point of the target is smooth;
the converse holds for every complex variety.  This proves both
directions without replacing a nonreduced fibre by its set of points.
\end{proof}

\begin{corollary}[Universality of the exact stratum]
\label{cor:canonical-exact-v151}
The open subscheme \(U_g\subset Y_g\) represents the unmarked
relation families in the divisor incidence whose residual space is
gcd-free on every geometric fibre.  Over it the divisor line and
residual subbundle are unique and commute with arbitrary base change.
Moreover \(U_g\), as a locally closed subscheme of \(\mathscr G\),
is the reduced exact-gcd stratum of Lemma~\ref{lem:universal-gcd}.
\end{corollary}
\begin{proof}
The inverse image of \(U_g\) is
\(\Pj(S_g)\times X_{h,r}(E)\), and the proof above shows that its
finite map to \(U_g\) is an isomorphism.  Its source is reduced,
so \(U_g\) has the reduced induced structure on its locally closed
support.  The incidence functor supplies the asserted universal
property.  A family in the ambient Grassmannian which only has the
right geometric gcd need not factor through \(Y_g\).  In particular,
the transverse dual-number family of
Remark~\ref{rem:gcd-scheme-condition} is not admitted by this
corollary.  The corollary is a scheme-theoretic identification, not a
claim that fibrewise gcd tests detect nilpotent equations.
\end{proof}

\begin{example}[A nonreduced fibre at a divisor collision]
\label{ex:fat-divisor-fibre-v151}
Take \(E=\langle x,y,z\rangle\), \(g=1,h=2,r=2\), and
\(K=\langle x^2y,x^2z\rangle\).  Its gcd is \(x^2\), so the only
point of the normalization fibre has divisor \(x\).  In the chart
\(f=x+ay+bz\), division by this monic linear form amounts to
substitution \(x=-ay-bz\).  The two remainders are
\[
             (ay+bz)^2y,\qquad (ay+bz)^2z.
\]
Their coefficient ideal is \((a^2,ab,b^2)\).  Thus the full fibre is
\begin{equation}\label{eq:fat-fibre-v151}
                     \Spec\C[a,b]/(a,b)^2,
\end{equation}
of length three and tangent dimension two.  It is not a reduced
point, despite uniqueness of the geometric divisor.  The family
\[
 K_t=x\langle xy+t y^2,\ xz+t z^2\rangle
\]
has independent generators for every \(t\), exact gcd \(x\) when
\(t\ne0\), and the displayed collision at \(t=0\).  Consequently
this nonreduced fibre occurs on the boundary of the exact stratum,
not in an unrelated example.
\end{example}

\begin{example}[Several branches and a normal gcd jump]
\label{ex:branches-normal-jump-v151}
For the same \(g,h,r,E\), let
\(K=xy\langle x,z\rangle\).  The divisors \(x\) and \(y\)
give two points of the normalization fibre.  At both points the
overlap is one, so the fibre consists of two reduced points.
The point of \(Y_1\) is nonnormal.  By contrast, put
\(q=x^2+yz\) and \(K=xq\langle y,z\rangle\), now with \(h=3\).
The polynomial \(q\) is irreducible and is not divisible by \(x\).
There is exactly one linear divisor of \(\gcd K=xq\), and its
overlap with the residual gcd \(q\) is one.  This point of \(Y_1\)
is normal although its gcd degree is three.  These two examples
show why neither gcd degree nor the number of geometric lifts alone
is a sufficient normality criterion.
\end{example}

\begin{corollary}[An intrinsic boundary for first-relation algebras]
\label{cor:algebra-boundary-v151}
On the maximal-lower-Hilbert-function stratum of
Proposition~\ref{prop:relative-first-relation}, the algebras whose
intrinsic degree-\(D\) kernel has classifying map in \(Y_g\)
have, after removing only tangent-identity substitutions, stack
\([Y_g/\operatorname{GL}(E)]\).  Its normalization is the finite
representable morphism
\begin{equation}\label{eq:quotient-normalization-v151}
 [\mathscr X_g/\operatorname{GL}(E)]
                    \longrightarrow[Y_g/\operatorname{GL}(E)].
\end{equation}
All members of the universal relation family give finite locally
free algebras of rank \(\binom{e+D}{D}-r\).  This construction imposes
no pencil image equations and permits the divisor and residual
common factor to collide.
\end{corollary}
\begin{proof}
For every subbundle \(\mathcal K\) classified by the Grassmannian,
\(\Sym E/((\mathcal K)+\mathfrak m^{D+1})\) is, as a module,
the sum of \(S_j\) for \(j<D\) and \(S_D/\mathcal K\).
This proves finite local freeness and the rank.  The first-relation
proposition identifies the quotient after the specified nonlinear
inertia is removed.  Incidence and multiplication are
\(\operatorname{GL}(E)\)-equivariant.  Pulling
\eqref{eq:quotient-normalization-v151} back along the smooth frame
atlas \(Y_g\to[Y_g/\operatorname{GL}(E)]\) gives precisely
\(\nu_g\).  Finiteness, representability, and the normalization
identification follow by descent; equivalently normality and integral
closure may be checked on these smooth atlases.  This quotient is
not claimed to classify smoothings outside the specified Hilbert
stratum.
\end{proof}

For the pencil first relations one can take
\(e=n^2\), \(g=n(n-1)\), \(h=3n-4\), and
\(r=\binom{\binom{n+1}{2}}2\).  Every pencil then lies in the
exact-gcd open of this canonical boundary.  The boundary theorem
concerns the much larger space of all degree-\(g\) divisorial first
relations, rather than a closure defined by recovering a pencil and
reimposing its Pluecker equations.  It determines its normalization
and normal locus, but does not assert a classification of all its
linear orbits or of all finite-flat algebra deformations.

% BEGIN NEW MATHEMATICS v153
\section{Formal images and marked divisor germs}
\label{sec:formal-images-v153}
All rings and schemes in the following sections are over $\C$. Write
$S_j=\Sym^j E$, and let $Y_g\subset\Gr(r,S_D)$ be the
scheme-theoretic image of
\[
 \nu_g:\Pj(S_g)\times\Gr(r,S_{D-g})\longrightarrow\Gr(r,S_D),
 \qquad (f,J)\longmapsto fJ,
\]
where $\dim E\geq2$, $g\geq1$, and $2\leq r\leq\dim S_{D-g}$.
The source is smooth, and $\nu_g$ is its finite normalization: a fibre
records the degree-$g$ divisors of the common divisor, and on the
exact-degree-$g$ locus there is one reduced lift. We shall retain the
scheme structures of these fibres and of the image.

\begin{lemma}[Completed finite images]\label{lem:formal-image-v153}
Let $X\to M$ be a finite morphism of complex finite-type schemes, with
$X$ reduced and $M$ smooth, and let $Y$ be its scheme-theoretic image.
Fix a closed point $y\in Y$, put $R=\widehat{\OO}_{M,y}$, and write
$B_i=\widehat{\OO}_{X,x_i}$ for the points over $y$. Then
\[
 \widehat{\OO}_{Y,y}=R/\bigcap_i\ker(R\to B_i).
\]
This completed image is reduced. If the cotangent map from $R$ to
$B_i$ is surjective, then $R\to B_i$ is surjective. If there are two
such points, writing $I_i=\ker(R\to B_i)$ gives
\[
 \widehat{\OO}_{Y,y}=(R/I_1)\mathbin{\times}_{R/(I_1+I_2)}(R/I_2).
\]
The formal schemes at the different $x_i$ are disjoint. A local Artin
point of the finite inverse image belongs to exactly one of them.
\end{lemma}
\begin{proof}
On an affine neighbourhood the image ring is the image of the
ambient ring in a finite algebra. Exactness of completion on finite
modules preserves this kernel. The completed finite algebra is the
product of the $B_i$: the orthogonal idempotents of its Artin closed
fibre lift uniquely through the powers of the maximal ideal. Reduced
finite-type complex rings are excellent and analytically unramified,
so their completions are reduced; alternatively this applies directly
to the reduced image ring. For cotangent surjectivity, lift a basis
of the maximal ideal of $B_i$ modulo its square to $R$. The resulting
lifts generate every successive associated-graded piece. Subtracting
successive homogeneous lifts and taking their convergent sum expresses
any element of $B_i$ as the image of an element of $R$. The finite map
factors through the completed image, so this is also a closed immersion
into that image, not only into $M$. Finally the elementary exact sequence
\[
0\longrightarrow R/(I_1\cap I_2)\longrightarrow R/I_1\oplus R/I_2
 \longrightarrow R/(I_1+I_2)\longrightarrow0
\]
proves the fibre-product assertion. The idempotent decomposition gives
the last assertion, including over nonreduced Artin bases.
\end{proof}

\begin{lemma}[Coprime lifts and the split codimension]
\label{lem:coprime-lifts-v153}
If homogeneous forms $a,b$ have coprime reductions over a local Artin
base $A$, then $(a)\cap(b)=(ab)$ in $A[E]$. The relevant graded
quotients and Koszul kernels are $A$-flat. In the split-divisor model,
with $s_j=\binom{e+j-1}{j}$, $e\geq2$, $h>g\geq1$, and $r\geq2$,
\[
 c=r(s_h-s_{h-g})-(s_g-1)\geq(r-1)(s_g-1)\geq1.
\]
\end{lemma}
\begin{proof}
Over the residue field the two forms form a regular sequence. In
each degree their Koszul complex consists of finite free modules and
has its prescribed ranks. The acyclic-complex criterion over an Artin
local ring, proved by lifting splittings successively along its maximal
ideal, gives exactness and flatness of the lifted cokernels. Thus $b$
is a nonzerodivisor modulo $a$, proving the intersection identity.
For the inequality, express $s_h-s_{h-g}$ as the sum of $g$ consecutive
first differences of the sequence $s_j$. These differences are
nondecreasing for $e\geq2$. Their sum is at least $s_g-s_0=s_g-1$.
\end{proof}

In particular the two formal marked lifts in a coprime split model
are separated before their images are compared. Their intersection
functor is precisely $K_A=a_A b_A L_A$. Division gives a subbundle:
multiplication by a fibrewise nonzero form is locally split, and the
quotient of the ambient bundle by $K_A$ is locally free. Combining
this Artin-functor identification with Lemma~\ref{lem:formal-image-v153}
justifies the completed-image fibre product without any assumption
that a tangent-space comparison determines an intersection scheme.

\section{A global factorization theorem for binary normalization fibres}
\label{sec:binary-fibres-v153}
For this section $E$ has dimension two. Set $N_0=r(D+1-r)$ and
$m=r-1$. Let
\[
 \mathcal T_s=Y_s\setminus Y_{s+1},\qquad 1\leq s\leq D+1-r,
\]
with the reduced locally closed structure; an inadmissible $Y_j$ is
empty. Binary forms split into linear factors, so the $Y_j$ are
nested. The multiplication map identifies
\[
 \mathcal T_s\simeq\Pj(S_s)\times X_{D-s,r},
\]
where $X_{D-s,r}$ is the open Grassmannian of gcd-free residual spaces.
The first projection is the algebraic common-divisor morphism, not a
choice made independently on individual fibres.

For nonnegative $a,b$ define $H_{a,b}=\C$ when $ab=0$. Otherwise,
put $k=\min(a,b)$ and $l=\max(a,b)$, and set
\begin{equation}\label{eq:H-ab-v153}
 H_{a,b}=\C[q_1,\ldots,q_k]/(c_{l+1},\ldots,c_{l+k}),\quad
 \sum_{j\geq0}c_jz^j=(1+q_1z+\cdots+q_kz^k)^{-1}.
\end{equation}
The weight of $q_i$ is $i$. This grading is not the maximal-ideal
adic grading when $k>1$.

\begin{theorem}[Stratified finite flatness and all binary fibres]
\label{thm:binary-fibres-v153}
For $g\leq s$, the restriction of $\nu_g$ over $\mathcal T_s$ is the
base change, along the common-divisor morphism, of
\[
 \Pj(S_g)\times\Pj(S_{s-g})\longrightarrow\Pj(S_s),\qquad (f,H)\mapsto fH.
\]
It is finite locally free of degree $\binom{s}{g}$. If
$\gcd(K)=\prod_{i=1}^{b}\ell_i^{s_i}$, with distinct $\ell_i$ and
$\sum_i s_i=s$, its entire normalization fibre is
\begin{equation}\label{eq:binary-fibre-product-v153}
 \prod_{\substack{0\leq a_i\leq s_i\\\sum_i a_i=g}}
                   \bigotimes_{i=1}^{b} H_{a_i,s_i-a_i}.
\end{equation}
Every local factor is an Artin complete intersection and is Gorenstein.
Its length, weighted Hilbert series, and socle weight are, respectively,
\begin{gather*}
 \prod_i\binom{s_i}{a_i},\qquad
 \prod_i\prod_{j=1}^{\min(a_i,s_i-a_i)}
       \frac{1-z^{\max(a_i,s_i-a_i)+j}}{1-z^j},\qquad
 \sum_i a_i(s_i-a_i).
\end{gather*}
The total length is $\binom{s}{g}$, including at arbitrary simultaneous
root collisions. Over fixed $s$, the closure order of common-root
partition strata is the order generated by merging parts.
\end{theorem}
\begin{proof}
Over $\mathcal T_s$ the common divisor $G$ and residual subbundle $L$
are defined scheme-theoretically by the inverse to the exact-divisor
incidence. Locally on any base change, choose a member of $L$ coprime
to $G$ on each relevant geometric fibre. There is such a member over
$\C$, since only finitely many roots must be avoided. Regular-sequence
cancellation, or monic division and its flat quotient, gives
$f\mid GL$ if and only if $f\mid G$. This equivalence is an equality
of incidence functors over nonreduced bases. The residual subspace is
then forced. This proves the base-change description.

The displayed multiplication map is finite by properness and the
finiteness of the divisors of a fixed binary form. Its source and
target are smooth of dimension $s$. A finite map between these smooth
spaces is flat: a regular system of parameters in the target is a
system of parameters in each source local ring, hence a regular
sequence; the local flatness criterion applies. Its degree can be
computed at a squarefree form, where the $\binom{s}{g}$ divisors are
reduced and distinct.

Coprime root clusters split uniquely over Artin bases by Hensel
factorization. On the cluster with polynomial $x^{s_i}$, choosing a
degree-$a_i$ factor is the equation $f_i H_i=x^{s_i}$ with both factors
monic. Eliminating one factor by reciprocal expansion gives
\eqref{eq:H-ab-v153}. This proves the full product decomposition,
not only its closed points. Its equations have their sole common zero
at the origin: a monic factor of $x^{a+b}$ over $\C$ is a power of
$x$. Thus its $k$ equations in $k$ variables are a homogeneous system
of parameters and a regular sequence. The Hilbert series is the
quotient of the corresponding weighted polynomial Hilbert series.
Taking its value at one gives $\binom{a+b}{a}$; its top weight is
$ab$. A zero-dimensional complete intersection has a one-dimensional
socle. Equivalently $q_k^l$, up to a nonzero scalar, represents the
rectangular Schur class. The identification with the usual
Grassmannian cohomology presentation is classical; a purely algebraic
Schur-basis statement is \cite[Theorem~2.7]{Grinberg2019}.
Tensoring the local factors proves all three formulas. The total
length identity is the coefficient of $z^g$ in
$\prod_i(1+z)^{s_i}$. Finally the exact-gcd stratum is the product
with $X_{D-s,r}$ above. In the symmetric power of the projective line,
root clusters specialize by coalescence, and every merging is realized
by moving the corresponding roots together. This gives precisely the
stated closure order within fixed gcd degree.
\end{proof}

\begin{remark}
The finite flat statement is on the exact-gcd stratum with its indicated
scheme structure. It does not assert that normalization or formation of
scheme-theoretic images commutes with an arbitrary nonflat base change
of the ambient Grassmannian. It identifies a specific base change of
the already constructed incidence morphism.
\end{remark}

\section{Conductors and all branch intersections on the squarefree locus}
\label{sec:squarefree-v153}
Let $\mathcal U\subset\Gr(r,S_D)$ be the open set on which the common
divisor is squarefree. Its complement is the proper image of the
incidence $K=\ell^2L$. This condition concerns the common divisor,
not every member of the subspace.

\begin{theorem}[Squarefree conductor stratification]
\label{thm:squarefree-v153}
Let $K\in Y_g\cap\mathcal U$, let $s=\deg\gcd(K)$, and set
\[
 t=s-g,\qquad m=r-1,\qquad q=N_0-ms.
\]
Introduce $s$ disjoint blocks $u_i=(u_{i1},\ldots,u_{im})$, and let
$J_j$ be the ideal generated by products of one variable from each of
$j$ distinct blocks. Put $J_0=(1)$ and $J_{s+1}=0$. Then
\begin{equation}\label{eq:squarefree-ring-v153}
 \widehat{\OO}_{Y_g,K}\simeq
 \C[[z_1,\ldots,z_q,u_{11},\ldots,u_{sm}]]/J_{t+1}.
\end{equation}
The branches are indexed by the subsets $A\subset\{1,\ldots,s\}$
of cardinality $g$. The branch $A$ has ideal
$I_A=(u_{ij}:i\in A)$ in the ambient power-series ring. For any
collection of branches, their full scheme-theoretic intersection has
ideal $I_{\cup A}$; it is smooth of dimension
$q+m(s-|\cup A|)$.

The conductor in \eqref{eq:squarefree-ring-v153} is $J_t/J_{t+1}$.
Its restriction to branch $A$ is
$\prod_{i\notin A}(u_{i1},\ldots,u_{im})$. Consequently the global
conductor sheaf on $Y_g\cap\mathcal U$ is exactly
\begin{equation}\label{eq:global-conductor-v153}
 \mathfrak c_{Y_g}|_{\mathcal U}
        =\II_{Y_{g+1}/Y_g}|_{\mathcal U}.
\end{equation}
In particular it is radical there. These local rings are seminormal.
Their dimension, depth and multiplicity are
\[
 q+mt,\qquad q+t,\qquad \binom{s}{g}.
\]
They are Cohen--Macaulay exactly when $s=g$ or $r=2$.
The Hilbert series of the nonsmooth factor in its standard grading is
\begin{equation}\label{eq:block-hilbert-v153}
 \sum_{j=0}^{t}\binom{s}{j}\bigl((1-z)^{-m}-1\bigr)^j.
\end{equation}
The whole collection of branch intersections and conductor strata
is invariant under permutation of the roots and descends from any
ordered-root etale chart.
\end{theorem}
\begin{proof}
Choose $F_0\in K$ with a simple zero at each of the $s$ common roots,
and complete it to a frame $F_0,\ldots,F_m$. Hensel factorization gives
moving roots $x_1,\ldots,x_s$ in the completion. Use as coordinates
the $s$ root parameters and the $ms$ evaluations
$u_{ij}=F_j(x_i)$, with harmless invertible normalizing factors.
These are part of a regular parameter system: the evaluation map
$S_D/K\to\C^s$ is surjective because $K$ vanishes on the common-root
scheme, and $D\geq s-1$. Thus
$\Hom(K,S_D/K)$ allows independent variation of all $rs$ jets.
The passage to the displayed coordinates is triangular with invertible
diagonal, since $F_0$ has simple zeros. The other $N_0-rs$ parameters
and the $s$ root parameters give exactly $q=N_0-ms$ smooth parameters.

A marked degree-$g$ divisor near this fibre chooses a subset $A$ of
these roots. Its incidence equations are exactly $u_{ij}=0$ for
$i\in A$, with no nilpotents suppressed. There are no other divisor
lifts. Lemma~\ref{lem:formal-image-v153} therefore identifies the
completed image ideal with $\bigcap_{|A|=g}I_A$. A monomial belongs
to this intersection exactly when its block support meets every
$g$-element subset, or equivalently has at least $s-g+1=t+1$ blocks.
This proves \eqref{eq:squarefree-ring-v153}. Sums of the coordinate
ideals give all intersections, including intersections of three or
more branches.

For completeness compute the conductor inside the product of the
branch rings. Its component on branch $A$ consists of functions whose
extension by zero on every other branch belongs to the image ring.
The image of $\bigcap_{B\ne A}I_B$ modulo $I_A$ is precisely
$\prod_{i\notin A}(u_i)$. Indeed a surviving monomial has support
outside $A$, and must meet every $g$-subset other than $A$. Replacing
one element of $A$ by each $i\notin A$ forces all those blocks to
occur. Conversely their product meets every such subset. Hence a
conductor monomial has exactly $t$ surviving blocks; monomials on more
blocks are zero. This is $J_t/J_{t+1}$. The same coordinates identify
$Y_{g+1}$ with $J_t$, proving the sheaf equality by faithful flatness
of completion. The case $t=0$ means that the conductor is the unit
ideal and the next stratum is absent locally.

The ring is the ring of tuples on its coordinate branches that agree
on every pairwise intersection. To verify this directly, compare the
coefficient of each monomial on all branches where it survives; the
pairwise agreements give a unique common coefficient. If an element
$b$ of the total fraction ring has $b^2,b^3$ in the image, it first
belongs to the normalization, by integrality. Its two restrictions
on any intersection have equal squares and cubes, hence are equal
in that reduced ring. Thus $b$ lies in the image. This proves
seminormality by the square-and-cube criterion.

Formula~\eqref{eq:block-hilbert-v153} follows by sorting nonzero
monomials by their block support. The top-dimensional contributions
come from the $\binom{s}{t}$ supports of size $t$, giving dimension
$mt$ and multiplicity $\binom{s}{t}$. We supply the homological step
for the depth, since branch counting alone does not give it. The
simplicial complex of the squarefree monomial ideal $J_{t+1}$ consists
of sets using at most $t$ blocks. Any induced subcomplex meeting $b$
blocks is homotopy equivalent to the $(t-1)$-skeleton of the simplex
on those $b$ blocks: collapse each block to a representative vertex.
The two composites are contiguous to the identity, because adjoining
the representatives never increases the number of blocks. If $b\leq t$
the complex is a simplex; otherwise its only reduced homology is in
degree $t-1$, with rank $\binom{b-1}{t}$. This last assertion follows
also by the elementary induction that adds the last vertex to a
simplex skeleton.

The squarefree multidegree-$W$ part of the Koszul complex on all $sm$
variables is the augmented simplicial cochain complex of this induced
subcomplex, with degree shifted by $|W|-1$. Non-squarefree multidegrees
are contractible in the relevant summands, by pairing faces using a
variable whose exponent exceeds one. Consequently a nonzero positive
Betti number has homological degree $|W|-t$. The full vertex set gives
the nonzero degree $sm-t$, since $s>t$ when $g>0$. The projective
dimension is therefore $sm-t$, and the Auslander--Buchsbaum formula
gives depth $t$. When $t=0$ the quotient is simply $\C$. Adjoining
$q$ power-series parameters gives depth $q+t$ in every case, and
comparison with $q+mt$ proves the Cohen--Macaulay criterion.
All the constructions commute with relabelling roots, proving descent.
\end{proof}

\begin{corollary}[Binary divisor-image intersections]
\label{cor:binary-intersections-v153}
For admissible binary divisor degrees $a_1,\ldots,a_v$, their
scheme-theoretic intersection is $Y_{\max a_i}$. On the squarefree
locus all multiple intersections of normalization branches, rather
than only the degree images, are given by the union-of-subsets rule
in Theorem~\ref{thm:squarefree-v153}.
\end{corollary}
\begin{proof}
The reduced binary images are nested as closed subschemes: a binary
form of degree $b$ has a divisor of each degree $a\leq b$, and
reduced closed subschemes are determined by their closed points.
Their ideals are thus nested, so their sum is the largest ideal.
The assertion about branches has already been proved as an identity
of coordinate ideals, without taking radicals.
\end{proof}

\section{An incidence atlas for all binary root collisions}
\label{sec:collision-atlas-v153}
The squarefree theorem has a uniform replacement when any number of
root clusters collide. This replacement specifies the full completed
image, not merely a transverse family in it.

\begin{theorem}[Simultaneous collision atlas]\label{thm:atlas-v153}
Let $K\in Y_g$ be binary with exact common divisor
$\prod_{i=1}^b\ell_i^{s_i}$ and $s=\sum_i s_i$. There are formal
ambient coordinates consisting of the coefficients of monic
polynomials $P_i$ of degree $s_i$, polynomials $R_{ij}$ of degree
less than $s_i$ ($1\leq j\leq m$), and $N_0-rs$ smooth parameters
$z$. They are centred at $P_i=x_i^{s_i}$ and $R_{ij}=0$.
For every $\alpha=(a_i)$ with $0\leq a_i\leq s_i$ and $\sum a_i=g$,
let $B_\alpha$ be the regular power-series ring in the coefficients
of monic $A_i,B_i'$ of degrees $a_i,s_i-a_i$, polynomials $U_{ij}$
of degree less than $s_i-a_i$, and the same $z$. The completed
normalization and image are exactly
\begin{equation}\label{eq:collision-image-v153}
 T\longrightarrow\prod_\alpha B_\alpha,\qquad
 P_i\mapsto A_iB_i',\quad R_{ij}\mapsto A_iU_{ij},\qquad
 \widehat{\OO}_{Y_g,K}=T/\bigcap_\alpha\ker(T\to B_\alpha).
\end{equation}
Here $T=\C[[\operatorname{coeff}(P_i),\operatorname{coeff}(R_{ij}),z]]$.
Degree-zero monic factors are one, and negative-degree remainder
spaces are zero. The formula is independent, up to isomorphism of
completed incidence diagrams, of the frame, the selected $F_0$, and
the local root coordinates. It covers every higher and simultaneous
binary collision.
\end{theorem}
\begin{proof}
Choose $F_0\in K$ having exact order $s_i$ at every common root.
The required nonvanishing of its residual value excludes finitely
many hyperplanes in $K$, so a choice exists. Complete it to a frame.
Weierstrass preparation at the distinct roots produces $P_i$; division
of the remaining frame members produces $R_{ij}$. The map
\[
 S_D/K\longrightarrow\bigoplus_i\C[x_i]/(x_i^{s_i})
\]
is surjective: its kernel contains $K$, and the usual interpolation
map from $S_D$ is surjective. Thus the Grassmannian tangent space
permits independent variation of all $rs$ jets. Division and preparation
change these jets by an invertible triangular matrix, whose diagonal
contains the nonzero residual values of $F_0$. Their differentials
are independent. Formal inverse coordinates supply the remaining
$N_0-rs$ parameters; this integer is nonnegative since
$r\leq D-s+1$.

A marked divisor near a chosen normalization point splits uniquely
among the coprime root clusters. In cluster $i$ its monic representative
has degree $a_i$ and divides both $P_i$ and all $R_{ij}$. Monic division
is equivalent, over every local Artin base, to the unique equations
$P_i=A_iB_i'$ and $R_{ij}=A_iU_{ij}$. Conversely these equations give
the marked divisor and the residual subbundle. Thus the normalized
incidence functor is represented by the displayed regular ring.
It is finite over $T$: monic factor coefficients are integral over the
coefficients of $P_i$, while the remaining factors and quotients are
uniquely determined by division. Lemma~\ref{lem:formal-image-v153}
now gives the exact completed image and the intersection of kernels.

Changing a frame or a coordinate changes only a presentation of the
same functor of a subbundle with a marked divisor. The mutually inverse
changes of incidence data therefore induce isomorphisms of its complete
representing rings and of the resulting finite-image diagrams. This
also proves that the calculation describes the entire completed image.
A point of that image over an Artin ring need not lift to its
normalization; no such lifting assertion is used or implied.
\end{proof}

\section{The singular Fitting ideal of every double-root model}
\label{sec:fitting-v153}
The atlas specializes for $s=2$, $g=1$ to the following ring, after
completion of the square and suppression of smooth parameters:
\[
 A_m=R\oplus It\subset B=\C[[u_1,\ldots,u_m,t]],\qquad
 R=\C[[u_1,\ldots,u_m,\Delta]],\ I=(u_1,\ldots,u_m),\ \Delta=t^2.
\]
Put $w_i=u_it$. The defining equations are
$u_iw_j-u_jw_i$ and $w_iw_j-\Delta u_iu_j$.

\begin{theorem}[Closed singular-ideal formula]\label{thm:fitting-v153}
For every $m\geq1$, the intrinsic singular Fitting ideal is
\begin{equation}\label{eq:fitting-v153}
 \Fitt_{m+1}\Omega^{\mathrm{cont}}_{A_m/\C}
  =(I^{m+1}+\Delta I^m)\oplus I^m t
  = I^{m-1}(I^2,(w_1,\ldots,w_m),\Delta I).
\end{equation}
The first equality is a decomposition as an $R$-submodule of $B$;
the right side is an ideal of $A_m$. It is not merely an ideal of
Jacobian minors left unevaluated. For $m=1$ it is
$(u^2,w,\Delta u)$, and its radical is the conductor $(u,w)$ for
all $m$.
\end{theorem}
\begin{proof}
There are $2m+1$ ambient variables and dimension $m+1$, so the Fitting
ideal consists of the $m$-minors of the Jacobian of the displayed
relations. Evaluate this matrix in $B$ using $w_i=u_it$ and
$\Delta=t^2$. A minor not using the $\Delta$ column has $u$-degree
$m$. It vanishes at $t=0$: the only remaining rows are the Koszul
rows in the $w$ columns, of rank at most $m-1$. Hence its terms have
a positive power of $t$, and belong to $\Delta I^m\oplus I^m t$.
A minor using the $\Delta$ column has $u$-degree $m+1$, and hence
lies in $I^{m+1}\oplus I^{m+1}t$. This proves the first inclusion.

Use the $m$ rows
\[
 u_1w_j-u_jw_1\ (2\leq j\leq m),\qquad w_1^2-\Delta u_1^2.
\]
The columns $(w_1,w_2,\ldots,w_m)$, $(u_1,w_2,\ldots,w_m)$,
and $(\Delta,w_2,\ldots,w_m)$ give, up to sign and nonzero scalar,
$u_1^m t$, $\Delta u_1^m$, and $u_1^{m+1}$. The ideal is invariant
under simultaneous linear change of $u$ and $w$. In characteristic
zero the powers of linear forms span each symmetric power, so these
three minors generate the three indicated spaces of all $u$-monomials.
They generate the claimed $R$-module; it is an $A_m$-ideal since
multiplication by $w_i=u_it$ preserves it. This proves equality.
The formula also immediately gives its radical, without discarding
the nonreduced Fitting structure in the assertion itself.
\end{proof}

\section{A structural dichotomy for pure-power divisor fibres}
\label{sec:reciprocal-v153}
Let $E=\C a\oplus W$, $d=\dim W\geq1$, and $g,k\geq1$.
The factorization fibre of $a^{g+k}$ records monic factors of degrees
$g$ and $k$. Interchanging these factors identifies its two descriptions;
we may and shall take $g\geq k$. Write
\[
 Q_i\in\Sym^iW\quad(1\leq i\leq k),\qquad
 C_0=1,\quad C_j=-\sum_{i=1}^kQ_i C_{j-i},\quad C_j=0\ (j<0).
\]
Let $P=\Sym(\bigoplus_{i=1}^k(\Sym^iW)^*)$, with its generators
of weight $i$, and let $\mathfrak n$ be its variable ideal. Put
\begin{equation}\label{eq:reciprocal-algebra-v153}
 F_{g,k}(W)=P/\mathcal I,\qquad
 \mathcal I=(\operatorname{coeff}C_{g+1},\ldots,
                                    \operatorname{coeff}C_{g+k}).
\end{equation}
It is an Artin local algebra. Indeed its complex points are exactly
the factorizations of $a^{g+k}$, hence the single origin by unique
factorization. The polynomial and completed presentations coincide.
If $K=a^{g+k}L$ with $\gcd(L)=1$, cancellation by a member of $L$
not divisible by $a$, as in Lemma~\ref{lem:coprime-lifts-v153},
identifies this with the full degree-$g$ normalization fibre at $K$.

\begin{theorem}[Minimal equations and complete-intersection classification]
\label{thm:reciprocal-structure-v153}
The embedding dimension and minimal number of equations of
$F_{g,k}(W)$ in a minimal regular local presentation are
\begin{align}
 \operatorname{edim}F_{g,k}(W)&=\binom{d+k}{k}-1,\label{eq:edim-v153}\\
 \mu(\mathcal I)&=\sum_{j=g+1}^{g+k}\binom{d+j-1}{j}
                  =\binom{d+g+k}{d}-\binom{d+g}{d}.
                  \label{eq:minrels-v153}
\end{align}
More precisely the minimal relation module has the equivariant
weighted decomposition
\begin{equation}\label{eq:relation-module-v153}
 \mathcal I/\mathfrak n\mathcal I
                 \simeq\bigoplus_{j=g+1}^{g+k}(\Sym^jW)^*.
\end{equation}
Consequently this fibre is a complete intersection if and only if
$d=1$. For $d=1$, it has length $\binom{g+k}{k}$, weighted Hilbert
series $\prod_{i=1}^k(1-z^{g+i})/(1-z^i)$, and one-dimensional
socle generated by $Q_k^g$, of weight $gk$.
For $k=1$ and arbitrary $d$,
\begin{equation}\label{eq:k-one-v153}
 F_{g,1}(W)=\Sym(W^*)/\mathfrak n^{g+1},\quad
 \length F_{g,1}(W)=\binom{d+g}{g},\quad
 \operatorname{Soc}F_{g,1}(W)=\Sym^g(W^*).
\end{equation}
In particular the last algebra is Gorenstein precisely when $d=1$.
No Gorenstein classification for $d>1,k>1$ is needed for the theorem.
\end{theorem}
\begin{proof}
Since $g\geq k$, each defining form has no linear term in the
coefficient variables. Thus the presentation is minimal and
\eqref{eq:edim-v153} follows by summing $\dim\Sym^iW$.
For $d=1$, the regular-sequence argument in the proof of
Theorem~\ref{thm:binary-fibres-v153} applies to the $k$ reciprocal
relations. In particular none of them belongs to $\mathfrak n\mathcal I$.
It gives the stated Hilbert series, length, and the classical rectangular
socle class.

For general $W$, taking coefficients of $C_j$ gives an equivariant
map $(\Sym^jW)^*\to\mathcal I/\mathfrak n\mathcal I$.
Its source is irreducible. This map is nonzero: choose one coordinate
$x$ of $W$, keep only the coefficients of $x^i$ in every $Q_i$, and
set every other coefficient to zero. The coefficient of $x^j$ maps
to the $j$th binary reciprocal relation. Were it in
$\mathfrak n\mathcal I$ before specialization, its specialization
would be in the binary $\mathfrak n\mathcal I$, contradicting the
preceding minimality. Irreducibility now makes the entire coefficient
map injective. The different $j$ have different weights, so their
images are independent. They generate the relation module by the
chosen presentation, proving \eqref{eq:relation-module-v153} and
\eqref{eq:minrels-v153}.

If $d\geq2$, the sequence $\binom{d+j-1}{j}$ strictly increases in
$j$. Comparing its terms at $j=g+i$ and $j=i$ shows that the minimal
number of relations is strictly greater than the embedding dimension.
A zero-dimensional complete intersection in a minimal regular local
presentation has exactly that many minimal relations as variables;
hence these fibres are not complete intersections. This proves both
directions of the classification. Finally when $k=1$ the sole form
is $C_{g+1}=(-Q_1)^{g+1}$. Its coefficients are, with nonzero multinomial
factors, every monomial of degree $g+1$ in the $d$ variables. This gives
\eqref{eq:k-one-v153}, including the socle and length.
\end{proof}

\begin{corollary}[The ternary-pencil ramification fibre]
\label{cor:ternary-fibre-v153}
At the common limiting relation algebra with $e=9$, $g=6$ and $k=2$,
the normalization fibre has embedding dimension $44$ and exactly
\[
 \binom{14}{7}+\binom{15}{8}=9867
\]
minimal scalar equations. Their two weighted pieces have dimensions
$3432$ and $6435$, in weights seven and eight. In particular this
actual pencil-orbit boundary fibre is not a complete intersection.
\end{corollary}
\begin{proof}
Here $d=e-1=8$. Apply \eqref{eq:edim-v153} and
\eqref{eq:relation-module-v153}. These are minimal equation counts,
not the number of coefficient variables or the length of the fibre.
\end{proof}

The binary factorization rings in this section are classical. The
additional point of the multivariate argument is that polarization
preserves every full irreducible coefficient space in the minimal
relation module. Neither the Grassmannian identification nor a finite
calculation in a few variables is being used as a claim of priority
or as a substitute for that argument.

\section{Homological invariants of the divisor fibres}
\label{sec:homological-v154}
We first extract all graded Betti numbers from the squarefree local
model and then identify a determinant-apolar quotient of the
multivariate fibres. A separate exact calculation treats an entire
multivariate fibre, rather than only a quotient or its minimal
relations. The grading conventions will be specified in each case.

\begin{theorem}[Betti numbers of the squarefree block model]
\label{thm:block-betti-v154}
Let $S=\C[u_{ij}:1\leq i\leq s,1\leq j\leq m]$ with its standard
multigrading, let $0\leq t<s$, and put $A=S/J_{t+1}$, with the
block-support ideal of Theorem~\ref{thm:squarefree-v153}.
For a nonempty squarefree multidegree $F$ let $b(F)$ be the number
of blocks it meets. Besides $\beta_{0,0}=1$, the only nonzero
multigraded Betti numbers are
\begin{equation}\label{eq:block-multibetti-v154}
 \beta_{i,F}^S(A)=\binom{b(F)-1}{t}
       \quad\text{if }b(F)>t\text{ and }i=|F|-t.
\end{equation}
For $t=0$ this is the Koszul resolution of $\C$.
For all $t$, the singly graded numbers are
\begin{equation}\label{eq:block-betti-v154}
 \beta_{i,i+t}^S(A)=
 \sum_{b=t+1}^{\min(s,i+t)}\binom{s}{b}\binom{b-1}{t}
       [z^{i+t}]\big((1+z)^m-1\big)^b\quad(i\geq1).
\end{equation}
The ideal $J_{t+1}$ has a $(t+1)$-linear resolution. The projective
dimension of $A$ is $sm-t$. In the Cohen--Macaulay case $m=1$,
its type is $\binom{s-1}{t}$; it is Gorenstein exactly when
$t=0$ or $t=s-1$. The same Betti numbers over the completed ambient
regular ring apply to the local model after adjoining its smooth
parameters.
\end{theorem}
\begin{proof}
The Stanley--Reisner complex consists of sets meeting at most $t$
blocks. If $t\geq1$, its induced subcomplex on $F$ has, in each
nonempty block, a simplex all of whose vertices have identical
external links. Contract that simplex to one of its vertices: the
map and the inclusion are contiguous, because adjoining the chosen
vertex to any face does not change its block support. Applying this
one block at a time gives a homotopy equivalence with the
$(t-1)$-skeleton of the simplex on $b(F)$ vertices. If $b(F)\leq t$
this simplex is contractible. Otherwise its reduced homology is
concentrated in degree $t-1$, with dimension $\binom{b(F)-1}{t}$.
The latter assertion follows by attaching its top faces successively,
or from the augmented simplicial chain complex and the binomial
identity for its Euler characteristic. Hochster's multigraded
formula now gives \eqref{eq:block-multibetti-v154}; see
\cite[Chapter~5]{MillerSturmfels2005}. If $t=0$, every variable is
in the ideal and the Koszul formula gives the same result directly.

For a support of size $b$, the coefficient in
\eqref{eq:block-betti-v154} counts choices of $i+t$ vertices meeting
each selected block. Summing proves the formula. All ideal syzygies
therefore have the indicated linear degrees. Taking $F$ to be all
$sm$ vertices gives the nonzero last Betti number
$\binom{s-1}{t}$ in homological degree $sm-t$; no larger degree
occurs. When $m=1$, Auslander--Buchsbaum gives depth $t$, equal to
dimension, and the last Betti number is the Cohen--Macaulay type.
Completion and adjoining independent regular parameters are flat
operations preserving the minimal free resolution. These are
classical squarefree-monomial methods. The assertion here is their
full specialization to the normalization image model, not a claim
of priority for Hochster's formula or skeleton resolutions.
\end{proof}

For the remaining results use the weighted reciprocal algebra
$F_{g,k}(W)$ of \eqref{eq:reciprocal-algebra-v153}, where
$d=\dim W$ and $g\geq k\geq2$. Set all $Q_i$ except $Q_2$ equal
to zero and, in addition, impose every coefficient of $Q_2^2$.
Denote the resulting algebra by
\begin{equation}\label{eq:apolar-quotient-v154}
 B_d=\Sym((\Sym^2W)^*)/(\operatorname{coeff}Q_2^2).
\end{equation}
Its coefficient variables have reciprocal weight two. When referring
to its ordinary polynomial grading, we instead give them degree one.

\begin{theorem}[A determinant-apolar quotient of every higher fibre]
\label{thm:apolar-quotient-v154}
There is a $\operatorname{GL}(W)$-equivariant graded surjection
$F_{g,k}(W)\twoheadrightarrow B_d$ for every $g\geq k\geq2$.
After fixing the standard coefficient/differential pairing, $B_d$
is the apolar algebra of the determinant of a generic symmetric
$d\times d$ matrix. Its polynomial-degree Hilbert function is
\begin{equation}\label{eq:Narayana-v154}
 \dim(B_d)_p=\frac{1}{d+1}\binom{d+1}{p}\binom{d+1}{p+1}
                 \quad(0\leq p\leq d).
\end{equation}
It is Gorenstein, its reciprocal-weight socle is $2d$, and
\begin{equation}\label{eq:Catalan-v154}
 \length F_{g,k}(W)\ \geq\ \length B_d
       =\frac{1}{d+2}\binom{2d+2}{d+1}.
\end{equation}
For $k=2$ and $g=2$ or $3$, the quotient by the coefficients of
$Q_1$ alone is already $B_d$.
In particular the actual ternary-pencil fibre $F_{6,2}(\C^8)$ has
length at least $4862$, independently of its embedding dimension
$44$ and its $9867$ minimal equations.
\end{theorem}
\begin{proof}
Under the stated substitution, the reciprocal series becomes
$(1+Q_2z^2)^{-1}$. Its odd coefficients vanish and its even
coefficients are powers of $Q_2$. Since $g+1\geq3$, imposing
$Q_2^2=0$ annihilates all defining coefficients. This proves the
surjection, functorially for linear changes of $W$. When $k=2$ and
$g=2$ or $3$, the first even defining coefficient is $Q_2^2$ and
there are no further independent equations after $Q_1=0$.

Choose coordinates $x_1,\ldots,x_d$ on $W$ and write
$Q_2=\sum_i y_{ii}x_i^2+\sum_{i<j}y_{ij}x_ix_j$.
Let $X=(X_{ij})$ be symmetric, with independent entries for $i\leq j$,
and interpret $y_{ij}$ as $\partial/\partial X_{ij}$. Then
$Q_2(x)$ acts as directional differentiation in the rank-one
symmetric direction $xx^t$. The determinant of $X+zxx^t$ is linear
in $z$, so every coefficient of $Q_2(x)^2$ annihilates $\det X$.
These coefficients are, up to nonzero scalars, $y_{ii}^2$,
$y_{ii}y_{ij}$, $y_{ij}^2+2y_{ii}y_{jj}$,
$y_{ii}y_{jl}+y_{ij}y_{il}$ for distinct $i,j,l$, and
$y_{ij}y_{kl}+y_{ik}y_{jl}+y_{il}y_{jk}$ for four distinct indices.
They are exactly the quadratic generators in
\cite[Proposition~3.1 and Theorem~3.10]{ShafieiSymmetric}.
That theorem proves equality with the full apolar ideal, not merely
containment. Its Corollary~2.7 gives \eqref{eq:Narayana-v154} and
the Catalan length. Alternatively, that Hilbert function counts the
independent minors obtained by differentiating $\det X$.
Apolarity to one homogeneous form gives a perfect pairing between
complementary degrees and a one-dimensional socle in degree $d$.
Replacing polynomial degree by twice that degree proves the weight
statement. The length inequality follows from the surjection.
For $d=8$ the Catalan number is $4862$.

The apolar Hilbert function and its Gorenstein property are classical
results of the cited determinant theory. The new identification in
this setting is the functorial quotient of the entire family of
normalization fibres. A quotient's Gorenstein property is not being
attributed to the original fibre.
\end{proof}

\begin{proposition}[The complete first multivariate two-factor fibre]
\label{prop:F22-v154}
Let $A=F_{2,2}(\C^2)$ and write
$Q_1=ax+by$, $Q_2=cx^2+dxy+ey^2$.
Then
\begin{gather}\label{eq:F22-hilbert-v154}
 H_A^{\mathrm{wt}}(z)=1+2z+6z^2+6z^3+7z^4,\qquad
 \length A=22,\qquad\dim\operatorname{Soc}A=7,\\
 \label{eq:F22-local-hilbert-v154}
 H_{\gr_{\mathfrak m}A}(z)=1+5z+6z^2+5z^3+5z^4.
\end{gather}
The weight grading gives $a,b$ weight one and $c,d,e$ weight two.
In the minimal weighted free resolution over
$S=\C[a,b,c,d,e]$, the shift polynomials
$B_i(z)=\sum_j\beta_{i,j}^S(A)z^j$ are
\begin{align*}
 B_0&=1,& B_1&=4z^3+5z^4,\\
 B_2&=4z^5+12z^6+12z^7,&
 B_3&=21z^8+20z^9+z^{10},\\
 B_4&=21z^{10}+8z^{11},&B_5&=7z^{12}.
\end{align*}
In particular the total Betti numbers are $(1,9,28,42,29,7)$.
This fibre is not Gorenstein. The assertion concerns this whole
five-variable fibre, not the determinant-apolar quotient.
\end{proposition}
\begin{proof}
The nine initial equations are the coefficients of
$-Q_1^3+2Q_1Q_2$ and $Q_1^4-3Q_1^2Q_2+Q_2^2$.
For graded reverse lexicographic order $e>d>c>b>a$, a Groebner basis
is the following list (a scalar multiple does not change a member):
\begin{gather*}
 (c,d,e)^3,\quad b(c,d,e)^2,\quad a(c,d,e)^2,\\
 b^2e-e^2,\quad b^2d-de,\quad b^2c-a^2e,\quad b^3-2be,\\
 2abe-de,\quad 2a^2e+2abd-2ce-d^2,\quad 2abc-cd,\\
 3ab^2-2ae-2bd,\quad a^2d-cd,\quad a^2c-c^2,\\
 3a^2b-2ad-2bc,\quad a^3-2ac.
\end{gather*}
Here a monomial ideal in the list means all its minimal monomial
generators. There are $10+6+6+12=34$ polynomials. Polynomial
division gives both equality with the ideal of the nine equations
and zero remainders for all their basis S-pairs. This is a finite
exact arithmetic verification; the explicit input and all reduction
and rank tests are supplied in the accompanying calculation.
The standard monomials, grouped by weight, are
\begin{equation}\label{eq:F22-basis-v154}
\begin{array}{c|l}
0&1\\
1&a,b\\
2&a^2,ab,b^2,c,d,e\\
3&ac,bc,ad,bd,ae,be\\
4&c^2,cd,d^2,a^2e,ce,de,e^2.
\end{array}
\end{equation}
Reduction of products of these monomials gives
$\dim A/\mathfrak m=1$ and
$\dim(\mathfrak m^i)=(21,16,10,5,0)$ for $1\leq i\leq5$.
The common kernel of multiplication by $a,b,c,d,e$ is exactly the
last row of \eqref{eq:F22-basis-v154}. This proves both distinct
Hilbert series and the type.

For precision, the resolution calculation uses no numerical rank
thresholds. In weight $w$ and homological degree $i$, take basis
pairs $(J,m)$ with $J\subset\{a,b,c,d,e\}$ of size $i$, $m$ a
standard monomial, and $\operatorname{wt}(J)+\operatorname{wt}(m)=w$.
The differential is
\begin{equation}\label{eq:F22-koszul-v154}
 d(J,m)=\sum_{v\in J}(-1)^{\operatorname{pos}(v,J)-1}
                  (J\setminus\{v\},\operatorname{NF}(vm)).
\end{equation}
All entries are rational, explicitly determined by the displayed
34 polynomials. Exact row reduction and
$\beta_{i,w}=\dim C_{i,w}-\rank d_{i,w}-\rank d_{i+1,w}$ give the
six shift polynomials above. The calculation is fully specified by
\eqref{eq:F22-basis-v154}--\eqref{eq:F22-koszul-v154}; the archived
rank table makes each finite elimination independently repeatable.
As additional checks, their alternating sum is
$H_A^{\mathrm{wt}}(z)(1-z)^2(1-z^2)^3$ and the last Betti space is
the socle shifted by the total variable weight eight. These checks
do not replace the exact rank computations.
\end{proof}

\section{Quadratic sections of reciprocal normalization fibres}
\label{sec:quadratic-sections-v155}
The reciprocal presentation has a distinguished linear coefficient
section on which the entire algebra, rather than only its minimal
relations, can be determined.  Throughout this section a decomposition
$E=\C a\oplus W$ is fixed.  The section is equivariant for
$\operatorname{GL}(W)$; it is not asserted to be independent of an
arbitrary change of that decomposition.

Put $d=\dim W$, $g\geq k\geq2$, and $h=\lfloor g/2\rfloor$.
In the algebra $F_{g,k}(W)$ of \eqref{eq:reciprocal-algebra-v153},
let $\mathfrak l$ be the ideal generated by the coefficients of all
$Q_i$ with $i\ne2$.  Write $Q=Q_2$ and
\begin{equation}\label{eq:quadratic-section-v155}
 B_{d,h}=\Sym((\Sym^2W)^*)/
                    (\operatorname{coeff}Q^{h+1}).
\end{equation}
Here the ordinary degree of a coefficient of $Q$ is one; its
reciprocal weight is two.  These two gradings will always be
specified separately.  A partition $\lambda\subseteq(h^d)$ has
at most $d$ rows, each of length at most $h$.

\begin{theorem}[The entire quadratic coefficient section]
\label{thm:quadratic-section-v155}
There is an isomorphism of $\operatorname{GL}(W)$-equivariant graded
algebras
\begin{equation}\label{eq:exact-section-v155}
                 F_{g,k}(W)/\mathfrak l\simeq B_{d,h}.
\end{equation}
Under the coefficient/differential pairing, $B_{d,h}$ is the apolar
algebra of the $h$th power of the determinant of a generic symmetric
$d\times d$ matrix.  Its ordinary-degree decomposition is
\begin{equation}\label{eq:section-character-v155}
 (B_{d,h})_p\simeq
 \bigoplus_{\substack{\lambda\subseteq(h^d)\\|\lambda|=p}}
                         \mathbb S_{2\lambda}(W^*).
\end{equation}
In particular it is Gorenstein of socle degree $dh$, with
one-dimensional socle representation $(\det W^*)^{2h}$.  Its Hilbert
function and length are
\begin{align}
 H_{d,h}(p)&=\sum_{\substack{\lambda\subseteq(h^d)\\|\lambda|=p}}
       \prod_{1\leq i<j\leq d}
             \frac{2\lambda_i-2\lambda_j+j-i}{j-i},
                     \label{eq:section-hilbert-v155}\\
 L_{d,h}&=\prod_{i=1}^d\frac{h+d-i+1}{d-i+1}
       \prod_{1\leq i<j\leq d}
             \frac{2h+2d-i-j+2}{2d-i-j+2}.
                     \label{eq:section-length-v155}
\end{align}
The Hilbert function is symmetric.  A degree-one element is strong
Lefschetz precisely when its associated symmetric matrix has full
rank.  More generally, an element associated with a rank-$r$ matrix
has nilpotence index $hr+1$ (index one for the zero element).
Consequently
\begin{equation}\label{eq:full-fibre-bounds-v155}
 \length F_{g,k}(W)\geq L_{d,h},\qquad
 \mathfrak m^{dh}\ne0\text{ in }F_{g,k}(W).
\end{equation}
The Gorenstein and Lefschetz assertions concern the section
$B_{d,h}$, not the full fibre $F_{g,k}(W)$.
\end{theorem}
\begin{proof}
After setting the indicated coefficients to zero, the reciprocal
series is $(1+Qz^2)^{-1}$.  The first even integer in
$\{g+1,\ldots,g+k\}$ is $2(h+1)$; it belongs to this interval
because $k\geq2$.  Every odd coefficient vanishes, and the even
ones are signed powers of $Q$.  Thus the specialized ideal is
exactly $(\operatorname{coeff}Q^{h+1})$: all its higher powers have
coefficients in this ideal.  This proves \eqref{eq:exact-section-v155}
as an equality of ideals, not only an equality of radicals.

For clarity about the pairing, write
$Q(x)=\sum_i y_{ii}x_i^2+\sum_{i<j}y_{ij}x_ix_j$ and let $X$ have
independent symmetric entries $X_{ij}$, $i\leq j$.  Send $y_{ij}$
to $\partial/\partial X_{ij}$.  Then $Q(x)$ acts by differentiation
in the rank-one direction $xx^{\mathsf t}$.  Since
$\det(X+zxx^{\mathsf t})^h$ has degree at most $h$ in $z$, the
coefficients of $Q(x)^{h+1}$ annihilate $(\det X)^h$.
The reverse containment is the symmetric-matrix case of
\cite[Remark~11, pp.~8--9]{NagaokaWachi2024}: that apolar ideal is
generated by the full general-linear orbit of the $(h+1)$st power
of a rank-one differentiation direction.  Powers of linear forms
span their symmetric power, so this orbit span is precisely the
coefficient space just displayed.  No assertion that an arbitrary
quotient of an apolar algebra is Gorenstein is used.

We record explicitly how the classical representation formulas
translate here.  In \cite[Example~9 and Proposition~10]{NagaokaWachi2024},
the surviving highest weights are indexed by
$k_1,\ldots,k_d\geq0$ with $\sum k_i\leq h$, and their degrees are
$\sum i k_i$.  Set $\lambda_i=\sum_{j\geq i}k_j$.
These are exactly the partitions in $(h^d)$, of size $\sum i k_i$,
and their polynomial representations are $\mathbb S_{2\lambda}(W^*)$.
This proves \eqref{eq:section-character-v155}; the ordinary Weyl
dimension formula gives \eqref{eq:section-hilbert-v155}.
The same apolar module is the finite-dimensional highest-weight
$\mathfrak{sp}_{2d}$-module of weight $h\omega_d$, with its Levi
character shifted by $\det^h$; see
\cite[Proposition~16 and Lemma~13]{NagaokaWachi2024}.
For the positive roots $e_i-e_j$, $e_i+e_j$, and $2e_i$, its Weyl
dimension factors are respectively $1$,
$(2h+2d-i-j+2)/(2d-i-j+2)$, and
$(h+d-i+1)/(d-i+1)$.  Their product is
\eqref{eq:section-length-v155}.
Apolarity to one homogeneous polynomial of degree $dh$ gives the
perfect pairings and the socle.  Equivalently, the complementary
partition $(h-\lambda_d,\ldots,h-\lambda_1)$ gives the dual summand
with the indicated determinant twist.

The strong Lefschetz assertion is the symmetric-matrix case of
\cite[Theorem~4 and Example~5]{NagaokaWachi2024}.
There is a direct argument for the asserted nilpotence index.
If $\ell_A=\sum_{i\leq j}A_{ij}y_{ij}$, then
$\ell_A^p=0$ in the apolar algebra exactly when
$\partial_z^p\det(X+zA)^h=0$ as a polynomial in $X$.
Congruence reduces a rank-$r$ symmetric matrix to
$\diag(1,\ldots,1,0,\ldots,0)$.
The last determinant power has degree exactly $hr$ in $z$, as
is seen by taking $X$ diagonal with independent entries.
This proves the index.  Finally a surjective local algebra map
surjects on every power of its maximal ideal.  The section has
nonzero degree $dh$, proving both assertions in
\eqref{eq:full-fibre-bounds-v155}.
\end{proof}

\begin{corollary}[Size and Loewy depth of the actual common boundary fibre]
\label{cor:boundary-size-v155}
In the notation of Theorem~\ref{thm:universal-boundary-v152}, put
$d=n^2-1$ and $h=n(n-1)/2$.  The entire normalization fibre at
$K_\infty$ has length at least $L_{d,h}$ and its maximal ideal has
nonzero $dh$th power.  For the ternary pencil boundary,
\begin{equation}\label{eq:ternary-large-section-v155}
 \length F_{6,2}(\C^8)\geq922268360,
       \qquad\mathfrak m^{24}\ne0.
\end{equation}
Its quadratic coefficient section has exactly that length, ordinary
socle degree $24$, and reciprocal-weight socle $48$.
\end{corollary}
\begin{proof}
The common boundary fibre is $F_{g,k_n}(\C^{n^2-1})$, with
$g=n(n-1)\geq k_n\geq2$.  Apply the theorem.
For $n=3$, substitute $(d,h)=(8,3)$ into
\eqref{eq:section-length-v155}.
This computes a whole closed section, not the length of the original
fibre.  The previously obtained $44$ coefficient variables and
$9867$ minimal equations remain counts of different invariants.
\end{proof}

\section{Spectral Fitting schemes as power-zero loci in one reciprocal section}
\label{sec:power-spectral-v155}
The algebra just calculated supplies a single geometric object in
which all spectral Fitting schemes can be read simultaneously.
This avoids constructing a separate binary form after each Fitting
ideal has already been recovered.

Let $V$ have dimension $n$ and $R\subset\Sym^2V$ be any pencil.
Fix $h\geq1$ and put $\widetilde E=\C a\oplus V^*$.
Consider the divisor-incidence point
\begin{equation}\label{eq:canonical-reciprocal-point-v155}
 K_h=a^{2h+2}\widetilde E
         \subset\Sym^{2h+3}\widetilde E.
\end{equation}
The entire degree-$2h$ normalization fibre over $K_h$ is
$F_{2h,2}(V^*)$, because $\widetilde E$ is gcd-free.  Its section
by the linear-factor coefficients is $B_{n,h}$ and
$(B_{n,h})_1=\Sym^2V$.  Thus the given pencil defines a line
$\Lambda=\Pj(R)$ in the projectivized degree-one space of this
actual reciprocal section.  This construction uses a new formal
line $\C a$ and is functorial for isomorphisms of $V$.

For $0\leq p\leq hn$ let $\mathcal J_p$ be the ideal sheaf on
$\Lambda$ locally generated by the coordinates of
\begin{equation}\label{eq:power-map-v155}
       \ell^p\in(B_{n,h})_p\otimes\OO_\Lambda(p),
\end{equation}
where $\ell$ is the universal pencil member.  Multiplying a local
representative of $\ell$ by a unit does not change this ideal.
We set $\mathcal J_0=\OO_\Lambda$.
Let $\mathcal T_R$ be the cokernel of the universal symmetric
matrix $V^*\otimes\OO_\Lambda(-1)\to V\otimes\OO_\Lambda$.
This definition includes singular pencils.

\begin{lemma}[The local power-ideal formula]
\label{lem:power-valuation-v155}
Let $A(t)$ be a symmetric matrix over $\C[[t]]$ of rank $r$ over
$\C((t))$.  Write its finite elementary exponents as
$0\leq e_1\leq\cdots\leq e_r$ and put $e_i=+\infty$ for $i>r$.
For $1\leq p\leq hn$, write $p=hq+s$, $0\leq s<h$.
The ideal generated by the coordinates of
$\ell_{A(t)}^p$ in $(B_{n,h})_p\otimes\C[[t]]$ is
\begin{equation}\label{eq:valuation-power-v155}
 \left(t^{\,h\sum_{i=1}^q e_i+s e_{q+1}}\right).
\end{equation}
An empty sum is zero, the last term is omitted if $s=0$, and
$t^{+\infty}$ means the zero ideal.
\end{lemma}
\begin{proof}
A symmetric matrix over this discrete valuation ring is congruent
to
\[\diag(t^{e_1},\ldots,t^{e_r},0,\ldots,0).\]
Here is the required elementary reduction.  Divide by the least
valuation of an entry.  A diagonal unit is a pivot; if only an
off-diagonal entry is a unit, the value of the form on a sum of two
basis vectors supplies a unit because $2$ is invertible.  Split
off the resulting rank-one summand by elimination and continue.
Every remaining unit has a square root in $\C[[t]]$, so the units
can be removed from the diagonal.  The finite exponents are the
Smith exponents, as congruence is in particular equivalence.

Congruence acts on every graded piece of $B_{n,h}$ by an invertible
matrix over $\C[[t]]$.  It therefore preserves the ideal in question.
The apolar embedding
\[
 (B_{n,h})_p\longrightarrow
       \C[X_{ij}]_{hn-p},\qquad
       D\longmapsto D(\det X)^h
\]
is injective and splits as a map of complex vector spaces.  After
tensoring with $\C[[t]]$, the coordinates of a vector and of its
image generate the same ideal.  We may therefore calculate the
coefficient ideal of $[z^p]\det(X+zA(t))^h$.

For diagonal $A(t)$, expand each of the $h$ determinant factors
in its chosen diagonal entries.  Every term is divisible by
$t^{\sum c_i e_i}$ with $0\leq c_i\leq h$ and $\sum c_i=p$.
The least such exponent is obtained by filling the smallest $e_i$
to capacity $h$, giving \eqref{eq:valuation-power-v155}.
There is no cancellation at this least exponent: specialize $X$
to $\diag(x_1,\ldots,x_n)$.  Then
\[
 \det(X+zA(t))^h=\prod_i(x_i+zt^{e_i})^h
\]
(with $x_i^h$ for a zero diagonal entry of $A$).
The monomial $\prod x_i^{h-c_i}$ has coefficient
$\prod\binom h{c_i}\,t^{\sum c_i e_i}$, which is nonzero and has
exactly the required valuation.  If $p>hr$, every term is zero.
This proves all the conventions, including the zero-ideal case.
\end{proof}

\begin{theorem}[A reciprocal multiplication model for every spectral Fitting ideal]
\label{thm:power-Fitting-v155}
For every complex quadratic pencil, regular or singular, and for
$0\leq j\leq n$, there is an equality of coherent ideal sheaves
on its intrinsic line
\begin{equation}\label{eq:power-Fitting-v155}
          \mathcal J_{h(n-j)}=\Fitt_j(\mathcal T_R)^h.
\end{equation}
In particular, when $h=1$, the complete spectral Fitting schemes,
including their multiplicities and any whole-line components, are
exactly the power-zero schemes inside the one quadratic section
of the normalization fibre at $K_1$.

For a regular pencil and a spectral point $x$, set
$b_{x,j}=\length\OO_{\Lambda,x}/\Fitt_j(\mathcal T_R)_x$.
The local length of the power-zero scheme is $h b_{x,j}$, and,
for $0\leq j<n$,
\begin{equation}\label{eq:power-Segre-v155}
 \lambda_{x,j+1}
   =\frac1h\left(\length(\OO_{\Lambda,x}/\mathcal J_{h(n-j),x})
           -\length(\OO_{\Lambda,x}/\mathcal J_{h(n-j-1),x})\right).
\end{equation}
Thus its multiplication, together with the embedded pencil line,
recovers the full Segre data and their common projective support.
For $n\geq3$, the unmarked inverse makes this an invariant of the
finite failure algebra.  For a singular pencil, the construction recovers its
normal rank and all its spectral Fitting ideals; it does not assert
that these ideals alone recover the singular minimal indices.
\end{theorem}
\begin{proof}
At a point of the smooth curve $\Lambda$, pass to the completed
discrete valuation ring.  The ideal of $(n-j)$-minors of a matrix
with the stated Smith exponents is
$(t^{e_1+\cdots+e_{n-j}})$, with the zero-ideal convention.
Apply Lemma~\ref{lem:power-valuation-v155} with $p=h(n-j)$.
The equality of completed ideals descends by faithful flatness.
This proves \eqref{eq:power-Fitting-v155} at every stalk and hence
globally.  This proof does not replace an ideal by its radical.

For a regular pencil the positive Smith exponents at $x$, in
reverse order, are $\lambda_{x,1}\geq\lambda_{x,2}\geq\cdots$.
Consequently $b_{x,j}=\sum_{i>j}\lambda_{x,i}$.
Equation~\eqref{eq:power-Fitting-v155} multiplies these lengths by
$h$, so taking successive differences proves
\eqref{eq:power-Segre-v155}.  The powers that vanish at the generic
point similarly determine the normal rank in the singular case.
Frame changes act on $\Lambda$ by projectivities and congruences
act equivariantly on the reciprocal construction.  Theorem~
\ref{thm:artin-local-inverse-v146} recovers precisely the needed
pencil up to these changes, proving the intrinsic assertion.
\end{proof}

\begin{remark}[Families and the scope of the ideal identity]
\label{rem:power-basechange-v155}
For a vector bundle and pencil subbundle over any complex base,
the quadratic-section algebra is obtained by applying the fixed
$\operatorname{GL}(V)$-equivariant construction to local frames.
Its graded pieces are vector bundles, and the coefficient ideal of
each power map commutes with arbitrary base change: it is the image
of the dual of a map between vector bundles.  No flatness of its
zero scheme is needed for that assertion.  Formula
\eqref{eq:power-Fitting-v155} has been proved on each complex
pencil line (and on a regular one-dimensional base by the same
local argument).  We do not infer an equality of those relative
ideals on an arbitrary nonreduced parameter space from equality
on its geometric fibres.  Nor are the power-zero schemes asserted
to be flat across a change in their length.
\end{remark}

\begin{corollary}[The native specialization order in reciprocal coordinates]
\label{cor:power-order-v155}
On the fixed-discriminant, fixed-corank locus
$X_{\boldsymbol m,\boldsymbol c}(A)$ of
\eqref{eq:fixed-spectral-locus}, a specialization from the labelled
partition tuple $\boldsymbol\lambda$ to $\boldsymbol\mu$ exists
if and only if, at every labelled root and every $j$, the local
length of the corresponding power-zero scheme does not decrease.
Equivalently the special power-zero divisor contains the generic
one, using the common labelled spectral line.
The reduced closure of a full $O(A)$-orbit within this locus is
therefore specified by these length inequalities.  Each inequality
can be written as the vanishing of the requisite jets of the
coordinates of the power map \eqref{eq:power-map-v155}.
Every strict transition is realized by a flat family of the original
finite failure neighbourhoods and is first distinguished at order
$n^2+2n-4$.
\end{corollary}
\begin{proof}
For fixed multiplicity $m_x$, dominance
$\mu_x\preceq\lambda_x$ is equivalent to
$\sum_{i>j}\mu_{x,i}\geq\sum_{i>j}\lambda_{x,i}$ for every $j$.
Use \eqref{eq:power-Segre-v155} and
Theorem~\ref{thm:fixed-spectral-classification}.
In a local coordinate $t$ at a labelled root, a coefficient ideal
has colength at least $b$ exactly when each of its generators is
divisible by $t^b$; this is the asserted finite jet condition.
The closure order and the flat realization are the inherited
full-orthogonal results, ultimately using Ohta's theorem.  The
new assertion is their expression inside one reciprocal
multiplication algebra.  No reducedness assertion is made for the
scheme cut out by the jet equations before taking its reduction.
\end{proof}

\begin{example}[A rank profile with a genuinely different power-zero scheme]
\label{ex:power-profile-v155}
At a root of multiplicity four the partitions $(3,1)$ and $(2,2)$
have equal determinant and corank.  In $B_{4,h}$ their sixth powers
when $h=2$ (in general their $3h$th powers) cut out local schemes
of lengths $h$ and $2h$, respectively.  Their $4h$th powers have
length $4h$ in both cases.  Thus the nonreduced power-zero scheme,
not just the set of non-Lefschetz elements, distinguishes them.
The specialization $(3,1)\rightsquigarrow(2,2)$ increases the first
of these lengths.  Additional simple spectral blocks give the
same distinction in all larger dimensions, with power indices
$h(n-1)$ and $hn$.
\end{example}

The representation theory and strong Lefschetz property used above
are classical inputs, explicitly credited to Nagaoka--Wachi and,
for $h=1$, Shafiei.  The geometric statements here are the exact
section of the whole reciprocal family, its size and Loewy-depth
consequences for actual first-relation boundary fibres, and the
power-ideal identity on pencil lines.  The last construction is an
intrinsically induced reciprocal model; it is not a claim that the
common extreme $\operatorname{GL}(E)$ boundary point remembers the
starting pencil.  Its embedded pencil line is essential.

The preceding relative caution concerns the valuation-ring proof. The following polynomial identity establishes the equality over arbitrary complex base algebras.

\section{Universal power ideals}
\label{sec:universal-power-v157}
The power-zero construction admits an equality of ideals before any
restriction to a pencil or to a valuation ring. This is the point at
which the reciprocal section acquires a global boundary geometry.
Let $V$ have dimension $n\geq2$, put $M=\Sym^2V$, and use the algebra
$B_{n,h}$ of Theorem~\ref{thm:quadratic-section-v155} with $W=V^*$.
Thus $(B_{n,h})_1=M$. If $A$ is a symmetric matrix with entries in a
commutative $\C$-algebra $T$, denote its element of $M\otimes T$ by
$\ell_A$, and set
\begin{equation}\label{eq:universal-power-definition-v157}
 J_p(A)=\im\bigl((B_{n,h})_p^*\otimes T\longrightarrow T,
                  \quad\varphi\longmapsto\varphi(\ell_A^p)\bigr).
\end{equation}
Write $I_q(A)$ for the ideal of all $q$-minors, with $I_0(A)=T$.
In this section powers of ideals are ordinary powers unless parentheses
explicitly indicate symbolic powers.

\begin{theorem}[Universal determinantal formula]
\label{thm:universal-power-v157}
For every commutative $\C$-algebra $T$, every symmetric $n\times n$
matrix $A$ over $T$, and $0\leq p\leq nh$, write $p=hq+s$ with
$0\leq s<h$. Then
\begin{equation}\label{eq:universal-power-v157}
             J_p(A)=I_q(A)^{h-s}I_{q+1}(A)^s.
\end{equation}
When $s=0$ the second factor is omitted. For $p>nh$, $J_p(A)=0$.
These are equalities of ideals, not of radicals or integral closures,
and commute with every homomorphism of $\C$-algebras. In particular,
for every symmetric map $A:T^n\to T^n$,
\begin{equation}\label{eq:universal-Fitting-v157}
       J_{h(n-j)}(A)=\Fitt_j(\coker A)^h,\qquad 0\leq j\leq n.
\end{equation}
The formulas also hold locally, and hence as ideal sheaves, for
symmetric maps of vector bundles with a line-bundle twist on an
arbitrary $\C$-scheme.
\end{theorem}

We separate the classical determinantal input. For an arbitrary matrix
$X$ over a characteristic-zero polynomial ring, the balancing inclusion
\begin{equation}\label{eq:minor-balancing-v157}
              I_u(X)I_v(X)\subseteq I_{u+1}(X)I_{v-1}(X),
              \qquad 0\leq u\leq v-2,
\end{equation}
is the determinantal symmetrization lemma. See
\cite[Lemma~2.5 and Theorem~2.4]{BrunsConca2003}, which attributes the
characteristic-zero product theorem to De Concini--Eisenbud--Procesi.
Because this is a polynomial ideal inclusion, it remains true after
specialization to a symmetric matrix; no flatness of that specialization
is needed. This classical lemma is not a claim of the present paper.

\begin{lemma}[The bounded-row generating space]
\label{lem:bounded-row-v157}
Let $S=\Sym(\Sym^2V^*)$ and let $A$ be its universal symmetric matrix.
For $p=hq+s$, the degree-$p$ part generating
$I_q(A)^{h-s}I_{q+1}(A)^s$ is
\begin{equation}\label{eq:bounded-row-v157}
 \bigoplus_{\substack{|\lambda|=p,\ \ell(\lambda)\leq n\\
                         \lambda_1\leq h}}\mathbb S_{2\lambda}(V^*)
                         \ \subseteq S_p.
\end{equation}
\end{lemma}
\begin{proof}
The symmetric Cauchy decomposition of $S_p$ is multiplicity free,
with summands $\mathbb S_{2\lambda}(V^*)$, $|\lambda|=p$; see
\cite[\S2.4]{HenriquesVarbaro2016}. We may use the dual general
linear group to regard these as polynomial representations. Every
minor has torus weight at most two in any one coordinate: its row
set and column set each use that coordinate at most once. A product
of $h$ minors therefore has no weight with a coordinate exceeding
$2h$. Its invariant span cannot contain a highest weight
$2\lambda$ with $\lambda_1>h$. This proves one inclusion.

For the converse, let $c_1\geq\cdots\geq c_h\geq0$ be the column
lengths of a partition $\lambda$ occurring on the right, padded with
zeros. The product of the leading principal minors of orders $c_i$
is a nonzero highest-weight vector of weight $2\lambda$ in $S_p$.
It belongs initially to $\prod_{i=1}^h I_{c_i}(A)$.
Whenever two column lengths differ by at least two, apply
\eqref{eq:minor-balancing-v157} to increase the smaller length by one
and decrease the larger by one. The sum of the squares of the lengths
strictly decreases, so the procedure terminates. At termination the
lengths are $q+1$ repeated $s$ times and $q$ repeated $h-s$ times,
since their sum is $p$. Thus that highest-weight vector belongs to
$I_q^{h-s}I_{q+1}^s$. The latter ideal is invariant; its degree-$p$
part contains the full irreducible summand. Multiplicity freeness
proves the equality. The case $q=0$ includes zero-length minors,
which are the constant $1$.
\end{proof}

\begin{proof}[Proof of Theorem~\ref{thm:universal-power-v157}]
First work over $S$ with its universal matrix. The coefficient map of
$\ell_A^p$ is the dual of the multiplication quotient
$\Sym^p M\twoheadrightarrow (B_{n,h})_p$, up to nonzero multinomial
constants. Consequently its degree-$p$ coefficient span is exactly
$(B_{n,h})_p^*\subseteq S_p$. The character formula
\eqref{eq:section-character-v155} identifies that span with
\eqref{eq:bounded-row-v157}. Both ideals in
\eqref{eq:universal-power-v157} are generated in degree $p$, so the
lemma proves their equality in $S$. This step uses the full coefficient
space, not just evaluation on diagonal matrices or closed points.

For a matrix over $T$, substitute its entries in this polynomial
identity. Taking the image of the fixed finite-dimensional coefficient
map commutes with scalar extension, as does taking minors and products
of ideals. This proves the assertion even when $T$ is nonreduced,
non-Noetherian, or not a domain. The Fitting formula follows from its
minor presentation. Changing a vector-bundle frame applies an invertible
congruence to $A$ and an invertible linear transformation to each
coefficient space; changing a local frame of the twisting line multiplies
generators by a unit. Therefore the local ideal identities glue. Finally,
$(B_{n,h})_p=0$ for $p>nh$.
\end{proof}

\begin{example}[A nonreduced test base]
\label{ex:nilpotent-base-v157}
Take $T=\C[\epsilon]/(\epsilon^3)$, $n=h=2$, and
$A=\diag(\epsilon,0)$. Then $J_1(A)=(\epsilon)$ and
$J_2(A)=(\epsilon^2)\ne0$, whereas $J_3(A)=J_4(A)=0$.
Every closed-point matrix is zero. Thus equality on geometric fibres
would not establish the theorem. Nor does
\eqref{eq:universal-Fitting-v157} license taking unique ideal roots:
for example, in $\C[\epsilon]/(\epsilon^2)$ the distinct ideals
$(\epsilon)$ and $0$ have the same square. At $h=1$ the power ideals
are the Fitting ideals themselves over every base.
\end{example}

\section{The reciprocal power graph and complete quadrics}
\label{sec:power-graph-v157}
Let $P=\Pj(M)$, let $P^{\mathrm{nd}}$ be the nonsingular-quadric open,
and define $\mathcal G_h$ as the scheme-theoretic closure of the graph
\begin{equation}\label{eq:power-graph-v157}
 P^{\mathrm{nd}}\longrightarrow
 P\times\prod_{p=1}^{nh-1}\Pj((B_{n,h})_p),
 \qquad [A]\longmapsto([A],[\ell_A],\ldots,[\ell_A^{nh-1}]).
\end{equation}
The first factor records the base point. We use the reduced, integral
graph closure, equivalently its scheme-theoretic image; no additional
normalization is inserted in this definition. Let $\mathrm{CQ}(V)$
be the classical space of complete quadrics: the graph closure of
$[A]\mapsto([\bigwedge^2A],\ldots,[\bigwedge^{n-1}A])$ over $P$.
For $n=2$ the empty product gives $\mathrm{CQ}(V)=P$.

\begin{theorem}[A compactification determined by multiplication]
\label{thm:power-graph-v157}
For every $h\geq1$ there is a canonical $\operatorname{GL}(V)$-equivariant
isomorphism over $P$
\begin{equation}\label{eq:power-graph-CQ-v157}
                      \mathcal G_h\simeq\mathrm{CQ}(V).
\end{equation}
In particular the graph is smooth and projective, is independent of
$h$ as a scheme over $P$, and has a simple-normal-crossing boundary.
Write $\pi:\mathrm{CQ}(V)\to P$, and let $E_r$ be its boundary divisor
of rank break $r$, $1\leq r<n$; for $r<n-1$ it is exceptional over
rank at most $r$, and $E_{n-1}$ is the strict transform of the determinant
divisor. Then every power ideal has the exact total transform
\begin{equation}\label{eq:power-boundary-v157}
 J_p\OO_{\mathrm{CQ}(V)}=
 \OO_{\mathrm{CQ}(V)}\left(-\sum_{r=1}^{n-1}(p-hr)_+E_r\right),
                    \qquad 0\leq p\leq nh,
\end{equation}
where $(a)_+=\max(a,0)$. Removing this common divisor from the universal
$p$th power gives a nowhere-zero projective section on the whole
compactification.
\end{theorem}

The smoothness, boundary and flag interpretation of complete quadrics
are classical, due to the complete-form construction; see
\cite{Vainsencher1982,Massarenti2020}. The assertion here identifies
that particular compactification, with all of its scheme structure,
from the multiplication in one reciprocal normalization-fibre section.
It does not assert that complete quadrics is a newly constructed variety.

\begin{proof}
We first recall the graph calculation with its scheme structure.
If rational maps from an integral scheme are given by generators of
nonzero ideals $K_1,\ldots,K_a$, their simultaneous graph closure is
$\operatorname{Bl}_{K_1\cdots K_a}$ of the base. Indeed, after the Segre
embedding the coordinates are all products of one generator from each
ideal. The closure of this graph is the Proj of the Rees algebra of
that product: on a generator chart its coordinate ring is the subring
of the function field obtained by adjoining the generator ratios.
These are exactly the affine charts of both constructions. Tensoring
an ideal by an invertible ideal, or replacing it by a positive power,
does not change the blow-up, the latter by the Veronese description
of Proj.

Apply this calculation on $P$. The sheaf $I_1\OO_P$ is the unit ideal,
and $I_n\OO_P$ is invertible. For each $2\leq q\leq n-1$, the exponent
of $I_q$ in $\prod_{p=1}^{nh-1}J_p$ is
\[
 \sum_{s=0}^{h-1}(h-s)+\sum_{s=1}^{h-1}s=h^2.
\]
Theorem~\ref{thm:universal-power-v157} therefore gives, up to an
invertible ideal,
\[
       \prod_{p=1}^{nh-1}J_p
                =\left(\prod_{q=2}^{n-1}I_q\right)^{h^2}.
\]
The same graph calculation for the exterior-power maps identifies
$\mathrm{CQ}(V)$ with the blow-up of $\prod_{q=2}^{n-1}I_q$.
This proves \eqref{eq:power-graph-CQ-v157} as an isomorphism of schemes,
not only after normalization. Both isomorphisms are the identity on
the common dense open, hence are canonical and equivariant.

For completeness, the local calculation behind the boundary formula
is as follows. The complete-quadric construction blows up the strict
rank loci successively. On its Gaussian charts, after a congruence
and projective normalization, its matrix has the form
\begin{equation}\label{eq:CQ-chart-v157}
 U\diag(1,\tau_1,\tau_1\tau_2,\ldots,
                         \tau_1\cdots\tau_{n-1})U^{\mathsf t},
\end{equation}
with $U$ lower unitriangular and $E_r=(\tau_r=0)$.
One obtains these charts inductively by making the first pivot one,
taking a Schur complement, blowing up the zero residual matrix, and
repeating. At each step the next nonzero residual quadratic form has
a nonisotropic vector, so congruence translates of these charts cover
the construction. This also explains the smooth charts and the normal
crossings stated in the classical construction.

Cauchy--Binet, applied to \eqref{eq:CQ-chart-v157}, shows that every
$q$-minor is divisible by
$\tau_1^{q-1}\tau_2^{q-2}\cdots\tau_{q-1}$.
The leading principal $q$-minor is exactly that product. Thus, on
these charts and therefore globally,
\begin{equation}\label{eq:minors-on-CQ-v157}
 I_q\OO_{\mathrm{CQ}(V)}=
              \OO\left(-\sum_{r<q}(q-r)E_r\right).
\end{equation}
Using \eqref{eq:universal-power-v157}, its exponent along $E_r$ is
$(h-s)(q-r)_++s(q+1-r)_+=(p-hr)_+$.
This proves \eqref{eq:power-boundary-v157}. The coefficient sections
generate the transformed ideal; division by its local generator makes
them generate the unit ideal, proving the last assertion.
\end{proof}

\begin{corollary}[Native boundary strata]
\label{cor:native-strata-v157}
The power graph has one congruence orbit for every subset
$I\subseteq\{1,\ldots,n-1\}$. It is the stratum
$E_I^\circ=(\bigcap_{r\in I}E_r)\setminus\bigcup_{r\notin I}E_r$,
of codimension $|I|$, and its closure is $\bigcap_{r\in I}E_r$.
Its points are flags with nondegenerate projective quadratic forms on
successive quotients; the quotient dimensions are the successive gaps
in $\{0\}\cup I\cup\{n\}$. The simultaneous normalized power values
recover this complete-quadric datum. They are not the one common limit
for the full cotangent-coordinate group.
\end{corollary}
\begin{proof}
The flag description and boundary stratification are the classical
complete-quadric description. The general linear group acts transitively
on flags of fixed dimensions, and over $\C$ every nondegenerate
quadratic form on a fixed-dimensional quotient is congruent to the
identity. Independent diagonal scalings remove the projective scalars.
The chart \eqref{eq:CQ-chart-v157} realizes each subset of vanishing
$\tau_r$ and permits further coordinates to vanish, giving the stated
codimensions and closures. The graph is a closed subscheme of the
product in \eqref{eq:power-graph-v157}; its projections therefore
retain the point, including its flag and quotient forms. The group
here is the native congruence group of $V$.
\end{proof}

\begin{remark}[Relative construction and base change]
\label{rem:relative-CQ-v157}
For a vector bundle $\mathcal V$ on any $\C$-scheme $S$, the associated
bundle with fibre $\mathrm{CQ}(\C^n)$ is a relative smooth projective
scheme over $S$. Equivariance glues the power maps and the boundary
formula. Formation of this associated bundle and the polynomial ideal
identity commute with arbitrary base change. This does not say that
taking the graph closure of an arbitrary restricted family commutes
with nonflat base change. The relative ambient resolution and the
closure of an individual family's graph are different constructions.
\end{remark}

\section{Boundary contacts and spectral schemes}
\label{sec:boundary-contacts-v157}
Let $R\subset M$ be a regular pencil and $\Lambda=\Pj(R)$. Its map
to $P$ lifts uniquely to a morphism
$\widetilde\iota_R:\Lambda\to\mathrm{CQ}(V)$: it lifts on the
nonsingular open, and properness extends it uniquely across the missing
points of the smooth curve. Put
$D_r(R)=\widetilde\iota_R^*E_r$. These are effective divisors because
the generic member is nonsingular.

\begin{theorem}[Contact-divisor reconstruction of spectral data]
\label{thm:contact-Segre-v157}
For every $h\geq1$, every regular pencil, and $0\leq p\leq nh$,
\begin{equation}\label{eq:contact-ideals-v157}
 J_p|_\Lambda=\OO_\Lambda\left(-\sum_{r=1}^{n-1}(p-hr)_+D_r(R)\right).
\end{equation}
At a point $x\in\Lambda$, let $0=e_1\leq\cdots\leq e_n$ be the
symmetric Smith exponents of the local pencil matrix. Then
\begin{equation}\label{eq:contact-exponents-v157}
 \operatorname{mult}_xD_r(R)=e_{r+1}-e_r\quad(1\leq r<n).
\end{equation}
Consequently the supported contact divisors determine every elementary
divisor, hence the full Segre symbol. If
$v_p(x)=\length(\OO_{\Lambda,x}/J_p)$ and $v_0(x)=0$, then
\begin{equation}\label{eq:contact-second-difference-v157}
 \operatorname{mult}_xD_r(R)=
 \frac{v_{h(r+1)}(x)-2v_{hr}(x)+v_{h(r-1)}(x)}{h}.
\end{equation}
\end{theorem}
\begin{proof}
Pull back \eqref{eq:power-boundary-v157}; equality of extended ideals
under composition gives \eqref{eq:contact-ideals-v157}.
The pencil matrix is nonzero at every point, since $R$ is a genuine
two-dimensional subspace; hence its entries generate the unit ideal
locally and $e_1=0$. Over $\C[[t]]$ symmetric elimination gives a
congruence to $\diag(t^{e_1},\ldots,t^{e_n})$, absorbing unit factors
by square roots. Its unique lift has, in
\eqref{eq:CQ-chart-v157}, $\tau_r=t^{e_{r+1}-e_r}$ up to a unit.
This proves \eqref{eq:contact-exponents-v157}; faithful flatness of
completion gives the local divisor statement. Taking valuations in
\eqref{eq:contact-ideals-v157} and second differences at multiples of
$h$ proves \eqref{eq:contact-second-difference-v157}.
Recover $e_i$ by summing the first $i-1$ contact multiplicities. The
positive exponents, with their points on $\Lambda$, are precisely
the elementary-divisor partitions of the spectral cokernel.
\end{proof}

\begin{example}[A boundary point and its contact order are distinct data]
\label{ex:contact-profile-v157}
For a symmetric germ with exponents $(0,1,3)$ and $h=2$, the valuations
of $J_1,\ldots,J_6$ are $(0,0,1,2,5,8)$.
Its two boundary contacts are $1$ and $2$. The germs
$\diag(1,t,t^3)$ and $\diag(1,t^2,t^5)$ reach the same full-flag
complete-quadric point but have contacts $(1,2)$ and $(2,3)$.
The compactification retains the flag; the arc with its divisorial
contacts retains the elementary exponents. These assertions must not
be conflated. For singular pencils the universal ideal theorem still
holds and detects all Fitting schemes and the normal rank; Fitting
ideals alone do not determine Kronecker minimal indices.
\end{example}

\section{A projective incidence compactification for pencils}
\label{sec:pencil-compactification-v157}
The complete-quadric graph concerns members of pencils. There is also
a projective parameter space carrying both the original finite failure
algebra and flat limits of their resolved member curves. This construction
keeps the pencil and does not identify distinct pencils at a common
full-coordinate-change limit.

\begin{proposition}[Resolved pencil curves and the failure family]
\label{prop:pencil-compactification-v157}
Let $n\geq3$, $G=\Gr(2,M)$, and let $G^\circ$ be the open set of
regular pencils whose lines avoid the rank-at-most-$(n-2)$ locus in
$P$. There is a normal projective $\operatorname{GL}(V)$-equivariant
scheme $\widehat G$ with a projective birational morphism
$\rho:\widehat G\to G$, an isomorphism over $G^\circ$, and a flat
family of subschemes
$\mathcal C\subseteq\widehat G\times\mathrm{CQ}(V)$ with the following
properties. Over $G^\circ$ it is exactly the lifted pencil line. Every
fibre over $R$ lies scheme-theoretically in $\pi^{-1}(\Pj(R))$.
All the resolved power maps restrict to $\mathcal C$, for every $h$.
The sharp finite failure algebra on $G$ pulls back to a finite locally
free algebra on $\widehat G$, with its original multiplication.
\end{proposition}
\begin{proof}
The rank-at-most-$(n-2)$ locus has codimension three in $P$, so a general
line avoids it; a general pencil is regular. Thus $G^\circ$ is nonempty
and dense. Over that open, $\pi$ is an isomorphism along every pencil
line and these lifts form a flat family of embedded projective lines.

For $1\leq q<n$, let $H_q$ be the pullback to $\mathrm{CQ}(V)$ of the
hyperplane bundle for its $q$th exterior-power projection, taking $q=1$
to mean $P$. The bundle $L=\bigotimes_{q=1}^{n-1}H_q$ is very ample,
by the graph embedding and Segre embedding. On a lifted line over
$G^\circ$, $\deg H_q=q$, since its $q$-minors have no common zero.
The Hilbert polynomial with respect to $L$ is therefore
$\binom n2 m+1$. The lifted family defines a morphism
\[
 G^\circ\longrightarrow
       \Hilb^{\binom n2 m+1}(\mathrm{CQ}(V)).
\]
Take the reduced closure of its graph in $G$ times this projective
Hilbert scheme, and then its normalization; this is $\widehat G$.
Normalization is finite here, so the result is projective and normal.
The graph is already closed over $G^\circ$ and that open is smooth,
so $\rho$ is an isomorphism there. Equivariance of the family and the
uniqueness of normalization give the stated group action.

Pull back the universal Hilbert family to obtain $\mathcal C$; it is
flat. To verify the asserted scheme-theoretic incidence, pull back
the linear equations for the universal pencil line to $\mathcal C$.
They vanish over the dense open $G^\circ$. On an affine open of the
integral base $\widehat G$, flatness makes the coordinate modules of
$\mathcal C$ torsion free over the base. A section which becomes zero
at the generic point is therefore zero. Applied locally to these
linear equations, this proves the incidence on the whole family.
The power maps are everywhere defined on $\mathrm{CQ}(V)$ by
Theorem~\ref{thm:power-graph-v157}, so their restrictions exist.
Finally pull back the finite locally free algebra of
Corollary~\ref{cor:finite-flat-family} along $\rho$; arbitrary base
change preserves its rank and multiplication.
\end{proof}

This is an incidence compactification over the Grassmannian, with its
native congruence action, not a properness assertion for the quotient
stack by that action. A limiting Hilbert fibre may have extra or
nonreduced components; no classification of all such fibres is used.
The construction is independent of $h$: its target, polarization and
Hilbert polynomial were specified using the $h$-independent complete
quadric graph. The unmarked inverse in the companion reconstruction
paper recovers the underlying pencil, so the incidence construction is
compatible with that inverse rather than a replacement of it.

\section{Singularities of the universal power-zero schemes}
\label{sec:power-multiplier-v157}
There is an explicit further consequence for all the universal ideals,
which we record separately from the geometry of the full reciprocal
normalization fibre. Multiplier ideals are taken on the affine smooth
space $M$, and $I_{r+1}^{(a)}$ denotes symbolic power, with exponent
zero interpreted as the unit ideal.

\begin{corollary}[Multiplier ideals and threshold]
\label{cor:power-multiplier-v157}
For $1\leq p\leq nh$, put $w_r=(p-hr)_+$ and
$c_r=\binom{n-r+1}{2}$ for $0\leq r<n$. For every real $b\geq0$,
\begin{align}
 \mathcal J(M,b\cdot J_p)
   &=\bigcap_{r=0}^{n-1}
       I_{r+1}^{(\max\{0,\lfloor b w_r\rfloor+1-c_r\})},
                         \label{eq:power-multiplier-v157}\\
 \operatorname{lct}_M(J_p)
   &=\min_{\substack{0\leq r<n\\p>hr}}
                    \frac{\binom{n-r+1}{2}}{p-hr}.
                         \label{eq:power-lct-v157}
\end{align}
\end{corollary}
\begin{proof}
Blow up the affine origin and then the strict rank loci. This is the
affine complete-quadric resolution, whose charts are
\eqref{eq:CQ-chart-v157} with one additional scalar coordinate.
Let $E_0$ be the divisor over the origin. The same minor calculation
now includes $r=0$, and the total transform of $J_p$ is
$\sum_{r=0}^{n-1}w_rE_r$. The discrepancy along $E_r$ is $c_r-1$:
over the generic rank-$r$ point the earlier blow-ups are isomorphisms,
and the smooth rank-$r$ locus has codimension $c_r$; for $r=n-1$
this is the strict determinant divisor and the discrepancy is zero.
All boundary divisors have normal crossings, so this is a log resolution.

The defining pushforward formula for the multiplier ideal requires
$\operatorname{ord}_{E_r}(f)\geq\lfloor b w_r\rfloor+1-c_r$.
At the generic rank-$r$ point this valuation is the order in the
regular local ring with maximal ideal $I_{r+1}$, so the condition on
polynomials is exactly membership in the indicated symbolic power.
Intersecting the conditions proves the first formula; the first value
at which one becomes nontrivial proves the second. Equivalently this
is the specialization, via Theorem~\ref{thm:universal-power-v157}, of
the classical symmetric-determinantal multiplier-ideal calculation
\cite[\S2.4 and Theorem~4.8]{HenriquesVarbaro2016}.
\end{proof}

The ideal formula, graph identification and contact reconstruction are
statements about the entire universal coefficient space. They do not
supply the conductor of every multivariate common-divisor image or the
Hilbert function of every full $F_{g,k}(W)$. Those objects remain in
the preserved incidence theory; the present results connect their exact
quadratic sections to a native global boundary and its spectral schemes.

\section{Power coefficients and their twists}
\label{sec:power-conventions-v158}
For a vector space $V$, put $M(V)=\Sym^2V$, with congruence action
$A\mapsto gAg^{\mathsf t}$, and define
\begin{equation}\label{eq:envelope-presentation-v158}
 B_h(V)=\Sym(M(V))/\langle (v^2)^{h+1}:v\in V\rangle .
\end{equation}
The degree of $M(V)$ is one. The generating space of the ideal is the
Cartan inclusion $\Sym^{2h+2}V\hookrightarrow\Sym^{h+1}(\Sym^2V)$,
normalized by $v^{2h+2}\mapsto(v^2)^{h+1}$. This definition requires
neither a volume form nor a complementary space. It agrees with
$B_{n,h}$ in Theorem~\ref{thm:quadratic-section-v155} when $W=V^*$.

\begin{proposition}[The coefficient map and its twists]
\label{prop:coefficient-conventions-v158}
Let $n=\dim V$ and $N=nh$. The quotient multiplication $\mu_p$ gives
an equivariant injection
\begin{equation}\label{eq:coefficient-duals-v158}
 \mu_p^*:B_h(V)_p^*\hookrightarrow\Sym^p(M(V)^*),
 \qquad (\mu_p^*\varphi)(A)=\varphi(A^p).
\end{equation}
Its image is the full coefficient space defining $J_p$, without a
choice of apolar coordinates. The socle and perfect pairings are
\begin{equation}\label{eq:perfect-twist-v158}
 B_h(V)_N=(\det V)^{2h},\qquad
 B_h(V)_p^*\simeq B_h(V)_{N-p}\otimes(\det V^*)^{2h}.
\end{equation}
For a line $L$, $B_h(V\otimes L)_p=B_h(V)_p\otimes L^{2p}$,
compatibly with all products. On $\Pj(M(V))$, the $p$th power is a
section of $B_h(V)_p\otimes\OO(p)$, and its base ideal is the image
of $B_h(V)_p^*\otimes\OO(-p)\to\OO$.
\end{proposition}
\begin{proof}
A homogeneous polynomial of degree $p$ is identified with its symmetric
$p$-linear polarization, normalized so that its value on $A$ is its
value on $(A,\ldots,A)$. Dualizing the surjection $\mu_p$ gives exactly
\eqref{eq:coefficient-duals-v158}. Multinomial coefficients appearing
in a monomial basis express this normalization and are nonzero; no
dual representation is silently replaced by the original one.
Theorem~\ref{thm:quadratic-section-v155} identifies the quotient with
determinant apolarity. The determinant of $A:V^*\to V$ takes values
in $(\det V)^2$. Choosing a trivialization gives a scalar polynomial
whose inverse-system line has character $(\det V^*)^{2h}$; its dual
socle therefore has character $(\det V)^{2h}$. The Gorenstein pairing
gives the second formula. Each generator in $M(V)$ has scalar weight
two, proving the line-twist assertion directly from
\eqref{eq:envelope-presentation-v158}. The tautological line
$\OO(-1)\to M(V)\otimes\OO$ has $p$th product
$\OO(-p)\to B_h(V)_p\otimes\OO$, which proves the last assertion.
\end{proof}


\subsection{Graph schemes and classical inputs}
The projective graph argument uses the following precise statement.
It applies to the ambient space and also to the two-dimensional total
base of the degeneration computed below.

\begin{lemma}[Simultaneous graph with invertible twists]
\label{lem:simultaneous-graph-v158}
Let $X$ be an integral noetherian complex scheme and let
$W_i\otimes L_i^{-1}\to\OO_X$ be nonzero coherent systems of
sections, with finite-dimensional $W_i$ and invertible $L_i$.
Write $K_i$ for their image ideals, and let
$U=X\setminus\bigcup_i V(K_i)$. The scheme-theoretic graph closure
of $U\to\prod_i\Pj(W_i^*)$ is
$\operatorname{Bl}_{\prod_iK_i}X$ with its specified graph embedding.
This statement identifies the scheme and the maps, not only their
normalizations. Replacing the product ideal by an invertible twist
or a positive power does not change its blow-up over $X$.
\end{lemma}
\begin{proof}
The open $U$ is dense because all $K_i$ are nonzero and $X$ is integral.
After the Segre embedding the simultaneous map is given by all products
of one section from each system. On a trivializing affine open, these
products generate $K=\prod_iK_i$. The homogeneous coordinate algebra
of the graph is the image of the polynomial algebra on these
products in $\OO_X[t]$, namely the Rees algebra
$\bigoplus_{m\geq0}K^m t^m$. It is a subalgebra of the function-field
polynomial ring. Thus its Proj is exactly the schematic closure of
the graph, with no extra component. The Segre equations hold on the
dense graph and hence on this integral closure scheme, so the graph
lands in the product, not merely in its ambient projective space.
Local trivializations glue. An invertible ideal changes the degree-$m$
piece by its $m$th tensor power and does not change relative Proj;
a positive power replaces the Rees algebra by a Veronese subalgebra.
\end{proof}

For the power system, $L_i=\OO(p)$ and $W_i=B_h(V)_p^*$;
for the exterior system of order $q$, $L_i=\OO(q)$ and its sections
are the $q$-minors. The determinant system is invertible on the
integral ambient space. It may be omitted from the blow-up product,
although its zero divisor must still be retained in contact formulas.
Likewise the first-coordinate identity map has unit base ideal.
These observations specify the projective twists in the proof of
Theorem~\ref{thm:power-graph-v157}.

The precise classical scheme-level complete-quadric construction is
\cite[Theorem~6.3]{Vainsencher1982}; the graph of exterior-power maps
and its successive symmetric-rank blow-ups are stated in
\cite[Remark~2.5 and Construction~2.6]{Massarenti2020}.
The relative rank version is recalled in
\cite[Definition~2.5, Remark~2.6 and Theorem~2.14]{CasarottiCornianiMassarenti2023}.
Consequently smoothness, the flag description, and the orbit
stratification in Corollary~\ref{cor:native-strata-v157} are classical
consequences of the graph identification, not new classifications.

There is also a precise invariant-ideal predecessor to the universal
identity. Henriques--Varbaro \cite[\S2.4, Theorems~2.6--2.8]{HenriquesVarbaro2016}
recall Abeasis's classification for the symmetric coordinate ring,
its products of determinantal ideals and their integral closures,
and the passage between invariant generators and products indexed
by diagrams. Our coefficient-space comparison uses this same
multiplicity-free setting. Theorem~\ref{thm:universal-power-v157}
identifies a particular multiplication coefficient map with the
balanced determinantal product as an ordinary ideal; it is not
an independent classification of invariant ideals. The exact
coefficient identification in Proposition~\ref{prop:coefficient-conventions-v158}
explains which map is being specialized on nonreduced bases.
We assert no exhaustive historical nonanticipation of the displayed
identity in all invariant-ideal terminology. Similarly the
multiplier calculation in Corollary~\ref{cor:power-multiplier-v157} is the
substitution of its exponents in the symmetric determinantal formula
\cite[Theorem~4.8]{HenriquesVarbaro2016}, not an independent
multiplier-ideal theorem.

\section{The intrinsic source and its apolar envelope}
\label{sec:intrinsic-envelope-v158}
The source of a pencil and the coefficient complement in a divisor
normalization fibre are different spaces. We construct the native power
geometry from the first relation of the failure algebra, without taking
a coefficient section of that fibre. The construction produces a
canonical projective power diagram and an honest algebra bundle on the
intrinsic source projective space. The location of the algebra bundle
is essential: forgetting a scalar lift does not give a linear
representation of the projective general linear group.

\subsection{A construction from the first relation}
Let $A$ be an unmarked, ungraded complex algebra in the sharp pencil
image, and write $\mathfrak m$ for its maximal ideal. Set
\begin{equation}\label{eq:intrinsic-data-v158}
 E_A=\mathfrak m/\mathfrak m^2,\qquad
 H_A=\ker\{\Sym^dE_A\longrightarrow\mathfrak m^d\},
 \qquad d=n^2+2n-4.
\end{equation}
Here $n^2=\dim E_A$, so neither $n$ nor $d$ is extra marking.
The reduced support of the homogeneous ideal generated by $H_A$ is a
determinant hypersurface in $\Pj(E_A^*)$. Its successive reduced
singular loci give the rank-one Segre variety. The scheme of maximal
rank-one linear spaces has two components. Cancel the determinant
factor from $H_A$ and select the component singled out by coefficient
support ranks $(1,M)$ in Lemma~\ref{lem:coefficient-orientation-v146}.
Denote the resulting source projective space by $S_A$; with a pencil
presentation it is $\Pj(V)$, not $\Pj(\C^{n^2-1})$.

There is a direct jet-theoretic presentation which does not begin with
a linear source lift. Let $L$ be an ample generator of $\Pic(S_A)$,
and let $J^1L$ denote its first principal-parts bundle. Put
\begin{equation}\label{eq:jet-envelope-v158}
 \mathscr E_A=\mathcal Hom(J^1L,L),\qquad
 \mathscr B_{A,h}=
 \Sym(\Sym^2\mathscr E_A)/
     \langle(u^2)^{h+1}:u\text{ a local section of }\mathscr E_A\rangle.
\end{equation}
Scalar automorphisms of $L$ act trivially on $\mathscr E_A$, so this
bundle is independent of the representative of the ample generator.
For $p>0$ its projective coefficient space can be constructed as
$\Pj H^0(S_A,\mathscr B_{A,h,p}\otimes L^{-2p})$; changing $L$
changes the linear identification only by a scalar. Multiplication
supplies the projective power maps. Thus jets, a Cartan quotient, and
projectivization give the construction from the intrinsic source.

\begin{theorem}[Intrinsic projective power geometry]
\label{thm:intrinsic-envelope-v158}
On the groupoid of complex sharp pencil-failure algebras and their
ungraded isomorphisms, the following construction is functorial.
There is a source projective space $S_A$, a finite locally free graded
algebra $\mathscr B_{A,h}$ on $S_A$ for each $h\geq1$, and projective
spaces $P_{A,p}$ with a multiplication power diagram. In any source
lift $S_A=\Pj(V)$ they are
\begin{equation}\label{eq:normalized-envelope-v158}
 \mathscr B_{A,h}=
 \bigoplus_{p=0}^{nh}B_h(V)_p\otimes\OO_{\Pj(V)}(2p),\qquad
 P_{A,p}=\Pj(B_h(V)_p).
\end{equation}
The nondegenerate power graph on $P_{A,1}$ is the native
$\mathrm{CQ}(V)$, canonically as a projective diagram, and every
rank-$r$ graph is the relative complete-quadric space of
Theorem~\ref{thm:all-rank-graph-v158}. Over a general base the projective spaces mean the associated
projective schemes, which need not admit a global linear lift.
These constructions use only
\eqref{eq:intrinsic-data-v158}, determinant cancellation, the oriented
ruling, and \eqref{eq:envelope-presentation-v158}; no pencil coefficient
line or reciprocal-fibre splitting is needed to construct the ambient
graph. Extracting that line afterwards supplies the intrinsic pencil
line $\Lambda_A\subset P_{A,1}$.

All source-lift choices cancel. Once an oriented projective-source
bundle is given, the same associated constructions descend over any
complex base, including a nonreduced base, and commute with arbitrary
base change. This last assertion concerns the associated objects; it
does not assert that reduction or graph closure commutes with arbitrary
base change on unrelated degenerating families.
\end{theorem}
\begin{proof}
The first-relation construction is preserved by every ungraded
isomorphism: terms of degree at least two in a change of generators
cannot change the degree-$d$ relation modulo $\mathfrak m^{d+1}$.
The determinant and its rank-one locus are therefore intrinsic.
Lemma~\ref{lem:rank-one-fano-scheme} identifies the two ruling
components as schemes; coefficient-support asymmetry orients them.
These are the steps of the local inverse before the final recovery
of the pencil coefficient line. Consequently an isomorphism of
algebras induces a unique projective isomorphism of the selected
source spaces, and compositions give compositions. This establishes
the source functor without first choosing a congruence representative
of the pencil.

To verify the jet presentation, take $S_A=\Pj(V)$ and $L=\OO(1)$.
The evaluation map $H^0(L)\otimes\OO\to J^1L$ is an isomorphism:
a linear form whose value and first derivatives vanish at a point is
zero, and both bundles have rank $n$. Hence
$\mathscr E_A=V\otimes\OO(1)$, and the Cartan quotient in
\eqref{eq:jet-envelope-v158} is exactly
\eqref{eq:normalized-envelope-v158}. Its degree-$p$ twist by
$L^{-2p}$ is constant, with global section space $B_h(V)_p$;
this proves the stated projective coefficient construction. Relative
principal parts commute with base change on the locally trivial smooth
projective-space bundle. Replacing $L$ by its tensor product with a
line from the base cancels in $\mathcal Hom(J^1L,L)$.

Fix a source lift temporarily. A scalar $c$ acts on $B_h(V)_p$ by
$c^{2p}$ and on $\OO_{\Pj(V)}(2p)$ by $c^{-2p}$. Hence every summand
in \eqref{eq:normalized-envelope-v158}, together with its multiplication,
is naturally $\PGL(V)$-equivariant. If two linear lifts of a projective
isomorphism differ by a scalar, their induced algebra-bundle maps
are equal. This also proves the cocycle identity on three overlaps;
in a family, replacing $V$ by $V\otimes L$ cancels $L^{2p}$ against
$L^{-2p}$. Thus the algebra bundle descends independently of all
lifts. For a projective-source bundle, take fppf local frames of its
$\PGL_n$-torsor; the scalar-trivial action just proved supplies the
cocycle descent even when no global vector-bundle lift exists. The same scalar acts trivially on each $P_{A,p}$, and power
maps commute with scalar changes and congruence. They therefore
define a projective diagram on the source torsor, with the same
cocycle compatibility. Its first projective space can alternatively
be described as the parameter space of quadric hypersurfaces in the
dual projective space $S_A^\vee$.

Apply Theorem~\ref{thm:power-graph-v157} to identify the nondegenerate
graph. The rank-$r$ identification is proved independently in
Theorem~\ref{thm:all-rank-graph-v158}. Equivariance makes both
identifications independent of the source lift. Finally the residual
left coefficient line, by the multiplicity-one exterior-duality map
in the local inverse, gives the Pluecker point of a two-dimensional
subspace in $\Sym^2V$. Scalar changes do not change its projective
line. This is $\Lambda_A$. Associated bundles, their specified
equations, and their multiplication commute with any base change;
no claim about arbitrary restricted graph closures is required.
\end{proof}

\begin{remark}[Exact intrinsic status]
\label{rem:envelope-status-v158}
The envelope is a functor constructed from multiplication, not a claimed
subalgebra or preferred quotient of $A$. It is an algebra bundle on
$S_A$, not an untwisted algebra canonically based at $\Spec\C$.
This distinction records, rather than suppresses, the scalar ambiguity:
for $n\geq3$ the central element $\zeta I\in\mathrm{SL}(V)$ acts on
$B_h(V)_1=\Sym^2V$ by $\zeta^2$, which need not be one although it
acts trivially projectively. Thus the usual linear degree-one action
does not descend to $\PGL(V)$. The twist in
\eqref{eq:normalized-envelope-v158} removes exactly this obstruction.
There is no selection of an $n$-plane inside the auxiliary complement
of dimension $n^2-1$. The latter still defines its own, different
coefficient-section algebra in Corollary~\ref{cor:boundary-size-v155}.
\end{remark}

\section{Power graphs of every rank}
\label{sec:all-rank-graph-v158}
For $1\leq r\leq n$, let $P^{(r)}\subset\Pj(\Sym^2V)$ be the
rank-exactly-$r$ locus. Define $\mathcal G_{r,h}$ as the
scheme-theoretic closure of
\begin{equation}\label{eq:rank-graph-definition-v158}
 P^{(r)}\longrightarrow
 \Pj(\Sym^2V)\times\prod_{p=1}^{hr}\Pj(B_h(V)_p),
 \qquad[A]\longmapsto([A],[A],\ldots,[A^{hr}]).
\end{equation}
Keeping the final power is essential when $r<n$. Write $\mathcal U$
for the tautological rank-$r$ bundle on $\Gr(r,V)$ and
$\mathrm{CQ}(\mathcal U)$ for its relative complete-quadric space.
For rank one this means $\Pj(\Sym^2\mathcal U)=\Gr(1,V)$.

\begin{theorem}[The rank-stratified power graph]
\label{thm:all-rank-graph-v158}
There is a canonical isomorphism
\begin{equation}\label{eq:rank-graph-isom-v158}
 \mathcal G_{r,h}\simeq\mathrm{CQ}(\mathcal U)
                  \longrightarrow\Gr(r,V).
\end{equation}
The final power in \eqref{eq:rank-graph-definition-v158} recovers the
image $r$-plane by the degree-$2h$ Pluecker embedding. In particular
the graph is smooth and projective, independently of $h$. For its
relative boundary divisors $E_i$, $1\leq i<r$, its power base ideals
have total transforms
\begin{equation}\label{eq:rank-boundary-v158}
 J_p\OO_{\mathcal G_{r,h}}=
 \OO_{\mathcal G_{r,h}}\left(-\sum_{i=1}^{r-1}(p-hi)_+E_i\right),
 \qquad 0\leq p\leq hr.
\end{equation}
The projective power maps in this formula have their common divisors
removed. All powers $p>hr$ vanish on the rank-$r$ locus.
\end{theorem}
\begin{proof}
For $U\subset V$, the map $B_h(U)\to B_h(V)$ from the presentation
\eqref{eq:envelope-presentation-v158} is injective: any linear projection
$V\to U$ splitting the inclusion induces a retraction on the quotient
algebras. The injection itself is independent of that projection.
Locally this gives subbundles $B_h(\mathcal U)_p\subset B_h(V)_p$.
If $A$ has image $U$ and rank $r$, it is a nonsingular element of
$\Sym^2U$, and $A^{hr}$ generates the socle
$B_h(U)_{hr}=(\det U)^{2h}$. As $U$ varies these lines give the
homogeneous embedding
\[
 \Gr(r,V)\hookrightarrow\Pj(\mathbb S_{(2h)^r}V)
                         \hookrightarrow\Pj(B_h(V)_{hr}).
\]
Indeed the line at the coordinate $r$-plane is a highest-weight line
of weight $(2h)^r$; its orbit map is the Pluecker embedding followed
by its degree-$2h$ linear system. In particular this is a closed
immersion, not just a recovery of the image plane on closed points.

Over the recovered Grassmannian, the remaining graph is precisely the
relative simultaneous power graph for $B_h(\mathcal U)$. By
Theorem~\ref{thm:power-graph-v157} on each frame chart, it is
$\mathrm{CQ}(\mathcal U)$. The relative graph embedding, the subbundle
injections above, and the final-power Grassmannian embedding together
give a closed immersion into the product in
\eqref{eq:rank-graph-definition-v158}. Its dense open is the graph
of $P^{(r)}$. Both its source and that graph closure are integral,
so equality of their scheme-theoretic images proves
\eqref{eq:rank-graph-isom-v158} without a normalization step.
Applying \eqref{eq:power-boundary-v157} in rank $r$ proves the
transform formula. Local split inclusion of the coefficient bundles
ensures that ambient and relative power ideals agree. The vanishing
assertion follows from the socle degree of $B_h(U)$.
\end{proof}

The varieties of complete singular quadrics are classical: see
\cite[Definition~2.5, Remark~2.6 and Theorem~2.14]{CasarottiCornianiMassarenti2023}.
The point of the theorem is the exact recovery of their image-plane
projection and entire scheme structure from the power maps in one
fixed algebra $B_h(V)$, simultaneously for every $r$. The classical
orbit stratification is background, not an additional novelty claim.

\section{Contacts and minimal indices of arbitrary pencils}
\label{sec:singular-pencil-v158}
Let $R=\langle A_0,A_1\rangle\subset\Sym^2V$ be a genuine pencil,
$\Lambda=\Pj(R)$, and $r$ its normal rank. Its tautological symmetric
map is $A:V^*\otimes\OO_\Lambda\to V\otimes\OO_\Lambda(1)$.
Write
\[
 K=\ker A,\qquad U=K^\perp\subset V\otimes\OO_\Lambda.
\]
The kernel is saturated and hence a subbundle on the smooth curve;
$U$ is the saturation of the generic image. Thus $A$ is a generically
nonsingular section of $\Sym^2U\otimes\OO_\Lambda(1)$.
The map on the rank-$r$ open lifts uniquely to
$\widetilde\Lambda\to\mathcal G_{r,h}$ by properness. Its
Grassmannian projection is the map defined by $U$, and its source is
still $\Lambda$. Let $D_i$ be its contact divisor with $E_i$.

\begin{lemma}[Symmetric diagonalization with a kernel]
\label{lem:symmetric-dvr-v158}
Over $\C[[t]]$, a symmetric matrix of rank $r$ over the fraction field
is congruent to $\diag(t^{e_1},\ldots,t^{e_r},0,\ldots,0)$, with
$e_1\leq\cdots\leq e_r$. For the local matrix of a genuine pencil,
$e_1=0$. The exponents are its nonzero Smith exponents.
\end{lemma}
\begin{proof}
Choose a matrix entry of smallest valuation. If a diagonal entry
has that valuation, it is a pivot. Otherwise an off-diagonal entry
does, and the value of the quadratic form on $e_i+c e_j$ has the
same valuation for some $c\in\C$: its leading coefficient is a
nonzero polynomial in $c$, since $2$ is invertible. Extend this
primitive vector to a basis. The pivot divides every entry in its
row, so symmetric row-and-column elimination clears the row and
column over the ring. All remaining entries have valuation at least
that of the pivot. Repeat until the remaining block is zero.
Units have square roots in $\C[[t]]$ and can be absorbed into the
basis. Invertible congruence preserves determinantal ideals, which
identify the resulting exponents with Smith exponents. A genuine
pencil has a nonzero matrix at every projective parameter, giving a
unit entry and $e_1=0$. This proof supplies the local congruence step
rather than relying on left-right Smith form alone; compare the
classical pencil congruence treatment in \cite{Thompson1991}.
\end{proof}

For $j\geq0$ define the finite multiplication-syzygy map
\begin{equation}\label{eq:toeplitz-map-v158}
 C_j(R):V^*\otimes\Sym^jR^*\longrightarrow
                         V\otimes\Sym^{j+1}R^*,\qquad q\longmapsto A q,
 \qquad\kappa_j=\dim\ker C_j(R).
\end{equation}
Set $\kappa_{-1}=\kappa_{-2}=0$. Under a basis of $R$ this is the block
Toeplitz map with rows $A_0q_0$, $A_0q_i+A_1q_{i-1}$, and $A_1q_j$;
its definition is independent of that basis.

\begin{theorem}[Complete power-and-syzygy data]
\label{thm:singular-complete-data-v158}
For every complex symmetric pencil, including every singular pencil,
the following formulas recover its complete congruence data.
If $\varepsilon_1,\ldots,\varepsilon_c$ are its minimal indices,
where $c=n-r$, then
\begin{align}
 K&\simeq\bigoplus_{a=1}^c\OO_\Lambda(-\varepsilon_a),
 &\kappa_j&=\sum_{a=1}^c(j-\varepsilon_a+1)_+,
                       \label{eq:kernel-splitting-v158}\\
 \#\{a:\varepsilon_a=j\}
   &=\kappa_j-2\kappa_{j-1}+\kappa_{j-2}.
                       \label{eq:minimal-differences-v158}
\end{align}
It is enough to use $0\leq j\leq\lfloor(n-1)/2\rfloor$.
At $x\in\Lambda$, let $0=e_1(x)\leq\cdots\leq e_r(x)$ be the
nonzero local Smith exponents. With $v_0=0$ and
$v_p(x)=\length(\OO_{\Lambda,x}/J_p)$ for $p\leq hr$, one has
\begin{align}
 v_{hq+s}(x)&=h\sum_{i=1}^{q}e_i(x)+s e_{q+1}(x),
                         \label{eq:singular-valuations-v158}\\
 e_i(x)&=\frac{v_{hi}(x)-v_{h(i-1)}(x)}h,
 &\operatorname{mult}_xD_i&=e_{i+1}(x)-e_i(x).
                         \label{eq:singular-contacts-v158}
\end{align}
In the first formula $s=0$ at $q=r$ and the last term is omitted.
The positive exponents at each supported point, together with the
minimal indices, determine the pencil up to congruence and projective
reparametrization of $\Lambda$. The spectral datum includes the
positions of the points, not only their local partitions.

There is also a global conservation law:
\begin{equation}\label{eq:conservation-v158}
 \sum_{i=1}^{r-1}(r-i)\deg D_i
       +2\deg\bigl(\Lambda\longrightarrow\Gr(r,V)\bigr)
       =r,\qquad
 \deg\bigl(\Lambda\longrightarrow\Gr(r,V)\bigr)
       =\sum_a\varepsilon_a.
\end{equation}
The Grassmannian degree uses its Pluecker polarization. The first
summand is the dimension of the regular summand in the Kronecker
congruence form. Every object and number in the theorem is invariant
under the ungraded failure-algebra isomorphisms of
Theorem~\ref{thm:intrinsic-envelope-v158}.
\end{theorem}
\begin{proof}
We use the classical complex symmetric Kronecker congruence theorem
\cite{Thompson1991}. A singular block of minimal index $\varepsilon$
has size $2\varepsilon+1$ and form
\[
 S_\varepsilon(s,t)=
 \begin{pmatrix}0&L_\varepsilon(s,t)^{\mathsf t}\\
                 L_\varepsilon(s,t)&0\end{pmatrix},\qquad
 L_\varepsilon=
 \begin{pmatrix}s&t& &0\\ &s&t& \\ &&\ddots&t\end{pmatrix},
\]
where $L_\varepsilon$ has $\varepsilon$ rows and $\varepsilon+1$
columns; $S_0$ is a zero $1\times1$ block. Its kernel is generated
by $((-t)^\varepsilon,s(-t)^{\varepsilon-1},\ldots,s^\varepsilon)$
in the first block. These components have no common projective zero,
so it gives $\OO(-\varepsilon)$ as a kernel subbundle. A regular
block has zero kernel as a sheaf. Direct sums prove the first formula
in \eqref{eq:kernel-splitting-v158}. Taking sections after twist by
$\OO(j)$ identifies them with \eqref{eq:toeplitz-map-v158}, proving
the second formula. The second finite difference of
$(j-\varepsilon+1)_+$ is one exactly at $j=\varepsilon$ and zero
otherwise. Each singular block uses $2\varepsilon+1\leq n$
dimensions, proving the stated finite bound. The canonical-form
classification is used here as a classical input, not reproved or
claimed as new.

Locally, symmetric diagonalization factors the matrix through the
saturated $r$-plane $U$ with an $r\times r$ symmetric block. Theorem
\ref{thm:universal-power-v157}, or its rank-$r$ version via the split
inclusion $B_h(U)\hookrightarrow B_h(V)$, yields
\eqref{eq:singular-valuations-v158}. Finite differences recover the
$e_i$. In the relative complete-quadric chart the successive ratios
of diagonal entries have valuations $e_{i+1}-e_i$, proving the contact
formula. The torsion of the pencil cokernel has precisely the positive
Smith exponents. The regular blocks of the classical congruence form
are specified by these elementary divisors, and its singular blocks
by the minimal indices; over $\C$ there are no additional signature
parameters. Retaining the supported points on $\Lambda$ makes the
statement compatible with arbitrary projective changes of the pencil
basis.

The exact sequence $0\to U\to V\otimes\OO\to K^*\to0$ gives
$\deg U=-\sum\varepsilon_a$, so the Grassmannian degree is
$\sum\varepsilon_a$. The determinant of the induced symmetric map
$U^*\to U(1)$ is a nonzero section of
$(\det U)^2\otimes\OO(r)$; its zero divisor has degree
$r-2\sum\varepsilon_a$. At each point its order is
$\sum e_i=\sum_{i<r}(r-i)\operatorname{mult}_xD_i$.
This proves \eqref{eq:conservation-v158}. The block-size identity
$n=n_{\mathrm{reg}}+\sum(2\varepsilon_a+1)$ identifies the first
summand with $n_{\mathrm{reg}}$. Finally the intrinsic first relation
extracts $\Lambda_A$ on its source and preserves all the maps involved;
changes of source lift or pencil basis give isomorphic kernel maps
and equal ideals, hence the asserted invariance.
\end{proof}

\begin{example}[Power contacts alone miss a singular invariant]
\label{ex:singular-separation-v158}
The size-six pencils $S_0\oplus S_2$ and $S_1\oplus S_1$ both have
constant rank four and no spectral torsion. All their power ideals
are the unit ideal for $p\leq4h$ and zero for $p>4h$, and every
contact divisor is empty. Their kernel splittings are respectively
$\OO\oplus\OO(-2)$ and $\OO(-1)^2$. Thus their first three
syzygy nullities are $(1,2,4)$ and $(0,2,4)$, distinguishing the two
congruence types. Both have Grassmannian degree two, as required by
\eqref{eq:conservation-v158}; the full splitting, rather than just
its degree, is necessary.
\end{example}

\begin{proposition}[Finite closed tests in families]
\label{prop:syzygy-family-v158}
For $j\geq0$ and $0\leq a\leq n(j+1)$, on the pencil
Grassmannian each condition $\kappa_j\geq a$ is the
closed determinantal condition on the bundle map $C_j$ given by minors
of size $n(j+1)-a+1$. These conditions commute with arbitrary base
change as schemes. Their fibrewise ranks need not be constant, so no
unqualified assertion that the kernel bundles commute with every
base change is made. Along a specialization over an integral curve,
$\kappa_j$ can only increase.
\end{proposition}
\begin{proof}
Use the tautological rank-two subbundle in
\eqref{eq:toeplitz-map-v158}. Its source has rank $n(j+1)$ and its
target rank $n(j+2)$. Nullity at least $a$ is equivalent to rank at
most $n(j+1)-a$, giving the stated minors. The ideals are polynomial
in the bundle map and therefore commute with arbitrary substitution.
Upper semicontinuity of nullity, or specialization of these minors,
gives the last assertion.
\end{proof}

\section{A computed boundary of the pencil incidence space}
\label{sec:computed-boundary-v158}
We give the incidence compactification a precise universal property and
compute a family of its boundary fibres in every dimension. The new
boundary is selected by one power base ideal, not by pulling an
unrelated algebra back to a parameter space.

\begin{proposition}[The normal main-component flattening property]
\label{prop:flattening-property-v158}
Let $G^\circ\subset G=\Gr(2,\Sym^2V)$ and
$\widehat G\to G$ be as in Proposition~\ref{prop:pencil-compactification-v157}.
Among normal integral schemes $T$ with a dominant morphism $f:T\to G$
for which the closure of the lifted pencil family over $f^{-1}(G^\circ)$
is an embedded flat family in $\mathrm{CQ}(V)\times T$ with Hilbert
polynomial $\binom n2 m+1$, there is a unique morphism $T\to\widehat G$
over $G$, and that family is the pullback of its universal family.
Consequently $\widehat G$ is the minimal normal dominant modification
with this specified embedded-flat-closure property.
\end{proposition}
\begin{proof}
The flat embedded family gives a morphism to the indicated Hilbert
scheme. Together with $f$ it agrees, over a dense open of $T$, with
the graph defining the reduced graph closure. Since $T$ is integral,
its morphism factors through this closed graph scheme: each equation
vanishing on the dense open vanishes on $T$. The induced morphism is
dominant. A dominant morphism from a normal integral scheme to an
integral scheme factors uniquely through the normalization: locally,
an element integral over the target ring becomes integral over the
source ring and lies in its fraction field, hence belongs to the
normal source ring. Thus the morphism factors uniquely through
$\widehat G$. The Hilbert-scheme representing property identifies
the families. Conversely its pulled-back universal family is flat;
flatness over an integral base makes it the schematic closure of its
restriction to the dense open. Uniqueness follows also from separatedness.
\end{proof}

The hypothesis of dominance is part of this property; this proposition
does not assert that normalization commutes with arbitrary base change.
The property is a standard consequence of the Hilbert functor and
normalization. Its usefulness here is that the following calculation
identifies a nontrivial fibre of the resulting modification.

Assume $n\geq3$. Choose distinct nonzero numbers $a_3,\ldots,a_n$ and put
$D=\Spec\C[[\tau]]$. On $\Pj^1_D$, consider the symmetric pencil
\begin{equation}\label{eq:collision-family-v158}
 A_\tau(s,t)=\diag\bigl(s,s+\tau t,
                           t-a_3s,\ldots,t-a_ns\bigr).
\end{equation}
For $n=3$ there is just one number $a_3$. Write $X=\Pj^1_D$ and let
$o$ be the point $s=\tau=0$ in the special fibre. Its generic pencil
has $n$ distinct rank-$(n-1)$ spectral points, so defines a point of
$G^\circ$.

\begin{theorem}[A transverse corank-two tail and its power maps]
\label{thm:hilbert-tail-v158}
The closure of the generic lifted curve of
\eqref{eq:collision-family-v158} in $\mathrm{CQ}(V)\times D$ is
\begin{equation}\label{eq:tail-blowup-v158}
                         Z=\operatorname{Bl}_o X.
\end{equation}
It is flat over $D$, and its special fibre is the reduced nodal union
$C\cup F$, where $C\simeq\Pj^1$ is the lifted special pencil and
$F\simeq\Pj^1$ is an exceptional tail. There are no embedded components.
If $H_q$ is the exterior-power hyperplane bundle of order $q$, then
\begin{equation}\label{eq:tail-degrees-v158}
 (\deg_C H_q,\deg_F H_q)=
 \begin{cases}(q,0),&1\leq q\leq n-2,\\
              (n-2,1),&q=n-1.
 \end{cases}
\end{equation}
Thus the special fibre has the same Hilbert polynomial
$\binom n2m+1$ as the generic lifted line.

The tail lies in $E_{n-2}$, its two intersections with $E_{n-1}$
are distinct and reduced, and its attachment to $C$ is neither of
those two points. For any $h\geq1$, let $L_p$ be the hyperplane
bundle of the resolved $p$th power map. Then
\begin{equation}\label{eq:tail-power-degree-v158}
 \deg_F L_p=
 \begin{cases}
 0,&p\leq h(n-2),\\
 \min\{b,2h-b\},&p=h(n-2)+b,\quad 0\leq b\leq2h.
 \end{cases}
\end{equation}
When this degree is positive, the restricted map on $F$ is the
complete Veronese map of that degree into its linear span.
Most directly, in the local coordinate $z=s/t$ at $o$,
\begin{equation}\label{eq:one-power-centre-v158}
 J_{h(n-2)+1}(A_\tau)=(z,\tau).
\end{equation}
Hence a single multiplication power ideal selects the blow-up centre
and its exceptional component.
\end{theorem}
\begin{proof}
The rank is at least $n-2$ on $X$. In a neighbourhood of $o$ the last
$n-2$ diagonal entries are units; all other spectral points have rank
$n-1$ and are pairwise distinct. Consequently the extended minor ideals
$I_q$ are the unit ideal for $q\leq n-2$, while
\[
 I_{n-1}=(z,z+\tau)=(z,\tau),\qquad
 I_n=(u z(z+\tau))
\]
near $o$, with $u$ a unit. Away from $o$, $I_{n-1}$ is a unit ideal;
$I_n$ is everywhere invertible. The simultaneous exterior-power graph
on the integral surface $X$ is therefore the blow-up of $(z,\tau)$.
This is an application of the graph-by-generators lemma, not a claim
that arbitrary base change of a blow-up is a blow-up. Since the
universal family of pencil lines embeds $X$ into $\Pj(\Sym^2V)\times D$,
this graph embeds into $\mathrm{CQ}(V)\times D$ and is the required
scheme-theoretic closure.

On the chart with $\tau=zv$, the special fibre has equation $zv=0$;
on the other chart, $z=\tau w$, it has equation $\tau=0$.
These charts prove flatness over the DVR, multiplicity one of both
components, their transverse intersection, and absence of embedded
components. The map on $C$ is the unique lift of the special pencil.
The $q$-minors of that pencil have degree $q$. They have no common
factor for $q\leq n-2$, and precisely one factor $s$ for $q=n-1$.
This proves their degrees on $C$. On $F$, all lower maps are constant,
and the $(n-1)$st exterior map is linear. This gives
\eqref{eq:tail-degrees-v158}. For $L=\bigotimes H_q$, the degrees
on $C$ and $F$ add to $\binom n2$. The normalization sequence of
$C\cup F$ subtracts one point, so its Euler characteristic is one
and its Hilbert polynomial is as stated. Thus it is the Hilbert
limit used by $\widehat G$; equivalently the generic morphism from
$D$ extends to $\widehat G$ by properness, with this universal fibre.

Use homogeneous coordinates $[\alpha:\beta]=[z:\tau]$ on $F$.
After removing the exceptional factor, the residual quadratic form
on the two-dimensional kernel is
$\diag(\alpha,\alpha+\beta)$. This is a line in
$\mathrm{CQ}(\C^2)=\Pj(\Sym^2\C^2)$, and gives the claimed image
of $F$ in the fibre of $E_{n-2}$. Its determinant $\alpha(\alpha+\beta)$ has two
simple zeros $[0:1]$ and $[-1:1]$, which are its intersections with
$E_{n-1}$. The main component attaches at $[1:0]$, where this residual
form is nondegenerate.

For clarity, the whole list of extended power ideals near $o$ is
\begin{equation}\label{eq:tail-all-ideals-v158}
 J_p=
 \begin{cases}
 \OO,&p\leq h(n-2),\\
 (z,\tau)^b,&p=h(n-2)+b,\ 0\leq b\leq h,\\
 (z,\tau)^{2h-b}[z(z+\tau)]^{b-h},
       &p=h(n-2)+b,\ h\leq b\leq2h,
 \end{cases}
\end{equation}
up to an invertible factor. These are actual ideals, by
Theorem~\ref{thm:universal-power-v157}. In particular $b=1$ gives
\eqref{eq:one-power-centre-v158} for every $h$.

It remains to check the linear system on the tail, not only its
degree. On diagonal matrices the coefficient polynomials of a $p$th
power span exactly the monomials of total degree $p$ in the diagonal
entries with each exponent at most $h$. For containment, use the
torus-weight bound in Lemma~\ref{lem:bounded-row-v157}; for the reverse
inclusion, take the coefficients of the independent diagonal entries
of $X$ in $\det(X+uA)^h$. Thus, for $p=h(n-2)+b$, the leading normal
coefficients use exponent $h$ in each of the $n-2$ invertible entries
and span
\[
  \alpha^j(\alpha+\beta)^{b-j},\qquad
           \max\{0,b-h\}\leq j\leq\min\{b,h\}.
\]
If $b\leq h$ these span every
binary form of degree $b$. If $b>h$, remove the common factor
$[\alpha(\alpha+\beta)]^{b-h}$; the remaining monomials span every binary form of
degree $2h-b$. Lower powers have nonzero constant normal term and
are constant on $F$. This proves the complete-linear-system assertion
and \eqref{eq:tail-power-degree-v158}.
\end{proof}

\begin{remark}[What the failure family contributes]
\label{rem:tail-integration-v158}
Along \eqref{eq:collision-family-v158}, the original finite failure
algebras form the finite locally free family already constructed in
Corollary~\ref{cor:finite-flat-family}. Its pullback to $F$ is constant;
we do not count that pullback as a new interaction. Instead, its first
relation canonically recovers the source and the pencil line, and the
source-normalized envelope has the total-base ideal
\eqref{eq:one-power-centre-v158}. Resolving that multiplication ideal
creates $F$, and the other products prescribe all the linear systems
\eqref{eq:tail-power-degree-v158} on it. A degeneration arc, not only
its closed failure algebra, supplies this total-base ideal. Thus the
calculation distinguishes information in the family from information
in the special point and supplies an explicit modular Hilbert limit.
\end{remark}

\section{Universal property of the intrinsic envelope}
\label{sec:universal-envelope-v159}
The intrinsic envelope is a construction from the multiplication table,
not an identification of an auxiliary coefficient section with the
original algebra. Its universal property can be stated before choosing
a linear source or recovering a representative of the pencil.

\begin{proposition}[The source-normalized nilpotence envelope]
\label{prop:universal-envelope-v159}
Let $S_A$ and $\mathscr E_A=\mathcal Hom(J^1L,L)$ be as in
Theorem~\ref{thm:intrinsic-envelope-v158}. The algebra
$\mathscr B_{A,h}$ represents, on $S_A$, graded commutative algebras
$C$ equipped with a linear map $\Sym^2\mathscr E_A\to C_1$ for which
\[
                  \bigl(\text{image of }u^2\bigr)^{h+1}=0
\]
for every local section $u$, also after base change. More precisely,
giving such a map is equivalent to giving a unique graded algebra map
$\mathscr B_{A,h}\to C$. The degree-zero part maps by the structure
map. This property, its representing algebra, and its projectivized
power maps descend along the oriented source torsor. The case $h=1$
requires no additional integer choice and gives the same native graph
as every $h\geq1$.
\end{proposition}
\begin{proof}
The symmetric algebra is free on its degree-one module. In
characteristic zero the span of $(u^2)^{h+1}$ is exactly the image of
\[
 \Sym^{2h+2}\mathscr E_A\longrightarrow
                   \Sym^{h+1}(\Sym^2\mathscr E_A),\qquad
 u^{2h+2}\longmapsto(u^2)^{h+1}.
\]
Polarization verifies this assertion on every trivializing open;
the required denominators are nonzero integers, hence units on every
complex base. Factoring a free-algebra map through the ideal generated
by this image is equivalent to the displayed identities. This proves
existence and uniqueness and their compatibility with base change.

On $\Pj(V)$, evaluation of linear sections identifies
$J^1\OO(1)$ with $V^*\otimes\OO$. Thus
$\mathscr E_A=V\otimes\OO(1)$. A scalar on $V$ acts oppositely on
$\OO(1)$, so it acts trivially on $\mathscr E_A$, the Cartan map,
and the quotient. This proves descent without choosing a global
$\OO(1)$ on a nonsplit source bundle. The scalar action on each
untwisted coefficient space is scalar, and hence also disappears on
its projectivization. The graph assertions follow from
Theorems~\ref{thm:power-graph-v157} and
\ref{thm:all-rank-graph-v158}.
\end{proof}

This universal property distinguishes an algebra bundle constructed
from the intrinsic source from a preferred linear subquotient of $A$.
The former is proved; the latter is neither needed nor asserted.
The source is extracted by the determinant and oriented-ruling steps
of the first-relation inverse, before its final pencil-coefficient step.

\section{Ordinary collision strata and their Hilbert fibres}
\label{sec:ordinary-collisions-v159}
We compute simultaneous collisions of arbitrary multiplicities, and
allow higher-order perturbations of the total family. Besides describing
its reduced components, we determine every resolved power system and
exhibit positive-dimensional families in a fibre of the pencil
incidence modification. The calculation uses the ordinary power ideals
on the two-dimensional total base.

\subsection{A local ideal calculation}
\begin{lemma}[Distinct tangent factors]
\label{lem:tangent-factors-v159}
Let $\ell_1,\ldots,\ell_c$ be pairwise nonproportional linear forms
in two variables. For $0\leq j<c$, the products of $j$ distinct
$\ell_i$ span all binary forms of degree $j$. Consequently, if
$R$ is a regular local complex surface ring with maximal ideal
$\mathfrak n=(z,\tau)$, and a symmetric $c\times c$ matrix $S$ satisfies
\[
 S=zC+\tau D+O(\mathfrak n^2),
 \qquad C\text{ nonsingular},\qquad
 \det(\alpha C+\beta D)\text{ squarefree},
\]
then $I_j(S)=\mathfrak n^j$ for $j<c$, while
$I_c(S)=(\det S)$.
\end{lemma}
\begin{proof}
Choose $j+1$ of the linear forms. Their products obtained by omitting
one factor are linearly independent: evaluation at the zero of the
omitted factor kills all other products and not that one. Their
number is $j+1$, the dimension of the required binary-form space.
For $j=0$ the assertion is immediate.

The endomorphism $C^{-1}D$ is self-adjoint for $C$ and has distinct
eigenvalues. Its eigenspaces are orthogonal and none is isotropic,
since otherwise $C$ would be degenerate. A constant congruence
therefore makes $\alpha C+\beta D$ diagonal with distinct linear
entries. Its $j$-minors span $\Sym^j(\mathfrak n/\mathfrak n^2)$.
The entries of $S$ lie in $\mathfrak n$, so
$I_j(S)\subset\mathfrak n^j$ and
$I_j(S)+\mathfrak n^{j+1}=\mathfrak n^j$. Nakayama's lemma applied
to $\mathfrak n^j/I_j(S)$ proves equality. The argument may be made
after completion and descended by faithful flatness. The determinant
assertion is its definition.
\end{proof}

\begin{lemma}[Block products and the first normal term]
\label{lem:block-products-v159}
For $V=U\oplus W$, multiplication gives an injection
\[
                 B_h(U)\otimes B_h(W)\longrightarrow B_h(V).
\]
If $\dim U=r$, $Q\in\Sym^2U$ is nonsingular, and
$S\in\mathfrak n\otimes\Sym^2W$ has a generically nonsingular linear term $S_1$, the first
nonzero normal term of $(Q+S)^p$, for $hr<p\leq h\dim V$, is
\[
                  \binom p{hr}Q^{hr}S_1^{p-hr}.
\]
For $p\leq hr$ its order-zero term $Q^p$ is nonzero. These statements
remain valid with a nonsingular formal unit block and after a formal
congruence, with the corresponding invertible changes of coefficients.
\end{lemma}
\begin{proof}
The Cartan relations on each summand vanish in $B_h(V)$, so the
algebra map is defined. Its source is Artin Gorenstein. Its socle is
the product of the two socles, and that product is nonzero in $B_h(V)$:
in determinant apolarity this is the nonzero mixed top derivative
of $\det(X_U)^h\det(X_W)^h$, obtained by restricting to block
diagonal matrices. A nonzero ideal in an Artin local Gorenstein
algebra meets its socle: successive multiplication of a nonzero
ideal by the maximal ideal ends in a nonzero subspace annihilated
by that maximal ideal. The kernel is therefore zero.

Expand $(Q+S)^p$ in the injected tensor product. Powers of $Q$
above $hr$ vanish, whereas $Q^{hr}$ is a nonzero socle generator.
For $p>hr$ the term with the fewest factors of $S$ consequently
has precisely $p-hr$ such factors and is the displayed term.
The unit block and a formal change of frame specialize to invertible
changes on the associated graded coefficient space. This proves the
last assertion as well.
\end{proof}

\subsection{The simultaneous collision theorem}
Let $D_0=\Spec\C[[\tau]]$, let $\mathcal R$ be a rank-two pencil
subbundle of $\Sym^2V\otimes\OO_{D_0}$, and put
$X=\Pj(\mathcal R)$. Assume its generic pencil has $n$ distinct
spectral points. Its special pencil is assumed regular. Apart from
simple spectral points, suppose it has distinct points
$x_1,\ldots,x_k$ of coranks $c_a$, where $2\leq c_a\leq n-1$.
Write $r_a=n-c_a$.

At $x_a$, choose a parameter $z_a$ on the special pencil and a unit
$r_a\times r_a$ block. Its symmetric Schur complement is required to
satisfy
\begin{equation}\label{eq:ordinary-cluster-v159}
 S_a=z_a C_a+\tau D_a+O((z_a,\tau)^2),\qquad
 C_a\text{ nonsingular},\qquad
 \det(\alpha C_a+\beta D_a)\text{ has }c_a\text{ simple zeros}.
\end{equation}
We call this an ordinary transverse cluster. Intrinsically,
$C_a,D_a\in\Sym^2W_a$, where
$W_a=(\ker(A(x_a):V^*\to V))^*$. Their pencil is the first derivative
of the form restricted to its radical. Taking a Schur complement does
not alter that first derivative. Replacing the parameter changes the
basis of this residual pencil and preserves its marked member $[C_a]$.
Thus the hypothesis does not depend on the chosen complement or frame.

\begin{theorem}[Simultaneous ordinary collision fibres]
\label{thm:ordinary-fibres-v159}
Under \eqref{eq:ordinary-cluster-v159}, the schematic closure of the
generic lifted pencil in $\mathrm{CQ}(V)\times D_0$ is
\begin{equation}\label{eq:ordinary-blowup-v159}
               Z=\operatorname{Bl}_{\{x_1,\ldots,x_k\}}X.
\end{equation}
It is regular and flat over $D_0$. Its special fibre is the reduced
nodal tree
\[
                          C\cup F_1\cup\cdots\cup F_k,
\]
where $C$ is the lift of the special pencil and each $F_a\simeq\Pj^1$
meets $C$ in one point. There are no embedded components and the tails
are pairwise disjoint. The tail $F_a$ is the complete-quadric lift of
the pointed residual pencil
$([\alpha C_a+\beta D_a],[C_a])$ on $W_a$. It lies in $E_{r_a}$,
meets $E_{n-1}$ in $c_a$ distinct reduced points, meets no other
boundary divisor, and attaches at the nonsingular member $[C_a]$.

Let $H_q$ be the hyperplane bundle of the resolved $q$th exterior map,
$1\leq q\leq n-1$. Then
\begin{equation}\label{eq:ordinary-multidegrees-v159}
 \deg_{F_a}H_q=(q-r_a)_+,\qquad
 \deg_C H_q=q-\sum_a(q-r_a)_+.
\end{equation}
For the product polarization $\bigotimes_{q=1}^{n-1}H_q$ the Hilbert
polynomial is $\binom n2 m+1$. In particular
$\deg F_a=\binom{c_a}{2}$ and
$\deg C=\binom n2-\sum_a\binom{c_a}{2}$.

For every $h\geq1$ and $1\leq p\leq hn$, set
\begin{equation}\label{eq:ordinary-power-numbers-v159}
 b_a=(p-hr_a)_+,\qquad s_p=(p-h(n-1))_+,\qquad
 \delta_{a,p}=b_a-c_a s_p.
\end{equation}
These integers are nonnegative. The resolved $p$th power restricts
to the complete linear system of degree $\delta_{a,p}$ on $F_a$.
It is constant if that degree is zero, and otherwise the complete
Veronese map into its linear span. Near $x_a$, with
$\mathfrak n_a=(z_a,\tau)$ and $f_a=\det S_a$, the exact ideals are
\begin{equation}\label{eq:ordinary-power-ideals-v159}
 J_p(A)=\mathfrak n_a^{\delta_{a,p}}(f_a)^{s_p},\qquad
                   J_{hr_a+1}(A)=\mathfrak n_a.
\end{equation}
Equalities are up to multiplication by a unit after a local frame
is fixed, not up to radical or integral closure. Consequently the
multiplication powers select the graph modification and every tail
linear system. This description depends only on the first normal jet
in \eqref{eq:ordinary-cluster-v159}; with the special pencil fixed,
higher-order terms do not change its embedded special fibre or its
power maps.
\end{theorem}
\begin{proof}
Near $x_a$ the unit block and its Schur complement give
$I_q(A)=\OO$ for $q\leq r_a$ and
$I_{r_a+j}(A)=I_j(S_a)$. Lemma~\ref{lem:tangent-factors-v159}
therefore gives
\[
 I_{r_a+j}(A)=\mathfrak n_a^j\quad(1\leq j<c_a),\qquad
 I_n(A)=(u_a f_a),\quad u_a\in\OO^\times.
\]
At every other point the rank is at least $n-1$, so all minor ideals
of order at most $n-1$ are units. The determinant ideal is invertible.
The simultaneous exterior graph is the blow-up of the product of
its base ideals by Lemma~\ref{lem:simultaneous-graph-v158}.
Near $x_a$ this product is
$\mathfrak n_a^{\binom{c_a}{2}}$. A positive power has the same
blow-up as the ideal itself, and the points are disjoint. These local
identifications glue uniquely over $X$ to \eqref{eq:ordinary-blowup-v159}.
This computes the graph on the total surface; it is not an assertion
that graph closure commutes with arbitrary base change.

The tautological family of genuine pencil lines embeds $X$ into
$\Pj(\Sym^2V)\times D_0$. Hence the simultaneous graph embedding
places $Z$ in $\mathrm{CQ}(V)\times D_0$, with exactly the prescribed
generic curve. The two charts $\tau=z_av$ and $z_a=\tau w$ of each
point blow-up show that $Z$ is regular. In the first chart its special
fibre is $z_av=0$. Both branches have multiplicity one and intersect
transversally. In the second chart only the exceptional branch occurs.
The parameter $\tau$ is a nonzerodivisor, so $Z$ is flat over the DVR.
The charts also rule out embedded components.

On $F_a$, divide the Schur block by the exceptional parameter.
Its value is $\alpha C_a+\beta D_a$. All earlier unit-block data
specialize to their values at $x_a$, independently of the direction.
The complete-quadric chart thus identifies $F_a$ with the lifted
residual pencil in the fibre over $[A(x_a)]$. Its only residual rank
drops have corank one and occur at the $c_a$ simple determinant zeros.
This proves the boundary assertions, including the attachment at
$[\alpha:\beta]=[1:0]$. For $q=r_a+j<n$, the leading $q$-minors
are the $j$-minors of this residual pencil. They have no common zero
and have degree $j$. The lower maps are constant. On $C$, the common
factor in the $q$-minors has order $(q-r_a)_+$ at $x_a$ and no other
zeros. This proves \eqref{eq:ordinary-multidegrees-v159}.
The sum of the component degrees is $\binom n2$. The normalization
sequence for a tree of $k+1$ projective lines subtracts its $k$ nodes,
so its Euler characteristic is one. This proves the Hilbert polynomial
and identifies the curve as the fibre in the Hilbert incidence space.

For the power ideals apply Theorem~\ref{thm:universal-power-v157}.
If $p\leq hr_a$, the ideal is the unit ideal. Otherwise write
$p=hr_a+hj+s$ with $0\leq s<h$. For $j<c_a-1$, the two minor
ideals are $\mathfrak n_a^j$ and $\mathfrak n_a^{j+1}$, so the
power ideal is $\mathfrak n_a^{hj+s}$. For $j=c_a-1$ it is
\[
        \mathfrak n_a^{(c_a-1)(h-s)}(f_a)^s.
\]
At $p=hn$ it is $(f_a)^h$. These cases are exactly
\eqref{eq:ordinary-power-ideals-v159}. The value $p=hr_a+1$ gives
the point ideal for every $h$.

It remains to identify the complete coefficient space on $F_a$.
By Lemma~\ref{lem:block-products-v159}, its first normal coefficients,
for $p>hr_a$, are exactly those of the $b_a$th power of the residual
pencil, multiplied by a nonzero unit-block socle vector. In the
polynomial ring $\C[\alpha,\beta]$, Lemma
\ref{lem:tangent-factors-v159} gives $I_j=(\alpha,\beta)^j$ for
$j<c_a$. The universal power identity gives the homogeneous ideal
\[
 (\alpha,\beta)^{\delta_{a,p}}
                    \det(\alpha C_a+\beta D_a)^{s_p}.
\]
All its original power generators have degree $b_a$; therefore their
span is its entire degree-$b_a$ part. Remove the common determinant
factor. The remaining space is exactly
$H^0(\Pj^1,\OO(\delta_{a,p}))$, proving the assertion about the
linear system. This argument determines the map as well as its degree.
All surviving leading coefficients use only the first normal jet;
the higher-order terms vanish on the exceptional divisor.
\end{proof}

\subsection{A modular family over a fixed pencil}
Let $R_0$ be a regular pencil whose spectral elementary divisors all
have exponent one, with repeated spectral multiplicities
$c_1,\ldots,c_k\geq2$. Thus these are semisimple collisions, not
nontrivial Jordan blocks. At $x_a$ keep the residual space $W_a$ and
the nonsingular tangent member $[C_a]$. Put
\begin{equation}\label{eq:residual-parameter-v159}
 U_a=\{\ell\text{ a line through }[C_a]\text{ in }\Pj(\Sym^2W_a):
             \ell\text{ meets the determinant in }c_a\text{ simple points}\}.
\end{equation}
This is a nonempty open in
$\Pj(\Sym^2W_a/\langle C_a\rangle)$, of dimension
$\binom{c_a+1}{2}-2$.

\begin{theorem}[Residual-pencil moduli in an incidence fibre]
\label{thm:boundary-moduli-v159}
The product $U=\prod_aU_a$ parametrizes an algebraic flat family of
distinct embedded Hilbert limits over the fixed pencil $R_0$.
Each is the fixed main curve with the tails prescribed by its
residual lines, and each is realized by an ordinary transverse
one-parameter smoothing of $R_0$. Hence, for the normal incidence
modification $\rho:\widehat G\to G$, one has
\begin{equation}\label{eq:boundary-fibre-dimension-v159}
                \dim\rho^{-1}(R_0)\geq
                    \sum_a\left(\binom{c_a+1}{2}-2\right).
\end{equation}
Distinctness is detected by the resolved multiplication power
$J_{h(n-c_a)+1}$ on the $a$th tail. The statement describes a specified
family in the fibre, not the entirety of the fibre or its normalization.
\end{theorem}
\begin{proof}
Since $C_a$ is nonsingular, lines through $[C_a]$ are parametrized
by the displayed projective quotient. The condition on determinant
zeros is open and nonempty, as is seen by choosing $D_a$ diagonal
with distinct eigenvalues relative to $C_a$.

The complete-quadric morphism is an isomorphism over forms of rank
at least $c_a-1$. Each line in $U_a$ stays in this open, and thus
its universal line has an algebraic lift. The fibre of
$\mathrm{CQ}(V)\to\Pj(\Sym^2V)$ over $[A(x_a)]$ contains the
residual $\mathrm{CQ}(W_a)$: the nondegenerate first quotient form
is fixed, and its radical carries the remaining complete form.
This follows also by taking the unit block in the Schur charts used
in the preceding proof. Thus the universal lifted lines give an
algebraic family of embedded tails, each passing through the fixed
attachment corresponding to $[C_a]$.

Adjoin the fixed main curve. The scheme-theoretic intersection of
each tail family with that main curve is just this reduced section.
Indeed the main curve has contact order one with $E_{r_a}$, since
its local nonzero Smith exponents are $0$ and $1$, whereas the tail
lies inside that divisor. The tails lie over different points, hence
do not intersect each other. The union exact sequence
\[
 0\longrightarrow\OO_{C\cup\bigcup F_a}
 \longrightarrow\OO_{C\times U}\oplus\bigoplus_a\OO_{F_a}
 \longrightarrow\bigoplus_a\OO_{\{\mathrm{attachment}_a\}\times U}
 \longrightarrow0
\]
proves flatness over $U$: the two terms on the right are flat, so
its kernel is flat. The degrees from
\eqref{eq:ordinary-multidegrees-v159} give the required fixed Hilbert
polynomial. We therefore obtain a morphism $U$ to the Hilbert scheme.
The embedded fibre over $[A(x_a)]$ recovers its residual tail. Its
first nonconstant exterior map, equivalently the power of index
$h(n-c_a)+1$, has degree one and recovers the line $\ell$ itself.
Distinct parameters give distinct embedded curves.

To see that these Hilbert points belong to the graph closure defining
$\widehat G$, choose a nonsingular member of $R_0$. Semisimplicity
makes its spectral eigenspaces orthogonal, so in a congruence frame
the pencil is a direct sum of blocks $(z-x_a)C_a$ and simple blocks.
Choose any representative $D_a$ of each residual line and replace
these blocks by $(z-x_a)C_a+\tau D_a$. Each cluster splits into
$c_a$ simple spectral points and distinct clusters stay separated.
Theorem~\ref{thm:ordinary-fibres-v159} identifies the Hilbert limit
with the constructed curve. This realizes every parameter by a
smoothing, independently on all the blocks.

Let $\overline G$ denote the reduced Hilbert graph before normalization.
The morphism from the smooth, reduced scheme $U$ factors through its
fibre over $R_0$: the defining equations vanish at every closed point
by the smoothing construction, hence vanish on $U$. Its geometric
fibres are zero-dimensional by the distinctness just proved; its
image consequently has dimension $\dim U$. Normalization
$\widehat G\to\overline G$ is finite and surjective, and the preimage
of this image has the same dimension. This gives
\eqref{eq:boundary-fibre-dimension-v159}.
We use this finite preimage, rather than asserting that a nondominant
map from a normal scheme must lift to a normalization. Such an
assertion would not be a valid normalization universal property.
\end{proof}

\begin{example}[Triple collision]
\label{ex:triple-collision-v159}
For one semisimple triple collision, $c=3$, the tail has exterior
multidegree $(0,\ldots,0,1,2)$ and total degree three. At $h=2$ its
successive residual powers $b=0,\ldots,6$ have degrees
\[
                             0,1,2,3,4,2,0.
\]
The residual lines through a fixed nonsingular ternary quadratic
form form an open in $\Pj^4$, and give a four-dimensional family of
distinct Hilbert limits over the same special pencil. Every member
has three distinct spectral points on its tail and one nonspectral
attachment. This is a family of embedded limits, not just a count
of branches of a normalization fibre.
\end{example}

The special finite failure algebra remains $A_{R_0}$ throughout this
family of limits. What varies is the total-family first-relation
geometry: its intrinsic envelope produces the maximal-ideal centre,
its successive products give the tail systems, and their degree-one
member recovers the residual line. Thus the family supplies information
not contained in its closed algebra alone. The construction is
equivariant under source congruence and under reparametrization of the
pencil, but the dimension bound concerns the embedded incidence space,
not a quotient by its stabilizer.

\section{The simple corank-two incidence modification}
\label{sec:incidence-blowup-v160}
The ordinary-collision calculation describes specified families of
limits. We now identify the entire normalized Hilbert incidence space
on a natural open of the pencil Grassmannian. The result holds in every
dimension. It does not require a semisimple collision or a squarefree
residual pencil.

Let $n\geq3$, $M=\Sym^2V$, $P=\Pj(M)$, and $G=\Gr(2,M)$.
Write $Z_j\subset P$ for the symmetric rank-at-most-$j$ scheme, and
put $Z_0=\varnothing$ when $n=3$. Let $\mathcal U\subset G$ consist
of regular pencils $R$ such that their lines $L_R=\Pj(R)$ avoid
$Z_{n-3}$ and their scheme-theoretic intersections with $Z_{n-2}$
are either empty or a single reduced point. Write $\mathcal P$ for
the universal line over $\mathcal U$. We use projective spaces of
lines, so that the exceptional fibre of a smooth blow-up is the
projectivization of the normal space.

\begin{lemma}[The incidence centre and its normal space]
\label{lem:incidence-centre-v160}
The locus $\mathcal U$ is open and contains $G^\circ$ of
Proposition~\ref{prop:pencil-compactification-v157}. Its incidence
scheme
\[
 \mathcal D=\mathcal P\times_P Z_{n-2}\longrightarrow\mathcal U
\]
is a closed immersion with smooth image $\Sigma_2$ of codimension two.
At $R\in\Sigma_2$, let $p=L_R\cap Z_{n-2}$ and let
$a_R\subset N_{Z_{n-2}/P,p}$ be the nonzero image of $T_pL_R$.
There is a canonical isomorphism
\begin{equation}\label{eq:incidence-normal-v160}
 N_{\Sigma_2/\mathcal U,R}
       \simeq N_{Z_{n-2}/P,p}/a_R.
\end{equation}
Near $(p,R)$, after an etale change of base and coordinates, the
ideal of $\mathcal D$ in $\mathcal P$ is
\begin{equation}\label{eq:incidence-local-v160}
       I_{n-1}(A)=(x,u,v),\qquad I_{\Sigma_2}=(u,v),
\end{equation}
where $x$ is a relative line coordinate and $u,v$ are part of a
regular system of parameters on the base. The generators may be
changed by an invertible matrix; the asserted equality is an equality
of ideals.
\end{lemma}
\begin{proof}
On $P\setminus Z_{n-3}$ the scheme $Z_{n-2}$ is smooth of codimension
three. A nonsingular $(n-2)$-block reduces its equations to the three
entries of a symmetric $2\times2$ Schur complement. In particular
these are regular local equations for its given determinantal scheme
structure.

The universal line evaluation $\mathcal P_G\to P$ is smooth: over a
point of $P$ its fibre is the projective space of lines through that
point. Hence its inverse image of the smooth rank locus is smooth
and has dimension $\dim G-2$. Lines contained in the rank locus, or
meeting the lower rank locus, form closed exceptional sets after the
usual proper incidence projection. On the complementary open the
incidence projection is finite. The condition that its fibre have
length at most one is open, by upper semicontinuity for a finite
module. This proves the asserted openness. On this open the unit map
$\OO_{\mathcal U}\to q_*\OO_{\mathcal D}$ is surjective: on each
geometric fibre its target is either zero or a one-dimensional
algebra with unit, and Nakayama's lemma applies to its coherent
cokernel. Thus $q$ is a closed immersion. Its source is smooth by the
smooth evaluation just noted; its codimension is two.

First-order motion of a line evaluates surjectively onto
$T_pP/T_pL_R$: in the description
$T_RG=\Hom(R,M/R)$ this is restriction to the one-dimensional
subspace representing $p$. To remain incident to the rank locus,
this evaluation must vanish in $T_pP/(T_pZ_{n-2}+T_pL_R)$.
A reduced intersection means that the normal image of $T_pL_R$ is
nonzero. The kernel is precisely $T_R\Sigma_2$, since the incidence
point may move along both the line and the rank locus. The quotient
is therefore \eqref{eq:incidence-normal-v160}.

For the local assertion choose a rank-locus equation with nonzero
relative derivative. Its zero hypersurface is etale over the base
near $(p,R)$; after an etale base change it gives a section. Use that
equation as the coordinate $x$. Subtract multiples of $x$ from the
other two equations. Their restrictions to the section are base
functions $u,v$, and the ideal becomes $(x,u,v)$. The normal-space
calculation proves that $du,dv$ are independent. This proves
\eqref{eq:incidence-local-v160}. The argument is the smooth-centre
coordinate calculation, not an assertion that arbitrary graph
closures commute with base change.
\end{proof}

The open is nonempty along its centre. For example, a diagonal pencil
with two equal linear entries and all other entries having distinct
zeros away from their common zero belongs to $\Sigma_2$.
Although $a_R\ne0$, its residual binary quadratic form may have rank
one. This case is allowed throughout the following theorem.

\begin{theorem}[The full simple-incidence Hilbert modification]
\label{thm:incidence-blowup-v160}
For the normalized Hilbert incidence modification
$\rho:\widehat G\to G$ of
Proposition~\ref{prop:pencil-compactification-v157}, there is a
canonical isomorphism over $\mathcal U$
\begin{equation}\label{eq:incidence-blowup-v160}
 \widehat G\times_G\mathcal U
       \simeq B:=\operatorname{Bl}_{\Sigma_2}\mathcal U.
\end{equation}
It is equivariant under source congruence. In particular this part
of the incidence space is smooth. Its exceptional divisor
$F=\Pj(N_{\Sigma_2/\mathcal U})$ is a $\Pj^1$-bundle, and
\begin{equation}\label{eq:incidence-canonical-v160}
 K_B=\beta^*K_{\mathcal U}+F,
 \qquad \OO_B(F)|_{F_R}=\OO_{\Pj^1}(-1),
\end{equation}
where $\beta:B\to\mathcal U$ is the blow-up.

For every $R\in\Sigma_2$ the entire fibre of $\rho$ is $F_R\simeq
\Pj^1$. Its point $[\bar b]\in\Pj(N_{Z_{n-2}/P,p}/a_R)$ represents
the reduced nodal curve
\begin{equation}\label{eq:incidence-fibre-v160}
 C_R\cup\ell_{[\bar b]},\qquad
 \ell_{[\bar b]}=\Pj(a_R+\C b)
                  \subset\pi^{-1}(p)=\Pj(N_{Z_{n-2}/P,p}).
\end{equation}
Here $C_R$ is the strict transform of $L_R$ and $b$ is any lift of
$\bar b$. The tail attaches to $C_R$ at $[a_R]$. Every line through
$[a_R]$ in this exceptional $\Pj^2$ occurs, and no further components,
embedded points, or nonreduced fibres occur over this open. The two
components have exterior multidegrees
\begin{equation}\label{eq:incidence-multidegrees-v160}
 \deg_{C_R}H_q=q-\mathbf1_{q=n-1},\qquad
 \deg_{\ell}H_q=\mathbf1_{q=n-1}\quad(1\leq q<n).
\end{equation}
Their common Hilbert polynomial is $\binom n2m+1$ for
$\bigotimes_{q=1}^{n-1}H_q$.

The universal curve over $B$ is the simultaneous multiplication graph
on $\mathcal P\times_{\mathcal U}B$. Near a node it is etale locally
$xy=t$, where $t=0$ defines $F$. Thus it is a flat family of nodal
curves, including along directions for which the residual determinant
has a repeated zero. It commutes with base change as this specified
flat family. The blow-up is characterized by the universal property
that $I_{\Sigma_2}\OO_B$ is an invertible ideal; the corresponding
factorization property holds for every $T\to\mathcal U$ on which
that image ideal is an effective Cartier ideal.
\end{theorem}
\begin{proof}
On $P\setminus Z_{n-3}$, the complete-quadric morphism is exactly the
blow-up of the smooth centre $Z_{n-2}$: all earlier rank blow-ups
are isomorphisms there and the determinant divisor requires no
blow-up. We use the scheme-level classical construction
\cite[Remark~2.5 and Construction~2.6]{Massarenti2020}, originating
in \cite[Theorem~6.3]{Vainsencher1982}. Nothing in this step declares
the space of complete quadrics itself new.

Pull the universal line back to $B$. In a base blow-up chart
$u=t$, $v=tw$, the ideal \eqref{eq:incidence-local-v160} becomes
$(x,t)$. Its blow-up has the two charts
\[
        \OO_B[x,y]/(xy-t),\qquad \OO_B[z],\quad x=tz.
\]
They glue by the usual ratio change. Their total spaces are smooth,
and they are flat over $B$. For example
$\OO_B[x,y]/(xy-t)$ is free as an $\OO_B$-module with basis
$1,x,x^2,\ldots,y,y^2,\ldots$; reduction by the monic relation $xy-t$
gives this basis. Its fibres at $t=0$ have two reduced branches and
one ordinary node. The other base chart has the same description.
Outside the incidence section the pulled ideal is the unit ideal.
Blowing up the globally defined pulled ideal therefore gives a
projective flat family $\mathcal C_B$ with these local charts.
All smaller minor ideals are units on this open. The universal power
identity makes the product of the nontrivial power ideals an
$h^2$th power of $I_{n-1}$, up to an invertible determinant ideal.
Thus this same blow-up is the simultaneous multiplication graph
for every $h$, not only the exterior graph.

There is a closed immersion
$\mathcal C_B\hookrightarrow\mathrm{CQ}(V)\times B$. Indeed the
base-changed Rees algebra of the rank-locus ideal surjects onto the
Rees algebra of its image on the universal line; taking Proj gives
a closed immersion into the base-changed complete-quadric graph.
The universal line is already closed in $P\times B$. Composing these
two closed immersions proves the assertion. Equivalently, use the
simultaneous graph lemma on the integral total space. This argument
does not identify a general base change of a blow-up with the blow-up
of the image ideal without the indicated Rees quotient.

In normal coordinates at $p$, the exceptional map is
$[x:t]\mapsto[x a_R+t b]$, with $b$ representing the normal direction
of the base motion. Formula \eqref{eq:incidence-normal-v160} implies
that $a_R,b$ are independent. This is a line embedded in
$\Pj(N_{Z_{n-2}/P,p})$, not a multiple cover. The strict-transform
branch meets it at $[a_R]$. This proves
\eqref{eq:incidence-fibre-v160}, including the converse parametrization
of all lines through that point. The local charts show the absence
of embedded components and multiplicities.

The minor systems of orders $q<n-1$ have no base point on $L_R$.
The $(n-1)$-minors have one simple common zero, by
\eqref{eq:incidence-local-v160}. Their degrees on $C_R$ are therefore
$q$ and $n-2$, respectively. On the exceptional $\Pj^2$ only the
last exterior system has positive degree, namely one. This proves
\eqref{eq:incidence-multidegrees-v160}. The two projective lines
have Euler characteristic $2-1=1$, since their intersection is one
reduced point. Their total degree is $\binom n2$.

The flat embedded family defines a morphism from $B$ to the reduced
Hilbert graph closure over $\mathcal U$. It is proper because $B$ is
proper over $\mathcal U$ and the Hilbert graph is separated over that
base. Its image is exactly the graph closure: it is closed and contains
the dense graph over $\mathcal U\setminus\Sigma_2=G^\circ$, while
$B$ is integral. Over an incident pencil the displayed exceptional
lines are pairwise distinct as embedded curves; over a nonincident
pencil the fibre is a point. The morphism thus has finite geometric
fibres and is finite by \cite[Tag~02LS]{StacksProject160}. It is
birational, and $B$ is normal because its centre is smooth. A finite
birational map from a normal integral scheme identifies it with the
normalization: on affine charts its ring is the integral closure in
the common function field. Normalization commutes with restriction
to the open $\mathcal U$. This proves
\eqref{eq:incidence-blowup-v160}, without assuming that an arbitrary
nondominant normal test scheme lifts to a normalization.

All maps are specified by the universal incidence and agree on the
dense nonincident locus. Their uniqueness proves canonicity and
congruence equivariance. For completeness, in coordinates $v=uw$,
$du\wedge dv=u\,du\wedge dw$, giving the discrepancy one in
\eqref{eq:incidence-canonical-v160}. The exceptional normal bundle
is tautological of degree $-1$. The universal property is the usual
Rees-algebra factorization for an effective Cartier image ideal
\cite[Tag~0806]{StacksProject160}; here the preceding argument
identifies the actual Hilbert incidence modification with that
universal object. Flat pullback preserves the explicitly constructed
nodal family. No assertion about unrestricted graph-closure base
change is used.
\end{proof}

\subsection{Every power on every exceptional direction}
\begin{proposition}[Complete exceptional power systems]
\label{prop:exceptional-powers-v160}
Fix $R\in\Sigma_2$ and $p=L_R\cap Z_{n-2}$. For $h\geq1$ and
$1\leq j\leq hn$, define
\begin{equation}\label{eq:exceptional-delta-v160}
 \delta_j=(j-h(n-2))_+-2(j-h(n-1))_+.
\end{equation}
On the full exceptional plane $\pi^{-1}(p)\simeq\Pj^2$, the resolved
$j$th multiplication map is the complete degree-$\delta_j$ Veronese
map into its linear span, with degree zero interpreted as constant.
Its restriction to every tail in \eqref{eq:incidence-fibre-v160} is
therefore the complete degree-$\delta_j$ system on $\Pj^1$.
In particular $j=h(n-2)+1$ detects the tail as an embedded line.
These assertions include tangent lines to the residual determinant
conic and attachments on that conic.
\end{proposition}
\begin{proof}
A unit $(n-2)$-block leaves a residual two-dimensional space $W$.
Lemma~\ref{lem:block-products-v159} identifies the first nonzero
normal coefficient at degree $j>h(n-2)$ with the power of order
$b=j-h(n-2)$ in $B_h(W)$, multiplied by a nonzero unit-block socle.
For $j\leq h(n-2)$ the order-zero coefficient is constant on the
exceptional fibre. In the three-variable coordinate ring of
$\Sym^2W$, put $\mathfrak a=I_1$ and $f=\det$. The universal ideal
formula gives
\[
 J_b=\mathfrak a^b\quad(b\leq h),\qquad
 J_b=\mathfrak a^{2h-b}f^{b-h}\quad(h\leq b\leq2h).
\]
The generators all have degree $b$. Their span is the entire degree
$b$ part of the displayed ideal. After removal of the common factor
$f^{(b-h)_+}$ it is the full space of ternary forms of degree
$b-2(b-h)_+=\delta_j$. Hence it defines the complete Veronese system
on the exceptional plane. Restriction of all ternary forms to any
line gives all binary forms of that degree, independently of its
intersection multiplicities with the conic $f=0$. This proves all
assertions, including the endpoint $j=hn$ where the map is constant.
\end{proof}

\section{The incidence centre from failure multiplication}
\label{sec:failure-incidence-v160}
We specify how the preceding modification is recovered from a family
of sharp failure algebras. This is a construction of its centre and
Rees algebra, not merely a pullback of the existing finite algebra.
For a framed source let $\Phi:G\to\mathcal H$ be the first-relation
closed immersion of Theorem~\ref{thm:finite-parameter-immersion}.
Here $\mathcal H$ denotes its relation-space Grassmannian. We use
$\Phi(\mathcal U)$ with its induced scheme structure, not an arbitrary
neighbourhood in $\mathcal H$.

\begin{theorem}[Intrinsic incidence ideal and normal direction]
\label{thm:failure-incidence-v160}
On the effective family of sharp pencil-failure algebras corresponding
to $\mathcal U$, its first relation, oriented source, and multiplication
diagram determine $I_{\Sigma_2}$ and the graded Rees algebra
\begin{equation}\label{eq:failure-rees-v160}
                \bigoplus_{m\geq0}I_{\Sigma_2}^{\,m}.
\end{equation}
Their Proj is the modification in
Theorem~\ref{thm:incidence-blowup-v160}. More explicitly, restrict
the power ideal $J_{h(n-2)+1}$ to the intrinsic pencil line and let
$q:\mathcal D\to\mathcal U$ be its zero scheme. Then
\begin{equation}\label{eq:failure-incidence-ideal-v160}
 I_{\Sigma_2}=\ker\{\OO_{\mathcal U}\longrightarrow q_*\OO_{\mathcal D}\}.
\end{equation}
The construction is independent of $h$. It is equivariant for ungraded
failure-algebra isomorphisms and descends along an oriented projective
source torsor.

An arc through $R\in\Sigma_2$ whose first normal derivative is nonzero
has exceptional limit given by its line in
$N_{\Sigma_2/\mathcal U,R}$. Thus its first-order failure-algebra
family determines the embedded Hilbert limit. A fixed closed algebra
$A_R$ alone determines the pencil and its full exceptional parameter
space, but not a preferred point of that space. For higher-contact
arcs the leading nonzero normal ideal quotient replaces the first
normal derivative.
\end{theorem}
\begin{proof}
The first-relation inverse recovers the oriented projective source
and the intrinsic pencil, with its arbitrary-base closed immersion
$\Phi$. The source-normalized envelope in
Theorem~\ref{thm:intrinsic-envelope-v158} supplies its power diagram.
The universal power identity gives
\[
 J_{h(n-2)+1}=I_{n-2}^{h-1}I_{n-1}=I_{n-1}
\]
on $P\setminus Z_{n-3}$, since $I_{n-2}$ is the unit ideal there.
For $h=1$ the same formula is $J_{n-1}=I_{n-1}$. Lemma
\ref{lem:incidence-centre-v160} identifies its zero scheme on the
universal line with $\Sigma_2$ by a closed immersion. Taking the
kernel of the unit map gives its exact scheme ideal, establishing
\eqref{eq:failure-incidence-ideal-v160}. Multiplication of this
ideal produces \eqref{eq:failure-rees-v160}, whose Proj was identified
with the normalized Hilbert graph in the preceding theorem. Because
$\Phi$ is a closed immersion onto its scheme image, this also gives
\[
 \operatorname{Bl}_{\Phi(\Sigma_2)}\Phi(\mathcal U)
                 \simeq\operatorname{Bl}_{\Sigma_2}\mathcal U.
\]
This is an identity on the effective image, not an unjustified
assertion about blow-ups of the whole ambient Grassmannian.

In the local coordinates of \eqref{eq:incidence-local-v160}, an arc
with $u(\tau)=\tau b_1+O(\tau^2)$ and
$v(\tau)=\tau b_2+O(\tau^2)$, $(b_1,b_2)\ne(0,0)$, lifts to the
exceptional point $[b_1:b_2]$. Under
\eqref{eq:incidence-normal-v160} this gives exactly the plane
$a_R+\C b$ and hence the embedded tail. The arbitrary-base inverse
identifies first-order pencil deformations with their effective
first-relation deformations. Terms of degree at least two in an
ungraded generator change do not change the sharp relation, so they
do not change this normal direction. Higher-contact arcs are handled
by dividing $u,v$ by their common smallest power of $\tau$ before
specializing. This calculation explains why no preferred tail can
be assigned to the closed algebra without its normal direction.

Every operation is invariant under the congruence and source-lift
changes already treated in the intrinsic-envelope theorem. Its scalar
normalization and projectivization remove the linear-lift ambiguity;
fppf descent therefore applies to the incidence ideal, its Rees
algebra, and the resulting projective family. The universal Cartier
factorization property, rather than an assertion of arbitrary
blow-up base change, specifies their functorial meaning on test bases.
\end{proof}

\begin{example}[A complete fibre with a tangential residual limit]
\label{ex:tangent-limit-v160}
For $n=3$, take the affine two-parameter family of genuine pencils
\[
 A_{u,v}(s,t)=
 \begin{pmatrix}s&0&0\\0&t+us&vs\\0&vs&t-us\end{pmatrix}.
\]
Near the pencil $R_0=\langle\diag(1,0,0),\diag(0,1,1)\rangle$,
the incidence centre is $(u,v)=(0,0)$. Indeed in the chart $s=1$
its minor ideal is $(t,u,v)$, while near $s=0$ the rank is two.
The exceptional fibre is the full projective line $[u:v]$.
The arc $(u,v)=\tau(\alpha,\beta)$ gives the residual tail
\[
 [x:y]\longmapsto
 \begin{pmatrix}x+\alpha y&\beta y\\
                 \beta y&x-\alpha y\end{pmatrix},
 \qquad \det=x^2-(\alpha^2+\beta^2)y^2.
\]
All directions $(\alpha,\beta)\ne(0,0)$ occur. For
$\alpha^2+\beta^2\ne0$ the two residual determinant zeros are distinct.
For $[\alpha:\beta]=[1:i]$ or $[1:-i]$ they coalesce into a double
zero. The tail remains an embedded reduced line and the total curve
remains nodal. These two points of the full incidence fibre are not
ordinary transverse residual pencils, and so are not covered by the
squarefree hypothesis of Theorem~\ref{thm:ordinary-fibres-v159}.
Every power system at these points is still the complete system of
Proposition~\ref{prop:exceptional-powers-v160}. The example separates
nodality of the Hilbert fibre from transversality of its boundary
contact.
\end{example}

\section{Multiple reduced incidence and its complete Hilbert boundary}
\label{sec:multiple-incidence-v161}
We remove the restriction to a single incidence point. The independence
of the incidence conditions is proved from the radicals of distinct
members of a regular pencil; it is not added as a transversality
hypothesis. The resulting modification has a global Fitting-ideal
description even though the incidence points need not be globally
labelled.

Put $M=\Sym^2V$, $P=\Pj(M)$ and $G=\Gr(2,M)$, with $n=\dim V\geq3$.
Let $Z_j\subset P$ have its symmetric determinantal scheme structure.
When $n=3$ put $Z_0=\varnothing$. Let $G_{\rm reg}$ be the open of
regular pencils and let $X_G=\Pj(\mathcal R)$ be the universal line.
The image of $X_G\times_P Z_{n-3}$ is closed, since the projection
from this incidence is proper. Remove it and call the resulting open
$G_*\subset G_{\rm reg}$. Every line over $G_*$ meets $Z_{n-2}$ in
a finite scheme: a positive-dimensional intersection would contain
the line and contradict regularity. Consequently
\[
 q_*^{\rm inc}:D_*:=X_{G_*}\times_P Z_{n-2}\longrightarrow G_*
\]
is proper and quasi-finite, hence finite. Define
\begin{equation}\label{eq:reduced-open-v161}
 U_{\rm red}=G_*\setminus
 q_*^{\rm inc}\bigl(\Supp\Omega_{D_*/G_*}\bigr),\qquad
 D=D_*\times_{G_*}U_{\rm red},\qquad q:D\to U_{\rm red}.
\end{equation}
This is open because $q_*^{\rm inc}$ is finite. Its geometric
incidences are precisely the finite reduced intersections: for a
finite algebra over $\C$, vanishing of relative differentials is
equivalent to being a product of copies of $\C$. Unlike the
single-incidence open, it allows every possible number of such points.
Set
\begin{equation}\label{eq:fitting-centre-v161}
 \mathcal Q=q_*\OO_D,\qquad
 \mathcal I=\Fitt_0(\mathcal Q),\qquad
 B_{\rm red}=\operatorname{Bl}_{\mathcal I}U_{\rm red}.
\end{equation}
The zeroth Fitting ideal of the zero module is the unit ideal.

\subsection{The general incidence calculation}
\begin{lemma}[Transverse incidence of a family of curves]
\label{lem:general-incidence-v161}
Let $S,Y$ be smooth complex varieties, let $Z\subset Y$ be smooth of
codimension $c\geq2$, and let $X\to S$ be a smooth family of embedded
curves in $Y$. Suppose $D=X\times_Y Z\to S$ is finite and unramified.
Fix $s\in S$ with incidence points $p_1,\ldots,p_k$. Write
$a_i\subset N_{Z/Y,p_i}$ for the normal image of $T_{p_i}X_s$.
Assume $a_i\ne0$ and that the joint normal evaluation map
\[
 T_sS\longrightarrow\bigoplus_{i=1}^k N_{Z/Y,p_i}/a_i
\]
is surjective. Etale locally at $s$, the incidence splits into closed
immersions $D_i\simeq\Sigma_i\hookrightarrow S$, with independent
smooth centres of codimension $c-1$. Near $p_i$ there are coordinates
in which
\begin{equation}\label{eq:general-normal-v161}
 I_D=(x_i,u_{i1},\ldots,u_{i,c-1}),\qquad
 I_{\Sigma_i}=(u_{i1},\ldots,u_{i,c-1}),
\end{equation}
where all base parameters $u_{ij}$ are independent. If
$\mathcal Q_S=q_*\OO_D$, then
\[
 \Fitt_0(\mathcal Q_S)=\prod_i I_{\Sigma_i}
                       =\bigcap_i I_{\Sigma_i}.
\]
The blow-up of this ideal is the fibre product of the individual
blow-ups. It is smooth, with etale normal-crossing exceptional boundary
and fibre $\prod_i\Pj^{c-2}$ at $s$.
\end{lemma}
\begin{proof}
Over the strict henselization at $s$, the finite algebra splits into
factors indexed by its geometric fibre. These factors and their
idempotents descend to an etale neighbourhood. For each factor the
unit map from the base ring has cokernel whose fibre at $s$ is zero:
finite base change identifies that fibre with the corresponding
one-dimensional residue algebra. Nakayama's lemma, followed by
shrinking, makes the unit map surjective. Each factor is therefore a
closed immersion. This argument uses the finite algebra, not just the
set of its geometric points.

The map of the curve tangent into the normal space is nonzero, so one
local equation of $Z$ has nonzero relative derivative. Its zero set
is etale over the base. After an etale base change choose this as a
section and use that equation as the relative coordinate $x_i$.
Modulo $x_i$, the other equations become base functions $u_{ij}$.
Their differentials give exactly the joint normal evaluation in the
statement. Surjectivity makes them independent regular parameters,
proving the normal form and smoothness of all intersections of the
centres. In particular the length-one factors in this argument are
reduced on complex geometric fibres; reducedness has not been inferred
from an arbitrary higher fibre length.

Now $\mathcal Q_S=\bigoplus_i\OO_S/I_{\Sigma_i}$. Fitting ideals of
a direct sum multiply. In the displayed independent coordinates,
intersection equals product: in a polynomial coordinate ring a
monomial in every ideal uses a variable from each disjoint group, and
the equality persists under etale flat extension. Equivalently one
may verify it in the faithfully flat completion. The simultaneous
blow-up is covered by charts in which, for each $i$, one $u_{ij}$ is
retained and the others are its multiples by new ratios. These charts
are smooth and identify the fibre product with the blow-up of the
product ideal. The latter equality also follows from the simultaneous
Rees graph lemma. There are $k$ independent exceptional parameters,
and their zero sets have normal crossings. The fibre consists of
independent choices of a line in each normal space of dimension
$c-1$, proving the last assertion.
\end{proof}

The arrangement blow-up mechanism is classical. Li's wonderful-model
Theorems~1.2 and 1.3 \cite{Li2009v161} give the normal-crossing and
product-ideal description for a building set. Here we have proved the
independent-centre case explicitly. The pencil-specific assertion
needed to apply it at every point of $U_{\rm red}$ is the following.

\begin{lemma}[Independence at distinct spectral points]
\label{lem:independent-radicals-v161}
For every $R\in U_{\rm red}$, with incidence points $p_1,\ldots,p_k$,
the joint normal evaluation of Lemma~\ref{lem:general-incidence-v161}
is surjective. In particular the local incidence centres have
codimension two and meet transversely, without any additional genericity
assumption on the other elementary divisors of $R$.
\end{lemma}
\begin{proof}
Choose a nonsingular member $Q_0$ of the pencil and write its other
members as $Q_1-\lambda Q_0$. Put $T=Q_0^{-1}Q_1$ on $V^*$.
The radicals $K_i=\ker(T-\lambda_i)$ at distinct spectral points
lie in distinct generalized eigenspaces. Their sum is direct.
Restriction therefore gives a surjection
\[
 \Sym^2V\longrightarrow\bigoplus_i\Sym^2(K_i^*):
\]
extend a basis of the direct sum of the $K_i$ to a basis of $V^*$ and
prescribe the indicated diagonal blocks of a symmetric form.
Varying $Q_1$ with $Q_0$ fixed realizes every one of these restrictions.
At a corank-two point the normal space to the rank locus is the
three-dimensional space of forms on $K_i$, with its harmless local
projective twist. The tangent of the pencil maps to the nonzero line
$a_i$, because the incidence is reduced. Quotienting by $a_i$ gives
the joint evaluation in the statement. A change of basis of the
pencil contributes only multiples of its tangent image in each
summand, so the map factors through $T_RG=\Hom(R,M/R)$. It remains
surjective after this quotient. This proves the claim also when the
tangent form on $K_i$ has rank one rather than rank two.
\end{proof}

\subsection{Identification of the full Hilbert graph}
Let $G^\circ\subset U_{\rm red}$ be the open with no corank-two
incidence. A pencil over it has a unique lifted embedded curve in
$Q=\mathrm{CQ}(V)$. Use the polarization
$H=\bigotimes_{q=1}^{n-1}H_q$, where $H_q$ is the resolved exterior
hyperplane bundle, and put $P_H(l)=\binom n2 l+1$.
Let $\Gamma_{\rm red}$ be the reduced closure of this graph in
$U_{\rm red}\times\Hilb^{P_H}(Q)$, and let
$\widehat\Gamma_{\rm red}$ be its normalization. Thus
$\widehat\Gamma_{\rm red}=\widehat G\times_G U_{\rm red}$ with the
conventions of Proposition~\ref{prop:pencil-compactification-v157}.
This equality uses localization of integral closure
\cite[Section~29.55, Lemmas~29.55.3--4]{Stacks161}; it is an open
restriction, not an arbitrary base-change assertion.

\begin{theorem}[The complete multiple-incidence modification]
\label{thm:multiple-incidence-v161}
There is a canonical congruence-equivariant isomorphism
\begin{equation}\label{eq:global-multiple-v161}
 \widehat\Gamma_{\rm red}\simeq B_{\rm red}
                =\operatorname{Bl}_{\Fitt_0(q_*\OO_D)}U_{\rm red}.
\end{equation}
The space is smooth, and its exceptional boundary has normal crossings
etale locally. Over a pencil with $k$ distinct reduced corank-two
incidences the whole scheme fibre is $(\Pj^1)^k$. A point of that fibre
specifies independently, at each $p_i$, a line through $[a_i]$ in
\[
 \pi^{-1}(p_i)=\Pj(N_{Z_{n-2}/P,p_i})\simeq\Pj^2.
\]
Its Hilbert curve is the strict transform $C_R$ of the pencil line,
with these $k$ embedded tails attached at their respective points.
It is reduced and nodal, with no embedded components. For each tail
$F_i$ and each $1\leq q<n$,
\begin{equation}\label{eq:multiple-degrees-v161}
 \deg_{F_i}H_q=\mathbf1_{q=n-1},\qquad
 \deg_{C_R}H_q=q-k\mathbf1_{q=n-1}.
\end{equation}
The incidence scheme satisfies $k\leq\lfloor n/2\rfloor$.

On a labelled etale neighbourhood the exceptional divisors $E_i$
meet like coordinate hyperplanes; every subset occurs locally, and
\[
 K_{B_{\rm red}}=\beta^*K_{U_{\rm red}}+\sum_i E_i.
\]
No global labelling, or global simplicity of each irreducible boundary
component, is asserted. The universal curve is a flat nodal family;
near its $i$th node its equation is $x_i y_i=t_i$, with $t_i=0$
the corresponding exceptional divisor. The base blow-up has its
usual effective-Cartier image-ideal universal property. These
statements classify all fibres over $U_{\rm red}$, not only the fibres
of selected smoothing arcs.
\end{theorem}
\begin{proof}
Apply the preceding two lemmas with $c=3$. Locally write
$I_i=(u_i,v_i)$. In the simultaneous blow-up chart put
$u_i=t_i$, $v_i=t_iw_i$. Near the $i$th incidence the image of the
ambient rank-locus ideal on the pulled-back universal line is
$(x_i,t_i)$. Its blow-up has charts
\[
 \OO_B[x_i,y_i]/(x_i y_i-t_i),\qquad
 \OO_B[z_i],\quad x_i=t_i z_i.
\]
The incidence points have disjoint neighbourhoods on the curve, so
these modifications do not interfere. Outside them the ideal is the
unit ideal. The local hypersurfaces are flat over $B$: they are
relative Cartier divisors in a smooth relative plane and their
nonzero defining equations remain nonzerodivisors on every fibre.
The fibrewise Cartier flatness criterion applies. These are the
standard nodal local models, so all the stated fibre properties
follow. Their total spaces are smooth in the displayed independent
coordinates.

Here is the scheme-level embedding needed for the Hilbert argument.
If $\mathcal K$ is the rank-locus ideal on the open ambient space and
$f:X_B\to P$ is evaluation, there is a surjection of graded algebras
\begin{equation}\label{eq:rees-surjection-v161}
 f^*\!\left(\bigoplus_{d\geq0}\mathcal K^d\right)
       \twoheadrightarrow
             \bigoplus_{d\geq0}(\mathcal K\OO_{X_B})^d.
\end{equation}
Its kernel removes precisely the relations annihilated on passing to
the image ideal. The target is the image Rees algebra inside
$\OO_{X_B}[T]$; it is not an assumed tensor-product identification.
Taking Proj gives a closed immersion into the base-changed ambient
blow-up, which is $Q$ over rank at least $n-2$. The universal line
is already embedded in $P\times B$. We obtain a flat closed family
of curves in $Q\times B$. This agrees with the lifted pencil on the
dense nonincident open.

In normal coordinates its exceptional map at $p_i$ is
$[x_i:t_i]\mapsto[x_i a_i+t_i b_i]$. The vectors $a_i,b_i$ are
independent, and their span determines an embedded line, whose
homogeneous ideal is its linear annihilator. Distinct directions
modulo $a_i$ therefore give distinct closed subschemes of the
exceptional plane, not merely different parametrizations of one
curve. Tails at different $p_i$ cannot be confused. The induced
proper map $B\to\Gamma_{\rm red}$ has finite geometric fibres and
is birational. It is consequently finite, and the smooth source is
the normalization of the reduced integral graph. This proves
\eqref{eq:global-multiple-v161}, including its scheme fibres. Both
maps restrict to the identity on the dense graph, giving canonicity
and equivariance and gluing the etale calculations.

The smaller minor systems have no base point, while the last minor
system has one simple base point at each $p_i$. Subtracting these
base points and adding one unit of degree on each tail gives
\eqref{eq:multiple-degrees-v161}. The determinant has vanishing order
at least two at each distinct corank-two point, so $2k\leq n$.
The nodal tree has Euler characteristic one and total $H$-degree
$\binom n2$. This proves the Hilbert polynomial used above. Finally
the product blow-up charts give independent parameters $t_i$; the
Jacobian for each codimension-two blow-up contributes one $E_i$ to
the canonical divisor. These computations also prove the local
intersection statement.
\end{proof}

\begin{corollary}[Multiplication systems at every node]
\label{cor:multi-powers-v161}
The ideal $\mathcal I$ in \eqref{eq:fitting-centre-v161} is also obtained
by taking the zeroth Fitting ideal of the pushforward of the zero
scheme of $J_{h(n-2)+1}$ on the universal line, for every $h\geq1$.
At $h=1$ this is the distinguished construction without an exponent
choice. On every tail the resolved $j$th power is the complete system
of degree
\[
 \delta_j=(j-h(n-2))_+-2(j-h(n-1))_+.
\]
The exceptional parameter space at a $k$-incidence pencil has Chow
ring $\mathbb Z[\xi_1,\ldots,\xi_k]/(\xi_1^2,\ldots,\xi_k^2)$.
In particular at a double incidence the full fibre is a smooth
quadric surface $\Pj^1\times\Pj^1$, not merely two independent
one-parameter examples.
\end{corollary}
\begin{proof}
On this rank open the universal formula gives
$J_{h(n-2)+1}=I_{n-1}$ because $I_{n-2}=\OO$. Its zero scheme is
exactly $D$. Finite pushforward and Fitting ideals give the specified
centre. The power assertion is Proposition~\ref{prop:exceptional-powers-v160}
at each disjoint tail. For clarity, a coefficient ideal generated in
degree $b$ has degree-$b$ part equal to the linear span of its actual
generators. In the residual binary ring this part is
$f^{(b-h)_+}\Sym^{b-2(b-h)_+}(\Sym^2W^*)$; division by the common
factor therefore leaves all forms, not merely a subspace with the
same base locus. The Chow ring is the projective bundle formula
applied to the identified scheme fibre $(\Pj^1)^k$.
\end{proof}

\section{Higher contact, thick tails, and singular multiplication graphs}
\label{sec:higher-contact-v161}
Nonreduced incidence has a different Hilbert boundary from reduced
multiple incidence. We give a two-parameter family of genuine symmetric
pencils for every contact order $m\geq2$, determine its entire normalized
Hilbert graph, and compute the nonreduced fibres and the singularities
of the universal curve. The result is a slice theorem, not an assertion
that the whole ambient Hilbert modification commutes with restriction
to that slice.

\subsection{Symmetric linear pencils of every contact order}
Let $E_m$ be the reversal matrix, $N_m$ the nilpotent matrix with ones
on its superdiagonal, and put
\[
 L_m(x)=E_m(xI_m-N_m),\qquad e=e_1.
\]
This is symmetric. In dimension $n=2m$, consider the affine family
of pencils, with $x$ the pencil parameter and $(u,v)\in S=\A^2$,
\begin{equation}\label{eq:jordan-slice-v161}
 A_{u,v}(x)=
 \begin{pmatrix}
 L_m(x)+u ee^{\mathsf t}&v ee^{\mathsf t}\\
 v ee^{\mathsf t}&L_m(x)-u ee^{\mathsf t}
 \end{pmatrix}.
\end{equation}
The projective pencil is obtained by homogenizing every entry to degree
one. Its leading matrix is $E_m\oplus E_m$, so infinity is nonsingular.
For $n>2m$ one may adjoin fixed scalar blocks $x-\lambda_j$ with
distinct nonzero $\lambda_j$ and restrict the parameter base to a
neighbourhood where these additional simple roots stay disjoint from
the cluster. All assertions below then hold near the same cluster;
we state the global slice formula first for $n=2m$.

\begin{lemma}[Exact contact slice]
\label{lem:jordan-slice-v161}
For the family \eqref{eq:jordan-slice-v161}, the complement of the two
$e_1$ directions is a nonsingular block over $\C[x]$, of constant
nonzero determinant. Its Schur complement is exactly
\begin{equation}\label{eq:schur-contact-v161}
 \begin{pmatrix}x^m+u&v\\v&x^m-u\end{pmatrix}.
\end{equation}
Consequently, on the affine pencil chart,
\[
 I_{n-1}(A)=(x^m,u,v),\qquad
 \det A=c_m\bigl(x^{2m}-u^2-v^2\bigr),\quad c_m\in\C^*.
\]
All pencils are regular and have rank at least $n-2$ everywhere. Only
the origin of $S$ has a corank-two incidence; its intersection scheme
has length $m$. Its two positive local Smith exponents are $m,m$.
The finite incidence algebra pushed to $S$ is
$(\OO_S/(u,v))^{\oplus m}$, so its zeroth Fitting ideal is $(u,v)^m$.
\end{lemma}
\begin{proof}
Delete row and column one from $L_m(x)$. The remaining matrix is
anti-triangular with $-1$ on its antidiagonal, hence has constant
nonzero determinant. Since
$\det L_m(x)=\det(E_m)x^m$, its scalar Schur complement is $x^m$;
the ratio of the two constant determinants is one. This can also be
read directly by eliminating along the superdiagonal. The two kernel
directions receive exactly the displayed perturbation, proving
\eqref{eq:schur-contact-v161}. The unit-block minor identity then gives
$I_{n-1}=I_1(\text{Schur complement})=(x^m,u,v)$ in characteristic
zero. All smaller required minors are units. The determinant and rank
claims follow from the same block decomposition, including infinity.
At $u=v=0$ the residual matrix is $x^m I_2$, which proves both the
intersection length and the Smith exponents. The incidence algebra is
$\C[x,u,v]/(x^m,u,v)$; its $\C[u,v]$-module basis is
$1,x,\ldots,x^{m-1}$ with $(u,v)$ annihilating every basis vector.
The direct-sum formula for Fitting ideals proves the last assertion.
\end{proof}

Let $\Gamma_m$ be the reduced closure of the lifted-pencil graph over
$S\setminus\{0\}$ in $S\times\Hilb^{P_H}(\mathrm{CQ}(V))$, with
$P_H(l)=\binom n2l+1$, and let $\widehat\Gamma_m$ be its normalization.
This definition is independent of whether normalization of the global
pencil graph commutes with this particular slice; no such commutation
is assumed.

\begin{theorem}[The full contact-slice Hilbert modification]
\label{thm:contact-hilbert-v161}
For every $m\geq2$ the slice graph is
\begin{equation}\label{eq:contact-base-blowup-v161}
 \widehat\Gamma_m\simeq\operatorname{Bl}_0\A^2
     =\operatorname{Bl}_{\Fitt_0(q_*\OO_D)}S.
\end{equation}
Its whole exceptional fibre is $\Pj^1$; the Hilbert curves over
that exceptional fibre are nonreduced. In a base chart $u=t$, $v=tw$,
the total multiplication graph is the blow-up of $(x^m,t)$ on the
pulled-back pencil surface. With $R=\C[x,t,w]$, its Rees presentation
and its two projective charts are
\begin{equation}\label{eq:contact-rees-v161}
 \begin{aligned}
 \mathcal R(x^m,t)&=R[U,V]/(tU-x^m V),\\
 D_+(U):\quad t&=x^m y,\qquad D_+(V):\quad x^m=t z.
 \end{aligned}
\end{equation}
Here $U$ represents $x^mT$ and $V$ represents $tT$ in $R[T]$.
The charts glue by $y=z^{-1}$. The universal curve is flat over the
base blow-up, is normal, and has a transverse $A_{m-1}$ surface
singularity along the section $x=t=z=0$ of its exceptional divisor.

For every direction $[\alpha:\beta]\in\Pj^1$, its special Hilbert
curve is
\begin{equation}\label{eq:thick-fibre-v161}
 C_0\cup T_m,\qquad
 T_m\simeq\Pj^1_{\C[x]/(x^m)},\qquad
 C_0\cap T_m\simeq\Spec\C[x]/(x^m).
\end{equation}
The intersection is at one section of $T_m$; the support of $T_m$
is the residual line through $[I_2]$ with second member
$\left(\begin{smallmatrix}\alpha&\beta\\\beta&-\alpha\end{smallmatrix}\right)$
in the exceptional plane. The fibre has no embedded associated points.
Its cycle is $[C_0]+m[F]$, where $F=(T_m)_{\rm red}$. Every direction
occurs, and these embedded thickened curves are pairwise distinct.
The multidegrees are
\[
 \deg_{C_0}H_q=q-m\mathbf1_{q=n-1},\qquad
 \deg_F H_q=\mathbf1_{q=n-1}.
\]
For $\delta_j=(j-h(n-2))_+-2(j-h(n-1))_+$, the $j$th resolved
multiplication system on $T_m$ is, over $\C[x]/(x^m)$, the complete
system $\OO_{T_m}(\delta_j)$. In particular its reduced degree is
$\delta_j$ and its degree as a cycle is $m\delta_j$.
\end{theorem}
\begin{proof}
All rank loci below corank two are avoided, so the relevant ambient
complete-quadric map is the blow-up of the rank-$(n-2)$ locus.
Pull the exact ideal in Lemma~\ref{lem:jordan-slice-v161} to
$\operatorname{Bl}_0 S$. In the stated chart it is $(x^m,t)$.
These two elements are a regular sequence. The only homogeneous Rees
relation is $tU-x^m V$: successive comparison of powers of $T$,
using that $x^m$ and $t$ are relatively prime, expresses every relation
as a multiple of this one. Thus the relation is saturated in the image
subalgebra of $R[T]$, and the stated Proj charts follow. The other base
chart is identical with the roles of $u,v$ exchanged.

Both chart equations remain nonzerodivisors on every fibre of their
smooth relative affine planes, so the fibrewise Cartier flatness
criterion proves flatness. The first chart eliminates $t$ and is
smooth. The second is a hypersurface with singular locus
$x=t=z=0$, of codimension two in its total space. It is a domain and
Cohen--Macaulay, so regularity in codimension one proves normality.
Its transverse completed equation is $x^m=tz$, the $A_{m-1}$ surface
singularity. This describes a singular universal curve over a smooth
parameter blow-up; the two smoothness assertions are not conflated.

At $t=0$ the first chart has ring $\C[x,y]/(x^m y)$ and the second
has ring $\C[x,z]/(x^m)$. In the first ring the ideals $(y)$ and
$(x^m)$ intersect in $(x^m y)$. These pieces glue to the strict
transform $C_0$ and to the thick tail whose two rings are
$(\C[x]/(x^m))[y]$ and $(\C[x]/(x^m))[z]$, with $y=z^{-1}$ and
$x$ unchanged. This proves \eqref{eq:thick-fibre-v161}, including its
length-$m$ intersection. The local hypersurfaces of the fibre are
Cohen--Macaulay of dimension one and have only their component primes
as associated primes. There are no embedded points.

The Rees quotient \eqref{eq:rees-surjection-v161} embeds this flat
family into the fixed complete-quadric space times the base blow-up.
Its generic curve is the prescribed lift. The support of each tail
is the indicated residual line. Distinct $[\alpha:\beta]$ give
distinct residual lines, hence distinct embedded closed subschemes.
The resulting proper map to $\Gamma_m$ is quasi-finite and birational,
so the smooth base blow-up is its normalization. This proves the
whole fibre statement in \eqref{eq:contact-base-blowup-v161}. The
Fitting formula follows from $(u,v)^m$ and invariance of Proj under
a positive Veronese power.

The last minor system loses a base divisor of length $m$ on $C_0$,
while the smaller systems lose none. The graph line bundle restricts
to $\OO(1)$ on the residual projective line over $\C[x]/(x^m)$.
It follows that $\chi(\OO_{T_m}\otimes H^l)=m(l+1)$ and the
intersection has Euler characteristic $m$. Adding to the main
component gives $\binom n2l+1$, as required. More generally the
unit-block congruence is invertible over $\C[x]$ and thus over
$\C[x]/(x^m)$. The binary multiplication coefficient calculation of
Corollary~\ref{cor:multi-powers-v161}, extended to this Artin ring,
spans every degree-$\delta_j$ binary form. Multiplication by the
unit-block socle is a split injection on these coefficient modules:
it has a unit coordinate in a trivializing frame. This proves the
relative complete-system assertion, not just its reduced-fibre version.
The degree-$j$ coordinates originally have the twist $\OO(j)$ on
the pencil line; it is trivial on this local thickening. Removing
their Cartier common factors produces the stated graph line bundle.
\end{proof}

\subsection{The singularity and the stable-map comparison}
\begin{proposition}[Local deformation module of the multiplication graph]
\label{prop:contact-deformations-v161}
At a singular point of a one-parameter direction in
Theorem~\ref{thm:contact-hilbert-v161}, let
$A_m=\C[[x,t,z]]/(x^m-tz)$. The abstract local cotangent deformation
modules are
\begin{equation}\label{eq:contact-T1-v161}
 T^1(A_m)=\C[[x]]/(x^{m-1}),\qquad T^i(A_m)=0\quad(i\geq2).
\end{equation}
The first equality has dimension $m-1$. A formally versal family of
hypersurface deformations is obtained by adding
$\sum_{a=0}^{m-2}c_a x^a$ to $x^m-tz$. These are deformation modules
of the total graph singularity, not of its fixed-base relative curve
or of the raw failure algebra.

In a transverse surface the reduced tail is $\Q$-Cartier with
$F^2=-1/m$. Its minimal resolution introduces the $A_{m-1}$ chain;
the total fibre has multiplicities $m$ on the strict transform of
$F$ and $m-1,m-2,\ldots,1$ along the consecutive chain, and
multiplicity one on $C_0$.
\end{proposition}
\begin{proof}
The cotangent complex of a hypersurface is represented by the
conormal free module mapping to the three ambient differentials.
Dualizing gives cokernel by the partial derivatives
$mx^{m-1},-z,-t$, hence \eqref{eq:contact-T1-v161}; there are no higher
terms. Lifting an equation across a small extension has no obstruction,
and the displayed monomials form a basis of the cokernel, giving the
formal versality assertion. This is the classical hypersurface
calculation applied to the exact Rees equation.

In the first Rees chart the exceptional Cartier ideal is generated
by $x^m$, so its Cartier divisor is $mF$. Its associated ideal bundle
restricts to $\OO_F(1)$; thus $mF^2=-1$. The standard cyclic quotient
presentation can be checked directly:
\[
 x=ab,\quad t=a^m,\quad z=b^m,
 \qquad A_m=\C[[a,b]]^{\mu_m},\quad
 (a,b)\mapsto(\zeta a,\zeta^{-1}b).
\]
Subdivide its two-dimensional cone by the primitive rays
$((m-i)/m,i/m)$, $1\leq i<m$, in the quotient lattice. Consecutive
cones are smooth. The internal curves form the chain with
self-intersection $-2$, and the orders of $t=a^m$ are $m-i$.
The end attached to the tail has multiplicity $m-1$; the next ones
decrease by one. The main component meets the other, already smooth,
end of the tail, where $t=x^m y$ gives multiplicities $m,1$.
These charts prove the stated fibre description.
\end{proof}

\begin{theorem}[A thick Hilbert tail and a covered stable-map tail]
\label{thm:stable-tail-v161}
Fix a nonzero direction $(u,v)=\tau(\alpha,\beta)$ in the slice,
with contact order $m$. After the base change $\tau=s^m$, the
normalization of its total Hilbert graph is the blow-up of the smooth
pencil surface at $(x,s)=(0,0)$. Its special fibre is a reduced nodal
union $C_0\cup F'$. The induced map to complete quadrics is an
embedding on $C_0$ and a degree-$m$ map on $F'$ onto $F$,
\begin{equation}\label{eq:covered-tail-v161}
 [X:S]\longmapsto[X^m:S^m]
\end{equation}
in the residual line coordinates. It is a stable map with tail deck
group $\mu_m$. It has the same cycle $[C_0]+m[F]$ as the Hilbert
limit, but the latter has the nonreduced scheme structure
\eqref{eq:thick-fibre-v161}, not the reduced image of this map.
\end{theorem}
\begin{proof}
The base-changed graph is the blow-up of $(x^m,s^m)$. This particular
base change is flat, and the Rees presentation also verifies it
directly. On the chart $x^m=s^m z$ adjoining $w=x/s$ gives the
integral element $w^m=z$ and the normal chart $\C[s,w]$, with
$x=sw$. On the other chart $s^m=x^m y$ adjoining $r=s/x$ gives
$r^m=y$ and the normal chart $\C[x,r]$, with $s=xr$. These charts
are finite and birational over the original charts and glue by
$w=r^{-1}$. They are exactly the point blow-up charts. They also
show a reduced nodal fibre after normalization.

The old graph coordinates $[x^m:s^m]$ restrict to
\eqref{eq:covered-tail-v161} on the exceptional line of the point
blow-up. Away from the contact point normalization does not alter
the original lifted line. Both components are nonconstant; the
main component is embedded and the nonconstant tail has precisely
the finite deck group $\mu_m$, which fixes the attachment point.
Thus the map is stable. The image of $F'$ is reduced while the flat
Hilbert fibre has the explicit nilpotent $x$ of order $m$, so the
two objects differ scheme-theoretically even though pushforward
cycles and degrees agree. This comparison is calculated here; it is
not deduced from a general identification of Hilbert and stable-map
compactifications.
\end{proof}

For the reduced-incidence local model ($m=1$), the two local
descriptions coincide with the nodal model. For every $m\geq2$,
the length of the contact, the tail multiplicity, the transverse
graph singularity $A_{m-1}$, and the stable tail covering degree
record the same integer in four different objects. The assertions
concern all orders in the specified symmetric Jordan slices. They
do not classify arbitrary higher-corank contacts or the entire
ambient Hilbert fibre over a nonreduced-incidence pencil.

\section{Effective boundary groupoids and relative lifting obstructions}
\label{sec:effective-boundary-v161}
This section places the intrinsic envelope and the boundary modification
in a single category. It also distinguishes three deformation questions:
deforming the effective failure object, lifting a prescribed base
deformation to its boundary modification, and deforming a singularity
of the universal curve. These questions have different answers.

\subsection{The category and the envelope functor}
Let $\mathscr F_{\rm pen}^{\rm eff}$ denote the fppf stack of
Theorem~\ref{thm:pencil-stack-equivalence}. Starting with the sharp
pencil-failure families, its arrows are first divided by the
tangent-identity nonlinear generator kernel and then by the ineffective
right tensor-factor action, and the resulting groupoid is stackified.
Its oriented source is specified by coefficient support, rather than
by an arbitrary choice between the determinant rulings. The theorem
gives, on every complex base,
\begin{equation}\label{eq:effective-stack-v161}
       \mathscr F_{\rm pen}^{\rm eff}\simeq[G/\operatorname{PGL}(V)].
\end{equation}
The notation does not denote the stack of all Artin algebras of the
same length. In particular deformations leaving the effective
pencil image are not included in this equivalence.

\begin{proposition}[The envelope as a morphism of stacks]
\label{prop:envelope-stack-v161}
Let $\mathscr A_h$ be the stack whose objects over $T$ are a
projective-space bundle $S\to T$ of relative dimension $n-1$ and a
finite locally free graded algebra bundle on $S$ locally isomorphic
to $\bigoplus_{p=0}^{nh}B_h(V)_p\otimes\OO(2p)$, with its
projectivized power diagram. Isomorphisms preserve the grading,
algebra structure, and diagram. The construction of
Theorem~\ref{thm:intrinsic-envelope-v158} defines a morphism
\[
                \mathscr F_{\rm pen}^{\rm eff}\longrightarrow\mathscr A_h.
\]
It is compatible with arbitrary complex base change and does not
require a global line bundle $\OO_S(1)$. The case $h=1$ is
distinguished without an additional exponent choice. The envelope
is a functorial algebra bundle on the recovered source, not a
claimed subalgebra or quotient of the original finite algebra.
\end{proposition}
\begin{proof}
An object on the right side of \eqref{eq:effective-stack-v161} is a
$\operatorname{PGL}(V)$-torsor with an equivariant map to $G$.
On a lift of its source to $\Pj(V)$, the first-jet construction is
$\mathcal Hom(J^1\OO(1),\OO(1))=V\otimes\OO(1)$. A scalar
$\lambda$ acts on $B_h(V)_p$ with weight $2p$ and on $\OO(2p)$
with weight $-2p$. It therefore acts trivially on their product,
on its multiplication, and on the polarized Cartan relation maps.
The algebra bundle is genuinely $\operatorname{PGL}(V)$-equivariant,
not only equivariant up to a scalar. Associate it to the torsor and
apply fppf descent. This defines both objects and arrows of the
asserted morphism and gives its compatibility with pullback. This
argument uses local $\OO(1)$ only to exhibit a scalar-trivial
representation; it never posits such a global line bundle on a
Brauer-nonsplit projective-space bundle. The universal property of
Proposition~\ref{prop:universal-envelope-v159} verifies the same
construction independently of the chosen trivializations.
\end{proof}

\subsection{A modular modification and its defining equations}
Restrict \eqref{eq:effective-stack-v161} to $U_{\rm red}$ of
\eqref{eq:reduced-open-v161}. With $B_{\rm red}$ and its invariant
Fitting ideal as in \eqref{eq:fitting-centre-v161}, put
\begin{equation}\label{eq:boundary-stack-v161}
 \mathscr B_{\rm red}=[B_{\rm red}/\operatorname{PGL}(V)],\qquad
 \mathscr F_{\rm red}^{\rm eff}=[U_{\rm red}/\operatorname{PGL}(V)].
\end{equation}

\begin{theorem}[The effective multiple-incidence boundary]
\label{thm:boundary-stack-v161}
The map $\mathscr B_{\rm red}\to\mathscr F_{\rm red}^{\rm eff}$
is representable, proper, birational, and carries the flat embedded
curve family of Theorem~\ref{thm:multiple-incidence-v161}. Both stacks
are smooth. This is a modular modification over $U_{\rm red}$, not
an assertion that either stack is proper over $\C$ or has finite
stabilizers. On an etale neighbourhood with labelled incidence
centres $I_i=(u_i,v_i)$, the multi-Rees algebra has presentation
\begin{equation}\label{eq:multi-Rees-v161}
 \mathcal R(I_1,\ldots,I_k)
 =R[U_1,V_1,\ldots,U_k,V_k]/(v_iU_i-u_iV_i)_{i=1}^k,
\end{equation}
where $\deg U_i=\deg V_i$ is the $i$th coordinate vector and
$U_i=u_iT_i$, $V_i=v_iT_i$. There is no torsion kernel left in this
presentation. Its diagonal subalgebra is
\begin{equation}\label{eq:diagonal-Rees-v161}
 \bigoplus_{d\geq0}(I_1\cdots I_k)^d(T_1\cdots T_k)^d,
\end{equation}
so its diagonal Proj is the computed Hilbert modification.
All these objects descend without a global labelling of the points.
\end{theorem}
\begin{proof}
For a test scheme $T\to[U_{\rm red}/\operatorname{PGL}(V)]$, take
its defining torsor $P_T$. The pullback is the associated space
$(P_T\times_{U_{\rm red}}B_{\rm red})/\operatorname{PGL}(V)$.
The projective equivariant morphism descends to a projective
morphism over $T$. This proves representability and properness;
birationality holds over the dense nonincident open. Smoothness of
the two atlases and of the acting group proves smoothness of the
stacks. The curve family descends by the same construction. No
quotient by the stabilizer of a pencil has been mistaken for its
fixed-frame scheme fibre.

For the algebra calculation begin in a polynomial ring in the
independent parameters $u_i,v_i$ and complementary smooth
coordinates. For one pair, the kernel of
$R[U_i,V_i]\to R[T_i]$ is generated by $v_iU_i-u_iV_i$:
comparison of successive powers and cancellation using the regular
sequence $(u_i,v_i)$ proves the assertion. For several pairs the
variables are disjoint. Tensoring the one-pair presentations over
the ground field preserves their injections into the polynomial
ring in the $T_i$, giving \eqref{eq:multi-Rees-v161}. The result
persists after the etale flat extension supplied by
Lemma~\ref{lem:general-incidence-v161}; flat base change preserves
both the coefficient ideals and the displayed kernel. This proves
the absence of an additional saturation or torsion relation.
The multidegree $(d,\ldots,d)$ part is exactly
$I_1^d\cdots I_k^d(T_1\cdots T_k)^d$, which proves
\eqref{eq:diagonal-Rees-v161}. It gives the actual defining equations
before taking the ordinary-power diagonal, rather than treating the
formation of ordinary powers as a second recovery theorem.

The ideal itself is $\Fitt_0(q_*\OO_D)$, where $D$ is the zero
scheme of the distinguished power $J_{n-1}$ on the intrinsic pencil
line. Finite base change identifies the pushforward module, and
Fitting ideals commute with arbitrary base change. Its formation
is invariant under the group action. The local multi-Rees
presentations are labelled charts of this global ideal blow-up;
permutation of the labels does not affect their diagonal. Descent
therefore follows from the intrinsic Fitting construction, without
claiming that graph closure or blow-up commutes with every base
change. A pulled-back curve family is defined by its flat pullback;
its existence does not require such a graph-closure assertion.
\end{proof}

\subsection{Relative lifting is not absolute obstruction}
For precision a deformation of an effective object at $R$ over a
local Artin complex algebra $A$ is an object of
$\mathscr F_{\rm pen}^{\rm eff}(A)$ with an identification of its
closed fibre with $R$. A source framing identifies it with a point
of the completed pencil parameter scheme; changing the framing
acts by $\operatorname{PGL}(V)(A)$. A boundary deformation also
specifies a lift to $B_{\rm red}$ of the closed boundary point.

\begin{proposition}[Exact relative lifting equations]
\label{prop:relative-obstructions-v161}
Fix a boundary point with $k$ incidences. In a labelled chart of the
base blow-ups write $v_i=u_iw_i$. Let $A'\twoheadrightarrow A$ be
a square-zero extension with kernel $J$. Fix a lift of the base
object to $A'$, with coordinates $u_i',v_i'$, and an existing lift
to the boundary over $A$, with slopes $w_i$. For any lifts
$\widetilde w_i\in A'$, define
\begin{equation}\label{eq:relative-obstruction-v161}
 o_i=[v_i'-u_i'\widetilde w_i]\in J/u_i'J.
\end{equation}
These classes do not depend on the lifts $\widetilde w_i$.
The prescribed base deformation lifts in this chart if and only if
all $o_i$ vanish. When nonempty, its set of framed lifts is a torsor
under $\bigoplus_i\ker(u_i':J\to J)$.
At a closed boundary point with slopes $w_i^0$, a tangent vector
$u_i=\epsilon a_i$, $v_i=\epsilon b_i$ lifts if and only if
$b_i=w_i^0a_i$. The kernel and cokernel of the differential of the
fixed-frame blow-up map both have dimension $k$ there.

The absolute effective and boundary stacks are smooth, with tangent
complexes $[\mathfrak{pgl}(V)\to T_RU_{\rm red}]$ and
$[\mathfrak{pgl}(V)\to T_bB_{\rm red}]$ in degrees $-1,0$.
Thus the relative conditions above do not assert an absolute
obstruction to deforming an effective failure object. Nor do they
identify all deformations of its raw unmarked Artin algebra.
\end{proposition}
\begin{proof}
The defining equation over $A'$ is $v_i'=u_i'w_i'$.
Its defect for a chosen lift lies in $J$ because it vanishes over
$A$. Replacing that lift by $\widetilde w_i+\eta_i$, with
$\eta_i\in J$, changes the defect by $-u_i'\eta_i$.
This proves the well-defined obstruction class, its necessity and
sufficiency, and the stated torsor of solutions. The independent
coordinate charts introduce no further equations. For dual numbers
at the closed point, multiplication of $\epsilon a_i$ by a change
$\epsilon\eta_i$ is zero, giving precisely the displayed tangent
condition. In the chart $(u_i,w_i)\mapsto(u_i,u_iw_i)$ the
differential has rank one for each normal parameter pair and kills
its slope direction. Complementary coordinates are unchanged, so
both kernel and cokernel have dimension $k$.

The tangent complexes are those of the smooth quotient atlases in
\eqref{eq:boundary-stack-v161}, and
Theorem~\ref{thm:pencil-stack-equivalence} identifies the effective
failure deformation groupoid with the first of them over every
Artin base. To check explicitly the nonlinear kernel used here,
let the coefficient ring be $\C[\epsilon]/(\epsilon^2)$ and replace
a generator $x$ by $gx+q(x)$, with $q$ of degree at least two. In a
homogeneous relation of sharp degree $d$, every term containing a
factor from $q$ has degree at least $d+1$ and vanishes in the sharp
truncation. The surviving relation is exactly the linear action
$\Sym^d(g)$. This argument uses degrees in the generators, not
powers of $\epsilon$, and is valid for coefficients in any base
algebra. It therefore separates the tangent-identity generator
kernel from an actual infinitesimal change of the pencil over the
base. Passing to the two prescribed rigidifications gives the
claimed categorical interpretation of the coordinate calculation.
\end{proof}

The ideal cutting out the incidence is scheme-theoretically intrinsic
to the effective failure family \emph{with its recovered pencil line}.
That line is an essential output of the inverse; no bypass of it, and
no canonical choice of a slope from a closed algebra, is asserted.
The additional content is the complete multi-Rees presentation and
relative lifting equations, and, on the nonreduced-contact slices,
the singular graph and its computed cotangent module in
Proposition~\ref{prop:contact-deformations-v161}. These calculations
specify the interaction with boundary singularities rather than
merely transporting a pre-existing finite algebra by pullback.

\subsection{An accessible completeness input for singular pencils}
The congruence classification used by
Theorem~\ref{thm:singular-complete-data-v158} is recalled with explicit
blocks in \cite[Theorem~2.1]{DeTeranDmytryshynDopico2018v161}.
In that statement a symmetric minimal-index block with index
$\varepsilon$ has dimension $2\varepsilon+1$; hence no index in an
$n\times n$ symmetric pencil exceeds $\lfloor(n-1)/2\rfloor$.
The Toeplitz nullities in that finite range recover the
multiplicities by the second-difference formula already proved in
our theorem. The regular blocks are indexed by their projective
spectral points and local exponent partitions. Accordingly
``complete data'' includes both the points and the partitions,
modulo one common projective reparametrization of the pencil line.
This is an accessible theorem-level classical completeness input,
not a new symmetric Kronecker classification. It is distinct from
the local congruence diagonalization proved in the paper.

\section{Versal contact coordinates for regular symmetric pencils}
\label{sec:versal-contacts-v162}
The contact order does not determine a neighbourhood of a pencil.
We describe the missing parameters and then use them to calculate
open charts of the ambient Hilbert boundary. The normal-form input
belongs to the classical deformation theory of symmetric pencils
\cite[Theorem~2.1, Corollary~2.1 and Lemma~3.1]{Dmytryshyn162}.
We give the tangent calculation and the formal reduction needed here,
because a normal form for a pencil and a normal form for its embedded
graph are different assertions.

All completions in this section are over $\C$. The pencil parameter
is fixed; reparametrizing it is not part of the local congruence action.
A source-and-parameter framing is a smooth atlas of the Grassmannian
of genuine pencils. Statements on this atlas descend by change of
frame, rather than by identifying different embedded Hilbert points.

\begin{lemma}[The contact deformation module]
\label{lem:contact-module-v162}
Put $D=\diag(x^{d_1},\ldots,x^{d_c})$, where $d_i\geq1$.
For deformations of this symmetric matrix germ under congruences
which reduce to the identity, the tangent quotient is
\begin{equation}\label{eq:contact-module-v162}
 \mathcal T_D=
 \bigoplus_i\C[x]_{<d_i}\ \oplus\!
 \bigoplus_{i<j}\C[x]_{<\min(d_i,d_j)}.
\end{equation}
A formally miniversal family is the symmetric polynomial matrix
\begin{equation}\label{eq:versal-matrix-v162}
 S_{ii}=x^{d_i}+p_{ii}(x),\qquad
 S_{ij}=p_{ij}(x),\quad
 \deg p_{ii}<d_i,\quad\deg p_{ij}<\min(d_i,d_j).
\end{equation}
Its coefficients vanish at the origin. Miniversality means existence
of a formal classifying map and a congruence after every complete
local base change, with an isomorphism on the tangent quotient. It
does not assert uniqueness of that map in the presence of stabilizers.
\end{lemma}
\begin{proof}
An infinitesimal change $1+\epsilon H$ adds $H^{\mathsf t}D+DH$.
Its diagonal entries generate $(x^{d_i})$, and its $(i,j)$ entry
ranges over $(x^{d_i},x^{d_j})=(x^{\min(d_i,d_j)})$.
Different unordered entries can be prescribed independently. This
proves \eqref{eq:contact-module-v162}, including the factor two on
the diagonal, which is invertible.

For completeness let $(R,\mathfrak m)$ be complete local and let
$S(x)\equiv D\pmod{\mathfrak m}$. Work successively modulo
$\mathfrak m^{q+1}$. Divide each coefficient of the degree-$q$
error in an off-diagonal position by $x^{\min(d_i,d_j)}$, and each
diagonal error by $x^{d_i}$. Put the remainders into the indicated
polynomials. Remove the quotients with $H$ of order $q$. On an
off-diagonal entry use the side with the smaller exponent; on the
diagonal divide by two. Terms involving an already deformed
coefficient and $H$ have order at least $q+1$. The compatible products
of these changes converge $\mathfrak m$-adically, as do the finitely
many polynomial coefficients. This constructs the formal reduction.
The tangent quotient proves minimality of the number of parameters.
The choices made by this algorithm are coordinates, not invariants.
\end{proof}

\begin{theorem}[Ambient realization of the contact parameters]
\label{thm:ambient-versal-v162}
Let $R_0$ be a regular complex symmetric pencil. At each spectral
point $\lambda$, let $d_{\lambda,1},\ldots,d_{\lambda,c_\lambda}$
be its positive Smith exponents. On a source-and-parameter frame atlas,
the completed pencil base has a formally smooth map to the product
of the coefficient spaces \eqref{eq:versal-matrix-v162}, with
independent coordinates in all factors. After a congruence on each
completed parameter disc, its residual symmetric matrix is exactly
\eqref{eq:versal-matrix-v162}; the complementary block is a unit.
The number of essential congruence parameters is
\begin{equation}\label{eq:versal-dimension-v162}
 N=\sum_{\lambda,i}d_{\lambda,i}
       +\sum_\lambda\sum_{i<j}\min(d_{\lambda,i},d_{\lambda,j}).
\end{equation}
The remaining parameters are smooth frame-orbit directions. Thus
these models describe deformations in the full regular-pencil
neighbourhood, not only a selected two-parameter slice.

For one corank-two point with exponents $a\leq b$, the residual
matrix can be written
\begin{equation}\label{eq:Euclidean-normal-v162}
 S(x)=\begin{pmatrix}f(x)&g(x)\\g(x)&f(x)k(x)+r(x)\end{pmatrix},
\end{equation}
where $f$ is monic of degree $a$, $g,r$ have degree less than $a$,
and $k$ is monic of degree $b-a$. For $a=b$, $k=1$.
The central values are $f=x^a$, $g=r=0$, and $k=x^{b-a}$.
These are $2a+b$ independent formal coordinates. With
$A_f=\OO[x]/(f)$ and multiplication matrices $M_g,M_r$ on its
basis $1,x,\ldots,x^{a-1}$, the exact incidence module and ideal are
\begin{equation}\label{eq:versal-incidence-v162}
 q_*\OO_D=\operatorname{coker}[M_g\ M_r],\qquad
 \Fitt_0(q_*\OO_D)=I_a([M_g\ M_r]).
\end{equation}
These identities hold over arbitrary base algebras in this polynomial
model. They do not identify its nonreduced Hilbert modification with
the blow-up of that Fitting ideal.
\end{theorem}
\begin{proof}
Choose a nonsingular pencil member. Fixing it by congruence identifies
the other member with a self-adjoint endomorphism $T$ for a fixed
nondegenerate symmetric form. Its primary subspaces at distinct
eigenvalues are a direct sum, even when $T$ is not semisimple.
Within one primary subspace the symmetric Jordan presentation is
$\bigoplus_i L_{d_i}(x)$, where
$L_d(x)=E_d(xI-N_d)$ as in Lemma~\ref{lem:jordan-slice-v161}.
If $v_d(x)=(1,x,\ldots,x^{d-1})^{\mathsf t}$, then
\begin{equation}\label{eq:kernel-vector-v162}
                  L_d(x)v_d(x)=x^d e_1.
\end{equation}
The complement to $e_1$ is a unit block. Differentiating its Schur
complement in the direction of a constant symmetric perturbation
$H$ gives, on an $(i,j)$ block,
$v_{d_i}^{\mathsf t}H_{ij}v_{d_j}$.
For $0\leq l<\min(d_i,d_j)$ the perturbation
$H_{ij}=e_1e_{l+1}^{\mathsf t}$ gives $x^l$; insert its transpose
in the opposite block. On a diagonal block use
$(e_1e_{l+1}^{\mathsf t}+e_{l+1}e_1^{\mathsf t})/2$ for $l>0$
and $e_1e_1^{\mathsf t}$ for $l=0$. They give every coefficient
in \eqref{eq:contact-module-v162}. Block-diagonal perturbations
prescribe these coefficients independently at distinct eigenvalues.
This proves surjectivity from actual pencil deformations, not merely
from arbitrary matrix-valued power series.

Here is also the dimension check against the congruence orbit.
The centralizer between Jordan blocks of sizes $d_i,d_j$ at the
same eigenvalue has dimension $\min(d_i,d_j)$, by solving
$N_{d_i}X=XN_{d_j}$ along its diagonals. A commuting endomorphism
on a single block is polynomial in $N_d$, hence self-adjoint.
A skew-adjoint commuting endomorphism has no diagonal contribution;
its two opposite off-diagonal blocks determine one another and
contribute $\min(d_i,d_j)$. Distinct primary eigenvalues contribute
zero, since the corresponding Sylvester map is invertible. Therefore
the orthogonal stabilizer has the second sum in
\eqref{eq:versal-dimension-v162} as its dimension. Subtracting its
orbit dimension from $n(n+1)/2$ gives $n$ plus that sum, namely $N$.
Congruence directions are killed by the contact quotient. Surjectivity
and equality of dimensions identify their kernel exactly.

Apply the formal reduction of Lemma~\ref{lem:contact-module-v162}
to the universal pencil. Its classifying coefficient map has the
surjective differential just computed. Complete its coefficients
by coordinates in the kernel. The formal inverse function theorem
identifies the completed smooth base with these coefficients and
the complementary power-series variables. This is the claimed formal
smoothness. Equivalently a transverse affine pencil family is
miniversal by the congruence tangent-space criterion in
\cite[Lemma~3.1]{Dmytryshyn162}. Normalizing the chosen nonsingular
member and then undoing that normalization restores the smooth
source-and-parameter framing directions.

For two blocks the first normal form is
$\left(\begin{smallmatrix}f&g\\g&h\end{smallmatrix}\right)$ with
$h$ monic of degree $b$. Euclidean division $h=fk+r$ is a polynomial
change of coefficient coordinates with polynomial inverse. This
proves \eqref{eq:Euclidean-normal-v162} and its dimension. The
submaximal-minor ideal on the residual chart is $(f,g,h)=(f,g,r)$.
Quotienting first by the monic polynomial $f$ makes a free module
of rank $a$, and quotienting further by $g,r$ gives the displayed
cokernel. The Fitting formula follows directly from this presentation.
No reduced-fibre argument or extraction of an ideal root is involved.
\end{proof}

The theorem includes higher corank in its normal-coordinate statement.
The explicit Hilbert charts below concern corank two. The all-pencil
inverse and the separate singular-pencil invariant theorem are not
restricted by this choice. In particular the balanced two-parameter
Jordan slice has only two directions inside a $3a$-parameter contact
base; it is not a versal neighbourhood when $a>1$.

\section{Euclidean charts of the ambient Hilbert boundary}
\label{sec:Euclidean-boundary-v162}
We first calculate a universal polynomial graph in a fixed target.
This separates the Hilbert calculation from the choice of pencil
normal coordinates. Write $\C[x]_{<a}$ for polynomials of degree
less than $a$, including zero. Let
\[
 B_a=\{(f,g,r): f\text{ monic of degree }a,\quad
                          g,r\in\C[x]_{<a}\}\simeq\A^{3a}.
\]
Homogenize the three polynomials to degree $a$. They have no common
zero at infinity, where their value is $[1:0:0]$. Over the dense open
where they have no common zero, their graph in $\Pj^1\times\Pj^2$
has Hilbert polynomial $(a+1)l+1$ for $\OO(1,1)$. Denote its reduced
Hilbert graph closure over $B_a$ by $\Gamma_a$. It is integral,
since it is the reduced closure of the graph over an integral dense
open. We never quotient its embedded curves by target automorphisms.

\begin{theorem}[An open division chart]
\label{thm:division-chart-v162}
Let $W_a=\{(f,g,w): f\text{ monic of degree }a,
                   \ g,w\in\C[x]_{<a}\}\simeq\A^{3a}$ and define
\begin{equation}\label{eq:division-map-v162}
 \mu:W_a\longrightarrow B_a,
       \qquad (f,g,w)\longmapsto(f,g,r=\operatorname{rem}_f(gw)).
\end{equation}
Put $q=(gw-r)/f$. The graph modification on the universal parameter
line over $W_a$ is the blow-up of $(f,g)$, and its two-generator
Rees equation is
\begin{equation}\label{eq:division-Rees-v162}
                  gU-fV=0.
\end{equation}
It is a projective flat family embedded in
$\Pj^1\times\Pj^2\times W_a$, by the third graph coordinate
\begin{equation}\label{eq:third-coordinate-v162}
                  [F:G:R]=[U:V:wV-qU]
\end{equation}
on the finite parameter chart. At infinity it is the unique original
graph. There is an open neighbourhood of
\[
         \{f=x^a,\ g=0,\ w\in\C[x]_{<a}\}\subset W_a
\]
on which this family identifies $W_a$ with an open subscheme of
$\Gamma_a$. In particular $\Gamma_a$ is smooth and already normal
there. The assertion is an open chart of the Hilbert graph, not just
a generically injective parametrization of some limits.

For any point of this chart, let $d=\gcd(f,g)$ be monic, let
$r_0=\deg d$, and put $f=df_1$, $g=dg_1$. Its embedded fibre has
the scheme decomposition
\begin{equation}\label{eq:gcd-fibre-v162}
 C\cup T_d,\qquad T_d=\Pj^1_{\C[x]/(d)},\qquad
 C\cap T_d=\Spec\C[x]/(d).
\end{equation}
The attaching section on $T_d$ is $[U:V]=[f_1:g_1]$.
It has no embedded associated points. If
$d=\prod_i(x-\alpha_i)^{m_i}$, this is one main component and
thick tails of multiplicities $m_i$ at the distinct $\alpha_i$.
Its graph degree on the main component is $a-r_0$, and its tail
cycle degree is $r_0$. Thus the same open chart describes splitting,
collision, and simultaneous unequal thicknesses by an exact gcd law.
\end{theorem}
\begin{proof}
In the polynomial ring of the universal line and $W_a$, the elements
$f,g$ form a regular sequence: $f$ is monic, and modulo $f$ the
coefficient $g_0$ still occurs freely and linearly in $g$. The Rees
algebra is consequently the symmetric algebra modulo the single
syzygy $gU-fV$. One can alternatively compare coefficients in its
injection into the polynomial ring in a Rees variable. On every
geometric fibre the homogeneous equation is nonzero, because $f$
is monic. The fibrewise Cartier criterion therefore proves flatness
\cite[Tag~062Y, part~(3)]{Stacks162}. The original ideal $(f,g,r)$ is exactly
$(f,g)$, since $r=gw-fq$. Formula~\eqref{eq:third-coordinate-v162}
is a rank-two linear embedding on every finite parameter chart;
the first two coordinates are unchanged. Thus the Rees graph is
closed in the three-coordinate graph. Infinity is disjoint from
the base ideal and needs no blow-up. This constructs the claimed
flat embedded family and its morphism to the Hilbert scheme.

Its image lies in $\Gamma_a$: over $\operatorname{Res}(f,g)\ne0$
it is the original graph, that open is dense in the integral scheme
$W_a$, and the reduced graph is closed. To prove that the resulting
map is an open chart we construct a regular inverse. This step
cannot be replaced by pointwise injectivity alone.

Let $Z_w$ be a curve with $f=x^a,g=0$ in the family, and let
$L=\OO_{\Pj^1\times\Pj^2}(a-1,1)|_{Z_w}$. On its main component
$L$ has degree $a-1$. On its thick tail it is $\OO(1)$ over
$A=\C[x]/(x^a)$, and their intersection has length $a$. The
restriction map from tail sections onto the intersection is
surjective. The union exact sequence therefore gives
\begin{equation}\label{eq:chart-cohomology-v162}
                    H^1(Z_w,L)=0,\qquad h^0(Z_w,L)=2a.
\end{equation}
These statements hold in a Hilbert neighbourhood by cohomology and
base change. On that neighbourhood consider the linear map to the
rank-$2a$ bundle of sections
\begin{equation}\label{eq:division-recovery-v162}
 \Psi:\C[x]_{<a}\oplus\C[x]_{<a-1}\longrightarrow H^0(L),
                   \quad (w,q)\longmapsto wG-qF.
\end{equation}
The expressions are homogenized to degree $a-1$ in the first factor;
in particular the second polynomial has zero leading coefficient.
For $a=1$ its second summand is zero. At $Z_w$ this map is injective.
Indeed restriction to the main component makes $q$ zero. Restriction
to the tail then makes $w$ zero modulo $x^a$, hence zero as a
polynomial of degree less than $a$. After shrinking, $\Psi$ is a
split injection with line-bundle cokernel.

The section to be represented is $R$, multiplied by the homogeneous
coordinate vanishing at infinity to the power $a-1$. On the dense
open $\operatorname{Res}(f,g)\ne0$, Euclidean inversion of $g$ in
$\C[x]/(f)$ gives its unique representation $R=wG-qF$ with the
bounds in \eqref{eq:division-recovery-v162}. Its class in the
line-bundle cokernel therefore vanishes on a dense open of the
integral reduced Hilbert graph, and hence vanishes identically.
The split injection recovers regular coefficients $w,q$.
The polynomial equality $r=gw-fq$ follows on the dense open and
therefore everywhere. This defines the inverse of
\eqref{eq:division-map-v162}. Conversely the reconstructed curves
agree with the universal Hilbert curve on a dense open; the two
maps to the separated Hilbert scheme agree on the entire reduced
neighbourhood. This proves the open-chart assertion, including
scheme structure and normality.

For the fibre calculation, factor \eqref{eq:division-Rees-v162}
as $d(g_1U-f_1V)$. The two factors are relatively prime on each
affine chart: $f_1,g_1$ are coprime. Their principal ideals
intersect in their product. The second factor is the graph of the
coprime pair, and the first is $\Pj^1$ over $\C[x]/(d)$.
Since $f_1,g_1$ generate the unit ideal modulo $d$, their intersection
is exactly the displayed section over that Artin ring. Each chart
is a hypersurface in a regular surface, hence Cohen--Macaulay;
its associated primes are precisely its component primes. This
also proves absence of embedded points. The Chinese remainder
theorem separates the distinct roots of $d$ with their full
multiplicities. Finally, $\chi(\OO_{T_d}(l))=r_0(l+1)$ and the
attaching scheme has Euler characteristic $r_0$. Adding the main
component gives $(a+1)l+1$ for $\OO(1,1)$, as required.
\end{proof}

\subsection{From the universal division chart to actual pencils}
For a corank-two incidence $p$ of length $a$, let $D_p$ be its
scheme-theoretic intersection with the pencil. The strict transform
of the pencil meets the exceptional projective plane bundle over
$D_p$ in a section $s_p$. A \emph{primitive line tail} means a line
subbundle in that plane bundle containing $s_p$, over all of $D_p$,
not merely a line on its reduced support. Its union with the strict
transform is taken with their common attaching subscheme $D_p$.
This definition concerns a fixed embedded curve in complete quadrics.

\begin{theorem}[Ambient jet-tail charts and their collision law]
\label{thm:ambient-jet-tails-v162}
Let $R_0$ be a regular pencil with no corank at least three. Suppose
its corank-two points have Smith exponents $(a_i,b_i)$, $a_i\leq b_i$;
points of corank one are unrestricted. In the full normalized pencil
Hilbert graph, the locus over $R_0$ with a primitive line tail at
each corank-two point is open and is
\begin{equation}\label{eq:jet-tail-fibre-v162}
 \prod_i\operatorname{Res}_{\C[x_i]/(x_i^{a_i})/\C}\Pj^1
                  =\prod_i J_{a_i-1}(\Pj^1).
\end{equation}
It is smooth, irreducible, and has dimension $\sum_i a_i$.
For $a_i=1$ it recovers the reduced-incidence fibre factor $\Pj^1$;
for $a_i>1$ constant directions form only its zero-jet section.
All points of \eqref{eq:jet-tail-fibre-v162} are limits of genuine
pencils. No extra framing or congruence quotient is taken in this
scheme fibre.

More precisely, after a smooth frame atlas, the completed local map
at any such lift is the product of the division maps
\eqref{eq:division-map-v162}, the $k_i$ coordinates in
\eqref{eq:Euclidean-normal-v162}, and complementary smooth
coordinates. Thus, in each chart, its nonincident open is
\begin{equation}\label{eq:resultant-boundary-v162}
                  \prod_i\operatorname{Res}_{x_i}(f_i,g_i)\ne0.
\end{equation}
The modification is smooth in these charts, but the resultant
boundary need not have normal crossings. Its relative curve has
local equation $g_iU_i-f_iV_i=0$. Formula~\eqref{eq:gcd-fibre-v162}
classifies all fibres in the chart, including mixed nonreduced
contacts and collisions, with multiplicities and attaching lengths.
On the actual complete-quadric curve the exterior multidegrees are
\begin{equation}\label{eq:jet-degrees-v162}
 \deg_C H_q=q-\Bigl(\sum_i\deg\gcd(f_i,g_i)\Bigr)
                   \mathbf1_{q=n-1},\qquad
 \deg_{(T_i)_{\rm red}}H_q=\mathbf1_{q=n-1}.
\end{equation}
Each local thick factor has its corresponding cycle multiplicity.
The polarization is $\bigotimes_{q=1}^{n-1}H_q$ and the Hilbert
polynomial is $\binom n2l+1$.
\end{theorem}
\begin{proof}
We first specify the formal comparison which makes this an ambient
statement. Work on the smooth atlas of framed pencils and choose
a parameter disc at each point in question. Theorem
\ref{thm:ambient-versal-v162} supplies independent formal coordinates
$(f_i,g_i,r_i,k_i)$ and smooth complementary coordinates. On such
a disc a unit $(n-2)$-block reduces the submaximal-minor ideal to
$(f_i,g_i,r_i)$. The rank locus is smooth of codimension three, and
its complete-quadric modification is its blow-up. Equivariant
Gaussian elimination identifies its graph with the three-coordinate
normal graph; the extra unit-block and target coordinates are fixed
functions of the pencil and of the disc parameter.

This comparison is also a comparison of the embedded Hilbert
problems, for the following reason. The parameter $x_i$ is retained:
we work in the incidence target over the universal pencil line,
not in an abstract unparametrized normal plane. The first
complete-quadric projection recovers that line coordinate, so adding
it does not change any embedded Hilbert point in the pencil graph.
A source congruence depending on $x_i$ is an automorphism of this
incidence target over the disc. It need not, and is not claimed to,
be one constant automorphism of the fixed projective target.
Away from the finitely many contact supports the submaximal ideal
is the unit ideal; the curve there is forced to be the original
lift, with no additional embedded choices. On the punctured discs
the two descriptions are this same forced lift. The local ideals
therefore glue uniquely. For Artin base extensions this can be
checked in affine charts: equality and flatness of their coherent
ideal quotients are detected after the faithfully flat completion
at each remaining support. Thus the completed embedded deformation
functors on the graph closure agree with the polynomial normal-graph
functors and the complementary smooth factors. This uses the
actual Rees image, not a general assertion that blow-up commutes
with nonflat base change. The same argument applies to families
specializing the dense nonincident graph, so their reduced closures
and their completed local rings agree. Excellence over $\C$ and
normality of the smooth charts identify the normalization there.

Apply Theorem~\ref{thm:division-chart-v162} to each factor. In the
central fibre, $f_i=x_i^{a_i}$ and $g_i=r_i=0$, and $w_i$ is an
arbitrary polynomial of degree less than $a_i$. The opposite chart
interchanges $g_i$ and $r_i$. On the overlap its slope is
$w_i'=w_i^{-1}$ in $\C[x_i]/(x_i^{a_i})$. These are exactly the
two standard charts for lines through the section $s_p$ over that
local Artin algebra. They represent
$\operatorname{Res}_{D_p/\C}\Pj^1$, equivalently the indicated jet
scheme.

There is an algebraic morphism from this jet space to the actual
Hilbert fibre, not only a formal map. Pull back the universal line
subbundle over $D_p$, take its projectivization in the exceptional
plane bundle, and unite it with the fixed strict transform along
its attaching section $D_p$. The union exact sequence has flat
outer terms over the jet space and a surjective restriction to the
attaching scheme, so the union is flat. Its line subbundle is
recovered from the embedded tail, also after base change; hence this
morphism is a monomorphism onto its image. The completed local
calculation above and the inverse in
\eqref{eq:division-recovery-v162} identify its completed local rings
with those of the Hilbert fibre. A finite-type morphism with these
completed local isomorphisms and equal residue fields is etale there;
an etale monomorphism is an open immersion. This proves the claimed
openness algebraically. The charts of the total Hilbert graph are
regular, so its normalization changes neither their rings nor this
open fibre. The construction descends from the frame atlas and
preserves distinctions among tails in the fixed target, rather than
quotienting by automorphisms of $R_0$.

Every such point specializes from the dense open in its smooth
chart, hence occurs from a formal one-parameter family of pencils;
properness of the Hilbert graph gives the same embedded limit.
Alternatively choose $g_i=t$ and lift its prescribed slope by
$r_i=\operatorname{rem}_{f_i}(g_iw_i)$, then use the formally smooth
pencil-coordinate map. Its local ideal is $(x_i^{a_i},t)$.
The two jet charts are affine spaces of dimension $a_i$, with the
usual open overlap. They are smooth and irreducible. Products at
distinct supports follow from the independent contact coordinates.

In a chart $r_i$ is a multiple of $g_i$ modulo $f_i$, so a common
zero of the triple is exactly a common zero of $f_i,g_i$. This gives
\eqref{eq:resultant-boundary-v162}. The fibre and multiplicity
statements follow from the scheme calculation in
\eqref{eq:gcd-fibre-v162}. Only the last exterior system has a base
ideal on the corank-two open. Its base length is the degree of the
indicated gcd; the smaller systems are units. This proves
\eqref{eq:jet-degrees-v162} and the Hilbert polynomial, with no
assumption that the gcd is squarefree.
\end{proof}

The theorem determines a full neighbourhood of every primitive
line-tail lift. It does not identify this open with the whole
proper fibre. In fact the first nonreduced case already has another
kind of component, as the next result shows.

\section{A non-equidimensional ambient Hilbert fibre}
\label{sec:nonequidimensional-v162}
\begin{theorem}[Two different ambient boundary components]
\label{thm:nonequidimensional-v162}
Let $n\geq4$ and let $R_0$ be a regular symmetric pencil with one
corank-two point having Smith exponents $(2,2)$ and otherwise only
corank-one points. Its normalized Hilbert fibre has an irreducible
component of dimension two, with open part $J_1(\Pj^1)$. It also
has a distinct component of dimension at least four. In particular
the fibre is reducible and not equidimensional.

The second assertion is witnessed by every smooth conic through
the attaching point in the exceptional $\Pj^2$: union with the
fixed main strict transform gives a four-dimensional locally closed
family of embedded Hilbert limits. These are reduced conic tails,
not the thick line tails parametrized by the jet open. The assertion
is about the actual fibre over $R_0$, not the fibre of a chosen
parameter slice. It does not assert that these exhaust its components.
\end{theorem}
\begin{proof}
Theorem~\ref{thm:ambient-jet-tails-v162} gives the smooth irreducible
open $J_1(\Pj^1)$ of dimension two in the normalized fibre. Its
closure is an irreducible component of dimension two.
For the other family use the full normal coordinates with $a=b=2$,
$k=1$, and take
\begin{equation}\label{eq:conic-arcs-v162}
 \begin{split}
 f&=x^2+t\alpha_1x+t^2\alpha_0,\\
 g&=t\beta_1x+t^2\beta_0,\qquad
 r=t\gamma_1x+t^2\gamma_0,
 \end{split}\qquad
       \beta_1\gamma_0-\beta_0\gamma_1\ne0.
\end{equation}
The three coefficient quadrics span all quadratic forms in $x,t$.
Their ideal is exactly $(x,t)^2$. The simultaneous graph on the
parameter surface is therefore $\operatorname{Bl}_{(x,t)}$, and
its exceptional line maps as
\begin{equation}\label{eq:conic-tail-v162}
 [X:T]\longmapsto
 [X^2+\alpha_1XT+\alpha_0T^2:
             \beta_1XT+\beta_0T^2:
             \gamma_1XT+\gamma_0T^2].
\end{equation}
This is a complete Veronese embedding, giving a smooth conic
through $[1:0:0]$. The graph is a closed embedding in the incidence
target, since these three generators define its Rees graph. Its
flat special curve is the fixed main component and this reduced
conic, meeting at the point corresponding to $[X:T]=[1:0]$.
The intersection has length one. Their last-exterior degrees are
$n-3$ and two, respectively, and their Euler characteristic is one.
The other exterior degrees are unchanged. Thus these are limits
with the required Hilbert polynomial.

Conversely every smooth conic through $[1:0:0]$ has a Veronese
parametrization of the displayed form: choose its preimage of that
point as $[1:0]$, and normalize the first leading coefficient to
one. Nonsingularity gives independence of the three quadrics.
The normal-coordinate realization theorem lifts these arcs to
actual pencil arcs. The congruence on the parameter disc specializes
on the exceptional divisor to its fixed value at the contact point;
it consequently gives the fixed exceptional plane identification,
not a quotient identifying distinct conics.

Smooth conics through a fixed point form an open of a hyperplane
in $\Pj^5$, of dimension four. Their universal conic is flat and
its attachment to the fixed main component is one reduced point;
the union construction is therefore a flat family. It is a locally
closed immersion into the Hilbert scheme, as the conic is recovered
as the residual component to the fixed main curve. Equivalently its
homogeneous quadratic ideal in the exceptional plane determines it
uniquely, in families. Since all its closed points are limits from
\eqref{eq:conic-arcs-v162}, this reduced family factors through the
reduced Hilbert graph fibre. The normalization of that graph is
finite and surjective, and its inverse image over this family is
finite and surjective. Some component of this inverse image has
dimension four. Hence the normalized fibre has a component of
dimension at least four. We use this finite inverse image, not a
purported automatic lifting of a nondominant normal family to a
normalization. This component is distinct from the dimension-two
component already obtained. The asserted non-equidimensionality
follows.
\end{proof}

The distinction has a concrete source. For a thick line the length-two
parameter algebra acts nontrivially on the tail. For a conic tail the
parameter acts by its closed-point value, and degree two is carried
by a reduced curve in the exceptional plane. Smith data of the closed
pencil do not choose between these Hilbert objects. The full versal
neighbourhood contains both, whereas the two constant transverse
parameters of the original Jordan slice see only a section of the
jet-tail locus.

\section{Incidence descent and the local graph calculations}
\label{sec:technical-completions-v162}
We record the scheme and deformation details used by the reduced and
nonreduced boundary theorems. These also fix the relation between the
local coordinates and their invariant meanings.

\begin{proposition}[The reduced open and its tangent sequence]
\label{prop:reduced-details-v162}
In the regular-pencil locus remove lines meeting rank at most $n-3$.
The universal corank-two incidence $D\to U$ is finite. Its open locus
of reduced geometric fibres includes empty fibres and is precisely
\[
 U_{\rm red}=U\setminus q(\operatorname{Supp}\Omega_{D/U}).
\]
At $(p,R)$ in this open, with $Z=Z_{n-2}\subset P=\Pj(M)$,
there is an exact tangent-normal sequence
\begin{equation}\label{eq:tangent-normal-v162}
 0\longrightarrow T_R\Sigma_p\longrightarrow T_RG
 \xrightarrow{\ e_p\ }N_{Z/P,p}/\overline{T_pL_R}
 \longrightarrow0.
\end{equation}
For distinct incidence points the direct sum of the maps $e_p$ is
surjective, without a semisimplicity assumption. Each $\Sigma_p$ is
smooth of codimension two. The local labelling of these centres is
etale; permutations on overlaps give the unlabelled boundary and
its Fitting ideal.
\end{proposition}
\begin{proof}
The determinant is a nonzero polynomial on a regular pencil, so its
corank-two incidence is finite on every fibre. The incidence projection
is proper, and is therefore finite by the proper quasi-finite theorem
\cite[Tag~02LS]{StacksProject160}. The rank-$(n-2)$ locus on this open
is smooth of codimension three, by its three Schur-complement entries.
A finite complex algebra has zero Kahler differentials if and only if
it is reduced: after decomposition into local Artin factors, a
nontrivial maximal ideal has nonzero cotangent space by Nakayama;
a residue field $\C$ contributes zero. The same holds geometrically
after extension of the residue field in characteristic zero. Formation
of relative differentials commutes with that extension. Their support
has closed finite image. Its complement is the asserted open, and
the zero algebra, representing an empty fibre, is allowed. Accordingly
$\Fitt_0(0)=\OO$.

Let $\ell\subset R\subset M$ represent $p$. Then
$T_RG=\Hom(R,M/R)$ and
$T_pP/T_pL_R=\Hom(\ell,M/R)$. Restriction from $R$ to $\ell$
is surjective. Compose it with the quotient to
$T_pP/(T_pZ+T_pL_R)$; this is exactly $e_p$ in
\eqref{eq:tangent-normal-v162}. Moving the incidence point along
the line and inside $Z$ shows that its kernel is $T_R\Sigma_p$.
For the convention $A:V^*\to V$, put $K=\ker A\subset V^*$.
Restriction of the represented bilinear form to $K$ identifies the
normal space with $\Hom(\ell,\Sym^2K^*)=
\ell^*\otimes\Sym^2K^*$. Thus the projective factor $\ell^*$ occurs on both sides
of the restriction map; it is not silently discarded. A frame of
$\ell$ identifies the map with restriction of the symmetric variation
to the radical. Changing that frame multiplies both maps by a unit.
A reduced intersection makes the normal image of $T_pL_R$ a line,
so the final quotient has dimension two.

For several points choose a nonsingular member and its self-adjoint
operator. Radicals at distinct eigenvalues have direct sum, even
when Jordan blocks are present. A symmetric form can prescribe its
restrictions on this direct sum independently. Choose a pencil
coefficient combination nonzero at every one of the finitely many
points; varying that coefficient realizes all restrictions. Variations
in the pencil itself land in their respective tangent-normal lines,
so the joint map descends to $\Hom(R,M/R)$ and is still surjective.
Smooth evaluation of the universal line to $P$ also gives the
incidence dimension $\dim G+1-3=\dim G-2$.

Finally strict henselization of the base splits the finite algebra
into its finitely many local idempotent factors. Every idempotent
and its finitely many equations descend to an etale neighbourhood,
since the strict henselization is a filtered limit of such
neighbourhoods. Each factor has one-dimensional closed fibre, so its
unit map has zero cokernel after tensoring with that residue field.
Nakayama, followed by shrinking, makes the unit map surjective.
It is therefore a closed immersion, not just a labelled point set.
The tangent sequence identifies its two local equations, and joint
surjectivity makes all equation pairs independent. Labels on etale
overlaps are permuted. This proves the claimed descent description.
\end{proof}

\begin{lemma}[Independent-centre Proj and normalization]
\label{lem:Proj-details-v162}
Suppose etale locally a smooth base has independent parameter pairs
$I_i=(u_i,v_i)$, with any number of complementary parameters. Then
\[
 \operatorname{Bl}_{\prod_i I_i}U
     =\prod_U\operatorname{Bl}_{I_i}U.
\]
The multi-Rees algebra has the relations
$v_iU_i-u_iV_i$, and its diagonal subalgebra is standard graded and
is the ordinary Rees algebra of $\prod_i I_i$. Its ordinary Proj
is the displayed scheme; the full multigraded algebra uses MultiProj.
\end{lemma}
\begin{proof}
In polynomial coordinates tensor the one-pair Rees injections over
$\C$; all variables are independent. The same relations persist
under etale flat extension. On a chart choose for each $i$ either
$u_i$ or $v_i$ as denominator. Localizing by the product of the
chosen generators and taking diagonal degree zero adjoins exactly
the individual ratios, with equations $v_i=u_iw_i$ or their opposites.
These are the fibre-product blow-up charts. The diagonal degree-one
space consists of all products of one generator from every pair;
its $d$th products give precisely $(\prod I_i)^d$. This proves the
standard-grading and ordinary-Proj assertions. The equality of
intersection and product of the ideals is checked on disjoint
monomial variable blocks and descends by faithful flatness.

In its Hilbert application the reduced graph closure is integral
and has the same function field as the pencil base, since both
contain the same dense nonincident open. Over a fixed pencil distinct
choices of exceptional lines produce distinct embedded subschemes:
restriction to each fixed exceptional plane recovers its chosen
line. Hilbert points do not identify them by an ambient congruence.
Thus the proper map from the smooth product blow-up has finite
geometric fibres, is quasi-finite, and is finite by
\cite[Tag~02LS]{StacksProject160}. Only after that finiteness step
does normality identify the source with the normalization of the
integral graph in their common function field.
\end{proof}

\begin{lemma}[Jordan constants and normal thick charts]
\label{lem:Jordan-details-v162}
For $L_m=E_m(xI-N_m)$ the determinant of its complementary
$(m-1)$-block and that of $E_m$ are both
$(-1)^{m(m-1)/2}$. Thus its scalar Schur complement is exactly
$x^m$. In the two-equal-block family
\eqref{eq:jordan-slice-v161}, the determinant constant $c_m$ is one.
The homogenized pencil at infinity is $E_m\oplus E_m$ and is
nonsingular. Its thick graph charts are flat and normal; their
special fibres have only their displayed component primes.
\end{lemma}
\begin{proof}
The complementary block is anti-triangular of size $m-1$ with
antidiagonal $-1$. Its sign is
$(-1)^{m-1+(m-1)(m-2)/2}=(-1)^{m(m-1)/2}$. The determinant
ratio proves the claim, and the two signs multiply to one.
Homogenization gives $E_m\oplus E_m$ when the finite-coordinate
homogenizing variable is zero.

In the chart $x^m=tz$, the coordinate ring is free of rank $m$
over the base polynomial ring in $z$, with basis
$1,x,\ldots,x^{m-1}$. In the other chart $t=x^my$, the
hypersurface equation is a nonzerodivisor in every geometric fibre
of its smooth ambient relative affine plane. The fibrewise Cartier
criterion \cite[Tag~062Y, part~(3)]{Stacks162} proves its flatness. The chart
$t=x^my$ is smooth after eliminating $t$. The chart $x^m=tz$ is a
domain and a hypersurface, hence Cohen--Macaulay. Its singular locus
$x=t=z=0$ has codimension two in the total space. It satisfies
$S_2$ and $R_1$, hence is normal by Serre's criterion. At $t=0$
the $t=x^my$ chart has the irredundant primary decomposition
\[
                 (x^my)=(x^m)\cap(y)
                     \quad\hbox{in }\C[x,y].
\]
Its associated primes are $(x)$ and $(y)$. The other chart
$(x^m)$ has the sole associated prime $(x)$. Polynomial smooth
parameters do not change these conclusions. This gives the exact
absence of embedded points, not just the support of the fibre.
\end{proof}

\begin{proposition}[Power coefficient modules on Artin tails]
\label{prop:Artin-powers-v162}
Let $A=\C[x]/(d)$ for a monic polynomial of degree $r$, and let
$T=\Pj^1_A$ be a primitive line tail in the corank-two normal plane.
For $1\leq j\leq hn$, put
\[
 \delta_j=(j-h(n-2))_+-2(j-h(n-1))_+.
\]
After removal of the common Cartier factors, the actual coefficient
map of the $j$th power surjects onto
$H^0(T,\OO_T(\delta_j))=\Sym^{\delta_j}(A^2)^*$.
In particular it is the full relative system, including nilpotent
coefficients; its cycle degree is $r\delta_j$.
\end{proposition}
\begin{proof}
The unit $(n-2)$-block supplies a socle factor with a unit coordinate,
so its multiplication inclusion on the residual coefficient module
is split over $A$. Put $b=j-h(n-2)>0$ when the normal power is
nonconstant. On the residual symmetric two-space its degree-$b$
coefficient span is, by the universal identity,
\[
 A[\alpha,\beta,\gamma]_b\quad(b\leq h),\qquad
 (\alpha\gamma-\beta^2)^{b-h}
                 A[\alpha,\beta,\gamma]_{2h-b}\quad(b\geq h).
\]
Remove the displayed determinant factor. A primitive line is a
split rank-two submodule of the three-coordinate module over the
semilocal Artin algebra $A$. Restriction of all degree-$\delta_j$
ternary forms therefore surjects onto all binary forms of that
degree. Composition with the split unit-block map is the required
explicit coefficient-module surjection. For $j\leq h(n-2)$ the
normal system is constant, and for $j=hn$ the same is true after
removing the full determinant factor. Invertible Gaussian changes
act invertibly on these modules. The original $\OO(j)$ twist on
the pencil line is trivial over the affine Artin support; this
trivialization is separate from the graph line bundle $\OO_T(\delta_j)$.
Thus no reduced-point argument or ignored twist occurs.
\end{proof}

\begin{proposition}[The transverse surface and ramified normalization]
\label{prop:transverse-details-v162}
For the intersection calculation in Proposition
\ref{prop:contact-deformations-v161}, take the specified normal surface
$S_m=\operatorname{Bl}_{(x^m,t)}\A^2_{x,t}$ and its proper exceptional
curve $F$. Then $mF$ is Cartier, $F^2=-1/m$, and the minimal
resolution chain starts at $F$ with multiplicities
$m-1,m-2,\ldots,1$. The main component meets $F$ at the other,
nonsingular end point of its parameter line; it does not meet the
far end of that exceptional chain.

After $t=s^m$, the normalization is
$\operatorname{Bl}_{(x,s)}\A^2$. Its reduced tail maps to the original
residual line by $[X:S]\mapsto[X^m:S^m]$. With its attachment marked
at $[1:0]$, its automorphism group over that fixed map is $\mu_m$;
it is a stable nonconstant component. No other marked points are
required for this assertion.
\end{proposition}
\begin{proof}
On $S_m$ the exceptional Cartier ideal is $\OO(-mF)$ and restricts
to $\OO_F(1)$ on the proper curve $F\simeq\Pj^1$. Intersection
with that proper curve is well-defined even though the surface is
not proper. Hence $(mF)\cdot F=-1$. In the singular chart
$x^m=tz$ the quotient parametrization is
$t=A^m,z=B^m,x=AB$. The intermediate primitive rays of the
$A_{m-1}$ resolution have orders of $t$ equal to
$m-1,\ldots,1$ in order away from the ray $A=0$, which is $F$.
The main component meets $F$ at $z=\infty$ in the other smooth chart.
The far end of the resolution chain faces the strict transform of
$z=0$, which is horizontal for the map to the $t$-line, not the
main special-fibre component.

After ramification the two rings are
$\C[x,s,y]/(s^m-x^my)$ and
$\C[x,s,z]/(x^m-s^mz)$. Their normalizations are obtained by
adjoining respectively $s/x$ and $x/s$. These elements are integral,
with $m$th powers $y$ and $z$, and belong to the common function
field $\C(x,s)$. The resulting rings $\C[x,s/x]$ and
$\C[s,x/s]$ are regular and finite over the displayed rings,
and glue to the ordinary point blow-up. This proves the normalization
claim with its common function field. The graph coordinates are the
$m$th powers of the exceptional coordinates, as stated. Its deck
transformations multiply $X/S$ by an $m$th root of unity; they fix
the attachment and form a finite group. The local cotangent module
$T^1$ computed previously refers only to abstract total-surface
hypersurface deformations. Neither this normalization nor that
calculation identifies embedded, fixed-base relative, or raw
Artin-algebra deformations with those abstract deformations.
\end{proof}

\section{Algebraic envelope forms and invariant relative deformations}
\label{sec:invariant-deformations-v162}
\begin{proposition}[An algebraic target for the envelope]
\label{prop:algebraic-envelope-v162}
The envelope functor of Proposition~\ref{prop:envelope-stack-v161}
factors through an algebraic stack of finite type over $\C$: the
stack of fppf forms, over the base, of the fixed projective-space
bundle with its graded algebra and power diagram. It is the classifying
stack $BH_h$ of a finite-type affine group $H_h$. This statement is
about forms of the specified bundle; it is not a boundedness theorem
for all graded algebra bundles locally trivial only on their total
source spaces.
\end{proposition}
\begin{proof}
On the fixed $\Pj(V)$ the degree-$p$ summand is
$B_h(V)_p\otimes\OO(2p)$. Its degree-preserving bundle automorphisms
are constant matrices in $\operatorname{GL}(B_h(V)_p)$, since
$H^0(\OO)=\C$. Inside their finite product, preservation of the unit,
multiplication and indicated polarized maps is a finite set of
polynomial equations. It defines a closed affine group $K_h$ of
finite type. The proved $\operatorname{PGL}(V)$ linearization of
the entire bundle lifts every projective source automorphism.
Consequently its full automorphism group is
$H_h=K_h\rtimes\operatorname{PGL}(V)$, with the action specified by
that linearization. Its classifying stack is algebraic of finite
type. Descent identifies its objects with the stated fppf forms.
The envelope associated to a projective-source torsor is such a form,
so its functor lands here. No global $\OO(1)$ on a Brauer-nonsplit
source is used. The exponent-free choice remains $h=1$.
\end{proof}

\begin{proposition}[The invariant lifting class]
\label{prop:invariant-lifting-v162}
For the effective boundary map of Theorem
\ref{thm:boundary-stack-v161}, the chart classes in
\eqref{eq:relative-obstruction-v161} are the local representatives of
\[
           \operatorname{Ext}^1(b^*L_{B/U},J).
\]
Here a deformation of the base and a square-zero extension with
kernel $J$ have been fixed. On a chart $v=uw$, the relative cotangent
complex is $[\OO e\xrightarrow{-u}\OO\,dw]$ in degrees $-1,0$.
On the opposite chart $u=vz$, its frames change by
\begin{equation}\label{eq:cotangent-transition-v162}
           z=w^{-1},\qquad e'=-z e,\qquad dz=-z^2dw.
\end{equation}
These formulas identify the local obstruction cokernels and solution
torsors. For several independent incidences the complexes form a
direct sum, permuted by etale monodromy, and descend to the effective
quotient stack.
\end{proposition}
\begin{proof}
The one-equation relative presentation has relation $v-uw$; its
relative differential is $-u\,dw$. Its dual with coefficients in
$J$ has cokernel $J/uJ$, where the equation defect is
$v'-u'\widetilde w$. On the overlap the other relation is
$u-vz=-z(v-uw)$ and $dz=-z^2dw$. These give
\eqref{eq:cotangent-transition-v162} and commute with the two
differentials because $v=uw$. Choosing inverse lifts of $w,z$ makes
the new defect $-z$ times the old one; any other choice changes it
by the corresponding differential image. Tangent corrections change
by $\eta_z=-z^2\eta_w$. Thus the obstruction and torsor are
coordinate-independent. The multiple-centre equations have independent
variables, so their relative complexes are a direct sum. Changes of
labels and source frames act on this complex and its classes, giving
descent. The common Lie-algebra directions cancel in the relative
cotangent complex of the two quotient atlases. This is a relative
lifting obstruction for a prescribed base deformation, not an
absolute obstruction on the smooth effective stack.
\end{proof}

\begin{corollary}[The ambient contact boundary for effective failure families]
\label{cor:effective-contact-v162}
The versal normal-coordinate, jet-tail and non-equidimensional-fibre
statements transfer to the effectively rigidified sharp pencil-failure
image via Theorem~\ref{thm:pencil-stack-equivalence}. In particular,
a family with contact exponents $(2,2)$ has both the primitive thick
line-tail locus and reduced conic-tail limits described in
Theorem~\ref{thm:nonequidimensional-v162}. The closed algebra does
not choose one of these limits.
\end{corollary}
\begin{proof}
The equivalence on arbitrary complex bases identifies the framed
completed pencil atlas with the effective first-relation image.
Pull back the explicit normal-coordinate maps and the embedded
Hilbert families along that equivalence; their source torsors and
changes of frame agree. This identifies deformations, not merely
closed points, in the stipulated effective category. The multiplication
power $J_{n-1}$ on the intrinsic pencil gives the same exact contact
ideal $(f,g,r)$ and its finite presentation in
\eqref{eq:versal-incidence-v162}. The distinguished exponent is one;
higher powers give the same graph with additional coefficient systems.
The two types of limits have already been realized by genuine pencil
arcs, so their effective failure families realize them as well.
This argument explicitly uses the sharp inverse. It does not claim
a new operation internal to the closed finite algebra before source
reconstruction, or an equivalence with the deformation stack of all
raw Artin algebras of its length.
\end{proof}

\section{Fixed-target comparison of contact graphs}
\label{sec:embedded-comparison-v163}
This section makes the change of category in the contact construction
explicit. A formal congruence depending on the pencil parameter acts
on an incidence bundle; it is not an automorphism of the fixed target
whose Hilbert scheme is under consideration. We retain the incidence
projection and undo every such change before taking a Hilbert point.

\subsection{Objects, targets, and the gluing statement}
Work on a finite-type framed atlas $S$ of regular pencils avoiding
corank at least three, with central pencil $R_0$. A frame here means
a lift of the projective source, an ordered pencil basis, and a local
parameter at each of the finitely many central corank-two points.
These choices are not quotiented in an embedded Hilbert fibre. Let
$\mathcal L\subset\Pj(\Sym^2V)\times S$ be the universal pencil
line, $X=\mathrm{CQ}(V)$, and
\[
 \mathcal Y=\mathcal L\times_{\Pj(\Sym^2V)}X.
\]
The targets and their maps are
\begin{equation}\label{eq:incidence-diagram-v163}
 \begin{array}{ccccc}
 \mathcal Z&\lhook\joinrel\longrightarrow&\mathcal Y
             &\lhook\joinrel\longrightarrow&X\times S\\
 &&\downarrow&&\downarrow\\
 &&\mathcal L&\longrightarrow&S .
 \end{array}
\end{equation}
Here $\mathcal Z$ denotes an embedded flat family, and the last
vertical arrow is projection. The middle map is a closed immersion
because $\mathcal L$ is the universal family of actual lines.
An object in the graph retains its point of $S$, not just the curve
in $X$. On the complement of the central contact supports, the
complete-quadric lift is the forced section of this incidence target.

Write $\widehat S=\operatorname{Spf}\widehat\OO_{S,R_0}$.
For a local Artin complex algebra $T$ with a specified map to
$\widehat S$, define $\mathcal D_X(T)$ to be the set of embedded
$T$-flat closed subschemes of \eqref{eq:incidence-diagram-v163},
with the fixed central curve, the specified Hilbert polynomial, and
the forced lift off the contacts. Isomorphism means equality as
closed subschemes, with the identity on the retained base. Restrict
this functor, when needed, to the completed reduced closure of the
nonincident graph. We do not use the abstract deformation functor
of the curve or identify objects by an ambient automorphism.

\begin{lemma}[Patching closed ideals with their quotients]
\label{lem:ideal-patching-v163}
Let $T$ be a local Artin complex algebra, let $W$ be a Noetherian
ambient scheme over a neighbourhood of a point of $\mathcal L_T$,
and complete along that point's inverse image. Let $W^\circ$ be
its complement. Compatible coherent ideal quotients on $W^\circ$
and the completion, with their identification on the punctured
completion, glue uniquely to a coherent ideal quotient on $W$.
The assertion applies to quotients with torsion supported at the
contact and to their flat families over $T$. Flatness, equality of
ideals and the specified closed incidence equations can be checked
on this cover. The same statement holds for finitely many disjoint
contact supports and after any Artin base change.
\end{lemma}
\begin{proof}
Use the ring-theoretic completion on affine neighbourhoods and its
pullback on projective charts. The map from this completion is flat,
Noetherian, and is an isomorphism over the closed support. The
formal patching equivalence for quasi-coherent sheaves is therefore
\cite[Proposition~5.6(4), Example~5.8]{Bhatt163}; its hypotheses are
also covered by Theorem~1.4 and Lemma~5.12(2) there. Flatness of the
completion is essential: the unrestricted module assertion for a
nonflat completion is not the statement being invoked. In
particular no torsion-freeness at the contact is imposed.

Apply this symmetric monoidal equivalence to the quotient with its
unit and multiplication. The ambient structure map and the ideal
kernel are retained as part of the data. Surjectivity and the
ordinary degree-zero algebra property are detected on the jointly
surjective flat cover consisting of the completion and the open
complement. Thus the glued object is an ordinary ideal quotient.
Finite presentation descends on this faithfully flat cover; the
ambient rings are Noetherian. The same cover detects base flatness
by the local criterion for flatness, or by vanishing of the relevant
Tor modules. It detects all the additional incidence equations.
Patching on the projective affine cover then glues the coherent
quotients. Iterating over disjoint supports gives the last assertion.
This argument supplies descent data on the punctured completions;
faithfulness of completion alone would only give uniqueness, not
existence.
\end{proof}

At the $i$th contact split a nonsingular block of size $n-2$. The
remaining symmetric Schur complement has entries $a_i,b_i,c_i$.
On this rank open the only noninvertible minor ideal below the
determinant is
\[
                  J_i=(a_i,b_i,c_i)=I_{n-1}.
\]
The determinant factor is an invertible ideal on the integral total
space and does not change the simultaneous graph. A congruence
$M_i(x_i)$ and the Schur splitting replace this generating triple
by an invertible linear transformation over the completed base
ring. Consequently they give an isomorphism between the corresponding
projective graph targets over the formal parameter disc. Denote it
by $\Theta_i$; it preserves $x_i$ and is part of the comparison data.
In normal coordinates the triple is $(f_i,g_i,r_i)$, the fourth
polynomial $k_i$ having been removed by the invertible row operation
$c_i\mapsto c_i-k_i a_i$.

\begin{proposition}[Completed embedded comparison]
\label{prop:embedded-functors-v163}
Fix the framed atlas, formal Schur splittings and normalizing
congruences just described. Let $\mathcal D_{\rm nor}$ be the
functor of embedded normal-graph quotients on the formal contact
discs, together with their identification with the forced lift on
the punctured discs. Then there is an equivalence of functors
\begin{equation}\label{eq:functor-comparison-v163}
 \mathcal D_X\ \simeq\ \mathcal D_{\rm nor}
\end{equation}
obtained by restriction and $\Theta_i$, with inverse given by
$\Theta_i^{-1}$ and ideal patching. No congruence quotient is taken
on either side. Changing the chosen frames or congruences conjugates
these functors and their inverse maps; it does not identify distinct
fixed-target Hilbert points.

On the reduced graph component, the complete local rings are the
complete local rings of the corresponding polynomial graph charts,
with the smooth parameters of the contact atlas adjoined. In
particular this assertion holds both at the division charts of
Theorem~\ref{thm:division-chart-v162} and at every determinantal
conic chart of Lemma~\ref{lem:HB-family-v163}. The comparison is
compatible with restriction to fibres. Where these total charts
are regular, normalization of the total graph changes neither them
nor their scheme fibres.
\end{proposition}
\begin{proof}
The completed normal form of
Theorem~\ref{thm:ambient-versal-v162} is used over
$T[[x_i]]$, with the base maximal-ideal filtration for its inductive
construction. The parameter $x_i$ is fixed; its powers specify the
remainder degrees, not the topology in that induction. Since $T$
is Artin the base induction is finite. For the complete universal
base it is the inverse limit of these finite inductions, so no
interchange of two unrelated convergence assertions is required.

An embedded family restricts to a coherent quotient on each
completed contact neighbourhood and to the forced quotient on the
complement. Apply $\Theta_i$ on the first and the induced
identification on the punctured neighbourhood. Conversely undo
$\Theta_i$, impose the original incidence equations, and apply
Lemma~\ref{lem:ideal-patching-v163}. These two operations are inverse
on quotients and their structure maps, proving \eqref{eq:functor-comparison-v163}
for every $T$, including nonreduced bases. The use of the image Rees
algebra can be recorded explicitly: a homomorphism $R\to R'$ gives
\[
        \mathcal R_R(J)\otimes_RR'\twoheadrightarrow
                          \mathcal R_{R'}(JR').
\]
Proj of this surjection embeds the actual image graph into the
base-changed incidence target. We do not assert that this
surjection is always an isomorphism. The closed ideals compared
in the proposition are these actual embedded ideals, rather than
an arbitrary base change of a blow-up.

For precision about the graph component, work first on the reduced
finite-type Hilbert graph over the atlas. The closure ideal of its
nonincident open can be computed in each finite presentation by
saturating with the ideal of the excluded base locus. Noetherianity
makes the increasing sequence of ideal quotients stationary; flat
completion commutes with each finite ideal quotient. It therefore
commutes with this closure computation. The rings are excellent
complex algebras, and their reduced completions stay reduced.
The patching equivalence preserves the nonincident section and the
retained base, hence also this specified reduced closure. It is
not a statement that unrelated components of an abstract curve
Hilbert scheme have disappeared.

The contact map on the completed smooth framed base is formally
smooth, with the bounded coefficient parameters as independent
coordinates and the remaining directions as power-series variables.
The scalar $k_i$ is one of these independent smooth factors for
the graph problem: the row operation above removes it without
changing the embedded ideal when the operation is undone.
Different contact discs patch independently. The polynomial graph
charts give algebraic representatives for their contact quotients;
their recovery maps retain the base coefficients. Applying the
functor equivalence gives the asserted isomorphism of completed
local rings with the smooth factors adjoined.

This is not yet an algebraization by formal assertion. In each
application below the morphism on the special fibre is constructed
algebraically first: a primitive line over the finite contact
algebra, or a conic together with the fixed main curve, supplies
its explicit coherent ideal. The finitely many coefficients of
the central gauge on that contact algebra are algebraic. Thus
these are finite-presentation morphisms into the actual Hilbert
graph. The completed comparison proves their local isomorphism
on the specified fibre charts. The residue fields at closed
points are $\C$ and the local rings are Noetherian; the formal
criterion for etaleness, \cite[Lemma~41.11.3]{Stacks163}, applies.
The coefficient or ideal recovery gives a monomorphism. Such a
morphism is an open immersion on these neighbourhoods. This orders
the argument as algebraic construction, completed comparison, then
etale monomorphism, rather than treating formal normal coordinates
as an algebraic family without descent.

An alternative normalizing congruence acts through a transformation
of the same relative incidence target. Both directions of the
comparison compose with that transformation and its inverse. The
composition in the fixed target remains the identity. The same
identity is the cocycle on overlaps of the source-frame atlas, so
the comparison descends. Finally a regular total chart is normal;
the finite normalization of the total graph is its identity there.
Taking its fibre retains the original base equations, possibly
with nilpotents. No commutation of fibre and normalization is used.
\end{proof}

\subsection{Cohomological recovery along the primitive locus}
\begin{lemma}[The union sequence and the retained coefficient base]
\label{lem:recovery-details-v163}
Let $D=\Spec\C[x]/(x^a)$ and let $Z_w=C_0\cup T$, with
$T=\Pj^1_D$ the primitive tail and intersection the attaching
section $D$. For $L=\OO(a-1,1)|_{Z_w}$ there is an exact sequence
\begin{equation}\label{eq:union-cohomology-v163}
 0\to L\to L|_{C_0}\oplus L|_T\xrightarrow{\operatorname{res}_C-
 \operatorname{res}_T}L|_D\to0.
\end{equation}
The restriction from $T$ is onto, $H^1(Z_w,L)=0$, and
$h^0(Z_w,L)=2a$. For each such $w$ these assertions persist on a
neighbourhood in the integral graph $\Gamma_a$ over $B_a$. On that
neighbourhood the recovery map has a line-bundle cokernel, and
vanishing of the retained third coordinate in that cokernel is a
scheme-theoretic equality, not a fibrewise assertion.
\end{lemma}
\begin{proof}
The ideals of a scheme-theoretic union give
$0\to\OO_{Z_w}\to\OO_C\oplus\OO_T\to\OO_D\to0$; tensor by
$L$. A line bundle on the local Artin scheme $D$ is free, and the
evaluation $H^0(\Pj^1_D,\OO(1))=\OO_D^2\to\OO_D$ at its section
has a unit coefficient, hence is surjective. On the main component
$L=\OO(a-1)$ and on the tail it is $\OO(1)$ up to an invertible
Artin factor. The cohomology sequence gives $a+2a-a=2a$ and the
claimed vanishing, over the Artin ring rather than only its residue
field.

To specify homogenization, write $W$ for degree-$(a-1)$ homogenization
of $w$, and $Q_0$ for degree-$(a-2)$ homogenization of $q$; when
$a=1$ the latter space is zero. The map is
\[
 (W,Q_0)\longmapsto WG-sQ_0F
       \quad\hbox{in }H^0(Z,\OO(a-1,1)).
\]
The section to be recovered is $s^{a-1}R$. There are $a+(a-1)$
unknown coefficients. Restricting to the main component, then to
the tail, proves injectivity at every central $w$. Upper
semicontinuity and cohomology and base change give a pointwise open
neighbourhood on which the map is a split injection of ranks
$2a-1$ and $2a$. These neighbourhoods cover the whole central
locus; no single uniform affine shrink is asserted. Its cokernel
is a line bundle. The class of $s^{a-1}R$ is zero on the dense
nonincident graph. A section of a line bundle on an integral
reduced neighbourhood that vanishes densely is zero. The recovered
coefficients are thus regular and satisfy the division identity.
The projection to $B_a$ is indispensable: it already supplies
$f,g,r$. All comparisons retain it.
\end{proof}

\section{The complete first nonreduced contact fibre}
\label{sec:complete-contact-v163}
We determine the entire fibre, including its nilpotents, at the first
nonreduced contact. The calculation also determines the common boundary
of the two components found in Theorem~\ref{thm:nonequidimensional-v162}.
It takes place in the graph over the coefficient base: the coefficients
are part of a point throughout.

Put $B_2=\mathbb A^6$, with coordinates the coefficients of a monic
quadratic $f$ and of two polynomials $g,r$ of degree at most one. Let
$\Gamma_2$ be the reduced closure in
$B_2\times\operatorname{Hilb}_{3l+1}(\Pj^1\times\Pj^2)$ of the graphs
of $[f:g:r]$, on their common nonvanishing locus. The polarization is
$\OO(1,1)$. Write $\nu:\widetilde\Gamma_2\to\Gamma_2$ for the
normalization of the total graph, and set
\[
 o=(x^2,0,0),\qquad X_2=(\widetilde\Gamma_2)_o.
\]
The scheme $X_2$ is not defined by normalizing a fibre. In the target
use coordinates $[F:G:R]$ and put $P=[1:0:0]$. The main curve at $o$
is $C_0=\Pj^1\times\{P\}$.

\subsection{An embedded chart containing all conic tails}
Let $[s:t]$ be coordinates on the first projective line, with $x=t/s$.
Consider the affine space with coordinates $(c,e,\lambda,A,B,C)$, and put
$z=t-cs$, $H=R-\lambda G$. In the product with $\Pj^1\times\Pj^2$
form the matrix
\begin{equation}\label{eq:HB-matrix-v163}
 \mathcal M=
 \begin{pmatrix}
 H&AF+CG\\
 -G&H+BF\\
 es&-z
 \end{pmatrix}.
\end{equation}
Its minors generate the bihomogeneous ideal
\begin{equation}\label{eq:HB-ideal-v163}
 \begin{split}
 Q&=H^2+BFH+AFG+CG^2,\\
 L_1&=zG-es(H+BF),\\
 L_2&=zH+es(AF+CG).
 \end{split}
\end{equation}
The associated map to $B_2$, written on $s=1$ with $z=x-c$, is
\begin{equation}\label{eq:HB-base-v163}
 \begin{split}
 f&=z^2+e^2C,\qquad g=eBz-e^2A,\\
 r&=\lambda g-eAz-e^2BC.
 \end{split}
\end{equation}

\begin{lemma}[The determinantal family and its recovery maps]
\label{lem:HB-family-v163}
The minors in \eqref{eq:HB-ideal-v163} define a flat family of embedded
curves of Hilbert polynomial $3l+1$. It defines a morphism to $\Gamma_2$.
Near every point with $e=0$, this morphism is an open immersion. Its
image is regular and is unchanged by normalization of the total graph.
At $e=0$ the curve is the scheme-theoretic union of $C_0$ and the plane
conic $Q=0$ over $x=c$, with the attachment placed at that base
point. The intersection is the reduced point $(c,P)$, even when $Q$
is a doubled line. Constant changes of target coordinates fixing $P$
provide such charts around every plane conic through $P$.
\end{lemma}
\begin{proof}
The minors have height two on every geometric fibre. When $e\ne0$,
one eliminates $H$ locally from $L_1,L_2$; equivalently the rank-one
condition is the two-generator graph of $f,g$, with the indicated third
coordinate. At $e=0$ the ideal is
\[
 (zG,zH,Q)=(G,H)\cap(z,Q),
\]
where $Q\in(G,H)$ is nonzero. The equality follows by reducing modulo
$z$ and then using that $z$ is regular on the polynomial ring in
$F,G,H$ modulo $Q$. This gives the asserted union and intersection.
In particular the height is two also for a singular conic. The
Hilbert--Burch complex of \eqref{eq:HB-matrix-v163} is consequently
exact on every fibre. Its sheaf form is
\begin{equation}\label{eq:HB-resolution-v163}
 0\longrightarrow\OO(-1,-2)^2\longrightarrow
 \OO(-1,-1)^2\oplus\OO(0,-2)\longrightarrow
 \OO\longrightarrow\OO_{\mathcal C}\longrightarrow0.
\end{equation}
The fibrewise exactness criterion for a complex of base-flat locally
free sheaves gives base flatness of the cokernel. Computing the Euler
polynomial from this complex gives $3l+1$. This is an application of
the classical determinantal resolution, not an assertion that the
resolution theorem is new.

Substitution of \eqref{eq:HB-base-v163} into the three minors gives
zero. On the dense open
\[
             e(A^2+CB^2)\ne0
\]
the three binary quadratics $f,g,r$ are a basis, the curve is their
graph, and the image is a smooth conic. Thus the flat family gives a
morphism into the reduced graph closure. The assertion at infinity
is direct: $s=0$ implies $z=t\ne0$, so $L_1,L_2$ give $G=H=0$.
There is no additional component or choice at infinity.

We prove the open-immersion assertion by regular recovery, not by
counting parameters. At a central curve $C_0\cup Q$, the union exact
sequence and its twists give
\[
 H^1(\OO(0,2)|_{C_0\cup Q})=0,\quad h^0=5,
 \qquad H^1(\OO(1,1)|_{C_0\cup Q})=0,\quad h^0=4.
\]
Indeed $H^0(Q,\OO_Q(2))$ has dimension five and
$H^0(Q,\OO_Q(1))$ dimension three for every nonzero plane conic,
including a doubled line, by
$0\to\OO_{\Pj^2}(j-2)\to\OO_{\Pj^2}(j)\to\OO_Q(j)\to0$.
The restriction to $P$ is surjective. Combining it with the main
component proves both equalities and the surjectivity of the ambient
restriction maps. Cohomology and base change therefore give, in a
neighbourhood of this point in the integral graph $\Gamma_2$, a line
of quadratic equations and a rank-two bundle of equations of bidegree
$(1,1)$.

Normalize the coefficient of $R^2$ in the quadratic equation to one.
Its $F^2$ coefficient is zero on the dense graph, because infinity
maps to $P$, and hence is zero throughout this neighbourhood.
Completing the square in $R,G$ uniquely writes the equation as $Q$
in \eqref{eq:HB-ideal-v163}; it recovers
$\lambda,A,B,C$ regularly. The projection to $B_2$ supplies $f$,
so $c$ is minus half its linear coefficient. Normalize the two
bilinear equations by their coefficients at $tG,tH$; the relevant
minor is a unit at the central curve. The negative coefficient of
$sH$ in the equation with $tG$ coefficient one recovers $e$.

Here is an explicit verification on the dense graph, which also
checks the normalizations. Write $f=(x-c)^2+d$,
$g=eB(x-c)+u$, and $r-\lambda g=-eA(x-c)+v$. The coefficient of
$x^3$ in $Q(f,g,r-\lambda g)=0$ gives the displayed opposite linear
coefficients. When $A^2+CB^2\ne0$, the remaining coefficients give
\[
 u=-e^2A,\qquad v=-e^2BC,\qquad d=e^2C.
\]
The normalized bilinear equations are then precisely $L_1,L_2$.
All these are equalities of regular functions or sections on an
integral neighbourhood, so they extend from its dense graph open.
The recovered parameters satisfy \eqref{eq:HB-base-v163}, and the
universal curve equals \eqref{eq:HB-ideal-v163}, first densely and
then by separatedness of the Hilbert scheme. Recovery is also the
identity on the determinantal parameter space. These two morphisms
are inverse after restricting to their common open neighbourhoods.
The parameter space is smooth; its image is therefore regular and
normalization does not change it.

Finally any nonzero quadratic form vanishing at $P$ is nonzero on
some vector not proportional to $P$. Choose that vector as the $R$
direction. This is a constant target change fixing $P$, followed by
rescaling. The preceding argument applies. Such changes are changes
of chart, not an equivalence relation on Hilbert points.
\end{proof}

\subsection{The full fibre and its nilpotents}
Let $\mathbb F_2=\Pj(\OO_{\Pj^1}\oplus\OO_{\Pj^1}(2))$ be the
compactification of $\operatorname{Tot}(\OO(2))$ whose section at
infinity has self-intersection $-2$. Denote that section by $D_\infty$.
Let $\mathcal Q_P\simeq\Pj^4$ be the projective space of plane conics
through $P$, and let $D_{\rm dbl}\simeq\Pj^1\subset\mathcal Q_P$
be its locus of doubled lines through $P$.

\begin{theorem}[Complete first nonreduced fibre]
\label{thm:complete-fibre-v163}
The scheme $X_2$ has exactly two irreducible components, with reduced
structures $\mathcal Q_P\simeq\Pj^4$ and $\mathbb F_2$. Their
scheme-theoretic intersection, when each is given its reduced closed
structure, is the reduced projective line
\[
                  D_{\rm dbl}=D_\infty.
\]
The open complement of $D_\infty$ in the second component is
$J_1(\Pj^1)=\operatorname{Tot}(T\Pj^1)$, and the first component
parametrizes exactly the curves $C_0\cup Q$, $Q\in\mathcal Q_P$.
Thus the second component has dimension exactly two and the conic
component has dimension exactly four.

The fibre is connected and generically reduced on both components,
but is not reduced. In the charts of Lemma~\ref{lem:HB-family-v163}
its scheme structure is
\begin{equation}\label{eq:full-fibre-ring-v163}
 \C[\lambda,e,A,B,C]/(eA,eB,e^2C).
\end{equation}
At a doubled line with slope $\lambda_0$ the completed local ring is
obtained by replacing the polynomial ring by
$\C[[\lambda-\lambda_0,e,A,B,C]]$ in this formula. The nilradical
is generated by $eC$, has square zero, and is supported on the locus
of conics singular at $P$. The fibre is not Cohen--Macaulay at that
locus. Its reduced support is the union of the two indicated smooth
components, and the normalization of that reduced support is their
disjoint union. This last normalization is different from the
normalization of the total Hilbert graph used to define $X_2$.
\end{theorem}
\begin{proof}
First take the fibre of \eqref{eq:HB-base-v163} at $o$. Its coefficient
equations give $c=0$, $e^2C=0$, $eB=0$, $e^2A=0$ and
$e(\lambda B-A)=0$, together with a redundant constant equation.
The resulting ideal is exactly
\[
                    (c,eA,eB,e^2C).
\]
No radical has been taken. In particular the charts compute the
scheme fibre of a regular total graph and therefore of its
normalization as well.

The conic family is algebraic over all of $\mathcal Q_P$. Its ideal
is $I_{C_0}\cap I_Q$, with $Q$ placed at $x=0$; the intersection
with $C_0$ is the constant reduced point $P$. The union exact
sequence proves flatness and the Hilbert polynomial. The family
is a closed immersion into the Hilbert fibre: its inverse on its
image recovers the quadratic equation as the kernel of the ambient
restriction in degree $(0,2)$. This also gives the relative residual
construction; locally
\[
              (I_{C_0}\cap I_Q):I_{C_0}=I_Q.
\]
For the displayed ideals this follows by dividing by $G$ and $H$:
no nonzero element in the hypersurface ring of a conic is annihilated
by both $G$ and $H$. This remains true in the universal conic family
by its same presentation. Alternatively the kernel-line construction
alone proves the asserted relative closed immersion. The regular
charts of Lemma~\ref{lem:HB-family-v163} show that this family lifts
uniquely to the normalization near every one of its points; we do
not use a general lifting assertion for maps from normal schemes.

The primitive open of the other component is the established jet
chart. Write a jet as $w=w_0+w_1x$ modulo $x^2$. Near its direction
at infinity use
\[
             \lambda=w_0,\qquad e=1/w_1.
\]
In \eqref{eq:HB-ideal-v163} put $c=A=B=C=0$. This gives
\begin{equation}\label{eq:adjacency-family-v163}
         (xG-eH,\ xH,\ H^2),\qquad H=R-\lambda G.
\end{equation}
For $e\ne0$ it is the primitive thick-line curve with slope
$w=\lambda+x/e$. At $e=0$ it is
\[
               (xG,xH,H^2)=(G,H)\cap(x,H^2),
\]
the main curve together with the planar doubled line $H^2=0$ at
$x=0$. Under the other affine coordinate of the direction line,
$w_0'=1/w_0$ and $w_1'=-w_1/w_0^2$; hence
$e'=-w_0^2e$ on the boundary chart. These are exactly the transition
functions of the compactification of $T\Pj^1=\OO(2)$ described
above. Formula \eqref{eq:adjacency-family-v163} and the jet families
glue to a morphism from $\mathbb F_2$ into $X_2$. The local inverse
recovers $\lambda,e$; the original jet chart supplies the inverse
on its complement. Thus this is a closed immersion, since its
source is proper. Its infinity section is precisely $D_{\rm dbl}$.
It has self-intersection $-2$ by the transition functions. The sum
of the two reduced component ideals in \eqref{eq:full-fibre-ring-v163}
is $(e,A,B,C)$, giving the stated reduced intersection.

It remains to prove that these two families exhaust the fibre.
Their union $K$ is closed and proper. It is also open in the
underlying fibre: at a point of the conic family the local fibre
\eqref{eq:full-fibre-ring-v163} has support
\[
                    V(e)\ \cup\ V(A,B,C),
\]
which is exactly the union just constructed; at a jet point the
original open jet chart has no other branch. The scheme
$\widetilde\Gamma_2$ is integral and normal, and its projection
to the normal integral base $B_2$ is proper and birational.
Consequently all its geometric fibres are connected by Stein
factorization \cite[Lemma~37.53.6]{Stacks163}. The nonempty subset
$K$ is open and closed in $X_2$, so it is all of its support.
This is an exhaustion argument for the proper fibre, not an
inference from the existence of two families of curves.

For clarity, the local primary calculation is
\begin{equation}\label{eq:primary-fibre-v163}
 (eA,eB,e^2C)=(e)\cap(A,B,C)\cap(A,B,e^2).
\end{equation}
It is irredundant wherever all the displayed coordinate functions
vanish. It follows either by divisibility of monomials, or by
intersecting first $(e)$ with $(A,B,C)$ and then with $(A,B,e^2)$.
The two minimal primes are $(e)$ and $(A,B,C)$; $(e,A,B)$ is the
embedded associated prime when present. The radical is
$(eA,eB,eC)$, so the nilradical is generated by $eC$ and
\[
              \operatorname{Ann}(eC)=(e,A,B),\qquad (eC)^2=0.
\]
In conic coordinates the equations $A=B=0$ say exactly that $Q$
is singular at $P$. These calculations also apply after localization
and completion and identify the claimed support intrinsically.
They prove generic reducedness, nonreducedness and failure of the
Cohen--Macaulay property. The two reduced components are smooth and
are distinct; the final assertion about their normalization follows.
\end{proof}

\begin{corollary}[Explicit adjacency and a numerical consequence]
\label{cor:adjacency-v163}
Every doubled line through the attaching point is the common limit of
primitive thick-line jets and of smooth conic tails. At such a point
there is no extra component beyond the two in
Theorem~\ref{thm:complete-fibre-v163}. The Euler characteristic of the
underlying complex fibre is seven.
\end{corollary}
\begin{proof}
The first specialization is \eqref{eq:adjacency-family-v163}. For the
second, set $e=0$, $\lambda$ constant, $B=C=0$, $A=\tau$; then
$Q=H^2+\tau FG$ is nonsingular for $\tau\ne0$ and specializes to
the doubled line. These are algebraic one-parameter families.
Exhaustion was proved in the theorem. Euler characteristic is
additive for the closed union, giving $5+4-2=7$.
\end{proof}

\section{Products and adjacency for the first nonreduced boundary}
\label{sec:products-v163}
The local fibre calculation applies to arbitrary larger Smith exponents
and to any number of distinct contacts. This gives a complete boundary
class, rather than a primitive open in a class whose complement has
not been identified.

Let $R_0$ be a regular symmetric pencil whose line avoids corank at
least three. At its distinct corank-two points write the positive
Smith exponents as $(a_i,b_i)$, $a_i\leq b_i$. Assume
\begin{equation}\label{eq:first-boundary-class-v163}
                      a_i\in\{1,2\}\quad\hbox{for all }i.
\end{equation}
There is no bound on $b_i$ except the dimension constraint on the
pencil. Let $r$ count the points with $a_i=2$ and $s$ those with
$a_i=1$. Denote by $\widehat G$ the normalization of the total reduced
Hilbert graph of lifted pencil lines and by $\mathfrak X_{R_0}$ its
scheme fibre. The polarization on complete quadrics is
$\bigotimes_{q=1}^{n-1}H_q$, with Hilbert polynomial $\binom n2l+1$.

\begin{theorem}[Complete fibres and their component intersections]
\label{thm:product-fibres-v163}
After choosing parameters and residual frames at the distinct contact
points, there is an isomorphism of schemes
\begin{equation}\label{eq:product-fibres-v163}
              \mathfrak X_{R_0}\simeq X_2^{\,r}\times(\Pj^1)^s.
\end{equation}
Frame changes act by the corresponding transition isomorphisms; the
construction is intrinsic without an ordering of the supports.
In particular the reduced fibre has exactly $2^r$ irreducible
components. For a subset $S$ of the $r$ length-two contacts, the
component is
\begin{equation}\label{eq:component-product-v163}
 (\Pj^4)^{|S|}\times(\mathbb F_2)^{r-|S|}\times(\Pj^1)^s,
 \qquad \dim=s+2r+2|S|.
\end{equation}
The intersection of two reduced components is obtained by replacing
each factor where their choices differ by the doubled-line
$\Pj^1$, and leaving the common factors unchanged. These intersections
are reduced. This describes all multiple intersections as well.

The fibre is connected and is reduced exactly when $r=0$. For $r>0$
its nilradical has nilpotence order exactly $r+1$: its $(r+1)$st
power is zero, and its $r$th power is nonzero. Locally its equations
are tensor products of \eqref{eq:full-fibre-ring-v163} and smooth
jet charts. In particular the reduced support is not a replacement
for its Hilbert fibre scheme. The Euler characteristic of the
underlying complex space is $7^r2^s$.
\end{theorem}
\begin{proof}
We construct the morphism before using completed rings. At a
length-two support, the universal curve over $X_2$ contains its
fixed main curve. In the charts \eqref{eq:HB-ideal-v163}, the fibre
relations $eA=eB=e^2C=0$ give
\[
             xG=eH,\qquad xH=-eCG,
             \qquad x^2G=x^2H=0.
\]
Consequently the quotient of the fixed main ideal by the curve ideal
is annihilated by $x^2$. Thus the entire modification, including the
nilpotent directions in its parameter scheme, is specified by a
coherent ideal quotient on the first infinitesimal neighbourhood of
the support. At a length-one contact the corresponding quotient is
annihilated by $x$. These are finite algebraic neighbourhoods, not
unbounded formal data.

Use the central Schur splitting and undo its congruence on these
finite quotients. Only finitely many coefficients of the central
gauge are involved. Include the fixed main lift, and extend by its
forced ideal off the supports. This gives algebraic closed ideal
families in the fixed complete-quadric target. Multiplicative
stability of the ideals is preserved by the target coordinate
isomorphisms. Distinct supports are disjoint, so the constructions
combine independently. Flatness follows from the local flat
families and the patching lemma; their Hilbert polynomial is the
one stated above. Indeed each contact of length $a_i$ removes
$a_i$ from the last exterior degree of the main curve and puts
that degree on the tail. The attachment lengths give Euler
characteristic one for the total curve. The smaller exterior
systems have no base point on this rank open.

This constructs a finite-presentation morphism from the right side
of \eqref{eq:product-fibres-v163} to the fibre of the reduced Hilbert
graph. Each point is a specialization of nonincident graph points:
for primitive tails this is the division family, and for conic tails
it is the dense open $e(A^2+CB^2)\ne0$ in
\eqref{eq:HB-base-v163}. A polynomial line through a specified
point with generic direction meets that open; normalization of
its closure gives an algebraic one-parameter specialization.
The nilpotent part of the constructed map also factors through
the reduced total graph, by the completed comparison in
Proposition~\ref{prop:embedded-functors-v163}, not merely because
its closed points do. That proposition identifies its completed
total neighbourhoods with the regular polynomial graph charts
and the smooth contact factors. Hence the total graph is regular
there, the morphism lifts uniquely to $\widehat G$, and the
completed local maps on the fibre are isomorphisms.

Recovery of the embedded ideal on each disjoint finite neighbourhood
makes this morphism a monomorphism. The formal etaleness criterion
therefore makes it an open immersion. Its source is proper, so it
is also a closed immersion. The fibre of $\widehat G\to G$ is
connected: this is a proper birational morphism from an integral
normal variety to the smooth Grassmannian, and Stein factorization
applies. The nonempty open-and-closed image is the entire fibre.
This proves the scheme isomorphism, not just an equality of reduced
supports. It also explains why the arbitrary values of $b_i$ make
no difference: $k_i$ is removed by a generating-row transformation
which is undone in the fixed incidence target, while the contact
algebra and the resulting finite quotient have length $a_i$.

The reduced components and intersections follow from
Theorem~\ref{thm:complete-fibre-v163} by products. Tensor products
over $\C$ preserve the displayed reduced smooth components and
their reduced intersections. The nilradical is the sum of the
pullbacks $N_i$ of the square-zero nilradicals of the $r$ factors.
The quotient by this sum is reduced, so there is no additional
nilradical. Every monomial of degree $r+1$ repeats a factor and
vanishes. The tensor product of nonzero sections of all $N_i$ is
nonzero on a product of their nonzero affine charts, so the $r$th
power is nonzero. The Euler characteristic follows by multiplicativity
and Corollary~\ref{cor:adjacency-v163}.

Finally the construction uses the actual support schemes and the
fixed embedded ideals. Permuting the support labels permutes the
factors; changing local source frames applies their functorial
transition maps. These satisfy the cocycle identity because the
coordinate changes are undone in the fixed target. This gives the
asserted unlabelled interpretation, rather than a preferred global
ordering of contacts.
\end{proof}

\subsection{Resultant strata on the higher-contact open}
For larger $a_i$ the division chart remains useful even though the
whole proper fibre is not classified here. The following records
exactly its boundary stratification. It is also a check that repeated
roots are not being mistaken for multiple reduced components.

\begin{proposition}[Gcd partitions, codimensions and multiplicity]
\label{prop:resultant-strata-v163}
In $W_a=\{(f,g,w)\}$, with $f$ monic of degree $a$ and
$\deg g,\deg w<a$, fix a partition
$\mu=(\mu_1,\ldots,\mu_k)$ of $d\leq a$. The locally closed locus
where the monic gcd of $f,g$ is
$\prod_{j=1}^k(x-\xi_j)^{\mu_j}$ with distinct $\xi_j$ has dimension
\[
                         3a-2d+k
\]
and codimension $2d-k$. Its finite attaching scheme is the disjoint
union of the primary schemes
$\Spec\C[x]/((x-\xi_j)^{\mu_j})$, and its tails are projective
lines over those schemes. A multiple root is one primary thickening,
not several components. At a point with squarefree $f$ and precisely
$d$ common roots, the resultant divisor is locally a product of
$d$ independent parameters and has multiplicity $d$.
\end{proposition}
\begin{proof}
Parametrize the gcd by its distinct roots, giving $k$ parameters,
and write $f=hf_1$, $g=hg_1$. The monic polynomial $f_1$ has
$a-d$ coefficients and $g_1$ has $a-d$ coefficients; impose the
open coprimality condition on this residual pair and the conditions
specifying the exact gcd partition. Recovering the monic gcd is
regular on the constant-gcd-degree stratum by the fixed-rank
Euclidean linear equations. Up to the finite permutations of equal
parts, this gives dimension $k+2(a-d)$; $w$ contributes $a$.
The statement includes $d=a$, where $g=0$ and $f=h$.
The primary decomposition is the Chinese remainder theorem for
distinct roots and the gcd fibre equation of
Theorem~\ref{thm:division-chart-v162}. At a squarefree $f$, its
roots are etale local coordinates and
$\operatorname{Res}(f,g)=\prod_{f(\xi)=0}g(\xi)$.
The evaluations of a polynomial of degree $<a$ at these roots are
independent by the Vandermonde matrix. The nonzero factors are
units, proving the final assertion. No normal-crossing statement
is made at a multiple root of $f$.
\end{proof}

\subsection{A testable extension beyond length two}
The following conjecture concerns the reduced components of a local
fibre, not the miniversal congruence quotient. It distinguishes two
questions: component geometry and a possible nilpotent structure on
their union.

\begin{remark}[Conjecture: rational components of the corank-two contact fibre]
\label{conj:contact-components-v163}
For every $a\geq1$, let $X_a$ be the fibre over $(x^a,0,0)$ of the
normalization of the total reduced graph $\Gamma_a\to B_a$ defined
in Theorem~\ref{thm:division-chart-v162}. Every irreducible component
of $(X_a)_{\rm red}$ has rational normalization. The closure of the
primitive jet open is an irreducible component of dimension $a$,
and its intersections with the other components contain a locus of
curves whose tail cycle has a repeated line through the attaching
point. For $a\geq2$ the scheme $X_a$ is not reduced.
\end{remark}
For $a=1$ the fibre is $\Pj^1$. For $a=2$ the conjecture follows
from the complete description above, including the doubled-line
intersection and the square-zero ideal. For $a\geq3$ it is a
conjecture, not a component classification. The statement is
falsifiable separately by the birational type of a component, its
adjacency to the primitive component, or reducedness of the fibre.
An eventual classification should refine it by the degrees and
primary thicknesses of tail cycles of total degree $a$. The present
product theorem makes no assertion about higher-corank contacts or
singular pencils in the Hilbert problem; the independent inverse
and singular-pencil invariant theorems continue to apply on their
original domains.

\section{Tangent maps, linearizations, and effective transfer}
\label{sec:tangent-effective-v163}
We separate the source of the contact tangent map from the category in
which the boundary is embedded. The elementary linearizations below
also specify the splitting used in the envelope automorphism group.

\begin{lemma}[Framed tangent map and stabilizer]
\label{lem:framed-tangent-v163}
Let $A(x)$ be a regular symmetric pencil, expressed as a symmetric map
from $E$ to $E^*$ with its parameter-line twist trivialized locally.
After splitting a unit block, let $j(x):W[[x]]\to E[[x]]$ be the
inclusion of its orthogonal residual summand, and let $D(x)$ be the
residual diagonal form. On the atlas of ordered pencil coefficients
the contact differential is
\begin{equation}\label{eq:contact-differential-v163}
 (\dot A_0,\dot A_1)\longmapsto
 \left[j(x)^*(\dot A_0+x\dot A_1)j(x)\right]
 \quad\hbox{in }\frac{\Sym^2W^*[[x]]}
 {\{H^{\mathsf t}D+DH:H\in\operatorname{End}(W)[[x]]\}}.
\end{equation}
The restriction to distinct contacts gives the direct sum of these
maps. Local frame changes act by congruence and by the corresponding
line-bundle unit; they identify these modules and differentials.
This formula does not discard the pencil-basis directions.

If a nonsingular member is fixed and $T$ is its self-adjoint pencil
operator, the Lie algebra of its orthogonal stabilizer is
\[
        \{K:KT=TK,\ K^*=-K\}.
\]
For Jordan block sizes $d_{\lambda,i}$ its dimension is
$\sum_{\lambda,i<j}\min(d_{\lambda,i},d_{\lambda,j})$. Hence on
this fixed-member framed slice the contact quotient has dimension
\[
 n+\sum_{\lambda,i<j}\min(d_{\lambda,i},d_{\lambda,j}).
\]
This is not asserted to be the dimension of the unframed pencil
quotient by source congruence and parameter reparametrization.
\end{lemma}
\begin{proof}
For a block matrix with invertible upper block $U$ and off-diagonal
block $B$, the inclusion is $j=(-U^{-1}B,1)^{\mathsf t}$. Differentiating
the Schur complement gives $j^*\dot A j$: the terms differentiating
$j$ cancel because its image is orthogonal to the unit block.
Infinitesimal residual frame changes contribute exactly
$H^{\mathsf t}D+DH$, proving \eqref{eq:contact-differential-v163}.
The polarizing line of the quadratic form has not disappeared. If
$\ell$ is the tautological line at the pencil point, the untrivialized
value is in $\Sym^2W^*\otimes\ell^*$, with the same congruence
quotient; choosing a frame of $\ell$ gives the displayed formula.
Changing it multiplies both numerator and congruence image by the
same unit.

The ordered-basis atlas has tangent space $\Hom(\C^2,\Sym^2E^*)$.
At its image pencil $R$, the sequence forgetting the basis is
\[
 0\longrightarrow\operatorname{End}(R)
 \longrightarrow\Hom(R,\Sym^2E^*)
 \longrightarrow\Hom(R,\Sym^2E^*/R)\longrightarrow0.
\]
The rightmost term is the Grassmannian tangent space. A splitting
specifies a local basis section; a different splitting may change
the parameter-fixed contact coordinates. Thus the differential is
first formulated on the framed atlas. The comparison theorem
retains these choices and descends by its transition maps, instead
of silently treating the four basis directions as zero contact
vectors. Source-frame directions contribute the infinitesimal
congruence image. The block computation in
Theorem~\ref{thm:ambient-versal-v162} proves surjectivity on this
atlas, so it proves formal smoothness of its completed contact map.

For the stabilizer, decompose the self-adjoint operator into
orthogonal primary cyclic blocks. Intertwiners between distinct
eigenvalues vanish. Between blocks of the same eigenvalue their
space has dimension $\min(d_i,d_j)$, as homomorphisms of
$\C[t]$-modules $\C[t]/(t^{d_i})$ and $\C[t]/(t^{d_j})$.
An off-diagonal intertwiner $K_{ij}$ forces
$K_{ji}=-K_{ij}^*$ and supplies one independent copy. On a single
cyclic block a commuting endomorphism is a polynomial in $T$,
hence is self-adjoint. It is skew-adjoint only when it is zero.
This proves the dimension formula without assuming semisimplicity.
Subtracting the orbit dimension
$n(n-1)/2-\dim\operatorname{Stab}$ from the dimension
$n(n+1)/2$ of self-adjoint operators gives the last formula.
\end{proof}

\begin{proposition}[The specified envelope has a split automorphism extension]
\label{prop:linearized-envelope-v163}
Fix $V$, $h\geq1$, and the graded algebra bundle on $\Pj(V)$ with
pieces $B_h(V)_p\otimes\OO(2p)$, including its multiplication and
power diagram. Its automorphism functor, allowing automorphisms of
$\Pj(V)$, is represented by a finite-type affine complex group $H_h$.
There is a split exact sequence
\[
 1\longrightarrow K_h\longrightarrow H_h\longrightarrow
                          \operatorname{PGL}(V)\longrightarrow1.
\]
The splitting is the one induced by the explicit source-normalized
linearization, not a general assertion about automorphism extensions.
The fppf stack of forms of this fixed object is the algebraic
classifying stack $BH_h$. This statement permits twisted projective
source bundles but does not classify all graded algebra bundles
of those ranks, or all raw failure algebras.
\end{proposition}
\begin{proof}
Start with the natural $\operatorname{GL}(V)$ action on the apolar
quotient and on the tautological bundle. A scalar $t$ acts on
$B_h(V)_p$ by $t^{2p}$ and on $\OO(2p)$ by $t^{-2p}$. It therefore
acts trivially on their tensor product and on its multiplication.
The action factors through $\operatorname{PGL}(V)$, giving a group
homomorphism into automorphisms of the specified object whose
projection to source automorphisms is the identity. This is the
required splitting and defines the conjugation action on $K_h$.
Without this chosen linearization there would only be an extension,
and a semidirect-product notation would require justification.

An automorphism over the identity of $\Pj(V)$ is, in each degree,
a constant invertible matrix, since
$H^0(\Pj(V),\mathcal End(B_h(V)_p\otimes\OO(2p)))
=\operatorname{End}(B_h(V)_p)$. There are finitely many nonzero
degrees. Preservation of the unit and multiplication is a finite
collection of polynomial equations in these matrices and their
inverses. The resulting closed subgroup $K_h$ is affine of finite
type, and $H_h=K_h\rtimes\operatorname{PGL}(V)$ has the same
properties. An isomorphism torsor for fppf-local forms of this one
object is an $H_h$-torsor, and effective descent gives the inverse
construction. Thus the category of forms is precisely $BH_h$.
Affine finite-type groups over $\C$ are smooth, and the usual
classifying-stack atlas establishes the asserted algebraicity.
\end{proof}

\begin{corollary}[Complete boundary fibres on the effective failure image]
\label{cor:effective-full-fibres-v163}
On the effectively rigidified sharp pencil-failure image, families
satisfying \eqref{eq:first-boundary-class-v163} have the complete
boundary fibres of Theorem~\ref{thm:product-fibres-v163}, including
their nilpotent structures, component dimensions and intersections.
The exponent $h=1$ defines the envelope without an additional integer
choice; other $h\geq1$ have the same graph. The assertion is a
functorial consequence of the sharp inverse on effective families.
It is not an internal boundary construction on a closed algebra
before reconstruction and is not a statement about the full raw
Artin-algebra deformation stack.
\end{corollary}
\begin{proof}
Use the family inverse to recover the oriented source and pencil,
then the source-normalized envelope to obtain its multiplication
ideals. On the corank-two rank open the relevant graph is determined
by $I_{n-1}$; the universal power identity gives the same graph for
every $h$. The completed comparison and the full-fibre theorem
identify its embedded family. Their transition maps commute with
the effective equivalence and the source linearization, so all the
scheme structures descend, including the local ideal
$(eA,eB,e^2C)$ and its products. This proves the claim in the stated
category and no larger one.
\end{proof}

\section{Collision of two incidence points}
\label{sec:collision-family-v164}
The complete fibre of Theorem~\ref{thm:complete-fibre-v163} does not
by itself describe which of its points arise by colliding two reduced
incidences. We determine this specialization as a family, including
its vertical excess and its scheme multiplicities. All spaces in this
section are parameter spaces for embedded curves. A semistable model
of such a parameter space is not a semistable model of every curve
parametrized by it.

Retain the coefficient graph $\widetilde\Gamma_2\to B_2$ and the
point $o$ of Section~\ref{sec:complete-contact-v163}. Set
\begin{equation}\label{eq:collision-base-v164}
 T=\Spec\C[\delta],\qquad
 \gamma:T\longrightarrow B_2,\qquad
 \delta\longmapsto(x^2-\delta,0,0),\qquad
 \mathcal H=\widetilde\Gamma_2\times_{B_2}T.
\end{equation}
Let $\mathcal Y$ be the schematic closure of
$\mathcal H|_{\delta\ne0}$ in $\mathcal H$. Thus the closure is
formed \emph{after} this base change. Write $\Lambda\simeq\Pj^1$
for the space of lines through $P=[1:0:0]$ in the residual plane,
and $\mathcal Q_P\simeq\Pj^4$ for its conics through $P$.
Its linear subspace
\[
 S=\{Q\in\mathcal Q_P:Q\hbox{ is singular at }P\}\simeq\Pj^2
       =\operatorname{Sym}^2\Lambda
\]
parametrizes unordered pairs of lines through $P$, with repetitions.
We distinguish this plane from the entire conic component.

\begin{theorem}[The horizontal collision space and its excess]
\label{thm:collision-family-v164}
The scheme $\mathcal Y$ is a smooth integral threefold, projective
and flat over $T$. Define
\begin{equation}\label{eq:ordered-blowup-v164}
 Z=\operatorname{Bl}_{\Delta_\Lambda\times\{0\}}
           (\Lambda\times\Lambda\times\mathbb A^1_t),\qquad
 \iota(\ell_+,\ell_-,t)=(\ell_-,\ell_+,-t).
\end{equation}
The involution lifts to $Z$, and there is an isomorphism over $T$
\begin{equation}\label{eq:quotient-model-v164}
             \mathcal Y\simeq Z/\langle\iota\rangle,
                         \qquad \delta=t^2.
\end{equation}
For $\delta\ne0$ the geometric fibre is $\Lambda\times\Lambda$,
with the two factors attached at $x=\sqrt\delta$ and $x=-\sqrt\delta$.
The identification with an ordered product requires the displayed
root cover; it is not a global ordering over $T\setminus\{0\}$.

There are smooth divisors $E\simeq\mathbb F_2$ and $S\simeq\Pj^2$
in $\mathcal Y$ such that
\begin{equation}\label{eq:horizontal-special-v164}
       \mathcal Y_0=E+2S,\qquad
       D=E\cap S\simeq\Pj^1,\qquad
       \OO_S(S)=\OO_{\Pj^2}(-1).
\end{equation}
The equality is an equality of effective Cartier divisors in the
smooth threefold. The reduced support has normal crossings.
The curve $D$ is the negative section in $E$ and the doubled-line
conic in $S=\operatorname{Sym}^2\Lambda$.

The full pullback $\mathcal H$ is reduced but is not flat over $T$.
It has precisely the horizontal component $\mathcal Y$ and the
vertical component $\mathcal Q_P$ over $0$. Their scheme intersection
is the reduced plane $S$, and their union is scheme-theoretic:
\begin{equation}\label{eq:collision-union-v164}
 0\longrightarrow\OO_{\mathcal H}\longrightarrow
 \OO_{\mathcal Y}\oplus\OO_{\mathcal Q_P}
 \longrightarrow\OO_S\longrightarrow0.
\end{equation}
Closed-immersion direct images are implicit in this formula. The
entire $\C[\delta]$-torsion is killed by $\delta$ and is
\begin{equation}\label{eq:collision-torsion-v164}
 \ker(\OO_{\mathcal H}\longrightarrow\OO_{\mathcal Y})
  =i_*\mathcal I_{S/\mathcal Q_P}
  =\operatorname{Tor}^{\C[\delta]}_1(\OO_{\mathcal H},\C),
\end{equation}
where $i:\mathcal Q_P\hookrightarrow\mathcal H$.
Consequently colliding two reduced incidences reaches exactly
$E\cup S$ in the support of $X_2=\mathcal H_0$, not every conic
in $\mathcal Q_P$.
\end{theorem}

\begin{proof}
The target and the coefficients are retained throughout. In a conic
chart of Lemma~\ref{lem:HB-family-v163}, taking the base change
\eqref{eq:collision-base-v164} gives exactly
\begin{equation}\label{eq:collision-chart-v164}
 \mathcal H:\quad
 \C[\lambda,e,A,B,C,\delta]/(eA,eB,e^2C+\delta).
\end{equation}
Indeed the linear coefficient of $f$ gives $c=0$, its constant
coefficient gives $e^2C=-\delta$, the linear coefficient of $g$
gives $eB=0$, and that of $r$ gives $eA=0$. The two remaining
constant equations are then redundant. These are equations in a
regular chart of the normalized total graph; no fibre normalization
or fibrewise radical is being taken.

After eliminating $\delta$, the ideal is $(eA,eB)$. Saturation by
$\delta=-e^2C$ gives $(A,B)$. The horizontal chart is therefore
\begin{equation}\label{eq:horizontal-chart-v164}
 \C[\lambda,e,C],\qquad \delta=-e^2C.
\end{equation}
The other reduced component is $e=0$, and their intersection is
$e=A=B=0$. On the conic component this is exactly the singular-at-$P$
plane. The identity
\[
            (eA,eB)=(e)\cap(A,B)
\]
proves reducedness of the pullback on this chart. It also gives
\eqref{eq:collision-union-v164} locally. Its annihilator calculation
is
\[
       ((eA,eB):e^2C)=(A,B),
\]
so every base-torsion section is already killed by $\delta$.

For completeness we construct the projective horizontal model,
rather than identify it from its closed points. Choose an affine
coordinate $w$ on $\Lambda$ containing both directions and put
\[
 w_+=\lambda+d,\quad w_-=\lambda-d.
\]
Near the diagonal the ordered centre has ideal $(d,t)$. The two
blow-up charts and their involutions are
\begin{align*}
 &t=ed: &&(\lambda,d,e)\longmapsto(\lambda,-d,e),\\
 &d=ut: &&(\lambda,t,u)\longmapsto(\lambda,-t,u).
\end{align*}
Their invariant rings, respectively, are
\begin{equation}\label{eq:invariant-charts-v164}
 \C[\lambda,e,C],\quad C=-d^2,\quad\delta=-e^2C;
 \qquad \C[\lambda,u,\delta],\quad\delta=t^2.
\end{equation}
On the overlap $e=1/u$ and $C=-\delta u^2$. These are regular
rings. Common affine coordinates cover all pairs in
$\Lambda\times\Lambda$: a point of $\Lambda$ outside either of
two specified directions may be chosen as infinity. The invariant
charts and their target-coordinate translates glue, giving a
scheme quotient. This also verifies smoothness at the fixed locus;
no theorem about quotients by a \emph{free} action is applied there.
Projectivity follows, for example, by tensoring the translates of
a relatively ample line bundle on the blow-up, taking an even power
to trivialize stabilizer actions on fibres, and descending it to the
finite quotient. Equivalently its invariant section algebra gives
this projective quotient. The affine calculations exhibit the same
quotient directly.

We next map it to the retained-coefficient Hilbert graph. On the
first invariant chart use \eqref{eq:HB-ideal-v163} with $c=A=B=0$:
\begin{equation}\label{eq:collision-Hilbert-ideal-v164}
 (H^2+CG^2,\ xG-eH,\ xH+eCG),\qquad H=R-\lambda G.
\end{equation}
The fixed coefficients are $(x^2-\delta,0,0)$ because
$\delta=-e^2C$. On the second chart use the division family with
$f=x^2-\delta$, $g=r=0$, and $w=\lambda+ux$; its equations include
$fG=0$ and $R=wG$. The original flat-family results apply after
these base changes. On $u\ne0$ these ideals agree with
\eqref{eq:collision-Hilbert-ideal-v164}, since $e=1/u$.
Under changes of the direction coordinate they give the same
embedded family: this holds on the dense locus $t\ne0$ and hence
everywhere by separatedness of the Hilbert scheme on each integral
overlap. Their maps to the coefficient graph therefore glue.
They land in its regular division and conic charts, where total
normalization is an isomorphism.

For $t\ne0$, the first ideal places the direction $w_+$ at $x=t$
and $w_-$ at $x=-t$. It records all choices of the two lines,
including equality of their directions, which is covered by the
second chart. The common affine-coordinate argument covers the
whole generic fibre. The recovery maps in the conic and division
charts identify the quotient with the horizontal charts of
$\mathcal H$, not just with their supports. The two resulting
morphisms are inverse on these charts, proving
\eqref{eq:quotient-model-v164}. This constructs the isomorphism
algebraically; formal normal forms are not used to manufacture it.

The strict transform of $t=0$ in $Z$ is
$\Lambda\times\Lambda$. Its quotient is
$\operatorname{Sym}^2\Lambda=S$. The exceptional divisor is
\[
 \Pj(N_{\Delta_\Lambda/(\Lambda\times\Lambda)}\oplus\OO)
       =\Pj(\OO(2)\oplus\OO)=\mathbb F_2.
\]
The involution acts trivially on this divisor, since it is $-1$
on both normal directions before projectivization. Their common
curve is the diagonal, whose image in $S$ is the doubled-line
conic. On the chart \eqref{eq:horizontal-chart-v164}, $S=(e=0)$,
$E=(C=0)$, and $\delta=-e^2C$. This proves the divisor multiplicities
and transverse reduced intersection. On $E$ the coordinate
$u=1/e$ is the affine tangent-jet coordinate, so $e=0$ is its
negative section, as in \eqref{eq:adjacency-family-v163}.
Since $2S+E$ is principal and $E|_S$ is the degree-two conic,
$2S|_S=-2H_S$. The Picard group of $S=\Pj^2$ has no torsion,
giving $\OO_S(S)=\OO_S(-1)$.

The quotient is integral and hence torsion-free as a module over
$\C[\delta]$. Localize this ring at each prime: its nonfield
localizations are discrete valuation rings. The valuation-ring
flatness criterion \cite[Tag~0539]{Stacks164} proves flatness at
every base point, and hence over $T$.
The complete first-contact theorem covers every point of
$\mathcal H_0$ by its conic charts or primitive charts. On a primitive
chart the pullback is $\C[\lambda,u,\delta]$ and has no vertical
part. Away from $0$ all fibres are the reduced-incidence product.
Thus the local descriptions just proved exhaust $\mathcal H$.
They glue \eqref{eq:collision-union-v164}; no unexamined component
can be supported over a different point of $T$. The kernel of the
projection to $\mathcal Y$ is, on the vertical component, exactly
its ideal of $S$. Resolving $\C$ over $\C[\delta]$ by multiplication
by $\delta$ identifies this kernel with the stated Tor sheaf.
It is nonzero, so $\mathcal H$ is not flat. The last assertion follows
from the proper flat horizontal model and
\eqref{eq:horizontal-special-v164}.
\end{proof}

\subsection{The ramified model and monodromy}
\begin{theorem}[Ordered resolution of the collision]
\label{thm:ordered-resolution-v164}
The normalization of
$\mathcal Y\times_{\mathbb A^1_\delta}\mathbb A^1_t$, with
$\delta=t^2$, is $Z$ in \eqref{eq:ordered-blowup-v164}.
It is smooth and its fibre over $t=0$ is the reduced normal-crossing
union of $\Lambda\times\Lambda$ and $\mathbb F_2$, meeting along
the diagonal section. In particular the multiplicity two in
\eqref{eq:horizontal-special-v164} is caused by forgetting the order
of the colliding supports.

On the punctured complex base the monodromy on
$H^2(\mathcal Y_\delta,\mathbb Q)$ interchanges the two ruling
classes. Its invariant and anti-invariant spaces both have dimension
one. These assertions concern the degeneration of the parameter
space. They do not replace thick embedded Hilbert tails by reduced
stable-map tails.
\end{theorem}
\begin{proof}
The first invariant chart after base change is
\[
 \C[\lambda,e,C,t]/(t^2+e^2C).
\]
Its normalization is obtained by adjoining $d=t/e$, with
$d^2=-C$, and is $\C[\lambda,e,d]$, $t=ed$. This extension is
finite because $d$ satisfies a monic equation, and it has the same
fraction field; the polynomial ring is normal. The second chart
becomes $\C[\lambda,t,u]$ with $d=tu$. These finite normal charts
glue by $e=1/u$, giving precisely the ordered blow-up. They prove
normalization of the base change itself, not of its special fibre.
The central equations $t=ed$ and $t=0$ in the two smooth charts
give the reduced normal crossings and the stated components.

Over $t\ne0$ the blow-up is the constant product
$\Lambda\times\Lambda\times\mathbb G_m$. Descent identifies
$(\ell_+,\ell_-,t)$ with $(\ell_-,\ell_+,-t)$.
A loop around $\delta=0$ exchanges the roots and hence the ruling
classes. Their sum and difference give the two claimed eigenspaces.
This calculation uses the ordering cover, not a claim that the
whole collision family is topologically locally trivial at $0$.
\end{proof}

\subsection{The nilpotent sheaf of the full contact fibre}
\begin{theorem}[Global nilradical and coherent cohomology]
\label{thm:nilradical-sheaf-v164}
Let $j:S=\Pj^2\hookrightarrow X_2$ be the singular-conic plane.
The nilradical $\mathcal N$ of the complete scheme fibre satisfies
\begin{equation}\label{eq:nilradical-sheaf-v164}
                \mathcal N\simeq j_*\OO_{\Pj^2}(-1),
                         \qquad\mathcal N^2=0.
\end{equation}
The isomorphism specifies the module; it does not assert that the
square-zero algebra extension is split. Moreover
\[
 H^0(X_2,\OO_{X_2})=\C,\qquad
 H^q(X_2,\OO_{X_2})=0\quad(q>0).
\]
The same coherent-cohomology statement holds for every complete
fibre $X_2^r\times(\Pj^1)^s$ of
Theorem~\ref{thm:product-fibres-v163}, although for $r>0$ that
fibre is nonreduced and is not Cohen--Macaulay on its nilpotent support.
\end{theorem}
\begin{proof}
Restrict the quotient $\OO_{X_2}\to\OO_{\mathcal Y_0}$ in the
conic charts. Its source and target have presentations
\[
 \frac{\C[\lambda,e,A,B,C]}{(eA,eB,e^2C)}
       \longrightarrow
 \frac{\C[\lambda,e,C]}{(e^2C)}.
\]
Both nilradicals are generated by $eC$, and the map identifies
them as modules on $S$. Off $S$ both vanish. This is consequently
a global isomorphism of nilradical sheaves.
In the smooth threefold $\mathcal Y$, the nilradical of $E+2S$ is
\[
 \OO_{\mathcal Y}(-E-S)/\OO_{\mathcal Y}(-E-2S)
            \simeq\OO_S(-E-S)\simeq\OO_S(S).
\]
The last isomorphism uses the principal divisor $E+2S$;
\eqref{eq:horizontal-special-v164} now gives
\eqref{eq:nilradical-sheaf-v164}. Thus the local nilpotent generator
has a nontrivial global twist; its support alone would not determine it.

The reduced union has the exact sequence
\[
 0\longrightarrow\OO_{(X_2)_{\rm red}}\longrightarrow
 \OO_{\Pj^4}\oplus\OO_{\mathbb F_2}
 \longrightarrow\OO_{\Pj^1}\longrightarrow0.
\]
Each of the three smooth terms has constants as its only coherent
cohomology: for $\mathbb F_2$ this follows from its projective-line
bundle over $\Pj^1$. The map on constants is the difference and is
surjective. Also $H^q(\Pj^2,\OO(-1))=0$ for every $q$.
Apply the cohomology sequences first to this union and then to
$0\to\mathcal N\to\OO_{X_2}\to\OO_{(X_2)_{\rm red}}\to0$.
They give the assertions. The product statement follows from
Kunneth over $\C$ and the already proved scheme-product theorem.
It is not an inference of reducedness or flatness from constant
coherent cohomology.
\end{proof}

\subsection{Several simultaneous collisions}
\begin{corollary}[Independent collision multiplicities]
\label{cor:multiple-collisions-v164}
For $r$ independent copies of \eqref{eq:collision-base-v164} and
$s$ additional simple contacts, the horizontal collision space is
$\mathcal Y^r\times(\Pj^1)^s$ over $\mathbb A^r_{\boldsymbol\delta}$.
Its central fibre has $2^r$ reduced components
\[
 S^{|J|}\times E^{r-|J|}\times(\Pj^1)^s,
                   \qquad J\subseteq\{1,\ldots,r\},
\]
each of dimension $2r+s$ and of cycle multiplicity $2^{|J|}$.
Intersections replace differing factors by $D$.
After $\delta_i=t_i^2$ and normalization the ordered model is
$Z^r\times(\Pj^1)^s$. Its total space is smooth, and the inverse
image of the union of the base coordinate hyperplanes is a reduced
simple-normal-crossing divisor. The fibre over the origin instead
has the product local equations $u_iv_i=0$ in disjoint coordinate
pairs, omitting one factor at a smooth point of a component.
The full central pullback is $X_2^r\times(\Pj^1)^s$, whose conic
factors are $\Pj^4$, not $S$. Thus the horizontal and full fibres
have different component dimensions even though both keep all
choices at separated supports.
\end{corollary}
\begin{proof}
Form the product of the explicit flat horizontal models. The
open where all $\delta_i$ are nonzero is dense in this integral
product. On each chart its closure is obtained by saturating with
$\prod_i\delta_i$, and removes precisely the vertical ideals in
the individual factors. Hence it is the horizontal closure of the
full product pullback. At the generic point of a component indexed
by $J$, the central equations have the forms $e_i^2$ for $i\in J$
and $C_i$ for the other indices, up to units. Their independent
local lengths multiply to $2^{|J|}$. The intersection descriptions
follow from the two smooth divisors in each factor. The finite
normalization charts used in Theorem~\ref{thm:ordered-resolution-v164}
tensor over the independent complex variables, giving the stated
smooth ordered model. Locally $t_i=u_iv_i$ at the two-component
locus and $t_i=u_i$ away from it. The equation of the full inverse
image of the coordinate boundary is therefore the product of the
appearing distinct coordinates, whereas the central fibre is cut
out by the individual $t_i$. The last assertion is the complete-fibre
product theorem, not a normalization-of-fibre claim.
\end{proof}

\section{Collisions under the effective inverse}
\label{sec:effective-collision-v164}
The collision calculation has an explicit realization by genuine
symmetric pencils in a fixed target. This realizes the entire family
of possible Hilbert limits over the chosen pencil family; a single
parameter value is not assigned a preferred tail.

\begin{proposition}[A genuine collision and its effective transfer]
\label{prop:genuine-collision-v164}
Let $V=\C^4$ and consider the family of homogeneous symmetric pencils
\begin{equation}\label{eq:genuine-collision-v164}
 A_\delta(s,t)=L_\delta(s,t)\oplus L_\delta(s,t),\qquad
 L_\delta(s,t)=
 \begin{pmatrix}-\delta s&t\\t&-s\end{pmatrix}.
\end{equation}
Every pencil is regular and has rank at least two. For $\delta\ne0$
it has two reduced corank-two incidences, and at $\delta=0$ it has
one corank-two incidence with Smith exponents $(2,2)$. The member
at infinity $s=0$ is nonsingular. The pullback of the ambient
normalized Hilbert modification along this family is isomorphic,
with the indicated residual coordinates, to $\mathcal H$ of
Theorem~\ref{thm:collision-family-v164}. In particular its horizontal
closure, its vertical excess and its torsion sheaf are those computed
there, not merely families of selected curves inside its fibres.

The same statements transport to the effectively rigidified sharp
failure-family image by the established inverse. The coherent
$\delta$-torsion and the horizontal closure are preserved under this
equivalence. No assertion is made about deformations leaving that
effective image or about a boundary operation preceding reconstruction
of the source and pencil.
\end{proposition}
\begin{proof}
On $s=1$, put $x=t$. Each block has determinant $\delta-x^2$ and
has the invertible lower-right entry $-1$. Its Schur complement is
$f=x^2-\delta$. The polynomial congruence
\[
 U(x)=\begin{pmatrix}1&0\\x&1\end{pmatrix},\qquad
 U(x)^{\mathsf t}L_\delta(1,x)U(x)=\diag(f,-1)
\]
has determinant one. Apply it on both blocks, and permute the unit
blocks to the front. The residual matrix is $fI_2$. Its triple of
normal coordinates is $(f,0,f)$; subtracting the first entry from
the third gives the retained-coefficient triple $(f,0,0)$.
The determinant of the full pencil is $(t^2-\delta s^2)^2$.
The unit blocks give the rank bound, and the affine Schur form
gives the stated incidences and exponents. At $s=0$ both blocks are
$\left(\begin{smallmatrix}0&t\\t&0\end{smallmatrix}\right)$,
so infinity contributes no incidence or Hilbert choice.

We detail why these calculations describe the pullback rather than
only its closed points. Over the finite support scheme
$\Spec\C[\delta,x]/(f)$, the congruence and the row subtraction are
algebraic invertible changes in the \emph{relative incidence target}
over the actual pencil line. They are undone before the curve is
viewed in the fixed space $\mathrm{CQ}(V)$; $U(x)$ is not being
used as one constant automorphism of that space. The forced main
lift exists algebraically on this family, since the normal ideal
$(f,0,0)$ is principal and the residual direction is fixed.

In the conic charts of the model $\mathcal H$ the equations
$eA=eB=0$ yield
\[
           xG=eH,\qquad xH=-eCG,\qquad fG=fH=0.
\]
In a division chart the same annihilation is $fG=0$ and $R=wG$.
Thus the quotient of the fixed main ideal by the curve ideal is
annihilated by $f$, including the parameter-scheme nilpotents in
special fibres. It is an algebraic quotient on a finite locally
free neighbourhood, not data requiring convergence of a formal
gauge. Undo the displayed polynomial congruence on this quotient
and patch with the forced main lift away from the support. This
constructs a morphism of the whole model to the actual fixed-target
Hilbert pullback. Flatness of the families is detected on these
local pieces, and their exterior degrees add to
$\binom42=6$, with Euler characteristic one.

Conversely restrict an actual family to this incidence target and
apply the inverse changes. Proposition~\ref{prop:embedded-functors-v163}
identifies its completed embedded quotients with the retained
coefficient model; its smooth coefficient factors are fixed by
\eqref{eq:genuine-collision-v164}. In particular the displayed
annihilation and all chart equations hold on every completed local
ring, not just on geometric fibres. Faithful completion detects
these finitely generated ideal containments and their equality.
The algebraic quotient construction above and its inverse are
therefore inverse on the division and conic charts. These cover
the fibres by the reduced-incidence theorem for $\delta\ne0$ and
the complete first-contact theorem for $\delta=0$. The ambient
total charts are regular there, so taking their normalization does
not alter them before restriction to the family. This proves the
scheme isomorphism of pullbacks with the fixed target retained.

Finally the effective inverse is an equivalence on the specified
pencil-family image, with the intrinsic oriented source included.
Applying the Hilbert construction through that equivalence preserves
its structure map and the actual local rings. The kernel of
multiplication by $\delta$, its saturation defining the horizontal
closure, and the finite quotient descriptions are consequently
preserved. This is an effective-family consequence of the inverse,
not a new assertion about the stack of all finite algebras.
\end{proof}

Distinct copies of \eqref{eq:genuine-collision-v164}, translated to
pairwise distinct affine spectral points, give simultaneous
collisions. Restrict the base so that the corresponding support
schemes do not meet. Algebraic changes of the source frame and
permutations of these supports act on the above descriptions by
transition isomorphisms. The disjoint-contact comparison then gives
Corollary~\ref{cor:multiple-collisions-v164}. This supplies actual
pencil families with the computed multiplicities without claiming a
classification of collisions passing through higher corank.

\section{Descent and complete-local details}
\label{sec:formal-details-v154}
The following statements supply the local-to-global steps used in
Sections~\ref{sec:formal-images-v153}--\ref{sec:reciprocal-v153}.
They also specify which assertions concern a completed image, a
normalization fibre, or a family on an exact-divisor stratum.

\begin{lemma}[Complete-local surjectivity and reduced completion]
\label{lem:complete-local-v154}
In the cotangent-surjectivity assertion of
Lemma~\ref{lem:formal-image-v153}, the rings are Noetherian complete
local $\C$-algebras with residue field $\C$, and the homomorphism
is local, continuous, and the identity on that residue field. Under
these hypotheses cotangent surjectivity implies ring surjectivity.
A reduced local ring essentially of finite type over $\C$ has
reduced completion.
\end{lemma}
\begin{proof}
Choose preimages in the source maximal ideal of generators of the
target cotangent space. Their degree-$q$ products span every quotient
$\mathfrak m_B^q/\mathfrak m_B^{q+1}$. Given an element of $B$, first
lift its residue, then its error successively in degrees $q=1,2,\ldots$.
The correction in degree $q$ lies in $\mathfrak m_R^q$. Their sum
converges in $R$, and continuity and separatedness imply that its
image is the desired element of $B$. This is not a statement about
arbitrary noncomplete local homomorphisms.

Finite-type complex rings are excellent. Quasi-excellent local
rings are Nagata and have geometrically reduced formal fibres;
see \cite[Tags 07QV and 0BJ0]{StacksLocalV154}. For a reduced local
ring $A$, the injection into the product of its finitely many
minimal-prime quotients remains injective after flat completion.
Each quotient is a Nagata domain and its completion is reduced.
Thus $\widehat A$ is a subring of a product of reduced rings.
This justifies the reduced completed-image step of the earlier lemma.
\end{proof}

\begin{lemma}[Coprime cluster products over Artin bases]
\label{lem:hensel-products-v154}
Let $A$ be a local Artin $\C$-algebra. If the reduction of a monic
polynomial has a factorization into pairwise coprime monic factors,
there is a unique factorization into monic lifts of the same degrees.
In a factorization fibre, the functor of divisors with prescribed
degrees in the distinct residual clusters is the product of the
cluster divisor functors. Its coordinate algebra is therefore the
tensor product of their Artin coordinate algebras.
\end{lemma}
\begin{proof}
For two coprime factors $p,q$, the linear correction map
\[
 (\dot p,\dot q)\longmapsto q\dot p+p\dot q,
 \quad\deg\dot p<\deg p,\quad\deg\dot q<\deg q,
\]
is an isomorphism on coefficient spaces over the residue field.
Its kernel is zero by coprimeness, and source and target have equal
dimension. It is therefore an isomorphism over $A$ as well. Across
a square-zero extension of Artin rings it uniquely corrects a
lifted product to the given polynomial. Induction on the nilpotence
order and then on the number of clusters proves existence and
uniqueness. Apply the same argument to both the marked divisor and
its complement. The cluster product and decomposition constructions
are mutually inverse on every Artin base; Yoneda gives the asserted
product of finite schemes. Passing to inverse limits gives the
corresponding assertion on complete local bases.
\end{proof}

\begin{lemma}[The local choice and descent in the binary stratum]
\label{lem:residual-selection-v154}
The base-change identification of
Theorem~\ref{thm:binary-fibres-v153} holds as an identification of
incidence functors, not only on closed fibres. The required choice
of a member of the residual space coprime to the common divisor can
be made on a Zariski open cover of the exact-gcd stratum; the choices
are auxiliary and the resulting identifications glue uniquely.
\end{lemma}
\begin{proof}
Write a point of the stratum as $(G,L)$ in
$\Pj(S_s)\times X_{D-s,r}$. The vanishing of a member of $L$ at
one of the finitely many roots of $G$ is a proper linear condition
on that member. Over $\C$ a finite union of these proper hyperplanes
does not cover $L$. Choose a linear combination avoiding them, lift
its coefficients on a bundle-trivializing neighbourhood, and invert
its resultant with $G$. This gives the required Zariski neighbourhood.
These neighbourhoods cover the stratum and may be pulled back by
any base change.

On the marked-divisor incidence, pass locally to a monic coordinate
for the proposed divisor $f$. The quotient by $f$ is finite free.
On every geometric fibre $f$ divides $G$, so the chosen residual
member is invertible in that quotient, by coprimeness. Its
multiplication matrix has invertible determinant locally, hence it
is invertible also over a nonreduced base. Thus $f\mid G L$ implies
$f\mid G$ as an equation on that base, and the converse is immediate.
Monic division forces the remaining subspace. The identification is
independent of which coprime residual member established cancellation,
so it glues on overlaps. No lifting of arbitrary image points to the
normalization is asserted.

The finite factorization map used there is flat by miracle flatness:
its target local rings are regular, its source local rings are
Cohen--Macaulay, it is finite, and dimensions agree. These are exactly
the hypotheses of \cite[Tag 00R4]{StacksLocalV154}. This names the
flatness input separately from the incidence-functor argument.
\end{proof}

\begin{lemma}[The coordinate equalizer and coherent conductor]
\label{lem:equalizer-v154}
For the block ideals $I_A$ of Theorem~\ref{thm:squarefree-v153},
the completed image ring is the equalizer of the branch rings with
their restrictions to all pairwise intersections. This assertion
uses the coordinate monomial structure. The completed conductor
calculation there implies equality of the corresponding coherent
ideal sheaves on the squarefree open.
\end{lemma}
\begin{proof}
A monomial has a coefficient on each branch where its support
survives. If it survives on two branches, it also survives on their
coordinate intersection. Pairwise agreement therefore assigns it
one well-defined coefficient. Collecting these coefficients gives
one formal series modulo the intersection of the $I_A$; conversely
such a series has compatible restrictions. Formal convergence causes
no issue, since there are finitely many monomials in each degree.
The claim would be false for general configurations of ideals without
this coordinate argument.

Write $A\subset B$ for the local finite normalization algebra.
The conductor is $\Ann_A(B/A)$, a coherent ideal. Flat completion
preserves the exact finite-module sequence and its annihilator:
compute the latter as the intersection of the kernels of multiplication
on a finite generating set. The completed finite algebra is the
product of the normal branch rings. Thus the conductor and the ideal
of $Y_{g+1}$ have equal completed stalks at every squarefree closed
point. Faithful flatness gives equal stalks, and coherence on a
finite-type complex scheme gives equality of the ideal sheaves.
\end{proof}

\begin{proposition}[Finiteness and uniqueness of the collision diagram]
\label{prop:atlas-functor-v154}
Every branch map in Theorem~\ref{thm:atlas-v153} is finite, and the
completed incidence diagram is determined, up to the unique
isomorphism respecting its incidence data, by its marked-divisor
functor. It is independent of the chosen frame and root coordinates.
\end{proposition}
\begin{proof}
For each cluster $P_i=A_iB_i'$, the coefficients of the monic $A_i$
are integral over the coefficients of $P_i$. This follows from the
finite factorization morphism, or by expressing them as elementary
symmetric functions of a selected subset of the roots in a finite
splitting algebra. Division by a monic polynomial expresses every
coefficient of $B_i'$ polynomially in those of $A_i$ and $P_i$.
Likewise $R_{ij}=A_iU_{ij}$ expresses all quotient coefficients in
terms of $A_i$ and $R_{ij}$. Consequently the branch algebra is
generated over the ambient coefficient algebra by finitely many
integral coefficients of the $A_i$, so it is finite. Completing
preserves this finite-module assertion.

The equations represent the functor of a subbundle with a marked
divisor and its uniquely determined quotient. A change of frame or
local coordinate gives inverse natural transformations of this same
functor and of its forgetful map to the ambient Grassmannian.
The complete representing rings and their maps are uniquely
isomorphic by Yoneda. Their scheme-theoretic images consequently
agree as completed diagrams. This proves coordinate independence
without identifying arbitrary Artin points of an image with lifted
Artin points of its normalization.
\end{proof}

\paragraph{Partition closures.}
For a fixed partition $(s_1,\ldots,s_b)$, the proper map
$(\Pj^1)^b\to\Pj(S_s)$ sending $(\ell_i)$ to
$\prod\ell_i^{s_i}$ has closed image. Its open set of distinct
$\ell_i$ is dense. At the boundary precisely some of these points
coincide, and the associated parts add. Every such coincidence can
be approached from distinct points. This proves exactly the merging
closure order in the fixed-degree symmetric product; compare the
coincident-root constructions of \cite{Chipalkatti2003,Kurmann2012}.
It is not a closure classification for all analytic singularities
of the completed collision atlas.

\paragraph{Equivariance of the singular Fitting calculation.}
The group used in Theorem~\ref{thm:fitting-v153} is
$\operatorname{GL}_m(\C)$, acting simultaneously on the columns
$u$ and $w$ and fixing $\Delta$. It preserves the presentation and
therefore its intrinsic Fitting ideal. In the normalization the
three displayed minors span, under this group, respectively
$\Sym^m(\C^m)t$, $\Delta\Sym^m(\C^m)$, and
$\Sym^{m+1}(\C^m)$, with the coordinate/dual convention determined
by the action on $u$. To see spanning directly, their translates
are nonzero scalar multiples of the corresponding powers of all
linear forms; polarization in characteristic zero spans each
symmetric power. Multiplication by $R$ gives exactly the modules
asserted in \eqref{eq:fitting-v153}. This is stronger than invoking
invariance without identifying the representations.

\paragraph{Specialization of minimal relations.}
For the surjective coefficient specialization $\pi:P\to P'$ to
one coordinate in Theorem~\ref{thm:reciprocal-structure-v153}, one
has $\pi(\mathfrak n)=\mathfrak n'$ and
$\pi(\mathcal I)=\mathcal I'$. Consequently
$\pi(\mathfrak n\mathcal I)=\mathfrak n'\mathcal I'$.
A coefficient whose image is a minimal binary relation cannot lie
in $\mathfrak n\mathcal I$. This explicitly justifies the nonzero
irreducible coefficient map used in that proof. All reciprocal
series and socle weights in these arguments use the weight of a
$Q_i$ coefficient equal to $i$. They do not claim ordinary
maximal-ideal Hilbert functions; \eqref{eq:F22-local-hilbert-v154}
gives one such function separately.

% BEGIN PRESERVED PART 44-conductor-and-collisions-v152.tex
\section{Conductors and local models of divisor collisions}
\label{sec:conductor-v152}

We retain the notation of Section~\ref{sec:canonical-boundary-v151}.
All completed local rings below are completed at closed complex points.
For a reduced local ring $A$ with finite normalization $B$, its
\emph{conductor} is $\mathfrak c=\Ann_A(B/A)$, regarded also as an
ideal of $B$.  We use the intrinsic singular subscheme defined by
$\Fitt_{\dim A}\Omega_{A/\C}$; on completions, differentials mean
continuous differentials.  The conductor and this Fitting ideal need
not agree.  Our first result determines both on a multivariate
stratum.  The second identifies the way two branches merge when a
common binary root becomes double.

The formal-image and separated-germ steps below are isolated in
Lemma~\ref{lem:formal-image-v153}. Lemma~\ref{lem:coprime-lifts-v153}
gives the Artin cancellation and a direct positive lower bound for the
split codimension. Theorem~\ref{thm:atlas-v153} supplies the full
Weierstrass incidence functor in every collision order;
Theorem~\ref{thm:fitting-v153} evaluates the singular ideal for all ranks.

\subsection{Two coprime divisor choices}
\begin{theorem}[The split-divisor stratum]\label{thm:split-conductor-v152}
Suppose $h>g$, $2\leq r\leq s_{h-g}$, and
\[
 K=abL,\qquad [a],[b]\in\Pj(S_g),\qquad
 L\in\Gr(r,S_{h-g}),\qquad \gcd(L)=1.
\]
Assume that $a,b$ are coprime and that the only degree-$g$ divisors of
$ab$, up to scalar, are $a$ and $b$.  Put
\begin{align*}
 t&=2(s_g-1)+r(s_{h-g}-r),\\
 c&=r(s_h-s_{h-g})-(s_g-1),\qquad d_g=t+c.
\end{align*}
Then $c\geq1$, and there are formal coordinates in which
\begin{equation}\label{eq:split-ring-v152}
 \widehat{\OO}_{Y_g,K}=
 \C[[z_1,\ldots,z_t,x_1,\ldots,x_c,y_1,\ldots,y_c]]/
                       (x_i y_j:1\leq i,j\leq c).
\end{equation}
Its normalization is
$\C[[z,x]]\times\C[[z,y]]$.  The conductor is
\begin{equation}\label{eq:split-conductor-v152}
 \mathfrak c=(x_1,\ldots,x_c,y_1,\ldots,y_c),\qquad
 \widehat{\OO}_{Y_g,K}/\mathfrak c=\C[[z]],
\end{equation}
and the intrinsic singular subscheme has ideal $\mathfrak c^{\,c}$.
The two smooth branches meet scheme-theoretically in a smooth
$t$-dimensional subscheme.  The local ring has multiplicity two and
depth $t+1$; it is Cohen--Macaulay exactly when $c=1$.
\end{theorem}

\begin{proof}
We first determine the intersection as a scheme, not by comparing
closed points or tangent dimensions.  The normalization has precisely
two points above $K$, namely $(a,bL)$ and $(b,aL)$.  At either point,
the overlap in Lemma~\ref{lem:fibre-equations-v151} is zero.  The
map from the completed local ring of the ambient Grassmannian to
each completed source ring is consequently surjective: its cotangent
map is surjective, and completeness and Nakayama's lemma lift
successive homogeneous terms.  Thus each source completion is a
closed smooth formal branch of the image.

Here is the required assertion over infinitesimal bases.  Let $A_0$
be a local Artin $\C$-algebra.  A point on both formal branches carries
unique marked lifts $a_{A_0},b_{A_0}$ near $a,b$.  If homogeneous
polynomials have coprime reductions, those reductions form a regular
sequence in the polynomial ring over $\C$.  Its graded Koszul
complex is exact.  In each degree the polynomial modules are finite
free; lifting a splitting of the complex, or applying the local
criterion for flatness successively, shows that its lift over $A_0$
is exact and that the quotient is $A_0$-flat.  In particular
\begin{equation}\label{eq:coprime-artin-v152}
 (a_{A_0})\cap(b_{A_0})=(a_{A_0}b_{A_0}).
\end{equation}
This also follows by noting that $b_{A_0}$ is a nonzerodivisor
modulo $a_{A_0}$ and solving a relation between them.  Applying
\eqref{eq:coprime-artin-v152} to the columns of the relation
subbundle gives, uniquely,
\[
                  K_{A_0}=a_{A_0}b_{A_0}L_{A_0}.
\]
The residual module is a subbundle: multiplication by
$a_{A_0}b_{A_0}$ is a split injection in degree $D$, and the
quotient of the ambient free module by $K_{A_0}$ is free.  The
kernel of the induced surjection onto the complementary quotient
is therefore free.  This proves that the formal intersection is
represented by the product of the two divisor charts and the
Grassmannian chart for $L$, with no nilpotents omitted.

The differential of this product parametrization is injective.
Indeed, an infinitesimally constant product gives
\[
       (\dot a b+a\dot b)l+ab\dot l\in abL\quad(l\in L).
\]
Taking valuations at the factors of $a$, and using
$\gcd(a,b)=\gcd(L)=1$, forces $a\mid\dot a$; similarly
$b\mid\dot b$.  These are the two projective scalar directions,
and then $\dot L=0$.  The intersection is thus a smooth formal
closed subscheme of dimension $t$.

Let $R_0$ be the ambient completed local ring and $I_a,I_b$ the
ideals of these branches.  Finiteness of the normalization, exactness
of completion on finite modules, and the definition of the
scheme-theoretic image give
\[
 \widehat{\OO}_{Y_g,K}=R_0/(I_a\cap I_b)
   = (R_0/I_a)\mathbin{\times}_{R_0/(I_a+I_b)}(R_0/I_b).
\]
This is the elementary conductor-square construction; see also
\cite[Lemma~1.3 and \S1.3.2]{Ferrand2003}.  All three quotient rings
are regular.  Choosing compatible regular parameters for the two
surjections gives \eqref{eq:split-ring-v152}.  The branch dimension
is $d_g=s_g-1+r(s_h-r)$, so its codimension along the intersection
is the asserted $c$.  It is positive: if $c=0$, both branch ideals
would coincide locally with the intersection ideal, contradicting
the two distinct normalization branches of a reduced local ring.

In the fibre-product description, an element belongs to the
conductor exactly when both its restrictions to the intersection
vanish.  This gives \eqref{eq:split-conductor-v152}.  The Jacobian
of the relations $x_i y_j$ has nonzero entries only among the $x_i$
and $y_j$.  The Fitting ideal in question is its ideal of $c$-minors.
Every such minor belongs to $\mathfrak c^{\,c}$.  Conversely,
choosing the $c$ columns indexed by the $y_j$ and one row $(i_j,j)$
for each $j$ gives every monomial $x_{i_1}\cdots x_{i_c}$.
The symmetric choice gives every degree-$c$ monomial in the $y$'s.
Mixed monomials vanish in the ring.  Thus the ideal is exactly
$\mathfrak c^{\,c}$, including its nonreduced structure.

Finally use
\[
 0\longrightarrow \widehat{\OO}_{Y_g,K}
 \longrightarrow\C[[z,x]]\oplus\C[[z,y]]
 \longrightarrow\C[[z]]\longrightarrow0,
\]
where the last arrow is the difference of restrictions.  The depth
lemma gives depth $t+1$; for $c=1$ this equals the dimension.
For $c>1$ it is strictly smaller.  Additivity of leading Hilbert
coefficients gives multiplicity two, since the last module has
smaller dimension than the two regular branches.
\end{proof}

\begin{remark}\label{rem:split-scope-v152}
The hypotheses are nonempty in every number of variables: for $e=2$
one may take $g=1$, and for $e\geq3$ one may take two distinct
irreducible forms of degree $g$.  The theorem treats an explicitly
specified stratum, not all pairs of overlapping degree-$g$ divisors.
In particular, it does not substitute a binary factorization rule
for multivariate divisibility.  Its high-codimension branch
intersections need not be normal-crossing divisors.
\end{remark}

\subsection{A double common root}
For $m\geq1$, put $R_m=\C[[u_1,\ldots,u_m,\Delta]]$ and
$I=(u_1,\ldots,u_m)\subset R_m$.  Define
\begin{equation}\label{eq:pinch-algebra-v152}
 A_m=R_m\oplus I t\ \subset\ B_m=\C[[u_1,\ldots,u_m,t]],
                    \qquad \Delta=t^2.
\end{equation}
The right side of \eqref{eq:pinch-algebra-v152} is a subring,
not a square-zero extension: $(It)^2=\Delta I^2$.  For $m=1$
it is the ordinary pinch-point ring.

\begin{theorem}[Collision of two binary divisor choices]
\label{thm:binary-pinch-v152}
Let $e=2$, $g=1$, $D=h+1$, and $2\leq r\leq D-1$.
At a point $K_0=\ell^2 L$ with
$L\in\Gr(r,S_{D-2})$ and $\gcd(L)=1$, set
\[
             m=r-1,\qquad q=r(D-1-r)+1.
\]
Then
\begin{equation}\label{eq:binary-completion-v152}
              \widehat{\OO}_{Y_1,K_0}\simeq A_m[[z_1,\ldots,z_q]].
\end{equation}
Writing $w_i=u_i t$, a complete presentation of $A_m$ is
\begin{equation}\label{eq:pinch-equations-v152}
 \C[[u_1,\ldots,u_m,\Delta,w_1,\ldots,w_m]]/
 \bigl(u_iw_j-u_jw_i,\ w_iw_j-\Delta u_i u_j\bigr),
\end{equation}
where the first relations have $i<j$ and the second have $i\leq j$.
Its normalization is $B_m$, and its conductor is
\[
 \mathfrak c=(u_1,\ldots,u_m,w_1,\ldots,w_m),\qquad
 \mathfrak c B_m=I B_m.
\]
The induced map on conductor quotients is
\begin{equation}\label{eq:conductor-cover-v152}
              \C[[\Delta]]\longrightarrow\C[[t]],\qquad\Delta\longmapsto t^2.
\end{equation}
Thus the two reduced divisor lifts merge through a ramified double
cover of the conductor.  The completed local ring has dimension
$q+m+1$ and depth $q+2$.  It is Cohen--Macaulay if and only if $r=2$.
For $r=2$, its intrinsic singular ideal, with the smooth variables
suppressed, is $(w,\Delta u,u^2)$, whereas its conductor is $(u,w)$.
\end{theorem}

\begin{proof}
Choose an affine coordinate $x$ with the root of $\ell$ at zero.
Some member $F_0\in K_0$ has multiplicity exactly two there, since
$\gcd(L)=1$.  Extend it to a frame $F_0,\ldots,F_{r-1}$.
In the completed local ring of the ambient Grassmannian the formal
Weierstrass theorem writes the universal $F_0$ as a unit times
\[
                         P(x)=x^2+a x+b.
\]
Divide each $F_i$, $i>0$, by this monic polynomial.  Its remainder
is $u_i x+v_i$.  The $2r$ functions
$a,b,u_1,v_1,\ldots,u_m,v_m$ are part of a regular parameter
system.  To see this rather than assume versality, note that
$K_0$ has zero two-jet at the root, while
$S_D/K_0\to\C[x]/(x^2)$ is surjective.  The Grassmannian tangent
space $\Hom(K_0,S_D/K_0)$ therefore allows independent variations
of both jets in every frame column.  Passing from these jets to
the displayed remainders is an invertible triangular operation,
with leading coefficient $F_0/x^2(0)\ne0$.  Their differentials
are consequently independent.

Set
\[
       t=x+a/2,\quad \Delta=a^2/4-b,\quad
       w_i=a u_i/2-v_i.
\]
The incidence of a common root near zero is now exactly
\[
                       \Delta=t^2,\qquad w_i=u_i t.
\]
There are no other normalization points above $K_0$, because its
exact common divisor is $\ell^2$.  Consequently the image of this
completed incidence is the completed image defining $Y_1$.
The ambient dimension is $N_0=r(D+1-r)$.  Of the $2r$ parameters
above, $a$ is free, as are the other $N_0-2r$ parameters.  This
accounts for $q=N_0-2r+1$ smooth variables in
\eqref{eq:binary-completion-v152}.

It remains to justify elimination scheme-theoretically.  The
relations in \eqref{eq:pinch-equations-v152} reduce any series to
$p+\sum_i w_i p_i$ with $p,p_i\in R_m$.  If its image in $B_m$
is zero, the decomposition $B_m=R_m\oplus R_m t$ gives
$p=0$ and $\sum_i u_i p_i=0$.  The first syzygies of the regular
sequence $u_1,\ldots,u_m$ are generated by the Koszul syzygies.
They are precisely the relations $u_iw_j-u_jw_i$.  This proves
the presentation without taking a radical.  The same argument
works first over polynomial rings and then by exact completion.

The extension $A_m\subset B_m$ is finite, since $t^2=\Delta$,
and birational, since $t=w_1/u_1$ in the fraction field.  The
regular ring $B_m$ is thus the normalization.  For
$p+v t\in R_m\oplus I t$, multiplication by $t$ stays in
$A_m$ if and only if $p\in I$.  As $1,t$ generate $B_m$,
the conductor is $I\oplus I t$, proving the quotient map
\eqref{eq:conductor-cover-v152}.  Equivalently,
\begin{equation}\label{eq:pinch-square-v152}
 A_m=B_m\mathbin{\times}_{\C[[t]]}\C[[\Delta]],
\end{equation}
where the first map sets $u=0$ and the second sends $\Delta$ to
$t^2$.  This is a specific conductor square, not an assumption
that the normalization descends through its reduced fibres.

As an $R_m$-module, $A_m=R_m\oplus I$.  For $m=1$, the ideal
$I$ is free.  For $m>1$, the exact sequence
$0\to I\to R_m\to\C[[\Delta]]\to0$ gives
$\operatorname{depth}_{R_m}I=2$.  Hence $A_m$ has depth two
for every $m\geq1$.  The maximal ideal of $R_m$ generates an
ideal with maximal radical in $A_m$, so this is also its local
ring depth.  Adding the $q$ smooth parameters gives the assertion.
When $m=1$, the equation is $w^2-\Delta u^2=0$; its partial
derivatives generate $(w,\Delta u,u^2)$ in characteristic zero.
\end{proof}

\begin{corollary}[The split and ramified parts of the same model]
\label{cor:collision-strata-v152}
On the binary locus where the common divisor has degree exactly two,
there are two strata.  Over two distinct roots the formal local ring
has the split form \eqref{eq:split-ring-v152}, with
$t=q+1$ and $c=r-1$.  Over a double root it has
\eqref{eq:binary-completion-v152}.  The support of the singular
locus in the latter model is the conductor, and the normalization
fibre at a collision is $\Spec\C[t]/(t^2)$.
For arbitrary $m$, its singular subscheme is given without
radicalization by the $m$-minors of the Jacobian of the displayed
$m^2$ relations in \eqref{eq:pinch-equations-v152}.
\end{corollary}
\begin{proof}
The assertion at distinct roots is Theorem~\ref{thm:split-conductor-v152}.
In the collision model the normalization is an isomorphism wherever
some $u_i$ is invertible.  On $u=w=0$, a nonzero value of $\Delta$
gives two distinct points after an \'etale square-root extension,
while $\Delta=0$ gives the stated length-two fibre.  The exact
normality criterion proves the assertion about singular support.
There are $2m+1$ nonsmooth coordinates and dimension $m+1$;
the cotangent presentation therefore gives the $m$-minors as
claimed.  This formula specifies the whole singular scheme, even
where it is thicker than its conductor support.
\end{proof}

\begin{remark}[Descent to intrinsic first relations]
\label{rem:conductor-stack-v152}
On the maximal-lower-Hilbert-function stratum of
Corollary~\ref{cor:algebra-boundary-v151}, these statements compute
the completed smooth-atlas singularities of the intrinsic
first-relation algebra stack after tangent-identity rigidification.
The conductor is equivariant, and formation of its finite-module
annihilator and of the Fitting ideals commutes with flat base change.
Thus the calculations descend through the frame atlas.  This does
not identify these rings with completed local rings of a coarse
orbit space, where stabilizers can introduce further singularities.
\end{remark}

% BEGIN PRESERVED PART 46-divisor-intersections-v152.tex
\section{Intersections of divisor-degree images}
\label{sec:degree-intersections-v152}

Here the total relation degree $D$ is fixed.  Write $Y_a\subset
\Gr(r,S_D)$ for the image admitting a divisor of degree $a$;
empty Grassmannians are omitted.  In several variables these images
are not generally nested.  Their reduced intersections nevertheless
have a finite description by common parts of two marked divisors.

\begin{proposition}[The common-part formula]
\label{prop:degree-intersections-v152}
For $1\leq a,b<D$ and $0\leq j\leq\min(a,b)$, let $M_j$ be the
scheme-theoretic image of
\begin{equation}\label{eq:intersection-map-v152}
\begin{split}
 \Pj(S_j)\times\Pj(S_{a-j})\times\Pj(S_{b-j})
       \times\Gr(r,S_{D-a-b+j})&\longrightarrow\Gr(r,S_D),\\
                   (c,u,v,L)&\longmapsto cuvL.
\end{split}
\end{equation}
A negative last degree or insufficient dimension gives the empty
image, and $\Pj(S_0)$ is a point.  Then
\begin{equation}\label{eq:intersection-reduced-v152}
             (Y_a\cap Y_b)_{\mathrm{red}}
                         =\left(\bigcup_j M_j\right)_{\mathrm{red}}.
\end{equation}
On the locus of two marked divisors $f_a,f_b$ with exact common
part of degree $j$, the corresponding factorization is
$c=\gcd(f_a,f_b)$, $f_a=cu$, $f_b=cv$, with $\gcd(u,v)=1$.
In particular, at a relation $K=a_0^mL$ with $a_0$ a nonzero
linear form and $\gcd(L)=1$, the
set of positive divisor degrees is exactly $\{1,\ldots,m\}$.
For the boundary point in Theorem~\ref{thm:universal-boundary-v152},
$m=D-q_n$.
\end{proposition}
\begin{proof}
Each map in \eqref{eq:intersection-map-v152} is defined on a
projective scheme: multiplication by a nonzero form has constant
rank, so the subspace construction is a morphism.  Its image is
closed and admits the two divisors $cu$ and $cv$, proving one
inclusion.  Conversely, if $K$ has divisors $f_a,f_b$, unique
factorization gives their common part $c$ of some degree $j$.
Their least common multiple is $cuv$ and divides every member of
$K$.  Division yields $L$ of degree $D-a-b+j$, of the same rank.
This proves equality on geometric points.  Over $\C$, a reduced
closed subscheme of a finite-type scheme is determined by its closed
points, giving \eqref{eq:intersection-reduced-v152}.  The last
assertion follows because every homogeneous divisor of a power
of one linear form is a power of that form.
\end{proof}

The reduction in \eqref{eq:intersection-reduced-v152} is essential
to the statement of this global formula.  It is not a prescription
to reduce the individual $Y_a$, their normalization fibres, or their
singular subschemes.  The full intersection of the two formal
branches on the coprime stratum was computed over Artin bases in
Theorem~\ref{thm:split-conductor-v152};
Theorem~\ref{thm:binary-pinch-v152} computes the thickening when
those branches collide.  Thus the degree-poset description and
the local scheme calculations have distinct, explicit roles.

% BEGIN PRESERVED PART 45-pencil-orbit-boundary-v152.tex
\section{A common ramified boundary of pencil failure orbits}
\label{sec:orbit-boundary-v152}

The fixed tensor presentation of a pencil has exact common divisor
$\delta^{n-1}$.  Varying the pencil alone therefore stays in the
exact-divisor locus.  The intrinsic finite algebra permits a larger,
linear coordinate-change action: changing the basis of its entire cotangent
space.  We now describe a boundary point produced by such changes.
The generic members are genuine pencil failure algebras.  The central
member has a larger common divisor, and its normalization fibre is
nonreduced.  Thus the boundary of Section~\ref{sec:canonical-boundary-v151}
occurs inside the orbit closures of the motivating algebras themselves.

Let $n=\dim V=\dim U\geq3$, $E=V^*\otimes U$, and put
\[
 e=n^2,\quad N=\binom{n+1}{2},\quad r=\binom N2,\quad
 g=n(n-1),\quad h=3n-4,\quad D=g+h.
\]
For a pencil $R$, write $K_R=\delta^{n-1}C_R\subset\Sym^D E$
as in Lemma~\ref{lem:primitive-pencil}.  Define the projective orbit
closure
\[
 Z_R=\overline{\operatorname{GL}(E)\cdot K_R}
                  \ \subset\ \Gr(r,\Sym^D E).
\]
It is independent of a chosen basis of the cotangent space.  Since
$Y_g$ is closed and invariant, $Z_R\subset Y_g$.

The group in $Z_R$ is the full $\operatorname{GL}(E)$, not the
$\PGL(V)$ congruence group of the pencil. Fixing the standard flag,
scalar direction and complement specifies the displayed common point;
no choice-independent distinguished moduli point is asserted.

\begin{theorem}[A common point in full first-relation orbit closures]
\label{thm:universal-boundary-v152}
Set
\[
 q_n=\begin{cases}3,&n=3,\\4,&n\geq4,\end{cases}
 \qquad k_n=h-q_n.
\]
There is an explicitly constructed point $K_\infty$ belonging to
$Z_R$ for every quadratic pencil $R$, with
\begin{equation}\label{eq:boundary-gcd-v152}
 K_\infty=a^{g+k_n}L_n,\qquad
 L_n\subset\Sym^{q_n}E,\quad \rank L_n=r,\quad\gcd(L_n)=1,
\end{equation}
where $a\in E$ is a nonzero linear form.  The flag pencil
$R_*=\langle x_1^2,x_1x_2\rangle$ admits a finite-flat family of
first-relation algebras over $\A^1$, with special relation space
$K_\infty$ and every nonzero fibre isomorphic to its failure algebra.
Together with Theorem~\ref{thm:closed-pencil-v151}, this gives a
chain of two explicit finite-flat specializations from the failure
algebra of every $R$ to the same boundary algebra.

The fibre of $\nu_g$ over $K_\infty$ has a single geometric point,
marked by $a^g$, and is nonreduced.  Its tangent dimension is
\begin{equation}\label{eq:boundary-tangent-v152}
                 \binom{e+k_n-1}{k_n}-1.
\end{equation}
In particular $K_\infty$ lies in the singular locus of $Y_g$ and
outside the exact-degree-$g$ locus.  For $n=3$ this tangent dimension
is $44$.
\end{theorem}

The construction requires the exact osculating order of one elementary
flag embedding.  We include the calculation to distinguish the theorem
from a computationally suggested degeneration.

\begin{lemma}[The flag coefficient filtration]
\label{lem:flag-jets-v152}
Write a matrix of independent linear coordinates as
\[
                    T=aI_n+B,\qquad B_{11}=0.
\]
For $C_*=C_{R_*}$, the filtration by vanishing order at $T=I_n$
on the chart $T_{11}=1$ has last nonzero graded piece in degree $q_n$.
It has a nonzero degree-zero piece.  Hence an adapted basis
$c_1,\ldots,c_r$ has orders $0\leq d_j\leq q_n$, attaining both
endpoints, and
\begin{equation}\label{eq:flag-initial-v152}
 \lim_{\tau\to0}\tau^{-d_j}c_j(aI_n+\tau B)
                           =a^{h-d_j}p_j(B),
\end{equation}
where the displayed initial forms are linearly independent.
\end{lemma}
\begin{proof}
The complementary-minor identity for $\Sym^2T$ identifies the vector
of primitive coefficients, up to a fixed invertible coordinate
change, with
\begin{equation}\label{eq:flag-coefficients-v152}
 \delta(T)^3\bigl(v^2\wedge vw\bigr),\qquad
                 v=T^{-1}e_1,\quad w=T^{-1}e_2.
\end{equation}
Indeed $\det(\Sym^2T)=\delta^{n+1}$, the complementary minors have
order two, and division of $\delta I_{N-2}$ by $\delta^{n-1}$ leaves
$\delta^3$.  This identity is initially on $\delta\ne0$; both sides
are the same polynomial vector of degree $h$.

A standard affine flag chart has representatives
\[
       v=e_1+\sum_{i=2}^n z_i e_i,
       \qquad w=e_2+\sum_{i=3}^n y_i e_i.
\]
Every coordinate of $v^2\wedge vw$ is a polynomial of degree at most
four.  The map from $\operatorname{GL}_n$ to this flag chart is
smooth, and restricting to $T_{11}=1$ remains smooth: the scalar
subgroup is contained in each flag stabilizer and is transverse to
this chart.  Normalizing $v,w$ only multiplies the whole wedge by
a common unit.  Smooth pullback preserves the order of a nonzero
function, as does multiplication by a unit.  The coordinates in
\eqref{eq:flag-coefficients-v152} are linearly independent by Lemma~\ref{lem:pencil-irreducibility} and the nonzero
matrix-coefficient inclusion in \eqref{eq:general-coefficient-map}.  Thus no nonzero linear
combination can have order exceeding four.  For $n\geq4$, the
coordinate indexed by $e_3^2\wedge e_3e_4$ is
\[
                       z_3^2(z_3y_4-z_4y_3),
\]
which is nonzero and has order exactly four.

For $n=3$, put $z=z_2$, $w=z_3$, $t=y_3$.  In the ordered basis
$e_1^2,e_1e_2,e_1e_3,e_2^2,e_2e_3,e_3^2$, the two factors are
\[
 (1,2z,2w,z^2,2zw,w^2),\qquad
 (0,1,t,z,zt+w,wt).
\]
Their fifteen wedge coordinates have the same linear span as
\begin{equation}\label{eq:ternary-flag-basis-v152}
\begin{gathered}
 1,\ z,\ w,\ t,\ z^2,\ zw,\ w^2,\ zt,\ wt,\ z^2t,\ zwt,\ w^2t,\\
 z^2w-z^3t,\quad zw^2-z^2wt,\quad w^3-zw^2t.
\end{gathered}
\end{equation}
This is obtained by subtracting the two coordinates with a common
quadratic term and rescaling by $2$.  Their lowest homogeneous terms
are distinct monomials, all of degrees at most three, and some have
degree three.  Thus the third jet is injective on this coefficient
space and the second is not.  This proves $q_3=3$.  The coordinate
$e_1^2\wedge e_1e_2$ is $1$ in every dimension, proving the
order-zero assertion.  Finally, homogeneity of degree $h$ gives
\eqref{eq:flag-initial-v152}; an adapted basis is by definition a
basis on each associated-graded piece, so its initial forms are
independent.
\end{proof}

\begin{proof}[Proof of Theorem~\ref{thm:universal-boundary-v152}]
The linear substitution $T\mapsto aI_n+\tau B$ is invertible for
$\tau\ne0$.  Divide the adapted coefficient $c_j$ by its exact
power $\tau^{d_j}$ after substitution.  There are no negative powers:
all its terms have $B$-degree at least $d_j$.  The resulting columns
span a rank-$r$ subbundle $J_\tau\subset\Sym^h E\otimes\C[\tau]$.
Independence at zero follows from the lemma, and away from zero from
the invertible substitution.  The ideal of its maximal minors in $\C[\tau]$ has no common
closed zero, and is therefore the unit ideal. On each maximal-minor
open a corresponding square submatrix is invertible and splits the
image. These opens cover $\operatorname{Spec}\C[\tau]$, proving the
subbundle assertion over the whole line, including all scheme points.  Set
\[
                 K_\tau=\det(aI_n+\tau B)^{n-1}J_\tau.
\]
Multiplication by the nonzero homogeneous polynomial in every fibre
again has constant rank.  Consequently $K_\tau$ is a subbundle, and
\[
 \mathcal A_\tau=\Sym E\otimes\C[\tau]\big/
                           (K_\tau,\mathfrak m^{D+1})
\]
is finite locally free of rank $\binom{e+D}{D}-r$.
For $\tau\ne0$, its relation space is a linear transform of $K_{R_*}$.

At zero the marked factor becomes $a^g$, while
$J_0=\langle a^{h-d_j}p_j(B)\rangle$.  It contains a nonzero multiple
of $a^h$, so its gcd is a power of $a$.  It also contains a polynomial
with exact $a$-adic order $h-q_n$, and every column is divisible by
that power.  Therefore $\gcd(J_0)=a^{k_n}$, proving
\eqref{eq:boundary-gcd-v152}.  Notice that $0<k_n<g$ for every
$n\geq3$.

The closed-flag specialization of Theorem~\ref{thm:closed-pencil-v151}
and its failure-algebra lift imply
$K_{R_*}\in Z_R$.  A closed invariant subset containing a point
contains that point's orbit closure.  Hence the just-constructed
$K_\infty\in Z_{R_*}$ belongs to every $Z_R$.  This proves the
assertion about the two-stage specialization without replacing it
by an unproved assertion about one fixed one-parameter subgroup
acting on every initial pencil.

The sole degree-$g$ divisor of $a^{g+k_n}$ is $a^g$.  Its overlap
with the residual gcd is $a^{k_n}$.  Lemma~\ref{lem:fibre-equations-v151}
gives \eqref{eq:boundary-tangent-v152}; the fibre is supported at one
point and has positive tangent space.  The normalization criterion
then proves singularity.  With $n=3$, $e=9$ and $k_n=2$, the dimension
is $\binom{10}{2}-1=44$.
\end{proof}

\subsection{The complete normalization fibre at the limiting algebra}
The last theorem is not restricted to the tangent space of its fibre.
The full local algebra can be given uniformly.

\begin{proposition}[Pure-power divisor fibre]
\label{prop:pure-power-fibre-v152}
Let $E=\C a\oplus W$, $g\geq k\geq1$, and
$K=a^{g+k}L$ with $\gcd(L)=1$.  Introduce coefficient variables for
$q_i\in\Sym^i W$, $1\leq i\leq k$, and define forms $c_j$ recursively by
\begin{equation}\label{eq:reciprocal-v152}
 c_0=1,\qquad c_j=-\sum_{i=1}^k q_i c_{j-i},\qquad c_j=0\ (j<0).
\end{equation}
The entire fibre of $\nu_g$ over $K$, supported at its unique point,
is the Artin algebra
\begin{equation}\label{eq:pure-power-fibre-v152}
 \C[[\text{coefficients of }q_1,\ldots,q_k]]\big/
     (\text{coefficients of }c_{g+1},\ldots,c_{g+k}).
\end{equation}
Its tangent dimension is $s_k(E)-1$.
\end{proposition}
\begin{proof}
Choose $p\in L$ not divisible by $a$.  Over any local Artin base,
a divisor $f$ near $a^g$ is monic in $a$, and the quotient by $f$
is flat over the base.  The reductions of $f,p$ are coprime.
The regular-sequence argument used in
\eqref{eq:coprime-artin-v152} makes $p$ a nonzerodivisor modulo $f$.
Thus $f\mid a^{g+k}p$ is equivalent to $f\mid a^{g+k}$, functorially
on these bases.  Once this condition holds, it holds for every
member of $K$.

Polynomial division gives a unique monic residual factor
$H=a^k+q_1a^{k-1}+\cdots+q_k$.  The equation $fH=a^{g+k}$ determines
the coefficients of $f$ successively as $c_1,\ldots,c_g$ in
\eqref{eq:reciprocal-v152}.  The remaining $k$ product coefficients
vanish if and only if $c_{g+1},\ldots,c_{g+k}$ vanish.  This follows
successively by comparing the product with the reciprocal series
$(1+\sum q_i a^{-i})^{-1}$.  Hence
\eqref{eq:pure-power-fibre-v152} represents the full infinitesimal
fibre, without reduction.  Finiteness of $\nu_g$ makes it Artin;
there is no other geometric divisor, so completion loses no part
of the fibre.  Since $g\geq k$, these equations have no linear
terms in the coefficient variables.  The number of coefficient variables is
$\sum_{i=1}^k\dim\Sym^iW=\dim\Sym^kE-1$, as asserted.
\end{proof}

\begin{example}[The nine-generator limit]\label{ex:n3-boundary-v152}
For $n=3$, the construction can be written with four of the nine
linear coordinates, denoted $a,u,v,w$.  The other five coordinates
remain generators of the algebra.  Up to nonzero scalar changes
of the relation basis, $K_\infty=a^8L_3$, where
\begin{equation}\label{eq:n3-limit-v152}
\begin{split}
 L_3=\langle& a^3,\ a^2u,\ a^2v,\ a^2w,\ au^2,\ auv,\ av^2,
             \ auw,\ avw,\\
            &u^2v,\ u^2w,\ uv^2,\ uvw,\ v^3,\ v^2w\rangle.
\end{split}
\end{equation}
Here one may take $u=B_{21}$, $v=B_{31}$, $w=B_{32}$; the monomials
are the lowest terms in \eqref{eq:ternary-flag-basis-v152}.
The relation space has dimension $15$ and exact gcd $a^8$.
The selected divisor has degree $g=6$, and its fibre has $44$
coefficient variables: eight coefficients of a linear form $Q_1$
and thirty-six of a quadratic form $Q_2$ in $W$.
Its defining ideal consists of all coefficients of the two forms
\begin{align*}
 c_7&=-Q_1^7+6Q_1^5Q_2-10Q_1^3Q_2^2+4Q_1Q_2^3,\\
 c_8&=Q_1^8-7Q_1^6Q_2+15Q_1^4Q_2^2-10Q_1^2Q_2^3+Q_2^4.
\end{align*}
This is a full presentation of the normalization fibre.  Its embedding
dimension $44$ is not being identified with its length, nor is this
fibre algebra being identified with the completed local ring of $Y_g$.
\end{example}

\begin{remark}[Relation to pencil invariants]\label{rem:segre-boundary-v152}
The action defining $Z_R$ is the intrinsic $\operatorname{GL}(E)$
action on first relations, not merely the congruence action on $V$.
The flag pencil is singular: every member has rank at most two.
It is therefore outside the regular-pencil open set treated by
\cite[Theorem~1.1 and \S5]{FevolaMandelshtamSturmfels}, where Segre
symbols organize classification and degenerations.  The theorem
above does not purport to classify that poset or all orbits in $Z_R$.
Its conclusion is instead a uniform common boundary and an exact
ramification algebra for genuine finite failure algebras.  The
increase from $g$ to $g+k_n$ in the intrinsic gcd is precisely what
is invisible in a fixed-determinant pencil specialization.
\end{remark}

\section{Universal horizontal models of the Hilbert boundary}
\label{sec:horizontal-v166}
A boundary fibre and a limit of boundary fibres are different objects.
The first is obtained by imposing the coefficient equations; the second
also removes the torsion created by the chosen specialization. We begin
with a construction that organizes this distinction over a whole base.
The representability input is the classical projective Hilbert scheme,
not a new representability theorem. Its application below supplies a
universal property for the horizontal models of the pencil boundary.

\subsection{The flat horizontal functor}
Let $p:W\to B$ be a projective morphism of finite presentation, with
$B$ an integral scheme of finite type over $\C$. Fix a relatively ample
line bundle and a nonempty open $U\subset B$ on which $p$ is flat with
constant Hilbert polynomial $P$. The scheme $W$ need not be reduced,
equidimensional, or flat over the boundary of $U$.

A locally Noetherian $B$-scheme $T$ is called $U$-admissible if
$U_T=U\times_B T$ is schematically dense in $T$. Write
\begin{equation}\label{eq:horizontal-definition-v166}
 W_T^{\mathrm h}=\overline{W_{U_T}}^{\,W_T},\qquad
 \mathcal J_T=\ker\bigl(\OO_{W_T}\longrightarrow
                       j_*\OO_{W_{U_T}}\bigr).
\end{equation}
The bar denotes schematic closure, and $j$ is the indicated open
immersion. On these test schemes define $\mathfrak F_{W,U}(T)$ to be
the singleton if $W_T^{\mathrm h}\to T$ is flat, and the empty set
otherwise. Morphisms are $B$-morphisms between admissible test schemes.
Whenever an object exists, its pullback is the horizontal closure and
is again flat: a flat scheme over $T$ has schematically dense inverse
image of $U_T$. Thus this is a functor on the stated category. This
last assertion does not concern arbitrary nonflat pullback of $W$.

\begin{theorem}[Universal horizontal modification]
\label{thm:horizontal-universal-v166}
There are an integral projective birational $B$-scheme $F_{W,U}$ and a
closed subscheme
\[
 \mathcal W\subset W\times_B F_{W,U}
\]
with the following properties. The morphism $F_{W,U}\to B$ is an
isomorphism over $U$. The family $\mathcal W\to F_{W,U}$ is flat,
and is the schematic closure of $W_U$. For every $U$-admissible $T$,
composition with $F_{W,U}\to B$ gives a natural bijection
\begin{equation}\label{eq:horizontal-universal-property-v166}
 \operatorname{Hom}_B(T,F_{W,U})\ \simeq\ \mathfrak F_{W,U}(T).
\end{equation}
Under this bijection $\mathcal W_T=W_T^{\mathrm h}$.
Consequently $F_{W,U}$ is terminal among $U$-admissible modifications
on which the horizontal closure is flat. It is independent of the
chosen relatively ample line bundle and of replacing $U$ by a smaller
nonempty open on which $p$ is flat.
\end{theorem}
\begin{proof}
Let $H=\operatorname{Hilb}^{P}(W/B)$, which is projective over $B$
by the projective Hilbert-scheme theorem \cite{Grothendieck166}.
The entire fibre $W_u$ defines a section $s:U\to H$. Put
$F_{W,U}=\overline{s(U)}\subset H$, with its schematic closure
structure, and restrict the universal quotient on $H$ to obtain
$\mathcal W$. Since $U$ is integral, its schematic closure is integral.
The section is a closed immersion over $U$, so the restriction of
$F_{W,U}$ over $U$ is exactly $s(U)$. This proves projectivity,
birationality, and the assertion over $U$.

Flatness of the universal Hilbert quotient proves flatness of
$\mathcal W$. A flat module over an integral base injects into its
localization at the generic point. Hence the open $W_U$ is
schematically dense in $\mathcal W$, and $\mathcal W$ is its closure
inside $W\times_B F_{W,U}$.

Suppose first that $W_T^{\mathrm h}$ is flat. This closed subscheme
is of finite presentation and projective over $T$. Its Hilbert
polynomial is locally constant and equals $P$ on $U_T$. Every
nonempty open-and-closed piece of $T$ meets $U_T$, so the polynomial
is $P$ everywhere. Representability gives $h:T\to H$. On $U_T$ this
is the section $s$. The pullback of the ideal of $F_{W,U}$ therefore
vanishes on $U_T$ and hence on $T$, by schematic density. Thus $h$
factors through $F_{W,U}$.

Conversely, pull the universal flat quotient back by any
$h:T\to F_{W,U}$. Flat base change preserves schematic density of
the inverse image of $U_T$: this can be checked on affine charts by
tensoring the injection into the finite product of localizations
associated to an affine cover of $U_T$. The pulled-back quotient is
therefore precisely the closure in \eqref{eq:horizontal-definition-v166}.
This also proves uniqueness. In particular two possible lifts give
the same embedded quotient, not merely isomorphic abstract fibres.
All constructions commute with morphisms of admissible tests.

The same universal property identifies the constructions from two
polarizations. For a smaller open $U'$, its section is schematically
dense in $s(U)$; their closures in the Hilbert scheme coincide.
Alternatively apply the universal property in both directions: the
universal flat closure over the integral modification is already
determined by the generic fibre. The two resulting morphisms are
inverse because they agree over the dense open and the targets are
separated.
\end{proof}

The theorem modifies the base to flatten a \emph{horizontal closed
subscheme}. It does not assert that the full pullback $W\times_B F$
is flat. Nor is $F\to B$ the monomorphism representing flatness of
the full pulled-back sheaf in the usual flattening stratification;
compare \cite[Tag 05PS]{Stacks166}. This distinction accounts for
the vertical component in Theorem~\ref{thm:collision-family-v164}.

\subsection{Finite saturation and base torsion}
Compatibility with completion reduces to the following finite kernel calculation.
\begin{lemma}[Flat change of rings and finite saturation]
\label{lem:saturation-v166}
Let $A$ be Noetherian, let $I,J\subset A$ be ideals, and let $A\to A'$
be flat with $A'$ Noetherian. Choose $N$ such that
$(I:J^N)=(I:J^{N+1})$. Then
\begin{equation}\label{eq:saturation-flat-v166}
 (I:J^\infty)A'=(IA':(JA')^\infty)
              =(IA':(JA')^N).
\end{equation}
This applies to Noetherian completion. If the map is faithfully flat,
the least stabilization exponent is unchanged.
\end{lemma}
\begin{proof}
For a finite set $j_1,\ldots,j_m$ of generators of $J^n$, the ideal
quotient modulo $I$ is the kernel of
$A/I\to(A/I)^m$, $a\mapsto(j_1a,\ldots,j_ma)$. Tensoring this
kernel sequence with $A'$ proves
$(I:J^n)A'=(IA':(JA')^n)$ for every $n$. The increasing sequence of
ideal quotients stabilizes by Noetherianity. Equality of consecutive
terms implies equality thereafter, since colon by $J$ sends one
term to the next. The same equality at $N$ holds after flat change
of rings. Taking the stabilized terms proves the formula. Faithful
flatness detects equality and inequality of the finite terms,
which gives the last assertion. There is no assertion of a uniform
exponent for all ideals, or for all arcs in a coefficient space.
\end{proof}

\begin{proposition}[The horizontal quotient over a discrete valuation ring]
\label{prop:dvr-horizontal-v166}
Let $R$ be a discrete valuation ring with uniformizer $\tau$, residue
field $k$, and fraction field $K$. Let $A$ be a Noetherian $R$-algebra.
Set
\[
 T(A)=\{a\in A:\tau^n a=0\text{ for some }n\},\qquad A^{\mathrm h}=A/T(A).
\]
Then $T(A)$ is an ideal killed by some $\tau^N$, and $A^{\mathrm h}$
is the unique flat quotient of $A$ with generic fibre $A_K$.
For a presentation $A=D/I$ over an $R$-algebra $D$ with no
$\tau$-torsion,
\begin{equation}\label{eq:dvr-saturation-v166}
 A^{\mathrm h}=D/(I:\tau^\infty),\qquad
 T(A)=(I:\tau^\infty)/I.
\end{equation}
Moreover
\begin{equation}\label{eq:all-tor-v166}
 \operatorname{Tor}_1^R(A,k)=\operatorname{Ann}_A(\tau),\qquad
 \operatorname{Tor}_i^R(A,k)=0\quad(i\geq2).
\end{equation}
Writing $T_n=\operatorname{Ann}_A(\tau^n)$ gives
$T_n/T_{n-1}\simeq\operatorname{Ann}_A(\tau)\cap\tau^{n-1}A$.
These statements glue for schemes of finite type over $R$.
\end{proposition}
\begin{proof}
The annihilator ideals stabilize in $A$. Their union is an ideal,
so stabilization supplies $N$. Multiplication by $\tau$ on
$A/T(A)$ is injective. A torsion-free module over a valuation ring
is flat \cite[Tag 0539]{Stacks164}; this proves flatness. If $A/J$
is another flat quotient with the same generic fibre, its
$\tau$-torsion is zero, so $T(A)\subset J$. The submodule
$J/T(A)$ of the torsion-free module $A^{\mathrm h}$ is generically
zero, hence zero. This proves uniqueness and the presentation.

Resolve $k$ by $0\to R\xrightarrow{\tau}R\to k\to0$ to obtain
all of \eqref{eq:all-tor-v166}. Multiplication by $\tau^{n-1}$ on
$T_n$ has kernel $T_{n-1}$ and image exactly the stated intersection.
Localization commutes with this finite filtration, so the formulas
glue. Flat change of discrete valuation rings commutes with the
horizontal quotient by Lemma~\ref{lem:saturation-v166}. In general
$T(A)$ is larger than $\operatorname{Ann}_A(\tau)$.
\end{proof}

For polynomial coefficient arcs the saturation is effective: in an
auxiliary variable $y$, it is the elimination ideal
\[
 (I:\tau^\infty)=(I,1-y\tau)D[y]\cap D.
\]
This gives a finite elimination procedure when the coefficients are
provided by polynomial or localized polynomial data. For an arbitrary
formal power series arc the formula remains exact, but it is not an
algorithm requiring only a fixed finite jet of unspecified series.

\subsection{A stratumwise compactification for every pencil specialization}
Let $G$ be the coefficient space of pencils under consideration and
let $\widetilde\Gamma\to G$ be the normalization of its integral
projective Hilbert graph. More generally the next assertion applies
to any fixed projective graph model of finite type. Normalization is
performed on the total graph before any specialization; it is finite
over a complex finite-type scheme, hence the map remains projective.
For every integral closed $B\subset G$ put
$W_B=\widetilde\Gamma\times_G B$, \emph{without} reducing this
pullback, and choose its dense generic-flat open $U_B$.

\begin{theorem}[Stratumwise universal boundary models]
\label{thm:stratum-flattening-v166}
The modification $F_{W_B,U_B}\to B$ and its universal family give a
canonical horizontal model on every such $B$. Every discrete
valuation ring arc $\gamma:\operatorname{Spec}R\to B$ whose generic
point lies in $U_B$ has a unique lift to this modification, and
\begin{equation}\label{eq:arc-universal-v166}
 \gamma^*\mathcal W=(\widetilde\Gamma\times_G\operatorname{Spec}R)^{\mathrm h}.
\end{equation}
There exists a finite stratification of $G$ by integral locally
closed strata $S_i$, with projective modifications over their reduced
closures $\overline S_i$, such that every discrete valuation ring
arc in $G$ is covered by one of these formulas, determined by the
stratum containing its generic point.
\end{theorem}
\begin{proof}
The first construction is Theorem~\ref{thm:horizontal-universal-v166}.
An arc whose generic point belongs to $U_B$ is admissible. Its
horizontal quotient is flat by
Proposition~\ref{prop:dvr-horizontal-v166}, so the universal property
gives the unique lift and the equality of embedded quotients.
Equivalently, extend the generic lift by properness and apply
uniqueness of a flat quotient with that generic fibre.

To obtain finitely many strata, take dense generic-flat opens in
each of the finitely many irreducible components of the reduced
base, shrinking them to be disjoint and to have constant Hilbert
polynomial. Apply the same procedure to their closed complement.
Its dimension strictly decreases, so the procedure terminates. Each
stratum is a dense open in its integral reduced closure after the
previous closed pieces are removed. Apply the first construction
to that closure and its generic-flat open. If the generic point of
an arc lies in $S_i$, the whole arc factors through $\overline S_i$:
the pullback of its defining ideal vanishes generically in the
domain $R$ and hence vanishes in $R$. The preceding argument applies.
The stratification is not claimed unique; the horizontal modification
attached to a specified integral closure is canonical.
\end{proof}

This supplies universal horizontal models for arbitrary contact order
where a projective Hilbert graph has been defined. It does not use
an a priori classification of its components. In particular it does
not identify higher-corank Hilbert fibres from Kronecker invariants.
The complete component formulas proved in
Theorems~\ref{thm:complete-fibre-v163} and
\ref{thm:product-fibres-v163} retain their hypotheses. A single
modification of the full base cannot be substituted without argument
for the stratumwise assertion: the generic boundary fibre on a
multiple-incidence stratum has a different Hilbert polynomial from
the point fibre on the nonincident open.

For the coefficient spaces $B_a$, this theorem turns an arcwise
saturation into the pullback of a projective universal family. For
$B_2$, Section~\ref{sec:all-arcs-v166} gives finite equations for \emph{every} arc
through the first nonreduced contact. Thus the universal property is
accompanied by an explicit specialization calculation, rather than
being used as a substitute for it.

\section{All arcs through a double contact}
\label{sec:all-arcs-v166}
The six coefficient directions of $B_2$ need not vanish in pairs,
and a collision need not have a chosen square root. In this section
we give a finite presentation of the full pullback, its horizontal
closure, and its vertical torsion for every discrete valuation ring
arc through $o=(x^2,0,0)$. The presentation is by a fixed collection
of coordinate systems with their recovery opens. It is not a list
of irreducible components indexed by the valuation of a single
chosen coefficient.

Let $R$ be a discrete valuation ring which is a $\C$-algebra, with
uniformizer $\tau$ and residue field $\C$. The same statements hold
after extending that residue field. Write an arbitrary arc centred
at $o$ uniquely as
\begin{equation}\label{eq:all-arc-coefficients-v166}
 f=(x-c)^2+d,\qquad g=u(x-c)+v,\qquad
 r=w(x-c)+z_0,
\end{equation}
where $c,d,u,v,w,z_0\in\tau R$. No relation between these six
elements is assumed. In particular neither $g=r=0$ nor a prescribed
ramification order is imposed. Set
\[
 H_R=\widetilde\Gamma_2\times_{B_2}\operatorname{Spec}R,
 \qquad Y_R=\overline{(H_R)_K}^{\,H_R},
\]
where $K=\operatorname{Frac}R$ and total-graph normalization precedes
the base change.

\subsection{Five conic coordinates and two division coordinates}
The vector space of conic equations through $P=[1:0:0]$ has basis
$FG,FR,G^2,GR,R^2$. Evaluation at the following five vectors gives
linearly independent functionals:
\begin{equation}\label{eq:five-vectors-v166}
 (0,0,1),\quad(0,1,0),\quad(0,1,1),\quad(1,1,0),\quad(1,0,1).
\end{equation}
Here $R$ in a target coordinate denotes the third homogeneous
coordinate, not the discrete valuation ring. To avoid confusion
below it will only occur in homogeneous equations.
Choose matrices $M_i$ whose first column is $(1,0,0)$ and whose
third columns are the vectors in \eqref{eq:five-vectors-v166}:
\begin{equation}\label{eq:five-matrices-v166}
\begin{gathered}
 M_0=I_3,\qquad
 M_1=\begin{pmatrix}1&0&0\\0&0&1\\0&1&0\end{pmatrix},\qquad
 M_2=\begin{pmatrix}1&0&0\\0&1&1\\0&0&1\end{pmatrix},\\
 M_3=\begin{pmatrix}1&0&1\\0&0&1\\0&1&0\end{pmatrix},\qquad
 M_4=\begin{pmatrix}1&0&1\\0&1&0\\0&0&1\end{pmatrix}.
\end{gathered}
\end{equation}
For coordinate system $i$ replace the coefficient column by
$M_i^{-1}(f,g,r)^{\mathsf t}$, and recompute the six coefficients
in \eqref{eq:all-arc-coefficients-v166}. The first polynomial stays
monic of degree two and the other two have degree at most one.
The central point remains $o$. These are changes of target
coordinates, always undone when the embedded curve is returned to
the fixed target; they are not quotient identifications of Hilbert
points.

For any one of these transformed six-tuples put
\begin{equation}\label{eq:all-arc-HB-v166}
\begin{split}
 D_i&=R[\lambda,e,A,B,C],\\
 K_i&=(e^2C-d,\ eB-u,\ e^2A+v,\
                 \lambda u-eA-w,\ \lambda v-e^2BC-z_0)\subset D_i.
\end{split}
\end{equation}
There are also two division coordinate systems. The first has
\begin{equation}\label{eq:all-arc-division-v166}
 D_+=R[\alpha,\beta],\qquad
 K_+=(\alpha u+\beta v-w,\ \alpha v-\beta ud-z_0).
\end{equation}
The second is obtained by interchanging the pairs $(u,v)$ and
$(w,z_0)$ and is denoted $(D_-,K_-)$.

\begin{theorem}[Finite-chart miniversal specialization]
\label{thm:all-arcs-v166}
For every arc \eqref{eq:all-arc-coefficients-v166}, the seven
coordinate systems \eqref{eq:all-arc-HB-v166} and
\eqref{eq:all-arc-division-v166}, restricted to their Hilbert recovery
opens, cover $H_R$. On each such open its ring is the corresponding
$D_\bullet/K_\bullet$, its horizontal ring is
\begin{equation}\label{eq:all-arc-horizontal-v166}
 D_\bullet/(K_\bullet:\tau^\infty),
\end{equation}
and its vertical torsion ideal is
$(K_\bullet:\tau^\infty)/K_\bullet$. These formulas retain the
scheme structures, including all embedded primes and nilpotents.
They commute with localization, completion, and flat extensions of
discrete valuation rings. They determine the entire specialization,
not merely a neighbourhood of one chosen conic or primitive tail.

The recovery opens are the opens of the total coefficient graph,
not extra conditions on the six coefficients of the arc. Thus the
theorem permits nonzero $g$ and $r$ jets, arbitrary relative orders
of vanishing, unequal residual directions, and arbitrary ramification.
It includes arcs entirely contained in an incidence stratum.
\end{theorem}
\begin{proof}
In the Hilbert--Burch coordinates of
Lemma~\ref{lem:HB-family-v163}, put $z=x-c$ and $H=R-\lambda G$.
The base map reads
\[
 f=z^2+e^2C,\quad g=eBz-e^2A,\quad
 r=\lambda g-eAz-e^2BC.
\]
Equating its coefficients with those of the retained arc gives
exactly the five equations in \eqref{eq:all-arc-HB-v166}; the sixth
coefficient eliminates $c$. In particular the constant equation
$\lambda v-e^2BC=z_0$ must not be discarded for a general arc.
No radical has been taken. The total chart is regular on its
recovery open, so the same computation applies to the normalization
of the total graph.

In the primitive chart write the recovered direction as
$\alpha+\beta z$. Division by $f=z^2+d$ gives
\[
 (\alpha+\beta z)(uz+v)-\beta u(z^2+d)
       =(\alpha u+\beta v)z+\alpha v-\beta ud.
\]
The condition that this equal $r$ proves
\eqref{eq:all-arc-division-v166}. Interchanging $g$ and $r$ gives
the other primitive chart. These are the degree-two division charts
of Theorem~\ref{thm:division-chart-v162}, after the invertible change
between $x$ and $z$ coordinates. Their equations also retain the
coefficient projection.

We specify the opens used here. On a conic chart the restriction maps
in degrees $(0,2)$ and $(1,1)$ have ranks five and four. Restrict to
the opens defined by these ranks, by the nonzero selected conic
coefficient, and by the $tG,tH$ pivot for the two bilinear equations.
The normalized inverse is precisely the recovery map in
Lemma~\ref{lem:HB-family-v163}; the matrices in
\eqref{eq:five-matrices-v166} transport these opens. They contain
every central conic on which the corresponding evaluation functional
is nonzero. On a division chart use the split-injection open in
Lemma~\ref{lem:recovery-details-v163}. These opens contain every
central primitive tail for which the selected homogeneous direction
coordinate is a unit. Localizing the displayed rings to these opens
is understood; no global affine-isomorphism claim outside the
recovery domains is needed.

A nonzero conic through $P$ cannot vanish on all five vectors in
\eqref{eq:five-vectors-v166}. Indeed their evaluations recover,
in order, the coefficients of $R^2$, $G^2$, $GR$, $FG$, and $FR$.
Thus the five conic coordinate systems cover the complete conic
component. A primitive jet over $\C[x]/(x^2)$ is represented by a
pair with at least one unit, and therefore belongs to one of the two
division coordinate systems. The complete-fibre theorem
\ref{thm:complete-fibre-v163} shows that these are all central
points, including the doubled-line intersection. The same charts
compute their scheme structures, not only their closed points.

Their union is an open neighbourhood of the entire special fibre
in $H_R$. Its complement is closed in the proper $R$-scheme $H_R$.
A nonempty closed subset of a proper scheme over a local discrete
valuation ring has image meeting its closed point. Since the
complement misses that fibre, it is empty. This proves exhaustion
of $H_R$, including its generic fibre. Now apply
Proposition~\ref{prop:dvr-horizontal-v166} on the seven coordinate
domains. Its uniqueness assertion glues the horizontal quotients.
Lemma~\ref{lem:saturation-v166} proves compatibility on overlaps
and with the stated flat changes of rings. The torsion formula
follows from the same proposition.
\end{proof}

\begin{corollary}[Universal interpretation of the seven presentations]
\label{cor:all-arcs-universal-v166}
Choose the stratum containing the generic point of the arc in
Theorem~\ref{thm:all-arcs-v166}. Its seven horizontal presentations
are the pullback of the universal family on the corresponding
modification in Theorem~\ref{thm:stratum-flattening-v166}. Hence
specialization does not require an additional choice of a component,
a normalization of the special fibre, or a stable-map replacement.
\end{corollary}
\begin{proof}
Both constructions are flat embedded quotients with the original
generic fibre. Uniqueness in
Proposition~\ref{prop:dvr-horizontal-v166} identifies them on every
coordinate domain and hence globally.
\end{proof}

The adjective miniversal here refers to the full six-dimensional
coefficient base and its framed fixed-target comparison. It does
not assert a finite list of valuation chambers, nor a classification
of the full fibre for $a\geq3$. The theorem gives the exact ideals
and their saturation for every arc through the double contact. Its
specialization to $g=r=0$ recovers the previously computed collision;
it is not restricted to that subbase.

\subsection{Which conics are horizontal?}
The vertical component in the symmetric collision has a concrete
modular explanation: it records limits approached in coefficient
directions other than the double-incidence subbase. Let
\[
 \mathcal Q_P=\Pj\langle FG,FR,G^2,GR,R^2\rangle,\qquad
 S=V(q_{FG},q_{FR})\subset\mathcal Q_P.
\]
Thus $S\simeq\Pj^2$ is the linear space of conics singular at $P$.
Its ideal is generated by the two nontrivial first derivatives at
$P$; this is an intrinsic first-jet condition on the conic equation.

\begin{theorem}[Arc realization and the meaning of vertical excess]
\label{thm:conic-realization-v166}
Every conic $Q\in\mathcal Q_P$, including a doubled line, occurs as
the tail of the special Hilbert curve of an arc in $B_2$ whose
generic point is nonincident. Among arcs in the double-incidence
subbase $g=r=0$ with $f$ generically squarefree, the conic tails
reached horizontally are exactly $S$.
\end{theorem}
\begin{proof}
Choose one of the five coordinates so that $Q$ is
\[
 (R-\lambda_0G)^2+B_0F(R-\lambda_0G)+A_0FG+C_0G^2=0.
\]
Take $e=\tau$, $c=0$, $\lambda=\lambda_0$, and choose polynomial
series $A(\tau),B(\tau),C(\tau)$ with prescribed constant terms
$A_0,B_0,C_0$ and
$A(\tau)^2+C(\tau)B(\tau)^2\ne0$ in $K$. Such a choice always
exists because this polynomial is not identically zero on the
three-dimensional coefficient space; a generic line through the
specified point is not contained in its zero locus. The
Hilbert--Burch base map produces an arc centred at $o$. On its
generic point $e(A^2+CB^2)\ne0$, so the three quadratics span the
space of binary quadratics and define the nonincident graph. The
flat determinantal family has special curve $C_0\cup Q$ by
Lemma~\ref{lem:HB-family-v163}. Undoing the constant coordinate
change gives the required embedded arc in the original target.

For the second assertion use any conic coordinate system for a
horizontal limit; the transformations in
\eqref{eq:five-matrices-v166} preserve the condition $g=r=0$.
On the generic fibre the coefficient equations give $eB=eA=0$.
Also $e\ne0$ there, since $e^2C=d$ and the centred quadratic is
generically squarefree, so $d\ne0$. Thus $A=B=0$ on the horizontal
closure. Its conic limit lies in $S$. Conversely, in a chart for
any $Q\in S$, set $A=B=0$, $e=\tau$, and take $C(\tau)$ with the
specified constant term and nonzero generic value. Then
$f=x^2+\tau^2C(\tau)$ is generically squarefree and $g=r=0$.
The displayed flat Hilbert family gives a point of the horizontal
closure: it agrees with that closure generically and factors through
its saturated ideal over the torsion-free ring $R$. Its conic limit
is $Q$. This includes doubled lines by taking $C(0)=0$ and
$C(\tau)\ne0$.
\end{proof}

The first assertion concerns closures of individual generic Hilbert
points on generically nonincident arcs. The second concerns points
of the horizontal family of complete incidence fibres. The two
notions agree on their common incidence interpretation but should
not be conflated. In particular the vertical $\Pj^4$ is not a
collection of spurious curves: its complement of $S$ consists of
valid limits that require different coefficient directions. No
variation of a stability parameter is involved.

\section{The extension class and ramified collision models}
\label{sec:extension-ramification-v166}
The nilradical as a module does not determine a nonreduced scheme.
We compute the missing algebra extension, and then determine how
arbitrary ramification changes the horizontal normalization and the
vertical torsion of the collision family. The two calculations
separate intrinsic nonreducedness of a fibre from torsion introduced
by a nonflat pullback.

\subsection{The square-zero algebra is not locally split along the doubled lines}
Write $X=X_2$, $X_{\mathrm{red}}=\mathcal Q_P\cup\mathbb F_2$,
$S\simeq\Pj^2\subset\mathcal Q_P$ for the singular-conic plane,
and $D\simeq\Pj^1$ for the doubled-line locus. The closed immersions
are denoted $j:S\hookrightarrow X_{\mathrm{red}}$ and
$i:D\hookrightarrow X_{\mathrm{red}}$. By
Theorem~\ref{thm:nilradical-sheaf-v164}, the nilradical is
$N=j_*\OO_S(-1)$. The square-zero extension of sheaves of algebras
\begin{equation}\label{eq:extension-v166}
 0\longrightarrow N\longrightarrow\OO_X
   \longrightarrow\OO_{X_{\mathrm{red}}}\longrightarrow0
\end{equation}
defines a class
$\xi\in\operatorname{Ext}^1_{X_{\mathrm{red}}}
(L_{X_{\mathrm{red}}/\C},N)$.
Here $L$ is the cotangent complex; the degree-one group can equally
be computed using the naive cotangent complex. It is not the
ordinary extension group of the sheaf of differentials.

\begin{theorem}[The nonsplitting locus and its obstruction sheaf]
\label{thm:extension-class-v166}
The algebra extension \eqref{eq:extension-v166} splits locally on
$X_{\mathrm{red}}\setminus D$ and does not split locally at any
point of $D$. Moreover
\begin{equation}\label{eq:extension-sheaf-v166}
 \mathcal E xt^1_{X_{\mathrm{red}}}
       (L_{X_{\mathrm{red}}/\C},N)\simeq i_*\OO_D,
\end{equation}
where the image of $\xi$ is the section $1$. In a conic chart the
class is represented by the conormal homomorphism
\begin{equation}\label{eq:extension-map-v166}
 eA\longmapsto0,\qquad eB\longmapsto0,\qquad eC\longmapsto eC.
\end{equation}
In particular the nilpotent algebra is not the trivial square-zero
extension of its reduced support, even Zariski locally along $D$.
\end{theorem}
\begin{proof}
Work first in the polynomial chart
\[
 P_0=\C[\lambda,e,A,B,C],\quad
 I=(eA,eB,e^2C),\quad J=(eA,eB,eC),\quad R_0=P_0/J.
\]
Then $\OO_X=P_0/I$ and $N=J/I$ is generated by $n=eC$ with
annihilator $(e,A,B)$. Hence, as a module,
$N=\C[\lambda,C]n$. Its extension class is the homomorphism
$J/J^2\to N$ induced by lifting the polynomial generators to
$P_0/I$, namely \eqref{eq:extension-map-v166}. This is the
presentation description of square-zero extensions
\cite[Section 91.2, especially Lemma 91.2.3 and Remark 91.2.4]{Stacks166}.

We compute the whole local degree-one group. The module of relations
on $eA,eB,eC$ over $P_0$ is generated by the three Koszul relations
on $A,B,C$, since multiplication by $e$ is injective in $P_0$.
For a homomorphism to $N$, the relations
$C(eA)-A(eC)=0$ and $C(eB)-B(eC)=0$ force the images of $eA$ and
$eB$ to vanish: $A$ and $B$ annihilate $N$, whereas multiplication
by $C$ on $N$ is injective. The image of $eC$ is unrestricted.
Consequently $\operatorname{Hom}_{R_0}(J/J^2,N)=N$.

The image of a derivation from the polynomial presentation is
$\delta(eC)=C\delta(e)+e\delta(C)=C\delta(e)$ in $N$; its values
on $eA,eB$ are zero. All multiples of $C$ occur. Therefore
\begin{equation}\label{eq:local-ext-v166}
 \operatorname{Ext}^1_{R_0}(L_{R_0/\C},N)
  =\frac{\operatorname{Hom}_{R_0}(J/J^2,N)}
         {\operatorname{im}\operatorname{Hom}_{R_0}
               (\Omega_{P_0/\C}\otimes R_0,N)}
  =N/CN\simeq\C[\lambda].
\end{equation}
The class of the given algebra maps to $n\bmod CN$, hence to $1$.
Only the truncation in degrees $-1,0$ is used in this computation;
we have not asserted that the full cotangent complex has length
two, or that $J$ is a complete intersection.

The annihilator of the last module is the ideal of
$D=V(e,A,B,C)$ in this chart. At every point of $D$ the class is a
nonzero generator, so an algebra section is impossible. Where
$C$ is invertible, $R_0$ is the polynomial algebra
$\C[\lambda,A,B,C,C^{-1}]$ and its variables give an explicit
section into $(P_0/I)[C^{-1}]$. Outside $S$, the module $N$ is zero.
These opens, together with the primitive charts away from $D$,
prove the asserted local splitting elsewhere.

The complete-fibre charts cover $X$. Their extension classes are
restrictions of the same global algebra extension. Thus their
local generators glue to a nowhere vanishing section of the
obstruction sheaf supported on $D$. The map $\OO_D$ sending $1$
to this section is an isomorphism on every displayed chart, and is
zero off $D$ on both sides. This proves
\eqref{eq:extension-sheaf-v166}. We identify the image of the global
class in the local obstruction sheaf; no assertion that the entire
global hyperextension group is one-dimensional is required.
\end{proof}

\begin{corollary}[Associated subvarieties of the contact fibre]
\label{cor:associated-v166}
The associated subvarieties of $X_2$ are its two reduced components
$\mathcal Q_P$ and $\mathbb F_2$, and the embedded subvariety $S$.
The embedded associated prime disappears precisely off $S$.
The support of the nilradical is $S$, whereas the locus where its
algebra extension fails to split locally is the strictly smaller
curve $D$.
\end{corollary}
\begin{proof}
The primary decomposition
\[
 (eA,eB,e^2C)=(e)\cap(A,B,C)\cap(A,B,e^2)
\]
gives the associated primes $(e)$, $(A,B,C)$, and $(e,A,B)$ on
conic charts. The last is the generic point of $S$; the first two
are those of the two components. The primitive charts away from
their boundary are reduced and smooth. Since these charts cover
all of $X_2$, there are no additional associated subvarieties.
The final distinction follows from
Theorem~\ref{thm:extension-class-v166} and $N=j_*\OO_S(-1)$.
\end{proof}

\subsection{Arbitrary ramification of the ordered collision}
For $m\geq1$ replace the coefficient parameter in
Theorem~\ref{thm:collision-family-v164} by $\delta=\tau^m$.
Let $H_m$ be the full pullback and $Y_m$ its horizontal closure.
The following statements are over $\C[\tau]$; they also hold after
localization or completion at $\tau=0$.

\begin{theorem}[Ramified torsion and toric normalization]
\label{thm:ramified-v166}
The horizontal space is the flat base change of the previous
horizontal model. Its conic coordinate ring is
\begin{equation}\label{eq:ramified-horizontal-v166}
 \C[\lambda,e,C,\tau]/(e^2C+\tau^m).
\end{equation}
The normalization of this chart is
\begin{equation}\label{eq:ramified-invariants-v166}
 \C[\lambda,s,r]^{\mu_m},\qquad
 \zeta(s,r)=(\zeta s,\zeta^{-2}r),\qquad
 e=s^m,\quad C=-r^m,\quad\tau=s^2r.
\end{equation}
The two central divisors in this normalization have multiplicities
$2/\gcd(m,2)$ and $1$, respectively. If $m=2q$, the normalized
chart can equivalently be written
\begin{equation}\label{eq:even-ramification-v166}
 \C[\lambda,e,h,\tau]/(eh-\tau^q),\qquad C=-h^2.
\end{equation}
It is smooth for $q=1$; for $q>1$ its transverse surface singularity
has equation $eh=\tau^q$.

The full pullback has the scheme-theoretic union sequence
\begin{equation}\label{eq:ramified-union-v166}
 0\to\OO_{H_m}\to\OO_{Y_m}\oplus
       \OO_{\mathcal Q_P}\otimes_\C\C[\tau]/(\tau^m)
   \to\OO_S\otimes_\C\C[\tau]/(\tau^m)\to0.
\end{equation}
Its base torsion, with closed-immersion direct images understood, is
\begin{equation}\label{eq:ramified-torsion-v166}
 T_m=\mathcal I_{S/\mathcal Q_P}\otimes_\C\C[\tau]/(\tau^m).
\end{equation}
It has exact torsion order $m$. In particular
$\operatorname{Tor}_1^{\C[\tau]}(\OO_{H_m},\C)=\tau^{m-1}T_m$,
and all higher Tor sheaves vanish. For $m>1$ the full pullback is
not reduced along the generic point of its vertical component.
\end{theorem}
\begin{proof}
The map $\C[\delta]\to\C[\tau]$, $\delta\mapsto\tau^m$, is
finite free. The union sequence and flat horizontal quotient in
Theorem~\ref{thm:collision-family-v164} therefore survive this base
change. Its conic chart is
\[
 \C[\lambda,e,A,B,C,\tau]/(eA,eB,e^2C+\tau^m).
\]
The horizontal ring is \eqref{eq:ramified-horizontal-v166}, either
by flat saturation or by uniqueness of the flat quotient with the
same generic fibre. Pulling back the original torsion ideal proves
\eqref{eq:ramified-torsion-v166}. Multiplication by $\tau^{m-1}$
is nonzero on this sheaf and multiplication by $\tau^m$ is zero.
Its annihilator by $\tau$ is its last nonzero layer. The two-term
resolution of the residue field proves the Tor statements. At the
generic vertical point $e=0$ and one of $A,B$ is a unit, leaving
the nonreduced equation $\tau^m=0$ when $m>1$.

For normalization, omit the independent coordinate $\lambda$.
The ring in \eqref{eq:ramified-horizontal-v166} is a domain: its
relation is linear in $C$, with relatively prime coefficients
$e^2$ and $\tau^m$. Embed it in $\C[s,r]$ by the displayed
monomials. Its exponent lattice is generated by
$(m,0),(0,m),(2,1)$. This lattice has index $m$ in $\mathbb Z^2$
and is exactly
\[
 \{(a,b)\in\mathbb Z^2:a-2b\equiv0\pmod m\},
\]
the character lattice of the invariant field in
\eqref{eq:ramified-invariants-v166}. Thus the invariant ring and
the original ring have the same fraction field. The invariant
ring is normal: an element of its fraction field integral over
it is integral over the normal polynomial ring, belongs to that
ring, and is fixed by the group. It is finite over the original
ring, because $s$ and $r$ satisfy the monic equations
$s^m=e$ and $r^m=-C$, and the invariant submodule of this finite
module is finite. These facts prove that it is exactly the
normalization, not a cover with a larger fraction field.

The inverse image of $\tau=0$ is supported on $s=0$ and $r=0$.
At the generic point of $s=0$ the stabilizer has order
$\gcd(m,2)$, so the quotient has this ramification index there.
Upstairs $\tau=s^2r$ vanishes to order two. Downstairs its order
is $2/\gcd(m,2)$. At the generic point of $r=0$ the stabilizer
is trivial and the order is one. There are no other central
prime divisors.

If $m=2q$, first take invariants by the order-two subgroup, giving
$\C[s^2,r]$. The remaining action has order $q$ and opposite
weights on $s^2,r$. Its invariant ring is generated by
$e=s^{2q}$, $h=r^q$, and $\tau=s^2r$, with the single relation
$eh=\tau^q$. Every invariant monomial is a product of these
three, as is seen by subtracting the smaller of its two exponents.
This proves \eqref{eq:even-ramification-v166}. The derivative
criterion gives smoothness at $q=1$ and the stated surface
singularity for $q>1$.
\end{proof}

These local normalizations glue uniquely because they are integral
closures in the common function field of $Y_m$. On the primitive
chart the pullback is already $\C[\lambda,u,\tau]$, so no additional
normalization is needed. For $m=2$ the result recovers the ordered
normal-crossing parameter space. For larger even $m$ reducedness of
the normalized central divisor does not imply smoothness of the
total space. For odd $m$ the doubled divisor remains. None of these
operations normalizes the individual embedded curve or asserts its
stable reduction as a map.

\section{Universal horizontal families on the effective image}
\label{sec:effective-horizontal-v166}
The sharp inverse precedes the boundary construction. This order
also determines the correct category for the new universal models.
We use the equivalence
$\mathscr F_{\rm pen}^{\rm eff}\simeq[G/\operatorname{PGL}(V)]$
of Theorem~\ref{thm:pencil-stack-equivalence}, not the stack of all
finite algebras of the same length.

\begin{proposition}[Descent of the horizontal universal property]
\label{prop:effective-horizontal-v166}
Let $B\subset G$ be an integral invariant coefficient subspace,
and let $U\subset B$ be an invariant dense generic-flat open for
the pulled-back projective Hilbert graph. The universal horizontal
modification of Theorem~\ref{thm:horizontal-universal-v166} is
$\operatorname{PGL}(V)$-equivariant. Its quotient gives a
representable proper birational morphism over
$[B/\operatorname{PGL}(V)]$, equipped with the descended flat
horizontal family. Through the sharp effective inverse this is a
modification of the corresponding effective failure-family image.
It represents the same horizontal lifting property on admissible
tests, checked on source-frame torsors.

At a regular corank-two contact with smaller Smith exponent two,
the completed local presentations for an arbitrary coefficient arc
are those of Theorem~\ref{thm:all-arcs-v166}, with the smooth frame
and larger-exponent factors restored. The nilpotent algebra class
and ramified torsion are the classes computed in
Theorems~\ref{thm:extension-class-v166} and
\ref{thm:ramified-v166}. The first assertion is global on the
specified coefficient subspace; the contact assertion describes
its completed fixed-target local equations.
\end{proposition}
\begin{proof}
Apply a group element to the universal embedded flat quotient on
the Hilbert graph. The resulting quotient has the same required
horizontal restriction on the invariant open. The uniqueness in
Theorem~\ref{thm:horizontal-universal-v166}, applied over the group
times the integral modification, lifts the action uniquely.
Uniqueness also proves the identity and associativity axioms, so
this is an action, not only pointwise equivariance. The universal
family has the same equivariance. Taking a source-frame torsor
and descending this projective equivariant family gives the
representable proper morphism of quotient stacks. Birationality
holds on the dense open where the modification is an isomorphism.
Admissibility and flatness can be checked on that faithfully flat
torsor, and uniqueness of embedded quotients descends the universal
property.

For the local assertion use the completed coefficient normal form
and the fixed-target comparison of
Proposition~\ref{prop:embedded-functors-v163}. Lemma
\ref{lem:larger-smith-v166} restores the smooth complementary
coefficients, and \eqref{eq:frame-cocycle-v166} restores the
original incidence embedding. The comparison is on all Artin
quotients, so it retains the square-zero extension and not merely
its support or its vector-space dimension. Saturation commutes
with the flat completions by Lemma~\ref{lem:saturation-v166}.
The seven coefficient presentations and the extension and torsion
calculations therefore give the asserted completed equations.
Finally transport this construction through the already established
effective stack equivalence.
\end{proof}

This proposition is a consequence of reconstruction and of the
geometric universal property. It does not assert a boundary
operation internal to the original closed multiplication table
before reconstruction, nor does it include deformations leaving
the effective pencil image. The sharp inverse for all pencils,
including singular ones, retains its original hypotheses and proof;
no higher-corank Hilbert classification is inferred from it.

\section{A determinantal model for every contact order}
\label{sec:determinantal-v167}
The distinction between a curve in a Hilbert scheme and a fibre of
its parameter space is essential. This section constructs the former
Hilbert graph over the whole polynomial coefficient space. Later we
return to flat limits of the fibres of that graph. The two operations
have different Hilbert polynomials and different universal properties.

Fix an integer $a\geq1$ and set
\[
 B_a=\{(f,g,r):f\text{ monic of degree }a,\quad
                         \deg g,\deg r<a\}.
\]
Write $[s:t]$ for the source coordinates and $x=t/s$. Homogenization
to degree $a$ gives forms $F_0,F_1,F_2$. Let $U_a\subset B_a$ be
the open locus where these forms have no common zero. The graph
$C_b\subset X=\Pj^1\times\Pj^2$ of $[F_0:F_1:F_2]$ has, for the
Segre polarization, polynomial
\[
 P_a(l)=(a+1)l+1.
\]
The reduced closure of $(b,[C_b])$, $b\in U_a$, is denoted by
$\Gamma_a\subset B_a\times\Hilb^{P_a}(X)$, as in
Section~\ref{sec:Euclidean-boundary-v162}. It is an integral graph,
not a quotient by projective changes of the target.

Put
\begin{equation}\label{eq:hilbert-degree-v167}
 d=a+1,\qquad m=\frac{d(d-1)}2+1,\qquad
 q=dm+1,\qquad N=(m+1)\binom{m+2}{2}.
\end{equation}
For $0\leq i\leq m$ and $j_0+j_1+j_2=m$, let the corresponding
column of the $q\times N$ matrix $M_a$ be the coefficient vector of
\begin{equation}\label{eq:state-matrix-v167}
 s^{m-i}t^i F_0^{j_0}F_1^{j_1}F_2^{j_2}
                   \quad\text{in }\C[s,t]_{dm}.
\end{equation}
Thus its entries are explicit polynomials in the $3a$ coefficients.
Order rows by the exponent of $t$, and order columns
lexicographically in $(i,j_0,j_1,j_2)$. For each $q$-element column
set $J$, write $\Delta_J=\det(M_a)_J$; identically zero determinants
may be omitted. Let
\begin{equation}\label{eq:hilbert-ideal-v167}
                         \mathfrak a_a=(\Delta_J)_J\subset\OO(B_a).
\end{equation}
The size of this presentation is not a complexity bound. In
particular, enumerating every maximal minor is not the recommended
way to specialize one matrix.

\begin{theorem}[Determinantal Hilbert graph]\label{thm:determinantal-v167}
For every $a\geq1$ there is a canonical isomorphism over $B_a$
\begin{equation}\label{eq:rees-hilbert-v167}
 \Gamma_a\simeq\operatorname{Bl}_{\mathfrak a_a}(B_a)
       =\Proj_{B_a}\bigoplus_{n\geq0}\mathfrak a_a^n.
\end{equation}
Under this isomorphism the rational Pl\"ucker coordinates of the
Hilbert map are $[\Delta_J]_J$. The universal curve on this blow-up
is the flat embedded extension of the graph on $U_a$.
The normalization of $\Gamma_a$ is the normalization of this Rees
model. This assertion does not normalize any scheme fibre.

Let $R$ be a discrete valuation ring containing $\C$, with
uniformizer $\tau$, fraction field $K$, and residue field $k$.
For a coefficient arc $\gamma:\Spec R\to B_a$ with
$\gamma(\Spec K)\subset U_a$, set
\begin{equation}\label{eq:determinantal-order-v167}
 \mu(\gamma)=\min_J\operatorname{ord}_{\tau}\Delta_J(\gamma),
 \qquad
 \ell_J(\gamma)=\left(\tau^{-\mu(\gamma)}
                                \Delta_J(\gamma)\right)\bmod\tau.
\end{equation}
The special embedded curve is the unique Hilbert point whose
Pl\"ucker vector is $[\ell_J(\gamma)]_J$. In particular this vector
specifies its scheme structure, not merely its support or cycle.
Two such arcs centred at the same coefficient point have the same
point of $\Gamma_a$ precisely when these nonzero vectors are
proportional over $k$.
\end{theorem}

\begin{proof}
We first check the degree at which the Hilbert point is taken.
The Gotzmann representation of $dl+1$ has $d$ terms of degree one
and $(d-1)(d-2)/2$ terms of degree zero, since
\[
 \sum_{i=0}^{d-1}\binom{l+1-i}{1}
                       =dl+\frac{d(3-d)}2.
\]
Its total number of terms is $m$ in
\eqref{eq:hilbert-degree-v167}. Gotzmann regularity and the
Grassmannian construction of the Hilbert scheme therefore give a
closed immersion into the Grassmannian of rank-$q$ quotients in
degree $m$; see \cite{Gotzmann167,Grothendieck166}. We use the
ordinary projective Hilbert scheme after the Segre embedding
$X\hookrightarrow\Pj^5$, not an unproved multigraded regularity
bound. The restriction
\[
 H^0(\Pj^5,\OO(m))\longrightarrow H^0(X,\OO(m,m))
\]
is onto. All kernels defining subschemes of $X$ contain its fixed
kernel. Hence this Hilbert immersion factors as a closed immersion
into the Grassmannian of quotients of $H^0(X,\OO(m,m))$, of
dimension $N$.

For $b\in U_a$, evaluation on the graph is exactly
\[
 H^0(X,\OO(m,m))\longrightarrow
 H^0(C_b,\OO(m,m))=H^0(\Pj^1,\OO(dm)).
\]
This map is onto by the same regularity theorem. Its matrix is
$M_a(b)$ and its Pl\"ucker coordinates are its $q$-minors. This
proves that $\mathfrak a_a$ is nonzero and has no base point on
$U_a$.

For an integral affine scheme $\Spec A$ and a nonzero ideal
$(h_0,\ldots,h_n)$, the graph closure of $[h_0:\cdots:h_n]$ is
$\Proj A[(h_0,\ldots,h_n)T]$. Indeed, on the chart indexed by a
nonzero $h_i$, both schemes have ring
$A[h_0/h_i,\ldots,h_n/h_i]\subset\operatorname{Frac}(A)$.
These charts prove the assertion scheme-theoretically, with the
Rees algebra rather than the symmetric algebra. Apply this to
$\mathfrak a_a$. The Hilbert scheme is closed in the Pl\"ucker
space, so the closure of its dense graph is this same blow-up.
Pulling back the universal Hilbert family proves the flat-family
assertion. Normalizing the identical integral schemes proves the
normalization assertion.

Over $R$, at least one $\Delta_J(\gamma)$ is nonzero. Dividing
all coordinates by their common valuation gives a primitive
vector in a finite free $R$-module. Its special projective point
is $[\ell_J]$. The Grassmannian and Hilbert equations hold over
$K$ and hence over $R$, since $R$ has no $\tau$-torsion. The
resulting $R$-point is the unique extension by separatedness and
properness. The Hilbert immersion makes equality of its special
projective coordinates equivalent to equality of the embedded
special curves. Keeping the coefficient point supplies the last
assertion for the retained-coefficient graph.
\end{proof}

\begin{remark}[What the determinant construction supplies]
\label{rem:scope-states-v167}
The existence of a graph blow-up and the Hilbert immersion are
classical. The explicit ideal \eqref{eq:hilbert-ideal-v167} identifies
their input for every $B_a$ and makes the following valuation
statements uniform in $a$. It is not a claim that arbitrary
higher-contact fibres have only finitely many isomorphism classes.
It does not identify points on the normalization solely from their
Hilbert coordinates: distinct normalization points can lie over
one graph point. A generically nonincident DVR arc has a unique
lift to that normalization by properness and the generic
isomorphism. Equality of closed Hilbert points is a statement on
$\Gamma_a$, not an assertion that the lifts coincide.
\end{remark}

\begin{proposition}[A finite determinacy bound for each arc]
\label{prop:jet-v167}
Keep the matrix and coordinates above fixed. If two $R$-valued
coefficient arcs agree modulo $\tau^{\mu(\gamma)+1}$ and the first
is generically in $U_a$, the second is also generically in $U_a$
and their embedded special curves agree. Thus a jet of length
$\mu(\gamma)+1$ suffices for this arc. The assertion concerns a
Hilbert point, not the whole formal family.
\end{proposition}
\begin{proof}
Polynomial evaluation preserves congruences. All maximal minors
are therefore congruent modulo $\tau^{\mu+1}$, and at least one
has the same nonzero coefficient in degree $\mu$. The primitive
vectors in \eqref{eq:determinantal-order-v167} agree. To justify
the generic-locus assertion, a triple with common homogeneous
factor of positive degree has all images in
\eqref{eq:state-matrix-v167} divisible by its $m$th power and
cannot span $\C[s,t]_{dm}$. Thus a nonzero maximal minor forces
absence of a common zero over an algebraic closure of $K$.
\end{proof}

\begin{proposition}[Smith data and specialization without enumerating minors]
\label{prop:smith-v167}
Let the nonzero invariant factors of $M_a(\gamma)$ over $R$ have
valuations $e_1\leq\cdots\leq e_q$. Then
\[
 \mu(\gamma)=\sum_{i=1}^q e_i.
\]
The cokernel of this degree-$m$ evaluation map is
$\bigoplus_i R/(\tau^{e_i})$. Its least annihilating exponent is
$e_q$, and its length is $\mu$. The saturated row lattice of
$M_a(\gamma)$ in $R^N$ has special exterior line $[\ell_J]$.
The kernel of the corresponding rank-$q$ quotient over $k$,
viewed in the Segre coordinate ring, is the degree-$m$ part of the
special ideal. It generates its ideal sheaf.
\end{proposition}
\begin{proof}
Smith normal form follows by choosing an entry of least valuation,
using it to clear its row and column, and repeating. Elementary
invertible operations preserve the ideal of maximal minors. In
normal form that ideal is generated by
$\tau^{e_1+\cdots+e_q}$, which proves the formula and the cokernel
statements. Over a DVR a saturated submodule of a finite free
module is a direct summand. Saturating the row lattice divides
its exterior line by exactly this common factor. Reduction gives
the stated quotient. Gotzmann regularity, applied to the Hilbert
point already constructed, says that its saturated homogeneous
ideal is generated in degrees at most $m$, and its sheaf is
generated by the degree-$m$ part.
\end{proof}
The exponent $e_q$ here is an exponent for a finite module in one
fixed projective degree. It is not a bound for saturation of every
affine presentation of a pulled-back parameter graph.

\section{Finite coefficient-state chambers}
\label{sec:state-fan-v167}
A valuation without its leading coefficient loses information.
The following finite statement keeps that information, including
cancellation. Its use of coefficient weights is different from a
one-parameter torus acting on the fixed target. Classical state
polytopes organize initial ideals under the latter action
\cite{BayerMorrison167}; here the weighted variables are the input
coefficients of \eqref{eq:state-matrix-v167}.

Write these coefficients as $z_1,\ldots,z_{3a}$ and expand the
finite set of nonzero minors as
\[
 \Delta_J(z)=\sum_{\alpha\in E}c_{J\alpha}z^\alpha,
\]
where $E$ is the union of their finite monomial supports. Allow a
specified subset of the $z_i$ to be identically zero, and omit the
corresponding monomials first. For this theorem take $R=k[[\tau]]$ (or $k[\tau]_{(\tau)}$),
where $k/\C$ is any field extension, so the residue coefficients
are actual constants. No coefficient-field choice was required
for the preceding general DVR statements. On the remaining
coordinates consider exact monomial arcs
\begin{equation}\label{eq:monomial-arcs-v167}
 z_i=\xi_i\tau^{w_i},\qquad \xi_i\in k^*,\quad
                              w_i\in\mathbb Z_{>0}.
\end{equation}
The constant monic coefficient of $f$ is not among these variables.

\begin{theorem}[Finite decorated chamber classification]
\label{thm:finite-fan-v167}
For each fixed $a$ and each choice of zero coordinates there is an
explicit finite rational polyhedral fan in the nonnegative weight
orthant, and for each relatively open cone a finite locally closed
partition of the residue torus, with the following properties.
The pairs for which \eqref{eq:monomial-arcs-v167} is generically
in $U_a$ are specified by these partitions. On each such pair of
a cone and a residue stratum, the embedded Hilbert limit is given
by a single displayed vector of polynomial functions of $\xi$,
independent of $w$ in that relatively open cone. These vectors
classify the special embedded curves by projective equality.
They include nonreduced curves and embedded scheme structure.

One such fan is the hyperplane arrangement
\begin{equation}\label{eq:fan-v167}
              (\alpha-\beta)\cdot w=0\qquad(\alpha,\beta\in E)
\end{equation}
intersected with the coordinate orthant and all its faces.
The residue partition and the vectors require only finite
polynomial operations on the coefficients $c_{J\alpha}$.
\end{theorem}
\begin{proof}
On the relative interior of a cone, the weak order of all numbers
$\alpha\cdot w$ is fixed. Let $E_1,\ldots,E_t$ be its ordered
equality classes. For each $J$ and $h$ form
\[
 b_{Jh}(\xi)=\sum_{\alpha\in E_h}c_{J\alpha}\xi^\alpha.
\]
Partition the residue torus by the first index $h(J)$ at which
$b_{Jh}$ is nonzero, allowing $h(J)=\infty$ when all are zero.
Each condition is a finite list of polynomial equalities and one
inequality, so there are finitely many locally closed pieces.
Discard the piece on which all $h(J)$ are infinite. On the others
let $h_* =\min_J h(J)$. The answer is the projective vector
\begin{equation}\label{eq:decorated-vector-v167}
 L_J(\xi)=
 \begin{cases} b_{J,h_*}(\xi),&h(J)=h_*,\\0,&h(J)>h_*.
 \end{cases}
\end{equation}
Indeed substitution into every minor is a finite sum of powers of
$\tau$, with its terms grouped exactly as above. Equalities in a
residue stratum account for every possible cancellation, not just
a generic noncancellation assumption. The least surviving exponent
is the common valuation in \eqref{eq:determinantal-order-v167}.
Theorem~\ref{thm:determinantal-v167} proves both classification and
retention of scheme structure. The same rank argument as in
Proposition~\ref{prop:jet-v167} identifies the discarded set with
the triples having a generic common zero.
\end{proof}

\begin{remark}[Finite classification and its scope]
\label{rem:finite-classification-v167}
For fixed residue coefficients, the theorem gives finitely many
embedded limits as the positive integral weights vary. As the
residue coefficients vary, it gives finitely many algebraic
families, not finitely many curves. The arrangement need not be
minimal and does not assert that every one of its walls changes
the curve. The degree bound is uniform for each $a$, but no
claim of efficient enumeration of its potentially enormous minors
is made. Arbitrary formal arcs are treated by the pointwise jet
bound, not by declaring their leading coefficient monomials to be
exact arcs. In particular, repeated cancellation in a formal arc
cannot be discarded by this theorem.
\end{remark}

The construction gives a practical distinction between two sorts
of audit. An individual polynomial arc can be evaluated and its
row lattice saturated, without computing the whole arrangement.
The finite arrangement is a proof of finite classification of
exact monomial profiles. The next section computes the entire
answer on a nontrivial two-dimensional slice, including its
minimal walls and the geometry on those walls.

\section{A two-wall classification on a nonsymmetric slice}
\label{sec:two-wall-v167}
The determinant construction can be evaluated without listing its
maximal minors. We do so on a slice on which both a jet direction
and a conic direction vary. Put
\[
 S=\Spec\C[b,c]\longrightarrow B_2,\qquad
                  (f,g,r)=(x^2,b,cx).
\]
Let $T$ be the schematic closure of $S_{b\ne0}$ in
$\Gamma_2\times_{B_2}S$. Thus $T$ is the horizontal main component
of this base change, not its full pullback. All target coordinates
$[F:G:R]$ in this section are fixed.

\begin{theorem}[The two-wall surface]\label{thm:two-wall-v167}
The surface $T$ is the blow-up of $S$ at $(b,c)=(0,0)$ followed
by the blow-up of the point $e=c=0$ in the chart $b=ec$.
Equivalently,
\[
                   T\simeq\operatorname{Bl}_{(b,c)(b,c^2)}S.
\]
The maximal-minor ideal in Theorem~\ref{thm:determinantal-v167}
restricts to the explicitly factored ideal
\begin{equation}\label{eq:slice-minors-v167}
          \mathfrak a_2\OO_S=b^{10}(b,c)^4(b,c^2)^6.
\end{equation}
The surface is smooth and has the following three affine charts:
\begin{equation}\label{eq:three-charts-v167}
 \begin{array}{c|c|c}
 \text{chart}&\text{coordinates}&(b,c)\\ \hline
 T_0&(b,k)&(b,bk)\\
 T_1&(e,h)&(e^2h,eh)\\
 T_2&(c,z)&(c^2z,c).
 \end{array}
\end{equation}
On $T_1$ and $T_2$ the universal curve has respectively the
bihomogeneous ideals
\begin{align}
 I_1&=(R^2-hFG,\ tG-esR,\ tR-ehsF),\label{eq:slice-I1-v167}\\
 I_2&=(FG-zR^2,\ tG-czsR,\ tR-csF).\label{eq:slice-I2-v167}
\end{align}
On $T_0$ its ideal sheaf is generated by
\begin{equation}\label{eq:slice-I0-v167}
 t^2G-bs^2F,\quad sR-ktG,\quad tR-bksF,\quad
                         R^2-bk^2FG.
\end{equation}
These equations define flat families of polynomial $3l+1$.
The fibre $T_{(0,0)}$ is a reduced chain $E_{11}\cup E_{21}$
of two projective lines. Their self-intersections are $-2$ and
$-1$, respectively. The first line parametrizes primitive tails,
and the second the conic pencil through the attaching point;
their intersection parametrizes the doubled conic $R^2=0$.
\end{theorem}

\begin{proof}
Perform the stated blow-ups. The first chart of the first blow-up
is $c=bk$ and is unchanged by the second blow-up. In the other
chart write $b=ec$ and blow up $(e,c)$. Its charts $c=eh$ and
$e=cz$ give exactly \eqref{eq:three-charts-v167}. The transition
maps are
\begin{align*}
 k=e^{-1},\quad b=e^2h,
 &\qquad e\ne0 \text{ on }T_1,\ k\ne0\text{ on }T_0,\\
 z=h^{-1},\quad c=eh,
 &\qquad h\ne0 \text{ on }T_1,\ z\ne0\text{ on }T_2.
\end{align*}
In the first overlap the inverse is $e=k^{-1}$, $h=bk^2$;
in the second it is $h=z^{-1}$, $e=cz$. The third overlap is
the composite, so the cocycle identity holds as an identity of
rings. To check the alternative blow-up description, first blow
up $(b,c)$; on the chart $b=ec$ the transform of $(b,c^2)$ is
$c(e,c)$, whereas on $c=bk$ it is the invertible ideal $(b)$.
Removing the invertible factors and applying the universal
property of a blow-up proves the product-ideal assertion.

We next construct the morphism to the Hilbert scheme. On $T_1$
the Hilbert--Burch chart of
Section~\ref{sec:complete-contact-v163} has centre and parameters
$\lambda=B=C=0$, $A=-h$, and retained coefficient $e$.
Its equations are \eqref{eq:slice-I1-v167}. On $T_2$ a
Hilbert--Burch matrix is
\[
 \begin{pmatrix} F&R\\-zR&-G\\-t&-cs\end{pmatrix}.
\]
Its minors are $-(FG-zR^2)$, $tR-csF$, and
$-(tG-czsR)$. On every geometric fibre the ideal has height two:
if $c\ne0$, the equation $tR=csF$ and the other equations give a
one-dimensional graph or its single-contact limit; if $c=0$,
they give the main section and the conic $FG-zR^2$ at $t=0$.
This conic is nonzero for every $z$. There is no common divisor
of the three minors. The fibrewise Hilbert--Burch resolution
\[
 0\longrightarrow\OO(-1,-2)^2\longrightarrow
 \OO(-1,-1)^2\oplus\OO(0,-2)\longrightarrow\OO
 \longrightarrow\OO_C\longrightarrow0
\]
is therefore exact, and its constant Euler characteristic gives
flatness and the polynomial $3l+1$. Equivalently one may apply
the fibrewise flatness criterion successively to this complex of
relatively locally free sheaves.

For $T_0$, on $s\ne0$ eliminate $R$ from
\eqref{eq:slice-I0-v167}. The remaining equation in the relative
$\Pj^1$ with coordinates $[F:G]$ is
\[
                     x^2G=bF,\qquad R=kxG.
\]
It is a relative effective Cartier divisor, since $x^2G-bF$ is
nonzero on every fibre. On $t\ne0$, with $y=s/t$, the equations
instead give $G=by^2F$ and $R=bkyF$, the graph of a section.
These two descriptions prove flatness and exclude any excess
component at infinity. This is also the primitive division family
of Theorem~\ref{thm:division-chart-v162}, with $w=kx$.
On the overlaps substitution of the displayed transitions makes
the ideal sheaves identical; on $T_1\cap T_2$ the quadratic
equations differ by the invertible factor $-h$. Thus the three
flat families glue in the fixed target. Their generic graph is
$[t^2:bs^2:cst]$.

It remains to prove that the constructed surface is the whole
horizontal graph, rather than a resolution mapping to it.
For $a=2$, the matrix $M_2$ has $13$ rows and $75$ columns.
After restricting to $F_0=t^2$, $F_1=bs^2$, $F_2=cst$,
every column has exactly one nonzero entry. Its row index is
$i+2j_0+j_2$ and its entry is $b^{j_1}c^{j_2}$. A nonzero
maximal minor chooses exactly one column from each row.
Consequently its ideal is the product of the ideals generated
by the entries in the individual rows. With $J_1=(b,c)$ and
$J_2=(b,c^2)$ those row ideals, in order from $0$ to $12$, are
\[
 b^4,\ b^3J_1,\ b^2J_2,\ bJ_1J_2,\ J_2^2,\ J_1J_2,
                  \ J_2,\ J_1,\ \OO_S,\OO_S,\OO_S,\OO_S,\OO_S.
\]
They follow by taking all $0\leq i\leq4$ and
$j_0+j_1+j_2=4$ with the indicated row index; divisibility removes
redundant monomials. Their product proves
\eqref{eq:slice-minors-v167}.

The graph closure of the restricted rational Pl\"ucker map is
therefore $\operatorname{Bl}_{b^{10}J_1^4J_2^6}S$, by the graph
argument in Theorem~\ref{thm:determinantal-v167}. This is not an
appeal to unrestricted base change for blow-ups. A principal
nonzero factor does not change the Rees blow-up. Moreover blowing
up $J_1^4J_2^6$ is the same as blowing up $J_1J_2$: on either
integral blow-up, invertibility of the product makes each factor
invertible as a fractional ideal, and conversely. For completeness,
if two nonzero finitely generated ideals have invertible product,
then locally, after dividing their product by its generator, they
are inverse fractional ideals. Thus each is invertible; apply
this repeatedly to the positive powers. The universal properties
supply inverse morphisms between the two blow-ups, identical on
the dense open. This proves that the twice blown-up surface is
exactly $T$. The flat curve constructed above equals its Hilbert
universal curve since the two Hilbert morphisms agree on the
dense open of an integral surface and the Hilbert scheme is
separated. In particular no unproved normality of a graph image
is used in the argument.

On $T_1$ the ideal of the origin pulls back to $(eh)$; on $T_0$
it is $(b)$, and on $T_2$ it is $(c)$. Hence the exceptional fibre
is reduced with a transverse intersection. The strict transform
of the first exceptional curve loses one from its self-intersection
when the second point is blown up. This gives $E_{11}^2=-2$,
$E_{21}^2=-1$, and $E_{11}E_{21}=1$. Substitution $h=0$ and $e=0$
in \eqref{eq:slice-I1-v167} gives its two modular descriptions and
their common doubled conic.
\end{proof}

\begin{corollary}[Complete valued-arc limits on the slice]
\label{cor:slice-chambers-v167}
Let $R$ be a DVR containing $\C$, with uniformizer $\tau$ and
residue field $k$. Let $b=\beta\tau^p$, $c=\gamma\tau^q$, where
$p,q>0$ and $\beta,\gamma\in R^*$; write bars for their residues.
The special embedded curve
has the following ideal sheaf. The displayed bihomogeneous
ideals are understood with irrelevant saturation.
\begin{equation}\label{eq:slice-table-v167}
\begin{array}{c|c|c}
 \text{weights}&\text{residue parameter}&\text{special ideal}\\ \hline
 p<q&k=0&(R,t^2G)\\
 p=q&k=\bar\gamma/\bar\beta&(sR-ktG,tR,R^2)\\
 q<p<2q&& (tG,tR,R^2)\\
 p=2q&h=\bar\gamma^2/\bar\beta&(tG,tR,R^2-hFG)\\
 p>2q&& (tG,tR,FG).
\end{array}
\end{equation}
These five cases exhaust all generically nonincident DVR arcs
centred at the origin with $bc\ne0$ in the fraction field,
including units with arbitrary higher coefficients. If $c=0$
identically, the first row applies. Moreover the determinantal
order on this slice is
\begin{equation}\label{eq:slice-jet-order-v167}
 \mu=10p+4\min(p,q)+6\min(p,2q),
\end{equation}
with $\mu=20p$ when $c=0$.
The two walls $p=q$ and $p=2q$ are actual walls for embedded
limits; coefficients on the walls cannot be omitted.
\end{corollary}
\begin{proof}
For $p\leq q$, use $k=c/b$ on $T_0$. For $q<p\leq2q$ use
$e=b/c$ and $h=c^2/b$ on $T_1$. For $p>2q$ use $z=b/c^2$
on $T_2$. In each case these coordinates belong to $R$, and their
reductions give the table. The formula for $\mu$ follows by taking valuations in
\eqref{eq:slice-minors-v167}; an ideal generated by two elements
has valuation equal to the lesser of their valuations. It is
unaffected by higher unit coefficients. For $k\ne0$ the first generator
$t^2G$ of the primitive central fibre follows from
$sR-ktG$ and $tR$ after taking ideal sheaves. For $k=0$ it must
be retained. This explains why setting $k=0$ in just the last
three generators would give the wrong scheme.
\end{proof}

\begin{proposition}[Components, multiplicities, and the parameter boundary]
\label{prop:slice-geometry-v167}
Every curve in \eqref{eq:slice-table-v167} is Cohen--Macaulay and
has no embedded associated points. In the first three rows the
minimal primes are $(G,R)$ and $(t,R)$; the main section has
multiplicity one and the tail multiplicity two. In the fourth
row the primes are $(G,R)$ and $(t,R^2-hFG)$, each of multiplicity
one. In the fifth they are $(G,R)$, $(t,F)$, and $(t,G)$,
each of multiplicity one. The fourth curve is a main section
attached to a smooth conic; the fifth is a chain of three reduced
lines. The normalization of each reduced curve is the disjoint
union of its smooth components. None of these attached curves
is normal as a whole.

The parameter surface $T$ is normal and smooth. Its exceptional
boundary has simple normal crossings, with
\[
 \begin{array}{c|cc} &E_{11}&E_{21}\\ \hline
 \operatorname{ord}(b)&1&2\\
 \operatorname{ord}(c)&1&1
 \end{array},\qquad
 K_{T/S}=E_{11}+2E_{21}.
\]
These are statements about a parameter-space modification, not a
stable reduction theorem for curves.
\end{proposition}
\begin{proof}
For the primitive rows, the affine presentation
$R=kxG$, $x^2G=0$ is a hypersurface in a smooth relative surface.
Near infinity it is the smooth main section. The two generic
lengths are read from $x^2G$; the residue $k$ changes the
embedding but not the lengths. For the conic rows the
Hilbert--Burch resolution above makes every local ring on the
curve Cohen--Macaulay of dimension one. At the attaching point
use $F=s=1$: the ideals are $(xG,xR,R^2-hG)$, with the analogous
$FG$ presentation in the last row. The ideal decompositions
\[
 (tG,tR,Q)=(G,R)\cap(t,Q),\qquad Q\in\{R^2,R^2-hFG,FG\},
\]
as ideal sheaves give the claimed minimal components and lengths.
A one-dimensional Cohen--Macaulay scheme has no embedded points.
When $h\ne0$ the conic is smooth. When $Q=FG$, its two lines
meet at $[0:0:1]$ and only the line $G=0$ contains the attaching
point $[1:0:0]$. This gives the chain description. The assertions
on the normalization of the reductions follow.

The valuations follow directly from $b=e^2h$, $c=eh$:
$E_{11}$ is $h=0$ and $E_{21}$ is $e=0$. For a blow-up of a
smooth surface at a point the relative canonical divisor is its
exceptional curve, as the Jacobian in a chart shows. Pulling the
first exceptional divisor through the second blow-up adds one
more copy of $E_{21}$. This proves the discrepancy formula.
\end{proof}

\begin{corollary}[Pointwise, but not uniform, finite determinacy]
\label{cor:no-uniform-v167}
Even on $B_2$ no integer $n$ has the property that the coefficient
$n$-jet determines every embedded graph limit. In contrast,
Proposition~\ref{prop:jet-v167} gives an explicit finite bound for
each generically nonincident arc.
\end{corollary}
\begin{proof}
Given $n$, choose $M>n$ and two distinct nonzero constants
$\beta_1,\beta_2$. The arcs
\[
 f=x^2,\qquad g=\beta_i\tau^{2M},\qquad r=\tau^M x
 \quad(i=1,2)
\]
agree modulo $\tau^{n+1}$ and are generically base-point free.
Their limits are the distinct fixed-target conics
$R^2-\beta_i^{-1}FG$ attached to the same main section, by the
fourth row of \eqref{eq:slice-table-v167}. This proves the claim.
\end{proof}

\begin{proposition}[No uniform saturation exponent for full pullbacks]
\label{prop:no-saturation-bound-v167}
The uniformizer saturation exponents of the full pulled-back
parameter graph at contact order two are unbounded. This is
compatible with the finite degree-module statement of
Proposition~\ref{prop:smith-v167}.
\end{proposition}
\begin{proof}
In the symmetric family of Theorem~\ref{thm:ramified-v166}, take
$\delta=\tau^m$. Its full vertical torsion is a nonzero ideal
module tensored with $\C[\tau]/(\tau^m)$. On a chart where that
ideal is nonzero it is killed by $\tau^m$ and not by
$\tau^{m-1}$. Thus the least exponent at which uniformizer
saturation stabilizes is at least $m$. Since $m$ is arbitrary,
there is no uniform bound with contact order fixed at two.
This assertion is about the full parameter-family pullback; it
does not replace the embedded-curve jet theorem with a negative
statement.
\end{proof}

\section{Base change and comparison of horizontal models}
\label{sec:comparison-v167}
We now distinguish the Hilbert graph of embedded curves from the
horizontal modification of a projective parameter family. In the
notation of Section~\ref{sec:horizontal-v166}, let $W\to B$ be
projective and of finite presentation, with $B$ integral of finite
type over $\C$. Fix a dense open $U\subset B$ over which $W$ is
flat with polynomial $P$. The admissible tests are locally
Noetherian morphisms $Y\to B$ for which $U_Y$ is schematically
dense. The horizontal object $W_Y^h$ is the schematic closure of
$W_{U_Y}$ in $W_Y$. Only tests for which this closure is flat
belong to the represented horizontal-flat subfunctor. Throughout,
\[
 H_B=\overline{s(U)}\ \subset\operatorname{Hilb}^{P}(W/B)
\]
means the \emph{horizontal Hilbert-main-component modification}.
Its terminal property refers to this admissible category and
this prescribed generic embedded quotient.

\begin{lemma}[Flat schematic density]\label{lem:density-v167}
Let $V\subset Y$ be a quasi-compact schematically dense open in a
locally Noetherian scheme. If $Z\to Y$ is flat, then $Z_V$ is
schematically dense in $Z$.
\end{lemma}
\begin{proof}
The assertion is local on $Y$ and $Z$. For $Y=\Spec A$, write
$V=\bigcup_{i=1}^rD(f_i)$. Schematic density is the injectivity
of $A\to\prod_i A_{f_i}$. Tensoring with any flat $A$-module
preserves this injection; finiteness of the product identifies
the tensor product with the product of the localizations. Apply
this to the rings of affine opens of $Z$. A section vanishing on
$Z_V$ is zero. This is precisely the required schematic density.
\end{proof}

\begin{theorem}[Strict base change and coherent comparison]
\label{thm:basechange-v167}
Let $B'\to B$ be an integral finite-type base change such that
$U'=U\times_BB'$ is nonempty. Use the pullback polarization and
write $H_{B'}$ for the horizontal Hilbert-main-component
modification of $W_{B'}$ relative to $U'$. There is a canonical
closed immersion
\begin{equation}\label{eq:comparison-closure-v167}
 H_{B'}=\overline{U'}^{\,H_B\times_BB'}
                    \lhook\joinrel\longrightarrow H_B\times_BB'.
\end{equation}
It is an equality if and only if $U'$ is schematically dense in
the full fibre product. In particular it is an equality for flat
base change. For composable base changes meeting the indicated
generic opens, these comparison maps compose to the comparison
for the composite. They obey the corresponding cocycle identity.

If $B^\nu\to B$ is the normalization, then $H_{B^\nu}\to H_B$
is finite and birational, and the normalizations of these two
integral schemes are canonically isomorphic over $B$.
\end{theorem}
\begin{proof}
The relative Hilbert scheme and its universal family commute with
base change. In the resulting Hilbert scheme over $B'$, the
closure of $s(U')$ is contained in the closed subscheme
$H_B\times_BB'$, since the ideal of $H_B$ vanishes on $s(U')$.
It is therefore exactly the schematic closure on the right of
\eqref{eq:comparison-closure-v167}. This proves the first
assertion and the density criterion. Flat base change preserves
the schematic density of $U\subset H_B$ by
Lemma~\ref{lem:density-v167}; note that the projection of the
fibre product to $H_B$ is flat.

For a further base change $B''\to B'$, both routes give the
schematic closure of $U''$ in the same relative Hilbert scheme.
The closed immersions consequently agree, including their scheme
structures. This proves transitivity and the cocycle, not merely
agreement on closed points. Finally, normalization of a finite-type
integral $\C$-scheme is finite. The map
$H_B\times_B B^\nu\to H_B$ is finite; its closed integral main
component is finite and birational because it agrees on a dense
normal open. Finite birational integral extensions have the same
integral closure in their common function field. This proves the
last assertion.
\end{proof}

\begin{example}[The full pullback need not commute]
\label{ex:basechange-v167}
Take $B=\A^2_{x,y}$, $W=\operatorname{Bl}_0B$, and
$U=B\setminus\{0\}$. The generic polynomial is $P=1$.
A flat family of length-one embedded subschemes of $W/B$ is a
section, so its relative Hilbert scheme is $W$ itself and
$H_B=W$. Restrict to $B'=\{y=0\}$. With homogeneous coordinates
$[u:v]=[x:y]$, the full pullback has equation $xv=0$ in
$\A^1_x\times\Pj^1$. It contains a vertical projective line.
The closure of $x\ne0$ is instead $v=0$, isomorphic to $B'$.
Thus \eqref{eq:comparison-closure-v167} is a strict inclusion.
The parameter $y=0$ is a nonflat base change; replacing the full
pullback by its reduction would not remove its extra component.
\end{example}

\begin{proposition}[Common refinements on one generic open]
\label{prop:refinement-v167}
Fix $B,U$. For two integral proper $U$-modifications $Y_1,Y_2$,
let $Y_1\star_UY_2$ be the schematic closure of the diagonal
$U$ in $Y_1\times_BY_2$. The projections are proper birational
maps and this object is terminal among integral $U$-modifications
equipped with maps to both $Y_1$ and $Y_2$ restricting to the
identity on $U$. Repeated products are canonically associative
and symmetric, and their comparison maps satisfy all product
cocycles. If a refinement of a stratum has the same reduced
closure and a smaller dense generic-flat open, its horizontal
model is canonically unchanged.
\end{proposition}
\begin{proof}
A pair of compatible maps gives a map into the product. Since its
integral source has dense $U$, the map factors through the
schematic closure of that open. The factorization is unique
because the closure is a closed subscheme. Properness follows
from properness of the product, and birationality from the common
dense open. For three or more factors every iterated object is
the closure of the same diagonal in the full product; this proves
associativity, symmetry, and their compatibility. Shrinking a
dense open of an integral section does not change its schematic
closure, proving the last assertion.
\end{proof}

The theorem supplies comparison for dominant restrictions and
strict base changes, and the proposition supplies common
refinements on a fixed generic problem. A stratum whose closure
lies entirely outside $U$ does not meet these hypotheses. In
particular these results do not identify Hilbert components with
different generic polynomials. For such strata the original
Noetherian-induction construction remains a choice of strata,
not a canonical incidence poset. If one asks for a flat family
whose fibres are the entire original fibres on two strata in one
connected base, equality of their Hilbert polynomials is necessary.
If $W$ is already flat over all of $B$, the horizontal model is
$B$ itself. These observations specify both the full-base flat
case and the precise limitation of the comparison theorem.

\subsection{Comparison with full flattening}
For an admissible $Y\to B$ on which $W_Y$ is flat,
Lemma~\ref{lem:density-v167} gives $W_Y=W_Y^h$. There is therefore
a natural inclusion of functors and a commutative diagram
\[
\begin{CD}
 \{Y:W_Y\text{ is flat}\}@>>>\{Y:W_Y^h\text{ is flat}\}\\
 @VVV @VVV\\
 \operatorname{Hom}_B(Y,B)@=\operatorname{Hom}_B(Y,B).
\end{CD}
\]
The top arrow is not in general surjective: the extra projective
line in Example~\ref{ex:basechange-v167} is exactly an obstruction
to full flattening. The usual flattening functor of a finitely
presented sheaf concerns its full tensor pullback and, when
represented, is represented by a monomorphism; see
\cite{Stacks166}, Tag~052F. By contrast, a horizontal modification
may have a positive-dimensional exceptional fibre. Flattening by
blow-up uses a pure transform and must likewise not be confused
with flatness of the full pullback. Theorem~\ref{thm:determinantal-v167}
computes a Rees ideal for the \emph{embedded-curve graph}; it does
not assert that the same ideal flattens every fibre of that
parameter graph.

The horizontal main component is integral, hence reduced. Neither
normality nor Cohen--Macaulayness follows from the construction:
for $W=B$ the construction returns $B$, which may be a nonnormal
integral curve or a non-Cohen--Macaulay integral higher-dimensional
scheme. The explicit surface in Theorem~\ref{thm:two-wall-v167}
has stronger properties because its charts have been computed.
Independence of polarization identifies representing objects for
the same horizontal functor; it does not identify the ambient
Hilbert schemes or their other components.

\section{Determinantal specialization after effective reconstruction}
\label{sec:effective-states-v167}
The inverse theorem and the Hilbert-limit theorem solve distinct
problems. We state their composition so that no operation on a
raw unframed multiplication table is inferred.

\begin{corollary}[Effective-family specialization]
\label{cor:effective-states-v167}
On the effectively rigidified failure-family image, suppose that
a reconstructed regular contact is given in the fixed-target
polynomial chart $B_a$ and that its DVR generic point lies in
$U_a$. The embedded specialization of that contact is determined
by the primitive maximal-minor vector of $M_a$, and by the
coefficient jet modulo $\tau^{\mu+1}$. This rule is compatible
with the fixed-target undo maps for changes of contact chart.
For exact monomial coefficient profiles the finite decorated
chamber theorem applies; on $f=x^2,g=b,r=cx$ its complete answer
is the two-wall classification.
\end{corollary}
\begin{proof}
The sharp inverse reconstructs the pencil on the effective image,
as used in Proposition~\ref{prop:effective-horizontal-v166}.
The fixed-target comparisons identify the local graph family
before and after undoing the contact coordinates. On each chart,
Theorem~\ref{thm:determinantal-v167} computes the same embedded
Hilbert point from its degree-$m$ quotient. Two presentations
of this quotient have the same Hilbert point because the Hilbert
immersion is injective as a morphism of functors. The finite jet
statement follows from Proposition~\ref{prop:jet-v167}, and the
last statements from Theorem~\ref{thm:finite-fan-v167} and
Corollary~\ref{cor:slice-chambers-v167}.
\end{proof}
This statement uses reconstruction first. Its finite jet is a jet
of the recovered coefficient family in the stated chart, not a
proved bound on the jet of an arbitrary multiplication table.
A closed finite algebra alone contains no specified specialization
arc. Nor does the corollary identify different normalization
branches lying over one embedded Hilbert point. The complete
proof of the inverse remains part of Paper I; no external
independent audit is asserted by this application.

\section{Evaluation ideals and their divisorial content}
\label{sec:content-v169}
Throughout the new results the ground field $k$ is algebraically closed
of characteristic zero. The equations, flat families, and valuation
formulas are defined over $\mathbb Z$ and remain valid over a field
whenever the displayed nonzero residues are retained. We distinguish a
rational projective map, a choice of homogeneous coordinates for that
map, and a choice of projective embedding of its target.

\subsection{A primitive ideal for a fixed projective map}
Let $B$ be a regular integral scheme of finite type over $k$, and let
$I\subset\OO_B$ be a nonzero coherent ideal. At a prime divisor $D$ put
$c_D=\min_{f\in I}\operatorname{ord}_D(f)$, and set
\[
 C(I)=\sum_D c_DD,\qquad
 I^{\mathrm{pr}}=I\otimes\OO_B(C(I))\subset\OO_B.
\]
Only finitely many coefficients are nonzero locally. Regularity makes
$C(I)$ Cartier. Equivalently, $I^{**}=\OO_B(-C(I))$ and
$I^{\mathrm{pr}}=I\otimes(I^{**})^{-1}$. The inclusion follows by checking
orders at prime divisors in the normal ring $\OO_B$.

\begin{proposition}[Content and the graph]\label{prop:content-v169}
Suppose $h_0,\ldots,h_N$ generate $I$ and define a rational map
$B\dashrightarrow\Pj^N$. Its schematic graph closure is
$\operatorname{Bl}_I B$, and
\[
 \operatorname{Bl}_I B\simeq\operatorname{Bl}_{I^{\mathrm{pr}}}B.
\]
Multiplication of the homogeneous coordinates by one nonzero rational
function, and an invertible change of their basis, leave
$I^{\mathrm{pr}}$ unchanged, whenever the resulting coordinates are
regular. These assertions concern the same projective embedding;
changing that embedding need not preserve the primitive ideal.
\end{proposition}
\begin{proof}
On an affine factorial chart write $I=gJ$, where $g$ is the greatest
common divisor of the generators. Division by $g$ identifies
$I^n$ with $J^n$ compatibly with multiplication, hence identifies the
two Rees algebras as graded algebras. The ratio chart with pivot $h_i$
is $A[h_0/h_i,\ldots,h_N/h_i]$ inside the fraction field. These charts
are exactly the closure of the generic graph. They also show the
invariance under a common rational scalar. An invertible matrix changes
the generating list but not its ideal. The local constructions agree
under the unit changes of $g$, and therefore glue. This is the classical
Rees graph construction, not a new representability assertion
\cite{Stacks169}.
\end{proof}

\begin{proposition}[Primitive finite determination]\label{prop:primitive-jet-v169}
On an affine factorial chart choose polynomial generators of
$I^{\mathrm{pr}}$. Let $\gamma:\Spec k[[\tau]]\to B$ have generic
point in the domain of the rational map, and put
$\nu=\operatorname{ord}_\gamma(I^{\mathrm{pr}})$. A second coefficient
arc congruent to $\gamma$ modulo $\tau^{\nu+1}$ has the same closed
projective limit. The bound is sufficient, and is not asserted to be
optimal. It determines the closed graph point, not the whole formal
family or a point of the normalization above it.
\end{proposition}
\begin{proof}
Polynomial evaluation preserves this congruence. At least one primitive
generator has order exactly $\nu$, so the second arc has the same
minimum order. After division by $\tau^\nu$, the two residue vectors
are equal and nonzero. The projective limits coincide. Changing a local
trivialization multiplies this vector by a unit. A change of regular
coordinates preserves the congruence relation on local-ring maps.
\end{proof}
The integer $\nu$ is posterior: an oracle for a formal series need not
reveal it in a predetermined number of queries. For polynomial or
localized-polynomial arcs, polynomial division and exact linear algebra
over $k(\tau)$ compute it in finitely many operations. No efficiency
claim follows from that termination statement.

\subsection{The evaluation module}
Write $V=k^2$, $W=k^3$, and $X=\Pj(V)\times\Pj(W)$. A triple of degree
$a$ binary forms gives, by substitution and multiplication, the map
\begin{equation}\label{eq:eval-covariant-v169}
 E_m:\operatorname{Sym}^m V^*\otimes\operatorname{Sym}^m W^*
       \longrightarrow\operatorname{Sym}^{(a+1)m}V^*.
\end{equation}
Our convention is a quotient matrix: columns are images of domain basis
vectors and rows are coefficients in the target. Thus
\[
 q=(a+1)m+1,\qquad N=(m+1)\binom{m+2}{2},\qquad M_{a,m}\in A^{q\times N}.
\]
Its maximal-minor ideal is $\operatorname{Fitt}_0(\operatorname{coker}E_m)$.
Source and target changes act through symmetric powers in
\eqref{eq:eval-covariant-v169}; in chosen bases the matrix changes by
invertible left and right matrices, together with the corresponding
base automorphism. The Fitting ideal is consequently covariant. Scalar
changes of a representative of the triple instead multiply all columns
by a common scalar to the $m$th power. This is a line-bundle twist on
projective coefficient space, not an invariant numerical order.

For $P_a(l)=(a+1)l+1$ put
\begin{equation}\label{eq:gotzmann-size-v169}
 m_a=\frac{a(a+1)}2+1.
\end{equation}
Indeed, with $d=a+1$, the $d$ linear terms
$\binom{l+1-i}{1}$, $0\leq i<d$, sum to
$dl+d(3-d)/2$. Adding $(d-1)(d-2)/2$ constant terms gives $dl+1$;
the number of terms is $d(d-1)/2+1$. Gotzmann regularity applies in
the ordinary projective space of the Segre embedding. The restriction
of its degree-$m$ section space to $X$ is surjective: a monomial of
bidegree $(m,m)$ is a product of $m$ Segre coordinates. The fixed
kernel of this restriction belongs to every curve ideal in question.
Thus the Hilbert immersion factors through the quotient space in
\eqref{eq:eval-covariant-v169}. For $m\geq m_a$ the degree-$m$ kernel
generates the ideal \emph{sheaf}; this is the truncation consequence of
regularity, not a claim that the homogeneous ideal has no smaller
minimal generators. These are the classical Hilbert-embedding inputs
\cite{Gotzmann167}.

\begin{proposition}[Rank, incidence, and content on the full chart]
\label{prop:incidence-content-v169}
Let $B_a$ be the affine space of triples $(f,g,r)$ with $f$ monic of
degree $a$ and $\deg g,\deg r<a$. For $m\geq m_a$, $E_m$ has rank $q$
exactly when the three homogenized forms have no common zero. The
rank-drop support is irreducible of codimension two. In particular,
$I_q(M_{a,m})$ has no divisorial content on $B_a$.

Let $C_f$ be the companion matrix for multiplication by $x$ on
$A[x]/(f)$. The smaller presentation
\begin{equation}\label{eq:basepoint-fitt-v169}
 A^{2a}\xrightarrow{[g(C_f)\ \ r(C_f)]}A^a
       \longrightarrow A[x]/(f,g,r)\longrightarrow0
\end{equation}
defines a natural incidence Fitting ideal with the same radical as
$I_q(M_{a,m})$. Equality of these two ideals is not asserted.
\end{proposition}
\begin{proof}
On the base-point-free locus restriction to the graph is the quotient
onto $H^0(\Pj^1,\OO((a+1)m))$ in a regularity degree. If a common form
of positive degree divides the triple, every substituted degree-$m$
section is divisible by its $m$th power, so the quotient cannot have
rank $q$. There is no common zero at infinity since $f$ is monic and
the other leading coefficients vanish. The incidence variety
\[
 \{(x,f,g,r):f(x)=g(x)=r(x)=0\}\subset\A^1\times B_a
\]
is affine space of dimension $3a-2$: solve for the three constant
coefficients. Its projection is finite, because $x$ satisfies the
monic equation $f$. Its image is therefore closed, irreducible, and
of that dimension. An ideal with codimension-two support in a regular
integral scheme has order zero at every prime divisor. Finally
\eqref{eq:basepoint-fitt-v169} is multiplication by $g,r$ in the finite
free algebra $A[x]/(f)$. Its cokernel is the stated quotient, whose
support is exactly the same incidence image.
\end{proof}
Content removal need not commute with nonflat restriction. The full
ideal in Proposition~\ref{prop:incidence-content-v169} is already
primitive; its restriction to the two-dimensional families below
has a large common factor. Thus the primitive construction does not
turn the original minor order into an embedding-independent invariant.

\begin{example}[Degree changes without graph changes]\label{ex:degree-content-v169}
For $a=1$, translation reduces the triple to $(x,v,w)$. In every degree
$m\geq1$ the maximal-minor ideal is
\[
 (v,w)^{m(m+1)/2}.
\]
To see this, the row indexed by $n$, $0\leq n\leq2m$, has ideal
$(v,w)^{\max(0,m-n)}$, and a full minor chooses one entry in every
row. The product gives the formula. All these ideals are primitive
and define the same blow-up, although their orders along an arc are
different. Their Smith lengths are lengths of different evaluation
modules, not lengths attached to the rational Hilbert map alone.
\end{example}

\begin{proposition}[What survives a change of Hilbert degree]
\label{prop:degree-change-v169}
For two degrees at least $m_a$, the graph modifications obtained from
$E_m$ are canonically isomorphic over $B_a$, with the same universal
embedded curve and the same normalization. Their exceptional prime
divisors therefore define the same divisorial valuations of $k(B_a)$.
The orders of the evaluation ideals on those divisors, the Smith
factors, and the primitive ideals themselves need not be equal.
There is in general no inclusion-minimal nonzero ideal defining a
given nontrivial blow-up: taking positive powers leaves its Proj
unchanged.
\end{proposition}
\begin{proof}
Both constructions are the closure of the same section in the same
Hilbert scheme; the two Hilbert immersions only embed that scheme in
different Grassmannians. The graph identification follows from
Proposition~\ref{prop:content-v169}. Normalization is intrinsic to the
integral graph. On the normalization, an exceptional divisor is a
prime divisor independently of its chosen relatively ample line
bundle. Those line bundles can have different orders, as the preceding
example demonstrates. Finally the Rees algebra of $I^n$ is a Veronese
subalgebra of that of $I$; a proper nonzero ideal in a domain has
strictly smaller powers at any point where it is proper.
\end{proof}

The worst-case presentation has $m_a=\tfrac12a^2+O(a)$,
$q=\tfrac12a^3+O(a^2)$, and $N=\tfrac1{16}a^6+O(a^5)$. There are at
most $\binom Nq$ candidate maximal minors, with
\[
 \log\binom Nq=\tfrac32a^3\log a+O(a^3+a^2\log a).
\]
This is a size bound, not a practical algorithm. The next theorem
replaces this list by $a$ two-generated ideals and at most $a+1$
curve equations on each chart of an all-order family.

\section{A chain model for contacts of arbitrary order}
\label{sec:chain-model-v169}
For $a\geq2$ consider the marked family over $S_a=\Spec k[b,c]$
\begin{equation}\label{eq:marked-family-v169}
 [F_0:F_1:F_2]=[t^a:b s^a:cst^{a-1}].
\end{equation}
Let $T_a$ be the closure of its Hilbert graph over $b\ne0$ in
$S_a\times\Hilb^{P_a}(X)$. This is a strict restriction of the
coefficient graph. The full fibre product with $\Gamma_a$ can have
additional vertical components and is not the definition of $T_a$.

\begin{theorem}[The all-order chain model]\label{thm:chain-model-v169}
For $m\geq m_a$ set
\[
 C=\frac{(a-1)m(m+1)}2,\quad
 E_j=m-j+1\ (1\leq j<a),\quad
 E_a=\frac{(m-a+1)(m-a+2)}2.
\]
The restriction of the evaluation ideal is exactly
\begin{equation}\label{eq:all-order-factor-v169}
 I_q(M_{a,m})\OO_{S_a}
   =b^C J_{a,m},\qquad
 J_{a,m}=\prod_{j=1}^a(b,c^j)^{E_j}.
\end{equation}
Its divisorial content is $b^C$. The ideals $J_{a,m}$ and all their
powers are integrally closed, and
\begin{equation}\label{eq:chain-blowup-v169}
 T_a=\operatorname{Bl}_{J_{a,m}}S_a
     =\operatorname{Bl}_{K_a}S_a,
 \qquad K_a=\prod_{j=1}^a(b,c^j).
\end{equation}
This surface is smooth. It is obtained by blowing up the origin and
then, $a-1$ times, the intersection of the newest exceptional curve
with the strict transform of $b=0$. Its exceptional divisors have
valuations
\[
 v_i(b)=i,\quad v_i(c)=1\qquad(1\leq i\leq a).
\]
The universal curve admits the explicit flat presentations below.
The theorem classifies the marked family \eqref{eq:marked-family-v169}
for every $a$; it does not identify the whole fibre of $\Gamma_a\to B_a$
with the fibre of this surface.
\end{theorem}

\subsection{A row factorization, without enumerating minors}
\begin{proof}[Proof of \eqref{eq:all-order-factor-v169}]
Every column has one nonzero entry. If $h,l$ are the exponents of the
second and third forms, that entry is $b^hc^l$ in row
$n=i+a(m-h-l)+(a-1)l$, $0\leq i\leq m$. A nonzero full minor chooses
one column in each row; columns belonging to distinct rows are
necessarily distinct. The determinant ideal is therefore the product
of the row ideals $Q_n$.

Put $d=am-n$. Rows with $n\geq am$ have $Q_n=(1)$. For $d\geq0$ set
\[
 h_0=\max\left(0,\left\lceil\frac{d-m}{a-1}\right\rceil\right),
 \qquad r=d-ah_0.
\]
For a fixed $h$, the smallest feasible $l$ is $\max(0,d-ah)$;
feasibility requires $h\geq h_0$. Removing dominated monomials gives
\begin{equation}\label{eq:row-factor-v169}
 Q_n=
 \begin{cases}
 (b^{h_0}),&r\leq0,\\
 b^{h_0}(b,c^a)^u(b,c^\rho),&r=au+\rho>0,\quad0\leq\rho<a,
 \end{cases}
\end{equation}
where the last factor is omitted for $\rho=0$. In fact the minimal
list in the second case is
$b^{h_0+v}c^{\max(0,r-av)}$, $0\leq v\leq\lceil r/a\rceil$.
Expanding the right side of \eqref{eq:row-factor-v169} and discarding
dominated terms gives precisely this list.

The sum of the $h_0$ is
$(a-1)(1+\cdots+m)=C$. Among all rows, a positive residual value
$r\in\{1,\ldots,m\}$ occurs
\[
 \nu(r)=\min(a,m-r+1)
\]
times. One occurrence comes from $h_0=0$; for $h_0=h\geq1$ it comes
from $m-a+2-h\leq r\leq m-h$. Counting the residual factors gives
\[
 \sum_{\substack{r\equiv j\pmod a\\1\leq r\leq m}}\nu(r)=m-j+1
 \quad(1\leq j<a),
 \qquad
 \sum_{r=1}^m\nu(r)\left\lfloor\frac ra\right\rfloor
       =\sum_{n=a}^m(n-a+1)=E_a.
\]
For the first equality write $n=r+v$, $0\leq v<a$, and count each
$n\in\{j,\ldots,m\}$ once. For the second, at a fixed $n$ sum
$\lfloor r/a\rfloor$ over
$\max(1,n-a+1)\leq r\leq n$; the answer is $\max(0,n-a+1)$.
This proves the ideal identity. The argument in fact works for every
$m\geq a$; the regularity bound is used only to identify its graph
with the Hilbert graph.
\end{proof}

\subsection{The Rees algebra and its rays}
\begin{proposition}[The normal Rees algebra]\label{prop:rees-v169}
Put $\eta_i=\sum_{j=1}^aE_j\min(i,j)$. Then
\begin{equation}\label{eq:rees-semigroup-v169}
 \mathcal R(J_{a,m})=
 k[b^u c^v T^n:\ u,v,n\geq0,\quad
                 iu+v\geq n\eta_i\ (1\leq i\leq a)].
\end{equation}
The right side is already normal, not merely the normalization of the
Rees algebra. For the raw ideal in \eqref{eq:all-order-factor-v169}
replace $u\geq0$ by $u\geq nC$ and the last inequalities by
$iu+v\geq n(iC+\eta_i)$. The exceptional Rees valuations of the
primitive ideal are exactly $v_1,\ldots,v_a$.
\end{proposition}
\begin{proof}
For a prescribed integer exponent $u$ of $b$, the least exponent of
$c$ in a product of the factors $(b,c^j)^{E_j}$ is obtained by replacing
$c^j$ by $b$ first for the largest $j$. Thus every integer abscissa
on the lower Newton boundary has an integer ordinate achieved by an
actual monomial generator. Its successive slopes are $-a,\ldots,-1$,
each with positive integral horizontal length $E_a,\ldots,E_1$.
The integral lattice points above this boundary are exactly the
monomials of the ideal. Replacing $E_j$ by $nE_j$ proves the same
assertion for every power. The half-space description of that boundary
is $iu+v\geq\eta_i$; adjoining the degree variable proves
\eqref{eq:rees-semigroup-v169}. A semigroup cut out by these integral
linear inequalities is saturated in its group. Its semigroup algebra
is normal. The compact facets have primitive inward normals $(i,1)$,
and their positive lengths show that none is redundant. They give
exactly the exceptional prime divisors of the normalized blow-up.
The shift by $b^{nC}$ proves the raw formula. The raw ideal also records
the nonexceptional strict transform of $b=0$; this is its content, not
an additional exceptional ray.
\end{proof}

\begin{proof}[Proof of the surface assertions in Theorem~\ref{thm:chain-model-v169}]
The graph Rees lemma and removal of $b^C$ give the first equality in
\eqref{eq:chain-blowup-v169}. The normal fan of the computed Newton
polygon has, in order, the rays
\[
 (1,0),(a,1),(a-1,1),\ldots,(1,1),(0,1).
\]
Every adjacent determinant is one. The same fan is obtained when all
$E_j$ are replaced by one, proving the second equality and smoothness.
Successive subdivisions insert $(1,1),(2,1),\ldots,(a,1)$ and give
the stated sequence of point blow-ups. Equivalently, a product of
positive powers of nonzero ideals is invertible on an integral test
scheme exactly when every factor is invertible as a fractional ideal;
the blow-up universal property gives the same identification.
\end{proof}

There are $a+1$ affine charts. The first and last are
\[
 T_0=\Spec k[b,k],\quad c=bk;
 \qquad T_a=\Spec k[c,z],\quad b=c^az.
\]
For $1\leq i<a$ the intermediate chart is
\begin{equation}\label{eq:chain-charts-v169}
 T_i=\Spec k[e,h],\qquad c=eh,\quad b=e^{i+1}h^i,
 \qquad e=b/c^i,\quad h=c^{i+1}/b.
\end{equation}
Adjacent transition formulas are
\begin{equation}\label{eq:chain-transitions-v169}
 \begin{gathered}
 k=e^{-1},\quad b=e^2h \quad(T_0,T_1),\\
 e_{i+1}=h_i^{-1},\quad h_{i+1}=e_i h_i^2
                  \quad(1\leq i<a-1),\\
 c=e_{a-1}h_{a-1},\quad z=h_{a-1}^{-1}
                  \quad(T_{a-1},T_a).
 \end{gathered}
\end{equation}
All further overlaps are compositions of these formulas.

\subsection{Low-degree equations for the universal curve}
In the source chart $s=1$, put $x=t/s$. Target coordinates $F,G,R$
remain homogeneous. On $T_0$ use
\begin{equation}\label{eq:curve-first-v169}
 R-kx^{a-1}G,\qquad x^aG-bF.
\end{equation}
On $T_i$, $1\leq i<a$, use
\begin{equation}\label{eq:curve-middle-v169}
 \begin{split}
 A&=xR-ehF,\\
 B_j&=x^{a-j}F^{j-1}G-e^{i+1-j}h^{i-j}R^j\quad(1\leq j\leq i),\\
 C_i&=R^{i+1}-h x^{a-i-1}F^iG.
 \end{split}
\end{equation}
On $T_a$ use
\begin{equation}\label{eq:curve-last-v169}
 A=xR-cF,\qquad
 B_j=x^{a-j}F^{j-1}G-zc^{a-j}R^j\quad(1\leq j\leq a).
\end{equation}
The ideal sheaves on $X$ are obtained by bihomogenizing these equations
and saturating by the source and target irrelevant ideals. In the
first chart a convenient global list adds $tR-bksF$ to the
bihomogenizations of \eqref{eq:curve-first-v169}; this prevents a
spurious component at $s=0$. In the other charts the listed
bihomogenizations already give the following two-source-chart cover.

\begin{proposition}[Flatness, exhaustion, and overlaps]\label{prop:curve-charts-v169}
Equations \eqref{eq:curve-first-v169}--\eqref{eq:curve-last-v169} and
\eqref{eq:chain-transitions-v169} define the flat universal embedded
curve over $T_a$. They cover all of that curve, including source
infinity. Each intermediate chart has a monic Gr\"obner basis over
$k[e,h]$, and the last chart has one over $k[c,z]$. Thus no
coefficient stratum is excluded by a division of a leading coefficient.
\end{proposition}
\begin{proof}
For the intermediate chart choose positive rational weights
\[
 \operatorname{wt}(F)=\operatorname{wt}(G)=1,\quad
 \operatorname{wt}(x)=\alpha,\quad
 \operatorname{wt}(R)=1+\rho\alpha,
 \qquad \frac{a-i-1}{i+1}<\rho<\frac{a-i}{i},\quad\alpha>0.
\]
Refine by any monomial order. The leading monomials of
\eqref{eq:curve-middle-v169} are respectively
\[
 xR,\qquad x^{a-j}F^{j-1}G\ (1\leq j\leq i),\qquad R^{i+1}.
\]
Buchberger reductions are uniform in the coefficient ring. The
$S$-polynomial of $A,B_j$ is $-ehB_{j+1}$ for $j<i$, and is $eC_i$
for $j=i$. The one for $A,C_i$ is $hFB_i$. For $j<l$, that of
$B_j,B_l$ is
\[
 e^{i+1-l}h^{i-l}R^j
       \big((xR)^{l-j}-(ehF)^{l-j}\big),
\]
and reduces by $A$. The remaining pairs have relatively prime leading
monomials. This proves the monic basis assertion and a free standard
monomial basis over $k[e,h]$.

For the last chart use lexicographic order $x>F>G>R$. The same
reductions replace $eh$ by $c$ and the coefficient of $R^j$ by
$zc^{a-j}$; the pair $A,B_a$ has relatively prime leading monomials.
The first chart eliminates $R$ and leaves the monic hypersurface
$x^aG-bF$. These bases make the target-graded coordinate rings flat
over their parameter rings. Localizing in a target coordinate and
taking the degree-zero part preserves flatness.

For completeness put $y=s/t$ on the other source chart. On every
parameter chart the bihomogeneous equations reduce to
\begin{equation}\label{eq:source-infinity-v169}
 G=b y^aF,\qquad R=c yF.
\end{equation}
For example $A$ gives the second equation, and $B_1$ gives the first;
the other equations then vanish. Here $F$ is a unit on the projective
target. This is the entire family near source infinity and is flat.
On $s,t\ne0$ the two displayed presentations agree by substitution.

Over the dense parameter torus all the equations are the graph of
\eqref{eq:marked-family-v169}. Their flat coordinate algebras are
torsion-free over the integral parameter ring, so each is the
schematic closure of that generic graph. Substitution of
\eqref{eq:chain-transitions-v169} therefore identifies the ideal
sheaves on each overlap, not merely their reduced supports: both are
the same closure in the fixed $X$. The substitutions of $b,c$ and the
identity on $[F:G:R]$ are explicit, so no target automorphism is hidden
in this argument. The glued family is projective, flat, and has generic
polynomial $P_a$; connectedness gives that polynomial on every fibre.
Its Hilbert map agrees with the defining map of $T_a$ on a dense open.
The graph identification proves that it is the universal curve.
\end{proof}

\begin{proposition}[Boundary geometry and contractions]\label{prop:boundary-v169}
The scheme fibre of $T_a\to S_a$ at the origin is the reduced chain
$E_1\cup\cdots\cup E_a$, with
\[
 E_i^2=-2\ (i<a),\qquad E_a^2=-1,\qquad
 E_iE_{i+1}=1,
 \qquad K_{T_a/S_a}=\sum_{i=1}^a iE_i.
\]
If $B',C'$ are the strict transforms of the two axes, then
\[
 \operatorname{div}(b)=B'+\sum iE_i,\qquad
 \operatorname{div}(c)=C'+\sum E_i.
\]
The Picard group is freely generated by the exceptional classes. A
relative divisor $D=\sum d_iE_i$ is nef precisely when
\[
 d_{i-1}-2d_i+d_{i+1}\geq0\ (1\leq i<a),\qquad
 d_{a-1}-d_a\geq0,\qquad d_0=0.
\]
Deleting rays gives the toric contraction morphisms. In particular,
contracting the intermediate chain between $(i,1)$ and $(j,1)$ gives
an $A_{j-i-1}$ singularity when $j-i>1$.
\end{proposition}
\begin{proof}
In an intermediate chart $(b,c)=(e^{i+1}h^i,eh)=(eh)$; at the end
charts its pullback is $(b)$ or $(c)$. Thus the fibre is reduced, with
ordinary crossings. The self-intersections follow either from the
successive point blow-ups or the relation between adjacent primitive
rays. The Jacobian in \eqref{eq:chain-charts-v169} is
$e^{i+1}h^i$, and the axis orders are the stated valuations. Each
point blow-up adds a free exceptional generator to the Picard group
of the preceding smooth surface, starting with $\Pic(\A^2)=0$.
The only contracted complete curves are the $E_i$, so intersecting
with them gives the nef inequalities. Coarsening a fan gives a proper
birational toric morphism. The two rays $(i,1),(j,1)$ have determinant
$j-i$ and lie at height one; the resulting affine surface is the
cyclic quotient, equivalently the hypersurface $uv=w^{j-i}$.
\end{proof}
The unweighted ideal $K_a$ is degree-independent and has exactly the
required Newton rays, but is not an inclusion-minimal graph ideal.
The powers $K_a^n$ already preclude such a minimality assertion. Its
boundary chain has a modular meaning for the marked family: it is the
universal modification on integral tests on which every $(b,c^j)$
becomes invertible. This is not a proposed intrinsic boundary complex
for all pencils of all contact types.

\section{Embedded limits and the punctual genus correction}
\label{sec:embedded-limits-v169}
The parameter fibre of $T_a$ is a reduced chain; its universal curves
need not be reduced or Cohen--Macaulay. These are different assertions
about different schemes. We now determine both the one-dimensional
parts and the embedded zero-dimensional structure of every curve over
that chain.

\subsection{All formal arcs in the marked family}
Let $b=\beta\tau^p$, $c=\gamma\tau^q$, with $p,q>0$ and
$\beta,\gamma\in k[[\tau]]^*$. The units may be arbitrary formal
series. The following ideals are written on $s=1$ with homogeneous
target coordinates. At source infinity they are glued to $G=R=0$;
equivalently, take the ideal sheaves specified by
Proposition~\ref{prop:curve-charts-v169}.

\begin{theorem}[The complete arc table]\label{thm:arc-table-v169}
There are $a$ walls $p=jq$, $1\leq j\leq a$, and $a+1$ open chambers.
The special embedded ideals are as follows:
\begin{align}
 \mathcal O_0\ (p<q):\quad& (R,x^aG),\label{eq:limit-O0-v169}\\
 \mathcal W_1\ (p=q):\quad&
       (R-\kappa x^{a-1}G,x^aG),\quad
       \kappa=\overline\gamma/\overline\beta,\label{eq:limit-W1-v169}\\
 \mathcal O_i\ (iq<p<(i+1)q):\quad&
       (xR,R^{i+1},x^{a-j}F^{j-1}G\ (1\leq j\leq i)),
       \quad1\leq i<a,\label{eq:limit-Oi-v169}\\
 \mathcal W_j\ (p=jq):\quad&
       (xR,x^{a-r}F^{r-1}G\ (1\leq r<j),
                   R^j-\lambda x^{a-j}F^{j-1}G),
       \quad2\leq j\leq a,\label{eq:limit-Wj-v169}\\
 \mathcal O_a\ (p>aq):\quad&
       (xR,x^{a-j}F^{j-1}G\ (1\leq j\leq a)).
       \label{eq:limit-Oa-v169}
\end{align}
In \eqref{eq:limit-Wj-v169}, $\lambda=\overline\gamma^j/\overline\beta$.
Each wall is a one-dimensional family with nonzero residue parameter;
each open chamber gives one fixed embedded curve. The closure of the
$j$th wall family is $E_j\simeq\Pj^1$, with endpoints
$\mathcal O_{j-1},\mathcal O_j$. These are all components, points, and
adjacencies of the parameter fibre of $T_a$. There is only one
normalization branch at every point of this smooth parameter surface.
The fan is minimal for these fixed-target embedded limits.
\end{theorem}
\begin{proof}
In the first chart $k=c/b$ has nonnegative valuation exactly when
$p\leq q$. Its residue is zero in the first chamber and is $\kappa$
on the first wall. In $T_i$ the two coordinates have valuations
$p-iq$ and $(i+1)q-p$. Setting their positive-valuation coordinates
to zero in \eqref{eq:curve-middle-v169} gives
\eqref{eq:limit-Oi-v169}. Setting $e=0$ and retaining
$h=\lambda\ne0$ gives the next wall; the lower wall is its expression
in the adjacent chart. In the last chart $z=b/c^a$ has nonnegative
valuation exactly when $p\geq aq$. Its nonzero residue gives
$\mathcal W_a$ with $\lambda=\overline z^{-1}$; $z=0$ gives
\eqref{eq:limit-Oa-v169}. Properness and the chart cover exhaust all
arcs. If $c$ is identically zero, the first chart gives
\eqref{eq:limit-O0-v169} as well.

Every nonzero wall residue is attained by an arc with $p=jq$.
The transition formulas identify its two endpoints as stated. The
Newton facets in Proposition~\ref{prop:rees-v169} are all present;
equivalently, the residue term in \eqref{eq:limit-W1-v169} or
\eqref{eq:limit-Wj-v169} changes the ideal sheaf as the residue varies.
A wall cannot therefore be removed while retaining the embedded limit
map. This minimal fan belongs to the marked two-parameter problem,
not to an arbitrary coefficient presentation on $B_a$.
\end{proof}

\begin{theorem}[A sharp, degree-independent jet bound on the family]
\label{thm:sharp-jet-v169}
Let a formal arc in $S_a$ be centred at the origin and satisfy
$\operatorname{ord}(b)=p<\infty$. Its closed embedded Hilbert limit is
determined by the coefficient arc modulo $\tau^{p+1}$. This bound
is independent of $a$ and of the Hilbert embedding degree. For every
$p\geq1$ it is optimal as a uniform sufficient bound over arcs with
that value of $p$.
\end{theorem}
\begin{proof}
If $\operatorname{ord}(c)>p$, including $c=0$, any congruent arc is
still in $\mathcal O_0$. Otherwise congruence preserves $p,q$ and the
leading coefficients of both $b,c$. It preserves every wall residue
in Theorem~\ref{thm:arc-table-v169}, which proves sufficiency.
For necessity take $c=\tau^p$ and $b=\beta\tau^p$ for two different
nonzero constants $\beta$. The arcs agree modulo $\tau^p$, but the
first-wall ideals differ: on $G=1$, the nonzero class $x^{a-1}$ has
$R=\beta^{-1}x^{a-1}$ and $x^a=0$. Thus that jet does not determine
the fixed-target limit.
\end{proof}
Here $p$ is the intersection order with the marked base-point divisor
$b=0$. It is not the raw Smith length. In degree $m$ the latter is
\begin{equation}\label{eq:raw-order-v169}
 \mu=Cp+\sum_{j=1}^aE_j\min(p,jq),
\end{equation}
whereas the primitive degree-$m$ ideal has order given by the sum
without $Cp$, and the unweighted graph ideal has order
$\sum_{j=1}^a\min(p,jq)$. These are three distinct quantities, none
of which should be substituted for the optimal family bound $p+1$.
For $a=3,m=7$, arcs $(b,c)=(\tau,\tau^2)$ and
$(\tau^2,\tau^3)$ have the same special Hilbert curve but raw Smith
lengths $84$ and $168$. Ramifying an arc multiplies all ideal orders,
while leaving its embedded closed limit unchanged.

\subsection{Local rings and associated points}
Let $L_0$ denote the main component $G=R=0$. Let $L$ denote the tail
$x=R=0$, and put
\[
 P=(x=0,[F:G:R]=[0:1:0]).
\]
All local ideals below are ordinary ideals in the indicated affine
rings, before passage to sheaves. Write $H^0_P(\OO_C)$ for the
largest submodule supported at the closed point $P$, not for the
whole nilradical of a curve with a generically thick tail.

\begin{theorem}[Pure tails and punctual excess]\label{thm:punctual-v169}
For $\mathcal O_0$ and $\mathcal W_1$ the curve is Cohen--Macaulay,
with components $L_0,L$ of generic multiplicities $1,a$.
For $1\leq i<a$, the curve $\mathcal O_i$ has the same two minimal
components and multiplicities. On $F=1$ its ideal is
\begin{equation}\label{eq:Oi-F-v169}
 (xR,x^{a-i}G,R^{i+1})
    =(G,R)\cap(xR,x^{a-i},R^{i+1}).
\end{equation}
On $G=1$ it is
\begin{equation}\label{eq:Oi-G-v169}
 (xR,R^{i+1},x^{a-i}(x,F)^{i-1}).
\end{equation}
There is precisely one embedded associated point when $i\geq2$, namely
$P$, and none when $i=1$. Its punctual module is
\begin{equation}\label{eq:Oi-torsion-v169}
 H^0_P(\OO_{\mathcal O_i})
   \simeq x^{a-i}k[x,F]\big/x^{a-i}(x,F)^{i-1},\qquad R\cdot H^0_P=0,
 \qquad\length H^0_P=\binom i2.
\end{equation}

For $2\leq j\leq a$ put $d=a-j$. The wall curve $\mathcal W_j$ has
on $F=1$ the complete-intersection ideal
\begin{equation}\label{eq:W-F-v169}
 (xR,R^j-\lambda x^dG).
\end{equation}
On $G=1$ its ideal is
\begin{equation}\label{eq:W-G-v169}
 (xR,x^{d+1}(x,F)^{j-2},R^j-\lambda x^dF^{j-1}).
\end{equation}
Its pure quotient on this chart is obtained by adding $x^{d+1}$.
The kernel is
\begin{equation}\label{eq:W-torsion-v169}
 H^0_P(\OO_{\mathcal W_j})
  \simeq x^{d+1}k[x,F]\big/x^{d+1}(x,F)^{j-2},\qquad R\cdot H^0_P=0,
 \qquad\length H^0_P=\binom{j-1}{2}.
\end{equation}
For $j<a$ the minimal components are again $L_0,L$ with multiplicities
$1,a$. For $j=a$ the tail is the reduced plane curve
$R^a=\lambda F^{a-1}G$. In both cases the only possible embedded point
is $P$, and it occurs exactly when $j\geq3$.
\end{theorem}
\begin{proof}
Eliminating $R$ in the first two cases leaves a hypersurface, and gives
the multiplicities directly. For $\mathcal O_i$, localization at $F$
reduces all the monomials in \eqref{eq:limit-Oi-v169} to
\eqref{eq:Oi-F-v169}. The intersection identity follows by monomial
divisibility. The second component is primary to $(x,R)$ and is free
of rank $a$ over $k[G]$, with transverse basis
\[
 1,x,\ldots,x^{a-i-1},R,\ldots,R^i.
\]
The entire ring in \eqref{eq:Oi-F-v169} has a length-two
Hilbert--Burch resolution; alternatively its three monomial generators
have the two syzygies coming from $xR^{i+1}$ and $x^{a-i}GR$.
It is one-dimensional Cohen--Macaulay and has no embedded point.

On $G=1$ the ideal is \eqref{eq:Oi-G-v169}. Adding $x^{a-i}$ gives
the pure primary ideal $(xR,R^{i+1},x^{a-i})$. The kernel is generated
by the image of $x^{a-i}$, with annihilator
$(R,(x,F)^{i-1})$ on that generator. This proves
\eqref{eq:Oi-torsion-v169}. For $i\geq2$ an explicit primary
decomposition is
\[
 (xR,R^{i+1},x^{a-i})
 \ \cap\ (xR,R^{i+1},(x,F)^{a-1}).
\]
The second factor is primary to $(x,F,R)$. The displayed length is
the number of monomials in two variables of total degree below $i-1$.

For a wall, the equations reduce to \eqref{eq:W-F-v169} and
\eqref{eq:W-G-v169}. The former is a regular sequence of length two.
In the latter, after adjoining $x^{d+1}$, there is a free $k[F]$-basis
\[
 1,x,\ldots,x^d,R,\ldots,R^{j-1};
\]
the relation $R^j=\lambda F^{j-1}x^d$ and $xR=0$ give its
multiplication. Hence this pure ring is Cohen--Macaulay of rank $a$
over $k[F]$. The kernel of the pure quotient is generated by
$x^{d+1}$ and has exactly the annihilator described in
\eqref{eq:W-torsion-v169}. No extra relation arises from multiplying
$R^j-\lambda x^dF^{j-1}$ by $x$: the resulting
$x^{d+1}F^{j-1}$ already lies in $x^{d+1}(x,F)^{j-2}$.
A monomial reduction in $x,F$, followed by reduction in $R$, proves
that these are all the relations. This also computes the length.
For $d>0$ the pure ideal is primary to $(x,R)$; for $d=0$ it is the
reduced cuspidal tail ideal $(x,R^a-\lambda F^{a-1})$.

There are no missing associated points. Away from $P$ the $F$ chart
or its overlap has just been treated. On $R=1$ the nonterminal cases
have empty special fibre, while on the last wall $F,G$ are units and
\eqref{eq:W-F-v169} applies. Source infinity is
\eqref{eq:source-infinity-v169}. These charts cover the curve.
\end{proof}

\begin{theorem}[The cuspidal wall and its nilradical]\label{thm:genus-correction-v169}
On the last wall, $\mathcal W_a$ consists of its reduced curve
$L_0\cup Q_a$ together with a punctual nilradical at $P$, where
\[
 Q_a:\ x=0,\quad R^a=\lambda F^{a-1}G.
\]
The normalization of $Q_a$ is $\Pj^1$, with map
\begin{equation}\label{eq:cusp-normalization-v169}
 [u:v]\longmapsto[u^a:\lambda^{-1}v^a:u^{a-1}v].
\end{equation}
The main component meets $Q_a$ transversely at its smooth point
$[1:0:0]$. Put $g_a=(a-1)(a-2)/2$. The nilradical $\mathcal N$ has
length $g_a$ and, for $a\geq3$, nilpotence order $a-1$. More precisely,
on $G=1$ it is the ideal generated by $x$ with
\begin{equation}\label{eq:genus-module-v169}
 \mathcal N\simeq xk[x,F]/x(x,F)^{a-2},\quad R\mathcal N=0,
 \qquad
 \length(\mathcal N^r)=\binom{a-r}{2}\quad(1\leq r\leq a-2).
\end{equation}
For $a=2$ it is zero. The Hilbert polynomial of the reduced curve is
$(a+1)l+1-g_a$; this exact nilradical restores $P_a(l)=(a+1)l+1$.
Thus the punctual part is forced by arithmetic genus, not removable
excess in a flat Hilbert limit.
\end{theorem}
\begin{proof}
Set $j=a,d=0$ in \eqref{eq:W-G-v169}. The ring is
\[
 k[x,F,R]/(xR,x(x,F)^{a-2},R^a-\lambda F^{a-1}).
\]
Its reduction is the domain $k[F,R]/(R^a-\lambda F^{a-1})$; the
nilradical is exactly $(x)$. The formula for its module follows from
Theorem~\ref{thm:punctual-v169}. A basis of its $r$th power consists
of $x^uF^v$ with $u\geq r$ and $u+v\leq a-2$. Counting these
monomials gives \eqref{eq:genus-module-v169} and proves the
nilpotence statement, including nonvanishing of $x^{a-2}$.
The parametrization \eqref{eq:cusp-normalization-v169} is birational;
on the cusp chart its local semigroup is generated by $a-1,a$.
Its missing nonnegative exponents number $g_a$. The only singularity
of the plane tail is that cusp, while at the attaching point $F=1$
the full reduced local ring is $k[x,R]/(xR)$ after eliminating $G$.
The normalization of the whole reduction is consequently two copies
of $\Pj^1$. A plane curve of degree $a$ has arithmetic genus $g_a$;
attaching the main line at one point gives the stated reduced Hilbert
polynomial. The exact sequence
$0\to\mathcal N\to\OO_{\mathcal W_a}\to\OO_{(\mathcal W_a)_{\rm red}}\to0$
then proves the correction formula.
\end{proof}

\begin{proposition}[The terminal chamber]\label{prop:terminal-v169}
The curve $\mathcal O_a$ has three minimal components: $L_0$ with
multiplicity one, the tail $x=G=0$ with multiplicity one, and the
tail $x=F=0$ with multiplicity $a-1$. Their reduced supports form a
chain of three lines. On $F=1$ the ideal is $(G,xR)$; on $G=1$ it is
\[
 (xR,(x,F)^{a-1}).
\]
The latter has pure quotient $(x,F^{a-1})$ and punctual kernel
$xk[x,F]/x(x,F)^{a-2}$ of length $g_a$. On $R=1$ the ideal is
$(x,F^{a-1}G)$, a hypersurface with no embedded point. Thus $P$ is
the only embedded associated point when $a\geq3$, and there is none
when $a=2$. The full nilradical is not just this punctual module:
the $F=0$ tail is generically nonreduced for $a>2$.
\end{proposition}
\begin{proof}
Localize \eqref{eq:limit-Oa-v169} in each target coordinate. On $G=1$
the monomials $x^{a-j}F^{j-1}$ generate $(x,F)^{a-1}$.
Adding $x$ yields the pure ring, free of rank $a-1$ over $k[R]$;
the kernel has the stated two-variable basis and is killed by $R$.
On $R=1$, $xR$ eliminates $x$ and leaves $F^{a-1}G$.
These ordinary affine calculations give the minimal multiplicities
and all associated points. The normalization of the reduction is the
disjoint union of its three lines. The pure plane tail
$F^{a-1}G=0$ has arithmetic genus $g_a$, so the same length correction
also recovers $P_a$ in this terminal chamber.
\end{proof}

\subsection{The first higher-contact instance}
For $a=3$ the Gotzmann degree is $7$, and
\begin{equation}\label{eq:cubic-factor-v169}
 I_{29}(M_{3,7})\OO_{S_3}
       =b^{56}(b,c)^7(b,c^2)^6(b,c^3)^{15}.
\end{equation}
The parameter fibre has three components, with self-intersections
$-2,-2,-1$. Every one is one-dimensional, their successive
intersections are single reduced points, and there are no other
intersections or normalization branches. The seven embedded curve
families are collected here; each ideal is on $s=1$.
\begin{center}
\small
\begin{tabular}{c|l|c}
region & ideal & $\length H^0_P$\\ \hline
$p<q$ & $(R,x^3G)$ & $0$\\
$p=q$ & $(R-\kappa x^2G,x^3G)$ & $0$\\
$q<p<2q$ & $(xR,x^2G,R^2)$ & $0$\\
$p=2q$ & $(xR,x^2G,R^2-\lambda xFG)$ & $0$\\
$2q<p<3q$ & $(xR,x^2G,xFG,R^3)$ & $1$\\
$p=3q$ & $(xR,x^2G,xFG,R^3-\lambda F^2G)$ & $1$\\
$p>3q$ & $(xR,x^2G,xFG,F^2G)$ & $1$
\end{tabular}
\end{center}
At the last wall this length-one module is square-zero and is supported
at the cusp of the cubic tail. At the preceding open chamber it is
an embedded point on a generically triple line. These are different
local algebras with the same required Euler-characteristic correction.
For $a\geq4$ the last-wall nilradical is no longer square-zero, as
\eqref{eq:genus-module-v169} shows. This is why the first-contact
extension algebra alone cannot describe higher contacts.

The entire $S_3$ parameter fibre has been classified above, including
its universal curves. The full $B_3$ parameter fibre has more
coefficient directions, and no equality with this three-component
chain is asserted. The distinction remains important when the
marked family is inserted into a local Smith normal form of a pencil.

\section{Ramification and the complete parameter boundary}
\label{sec:ramified-boundary-v170}
Let $a\geq2$ and retain the marked family and its Hilbert graph $T_a$
from Theorem~\ref{thm:chain-model-v169}. We work over an algebraically
closed field $k$ of characteristic zero. For positive integers
$\rho,\sigma$ consider the finite flat coefficient map
\[
 S'=\Spec k[u,v]\longrightarrow S_a,\qquad b=u^\rho,\quad c=v^\sigma.
\]
Its pullback family is $[t^a:u^\rho s^a:v^\sigma st^{a-1}]$. Put
\[
 Y=T_a\times_{S_a}S',\qquad Z=Y^\nu,\qquad
 \nu:Z\longrightarrow Y,\quad \pi:Z\longrightarrow S'.
\]
Here $Y$ is the retained-coefficient Hilbert graph, whereas $Z$ is its
normalization. The closed fibre $F=\pi^{-1}(0)$ is a scheme of
\emph{parameters}. It is not one of the embedded curves parameterized
by the Hilbert graph. This distinction matters: the conductor of
$\nu$, the nilradical of $F$, and the punctual nilradical of a universal
curve fibre are three different objects.

For $1\leq i\leq a$ define
\begin{equation}\label{eq:ramification-data-v170}
 g_i=\gcd(\rho,i\sigma),\quad
 n_i=(A_i,B_i)=\left(\frac{i\sigma}{g_i},\frac{\rho}{g_i}\right),
 \quad e_i=\frac{\rho\sigma}{g_i},\quad
 \ell_i=\min(A_i,B_i).
\end{equation}
Also set $n_0=(0,1)$ and $n_{a+1}=(1,0)$, ordered clockwise, and
$\Delta_i=|\det(n_i,n_{i+1})|$ for $0\leq i\leq a$.

\begin{theorem}[Ramified Hilbert boundary]\label{thm:ramified-boundary-v170}
The following assertions hold for every $a,\rho,\sigma$ as above.
\begin{enumerate}
\item The scheme $Y$ is integral, is a local complete intersection,
and is the full schematic Hilbert graph of the pulled-back family.
Its normalization $Z$ is the normal toric surface with maximal cones
$\langle n_i,n_{i+1}\rangle$, $0\leq i\leq a$. In particular its
exceptional prime divisors $D_1,\ldots,D_a$ are copies of $\Pj^1$.
\item The fibre $F$ is Cohen--Macaulay of pure dimension one, with
no embedded associated points. Its reduction is a nodal chain of
$a$ projective lines, and its generic length along $D_i$ is $\ell_i$.
Its nilradical $\mathcal N_F$ satisfies
\[
 \mathcal N_F^{L}=0,\qquad L=\max_i\ell_i;
 \quad \mathcal N_F^{L-1}\ne0\quad\text{if }L>1.
\]
It is reduced exactly when, for every $i$, either $\rho\mid i\sigma$
or $i\sigma\mid\rho$.
\item On the open orbit of $D_i$, the normalization maps to the
$i$th exceptional curve of $Y$ with residue degree $g_i$. The finite
map $Z\to T_a$ has ramification index $e_i$ there. A general closed
point of that exceptional curve of $Y$ has exactly $g_i$ distinct
normalization lifts. All torus-fixed boundary points have a single
normalization lift.
\item The conductor, regarded as an ideal on $Z$, is the divisorial
ideal
\begin{equation}\label{eq:conductor-divisor-v170}
 \mathfrak c_{Z/Y}=\OO_Z\left(-\sum_{i=1}^a\kappa_iD_i\right),
 \qquad
 \kappa_i=\frac{\rho i\sigma-\rho-i\sigma}{g_i}+1.
\end{equation}
The notation denotes a reflexive ideal and does not assert that the
divisor is Cartier. There is no additional condition supported only
at a boundary crossing. In particular, $Y$ is normal if and only if
$\rho=1$.
\end{enumerate}
The theorem concerns the complete boundary of this ramified marked
family, not the complete fibre of the full coefficient graph over
$B_a$. None of its multiplicities depends on a choice of Hilbert
embedding degree.
\end{theorem}

\subsection{The normalized Rees algebra}
For $m\geq m_a$ let $E_j$ and $C$ be the integers of
Theorem~\ref{thm:chain-model-v169}, and put
\[
 I=\prod_{j=1}^a(u^\rho,v^{j\sigma})^{E_j},\qquad
 \eta_i=\sum_{j=1}^a E_j\min(i,j).
\]
The raw evaluation ideal on $S'$ is $u^{\rho C}I$.

\begin{proposition}[All powers and all normalization charts]
\label{prop:ramified-rees-v170}
The integral closure of the Rees algebra of $I$ in $k(u,v)(Q)$ is
\begin{equation}\label{eq:ramified-rees-v170}
 \overline{\mathcal R(I)}
 =k\left[u^\alpha v^\beta Q^n:
 \begin{array}{l}
 \alpha,\beta,n\in\mathbb Z_{\geq0},\\
 i\sigma\alpha+\rho\beta\geq n\rho\sigma\eta_i
                  \quad(1\leq i\leq a)
 \end{array}\right].
\end{equation}
Thus its degree-$n$ part is $\overline{I^n}Q^n$, including $n=0$.
For the raw ideal one replaces $\alpha\geq0$ by
$\alpha\geq n\rho C$ and the right side of the $i$th inequality by
$n\rho\sigma(\eta_i+iC)$. The affine chart for the cone
$\tau_i=\langle n_i,n_{i+1}\rangle$ is exactly
\begin{equation}\label{eq:normal-chart-v170}
 Z_i=\Spec k[\tau_i^\vee\cap\mathbb Z^2].
\end{equation}
Adjacent charts glue by inverting the characters of their shared
face in the common Laurent ring $k[u^{\pm1},v^{\pm1}]$.
\end{proposition}
\begin{proof}
The coefficient map is free with basis $u^jv^k$, $0\leq j<\rho$,
$0\leq k<\sigma$. Flat base change of the Rees algebra consequently
gives $Y=\operatorname{Bl}_I S'$; see the classical flat-base-change
lemma for blow-ups \cite{Stacks169}. The open $uv\ne0$ is
schematically dense by flat base change from $T_a$. Its coordinate
ring is a domain, so $Y$ is integral. The same density identifies
$Y$ with the closure of the pulled-back Hilbert section. In particular
no separate vertical component has been discarded in this equality.

The Newton polygon of $I$ is the image, under
$(x,y)\mapsto(\rho x,\sigma y)$, of the polygon calculated in
Proposition~\ref{prop:rees-v169}. Its compact supporting inequalities
are precisely those in \eqref{eq:ramified-rees-v170}. All compact
edges have positive length. The semigroup of the Rees algebra
contains the two degree-zero basis vectors and a degree-one vector,
so its group is the whole $\mathbb Z^3$. Saturating that semigroup
therefore gives the integral lattice points in its rational cone.
Indeed, a rational positive combination becomes an element of the
original semigroup after clearing denominators; its monomial is
integral. Conversely every supporting valuation is nonnegative on
an integral element. The saturated semigroup ring is normal. This
is the standard semigroup-normalization theorem
\cite[Proposition~5 and Section~2.6]{GonzalezTeissier170}, here applied
to the explicit polygon. It proves the formula in every degree.

The normal fan is unchanged by the positive edge lengths $E_j$.
Its primitive compact rays are exactly $n_i$, giving
\eqref{eq:normal-chart-v170}. Normalization commutes with localization,
so the stated overlaps give the global normalization, not unrelated
normalizations of completed rings. The shift by $u^{n\rho C}$ proves
the raw formula and does not change Proj.
\end{proof}

\begin{proposition}[Singularities and divisor data]
\label{prop:ramified-singularities-v170}
The singularities of $Z$ are precisely those torus-fixed points whose
corresponding determinant in the following list is greater than one:
\begin{equation}\label{eq:cone-indices-v170}
 \Delta_0=\frac{\sigma}{g_1},\qquad
 \Delta_i=\frac{\rho\sigma}{g_i g_{i+1}}\ (1\leq i<a),\qquad
 \Delta_a=\frac{\rho}{g_a}.
\end{equation}
The affine chart has cyclic divisor class group
$\mathbb Z/\Delta_i\mathbb Z$. Its cyclic quotient type is determined
without ambiguity by the two primitive rays in
\eqref{eq:normal-chart-v170}: orient them positively, extend the
first to a lattice basis, and write the second as $(p,\Delta_i)$
in that basis. The chart is $\A^2/\mu_{\Delta_i}$ with weights
$(-p,1)$, up to swapping the coordinates. For $\Delta_i=1$ it is smooth.
Moreover
\[
 \operatorname{div}(u)=U'+\sum_i A_iD_i,
 \qquad \operatorname{div}(v)=V'+\sum_i B_iD_i,
 \qquad K_{Z/S'}=\sum_i(A_i+B_i-1)D_i,
\]
where $U',V'$ are the strict transforms of the axes. Rational
intersection numbers are
\[
 D_iD_{i+1}=\Delta_i^{-1},\qquad
 D_i^2=-\frac{|\det(n_{i-1},n_{i+1})|}{\Delta_{i-1}\Delta_i},
 \qquad D_iD_j=0\quad(|i-j|>1).
\]
\end{proposition}
\begin{proof}
The determinant formulas follow by substitution. For a cone generated
by primitive vectors $r_1,r_2$, its ray sublattice has index
$|\det(r_1,r_2)|$. The corresponding toric chart in that sublattice
is $\A^2$. The lattice quotient is cyclic because the first ray is
primitive; in the indicated basis its generator acts with weights
$(-p,1)$. Primitivity of the second ray implies
$\gcd(p,\Delta_i)=1$, so there are no pseudoreflections and the
quotient is smooth exactly for index one. The divisor class group
is the cokernel of the map from the character lattice to the free
group on the two rays; its Smith form is $(1,\Delta_i)$.
Normal surface charts are regular away from their closed toric point.

A character has order its pairing with a primitive ray. This proves
the two divisor formulas. The torus form
$(du/u)\wedge(dv/v)$ has simple poles on each invariant divisor,
which proves the relative canonical formula. The two adjacent
invariant divisors meet with rational intersection $1/\Delta_i$:
pull back to the finite affine cover, where the axes meet with
intersection one and the covering degree is $\Delta_i$. Intersecting
the principal-character relations with $D_i$ then gives the displayed
self-intersection. These calculations also specify a toric resolution
by subdivision into unimodular cones; its curves and intersection
numbers are obtained by the same determinant rule.
\end{proof}

\subsection{The scheme fibre, not only its cycle}
The following elementary lattice observation rules out hidden
zero-dimensional embedded components in the parameter fibre.

\begin{lemma}[The two coordinate ideal on a subdivided surface]
\label{lem:coordinate-ideal-v170}
Let $\tau$ be a two-dimensional cone inside the first quadrant with
primitive boundary vectors $l,h$, and suppose $\tau$ is not the
whole first quadrant. In $A=k[\tau^\vee\cap\mathbb Z^2]$ one has
\begin{equation}\label{eq:fibre-divisorial-v170}
 (u,v)=\left\langle u^xv^y:
 \langle l,(x,y)\rangle\geq\min(l_1,l_2),\quad
 \langle h,(x,y)\rangle\geq\min(h_1,h_2)\right\rangle_k.
\end{equation}
The exponents may be negative; the inequalities imply membership in
$A$. The ideal is divisorial, and $A/(u,v)$ is Cohen--Macaulay of
dimension one whenever its support is nonempty.
\end{lemma}
\begin{proof}
If both rays are on the same side of $(1,1)$, either $u/v$ or $v/u$
is regular on the chart, and the ideal is principal. Otherwise write
$l=(l_1,l_2)$, $h=(h_1,h_2)$ with
$l_1<l_2$ and $h_1>h_2$; equality cases belong to the preceding
argument. For a lattice point satisfying the displayed inequalities,
put $L=x+y$. If $L\leq0$ and $l_1,h_2>0$, the first inequality
forces $y>0$, while the second forces $y<0$. If one of $l_1,h_2$
is zero, its inequality instead forces the weak sign, and the other
strict inequality gives the contradiction. They cannot both vanish,
since the original quadrant was excluded. Thus $L\geq1$.

If neither $(x-1,y)$ nor $(x,y-1)$ belonged to the dual cone, the
only possible failed inequalities would be
\[
 l_1L+(l_2-l_1)y<l_2,
 \qquad h_1L-(h_1-h_2)y<h_1.
\]
The first implies $y\leq0$ and the second $y>0$, since $L\geq1$
and $y$ is integral. This is impossible. Thus the monomial is a
multiple of $u$ or $v$. The reverse inclusion is immediate.

The formula is the intersection of the symbolic valuation ideals
of the invariant height-one primes. All noninvariant height-one
conditions are supplied by membership in the normal ring $A$.
Consequently this rank-one ideal is reflexive. A reflexive module
over a normal surface has depth two at closed points. The exact
sequence with middle term $A$ now gives depth at least one for its
quotient, as asserted.
\end{proof}

\begin{proof}[Proof of the fibre assertions in Theorem~\ref{thm:ramified-boundary-v170}]
Every maximal cone of $Z$ is a proper subcone of the quadrant, so
Lemma~\ref{lem:coordinate-ideal-v170} applies. Globally it says
\begin{equation}\label{eq:complete-fibre-ideal-v170}
 (u,v)\OO_Z=\OO_Z\left(-\sum_i\ell_iD_i\right).
\end{equation}
The fibre thus has no embedded associated points. Its generic length
is the indicated DVR order $\ell_i$. Reducing the ideal replaces
each positive order by one, so the reduction is the chain of
invariant projective lines. At a cone with two exceptional edges,
the reduced union has ring $k[z_1,z_2]/(z_1z_2)$: on the cyclic
quotient cover the two axes have invariant coordinate rings generated
by the respective $\Delta_i$th powers. Thus the crossings really
are nodes, even when the ambient surface is singular.

The radical ideal has order one at every $D_i$. Its $L$th power lies
in \eqref{eq:complete-fibre-ideal-v170} by the divisorial description,
whereas localization at a divisor of length $L$ shows nonvanishing
of the preceding nilradical power if $L>1$. Finally
$\min(\rho,i\sigma)/g_i=1$ is equivalent to one of the two integers
$\rho,i\sigma$ dividing the other. This proves all claims, including
the exact nilpotence order rather than only a bound from a cycle.
\end{proof}

\section{Conductors, normalization lifts, and logarithmic compatibility}
\label{sec:conductors-v170}
We finish the ramified-boundary theorem by calculating the full
conductor ideal. The generic transverse calculation alone would not
exclude a further condition at the intersection of two exceptional
curves; the depth argument below is what excludes it.

\begin{lemma}[A reduced binomial curve]\label{lem:binomial-conductor-v170}
Let $r,q$ be positive integers, let $g=\gcd(r,q)$, and put
$r'=r/g$, $q'=q/g$. The normalization of
$k[[U,V]]/(U^r-V^q)$ is the product of $g$ power-series rings.
On each branch its conductor has order
\[
 c=g r'q'-r'-q'+1,
\]
and the length of the total normalization quotient is
\[
 \delta=\frac{rq-r-q+g}{2}.
\]
The same formulas hold with a nonzero scalar multiplying $V^q$.
\end{lemma}
\begin{proof}
Factor the equation into $g$ distinct factors
$U^{r'}-\zeta V^{q'}$. Each branch is parametrized, after multiplying
one parameter by a nonzero scalar, by
$U=t^{q'}, V=t^{r'}$. Its semigroup is
$\langle r',q'\rangle$. In each residue class modulo $r'$ its
smallest element is one of $j q'$, $0\leq j<r'$. This gives conductor
$(r'-1)(q'-1)$ and exactly $(r'-1)(q'-1)/2$ gaps, including the
cases $r'=1$ or $q'=1$.

A function supported on only one normalization branch belongs to the
original curve ring precisely when its restriction to that branch
is divisible, in the branch ring, by the product of the other factors.
Indeed their ideals in the two-variable regular ring intersect in
their product. That product has order $(g-1)r'q'$ on the selected
branch. To remain in the curve ring after multiplication by every
function of its normalization, a further factor of the branch
conductor is necessary and sufficient. This gives $c$.

Two distinct factors have intersection multiplicity $r'q'$.
The exact sequence for a union of two principal curves, with last
term the intersection ring, shows that their normalization defects
add with this intersection multiplicity. Induction over the $g$
factors yields
$g(r'-1)(q'-1)/2+\binom g2r'q'$, which is the stated $\delta$.
\end{proof}

\begin{lemma}[No isolated conductor condition]
\label{lem:conductor-depth-v170}
Let $A$ be a two-dimensional Cohen--Macaulay domain essentially of finite
type over a field, with finite normalization $\bar A$. Its conductor
$\{z\in\bar A:z\bar A\subset A\}$ is a divisorial ideal of $\bar A$.
It is determined by its orders at the height-one primes of $\bar A$.
\end{lemma}
\begin{proof}
Inside its fraction field one has
$A=\bigcap_{\operatorname{ht}\mathfrak p=1}A_{\mathfrak p}$.
For completeness, a failure of this equality would give a finite
nonzero quotient $(A+Az)/A$ supported in codimension two. The middle
term is torsion-free, hence has depth at least one, while $A$ has
depth two. The depth lemma excludes such a zero-dimensional quotient,
after localizing at a point of its support.

It follows that the condition $z\bar A\subset A$ can be checked
in all $A_{\mathfrak p}$. Their normalizations are semilocal
one-dimensional normal rings, and the conductor there is an ideal
specified by a nonnegative order at each maximal ideal. Finite
normalization identifies these with the height-one primes of
$\bar A$. Intersecting these valuation conditions with $\bar A$
gives exactly the conductor; by definition this is a divisorial
ideal. Thus no extra condition supported only in codimension two
can occur.
\end{proof}

\begin{proof}[Proof of the conductor and lifting assertions in
Theorem~\ref{thm:ramified-boundary-v170}]
The first, intermediate, and last charts of $Y$ are respectively
\begin{equation}\label{eq:raw-cover-charts-v170}
 \begin{gathered}
 k[u,v,\theta]/(v^\sigma-u^\rho\theta),\\
 k[e,h,u,v]/(u^\rho-e^{i+1}h^i,\ v^\sigma-eh)
                  \quad(1\leq i<a),\\
 k[u,v,z]/(u^\rho-v^{a\sigma}z).
 \end{gathered}
\end{equation}
The intermediate rings are free over $k[e,h]$ on the monomials
$u^jv^k$ with $j<\rho$, $k<\sigma$; the equations are a regular
sequence. The end charts are hypersurfaces. Integrality was proved
above. Thus every chart is Cohen--Macaulay and a local complete
intersection, as required in the first assertion of the theorem.

On the open orbit of the old exceptional curve $E_i$, the coordinates
$c$ and $w=b/c^i$ are regular, with $w$ invertible. The pulled-back
chart there has equation
\begin{equation}\label{eq:transverse-binomial-v170}
 u^\rho=w v^{i\sigma}.
\end{equation}
At a closed point $w=w_0\ne0$, a formal root of the unit $w$ removes
it from the equation. Lemma~\ref{lem:binomial-conductor-v170}, with
$r=\rho,q=i\sigma$, gives $g_i$ geometric branches and conductor
order $\kappa_i$. The transverse normalization defect is
\begin{equation}\label{eq:delta-defect-v170}
 \delta_i=\frac{\rho i\sigma-\rho-i\sigma+g_i}{2}.
\end{equation}
These formulas also descend to the generic point: the residue field
extension is generated by the character
\begin{equation}\label{eq:root-residue-v170}
 \zeta_i=\frac{u^{\rho/g_i}}{v^{i\sigma/g_i}},\qquad
 w=\zeta_i^{g_i}.
\end{equation}
The numerator and denominator have equal order on $D_i$, and their
exponent difference is primitive in the perpendicular lattice. Hence
$\zeta_i$ is a coordinate on its open orbit and the residue degree
is exactly $g_i$.

Under the lattice map $\operatorname{diag}(\rho,\sigma)$ the ray
$n_i$ maps to $e_i(i,1)$. This proves the ramification index, with
$e_i g_i=\rho\sigma$, in agreement with the total covering degree.
These are the usual lattice ramification rules for toric covers
\cite[Definition~1 and Lemma~1]{AlexeevPardini170}; here both the
ray and residue character are explicit. The power map
$\zeta_i\mapsto\zeta_i^{g_i}$ has exactly $g_i$ points above a
nonzero closed residue and a single point above each endpoint.
The chart semigroups in \eqref{eq:raw-cover-charts-v170} have the
same cones as their saturations, so each torus-fixed point also has
one preimage. This proves the assertions about lifts on the whole
boundary.

Off the exceptional chain, $Y$ and $Z$ both coincide with $S'$, so
there is no other conductor divisor. Lemma~\ref{lem:conductor-depth-v170}
applied to \eqref{eq:raw-cover-charts-v170} extends the computed
orders across every crossing and proves
\eqref{eq:conductor-divisor-v170}. For $\rho=1$ all orders vanish,
so the conductor is the unit ideal and $Y=Z$. If $\rho>1$, then
$a\sigma\geq2$ and the binomial formula gives $\kappa_a>0$.
Thus $Y$ is not normal. This completes the proof of the theorem.
\end{proof}

\subsection{Equal-power covers: an explicit conductor on every chart}
The conductor theorem includes unequal powers and singular normal
charts. In the equal-power case there is a simpler algebraic form
which also checks the global conductor at crossings directly.

\begin{corollary}[Equal-power conductor]\label{cor:equal-power-v170}
If $\rho=\sigma=d$, the normalized surface $Z$ is the smooth chain
model in the coordinates $u,v$. Its fibre over $0\in S'$ is reduced,
but $\nu$ has $d$ lifts over every nonendpoint point of each old
exceptional curve, and
\[
 \mathfrak c_{Z/Y}=\OO_Z\left(-(d-1)\sum_{i=1}^a iD_i\right),
 \qquad \delta_i=i\binom d2.
\]
On an intermediate normalization chart $k[E,H]$ one has
\[
 u=E^{i+1}H^i,\quad v=EH,\quad e=E^d,\quad h=H^d,
\]
and the conductor is the principal ideal $(u^{d-1})$. On the first
chart it is again $(u^{d-1})$; on the last chart, with
$u=v^a Z_0$, it is $(v^{a(d-1)})$.
\end{corollary}
\begin{proof}
Here $g_i=d$, so $n_i=(i,1)$ and $\ell_i=1$. The theorem proves
all numerical and global statements. To verify the interior ideal
without omitting a possible condition at $E=H=0$, write
\[
 A_i=k[E^d,H^d,E^{i+1}H^i,EH]\subset k[E,H].
\]
The exponent semigroup is the disjoint union
\[
 \coprod_{0\leq j,k<d}
 ((i+1)j+k,ij+k)+d\mathbb Z_{\geq0}^2.
\]
Indeed the monic relations in \eqref{eq:raw-cover-charts-v170}
reduce to this list, whose residue pairs modulo $d$ are distinct.
The exponent $c=((i+1)(d-1),i(d-1))$ plus any nonnegative lattice
point belongs to this union: choose $j$ from the coordinate
difference modulo $d$, then $k$ from the second coordinate modulo
$d$; each representative coordinate is at most the corresponding
coordinate of $c$ plus $d-1$. Congruence modulo $d$ therefore makes
the two remainders nonnegative multiples of $d$. Conversely if either coordinate is smaller than that of $c$,
increase the other coordinate until the residue pair has $j=d-1$.
The resulting point is outside the union. The conductor is therefore
exactly the rectangle $c+\mathbb Z_{\geq0}^2$, namely $(u^{d-1})$.

On the first chart the semigroup of
$k[u,uK,K^d]\subset k[u,K]$ is
$\{(x,y):x\geq y\bmod d\}$, giving conductor $(u^{d-1})$.
On the last chart that of $k[v,v^aZ_0,Z_0^d]$ satisfies
$x\geq a(y\bmod d)$, giving $(v^{a(d-1)})$.
The adjacent substitutions of \eqref{eq:chain-transitions-v169},
with $b,c$ replaced by $u,v$, make these generators differ by units
on overlaps. This verifies the claimed global ideal directly.
\end{proof}

\begin{example}[Third contact with unequal ramification]
\label{ex:third-ramified-v170}
For $a=3$, $\rho=2$, $\sigma=3$, the exceptional rays are
$(3,2),(3,1),(9,2)$. The chart indices, from the $v$-axis to the
$u$-axis, are $3,3,3,2$. Thus the normalized surface has four cyclic
quotient singularities. The boundary generic lengths are $2,1,2$,
so its nilradical is nonzero and square-zero, with no embedded
associated point. The residue degrees are $1,2,1$, the ramification
indices are $6,3,6$, and the conductor orders are $2,3,8$.
The corresponding transverse defects are $1,3,4$.
These numbers distinguish four different structures: the cyclic
quotient singularities of the normalization, the nonreduced
parameter fibre, the failure of normality of $Y$, and the branch
identifications made by its Hilbert coordinates.
\end{example}

\subsection{Compatibility as marked logarithmic geometry}
Let $\partial S'$ denote the two coordinate axes and let
$\partial Z$ be their strict transforms together with the full
reduced exceptional chain. These are the toroidal boundaries; their
logarithmic structures are the sheaves of functions invertible off
the boundaries, with their usual maps to the structure sheaves.

\begin{proposition}[Principalization and compatible root covers]
\label{prop:log-compatibility-v170}
The fan of $Z$ is the coarsest subdivision of the quadrant on which
each function
\[
 (p,q)\longmapsto\min(\rho p,j\sigma q),\qquad1\leq j\leq a,
\]
is linear on every cone. Equivalently, $Z$ is the terminal normal
toric modification making all $(u^\rho,v^{j\sigma})$ invertible.
The same terminal property holds among normal integral $S'$-schemes
whose inverse image of the coefficient torus is dense and on which
all these ideals become invertible.

The pair is log crepant:
\[
 K_Z+\partial Z=\pi^*(K_{S'}+\partial S').
\]
In particular it is log canonical. Replacing $u,v$ by
$U^r,V^s$ and normalizing the pullback gives, canonically,
$Z_{a,\rho r,\sigma s}$. These identifications are associative under
successive root covers. The construction also glues under changes
$u\mapsto\epsilon u$, $v\mapsto\eta v$ by units, or on a smooth
surface with two specified transverse Cartier divisors.
\end{proposition}
\begin{proof}
The bend of the $j$th function is the ray $n_j$. Every such ray is
necessary, and their ordered subdivision makes every function
linear. On a toric chart, linearity is precisely the condition that
one monomial generator divides the other. This proves the fan and
toric universal statements.

For a general test, the product ideal is invertible if and only if
all its factors are invertible: locally a factor of an invertible
product of nonzero fractional ideals has a fractional inverse.
The blow-up property therefore gives a unique map to $Y$. Pull
back the finite normalization $Z\to Y$ and close the unique generic
section over the dense coefficient torus. This is a finite birational
integral scheme over the normal test and is consequently isomorphic
to it. This proves the lift to $Z$ and its uniqueness.

The logarithmic canonical equality follows from the torus form
used in Proposition~\ref{prop:ramified-singularities-v170}. A toric
subdivision resolving the cyclic cones pulls back the same equality
to a smooth surface with a simple normal crossing boundary. That
pair is log canonical, which proves the assertion downstairs. This
is a statement about a parameter-space pair, not stable or semistable
reduction of the universal curves.

Normalization of the successive pullbacks is integral closure in
the same field $k(U,V)$. Transitivity of integral dependence gives
the canonical identification, and uniqueness gives associativity.
In lattice terms it is composition of the diagonal maps; no
arbitrary refinement is inserted. Unit changes preserve each ideal
$(u^\rho,v^{j\sigma})$. On a surface with ordered transverse
divisors the ideals are intrinsically
$\OO(-\rho U)+\OO(-j\sigma V)$, so their normalized blow-ups glue.
The conclusions are local in smooth or etale coordinate charts.
They do not identify fans after arbitrary unmarked coordinate changes.
\end{proof}

\begin{corollary}[Fibres of toric surface modifications]
\label{cor:toric-fibres-v170}
Let $W\to\A^2$ be any normal toric modification obtained by a finite
nontrivial subdivision of the first quadrant. If its interior
primitive rays are $(p_j,q_j)$, the fibre over the origin is a pure
one-dimensional Cohen--Macaulay scheme. Its reduction is a nodal
chain of projective lines, with generic lengths $\min(p_j,q_j)$
and nilradical of exact order $\max_j\min(p_j,q_j)$. In particular,
it is reduced exactly when every interior ray has a coordinate equal
to one. This assertion does not use the failure algebra or the
reconstruction theorem.
\end{corollary}
\begin{proof}
Apply Lemma~\ref{lem:coordinate-ideal-v170} on every cone of the
subdivision and then the argument proving
\eqref{eq:complete-fibre-ideal-v170}. Every maximal cone is proper
because the subdivision is nontrivial. The invariant divisors over
the origin are exactly the interior rays, and their ordered orbit
closures form the claimed chain.
\end{proof}

\subsection{The embedded curves and the information lost by normalization}
The flat universal curve on $Y$ is the pullback of the family on
$T_a$; its further pullback to $Z$ remains flat. This normalizes the
\emph{parameter space}, not the embedded curve. Consequently every
wall and chamber ring of Theorem~\ref{thm:arc-table-v169}, including
the genus-correction nilradical of
Theorem~\ref{thm:genus-correction-v169}, is retained unchanged at its
specified residue.

\begin{corollary}[Hilbert equality and normalization labels]
\label{cor:hilbert-lifts-v170}
For a coefficient arc with positive orders $p=\operatorname{ord}u$,
$q=\operatorname{ord}v$, the normalized wall rays are
$\rho p=i\sigma q$. On such a wall the additional normalization
coordinate is
\[
 \overline{u^{\rho/g_i}/v^{i\sigma/g_i}}\in k^*.
\]
The embedded curve remembers only its $g_i$th power. Thus if
$g_i>1$, equal retained coefficients and equal embedded Hilbert
limits need not give equal normalization points. Every such root
label is realized by an actual generically base-point-free arc.
In the interiors of the adjacent cones all arcs meet the unique
corresponding torus-fixed point.
\end{corollary}
\begin{proof}
The valuation vector lies in the unique cone of the normal fan
containing it in its relative interior. On a ray the primitive
perpendicular character is \eqref{eq:root-residue-v170}; its power
is the wall parameter $b/c^i$ of the old Hilbert graph. The map to
that graph therefore identifies exactly the $g_i$ roots on the
open orbit. Choose $u=\beta\tau^{A_i}$,
$v=\gamma\tau^{B_i}$; since
$\gcd(\rho/g_i,i\sigma/g_i)=1$, the expression
$\beta^{\rho/g_i}/\gamma^{i\sigma/g_i}$ takes every nonzero value.
These arcs have $u\ne0$ generically. The assertion for two-dimensional
cones follows directly from the affine character charts. This is an
explicit distinction between a Hilbert point and a normalization
lift, not a classification of target-automorphism orbits.
\end{proof}

\section{Coordinate descent and comparisons of compactifications}
\label{sec:equivariant-v169}
We first give an actual gluing statement at a fixed degree. It removes
the dependence on a monic pivot chart without identifying constructions
with different generic Hilbert polynomials.

\begin{theorem}[Equivariant gluing of the fixed-degree graphs]
\label{thm:global-gluing-v169}
Put
\[
 \mathcal P_a=\Pj(H^0(\Pj^1,\OO(a))\otimes k^3),\qquad
 H_a=\Hilb^{P_a}(\Pj^1\times\Pj^2),
\]
and let $\mathcal U_a\subset\mathcal P_a$ be the base-point-free
locus. The closure
\[
 \mathcal G_a=\overline{\{(F,[\operatorname{graph}(F)]):
                                F\in\mathcal U_a\}}
                    \subset\mathcal P_a\times H_a
\]
is an integral projective compactification of the graph of degree-$a$
morphisms with their coefficients retained. It carries a flat universal
embedded curve and a compatible $\PGL_2\times\PGL_3$ action. On every
monic pivot patch it is isomorphic, after the specified undoing of a
target change, to $\A^2\times\Gamma_a$. These patch isomorphisms
satisfy the cocycle identity on all overlaps, including boundary points.
The normalization also carries the action. The morphism
\[
 [\mathcal G_a/(\PGL_2\times\PGL_3)]
       \longrightarrow[\mathcal P_a/(\PGL_2\times\PGL_3)]
\]
is representable and proper, with projective pullback to a scheme
trivializing the group torsor.
\end{theorem}
\begin{proof}
The Hilbert scheme is projective and separated, so the closure is
projective and integral and is an isomorphism over $\mathcal U_a$.
Pullback of the universal Hilbert family is flat. A pair of source and
target automorphisms takes a graph to the graph of the transformed
map, preserves $\OO(1,1)$, and hence acts on its Hilbert scheme.
It preserves the closure and its universal curve.

Here are the patch maps explicitly. Choose a source point where some
coordinate form is nonzero, send it to $[0:1]$, and permute that target
coordinate into position zero. In the projective coefficient patch
normalize its $t^a$ coefficient to one and write
\[
 F_0=f^{\mathrm h},\qquad
 F_1=\alpha_1 f^{\mathrm h}+g^{\mathrm h},\qquad
 F_2=\alpha_2 f^{\mathrm h}+r^{\mathrm h},
 \qquad\deg g,\deg r<a.
\]
This identifies the coefficient patch with $\A^2\times B_a$.
The target matrix and its inverse are
\begin{equation}\label{eq:undo-global-v169}
 U(\alpha)=
 \begin{pmatrix}1&0&0\\-\alpha_1&1&0\\-\alpha_2&0&1\end{pmatrix},
 \qquad
 U(\alpha)^{-1}=
 \begin{pmatrix}1&0&0\\\alpha_1&1&0\\\alpha_2&0&1\end{pmatrix}.
\end{equation}
The first matrix identifies the graph with the polynomial-contact
graph; the second returns it to the original fixed target. Since the
product with $\A^2$ is flat, taking the closure of the generic graph
commutes with this product. Thus the closure on the patch is exactly
the asserted product, with its scheme structure.

The patches cover $\mathcal P_a$: a nonzero form cannot vanish at
$a+1$ distinct chosen source points, so finitely many such choices
and the three target pivots suffice. On an overlap the change is the
composition of the actual source matrices, the target permutations,
and \eqref{eq:undo-global-v169}. Its inverse reverses that composition.
The product of the matrices on a triple overlap is the identity as a
projective transformation. The induced Hilbert maps therefore satisfy
the same identity. No comparison based only on reduced closed points
is used.

For normalization, the product of the smooth connected group with
the normal variety $\mathcal G_a^\nu$ is normal. The action composed
with normalization is dominant and lifts uniquely to the normalization;
uniqueness proves the action identities. Finally an equivariant
morphism induces a representable map of the quotient stacks. After a
torsor is trivialized its pullback is the projective morphism
$\mathcal G_a\to\mathcal P_a$; properness descends.
\end{proof}
This theorem glues \emph{all} pivot charts of a fixed-degree rational-map
problem. Its construction is classical graph compactification; the
explicit patch maps state exactly how the contact graphs sit in it.
It is not a gluing of every Smith and Kronecker stratum of every
quadratic-pencil Hilbert problem. In particular, an adjugate tuple that
vanishes identically supplies no point of $\mathcal P_a$.

\begin{corollary}[Moving contacts and mixed target coefficients]
\label{cor:moving-v169}
Over $\A^1_\alpha\times S_a$ the family
\[
 ((x-\alpha)^a,b,c(x-\alpha)^{a-1})
\]
has strict Hilbert graph $\A^1_\alpha\times T_a$, with the universal
curve obtained by substituting $x-\alpha$ for $x$. Over the further
open parameter space $1-uv\ne0$, the same statement holds for
\begin{equation}\label{eq:mixed-family-v169}
 ((x-\alpha)^a,\ b+uc(x-\alpha)^{a-1},\ vb+c(x-\alpha)^{a-1}),
\end{equation}
after undoing the target matrix
$\left(\begin{smallmatrix}1&0&0\\0&1&u\\0&v&1\end{smallmatrix}\right)$.
For $a=2$ both latter coordinate polynomials can have nonzero constant
and linear terms. This family is an equivariant enlargement of the
marked slice, not the unrestricted six-coefficient contact family.
\end{corollary}
\begin{proof}
Translation is a source automorphism over $\A^1_\alpha$. The indicated
target matrix is invertible precisely on the stated open, with inverse
lower block $(1-uv)^{-1}\left(\begin{smallmatrix}1&-u\\-v&1\end{smallmatrix}\right)$.
They identify the entire ambient families, not just their generic
fibres. Applying these isomorphisms to the flat universal curve and
its graph closure proves the result.
\end{proof}

\begin{proposition}[Conservation of the Hilbert polynomial]
\label{prop:conservation-v169}
A projective flat finitely presented family over a connected base
cannot have, as its fibres, just the graphs of morphisms of two
different degrees. A degree drop from $a$ to $a-d$ in a limit of
$\mathcal G_a$ must be accompanied by additional scheme structure
whose Hilbert polynomial contributes the missing degree and Euler
characteristic. The extra structure may include punctual nilpotents,
as in Theorem~\ref{thm:genus-correction-v169}.
\end{proposition}
\begin{proof}
The Hilbert polynomial for $\OO(1,1)$ is locally constant in a
projective flat family, whereas the two graphs have polynomials
$(a+1)l+1$ and $(a-d+1)l+1$. These differ for $d>0$. In a fixed-degree
Hilbert compactification the polynomial is retained by the whole
limiting scheme, not by deleting its tails or embedded associated
points.
\end{proof}
Thus comparison across different generic polynomials cannot consist
of an equality of the unaugmented flat graph functors. The horizontal
comparison and common-refinement results in
Section~\ref{sec:comparison-v167} continue to apply on a fixed dense
open with a fixed generic embedded quotient. They are not changed into
a false full-pullback statement by the equivariant gluing above.

\subsection{Tropical data, residues, and coordinate changes}
\begin{proposition}[Comparison of coefficient presentations]
\label{prop:fan-comparison-v169}
For a finite list of polynomial projective coordinates, the finite
weight-layer and residue construction of
Theorem~\ref{thm:finite-fan-v167} extends to exact monomial profiles
with nonnegative weights, including specified unit coordinates, and
to polynomial coordinates translated about any fixed centre. It gives
finitely many algebraic families of closed projective points, not
finitely many embedded curves or abstract isomorphism types.

Under an invertible monomial change of coefficient torus coordinates,
the weights transform by the exponent matrix and the residues by the
same monomial substitution; the finite partitions have a common
pullback refinement. Under a general polynomial coordinate change,
the corresponding statement holds after composing the projective
coordinate polynomials with that change. An exact monomial arc need
not remain monomial in the new coordinates, and the two displayed
fans need not be isomorphic.
\end{proposition}
\begin{proof}
Expand each of the finitely many polynomials into its finite monomial
support, after any specified translation or polynomial substitution.
The equalities of weights partition the nonnegative orthant into a
finite rational polyhedral arrangement. Within a cone, group terms
of equal weight and partition the residue torus by the first nonzero
coefficient sum in each polynomial. Zero polynomials are kept as zero;
identically zero input coordinates are a separate designated pattern.
Every choice is finite, and the residual vector is polynomial on each
locally closed piece. Weight zero introduces constant residue terms
but does not change the argument. For a monomial substitution the
weight and residue formulas follow directly from substitution.
For a polynomial substitution repeat the finite expansion on the
original monomial arc; this compares the resulting points even when
there is no transformation of fans by a linear weight map. For
localized-polynomial changes work on an open where denominators are
units and clear them by a common denominator, which does not change
the projective map.
\end{proof}

Applied to maximal minors, valuations are realizable valuated-matroid
data: the minors satisfy all Pl\"ucker equations, and the minimum
valuation in each nonzero relation is attained at least twice. On the coordinate stratum of the nonzero minors these are points
of the tropicalization of that Grassmannian stratum, as in
Speyer--Sturmfels \cite[Theorem~2.1 and Section~3]{SpeyerSturmfels169}.
It is not enough to check only a selected set of tropical quadratic
relations and then claim realizability. Valuations forget residues;
our first wall already has the same valuations for every $\kappa\ne0$
but different embedded ideals. The primitive Hilbert vector retains
these residues and determines the degree-$m$ kernel. It does not
select among normalization branches over a graph point.

On $S_a$ the individual nonzero minors are monomials. The normal fan
of their Newton polygon is exactly the chain fan, after removal of
content; all its walls are effective. This identifies the computed fan
with the restriction to nonnegative weights of the state fan of this marked torus orbit, in the classical
initial-ideal setting of \cite{BayerMorrison167}. No such minimality
claim is made for the full coefficient-support arrangement. In
comprehensive Gr\"obner terminology, which organizes specializations
by leading-coefficient conditions \cite{ManubensMontes169},
Proposition~\ref{prop:curve-charts-v169} gives a particularly simple
system: a single monic basis on every affine parameter chart, with no
further coefficient conditions. Its explicit standard monomials,
rather than a Hilbert-scheme elimination, recover the curve ideals.

\begin{proposition}[Equality up to a target automorphism]
\label{prop:target-equivalence-v169}
Let two embedded curves with polynomial $P_a$ have degree-$m$
ideal kernels $K,K'\subset H^0(X,\OO(m,m))$, $m\geq m_a$.
They differ by a target automorphism exactly when there is
$A\in\PGL_3$ for which
\[
 (1\otimes\operatorname{Sym}^m A^*)K'=K.
\]
This is an invariant, finite algebraic criterion: equality of the
two fixed-rank subspaces is given by rank conditions, with $\det A\ne0$.
It does not classify abstract curve isomorphisms. On a fixed wall
of the marked family, all nonzero residue values are related by a
diagonal target automorphism, although their fixed-target Hilbert
points differ.
\end{proposition}
\begin{proof}
An automorphism acts on sections by pullback. The Hilbert immersion
in degree $m$ makes equality of the kernels equivalent to equality
of ideal sheaves. This proves both implications and the algebraic
criterion. Scaling $F,G,R$ independently scales $\kappa$, or scales
$\lambda$ by the ratio between $R^j$ and $F^{j-1}G$; every nonzero
ratio is available over $k$. No assertion about different walls is
needed for this last observation.
\end{proof}

\subsection{A concrete comparison with Quot and stable-map limits}
\begin{proposition}[A Quot contraction of the entire exceptional chain]
\label{prop:quot-contraction-v169}
The marked family defines a flat quotient on the fixed source
\begin{equation}\label{eq:quot-sequence-v169}
 0\longrightarrow\OO_{\Pj^1}(-a)
   \xrightarrow{(t^a,bs^a,cst^{a-1})}\OO_{\Pj^1}^{\oplus3}
   \longrightarrow\mathcal Q\longrightarrow0
\end{equation}
over all of $S_a$. The resulting morphism from $T_a$ to the fixed-source
Quot scheme is constant on its entire exceptional chain. The common
central quotient is
\[
 \mathcal Q_0\simeq\OO_{a[1:0]}\oplus\OO_{\Pj^1}^{\oplus2}.
\]
After adding three fixed source markings away from $[1:0]$, this also
gives stable quotients in the sense of Marian--Oprea--Pandharipande.
In contrast, the Hilbert map distinguishes every point of the chain.
\end{proposition}
\begin{proof}
The first coordinate $t^a$ makes the sheaf map fibrewise injective,
even when $b=c=0$. Both source and target bundles are flat over $S_a$;
the fibrewise injectivity criterion for this finite-presentation
sequence gives flatness of its cokernel over $S_a$. Thus it defines
a morphism to the Quot scheme already before the blow-up. At the
origin the sequence is multiplication by $t^a$ in its first summand,
which gives the stated quotient. Its pullback along $T_a\to S_a$
therefore contracts the chain. The Hilbert map on that chain is a
closed immersion because $T_a$ is a closed graph in $S_a\times H_a$.
For three markings the source is stable, the torsion is away from
the markings, and $\omega_{\Pj^1}(p_1+p_2+p_3)\otimes\OO(a\epsilon)$
has positive degree for every $\epsilon>0$. These verify the actual
stability conditions of \cite[Section~2]{MOP169}; no unstable
unmarked rational source is being called a stable quotient.
\end{proof}

\begin{corollary}[Non-Cohen--Macaulay limits of smooth rational graphs]
\label{cor:external-curves-v169}
For each $a\geq3$ the Hilbert closure of smooth rational graphs in
$\Pj^1\times\Pj^2$ with polynomial $P_a$ contains a curve whose
reduction is a line attached to a rational cuspidal plane curve of
degree $a$, and whose punctual nilradical has length $g_a$ and
nilpotence order $a-1$. This curve is invisible as a distinct punctual
structure to the associated-cycle map. It is not the scheme-theoretic
image of a map from a reduced nodal curve.
\end{corollary}
\begin{proof}
Take $b=\tau^a,c=\tau$. The generic graph is smooth because its
projection to the source is an isomorphism. Its limit is
$\mathcal W_a$ with $\lambda=1$, so
Theorem~\ref{thm:genus-correction-v169} proves the first assertions.
An associated one-cycle records generic lengths at one-dimensional
components and discards finite-length submodules. Finally the
scheme-theoretic image of a morphism from a reduced scheme is reduced:
if a local target function has a power mapping to zero, its pullback
is nilpotent and hence zero. The asserted nonreduced curve cannot
be such an image. This does not rule out comparison spaces between
stable-map and Hilbert compactifications; it rules out identifying
this Hilbert limit with the unmodified scheme image of a stable map.
\end{proof}
This application is independent of the finite failure algebra and of
Paper I. General comparisons between stable maps and stable quotients
already involve spaces dominating both constructions
\cite{Manolache169}; the explicit contraction above is a different,
fixed-source Hilbert-to-Quot calculation, not a replacement for those
virtual-class comparison theorems.

\subsection{Return to the reconstructed pencil}
A specified effective failure family first reconstructs its pencil by
Paper I. Wherever its fixed-target contact comparison lands in $B_a$,
the graph in Theorem~\ref{thm:global-gluing-v169} supplies compatible
models under source and target frame changes. If its coefficient map
factors through \eqref{eq:marked-family-v169} or
\eqref{eq:mixed-family-v169}, the chain, all local rings, and the
punctual modules above pull back as strict horizontal families.
The inherited complete-quadric comparison determines the target in
that application; forgetting that target and its coefficient map
would lose the meaning of the embedded limit. This statement does
not identify every complete-quadric Hilbert component with
$\mathcal G_a$, does not cover an identically zero adjugate tuple,
and does not turn Kronecker data alone into a higher-corank fibre
classification. No boundary operation on an isolated unframed
multiplication table is inferred.

\appendix
\section{Extension classes and normalized ramification charts}
\label{sec:local-algebra-v167}
We record two local consequences that separate properties of the
embedded curve from properties of its parameter space.

\subsection{The presentation complex and explicit splitting}
For a polynomial presentation $R=P/J$ and an $R$-module $N$,
square-zero $\C$-algebra extensions with specified kernel $N$
and specified quotient $R$ are classified by
$\operatorname{Ext}^1_R(L_{R/\C},N)$. We use cohomological degrees:
$J/J^2$ is in degree $-1$ and
$\Omega_{P/\C}\otimes_PR$ in degree $0$. The group is
\[
 \coker\bigl(\Hom_R(\Omega_{P/\C}\otimes_PR,N)
                       \longrightarrow\Hom_R(J/J^2,N)\bigr).
\]
Indeed a lift of polynomial generators gives a conormal map;
changing lifts changes it by a derivation. The part of the
cotangent complex in degrees at most $-2$ contributes only in
Ext degrees at least two when the target is in degree zero.
This justifies truncation for this particular group, not a claim
that the whole cotangent complex is two-term.

In the chart of Theorem~\ref{thm:extension-class-v166}, invert $C$.
Then the extension algebra is explicitly
\[
 \C[\lambda,A,B,C,C^{-1},n]/(n^2,An,Bn),\qquad e=n/C.
\]
Writing its elements as $r+an$ gives multiplication
$(r+an)(s+bn)=rs+(rb+sa)n$, where the coefficients of $n$ are
restricted modulo $(A,B)$. This is the actual split algebra on
this chart. On an overlap of two obstruction charts, let $\xi_i$
be the conormal representative obtained from the same algebra
extension. A change of polynomial lifts changes $\xi_i$ by a
derivation. Therefore $[\xi_i]=[\xi_j]$ under the induced overlap
map on the Ext sheaves. Taking the trivialization that sends
$1$ to $[\xi_i]$ makes the overlap function exactly $1$, rather
than an unspecified unit. This proves the asserted trivialization
of the local obstruction sheaf. It does not trivialize the
nilradical line bundle on the singular-conic plane, and it does
not identify the entire global hyper-Ext group with its sheaf of
local degree-one groups. The singular-conic plane is the locus
where the linear part of the conic equation at the attaching
point vanishes; the doubled-line curve is its discriminant-zero
locus. Thus the support of the embedded associated subvariety
and the obstruction to an algebra retraction are different loci.

\begin{proposition}[Singularities and divisor classes after ramification]
\label{prop:toric-classes-v167}
Let $m\geq1$ and put $g=\gcd(m,2)$. For the normalized transverse
surface in Theorem~\ref{thm:ramified-v166}, the character lattice
and its saturated semigroup are
\[
 M_m=\{(a,b)\in\mathbb Z^2:a-2b\equiv0\pmod m\},\qquad
                         M_m\cap\mathbb R_{\geq0}^2.
\]
The primitive rays of the dual cone in $N_m=\Hom(M_m,\mathbb Z)$
are $(1/g,0)$ and $(0,1)$. Its divisor class group is
$\mathbb Z/(m/g)$. For odd $m$ it is the cyclic quotient of type
$\frac1m(1,-2)$; it is smooth for $m=1$ and has a unique singular
point for every odd $m>1$. For even $m=2q$ it is of type
$\frac1q(1,-1)$, with equation $eh=\tau^q$; it is smooth for
$q=1$ and has a unique singular point otherwise. In the full
chart the additional coordinate $\lambda$ gives the product with
an affine line. The singular locus is correspondingly empty or
that line.
\end{proposition}
\begin{proof}
The exponent lattice and its saturation were computed in
Theorem~\ref{thm:ramified-v166}. Their dual lattice is
$N_m=\mathbb Z^2+\mathbb Z(1/m,-2/m)$. On the first coordinate
axis its primitive nonzero vector is $(1/g,0)$; on the second it
is $(0,1)$. This follows by solving the congruence in the unused
coordinate. The index of the lattice generated by these two rays
in $N_m$ is $m/g$. The torus-invariant divisor sequence is the
cokernel of
\[
 M_m\longrightarrow\mathbb Z^2,\qquad (a,b)\longmapsto(a/g,b).
\]
For odd $m$ its image is $a-2b\equiv0\pmod m$. For even
$m=2q$ write $a=2a'$; its image is
$a'-b\equiv0\pmod q$. In both cases the cokernel is the stated
cyclic group. This computes the class group of the normal affine
toric surface, including the smooth cases.

For odd $m$ neither weight has a nontrivial stabilizer on a
coordinate axis away from the origin. The quotient is therefore
smooth away from the origin. Its cone is unimodular exactly when
$m=1$, proving singularity at the origin otherwise. For even
$m$ taking invariants under the order-two pseudoreflection first
gives coordinates $s^2,r$ with residual opposite weights of order
$q$. The displayed hypersurface and its derivatives prove the
remaining assertions. Formation of the product with the
independent coordinate does not change this transverse analysis.
\end{proof}

\begin{lemma}[Gluing integral closures]\label{lem:normalization-gluing-v167}
Let an integral Noetherian scheme be covered by affine charts
$\Spec A_i$ in a fixed function field. Their integral closures
in that field glue to its normalization. In particular the
ramified invariant charts and the unchanged primitive charts
in Theorem~\ref{thm:ramified-v166} glue without a choice of a
cyclic cover or an ordering of branches.
\end{lemma}
\begin{proof}
Integral closure commutes with localization. One implication is
immediate by localizing a monic equation. For the other, clear
the finitely many denominators in a monic equation for an element
integral over a localization: a suitable multiple of that element
by an element of the multiplicative set is integral over the
original ring. Thus both chart normalizations restrict to the
same ring on every principal overlap. Affine refinements and the
cocycle identity of localization give the gluing. All rings here
are finite-type over $\C$, so these normalizations are finite.
\end{proof}


\section{Descent and the scheme structure of the contact charts}
\label{sec:comparison-details-v166}
This section records the descent, exactness, and overlap arguments
used by the contact theorems. In particular, formal coordinates are
never used as a substitute for an algebraic family in the fixed
target. The coefficient projection is retained in every graph point.

\subsection{The category of embedded deformations}
Fix a framed coefficient atlas $S$, a central point $s$, and a
central embedded curve. The test category has objects
$(A,\beta)$, where $A$ is a local Artin $\C$-algebra with residue
field $\C$ and $\beta:\widehat\OO_{S,s}\to A$ is a local map.
Its morphisms are local algebra maps commuting with $\beta$.
The Hilbert deformation functor is covariant on this algebra
category, or contravariant on the opposite category of test schemes.
It is \emph{set-valued}: an object is an embedded ideal quotient,
and two objects are equal only when their ideals in the fixed
target coincide. No automorphism group of the target is divided out.

Let $X=\mathrm{CQ}(V)$, let $\mathcal L$ be the universal pencil
line, and let $\mathcal Y=\mathcal L\times_{\Pj(\Sym^2V)}X$.
For each completed contact neighbourhood the comparison diagram is
\begin{equation}\label{eq:target-undo-v166}
\begin{array}{ccccc}
 \widehat{\mathcal Y}_{i,A}&\xrightarrow{\ \Theta_i\ }&
       \widehat{\mathcal N}_{i,A}
       &\xrightarrow{\ \Theta_i^{-1}\ }&\widehat{\mathcal Y}_{i,A}\\
 \downarrow&&\downarrow&&\downarrow\\
 \widehat{\mathcal L}_{i,A}&=&\widehat{\mathcal L}_{i,A}
       &=&\widehat{\mathcal L}_{i,A}.
\end{array}
\end{equation}
The last upper object maps to $X\times\operatorname{Spec}A$ by
the original incidence embedding. The middle upper object is the
normal-coordinate graph target, not a new fixed target. The
operation producing a Hilbert quotient always includes the last
arrow. On a common domain of two choices define
$\Phi_{ji}=\Theta_j\Theta_i^{-1}$. Then
\begin{equation}\label{eq:frame-cocycle-v166}
 \Phi_{kj}\Phi_{ji}=\Phi_{ki},\qquad
 \Theta_j^{-1}\Phi_{ji}\Theta_i=1.
\end{equation}
These identities hold on the ambient targets and hence on their
ideal quotients and on every Artin base change. They give the
triple-overlap cocycle and show why a parameter-dependent congruence
does not identify distinct curves in the fixed target.

\begin{proposition}[Precise patching and algebraic comparison]
\label{prop:patching-algebraic-v166}
On an affine contact chart let $D$ be a Noetherian ring and let $J$
cut out the closed contact support. Use the ring-theoretic completion
$\widehat D$, the open $U=\operatorname{Spec}D\setminus V(J)$,
and the overlap
$V=\operatorname{Spec}\widehat D\setminus V(J\widehat D)$.
Finitely presented ideal quotients on $U$ and
$\operatorname{Spec}\widehat D$, with compatible quotient algebra
structures on $V$, patch uniquely to an ideal quotient on
$\operatorname{Spec}D$. Quotients with contact-supported torsion
are allowed. Base flatness, equality, and the closed incidence
equations are detected by this patching cover.

Consequently, suppose an algebraic family of finite presentation
has first supplied a morphism $h:T\to H$ into a Hilbert scheme or
its fixed coefficient graph. If recovery makes $h$ a monomorphism
and the embedded comparison identifies completed local rings at
all closed points, then $h$ is an open immersion. This conclusion
uses the algebraic family, not only an isomorphism of formal functors.
\end{proposition}
\begin{proof}
Completion of a Noetherian ring is flat and induces an isomorphism
over $D/J$. Bhatt's flat formal-gluing statement
\cite[Proposition 5.6(4), Example 5.8]{Bhatt163} therefore applies
to quasi-coherent sheaves. Apply its monoidal equivalence to the
quotient, its unit, multiplication, and its structure map from the
ambient ring. Surjectivity and finite presentation descend on the
jointly surjective flat cover given by the completion and the open
complement. The glued algebra is thus an ordinary finitely presented
quotient, not a derived replacement. The same cover detects base
flatness and the vanishing of the prescribed incidence equations.
The gluing theorem supplies existence; faithful flatness alone
would not supply the missing overlap descent datum. Iteration over
disjoint supports and over projective affine charts gives the
claimed form for the incidence target.

For the last assertion, use the completed local criterion for
etaleness at closed points, whose residue fields are $\C$
\cite[Lemma 41.11.3]{Stacks163}. The morphism is already of finite
presentation. Its etale locus is open; its closed complement in a
finite-type complex chart would contain a closed point if nonempty.
Thus it is etale everywhere under consideration. An etale
monomorphism is an open immersion. Completed comparison here means
comparison on all Artin quotients, including nilpotent test rings;
an agreement of tangent dimensions or closed reduced points is
insufficient.
\end{proof}

The reduced graph closure is compatible with completion by
Lemma~\ref{lem:saturation-v166}, using the finite stabilized colon
ideal, not a general assertion that closure commutes with every base
change. Its completed rings remain reduced for the following precise
reason. A reduced local complex finite-type ring is excellent, so its
formal fibres are geometrically reduced; see
\cite[Tag 0BJ0]{Stacks166}. The injection into the finite product of
its minimal-prime fraction fields remains injective after flat
completion. Each resulting generic formal fibre is reduced.
The completion, as a subring of their product, is therefore reduced.
Normalization is unchanged on the regular total charts; this says
nothing about normalizing their possibly nonreduced fibres.

\begin{lemma}[The larger Smith exponent is a smooth graph factor]
\label{lem:larger-smith-v166}
In the formal corank-two model
$\left(\begin{smallmatrix}f&g\\g&fk+r\end{smallmatrix}\right)$,
the coefficient parameters of $k$ are independent smooth factors for
the relative embedded graph comparison. They are not forgotten in
the retained coefficient point.
\end{lemma}
\begin{proof}
The generating triples $(f,g,fk+r)$ and $(f,g,r)$ are related over
the completed parameter ring by
\[
 T_k=\begin{pmatrix}1&0&0\\0&1&0\\-k&0&1\end{pmatrix},
 \qquad T_k^{-1}=T_{-k},\qquad\det T_k=1.
\]
Thus their projective graph targets are isomorphic over the
parameter disc. Apply $T_k$ on that target and $T_k^{-1}$ before
patching into the fixed incidence target. This gives inverse
operations on embedded ideals over every Artin coefficient algebra.
The independence of the coefficients of $k$ is the coefficient
normal form of Theorem~\ref{thm:ambient-versal-v162}. Therefore
the completed graph rings acquire these parameters as power-series
variables. This is an isomorphism of relative graph presentations
with a specified inverse, not a constant ambient automorphism of
$X$ or an unlabelled Hilbert-point inversion.
\end{proof}

\subsection{Exactness and recovery at every conic}
\begin{lemma}[Height, flatness, and the recovery minors]
\label{lem:HB-details-v166}
The three minors in \eqref{eq:HB-ideal-v163} have codimension two
on every geometric fibre of the six-parameter family. Their
Hilbert--Burch complex is fibrewise exact, and its cokernel is flat
over that parameter space. At a conic-tail point the ambient
restriction maps in degrees $(0,2)$ and $(1,1)$ have respectively
ranks five and four, from ambient spaces of dimension six. Thus
recovery uses kernels of ranks one and two, including at doubled
lines.
\end{lemma}
\begin{proof}
At $e=0$ the ideal is
$(zG,zH,Q)=(G,H)\cap(z,Q)$, where $Q\ne0$ and is independent of
$z$. There is no common codimension-one factor: any factor common
to $zG,zH$ divides $z$, which does not divide $Q$. The displayed
union has pure support of codimension two. At $e\ne0$, on $s=1$
eliminate $H$ from $zG-e(H+BF)=0$. The other linear minor becomes
\[
 (z^2+e^2C)G+(e^2A-eBz)F=0.
\]
It is nonzero, with coefficient of $z^2G$ equal to one. Together
with the first linear equation it is a regular sequence in the
affine polynomial ring; the remaining quadratic follows from the
syzygies. At $s=0$ the two bilinear equations force $G=H=0$.
There is no component at infinity and no intermediate parameter
value at which the height drops. The perfect-height-two criterion
therefore gives the Hilbert--Burch resolution on every fibre.
The fibrewise exactness criterion for a complex of base-flat locally
free modules gives exactness and flatness of the cokernel
\cite[Lemma 10.99.5]{Stacks166}.

On $C_0\cup Q$ use its union sequence and the plane-hypersurface
sequence $0\to\OO_{\Pj^2}(j-2)\to\OO_{\Pj^2}(j)
\to\OO_Q(j)\to0$. They give $h^0=5,4$ and $H^1=0$ in the
respective degrees, and surjectivity of restriction from the
ambient six-dimensional spaces. This argument uses only that the
quadratic equation is nonzero, so includes reducible conics and
doubled lines. The nonzero $5$ by $5$ and $4$ by $4$ restriction
minors define the constant-rank opens. Normalize the selected
$R^2$ coefficient of the kernel quadratic to one, and normalize
the two bilinear equations by their $tG,tH$ coefficient matrix.
Its $2$ by $2$ pivot is invertible at the conic-tail point. These
are the open minors used by the recovery maps. The coefficient
projection supplies $f$ and hence $c$; recovery of $e$ from the
normalized bilinear equation would not replace that projection.
\end{proof}

\begin{lemma}[Connectedness, scheme exhaustion, and unlabelled descent]
\label{lem:exhaustion-v166}
Let $p:Z\to B$ be proper and birational, with $Z$ integral and
normal and $B$ integral and normal of finite type over $\C$.
Its geometric fibres are connected. A nonempty proper union of
constructed closed families which is open in the support of a
fibre therefore exhausts that support. If algebraic charts cover
those families and compute the fibre ideals, the same charts
determine the complete scheme fibre.

For the disjoint-contact product theorem, the map from the product
of local families to the Hilbert fibre is a representable algebraic
morphism of finite presentation. Completed comparison on all Artin
rings and ideal recovery make it an open immersion. Its properness
and the preceding connectedness then give the asserted isomorphism.
The construction descends without global spectral labels.
\end{lemma}
\begin{proof}
Stein factorization gives connected geometric fibres followed by
a finite birational map to $B$. Normality of $B$ makes the latter
an isomorphism; equivalently use
\cite[Lemma 37.53.6]{Stacks163}. The stated union is closed by
properness and open by the local chart descriptions, hence is the
whole nonempty connected support. Connectedness alone does not
compute nilpotents: those are supplied by the fibre ideals on the
covering charts.

For the product assertion, the contact-supported quotients are
finite-presentation algebraic quotients on disjoint finite contact
neighbourhoods. Patching them to the fixed main curve produces a
flat embedded family with the required polynomial, so Hilbert
representability supplies the morphism before formal comparison
is applied. Restriction to the supported quotients recovers each
factor on every test scheme, giving a monomorphism. Proposition
\ref{prop:patching-algebraic-v166} now applies to its completed
local maps, including nonreduced Artin tests. A proper open immersion
has open-and-closed image; connectedness supplies surjectivity.
On the finite etale cover ordering the distinct supports, changes
of order permute the corresponding factors and the supported
quotients. These permutation maps satisfy the cocycle identity,
so descent gives the unlabelled object. This does not impose an
ordering through a collision of supports.
\end{proof}

Our convention for $\mathbb F_2=\Pj(\OO\oplus\OO(2))$ is the
projectivization of one-dimensional subspaces. The chart $[1:v]$
is $\operatorname{Tot}\OO(2)$; its section at infinity is the
subline $\OO(2)$ and has normal bundle
$\operatorname{Hom}(\OO(2),\OO)=\OO(-2)$. This fixes the negative
section and the sign used in the component-intersection formulas.

\subsection{The quotient overlap maps and divisor classes}
Write the ordered directions as $w_\pm=\lambda\pm h$, and the
ordered root parameter as $t$. On the blow-up of $(h,t)$ use
$t=eh$ or $h=ut$. The involution sends $h$ to $-h$ in the first
chart and $t$ to $-t$ in the second. Its fixed locus is the
exceptional divisor. Its invariant coordinates are respectively
$(\lambda,e,C=-h^2)$ and $(\lambda,u,\delta=t^2)$; their rings
are polynomial even at the fixed locus.

For completeness, a change of direction coordinate
$\phi(w)=(aw+b)/(cw+d)$, with $\Delta_\phi=ad-bc\ne0$, gives
on the first invariant chart
\begin{equation}\label{eq:quotient-overlap-first-v166}
\begin{split}
 K&=(c\lambda+d)^2+c^2C,\\
 \lambda'&=\bigl((a\lambda+b)(c\lambda+d)+acC\bigr)/K,\\
 C'&=\Delta_\phi^2C/K^2,\qquad e'=Ke/\Delta_\phi.
\end{split}
\end{equation}
On the second chart it gives
\begin{equation}\label{eq:quotient-overlap-second-v166}
\begin{split}
 K&=(c\lambda+d)^2-c^2\delta u^2,\\
 \lambda'&=\bigl((a\lambda+b)(c\lambda+d)-ac\delta u^2\bigr)/K,\\
 u'&=\Delta_\phi u/K,\qquad\delta'=\delta.
\end{split}
\end{equation}
These formulas are used where both directions are in both affine
charts, so $K$ is invertible. They follow by taking the mean and
half-difference of $\phi(w_+),\phi(w_-)$. The overlap between the
two blow-up charts is $e=1/u$, $C=-\delta u^2$. These identities
supply every overlap, since any two directions admit a common
point at infinity outside them. The triple-overlap identities
follow from composition of the fractional linear maps, before
taking invariants.

\begin{lemma}[Projectivity and ramification of the collision quotient]
\label{lem:quotient-projective-v166}
The finite involution quotient of the ordered blow-up is projective
over the unordered coefficient line. If $T_0$ is the strict
transform of $t=0$ and $F_0$ its exceptional divisor, and if $S,E$
are their respective images, then for the quotient map $q$
\[
 q^*S=T_0,\qquad q^*E=2F_0,\qquad
 \operatorname{div}(t)=T_0+F_0,\qquad
 \operatorname{div}(\delta)=2S+E.
\]
In particular $\OO_S(S)=\OO_{\Pj^2}(-1)$.
\end{lemma}
\begin{proof}
Choose a relatively ample line bundle on the blow-up and tensor
its two translates. It has the permutation linearization. Squaring
kills the characters of every stabilizer on its fibres. It thus
descends to the finite quotient; equivalently, products of the
translates of sufficiently high-degree sections give invariant
sections and their affine projective charts. The descended line
bundle is ample, since its finite pullback is ample. This proves
projectivity without assuming that the action is free. In these
charts its overlap maps are those of the tensor product
linearization, followed by their squares, so they satisfy descent
on triple overlaps.

The invariant charts show ramification index two at the generic
exceptional divisor and index one at the generic strict transform.
They prove the first two identities. The blow-up of $(h,t)$ gives
$\operatorname{div}(t)=T_0+F_0$, and $\delta=t^2$ gives the last
identity after pullback by $q$. On $S=\operatorname{Sym}^2\Pj^1$
the divisor $E|_S$ is the doubled-line conic, of class $2H$.
Since $E+2S$ is principal, $2[S|_S]=-2H$ in
$\operatorname{Pic}(\Pj^2)=\mathbb Z H$. This group has no
two-torsion, so $[S|_S]=-H$ as claimed.
\end{proof}

In the full conic pullback the finite colon calculation is
$((eA,eB):e^2C)=(A,B)$ and further colons have the same value.
Thus for the original unramified coefficient line all torsion is
already annihilated by $\delta$. The horizontal ring is integral
and dominates $\C[\delta]$, hence is torsion-free over that base.
The discrete valuation rings used to infer flatness are the
nonfield localizations of \emph{$\C[\delta]$}, not localizations
of the threefold coordinate ring. The higher Tor groups vanish
by \eqref{eq:all-tor-v166}; for ramified arcs their first Tor is
only the last layer of the larger torsion ideal.

Finally choose $\delta_0>0$ and the positive root at that base point.
On the punctured analytic disc the local system
$R^2p_*\Q$ for the horizontal parameter space has basis the two
ruling classes of $\Pj^1\times\Pj^1$. The loop
$\delta(\theta)=\delta_0\exp(2\pi\sqrt{-1}\theta)$ lifts to an
exchange of the two roots, and acts by exchanging that basis.
Its invariant line is the sum and its anti-invariant line the
difference. In the normalization after $\delta=t^2$, the element
$h=t/e$ is taken in the fraction field of the integral base-changed
horizontal ring; adjoining it is finite and birational. This is
normalization of the base-changed total space. It is neither
normalization of the original total graph nor normalization of
its central fibre. The resulting normal crossings belong to the
parameter space, not to a stable-map model of every embedded curve.


% BEGIN PRESERVED PART 00-principal-introduction-v152.tex
\section{Detailed reconstruction statements and comparisons}
\label{sec:introduction-v132}

A degeneracy locus usually forgets the map that defines it.  We study
an inverse problem in which neither the ambient embedding nor the
normal directions are retained: can the abstract nonreduced failure
scheme recover a quadratic pencil?  The answer is affirmative for
every complex pencil, including singular pencils.  The first relation, rather than its reduced determinant support,
carries this information.  The inverse is also global: a thickening over a projective reduction
recovers an entire varying pencil subbundle and its actual source
bundle, up to one constant left transformation.  No fibrewise choice
of coordinates or projective line-bundle ambiguity is left.  A
single unmarked local algebra suffices for the constant-pencil
problem; a Schubert line gives another geometric realization.

Let \(V\) be a complex vector space of dimension \(n\ge3\), and put
\[
 W=\Sym^2V,\qquad N=\binom{n+1}{2},\qquad
 R\in\Gr(2,W),\qquad S_R=W/R,\qquad p=N-2.
\]
The cube-zero algebra \(A_R=\C\oplus V\oplus S_R\) has multiplication
\(\Sym^2V\twoheadrightarrow S_R\), with
\((V\oplus S_R)S_R=0\).  On \(X_R=\Gr(n,V\oplus S_R)\), let
\(\mathcal K\) denote the tautological bundle.  For any \(m\ge2\), set
\begin{equation}\label{eq:failure-intro}
 \widehat D_R=V\!\left(\Fitt_0\coker[
 \Sym^m(\OO\oplus\mathcal K)\longrightarrow A_R\otimes\OO]\right).
\end{equation}
This is the maximal-minor scheme of multiplication, not merely its
support.  It is independent of \(m\ge2\).  Its reduction is the
Schubert divisor where \(\mathcal K\to V\) drops rank.  Write
\(B_R=\Gr(n,S_R)\), let \(\mathfrak b_R\) be its ideal in
\(\widehat D_R\), and define
\[
 Z_R^{[k]}=V(\mathfrak b_R^{k+1})\subset\widehat D_R,
       \qquad d=n+2p=n^2+2n-4.
\]
An isomorphism of these neighbourhoods means an abstract complex
scheme isomorphism; it is supplied with no ambient or tensor marking.

\begin{theorem}[Sharp intrinsic reconstruction]
\label{thm:principal-finite-v143}
For all quadratic pencils \(R,R'\subset\Sym^2V\),
\[
 Z_R^{[d]}\simeq Z_{R'}^{[d]}
       \quad\Longleftrightarrow\quad
 R'=gR\quad\text{for some }g\in\PGL(V).
\]
For every \(0\le k<d\), the isomorphism class of \(Z_R^{[k]}\)
is independent of \(R\).  Thus \(d\) is the smallest uniform
reconstruction order.  The same statements hold after restriction,
using the graph projection, to any Schubert line in \(B_R\).
The resulting abstract curve-supported neighbourhood is independent
of the choice of that line.
\end{theorem}

The Grassmannian statement is Theorem~\ref{thm:sharp-finite-pencil}.
The curve-supported statement is
Corollary~\ref{cor:schubert-line-reconstruction-v145}.  Its more general
form, Theorem~\ref{thm:curve-reconstruction-v145}, applies over any
smooth connected projective curve with a rank-\(n\) bundle having
nontrivial projectivization.  For example,
\(\OO_C(-P)\oplus\OO_C^{n-1}\) works on every such pointed curve.
The curve is intrinsic as the reduction; a chosen embedding of the
curve is not part of the isomorphism data.

\subsection{A moving family, not only its individual fibres}
The same finite object also reconstructs a pencil that varies over its
projective reduction.  Let \(B\) be a smooth connected projective
complex variety, let \(\mathcal U\) have rank \(n\), and let \(\mathcal R\) be a rank-two subbundle of
\(\Sym^2V\otimes\OO_B\) with locally free quotient.  Use the local
ideal \((\det T)I_p(\gamma\Sym^2T)\) in
\(\Hom(\mathcal U,V)\) and truncate at order \(d\).
Theorem~\ref{thm:unrestricted-moving-v146} states that two resulting
abstract finite schemes are isomorphic exactly when their data are
related by an isomorphism of reduced bases, one constant linear left
map, and an actual isomorphism of the source bundles, carrying one
pencil subbundle to the other.  The projection and both tensor factors
are forgotten in the input.  The varying coefficient line is recovered
as a subbundle, not by pointwise choices of projective coordinates.

This strengthens the constant-pencil curve theorem in two directions:
the reduced base and source bundle may vary in the comparison, and the
pencil need not be constant on that base.  Example~\ref{ex:nonconstant-pencil-v146}
has a nonconstant unordered cross-ratio and six spectral collision
points, all retained by the same finite scheme.  Theorem~\ref{thm:automorphism-kernel-v146}
then identifies the remaining freedom in an isomorphism.  Its quotient
is the group of the recovered linear data; its kernel acts trivially
on the associated graded scheme and has a terminating commutator
filtration.  Explicit nonlinear shears show why that kernel cannot be
omitted.  These statements are consequences of the unmarked inverse,
not additional claims of novelty for the classical Pluecker map or
for spectral cross-ratios.

\subsection{A single unmarked local algebra suffices}
The strongest form of the inverse needs no positive-dimensional support.
Choose an \(n\)-dimensional source space \(U\), write \(T=(t_{ij})\)
for its universal map to \(V\), and let \(\mathfrak m=(t_{ij})\).
Define
\[
 \mathfrak A_R^{[d]}=
 \C[t_{ij}]/\bigl((\det T)I_p(\gamma_R\Sym^2T)
                                      +\mathfrak m^{d+1}\bigr).
\]
\begin{theorem}[Local algebraic inverse]\label{thm:principal-local-v146}
For every pair of complex quadratic pencils, an ungraded algebra
isomorphism \(\mathfrak A_R^{[d]}\simeq\mathfrak A_{R'}^{[d]}\)
exists if and only if \(R'=gR\) for some \(g\in\PGL(V)\).
The order \(d=n^2+2n-4\) is the smallest uniform order already for
these single-point objects.
\end{theorem}
The full statement and proof are
Theorem~\ref{thm:artin-local-inverse-v146}.  Proposition~\ref{prop:first-relation-orbit-v147}
separates the general truncation argument from the pencil-specific
orbit rigidity: nonlinear coordinate changes form an explicitly described
unipotent kernel, whereas the first relation space carries the
reconstruction data.  The extra local ingredient
is that the residual coefficient module has factor-support ranks
\((1,\binom N2)\).  Transposition reverses these unequal ranks and
therefore cannot carry one pencil ideal to another.  This orients the
unordered matrix factors without a nontrivial projective bundle.
Theorem~\ref{thm:unrestricted-moving-v146} consequently removes that
restriction from the moving-family theorem as well.  This is not a
claim that the reduced determinant cone alone remembers the pencil:
the nonreduced first relation kernel and its residual coefficients
are indispensable.  Proposition~\ref{prop:local-automorphisms-v146}
also identifies the entire tangent-identity kernel, including its
affine-space parametrization and nonadditive group law.

\subsection{Conductors, collision singularities, and actual failure orbits}
Normalization alone does not describe the gluing of its branches.
Theorems~\ref{thm:split-conductor-v152} and
\ref{thm:binary-pinch-v152} determine two successive local models.
On the coprime two-divisor stratum the completed image is
\[
       \C[[z,x_1,\ldots,x_c,y_1,\ldots,y_c]]/(x_i y_j).
\]
Its conductor is $(x,y)$, its intrinsic singular ideal is $(x,y)^c$,
and its depth is $\dim(z)+1$.  When two common binary roots collide,
the model is instead
\[
 \C[[u,\Delta]]\oplus(u_1,\ldots,u_m)t
          \ \subset\ \C[[u,t]],\qquad \Delta=t^2,
\]
up to a specified smooth factor.  Its conductor has a ramified
double cover in the normalization.  These formulas determine the
Cohen--Macaulay cases and distinguish the singular scheme from its
reduced support.  Proposition~\ref{prop:degree-intersections-v152}
organizes the reduced intersections of different divisor-degree
images by the common part of two chosen divisors.

The boundary also has a direct role in the pencil inverse.
Theorem~\ref{thm:universal-boundary-v152} constructs a common point
in every intrinsic first-relation orbit closure of a pencil failure
algebra.  A scalar-matrix contraction of the flag-pencil relation
has gcd $a^{D-q_n}$, where $q_3=3$ and $q_n=4$ for $n\geq4$.
The selected degree-$n(n-1)$ divisor is unique but ramified; its
normalization fibre has tangent dimension
$\binom{n^2+3n-5-q_n}{3n-4-q_n}-1$.
Proposition~\ref{prop:pure-power-fibre-v152} gives its entire Artin
algebra by reciprocal-polynomial equations.  This specializes
\emph{actual} failure algebras through flat families; it does not
merely place an unrelated common-factor example beside the inverse
theorem.  The exact order of the flag coefficient filtration, proved
in Lemma~\ref{lem:flag-jets-v152}, is the link between the pencil
representation and the divisor collision.

\subsection{Factorization loci and the geometric comparison}
The normalization mechanism has classical precedents.
Chipalkatti's coincident-root locus parametrizes one binary form
with a prescribed root partition.  His Proposition~5.1, Theorem~5.4,
and Corollary~5.8 distinguish multiple normalization points from
failure of tangent injectivity and describe the singular support
\cite{Chipalkatti2003}.  Kurmann gives local incidence equations,
proves their scheme-theoretic equality in Theorem~1.5, and expresses
smoothness through splittings of partitions in Proposition~2.1
\cite{Kurmann2012}.  In particular, uniqueness plus tangent
injectivity is not a new general principle of the present paper.

Here the datum is an $r$-plane of forms in $e$ variables, not one
binary form; the selected divisor has prescribed degree, not a
root partition.  The residual datum is a Grassmannian subspace,
and for $e>2$ divisor degrees need not fill an interval.  These
changes produce high-codimension conductor strata and non-Cohen--Macaulay
branch gluings.  Theorems~\ref{thm:split-conductor-v152} and
\ref{thm:binary-pinch-v152} compute the image rings, not merely
the incidence graph or singular support.  Ferrand's conductor-square
formalism \cite{Ferrand2003} is the standard algebraic framework;
the particular conductor ideals, singular Fitting ideals, and depths
are calculated here for these Grassmannian strata.

There is a second, closer binary comparison.  Hu--Lin--Shao
\cite[\S2.2, Corollary~3.6, and Theorem~1.1]{HuLinShao2007}
study tuples of binary forms, define common-factor strata using
resultant homomorphisms, and resolve the map-space boundary by
successive blow-ups.  On the open set of linearly independent
$r$-tuples, passing to their span is the frame quotient onto
$\Gr(r,S_D)$; the common-divisor condition is invariant under this
quotient.  Thus our binary loci genuinely meet that theory rather
than forming an unrelated example.  We do not identify its
resultant scheme structure with our image structure without a
scheme comparison, nor claim a new general resolution.  Our
multivariate local rings and the pencil-orbit boundary are the
specific additional conclusions.

Rigidification is used as infrastructure, not as an independent
headline innovation.  The appropriate primary general result is
Abramovich--Olsson--Vistoli, Appendix~A, Theorem~A.1
\cite{AOV2008}, for compatible flat finitely presented normal
inertia subgroups, without a centrality assumption.  Our argument
identifies the particular substitution subgroup and twisted
right-factor subgroup and retains their actual group laws.  The
boundary statement is on the maximal-lower-Hilbert-function
first-relation stratum throughout.

The comparison with regular-pencil Segre strata is made in
Remark~\ref{rem:segre-boundary-v152}.  The common closed flag orbit
is not presented as a new classification of congruence orbits.
The new pencil-side conclusion is its explicit ramified boundary
under intrinsic cotangent changes, together with the complete
normalization-fibre algebra.  The separate documentary status of
Ballico's 1993 paper, discussed below, remains unchanged: absent
its full theorem text, no nonanticipation conclusion is drawn.

\subsection{The intrinsic mechanism}
In the graph neighbourhood of \(B_R\), with tautological bundle
\(\mathcal U\) and graph map \(T:\mathcal U\to V\), the ideal is
\[
                 (\det T)I_p(\gamma_R\Sym^2T).
\]
Its generators have degree \(d\).  This degree calculation alone
only explains why lower orders are constant.  The inverse requires
three further intrinsic operations.  Multiplication in the finite
scheme recovers the first relation kernel and hence the whole
homogeneous cone.  The rank-one Fano scheme of its reduced determinant
cone then recovers the two tensor rulings.  On the socle Grassmannian
one may distinguish them by projective triviality, already on a
Schubert line.  In general the unequal coefficient-support ranks
orient them without that restriction.  Properness forces the induced
map to \(\PGL(V)\) to be constant.  After choosing its linear lift,
the actual normal-bundle map lies in the commutant of \(\End(V)\),
which recovers the right bundle without a line twist
(Lemma~\ref{lem:actual-factor-descent-v147}).  Finally, cancellation of the intrinsic
determinant ideal recovers the coefficient module.  For pencils its
exterior representation is irreducible, and exterior duality identifies
its single coefficient line with the Pluecker line of \(R\).

Theorem~\ref{thm:structural-coefficient-reconstruction} separates this
geometric mechanism from the particular multiplication algebra.  The
curve theorem shows why the mechanism is not tied to the dimension of
the original parameter Grassmannian.  The geometric orientation used here has a local replacement: coefficient-support asymmetry.  Theorem~\ref{thm:artin-local-inverse-v146} uses it to reconstruct the pencil from one unmarked Artin local algebra.

Theorem~\ref{thm:general-moving-coefficients-v147} gives the broader
principle.  It reconstructs simultaneously varying subspaces in
several Schur modules from the unmarked first-relation thickening,
provided one coefficient support is nonzero and proper.  The
coefficient spaces may have arbitrary prescribed ranks; they need
not be lines coming from quadratic pencils.  Its proof isolates the
actual factor-descent step from classical linear-preserver theory.
Example~\ref{ex:isotrivial-covers-v147} shows why the global descent is
not equivalent to fibrewise reconstruction: two Zariski locally
trivial families have the same Artin fibre and isomorphic nilradical
graded vector bundles in every degree, but their abstract total
schemes differ.  The inverse detects the inequivalent degree-three
coefficient maps.  This is the mechanism behind the moving-pencil
theorem, rather than
an assertion of novelty for canonical grading or for the numerical
value of the first relation degree.


\subsection{An exact image criterion and a pencil-native global consequence}
The coefficient inverse admits an intrinsic converse.
Theorem~\ref{thm:recognition-dichotomy-v148} starts with an
unmarked first relation whose homogeneous radical is a determinant
cone.  After recovering its unordered rulings, invariance under one
full factor group is necessary and sufficient for the residual
relation to come from coefficient subspaces in the Cauchy modules.
Both orderings satisfy this test precisely when every coefficient
support is zero or full.  Thus the strict-support condition is
exactly the criterion for a unique admissible orientation in the
recognized class.  Theorem~\ref{thm:exact-coefficient-stabilizer-v148}
computes its entire affine stabilizer, including the transpose
component and the universal nonlinear automorphism kernel.
Corollary~\ref{cor:intrinsic-orbits-v148} expresses the inverse as
an equivalence of geometric orbit sets, with those ambiguities
specified rather than suppressed.

A global consequence lies inside the quadratic-pencil problem
itself.  Fix a plane \(H=\langle x,y\rangle\subset V\).
For a degree-\(m\) map \(f=[a:b]:\Pj^1\to\Pj(H)\), take
the pencil subbundle
\[
                  \mathcal R_f=(ax+by)H.
\]
Every fibre is equivalent to \(\langle x^2,xy\rangle\).
Nevertheless, Theorem~\ref{thm:covering-pencil-moduli-v148} proves
\[
 Y_f\simeq Y_{f'}\quad\Longleftrightarrow\quad
 f'=\alpha\circ f\circ\sigma,\qquad
                         \alpha,\sigma\in\PGL_2(\C),
\]
for their unmarked order-\(d\) failure thickenings with trivial
source bundle.  Their fibre algebras and every nilradical graded
vector bundle depend only on \(n,m\), not on \(f\).
Their abstract pencil, quotient and first-relation bundles are
fixed as well.  Thus the entire covering equivalence problem, not
only one chosen pair, is retained in the global multiplication and
is invisible to those additive data.  The maps \(z^3\) and
\(z^3+z\) give an explicit separating pair of quadratic pencils
(Corollary~\ref{cor:pencil-native-pair-v148}).

Neither complete reducibility of Cauchy modules, linear preservers
of a Segre cone, nor division by the common factor of a pencil is
new in isolation.  Their synthesis here has a precise output:
recognition of the image and its exact orientation ambiguity after
all matrix markings are forgotten, followed by recovery of a
whole covering map from a locally isotrivial pencil thickening.
A projective factor map alone would not recover the actual source
bundle; the intrinsic normal map and its commutant do so.  A
pointwise pencil inverse would not recover a covering map; one
constant left transformation over the entire proper base does so.
These distinctions specify the geometric content beyond the
classical ingredients.  The recognized coefficient class is not the
class of all determinant-radical relations.  Its geometric orbit
statement alone does not describe nonreduced deformation groupoids;
the relative comparison below uses a further scheme-theoretic inverse.


\subsection{Intrinsic strata, infinitesimal parameters, and pencil moduli}
The relative inverse begins with a geometric condition on the first
relation, not invariance under a factor group.  For a space \(K\) of
homogeneous forms, the divisor-incidence functor defines a canonical
scheme-theoretic image in the relation Grassmannian.  Its exact-gcd open
is the locally closed stratum of Lemma~\ref{lem:universal-gcd}, now
identified by its incidence universal property in
Corollary~\ref{cor:canonical-exact-v151}.  Multiplication by the
universal common-divisor line is an isomorphism onto this open.
Consequently the common factor is recovered even over a nonreduced
parameter base.  The condition is scheme-theoretic: merely asking all
geometric fibres to have the specified common factor is insufficient.
When that divisor is a power of a determinant, the remaining
parameter is an arbitrary gcd-free residual subspace, modulo the
full determinant-preserver group
(Theorem~\ref{thm:determinant-stratum}).  This is a smooth quotient
stack, and neither of the two factor symmetries is required.

For pencils the exact common divisor is \(\delta^{n-1}\), and the
primitive residual relation has degree \(3n-4\)
(Lemma~\ref{lem:primitive-pencil}).  The Cauchy projection and the
Pluecker equations recognize its image as a locally closed scheme,
not just a set of complex points.  The resulting relative theorem is
\begin{equation}\label{eq:intro-effective-moduli}
 \mathscr F_{\rm pen}^{\rm eff}\simeq
                [\Gr(2,\Sym^2V)/\PGL(V)].
\end{equation}
Here the superscript records two explicit operations on morphisms:
first remove tangent-identity nonlinear substitutions, then remove
the ineffective right matrix factor, and fppf stackify.
Theorem~\ref{thm:pencil-stack-equivalence} proves this equivalence
over arbitrary complex parameter schemes, retaining all singular
pencils.  It does not identify the unrigidified algebra stack with
the pencil stack.  Its tangent complex is
\([\mathfrak{pgl}(V)\to\Hom(R,\Sym^2V/R)]\), so the effective
stack is smooth of dimension \(n-3\), with precisely the pencil
stabilizer jumps.  Rank-incidence and spectral Fitting strata retain
their scheme structures and deformation equations by
Corollary~\ref{cor:pencil-specialization-stack}.  This is an exact-image and rigidification comparison; the displayed
rank and spectral consequences are its functorial restrictions, not an
independent classification of those loci.  The two quotient operations
have the universal properties of
Theorem~\ref{thm:rigidification-v151} and
Corollary~\ref{cor:two-rigidifications-v151}.

The ambient determinant-divisor stratum is substantially larger.
Theorem~\ref{thm:transverse-pencil-directions} computes its normal
space to the pencil stratum in three parts: non-Pluecker lines,
non-scalar right-factor directions, and other Cauchy summands.
Theorem~\ref{thm:fixed-support-grassmannians} constructs through every
pencil a projective Grassmannian of relations with the same determinant
radical, containing a projective line whose other points admit neither
factor symmetry.  Thus the family inverse describes how the
coefficient and pencil loci sit in, and can be left within, a larger
intrinsic stratum.  It does not falsely infer deformation stability
from fibrewise recognition.

\subsection{The normalization of a natural common-divisor boundary}
There is a different geometric question beyond the pencil image:
what happens when a divisor of a first relation is no longer uniquely
recoverable in a family?  Let \(E\) have dimension \(e\geq2\), set
\(S_j=\Sym^jE\) and \(s_j=\dim S_j\), and fix
\(g,h\geq1\), \(D=g+h\), and \(2\leq r\leq s_h\).
Let \(Y_g\subset\Gr(r,S_D)\) be the scheme-theoretic image of the
incidence of a line of degree-\(g\) divisors of a rank-\(r\) relation
bundle.  This definition uses the vanishing of the universal quotient
map; it does not first recover a pencil, impose a factor symmetry,
or choose a reduced structure on a set.

\begin{theorem}[The divisorial boundary]
\label{thm:principal-boundary-v151}
The multiplication morphism
\[
 \Pj(S_g)\times\Gr(r,S_h)\longrightarrow Y_g,
                         \qquad ([f],J)\longmapsto fJ
\]
is finite and is the normalization.  The locus \(Y_g\) is integral,
of dimension \(s_g-1+r(s_h-r)\).  At a geometric relation \(K\)
it is normal exactly when \(F_K=\gcd(K)\) has one degree-\(g\)
divisor \([f]\) and \(\gcd(f,F_K/f)=1\).
At the preimage \(([f],K/f)\), if this last gcd has degree \(k\),
the fibre tangent dimension is \(s_k-1\).
\end{theorem}

The theorem is proved in Section~\ref{sec:canonical-boundary-v151}.
Lemma~\ref{lem:fibre-equations-v151} supplies equations for the
entire normalization fibre, not only its tangent space.  For example,
\(K=\langle x^2y,x^2z\rangle\) has a unique linear divisor, but the
normalization fibre is \(\Spec\C[a,b]/(a,b)^2\), of length three.
A different relation has two reduced divisor lifts; yet another has
a gcd-degree jump at a normal point.  The exact criterion distinguishes
these behaviours.  Corollary~\ref{cor:algebra-boundary-v151} carries
this intrinsic boundary to the first-relation algebra stack.  It is
not a classification of all finite-flat smoothings or of all linear
orbits in that boundary.

A direct consequence on the pencil side is proved separately in
Section~\ref{sec:closed-pencil-v151}.  Every quadratic pencil admits
an explicit congruence degeneration to
\(\langle x_1^2,x_1x_2\rangle\); its orbit is the closed flag
variety \(\operatorname{Fl}(1,2;V)\), and it is the unique closed
orbit in the full pencil Grassmannian.  The construction is uniform
for regular and singular pencils and lifts to a finite-flat family
of their failure algebras with exact common factor
\(\delta^{n-1}\).  This proves a particular universal orbit-closure
adjacency by equations, rather than by copying a prescribed invariant
locus through the effective-stack equivalence.  Its elementary
highest-weight mechanism is distinct from the normalization theorem,
which deals with a larger boundary of arbitrary first relations.

\subsection{Framed families and spectral consequences}
Once tensor coordinates are retained, the parameter map is simpler.
The first relation spaces give a closed immersion of
\(\Gr(2,W)\) into a Grassmannian, by
Theorem~\ref{thm:finite-parameter-immersion}.  Its proof factors through
the classical Pluecker embedding, tensoring a line with a fixed space,
and a fixed linear inclusion.  This is the framed encoding consequence
of the coefficient calculation, not a substitute for the unmarked
inverse.  Theorem~\ref{thm:universal-finite-neighbourhood} glues the
actual universal order-\(d\) neighbourhood.  It is finite locally
free over the relative socle of rank
\(\binom{n^2+d}{d}-\binom N2\), and projective and flat over the
pencil parameter scheme.  It commutes with arbitrary complex base
change with that socle retained.  Over a nonreduced base, the specified
relative socle ideal need not be the absolute nilradical.  These marked statements alone do not identify unmarked moduli
stacks or deformation groupoids.  The common-divisor descent and
ineffective-inertia analysis in
Theorem~\ref{thm:pencil-stack-equivalence} supply that separate step.

For regular pencils, the spectral sheaf and its Fitting incidence
schemes are read from the reconstructed pencil.  We retain these
consequences because they exhibit information genuinely lost by
coarser invariants: fixed discriminants and fixed reduced rank loci
can conceal different elementary divisors.  In
Theorem~\ref{thm:fixed-spectral-classification}, every classical
dominance specialization is transported to a flat family of the
finite neighbourhoods, with every strict transition detected at
order \(d\).  The spectral orbit classification itself is classical,
not a second independent Torelli theorem.

\subsection{Prior work and the publication object}
The Weierstrass--Segre classification is used in its root-labelled
form; compare \cite[Theorem~1.1 and Corollary~2.1]{FevolaMandelshtamSturmfels}.
The orthogonal nilpotent input is isolated in
Lemma~\ref{lem:orthogonal-orbits-v145}, with precise references to
\cite[Proposition~1 and Theorem~1]{Ohta1986}, an independent dimension
calculation, and an explicit distinction between full orthogonal
orbits and their special-orthogonal components.

Ballico's relative multiplication construction
\cite[Theorem~0.2; \S1, Remark~1.2 and equation~(6)]{Ballico96}
is a substantive antecedent for studying quadratic-normality failure
on finite linear sections of an embedded surface.  Here the defining
algebra varies, and the inverse starts from an unmarked infinitesimal
failure scheme.  The opening page of the distinct earlier paper \cite[p.~5]{Ballico93}
concerns higher-order embedding properties and restriction maps to
finite subschemes.  Its remaining theorem and proof text has not
been inspected; that first page does not establish the absence of
an inverse statement.  No anticipation or nonanticipation conclusion
relative to that article is asserted.

The Artin-local aspect is compared at theorem level with
\cite{EliasRossiShort,EliasRossiCompressed} in
Section~\ref{sec:first-relations-v147}, including the distinction
between canonical grading of a nonhomogeneous local algebra and
orbit rigidity of our already homogeneous first relation.  The
linear-preserver input is classical \cite{MarcusMoyls}; its precise
hypotheses and a proof in the required case are included in
Lemma~\ref{lem:linear-preserver-v147}.  The algebraic-group
quotient step is supplied by the homomorphism theorem
\cite[Theorem~5.39 and Remark~5.42]{Milne2017}; its use in the
exact stabilizer calculation distinguishes a scheme-theoretic
quotient from a bare bijection of complex points.

The present article consists of intrinsic reconstruction, the canonical
common-divisor boundary, relative rigidification, explicit pencil
specialization, and spectral readout.  The reciprocal-likelihood theorem and projective critical
schemes are preserved in a separate applications manuscript.  Its
finite-flat critical divisor requires supplied split semisimple data
and simple disjoint projective poles; its totally real cover requires
a supplied real definite realization and separated data.  None of
those hypotheses enters the all-pencil theorem here.  Independent
web, exterior-operator, K3, polarized, and boundary results are
preserved in an explicitly non-submitted research archive.  They are
not inputs to any proof in this article and are not part of its
significance claim.  Every proof required here is included here.

% BEGIN PRESERVED PART 35-coefficient-symmetries-v148.tex
\section{Recognition and the exact ambiguity of coefficient systems}
\label{sec:coefficient-symmetries-v148}

The support inequality in the preceding section is not merely a
sufficient test in the class under consideration.  We characterize
that class intrinsically and determine exactly when its orientation
is unique.  We then compute the full linear stabilizer, including
the cases in which an orientation is lost.  These are statements
about unmarked first relations, not about a supplied matrix
presentation.

Throughout this section \(n\ge2\), \(q\ge1\), and \(D=n+q\).
Let \(E\) have dimension \(n^2\), put \(S=\Sym E\), and let
\(0\ne K\subset S_D\).  The affine cone is in \(E^*\).
Suppose that the radical of the homogeneous ideal \((K)\) is a
reduced hypersurface linearly equivalent to the determinant
hypersurface in square \(n\)-by-\(n\) matrices.  We call this the
\emph{determinant-radical condition}.  The ideal \((K)\) is taken
in the full symmetric algebra, not in its Artin truncation.
Its reduced equation gives a line \(\C\delta\subset S_n\),
and unique factorization gives the intrinsic residual subspace
\begin{equation}\label{eq:residual-recognition-v148}
                   J=K/\delta\ \subset S_q.
\end{equation}
Changing the generator \(\delta\) does not change this subspace.
The reduced singular loci of the recovered cone determine its
unordered tensor rulings.  An ordering identifies
\(E=V^*\otimes U\), with \(\dim U=\dim V=n\), up to
separate factor changes.  The subgroup induced by
\(\operatorname{GL}(U)\) on the second factor is independent of
those changes.  An ordering is called \emph{admissible} when
\(J\) is invariant under this full subgroup.  Admissibility is
therefore an intrinsic condition on an ordered ruling, and does
not ask for a chosen basis of either factor.

\begin{theorem}[Intrinsic recognition and orientation dichotomy]
\label{thm:recognition-dichotomy-v148}
Under the determinant-radical condition, an ordering is admissible
if and only if its Cauchy decomposition has
\begin{equation}\label{eq:recognized-cauchy-v148}
 J=\bigoplus_{\lambda\vdash q,\,\ell(\lambda)\le n}
         L_\lambda\otimes\mathbb S_\lambda U,
 \qquad L_\lambda\subset\mathbb S_\lambda V^*.
\end{equation}
The coefficient spaces are unique and are recovered by projection
and contraction.  Suppose at least one ordering is admissible and
put \(r_\lambda=\dim L_\lambda\) and
\(m_\lambda=\dim\mathbb S_\lambda\C^n\).  Then exactly one
of the following alternatives occurs.

\smallskip
\noindent (i) Some \(0<r_\lambda<m_\lambda\).  Precisely one
of the two orderings is admissible, and every linear isomorphism
between such relation spaces preserves that ordering.

\smallskip
\noindent (ii) Every \(r_\lambda\) is either zero or
\(m_\lambda\).  Both orderings are admissible.  The residual
space is a sum of full Cauchy summands, and transposition preserves
\(K\) after identifications of the two matrix factors.

\smallskip
Conversely, every nonzero system in
\eqref{eq:recognized-cauchy-v148} satisfies the determinant-radical
condition after multiplication by \(\delta\).  Thus the theorem
characterizes the image of the coefficient construction, without
assuming a coefficient presentation as part of the input.
\end{theorem}
\begin{proof}
The determinant polynomial is irreducible for \(n\ge2\).
Every element of \(K\) belongs to the radical \((\delta)\),
so \(K\subset\delta S_q\) and the cancellation in
\eqref{eq:residual-recognition-v148} is legitimate.  The intrinsic
recovery of the unordered factors follows from
Lemma~\ref{lem:linear-preserver-v147} and the rank-one rulings.
Changing either factor frame conjugates its full general linear
group into itself; scalar changes act as scalars on \(S_q\).
Consequently the definition of admissibility is independent of
all these choices.

For an ordering, the Cauchy formula is
\[
 S_q=\bigoplus_{\lambda\vdash q,\,\ell(\lambda)\le n}
      \mathbb S_\lambda V^*\otimes\mathbb S_\lambda U.
\]
As representations of \(\operatorname{GL}(U)\), the nonzero
\(\mathbb S_\lambda U\) are irreducible and pairwise
nonisomorphic.  Finite-dimensional rational representations of
this group over \(\C\) are completely reducible.  Hence an
invariant subspace splits across its isotypic components.  In the
\(\lambda\)-component, evaluation identifies that subspace with
\[
 \Hom_{\operatorname{GL}(U)}(\mathbb S_\lambda U,J)
                              \otimes\mathbb S_\lambda U.
\]
Schur's lemma embeds the displayed multiplicity space uniquely in
\(\mathbb S_\lambda V^*\); this is \(L_\lambda\).
This proves necessity in \eqref{eq:recognized-cauchy-v148},
including the absence of unrecorded cross-summand graphs.
Sufficiency follows since the second factor is full.  Projection
onto a Cauchy summand followed by contraction with the dual of its
second factor recovers \(L_\lambda\).

If the reverse ordering is also admissible, \(J\) is invariant
under both full factor groups.  Their product acts irreducibly on
each Cauchy summand, and these product representations are pairwise
nonisomorphic.  Complete reducibility now makes \(J\) a sum of
whole summands.  Equivalently, all \(r_\lambda\) are zero or
full.  Conversely a sum of whole summands is invariant under both
groups and under the factor exchange.  This proves the dichotomy.
One may also see uniqueness directly from the contraction ranks
\((r_\lambda,m_\lambda)\): a factor exchange reverses them,
which is impossible at a nonzero proper support.  Every linear
isomorphism of relation spaces preserves the determinant cone;
Lemma~\ref{lem:linear-preserver-v147} leaves just the two
factor alternatives.  It must therefore preserve the unique
admissible ordering in case (i).

Finally choose a nonzero \(L_\lambda\).  At an invertible
matrix \(T:U\to V\), the map \(\mathbb S_\lambda T\) is
invertible.  A nonzero covector in \(L_\lambda\) consequently
has a nonzero matrix coefficient at \(T\).  Thus the residual
coefficients do not vanish simultaneously at any invertible
matrix.  All elements of \(\delta J\) vanish on the determinant
hypersurface.  Their common zero set is exactly that hypersurface;
the Nullstellensatz and its reduced irreducible equation give
\(\rad(\delta J)=(\delta)\), as asserted.
\end{proof}

The complete reducibility and Cauchy statements here are classical
representation theory; see \cite{ChengWang}.  Their role is to
supply an intrinsic image criterion \emph{after} the matrix factors
have been recovered.  For a relation space that fails the
admissibility criterion, the theorem does not assert a one-sided
coefficient description.  In particular it is not a classification
of arbitrary homogeneous ideals or arbitrary degeneracy schemes.

\subsection{Stabilizers and the full transpose ambiguity}
Use the normal action \(T\mapsto gTh^{-1}\).  Its action on
polynomial functions is inverse pullback, so it gives an honest
representation \(\rho\) on \(E=V^*\otimes U\).
The kernel of \(\rho:\operatorname{GL}(V)\times
\operatorname{GL}(U)\to\operatorname{GL}(E)\) is the diagonal
scalar subgroup.  For a recognized system define
\[
 G_L=\{g\in\operatorname{GL}(V):
          (\mathbb S_\lambda g^{-1})^*L_\lambda=L_\lambda
                                      \text{ for every }\lambda\},
 \qquad H_K=\operatorname{Stab}_{\operatorname{GL}(E)}(K).
\]
These are closed subgroup schemes, with the stabilizer scheme
structures from their actions on Grassmannians.

\begin{theorem}[Exact coefficient stabilizer]
\label{thm:exact-coefficient-stabilizer-v148}
In case (i) of Theorem~\ref{thm:recognition-dichotomy-v148},
\begin{equation}\label{eq:proper-stabilizer-v148}
                 H_K\simeq
        (G_L\times\operatorname{GL}(U))/\mathbb G_m.
\end{equation}
In case (ii), after choosing bases for the two factors,
\begin{equation}\label{eq:full-stabilizer-v148}
 H_K\simeq
 \bigl((\operatorname{GL}_n\times\operatorname{GL}_n)
                         /\mathbb G_m\bigr)\rtimes C_2,
\end{equation}
where \(C_2\) is generated by transposition.  These are
isomorphisms of affine algebraic groups over \(\C\), and hence
of their represented functors on commutative complex algebras.

For the ungraded algebra
\(A(K)=S/((K)+\mathfrak m^{D+1})\), they give the split exact
sequence
\begin{equation}\label{eq:complete-local-aut-v148}
 1\longrightarrow U_K\longrightarrow\operatorname{Aut}(A(K))
                      \longrightarrow H_K\longrightarrow1.
\end{equation}
Here \(U_K\) is the universal higher-order substitution group
from Proposition~\ref{prop:first-relation-group-v147}, of
dimension \(n^2(\length A(K)-1-n^2)\).
\end{theorem}
\begin{proof}
Every element of \(H_K(\C)\) preserves the radical of \((K)\),
and therefore the determinant cone.  In case (i), the unique
admissible orientation excludes its transposed component.  The
remaining map is \(\rho(g,h)\).  Cancellation of the determinant
line and projection onto each Cauchy summand show that it preserves
\(J\) if and only if \(g\in G_L\); the arbitrary right map
acts on the full factor and introduces no additional condition.
The kernel is the diagonal \(\mathbb G_m\).  In case (ii),
every separate factor transformation and transposition preserves
each active full summand.  These exhaust the determinant preservers,
so they exhaust \(H_K(\C)\).  The transposition has order two
and lies outside the preserving component for \(n\ge2\), giving
the asserted semidirect product.  Its action on the left-right
quotient is conjugation by the matrix-transpose map; no particular
formula for its action on chosen factor lifts is needed.

We justify the scheme statement as well.  The homomorphism theorem
for algebraic groups identifies the quotient by the kernel with
the closed image \cite[Theorem~5.39 and Remark~5.42]{Milne2017}.
Every finite-type group scheme over a characteristic-zero field is
smooth \cite[Lemma~39.8.2]{StacksGroups}.  Thus the stabilizer and
the image are reduced closed subgroup schemes.  The equalities on
complex points just proved imply equality of their defining radical
ideals.  This gives \eqref{eq:proper-stabilizer-v148} and
\eqref{eq:full-stabilizer-v148} as algebraic-group isomorphisms,
not only as bijections on complex points.  Finally
Proposition~\ref{prop:first-relation-group-v147} applies in degree
\(D\) and gives \eqref{eq:complete-local-aut-v148} with the
stated kernel, splitting and dimension.
\end{proof}

For completeness, the omitted zero system has \(K=0\).
Then \(A(K)\) is the free truncated symmetric algebra,
\(H_K=\operatorname{GL}(E)\), and its linear automorphisms
need not preserve any matrix factors.  The determinant-radical
condition deliberately excludes this third possibility.  A full
nonzero coefficient system and a zero system therefore lose
information in different ways: the former retains unordered
matrix geometry and has exactly a transpose ambiguity, while the
latter retains no intrinsic matrix geometry at all.

\subsection{The zero/full case over a projective reduction}
The local transpose ambiguity also has an exact global form.  Fix
a nonempty set of active partitions and take the corresponding
\(\mathcal L_\lambda\) to be full, with all other coefficient
bundles zero.  Use the same construction as in
\eqref{eq:coefficient-thickening-v147}, without its strict-support
hypothesis.

\begin{proposition}[Global ambiguity for full coefficient systems]
\label{prop:full-global-ambiguity-v148}
Let \(Y_i\) be two such schemes, for the same \(n,q\) and
active partitions, over smooth connected projective bases \(B_i\)
with source bundles \(\mathcal U_i\).  They are isomorphic as
abstract complex schemes if and only if there are an isomorphism
\(\sigma:B_1\to B_2\) and an actual bundle isomorphism
\(\mathcal U_1\simeq\sigma^*\mathcal U_2\).
An isomorphism exchanging the two ruling families over a fixed
\(\sigma\) exists precisely when, in addition, the common source
bundle is projectively trivial.  Equivalently,
\[
 \mathcal U_1\simeq\mathcal M^{\oplus n},\qquad
 \sigma^*\mathcal U_2\simeq\mathcal M^{\oplus n}
\]
for a line bundle \(\mathcal M\) on \(B_1\).
Thus even the full systems recover the actual source-bundle class;
the extra ambiguity is an exchange, not an undetected line twist.
\end{proposition}
\begin{proof}
An abstract isomorphism recovers the reduction, the actual normal
bundle map, the first relation, and the unordered determinant
rulings by the proof of
Theorem~\ref{thm:general-moving-coefficients-v147}, which uses
strict support only to orient the rulings.  Over the connected
base the choice to preserve or exchange the two components is
constant.  In the preserving case, projectivity makes the left
projective map constant.  A constant lift and
Lemma~\ref{lem:actual-factor-descent-v147} give an actual
isomorphism of the source bundles.

In the exchanging case each source projective bundle is identified
with the constant left projective bundle on the other side.  It
is therefore trivial.  A vector bundle with a trivialized
projectivization is a line twist of a constant vector space:
locally lift the given projective trivialization to linear frames;
on intersections their differences are scalar matrices, whose
scalar cocycle is the required line bundle.  After pullback by
\(\sigma\), write
\(\mathcal U_i^*=\mathcal M_i^*\otimes W_i\), with \(W_i\)
constant of dimension \(n\).  The normal bundles become
\(\mathcal M_i^*\otimes(V_i\otimes W_i)\).
Their actual isomorphism induces a morphism from \(B_1\) to
the projective-linear isomorphisms between two constant
\(n^2\)-spaces.  This target is affine; the morphism is constant.
Choose its constant linear lift.  On removing it, the actual
bundle map is scalar on each projective fibre.  In local line
frames its off-diagonal entries vanish and all diagonal entries
agree.  Reducedness makes these equalities hold as bundle maps.
Their common entry is a nowhere vanishing section of
\(\mathcal Hom(\mathcal M_1^*,\mathcal M_2^*)\), and hence
\(\mathcal M_1\simeq\mathcal M_2\).  This proves necessity
of the common actual line twist, and thus of the source-bundle
isomorphism.

Conversely any actual source-bundle isomorphism, together with a
constant left identification, preserves all active full Cauchy
summands, so it gives an isomorphism of the schemes.  If the
source bundles have the common line-twist presentation above,
choose identifications of the constant factor spaces and
transpose their tensor factors, leaving the line factor intact.
This is a global actual normal map.  It preserves the determinant
line, every active full summand, and the truncation, and supplies
an exchanging isomorphism.  These constructions prove all the
claims.
\end{proof}

\subsection{The intrinsic orbit space}
Fix a rank vector \(\boldsymbol r\) with at least one proper
nonzero entry and form the product of Grassmannians
\[
 X_{\boldsymbol r}=
   \prod_{\lambda\vdash q,\,\ell(\lambda)\le n}
       \Gr(r_\lambda,\mathbb S_\lambda V^*).
\]
Zero and full factors in this product are points.  Let
\(\mathcal Z_{\boldsymbol r}\) denote the set of subspaces of
\(\Sym^D\C^{n^2}\) satisfying the determinant-radical condition
and the intrinsic admissibility test, with ranks
\(\boldsymbol r\) in the unique admissible orientation.

\begin{corollary}[Geometric orbit classification]
\label{cor:intrinsic-orbits-v148}
The coefficient construction induces bijections
\begin{equation}\label{eq:intrinsic-orbits-v148}
 X_{\boldsymbol r}(\C)/\PGL(V)
  \ \simeq\ \mathcal Z_{\boldsymbol r}/\operatorname{GL}_{n^2}(\C)
  \ \simeq\ \{A(K):K\in\mathcal Z_{\boldsymbol r}\}/\simeq.
\end{equation}
Their stabilizers, including the ineffective right factor and the
nonlinear kernel of the last passage, are precisely those in
Theorem~\ref{thm:exact-coefficient-stabilizer-v148}.
For a nonempty zero/full rank pattern there is just one relation
orbit for that pattern, with the stabilizer in
\eqref{eq:full-stabilizer-v148}.
\end{corollary}
\begin{proof}
The recognition theorem gives surjectivity of the first map.
A linear isomorphism of two such relations is orientation-preserving;
its left factor carries every \(L_\lambda\) to the corresponding
coefficient space, and its right factor is unrestricted.  Scalars
act trivially on the Grassmannians.  This proves injectivity and
the first bijection.  The first-relation orbit proposition gives
the second.  For a zero/full pattern each coefficient Grassmannian
is a point, and the full-stabilizer theorem supplies the last claim.
\end{proof}

Equation~\eqref{eq:intrinsic-orbits-v148} is a statement about
geometric orbit sets with the explicitly computed affine
stabilizers.  It does not identify nonreduced moduli stacks or
assert that every infinitesimal deformation remains in this class.
For smooth projective reduced bases, its uniquely oriented case
combines with actual factor descent in
Theorem~\ref{thm:general-moving-coefficients-v147}.  The following
pencil-native application shows that this global step separates an
entire covering problem that fibrewise classification cannot see.

% BEGIN PRESERVED PART 29-curve-reconstruction-v145.tex
\section{Reconstruction on a projective curve}
\label{sec:curve-recognition-v145}

The global part of the inverse argument has a more economical geometric
realization.  It uses a proper base to make the projective left frame
constant, and a nontrivial right projective bundle to distinguish the
rulings.  Neither step requires the whole socle Grassmannian.  We now
replace that Grassmannian by a curve.  In particular, a neighbourhood
supported on one Schubert line already retains every pencil.

Let \(C\) be a smooth connected projective complex curve, and let
\(\mathcal U\) be a rank-\(n\) bundle on \(C\), where \(n\ge3\).
Assume that \(\Pj_C(\mathcal U^*)\) is not projectively trivial.
A sufficient numerical condition is
\begin{equation}\label{eq:curve-degree-obstruction-v145}
                 \deg\det\mathcal U\not\equiv0\pmod n.
\end{equation}
Indeed a projectively trivial rank-\(n\) bundle is a line-bundle twist
of a trivial bundle, so its determinant degree is divisible by \(n\).
Fix \(V\) of dimension \(n\), put \(E=\Hom(\mathcal U,V)\),
and let \(\mathfrak a\) be the zero-section ideal in
\(\Sym E^*\).  For a pencil \(R\subset\Sym^2V\), the quotient
\(\gamma_R:\Sym^2V\to S_R\) defines the homogeneous ideal
\begin{equation}\label{eq:curve-failure-ideal-v145}
 \mathcal I_{R,C}=(\det T)I_p(\gamma_R\Sym^2T),
 \qquad p=\binom{n+1}{2}-2,
\end{equation}
where \(T:\mathcal U\to V\) is the universal map on \(\Tot E\).
The expression is independent of local frames: both determinant
multiplication and the maximal-minor construction are bundle
constructions.  Define
\[
 Y_{R,C}^{[k]}=\Spec_C\bigl(\Sym E^*/
             (\mathcal I_{R,C}+\mathfrak a^{k+1})\bigr).
\]
Only its abstract complex scheme is retained, not its projection to
\(C\), its embedding in \(\Tot E\), or either tensor factor.

\begin{theorem}[Curve-supported reconstruction]
\label{thm:curve-reconstruction-v145}
For fixed \((C,\mathcal U,V)\) as above and
\(d=n^2+2n-4\), every pair of quadratic pencils satisfies
\begin{equation}\label{eq:curve-torelli-v145}
 Y_{R,C}^{[d]}\simeq Y_{R',C}^{[d]}
       \quad\Longleftrightarrow\quad R'=gR
              \quad\text{for some }g\in\PGL(V).
\end{equation}
All lower-order neighbourhoods are independent of \(R\).
The order-\(d\) schemes are finite locally free over \(C\) of rank
\(\binom{n^2+d}{d}-\binom{N}{2}\), where \(N=\binom{n+1}{2}\).
No regularity condition is imposed on the pencil.
\end{theorem}
\begin{proof}
On any frame chart, the ideal in
\eqref{eq:curve-failure-ideal-v145} is generated in degree \(d\).
Its degree-\(d\) piece \(K_R\) is a subbundle of rank
\(\binom N2\), by the fixed coefficient inclusion
\eqref{eq:finite-coefficient-factorization}; its quotient is locally
free.  The quotient algebra has its full symmetric powers in degrees
less than \(d\), has \(\Sym^dE^*/K_R\) in degree \(d\),
and is zero above that degree.  This proves the rank and the assertion
about lower orders.

Its reduction is \(C\), and its nilradical \(\mathfrak n\)
recovers \(L=\mathfrak n/\mathfrak n^2\) and
\[
       K_R=\ker[\Sym^dL\longrightarrow\mathfrak n^d].
\]
Thus the abstract finite-order scheme reconstructs the entire graded
algebra \(\Sym L/(K_R)\), by the same argument as
Lemma~\ref{lem:first-relation-cone}.  The ideal here means the
homogeneous ideal generated by the recovered kernel, not its radical.
No splitting of the reduction or supplied projection is used.

The reduction of this reconstructed cone is the determinant cone.
Indeed its ideal has the determinant factor, and when \(T\) is
invertible the quadratic quotient has full rank.  Its rank-one Fano
scheme has two ruling components,
\(\Pj(V)\times C\) and \(\Pj_C(\mathcal U^*)\).
The hypothesis distinguishes the first one intrinsically.  An abstract
isomorphism of neighbourhoods may induce a nonidentity automorphism
of \(C\); pull back the second bundle along that automorphism.
The oriented ruling map is nevertheless a morphism from \(C\) to
\(\PGL_n\).  It is constant, since \(C\) is connected projective
and \(H^0(C,\OO_C)=\C\).  Lemma~\ref{lem:oriented-segre-descent}
therefore supplies one constant projective left frame.

Cancel the intrinsic determinant ideal and apply the coefficient
construction of Theorem~\ref{thm:structural-coefficient-reconstruction}.
For pencils, Lemma~\ref{lem:pencil-irreducibility} leaves just the
single line \(\C\beta_R\).  Exterior duality
\eqref{eq:pencil-exterior-duality} recovers \(R\) with respect to
that same projective frame.  This proves the forward implication in
\eqref{eq:curve-torelli-v145}.  Conversely a linear lift of \(g\)
induces the bundle map on \(E\) and the quotient isomorphism on
\(S_R\), and carries all the displayed ideals and their truncations
to the corresponding ideals for \(R'\).  This proves the converse.
\end{proof}

\begin{corollary}[One Schubert line suffices]
\label{cor:schubert-line-reconstruction-v145}
Let \(\ell\simeq\Pj^1\) be any Schubert line in
\(B_R=\Gr(n,S_R)\), and use the canonical graph projection
\(Z_R^{[d]}\to B_R\) to form its restriction \(Z_{R,\ell}^{[d]}\).
Its abstract isomorphism class is independent of the choice of
\(\ell\), and
\[
 Z_{R,\ell}^{[d]}\simeq Z_{R',\ell'}^{[d]}
       \quad\Longleftrightarrow\quad R'=gR
                       \text{ for some }g\in\PGL(V).
\]
The smallest uniform reconstruction order for these curve-supported
neighbourhoods is again \(d\).  Their reductions all have dimension
one, instead of \(n(p-n)\).
\end{corollary}
\begin{proof}
The tautological bundle on a Schubert line is
\(\OO_{\Pj^1}(-1)\oplus\OO_{\Pj^1}^{n-1}\).
Pulling back the homogeneous ideal of the graph neighbourhood gives
exactly \eqref{eq:curve-failure-ideal-v145} with that bundle.  Its
isomorphism type depends on this bundle and \(R\), not on the
chosen inclusion of the line in the socle.  The determinant degree
is \(-1\), so Theorem~\ref{thm:curve-reconstruction-v145} applies.
Its lower-order schemes are all independent of the pencil, while
\(\langle x_1^2,x_1x_2\rangle\) and
\(\langle x_1^2,x_2^2\rangle\) are inequivalent.  This proves
the uniform minimality assertion.
\end{proof}

This restriction is made in the construction; it is not a claim that
an abstract neighbourhood canonically chooses a Schubert line.  No
line is supplied to an isomorphism in the corollary: the reduced curve
is recovered as the reduction of its source.  Nor is the corollary a
statement about one unmarked Artinian normal slice.  The projective
curve and the nontrivial ruling remain essential to this argument.
The same theorem applies, for example, to
\(\mathcal U=\OO_C(-P)\oplus\OO_C^{n-1}\) on every smooth
projective connected curve with a point \(P\).  It isolates the
properness and ruling obstruction used in reconstruction, independently
of the original Grassmannian realization.

% BEGIN PRESERVED PART 31-moving-pencils-v146.tex
\section{Reconstruction of moving pencils and their isomorphisms}
\label{sec:moving-pencils-v146}

The inverse need not be applied separately to constant pencils.  A single
unmarked finite scheme can retain a pencil varying over its entire
projective reduction.  We prove this assertion, including the recovery
of the normal bundle and the ambiguity in lifting the recovered data to
an isomorphism.  The coefficient lines in this section are subbundles;
constancy of the pencil is not assumed.

\subsection{The family and its first relation bundle}
Fix \(n\ge3\), and write
\[
 N=\binom{n+1}{2},\qquad p=N-2,\qquad
 d=n+2p=n^2+2n-4,\qquad M=\binom N2.
\]
An admissible datum consists of a smooth connected projective complex
variety \(B\), a rank-\(n\) vector bundle \(\mathcal U\) on \(B\), a
constant \(n\)-dimensional vector space \(V\), and a rank-two locally
direct-summand subbundle
\[
 \mathcal R\ \subset\ \Sym^2V\otimes\OO_B.
\]
We require that \(\Pj_B(\mathcal U^*)\) is not a trivial projective
bundle.  Put \(\mathcal S=(\Sym^2V\otimes\OO_B)/\mathcal R\) and
\(\mathcal E=\Hom(\mathcal U,V\otimes\OO_B)\).  On its total space,
let \(T\) be the universal map and let \(\gamma\) denote the quotient
onto \(\mathcal S\).  Define the homogeneous ideal and its truncations by
\begin{equation}\label{eq:moving-construction-v146}
 \begin{split}
 \mathcal I_{\mathcal R}&=(\det T)I_p(\gamma\Sym^2T),\\
 \mathcal A_{\mathcal R}^{[k]}
   &=\Sym(\mathcal E^*)/
       (\mathcal I_{\mathcal R}+\mathfrak a^{k+1}),\qquad
 Y_{\mathcal R}^{[k]}=\Spec_B\mathcal A_{\mathcal R}^{[k]},
 \end{split}
\end{equation}
where \(\mathfrak a\) is the positive-degree ideal.  Determinant-line
changes of frame multiply local generators by units, so this is an
ideal sheaf, independent of all frames.  This construction is the
pullback of the coefficient construction in
Theorem~\ref{thm:universal-finite-neighbourhood}, with an arbitrary
rank-\(n\) source bundle.  It is not necessary to choose an embedding
of \(B\) in a socle Grassmannian.

\begin{lemma}[A varying coefficient line]\label{lem:moving-coefficients-v146}
The ideal \(\mathcal I_{\mathcal R}\) has no component of degree less
than \(d\) and is generated by its degree-\(d\) component
\(\mathcal K_{\mathcal R}\).  This component is a subbundle of
\(\Sym^d\mathcal E^*\) of rank \(M\).  After cancellation of the
determinant line, its degree-\(2p\) coefficient module determines
\(\mathcal R\) as a subbundle, once the two tensor factors of
\(\mathcal E\) are supplied.  In particular,
\(Y_{\mathcal R}^{[d]}\to B\) is finite locally free of rank
\[
                   \binom{n^2+d}{d}-M.
\]
\end{lemma}
\begin{proof}
Put \(F(H)=\bigwedge^p\Sym^2H\) for rank-\(n\) vector spaces or
bundles \(H\).  This is an irreducible polynomial representation by
Lemma~\ref{lem:pencil-irreducibility}; its dimension is \(M\).  Write
\(\mathcal L_{\mathcal R}=\bigwedge^p\mathcal S^*\), viewed, by
\(\bigwedge^p\gamma^*\), as a line subbundle of
\(F(V)^*\otimes\OO_B\).  The residual minors are precisely the image of
\begin{equation}\label{eq:moving-coefficient-map-v146}
 \mathcal L_{\mathcal R}\otimes F(\mathcal U)
  \longrightarrow \Sym^{2p}(V^*\otimes\mathcal U),\qquad
 \beta\otimes\xi\longmapsto[T\mapsto\beta(F(T)\xi)].
\end{equation}
The full map \(F(V)^*\otimes F(U)\to
\Sym^{2p}(V^*\otimes U)\) is the nonzero inclusion of its unique
Cauchy summand: evaluation at \(T=1\) is nonzero, as in
\eqref{eq:general-coefficient-map}.  Restricting to a nonzero line on
the left therefore has rank \(M\) in every fibre.  The coefficients
vary algebraically.  The constant-rank criterion makes the image a
subbundle; equivalently one may choose a nonvanishing maximal minor
of this coefficient map locally on the Grassmannian of pencils and
pull it back to \(B\).  Multiplication by the determinant, a nonzero
polynomial in every fibre, again has constant rank \(M\).

All defining generators have degree \(d\), proving the assertion
about the ideal.  For the recovery assertion, the residual image,
inside the indicated Cauchy summand, is
\(\mathcal L_{\mathcal R}\otimes F(\mathcal U)\).  Contraction with
\(F(\mathcal U)^*\) has image exactly \(\mathcal L_{\mathcal R}\)
as a line subbundle of \(F(V)^*\otimes\OO_B\).  This statement is an
equality of images of bundle maps, not just equality at closed points.
Exterior duality identifies it with the Pluecker line
\(\bigwedge^2\mathcal R\), up to the fixed determinant line of
\(\Sym^2V\).  The Grassmannian Pluecker embedding consequently
recovers the morphism \(B\to\Gr(2,\Sym^2V)\), and its pullback
of the tautological bundle recovers \(\mathcal R\).
Finally, the truncated algebra is the sum of its free graded pieces
below \(d\) and the locally free degree-\(d\) quotient by
\(\mathcal K_{\mathcal R}\).  Its rank is the asserted binomial sum
minus \(M\).
\end{proof}

\subsection{The unmarked inverse}
\begin{theorem}[Reconstruction of a moving pencil]
\label{thm:moving-reconstruction-v146}
Let \((B_i,\mathcal U_i,V_i,\mathcal R_i)\), \(i=1,2\), be admissible
data of the same rank \(n\).  Then
\begin{equation}\label{eq:moving-inverse-v146}
 Y_{\mathcal R_1}^{[d]}\simeq Y_{\mathcal R_2}^{[d]}
\end{equation}
as abstract complex schemes if and only if there exist an isomorphism
\(\sigma:B_1\to B_2\), a constant linear isomorphism
\(g:V_1\to V_2\), and a bundle isomorphism
\(h:\mathcal U_1\to\sigma^*\mathcal U_2\) such that
\begin{equation}\label{eq:moving-data-v146}
                  \sigma^*\mathcal R_2=(\Sym^2g)\mathcal R_1.
\end{equation}
The reduced base, normal tensor orientation, isomorphism class of the
actual source bundle, and the varying pencil are recovered together.
Neither the projection onto the base nor the ambient bundle is
supplied in \eqref{eq:moving-inverse-v146}.  The constant left map
cannot be replaced by a separately chosen change of coordinates in
each fibre.  A pair \((g,h)\) defining the same normal-bundle map is
unique up to simultaneous multiplication by a nonzero constant.
\end{theorem}
\begin{proof}
Let \(\mathfrak n_i\) be the nilradical of
\(\OO_{Y_{\mathcal R_i}^{[d]}}\).  Since the positive-degree ideal
is nilpotent and \(B_i\) is reduced, the reduction is exactly \(B_i\),
and this ideal is exactly \(\mathfrak n_i\).  There are no relations
below degree \(d\), so multiplication intrinsically recovers
\begin{equation}\label{eq:moving-intrinsic-kernel-v146}
 \mathcal E_i^*=\mathfrak n_i/\mathfrak n_i^2,
 \qquad
 \mathcal K_i=\ker\left[
 \Sym^d(\mathfrak n_i/\mathfrak n_i^2)
       \longrightarrow\mathfrak n_i^d\right].
\end{equation}
The map is well defined since \(\mathfrak n_i^{d+1}=0\).
An abstract isomorphism gives \(\sigma\), an isomorphism of these
normal bundles, and the matching relation subbundles.  In particular
it gives an isomorphism of the untruncated homogeneous cones
\[
               \Spec_{B_i}\bigl(\Sym\mathcal E_i^*/(\mathcal K_i)\bigr).
\]
This cone is reconstructed from the first relation module; it is not
the associated graded algebra of the finite scheme, whose degrees
above \(d\) vanish.

The reduction of this cone is the determinant cone.  Indeed the
residual quadratic map is surjective wherever \(T\) is invertible,
whereas the determinant factor annihilates the ideal wherever
\(T\) is singular.  On a frame chart the determinant hypersurface
is reduced over the smooth base, so the radical ideal is precisely
the determinant ideal.  Successive relative reduced singular loci
give the lower matrix-rank cones.  The rank-one cone has the two
ruling parameter schemes
\[
             \Pj(V_i)\times B_i,\qquad \Pj_{B_i}(\mathcal U_i^*),
\]
by Lemma~\ref{lem:rank-one-fano-scheme}.  Only the first is
projectively trivial.  Thus an isomorphism cannot interchange them.
The induced map on the trivial ruling is a morphism
\(B_1\to\operatorname{Isom}(\Pj(V_1),\Pj(V_2))\).
After choosing one identification this target is the affine variety
\(\PGL_n\).  Since \(B_1\) is connected and projective, every such
morphism is constant.  Choose one constant linear lift \(g\).

We check that no line-bundle ambiguity remains in the other factor.
Locally trivialize both source bundles.  The linear rank-one
preserver with its left projective map fixed to \([g]\) is of the
form
\begin{equation}\label{eq:moving-normal-map-v146}
                           T\longmapsto gTh^{-1}
\end{equation}
by Lemma~\ref{lem:oriented-segre-descent}.  Once this particular
\(g\) is fixed, the local linear map \(h\) is unique: if two right
maps give the same value for every \(T\), they are equal.  The local
maps therefore agree on overlaps and glue to an actual bundle
isomorphism \(\mathcal U_1\to\sigma^*\mathcal U_2\).  Alternatively,
after removing \(g\) the normal-bundle map commutes with
\(\End(V_1)\); the commutant is the right-factor homomorphism bundle.
Rescaling the constant lift \(g\) rescales \(h\) by the same constant.
This proves both existence and the asserted uniqueness for a given
normal-bundle map.

Cancel the intrinsic determinant ideal in the recovered homogeneous
ideal.  On a frame chart the determinant polynomial is a
non-zero-divisor and
\(( (\det T)\mathcal J:(\det T))=\mathcal J\); the identities glue
without selecting a global determinant generator.  Apply
Lemma~\ref{lem:moving-coefficients-v146} to the degree-\(2p\) part.
The right map \(h\) acts on the full factor \(F(\mathcal U)\), so
contraction against its dual removes that change of frame.  The
remaining line subbundle transforms by the one constant left map
\(g\).  Exterior duality yields exactly
\eqref{eq:moving-data-v146}.  This uses the subbundle version proved
above, not the constant-coefficient hypothesis in
Theorem~\ref{thm:structural-coefficient-reconstruction}.

Conversely \(\sigma,g,h\) satisfying \eqref{eq:moving-data-v146}
identify the quotient bundles \(\mathcal S_i\), the quadratic maps,
and their maximal-minor ideals.  The determinant changes by a unit
in local frames.  The map \eqref{eq:moving-normal-map-v146} consequently
identifies both the homogeneous ideals and every truncation,
proving \eqref{eq:moving-inverse-v146}.
\end{proof}

\begin{corollary}[Uniform sharpness over an admissible base]
\label{cor:moving-sharpness-v146}
For fixed \(B,\mathcal U,V\), all schemes
\(Y_{\mathcal R}^{[k]}\) with \(k<d\) are independent of
\(\mathcal R\).  Order \(d\) is the smallest uniform order at which
the moving pencil is recoverable in the sense of
Theorem~\ref{thm:moving-reconstruction-v146}.
\end{corollary}
\begin{proof}
Below \(d\) the ideal \(\mathcal I_{\mathcal R}\) contributes no
relations, and the algebra is
\(\bigoplus_{j=0}^k\Sym^j\mathcal E^*\).  Constant pencils are
among the admissible subbundles.  The pencils
\(\langle x_1^2,x_1x_2\rangle\) and
\(\langle x_1^2,x_2^2\rangle\) are inequivalent: all members of
the first have a common linear factor and those of the second do
not.  A base automorphism does not change either constant pencil.
The theorem separates their order-\(d\) schemes, whereas every
lower truncation agrees.  This proves uniform, not pairwise, minimality.
\end{proof}

\subsection{The part of an isomorphism that the inverse does not record}
The recovered data classify isomorphism classes, but an isomorphism
of finite schemes need not itself be linear in the nilpotent
directions.  The following statement identifies this distinction
without replacing it by a claim about moduli stacks.

Fix an admissible datum, put \(Y=Y_{\mathcal R}^{[d]}\), and let
\(\mathscr L_{\mathcal R}\) be the group of triples
\((\sigma,g,h)\) satisfying \eqref{eq:moving-data-v146}, modulo
\((g,h)\sim(cg,ch)\) for \(c\in\C^*\).  Its group law is
composition of their induced normal-bundle maps.  Write
\(\operatorname{Aut}_{\C}(Y)\) for the group of complex scheme
automorphisms.  For \(j\ge1\), define \(G_j\) to consist of the
automorphisms inducing the identity on the reduction and satisfying
\begin{equation}\label{eq:automorphism-filtration-v146}
 (\alpha^*-1)(\mathfrak n^k)\subseteq\mathfrak n^{k+j}
 \quad\text{for every }k\ge0,\qquad \mathfrak n^0=\OO_Y.
\end{equation}
They act on sheaves on the same underlying topological space.

\begin{theorem}[Linear data and the nilpotent kernel]
\label{thm:automorphism-kernel-v146}
There is an exact sequence of groups
\begin{equation}\label{eq:automorphism-sequence-v146}
 1\longrightarrow G_1\longrightarrow
 \operatorname{Aut}_{\C}(Y)\longrightarrow
 \mathscr L_{\mathcal R}\longrightarrow1.
\end{equation}
The displayed homogeneous presentation of \(Y\) supplies a splitting.
Moreover \(G_{d+1}=1\),
\([G_i,G_j]\subseteq G_{i+j}\), and there are natural injections
\begin{equation}\label{eq:automorphism-symbol-v146}
 G_j/G_{j+1}\ \lhook\joinrel\longrightarrow\
 H^0\!\left(B,\operatorname{Der}_{\C}
       (\gr_{\mathfrak n}\OO_Y)_j\right)
\end{equation}
into additive groups of degree-\(j\) derivations.  Thus the kernel is
nilpotent and records precisely the automorphisms invisible on the
associated graded scheme.  No surjectivity onto the derivation spaces
is asserted.
\end{theorem}
\begin{proof}
An automorphism induces a base map and a degree-one bundle map.  The
proof of Theorem~\ref{thm:moving-reconstruction-v146} identifies these
with a unique class of triples in \(\mathscr L_{\mathcal R}\).
The graded algebra is generated in degree one over degree zero, with
relation subbundle \(\mathcal K_{\mathcal R}\).  Hence its entire
graded automorphism is already determined by that triple.  The kernel
is the identity on every graded piece, which is exactly condition
\eqref{eq:automorphism-filtration-v146} for \(j=1\).
Conversely the homogeneous maps \(T\mapsto gTh^{-1}\) give actual
automorphisms of \(Y\), respect composition, and induce the prescribed
triples.  This proves exactness and the splitting.  A choice of the
homogeneous presentation supplies that section; the assertion does
not endow an arbitrary unmarked presentation with a preferred section.

For \(\alpha\in G_j\), put \(D_\alpha=\alpha^*-1\).  It raises
filtration by \(j\), and is therefore nilpotent as an additive
operator.  The inverse of \(1+D_\alpha\) is the finite geometric
series.  Composition of filtration-raising operators shows that a
commutator of elements of \(G_i,G_j\) raises filtration by at least
\(i+j\).  Also \(G_{d+1}=1\) follows by taking \(k=0\), since
\(\mathfrak n^{d+1}=0\).

The leading term of \(D_\alpha\) defines maps
\(\mathfrak n^k/\mathfrak n^{k+1}\to
\mathfrak n^{k+j}/\mathfrak n^{k+j+1}\).
They are independent of lifts.  The identity
\[
 D_\alpha(ab)=aD_\alpha(b)+D_\alpha(a)b+
                           D_\alpha(a)D_\alpha(b)
\]
reduces to the Leibniz identity in the indicated graded degree, since
\(2j\ge j+1\).  Thus the leading term is a sheaf derivation over
\(\C\), including its action on degree zero; it need not be an
\(\OO_B\)-linear derivation.  Leading terms add under composition,
and their kernel is exactly \(G_{j+1}\).  They glue globally,
proving \eqref{eq:automorphism-symbol-v146} and all assertions.
\end{proof}

\begin{example}[A nontrivial invisible automorphism]
\label{ex:invisible-shear-v146}
Let \(B=\Pj^1\) and
\(\mathcal U=\OO(-1)\oplus\OO^{n-1}\).
The bundle \(\mathcal E^*=V^*\otimes\mathcal U\) has at least two
trivial direct summands.  Choose their degree-one generators \(x,y\).
The substitution
\[
                          x\longmapsto x+y^2
\]
fixing \(y\), the other summands and \(\OO_B\), defines a nonidentity
automorphism of \(Y_{\mathcal R}^{[d]}\) in \(G_1\), for every
admissible \(\mathcal R\).  Indeed each relation has degree \(d\);
its change under this substitution has degree at least \(d+1\) and
vanishes in the truncation.  The inverse is \(x\mapsto x-y^2\).
The automorphism is nonidentity because there are no degree-two
relations.  Therefore identifying all finite-scheme automorphisms
with the stabilizer of the pencil would be false, even in this simple
case.  The inverse theorem and the exact sequence retain, rather than
suppress, this higher-order freedom.
\end{example}

\subsection{A nonconstant family across spectral collisions}
\begin{example}\label{ex:nonconstant-pencil-v146}
Take \(n=4\), \(B=\Pj^1_{[s:t]}\),
\(\mathcal U=\OO(-1)\oplus\OO^3\), and \(V=\C^4\).  In diagonal
quadratic coordinates set
\[
 A=\operatorname{diag}(1,1,1,1),\qquad
 B_{s,t}=\operatorname{diag}(0,s,t,2s+3t).
\]
The map \(\OO\oplus\OO(-1)\to\Sym^2V\otimes\OO\) with columns
\(A,B_{s,t}\) has rank two everywhere.  The first diagonal entry
shows that any scalar relation involving \(A\) has zero coefficient
on \(A\); the next two entries then force \(s=t=0\), which is
impossible on \(\Pj^1\).  Its image \(\mathcal R\) is therefore a
locally direct-summand pencil bundle.  Since it contains the
nondegenerate form \(A\), it is regular at every point, including
all spectral collisions.

On the open set of four distinct eigenvalues, an ordered cross-ratio
of \(0,s,t,2s+3t\) is
\[
 \chi=\frac{t(s+3t)}{(2s+3t)(t-s)}.
\]
The rational function
\(J=(\chi^2-\chi+1)^3/(\chi^2(1-\chi)^2)\)
is invariant under the six cross-ratio transformations arising from
permutations of four points.  It is nonconstant: at \([s:t]=[1:2]\)
and \([1:3]\), the values of \(\chi\) are \(7/4\) and \(15/11\),
and the corresponding values of \(J\) differ.  Thus the pencil
family is not projectively isotrivial, even allowing a change of
parameter within each pencil.  The discriminant of
\(\det(zA-B_{s,t})\) is
\begin{equation}\label{eq:moving-discriminant-v146}
 4s^2t^2(2s+3t)^2(s-t)^2(s+3t)^2(s+t)^2.
\end{equation}
No factor is inverted in forming \(Y_{\mathcal R}^{[20]}\).
This finite scheme is locally free over its entire projective
reduction of rank \(\binom{36}{20}-45\).
By Theorem~\ref{thm:moving-reconstruction-v146}, its unmarked
isomorphism class recovers the whole morphism into the pencil
Grassmannian, including the six collision points and the pencil
there, up to one constant transformation of \(V\) and an automorphism
of the base.  This is a consequence about a genuinely moving family,
not just the pointwise Pluecker encoding.  The cross-ratio and
discriminant calculations are classical elementary readouts; their
role here is to exhibit a family to which the unmarked global inverse
applies.
\end{example}

\begin{remark}[Scope of the family statement]
\label{rem:moving-scope-v146}
The reduction in this section is smooth, projective and reduced.  The
source projective bundle is assumed nontrivial, which orients the two
rulings.  The theorem allows all fibre pencils, singular or regular;
regularity is used only in the spectral description of the last
example.  Theorem~\ref{thm:universal-finite-neighbourhood} separately
allows arbitrary complex base change with its specified relative
socle ideal.  The present intrinsic inverse does not replace that
ideal by an absolute nilradical over a nonreduced base.  Nor does
\eqref{eq:automorphism-sequence-v146} assert an equivalence of
unmarked moduli stacks or a description of every infinitesimal
deformation.  It is an exact classification of geometric
isomorphisms, their graded data, and the residual automorphism kernel
for the stated admissible objects.
\end{remark}

% BEGIN PRESERVED PART 38-pencil-deformation-equivalence-v149.tex
\section{The effective moduli stack of pencil failure algebras}
\label{sec:pencil-deformations}

We apply the common-divisor inverse to the actual first relations of
quadratic pencils.  Their common divisor has a higher multiplicity
than the single determinant factor used to recover their support.
Removing that multiplicity places every pencil, without a regularity
hypothesis, in the smooth stratum of the preceding section.

Keep \(W=\Sym^2V\), \(N=\binom{n+1}{2}\), \(p=N-2\), and put
\begin{equation}\label{eq:primitive-pencils}
 \begin{split}
 M&=\binom N2,\qquad s=n-1,\qquad h=3n-4,\\
 \mu&=(\underbrace{3,\ldots,3}_{n-2},2,0),\qquad
 A=\mathbb S_\mu V^*,\qquad F=\mathbb S_\mu U.
 \end{split}
\end{equation}
Both \(A\) and \(F\) have dimension \(M\).  The Cauchy inclusion
\(A\otimes F\subset S_h(E)\) will be fixed; its subspace and all
its projections are equivariant, independent of a choice of bases.

\begin{lemma}[Exact common multiplicity]\label{lem:primitive-pencil}
For every \(R\in\Gr(2,W)\), including every singular pencil, the
first relation space of its order-\(d\) failure algebra is
\begin{equation}\label{eq:primitive-factorization}
 K_R=\delta^{\,n-1}\widetilde J_R,\qquad
 \widetilde J_R=L_R^0\otimes F\subset S_h(E),
                   \qquad \gcd(\widetilde J_R)=1.
\end{equation}
Here \(L_R^0\subset A\) is the Pluecker line of \(R\), under the
canonical projective identification
\(\Pj(A)=\Pj(\bigwedge^2W)\).
The common divisor of \(K_R\) is therefore exactly
\(\delta^{n-1}\), up to scalar.  Its degree remains
\(d=sn+h=n^2+2n-4\).
\end{lemma}
\begin{proof}
Exterior duality and
\(\bigwedge^2\Sym^2H=\mathbb S_{(3,1)}H\) give
\[
 \bigwedge^{N-2}\Sym^2H
   =\mathbb S_{\lambda}H,\qquad
 \lambda=(\underbrace{n+1,\ldots,n+1}_{n-2},n,n-2)
   =\mu+(n-2)^n.
\]
Thus this representation is
\((\det H)^{n-2}\otimes\mathbb S_\mu H\).
Applied to the universal matrix, its coefficient polynomials factor
as \(\delta^{n-2}\) times the matrix coefficients of
\(\mathbb S_\mu T\).  The additional determinant in the failure
ideal gives \eqref{eq:primitive-factorization}.  On the left,
exterior duality also yields
\[
                A\simeq\bigwedge^2W\otimes(\det V)^{-3},
\]
which identifies the residual line with the Pluecker line.

We prove a slightly stronger assertion: the matrix coefficient
space \(\beta\otimes\mathbb S_\mu U\) has greatest common divisor
\(1\) for every nonzero \(\beta\in\mathbb S_\mu V^*\).
Its common divisor is taken to a scalar multiple of itself by the
full right group.  Its zero divisor is therefore invariant under
right multiplication.  A nonconstant invariant divisor cannot meet
the open orbit of invertible matrices.  Since the complement of
that orbit is the irreducible determinant hypersurface, any such
common divisor is a power of \(\delta\).

It remains to exclude \(\delta\).  Choose a rank-\(n-1\) map
\(T_0\).  The representation \(\mathbb S_\mu\) has length \(n-1\),
so \(\mathbb S_\mu T_0\ne0\).  Choose a nonzero vector in its image.
The span of its left \(\operatorname{GL}(V)\)-orbit is the entire
irreducible representation.  Consequently some left translate
\(gT_0\) has a coefficient not annihilated by \(\beta\).
Not all residual coefficients vanish on the rank-\(n-1\) locus.
They have no determinant factor, proving the assertion.  The same
argument at an invertible matrix shows that the coefficient vector
is nonzero everywhere on the invertible locus.
\end{proof}

Let \(\mathscr P=\Gr(2,W)\).  The morphism
\begin{equation}\label{eq:primitive-pencil-embedding}
 \iota:\mathscr P\longrightarrow X_{h,M}(E),
                 \qquad R\longmapsto L_R^0\otimes F
\end{equation}
is a closed immersion.  Indeed, the Pluecker embedding is followed
by the closed immersion \(\Pj(A)\to\Gr(M,A\otimes F)\),
\(L\mapsto L\otimes F\), and then by a linear Grassmannian
inclusion.  The inverse to the middle map is contraction with
\(F^*\).  This factorization proves the assertion on schemes,
including nonreduced test schemes, rather than only on points.
Lemma~\ref{lem:primitive-pencil} shows that its image lies in the
open set \(X_{h,M}(E)\).

Define the \emph{pencil relation stratum} \(Z_{\rm pen}\) inside
\(Z^{\det}_{n-1,3n-4,M}\) by the following intrinsic scheme
conditions.  Recover the common-divisor line, its determinant root,
and its unordered tensor rulings by
Theorem~\ref{thm:determinant-stratum}.  In an ordering, the primitive
residual relation is required to lie in the indicated Cauchy
summand, to have the form \(L\otimes F\), and to satisfy the
Pluecker equations on \(L\) under
\(A\simeq\bigwedge^2W\otimes(\det V)^{-3}\).
These are equations on subbundles; a test on the fibres alone is
not their definition.  Use both orderings and descend.  The two
ordered images are disjoint: their support ranks are respectively
\((1,M)\) and \((M,1)\), with \(M>1\).  Each is closed in the
residual Grassmannian.  Thus this prescription defines a locally
closed scheme, independent of the tensor frames.  It is an
intrinsic image criterion for a natural family of first relations,
not a definition by an algebra isomorphism to a preselected pencil.

\begin{theorem}[Effective all-pencil moduli and deformations]
\label{thm:pencil-stack-equivalence}
The stack of first relations in \(Z_{\rm pen}\) is
\begin{equation}\label{eq:pencil-stack-before-rigidification}
 [Z_{\rm pen}/\operatorname{GL}(E)]\simeq
                         [\mathscr P/G^\circ].
\end{equation}
There is an exact sequence of groups
\begin{equation}\label{eq:pencil-ineffective-group}
 1\longrightarrow\operatorname{GL}(U)\longrightarrow G^\circ
                 \longrightarrow\PGL(V)\longrightarrow1,
\end{equation}
and its kernel acts trivially on \(\mathscr P\).
Quotient the morphisms in the algebra stack first by the
nonlinear tangent-identity kernel of
Proposition~\ref{prop:relative-first-relation}, and then by this
ineffective right group, and fppf stackify.  The resulting
\emph{effective pencil failure stack} satisfies
\begin{equation}\label{eq:effective-pencil-stack}
       \mathscr F_{\rm pen}^{\rm eff}\simeq
                [\Gr(2,\Sym^2V)/\PGL(V)].
\end{equation}
This equivalence is valid over all complex parameter schemes.
It identifies automorphisms after the specified quotients,
infinitesimal deformation groupoids, and families of degenerating
pencils, including singular pencils.

The stack is smooth of dimension \(n-3\).  At \(R\), its tangent
complex is
\begin{equation}\label{eq:pencil-tangent-complex}
 [\,\mathfrak{pgl}(V)\xrightarrow{\ d_R\ }\Hom(R,W/R)\,],
 \qquad d_R(X)(r)=(X\otimes1+1\otimes X)r\bmod R.
\end{equation}
In particular, the dimension of its space of first-order
isomorphism classes is \(n-3+\dim\operatorname{Stab}_{\PGL(V)}(R)\).
Within this stratum, every infinitesimal pencil deformation and
its effective failure-algebra deformation lift together across a
small extension of local Artinian complex algebras.
\end{theorem}
\begin{proof}
In the determinant frame of
Theorem~\ref{thm:determinant-stratum}, the residual pencil locus is
\(\iota(\mathscr P)\amalg\tau\iota(\mathscr P)\), where
\(\tau\) exchanges the two factors.  The disjointness just proved
and the equivariance of the Cauchy and Pluecker maps give, as
schemes with group action,
\[
 \iota(\mathscr P)\amalg\tau\iota(\mathscr P)
                         \simeq G\times^{G^\circ}\mathscr P.
\]
Quotienting by \(G\) gives
\eqref{eq:pencil-stack-before-rigidification}.  All maps used here
are morphisms and all inverse operations are contractions,
common-divisor descent, or the inverse on the Pluecker image.
Consequently this calculation remains valid on a nonreduced base.
It is not an inference from equality of complex-point stabilizers.

The map in \eqref{eq:pencil-ineffective-group} sends the class of
\((g,h)\) to \([g]\).  If \(g=aI\), multiplication by the diagonal
scalar \(a^{-1}\) writes its class uniquely as \((I,a^{-1}h)\).
This proves the asserted kernel.  The action on the Pluecker line
is the projective action of \(g\); the full right factor has no
effect on \(R\).  Every \(\PGL(V)\)-valued isomorphism lifts fppf
locally to \(G^\circ\), and two lifts differ by that kernel.
The quotient of morphisms and stackification therefore gives
\([\mathscr P/\PGL(V)]\).  Combining this with
Proposition~\ref{prop:relative-first-relation} proves
\eqref{eq:effective-pencil-stack}.  No splitting of
\eqref{eq:pencil-ineffective-group} over the base is required.

A Grassmannian chart writes a nearby pencil as the graph of a map
\(R\to W/R\).  Infinitesimal coordinate changes act by \(d_R\),
giving \eqref{eq:pencil-tangent-complex}.  Grassmannian charts and
\(\PGL(V)\) are smooth, so the quotient is smooth and the graph
coefficients lift over every small extension.  A coordinate change
can also be lifted in the smooth group, and its action determines
a compatible lift of the corresponding object.  Torsors over a local
Artinian complex algebra can be trivialized, since smooth torsors
lift a point of their algebraically closed residue field.
This proves the lifting assertion for deformation groupoids.
Finally,
\[
 \dim\mathscr P-\dim\PGL(V)=2(N-2)-(n^2-1)=n-3.
\]
Taking kernel and cokernel of \(d_R\) gives the stated first-order
dimension, including stabilizer jumps at singular pencils.
\end{proof}

\begin{corollary}[Native specialization loci]\label{cor:pencil-specialization-stack}
Every \(\PGL(V)\)-invariant locally closed subscheme
\(Q\subset\Gr(2,\Sym^2V)\), with its specified scheme structure,
has a corresponding intrinsic effective failure stratum
isomorphic to \([Q/\PGL(V)]\).  Its infinitesimal equations and its
specializations are retained, not just its closed points.
In particular, this applies to rank-incidence and spectral Fitting
conditions on the universal pencil, including their intersections.
A family in such a locus specializes to a pencil in its closure
exactly when its effective failure family has the corresponding
specialization in the associated closure stratum.
\end{corollary}
\begin{proof}
Restrict the equivariant closed immersion
\eqref{eq:primitive-pencil-embedding} and its inverse to \(Q\).
The construction of \(Z_{\rm pen}\) and its two quotient steps
then restrict to give the stated isomorphism.  Equations for a
rank or a Fitting condition are equivariant equations in the
universal pencil, so they are included in this argument with their
scheme structures.  Apply the same construction to the closure of
\(Q\) to obtain the specialization assertion.
\end{proof}

\begin{remark}\label{rem:moduli-scope}
Theorem~\ref{thm:pencil-stack-equivalence} does not say that the
unrigidified algebra stack is the pencil stack.  At a complex
pencil its linear stabilizer remains
\((G_R\times\operatorname{GL}(U))/\mathbb G_m\), and its full
algebra automorphism group also has the nonlinear kernel in
Proposition~\ref{prop:first-relation-group-v147}.  Both are real
automorphisms, not redundant notation.  They are explicitly removed
only in the effective comparison.  Nor is the quotient asserted to
be proper, although its framed Grassmannian is projective.

The parameter base in this theorem is distinct from the intrinsic
projective reduction in Theorem~\ref{thm:unrestricted-moving-v146}.
Here projective coordinate changes may vary over the parameter
scheme.  In that moving theorem, properness of the intrinsic
reduction forces one constant left map and recovers an actual
source bundle.  Neither assertion is substituted for the other.
The present equivalence upgrades the geometric orbit statement to
a family statement without changing the hypotheses or conclusions
of the moving inverse.
\end{remark}

% BEGIN PRESERVED PART 41-rigidification-v151.tex
\section{Universal properties of the two rigidifications}
\label{sec:rigidification-v151}

The effective stacks used above have a universal characterization.
This specifies what is forgotten from isomorphisms, as opposed to what
is retained from objects.  We use rigidification in its standard sense
of removing a flat normal subgroup of inertia; see
\cite[Tag~04V2]{StacksRigidV151}.  The argument below gives the
universal property directly in the present situation, including the
nonconstant subgroup obtained by descent.

Let \(H_0\subset\Gr(r,S_D(E))\) be a
\(\operatorname{GL}(E)\)-invariant locally closed subscheme.  Let
\(\mathscr A_{H_0}\) be the stack of algebras satisfying the trace,
filtration and Hilbert-function conditions of
Proposition~\ref{prop:relative-first-relation}, whose intrinsic
kernel is in \(H_0\).  Write
\(\mathscr R_{H_0}=[H_0/\operatorname{GL}(E)]\).

\begin{theorem}[First-relation rigidification]
\label{thm:rigidification-v151}
Taking the first relation defines an fppf gerbe
\[
                      q:\mathscr A_{H_0}\longrightarrow\mathscr R_{H_0}.
\]
Its relative inertia \(\mathscr U\) consists exactly of the
substitutions which induce the identity on \(\mathcal N/\mathcal N^2\).
It is a smooth affine normal subgroup of the inertia.  In a homogeneous
presentation its underlying scheme is
\(\Hom(\mathcal E,\mathcal N^2)\), with the substitution group law.
For every stack \(\mathscr Z\) in groupoids on complex schemes,
precomposition with \(q\) identifies
\begin{equation}\label{eq:rigidification-universal-v151}
 \Hom(\mathscr R_{H_0},\mathscr Z)
 \simeq
 \{F\in\Hom(\mathscr A_{H_0},\mathscr Z):F(\mathscr U)=1\}.
\end{equation}
The right side is the full subgroupoid, including all natural
isomorphisms, of functors killing the indicated relative inertia.
The construction is compatible with base change on the relation stack.
\end{theorem}
\begin{proof}
A relation bundle has its homogeneous truncated algebra as a lift.
Every algebra over that relation bundle is locally isomorphic to this
lift.  The first-relation proposition shows that an isomorphism of
relation bundles lifts fppf locally, and any two lifts differ by a
unique element of the indicated kernel.  These assertions are exactly
the gerbe and relative-inertia assertions.  The kernel is intrinsic,
so it is stable under conjugation.  Its local affine-space description
shows smoothness and flatness; no additive group law is being asserted.

For clarity, we prove the descent implicit in
\eqref{eq:rigidification-universal-v151}.  Given \(F\) on the right
and an object \(y\) of \(\mathscr R_{H_0}\), choose local lifts
\(a_i\) over an fppf cover.  Locally choose comparison isomorphisms
on overlaps above the identity of \(y\).  Two choices differ by
\(\mathscr U\), hence have the same image under \(F\).  On triple
overlaps the failure of the chosen lifts to satisfy the cocycle
condition is again in \(\mathscr U\).  Their images therefore give
a genuine descent datum for \(F(a_i)\), which is effective because
\(\mathscr Z\) is a stack.  Refinements and different lift choices
give canonically isomorphic descended objects.  An arrow between two
objects of \(\mathscr R_{H_0}\) is treated in the same way, using
local lifts of that arrow.  This constructs the descended functor.
Natural isomorphisms descend by the same argument, proving full
faithfulness as well as essential surjectivity.  Conversely every
functor pulled back by \(q\) kills its relative inertia.  All the
constructions use covers and isomorphism sheaves, so commute with
base change.
\end{proof}

\begin{corollary}[The effective pencil stack is a successive rigidification]
\label{cor:two-rigidifications-v151}
On the pencil relation stratum the second morphism
\[
 [\mathscr P/G^\circ]\longrightarrow[\mathscr P/\PGL(V)],
                     \qquad \mathscr P=\Gr(2,\Sym^2V),
\]
is the rigidification by the normal relative inertia induced from
\(\ker(G^\circ\to\PGL(V))=\operatorname{GL}(U)\).
Consequently Theorem~\ref{thm:pencil-stack-equivalence} has the
universal property of first killing \(\mathscr U\) and then this
right-factor inertia.  It does not kill the stabilizer of the pencil
in \(\PGL(V)\).
\end{corollary}
\begin{proof}
The right factor acts trivially on \(\mathscr P\), and the quotient
map of groups is fppf surjective.  Objects and isomorphisms of the
second quotient therefore lift locally to the first, with differences
in that kernel.  Apply the descent proof of
Theorem~\ref{thm:rigidification-v151}.  Over an object represented by
a \(G^\circ\)-torsor the subgroup is its conjugation-associated
\(\operatorname{GL}(U)\)-bundle; it need not be a constant group
scheme on the base.  This describes the actual subgroup stack and
shows that no global splitting is assumed.  Composing with the first
rigidification gives the asserted successive universal property.
All remaining pencil automorphisms survive in the target.
\end{proof}

% BEGIN PRESERVED PART 42-pencil-specialization-v151.tex
\section{A common closed specialization of all quadratic pencils}
\label{sec:closed-pencil-v151}

We give an orbit-closure consequence directly on the pencil side.
Unlike restricting an equivalence to a previously specified locus, the
following argument constructs the specialization.  It works for every
pencil, including pencils with singular generic member.  Its mechanism
is elementary projective representation theory; the additional failure
algebra statement uses the all-pencil, exact-multiplicity construction.

\begin{theorem}[Explicit common closed orbit]
\label{thm:closed-pencil-v151}
Let \(\dim V=n\geq3\).  The action of \(\PGL(V)\) on
\(\mathscr P=\Gr(2,\Sym^2V)\) has a unique closed orbit, namely
\begin{equation}\label{eq:flag-orbit-v151}
 \mathscr C=\{\ell H:\ell\subset H\subset V,\
                      \dim\ell=1,\ \dim H=2\}.
\end{equation}
It is the flag variety \(\operatorname{Fl}(1,2;V)\), of dimension
\(2n-3\).  More precisely, after a fixed change of coordinates every
pencil has a degeneration along the one-parameter subgroup
\[
                  \lambda(t)=\diag(1,t,t^2,\ldots,t^2)
\]
to \(\langle x_1^2,x_1x_2\rangle\).  The pencil family is a
subbundle on the whole affine line.
\end{theorem}
\begin{proof}
Take independent forms \(q_0,q_1\).  The rational function
\(q_1/q_0\) on \(\Pj(V^*)\) is nonconstant.  In characteristic
zero its differential is not identically zero: on an affine coordinate
chart the common kernel of the coordinate derivations of the rational
function field is \(\C\).  Choose a point \([v]\) where \(q_0(v)\ne0\)
and this differential is nonzero.  Replace \(q_1\) by
\(q_1-(q_1(v)/q_0(v))q_0\).  Then \(q_1(v)=0\), whereas its
differential has a nonzero value on some vector \(w\) independent of
\(v\).  Euler's identity shows that its value on the radial vector
\(v\) is zero.  Extending \(v,w\) to a basis of \(V^*\) gives dual
coordinates in which
\[
 [x_1^2]q_0=a\ne0,\qquad [x_1^2]q_1=0,
                                  \qquad [x_1x_2]q_1=b\ne0.
\]
Here brackets denote a coefficient.

Apply \(\lambda(t)\) to the variables.  The forms
\[
 Q_0(t)=q_0(x_1,tx_2,t^2x_3,\ldots,t^2x_n),\qquad
 Q_1(t)=t^{-1}q_1(x_1,tx_2,t^2x_3,\ldots,t^2x_n)
\]
have polynomial coefficients: the only possible weight-zero term in
\(q_1\) was removed.  Their limits are \(a x_1^2\) and
\(b x_1x_2\).  For \(t\ne0\) they span a coordinate transform of
\(R\); at zero they are independent.  Thus their coefficient matrix
has rank two everywhere, defines a subbundle, and supplies the asserted
morphism from \(\A^1\) to \(\mathscr P\).

The orbit of \(\langle x_1^2,x_1x_2\rangle\) is exactly
\(\mathscr C\).  A pencil of the form \(\ell H\) recovers \(\ell\)
as its gcd and \(H\) by division.  These are scheme-theoretic inverse
operations on the exact-gcd-one-degree stratum by
Corollary~\ref{cor:canonical-exact-v151}, applied to \(g=h=1,r=2\).
Under that inverse the condition \(\ell\subset H\) is the usual
closed flag incidence.  The flag-to-pencil map is therefore a locally
closed immersion; since the flag variety is projective it is a closed
immersion.  It identifies \(\mathscr C\) with the flag variety, and
the dimension is \((n-1)+(n-2)\).  Every pencil orbit closure contains
\(\mathscr C\), by the constructed family and equivariance.  Any
closed orbit must consequently be \(\mathscr C\), proving uniqueness.
\end{proof}

\begin{corollary}[Flat closed specialization of failure algebras]
\label{cor:closed-failure-v151}
For every quadratic pencil \(R\) there is a finite locally free
family over \(\A^1\), of rank
\(\binom{n^2+d}{d}-\binom N2\), whose nonzero fibres are isomorphic
to \(\mathfrak A_R^{[d]}\) and whose zero fibre is
\(\mathfrak A_{\langle x_1^2,x_1x_2\rangle}^{[d]}\).
The intrinsic first relations throughout the family have exact gcd
\(\delta^{n-1}\).  In a fixed determinant frame, the two ordered
copies of the flag orbit form the unique closed \(G\)-orbit in the
pencil relation locus.  Equivalently the effective pencil failure
stack has the common closed specialization described in
Theorem~\ref{thm:closed-pencil-v151}.
\end{corollary}
\begin{proof}
Let \(R_t\) be the polynomial subbundle just constructed.  Its
Pluecker line gives, under the fixed exterior-duality identification,
a line bundle \(L_t^0\subset\mathbb S_\mu V^*\otimes\OO_{\A^1}\).
Set
\[
 \mathcal K_t=\delta^{n-1}(L_t^0\otimes\mathbb S_\mu U),\qquad
 \mathcal A_t=\Sym(E\otimes\OO_{\A^1})/
                       ((\mathcal K_t)+\mathfrak m^{d+1}).
\]
The first is a rank-\(M\) subbundle.  The second is finite locally
free by its graded direct-sum presentation, with the claimed rank.
Lemma~\ref{lem:primitive-pencil} gives the exact gcd at every fibre,
including zero; in fact this family lies scheme-theoretically in the
fixed-divisor exact stratum.  Equivariance and the local inverse
identify its nonzero fibres as asserted.  In a determinant frame the
pencil locus is \(G\times^{G^\circ}\mathscr P\).  The previous
closed-orbit calculation therefore gives the assertion for \(G\).
This is a statement inside that locus; it does not assert closedness
under arbitrary degenerations of the determinant form in the full
relation Grassmannian.
\end{proof}

The closed specialization is very coarse compared with the unmarked
local inverse: many pairwise nonisomorphic finite algebras have this
same limiting fibre.  The full first relation, not its common limit,
retains the original pencil.  The explicit family also differs from
transport of spectral rank loci: it crosses from arbitrary pencils to
a singular common-factor pencil without a regularity hypothesis.
Theorem~\ref{thm:boundary-normalization-v151}, on the other hand,
describes divisor collisions in the larger natural first-relation
boundary, outside the exact pencil image.  These are two distinct
geometric statements, and neither is a classification of every
congruence-orbit adjacency of quadratic pencils.

% BEGIN PRESERVED PART 39-transverse-relations-v149.tex
\section{Transverse deformations with unchanged determinant support}
\label{sec:transverse-relations}

The larger determinant-divisor stratum has deformations which are
not pencil deformations and need not satisfy either factor-invariance
condition.  We now identify their tangent directions and construct
complete families in which the homogeneous radical remains exactly
the determinant ideal.  Thus the intrinsic stratum theorem has content
outside the coefficient class, not merely a different description
of the same symmetry hypothesis.

Let \(J=L\otimes F\) be the primitive relation of a pencil
\(R\in\Gr(2,W)\).  In degree \(h=3n-4\), write the equivariant
Cauchy decomposition as
\[
                    S_h(E)=(A\otimes F)\oplus H_{\rm other},
\]
where \(A,F,L\) are as in \eqref{eq:primitive-pencils}.
The summand \(H_{\rm other}\) contains every partition different
from \(\mu\); this notation denotes its full vector space, not a
quotient chosen separately at each pencil.

\begin{theorem}[Normal directions to the pencil stratum]
\label{thm:transverse-pencil-directions}
The normal space at \(J\) of the pencil locus in \(X_{h,M}(E)\)
has the following equivariant direct-sum description:
\begin{equation}\label{eq:three-transverse-modes}
 \begin{split}
 N_J\simeq{}&
 \Hom(\bigwedge^2R,\bigwedge^2(W/R))\\
 &\oplus\bigl(\Hom(L,A/L)\otimes\mathfrak{sl}(F)\bigr)\\
 &\oplus\Hom(L\otimes F,H_{\rm other}).
 \end{split}
\end{equation}
The summands respectively record departure from the Pluecker locus,
departure from a full right multiplicity factor inside the same
Cauchy summand, and departure into other Cauchy summands.  They are
all transverse directions in the effective first-relation stack:
the infinitesimal \(G^\circ\)-orbit lies in the pencil tangent
space, so no vector in \(N_J\) is removed by that orbit.  Every
such direction is the first-order term of a local algebraic family
in the determinant-divisor stratum.  This last assertion refers to
the common-divisor stratum; preservation of the entire homogeneous
radical for complete families is supplied below.
\end{theorem}
\begin{proof}
The tangent space of a Grassmannian and the Cauchy splitting give
\begin{equation}\label{eq:ambient-cauchy-tangent}
 \Hom(J,S_h(E)/J)=
 \bigl(\Hom(L,A/L)\otimes\End(F)\bigr)
       \oplus\Hom(L\otimes F,H_{\rm other}).
\end{equation}
Under \(L\mapsto L\otimes F\), the tangent to \(\Pj(A)\) is the
first factor tensored with the identity of \(F\).
In characteristic zero the trace splits
\(\End(F)=\C I\oplus\mathfrak{sl}(F)\).
It remains to identify the normal space of the Pluecker embedding.
The natural map
\(\bigwedge^2W\to\bigwedge^2(W/R)\) has kernel \(R\wedge W\).
At the line \(\bigwedge^2R\), the tangent image is
\(\Hom(\bigwedge^2R,(R\wedge W)/\bigwedge^2R)\), which is the
image of \(\Hom(R,W/R)\).  Its quotient in the projective tangent
space is
\(\Hom(\bigwedge^2R,\bigwedge^2(W/R))\).
The common determinant twist in \(A\) cancels in this Hom space.
Substitution in \eqref{eq:ambient-cauchy-tangent} proves
\eqref{eq:three-transverse-modes} as a normal-space isomorphism.
No splitting of the tangent sequence of the Pluecker embedding is
being asserted.

The left group preserves the Pluecker locus, and the right group
fixes \(L\otimes F\); hence the differential of its orbit is in
the pencil tangent.  Passing to the quotient tangent complexes in
\eqref{eq:ambient-tangent-complex} does not change this normal
quotient.  Finally \(X_{h,M}(E)\) is open in a Grassmannian.  A
tangent vector is integrated by a graph with coefficients linear
in one parameter, after restricting to a neighbourhood of zero
which stays in this open set.  Multiplication by \(\delta^{n-1}\)
and Proposition~\ref{prop:relative-first-relation} give the claimed
family of finite flat algebras.
\end{proof}

\begin{theorem}[Complete fixed-support families through every pencil]
\label{thm:fixed-support-grassmannians}
Put \(e=n^2\).  For every pencil \(R\) there is an \(e\)-dimensional
subspace \(B\subset J=\widetilde J_R\) such that
\(\gcd(B)=1\) and the forms in \(B\) have no common zero on the
invertible matrix locus.  The projective Grassmannian
\begin{equation}\label{eq:fixed-support-grassmannian}
 \Sigma_B=\{J'\in\Gr(M,S_h(E)):B\subset J'\}
          \simeq\Gr(M-e,S_h(E)/B)
\end{equation}
contains \(J\), is contained in \(X_{h,M}(E)\), and satisfies
\begin{equation}\label{eq:fixed-radical}
              \sqrt{(\delta^{n-1}J')}=(\delta)
                         \quad\text{for every geometric }J'\in\Sigma_B.
\end{equation}
Its universal first relation gives a finite flat family of the
prescribed Hilbert function.  It contains a projective line through
\(J\) whose every other point is invariant under neither full
factor group.  Consequently determinant-radical relations through
each pencil admit positive-dimensional complete families outside
the recognized coefficient class, and their intrinsic deformation
problem is described by Theorem~\ref{thm:determinant-stratum}.
\end{theorem}
\begin{proof}
By Lemma~\ref{lem:primitive-pencil}, the residual coefficient vector
is nonzero on every invertible matrix.  Evaluation gives a morphism
\[
 \operatorname{GL}_n/\mathbb G_m\longrightarrow\Pj(J^*).
\]
Its projective image closure \(Y\) has dimension at most \(e-1\).
Since \(M>e\) for \(n\ge3\), a general projective linear subspace
of codimension \(e\) in \(\Pj(J^*)\) misses \(Y\).  For such a
subspace \(\Pj(B^\perp)\), the associated \(e\)-plane \(B\subset J\)
has no common zero on the invertible locus.  To justify the generic
avoidance directly, the incidence of pairs consisting of a point
of \(Y\) and an \(e\)-plane annihilated by it has dimension strictly
smaller than \(\Gr(e,J)\).  Its projection is closed because
\(Y\) is projective.

The condition \(\gcd(B)=1\) is a second nonempty open condition
on \(\Gr(e,J)\).  To see nonemptiness, take any nonzero \(b_1\in J\).
It has only finitely many irreducible factors.  Since \(\gcd(J)=1\),
a general \(b_2\in J\) avoids divisibility by each such factor;
then \(\gcd(b_1,b_2)=1\).  Extend their span to an \(e\)-plane.
Openness follows from Lemma~\ref{lem:universal-gcd}.  The two
nonempty opens intersect, as \(\Gr(e,J)\) is irreducible.

Every \(J'\) containing this \(B\) has greatest common divisor
\(1\), and its forms cannot vanish simultaneously on an invertible
matrix.  The common zero set of \(\delta^{n-1}J'\) is thus exactly
\(V(\delta)\).  The Nullstellensatz over an algebraically closed
field proves \eqref{eq:fixed-radical}.  This is a statement about
geometric homogeneous radicals.  The scheme-theoretic family used
here is the universal subbundle \(\delta^{n-1}\mathcal J'\), not
an operation of taking radicals on a nonreduced parameter ring.
Multiplication by this fixed nonzero form gives a subbundle;
its quotient in degree \(d\) and all the lower symmetric powers
are locally free.  The truncated algebras are therefore finite flat.

For the final assertion, choose \(j_0\in J\setminus B\), a
complement \(C\) to \(\C j_0\) in \(J\) containing \(B\), and
a nonzero decomposable vector
\[
 \eta\in\Sym^h V^*\otimes\Sym^h U\subset H_{\rm other}.
\]
The inequality \(M>e\) allows the choice of \(C\).
The family
\[
      J_{[a:b]}=C\oplus\C(aj_0+b\eta),\qquad[a:b]\in\Pj^1,
\]
is a projective line in \(\Sigma_B\) through \(J=J_{[1:0]}\).
For \(b\ne0\) its projection to the \((h)\)-Cauchy summand is
exactly the line \(\C\eta\).  A decomposable line in
\(\Sym^h V^*\otimes\Sym^h U\) is invariant under neither full
factor group: each of the two irreducible factors has dimension
greater than one.  Equivariance of the Cauchy projection therefore
excludes both possible admissible orientations of \(J_{[a:b]}\).
The determinant-preserver theorem says that these are the only
orientations after any change of the recovered matrix coordinates.
This proves the assertion intrinsically, not only in the chosen
matrix frame.
\end{proof}

\begin{remark}\label{rem:scope-transverse}
The three normal modes and the complete fixed-support families
separate two questions.  The effective pencil stratum has exactly
the deformation theory of quadratic pencils, including their
singular boundary.  The larger determinant-divisor stratum contains
additional directions, explicitly described here, which cease to
be pencil or coefficient relations.  Retaining the determinant
radical does not suppress these directions.  Conversely, no
classification of all ideals with a prescribed radical is needed
to describe this larger locally closed stratum and its family
inverse.  These assertions enlarge the reconstruction problem
without replacing the all-pencil theorem by a generic statement.
\end{remark}

% BEGIN PRESERVED PART 36-covering-moduli-v148.tex
\section{Covering maps inside one fibre orbit of quadratic pencils}
\label{sec:covering-moduli-v148}

We give a consequence native to the pencil construction.  It embeds
the equivalence problem for maps \(\Pj^1\to\Pj^1\) into
unmarked finite failure schemes, while keeping the fibre algebra
and every nilradical graded vector bundle fixed.  Unlike a varying
spectral invariant, the phenomenon occurs within a single orbit
of singular pencils.

Fix \(n\ge3\), a two-dimensional subspace
\(H=\langle x,y\rangle\subset V\), and the trivial source
bundle \(\mathcal U=\OO_{\Pj^1}^n\).  Let
\(f=[a:b]:\Pj^1\to\Pj(H)\) be a morphism of degree
\(m\ge1\), where \(a,b\) are homogeneous of degree \(m\)
with no common zero.  Its tautological line subbundle is
\[
 \mathcal L_f=\OO(-m)
       \xrightarrow{\ (a,b)\ }H\otimes\OO_{\Pj^1}.
\]
Multiplication in \(\Sym V\) defines a rank-two subbundle
\begin{equation}\label{eq:cover-pencil-v148}
 \mathcal R_f=\mathcal L_f H
       \subset\Sym^2V\otimes\OO_{\Pj^1},
 \qquad
 \mathcal R_f|_{[s:t]}=
   \langle (ax+by)x,(ax+by)y\rangle.
\end{equation}
Let \(Y_f=Y_{\mathcal R_f}^{[d]}\) be its order-\(d\)
failure thickening, with
\[
 N=\binom{n+1}{2},\qquad p=N-2,\qquad
 d=n^2+2n-4,\qquad M=\binom N2,\qquad e=n^2.
\]
All schemes below are compared without projection, grading,
embedding, matrix factors, or coefficient markings.

\begin{theorem}[Unmarked realization of the covering equivalence problem]
\label{thm:covering-pencil-moduli-v148}
For fixed \(n\ge3\) and \(m\ge1\), the schemes \(Y_f\)
have the following properties.

\smallskip
\noindent (i) They are finite locally free over their reduction
\(\Pj^1\), of rank \(\binom{e+d}{d}-M\), and are
Zariski locally trivial with the same Artin algebra as fibre.
Every fibre pencil is equivalent to \(\langle x^2,xy\rangle\).

\smallskip
\noindent (ii) Writing \(\mathfrak n_f\) for their nilradicals,
their graded vector bundles are independent of \(f\):
\begin{equation}\label{eq:cover-graded-bundles-v148}
 \begin{split}
 \mathfrak n_f^j/\mathfrak n_f^{j+1}
   &\simeq\OO^{\binom{e+j-1}{j}}\qquad(1\le j<d),\\
 \mathfrak n_f^d
   &\simeq\OO^{\binom{e+d-1}{d}-3M}
                              \oplus\OO(m)^{\oplus2M}.
 \end{split}
\end{equation}
Their abstract first-relation bundles are all
\(\OO(-2m)^{\oplus M}\), and their pencil and quotient bundles
are respectively
\(\OO(-m)^2\) and \(\OO(2m)\oplus\OO^{N-3}\).

\smallskip
\noindent (iii) For two degree-\(m\) maps,
\begin{equation}\label{eq:cover-equivalence-v148}
 Y_f\simeq Y_{f'}
 \quad\Longleftrightarrow\quad
 f'=\alpha\circ f\circ\sigma
 \quad\text{for some }\alpha,\sigma\in\PGL_2(\C).
\end{equation}
Thus the double-orbit set of degree-\(m\) morphisms
\(\Pj^1\to\Pj^1\) is realized faithfully by unmarked
quadratic-pencil failure schemes with all the additive data in
(ii) fixed.  This is a statement about isomorphism classes, not an
equivalence of their automorphism groupoids.
\end{theorem}
\begin{proof}
\emph{The pencil and local triviality.}
Multiplication by a nonzero linear form is injective on \(H\).
Since \((a,b)\) never vanishes, the map
\(\OO(-m)\otimes H\to\Sym^2H\otimes\OO\) has rank
two everywhere.  Its image is a subbundle with locally free
quotient.  In the basis \(x^2,xy,y^2\) its two columns are
\begin{equation}\label{eq:cover-columns-v148}
                 \begin{pmatrix}a&0\\ b&a\\0&b\end{pmatrix},
 \qquad\text{with Pluecker coordinates }(a^2,ab,b^2).
\end{equation}
It follows that \(\mathcal R_f\simeq\OO(-m)^2\).
The quotient of \(\Sym^2H\otimes\OO\) is a line bundle
of degree \(2m\).  A constant complement to \(\Sym^2H\)
in \(\Sym^2V\) gives
\(\mathcal S_f\simeq\OO(2m)\oplus\OO^{N-3}\).
On an open set where \(\mathcal L_f\) is trivialized, complete
its nonvanishing generator to a frame of \(H\otimes\OO\).
The inverse frame map takes \(\mathcal R_f\) to
\(\langle x^2,xy\rangle\).  Extend it by the identity on a
fixed complement of \(H\) in \(V\).  With the source bundle
trivial, it identifies the defining minors and the truncated
algebra with a product on that open set.  This proves local
triviality, not merely equality of fibre isomorphism classes.
Lemma~\ref{lem:moving-coefficients-v146} gives finite local
freeness and the asserted rank.

\emph{The entire graded bundle calculation.}
The first relations occur in degree \(d\), so all lower
positive graded pieces are the full constant symmetric powers of
\(E=V^*\otimes\C^n\).  Put
\(F(V)=\bigwedge^p\Sym^2V\), of dimension \(M\).
The coefficient line
\(\mathcal L_{\mathcal R_f}=\bigwedge^p\mathcal S_f^*\)
is \(\OO(-2m)\).  Under exterior duality
\[
 F(V)^*\simeq\bigwedge^2(\Sym^2V)
                              \otimes\det(\Sym^2V)^*,
\]
its inclusion is the Pluecker line of \(\mathcal R_f\).
The constant determinant twist does not affect the bundle
calculation.  Thus it lies in a fixed three-dimensional subspace
of \(F(V)^*\), with column \((a^2,ab,b^2)\).
There is an exact sequence
\begin{equation}\label{eq:conic-quotient-v148}
 0\longrightarrow\OO(-2m)
 \xrightarrow{\ (a^2,ab,b^2)^{\mathsf t}\ }\OO^3
 \xrightarrow{\left(\begin{smallmatrix}-b&a&0\\0&-b&a\end{smallmatrix}\right)}
                  \OO(m)^2\longrightarrow0.
\end{equation}
Indeed the composite is zero; the maximal minors of the second
map are \(b^2,-ab,a^2\), which have no common zero.  It is
therefore everywhere surjective.  Its rank-one kernel is the
image of the displayed nowhere vanishing column.  Taking a
constant complement proves
\begin{equation}\label{eq:coefficient-quotient-v148}
 (F(V)^*\otimes\OO)/\mathcal L_{\mathcal R_f}
                   \simeq\OO^{M-3}\oplus\OO(m)^2.
\end{equation}

The fixed Cauchy inclusion, followed by multiplication by the
fixed determinant, embeds
\(F(V)^*\otimes F(\C^n)\) as an \(M^2\)-dimensional
constant subspace of \(\Sym^d E\).
The first relation bundle is precisely
\(\mathcal L_{\mathcal R_f}\otimes F(\C^n)\) in that
subspace.  Choose a constant complement to the whole subspace,
not to its varying relation subbundle.  The top graded quotient is
therefore
\[
 \OO^{\binom{e+d-1}{d}-M^2}
 \oplus\left((F(V)^*\otimes\OO)/
                         \mathcal L_{\mathcal R_f}\right)^{\oplus M}.
\]
Equation~\eqref{eq:coefficient-quotient-v148} gives exactly
\eqref{eq:cover-graded-bundles-v148}.  It also shows that the
first relation bundle is \(\OO(-2m)^{\oplus M}\).
All assertions in (ii) concern vector bundles; their
multiplication maps have not been identified.

\emph{The global inverse is the covering equivalence.}
Suppose \(Y_f\simeq Y_{f'}\).  Their reductions give an
automorphism of \(\Pj^1\).  By
Theorem~\ref{thm:unrestricted-moving-v146}, or the actual
descent theorem with the pencil coefficient line, the varying
pencil subbundles are identified by that base map and a single
constant \(g\in\operatorname{GL}(V)\).
For a nonzero \(\ell\in H\), the common polynomial divisor
of the two-dimensional space \(\ell H\subset\Sym^2V\)
is precisely \(\ell\), up to scalar: \(x\) and \(y\)
have no common nonconstant divisor in \(\Sym V\).
Dividing the pencil space by this common factor recovers \(H\).
These operations are equivariant under \(\operatorname{GL}(V)\).
Consequently an equality between \(g(\ell H)\) and a pencil
of the form \(\ell'H\) forces \(gH=H\) and
\([g\ell]=[\ell']\).  The constant restriction \(g|_H\)
therefore supplies a single \(\alpha\in\PGL(H)\) for the
whole map.  Replacing the recovered base map by its inverse if
necessary gives the right-hand side of
\eqref{eq:cover-equivalence-v148}.

Conversely choose linear lifts of \(\alpha\) on \(H\) and
extend by the identity on a complement in \(V\).  Together with
the base map and the identity identification of the trivial source
bundles, this gives equivalent pencil data.  Their homogeneous
normal maps induce an isomorphism of the finite schemes.  This
proves (iii) and the theorem.
\end{proof}

\begin{corollary}[A pencil-native separating pair]
\label{cor:pencil-native-pair-v148}
For every \(n\ge3\), the degree-three maps
\[
          f_0=[s^3:t^3],\qquad f_1=[s^3+st^2:t^3]
\]
give nonisomorphic unmarked order-\(d\) failure schemes of
quadratic pencils, although they have the same Artin fibres, the
same conormal and first-relation bundles, the same pencil and
quotient bundles, and isomorphic nilradical graded vector bundles
in every degree.  For \(n=3\) their common finite rank is
\(167945\), and their top graded bundle is
\[
                  \OO^{75537}\oplus\OO(3)^{\oplus30}.
\]
\end{corollary}
\begin{proof}
The pairs of cubics have no common zero.  In the affine coordinate
\(z=s/t\) the maps are \(z^3\) and \(z^3+z\).
The first has two ramification points, both of index three.
The second has the two distinct zeros of \(3z^2+1\), each of
index two, and infinity of index three.  Independent projective
changes in source and target preserve this multiset of indices,
so the two maps are inequivalent.  The theorem proves all the
geometric assertions.  For \(n=3\),
\(e=9,d=11,M=15\), whence
\(\binom{20}{11}-15=167945\) and
\(\binom{19}{11}-45=75537\).
\end{proof}

This consequence uses the global inverse essentially.  Pointwise
local reconstruction returns the same pencil orbit at every
point for every \(f\); it cannot recover the map \(f\).
Even supplementing it with all the abstract bundles displayed in
the theorem leaves the covering equivalence problem unresolved.
It is the multiplication in the total unmarked thickening,
followed by descent to one constant left transformation, that
recovers that problem.  The double-orbit description is not an
assertion that the scheme has the automorphism group of the
cover: the unused complement of \(H\), the right factor, and
the higher-order kernel remain present, as recorded by the
previous automorphism theorems.

% BEGIN PRESERVED PART 21-spectral-specialization.tex
\section{Spectral torsion and rank-preserving specialization}
\label{sec:spectral-specialization}

The finite-neighbourhood theorem recovers more than the discriminant
or the corank of each member of a regular pencil.  We identify its
classical spectral invariant and then give a flat specialization on
which the entire discriminant and every reduced rank locus remain
constant, while the recovered invariant changes.  The two specializations
are separated at exactly the order of Theorem~\ref{thm:sharp-finite-pencil}.
The parameter space throughout is the usual Grassmannian of pencils;
the failure map is the original multiplication map, not an added
coefficient-evaluation block.

\subsection{The spectral torsion sheaf}

Let \(R\subset\Sym^2V\) be a regular pencil, meaning that it
contains a nonsingular quadratic form.  On its parameter line
\(L_R=\Pj(R)\), the tautological quadratic form gives an exact sequence
\begin{equation}\label{eq:spectral-torsion-sequence}
 0\longrightarrow V^*\otimes\OO_{L_R}(-1)
 \xrightarrow{\ q_R\ }V\otimes\OO_{L_R}
 \longrightarrow\mathcal T_R\longrightarrow0.
\end{equation}
The first arrow is injective because its determinant is nonzero at the
generic point of the smooth integral curve \(L_R\).  Thus
\(\mathcal T_R\) is a torsion sheaf of length \(n\), whose
zeroth Fitting divisor is the full discriminant.  The pair
\((L_R,\mathcal T_R)\) is considered up to isomorphism of the
parameter line.  Its definition involves no choice of a basis of \(R\).
The relation with elementary divisors is classical; see
\cite[\S2, Corollary~2.1 and equation~(5)]{FevolaMandelshtamSturmfels}.
We record the argument to make precise which datum is recovered.

\begin{lemma}[Sheaves on the affine spectral line]
\label{lem:spectral-sheaf-similarity}
Let \(A,B\) be symmetric complex matrices with \(A\) invertible.
On the affine line with coordinate \(z\), the cokernel of
\(zA-B\) is the finite-dimensional \(\C[z]\)-module on
\(\C^n\) for which multiplication by \(z\) is
\(C=A^{-1}B\).  An isomorphism of such sheaves, in this fixed
coordinate, is equivalent to similarity of the corresponding matrices.
\end{lemma}
\begin{proof}
Multiplication on the target by \(A^{-1}\) identifies the
presentation with \(zI-C\).  Reduction of a polynomial vector
modulo these relations sends \(\sum_jz^jv_j\) to
\(\sum_jC^jv_j\); constant vectors give the inverse identification.
The equivalence between coherent sheaves on an affine scheme and
finitely generated modules identifies a finite-length sheaf supported
in this affine line with this module.  The support is finite, so its
underlying complex vector space is finite dimensional.  A module
isomorphism is a linear isomorphism \(G\) satisfying
\(GC=C'G\), which is exactly similarity.  Extension by zero across
a point outside the support loses no sheaf data.
\end{proof}

\begin{lemma}[Polynomial square roots of invertible operators]
\label{lem:artinian-polynomial-square-root}
For every invertible complex matrix \(K\) there is
\(h\in\C[x]\) with \(h(K)^2=K\).  If \(K\) is
self-adjoint for a nondegenerate symmetric form and commutes with
\(C\), then \(h(K)\) has both properties.
\end{lemma}
\begin{proof}
Write the minimal polynomial as
\(\prod_\lambda(x-\lambda)^{a_\lambda}\), where
\(\lambda\ne0\).  In its local Artinian factor choose a square
root \(\mu_\lambda^2=\lambda\), and put
\[
 h_\lambda(x)=\mu_\lambda
 \sum_{j=0}^{a_\lambda-1}\binom{1/2}{j}
          \left(\frac{x-\lambda}{\lambda}\right)^j.
\]
The formal binomial identity, truncated modulo
\((x-\lambda)^{a_\lambda}\), gives \(h_\lambda(x)^2=x\).
The Chinese remainder theorem joins these polynomials to a single
class modulo the minimal polynomial.  Evaluation at \(K\) proves
the assertion even when \(K\) is not semisimple.  Polynomials in
one self-adjoint operator remain self-adjoint, and every polynomial
in \(K\) commutes with every operator commuting with \(K\).
\end{proof}

\begin{theorem}[Spectral reconstruction at the first relation]
\label{thm:spectral-finite-torelli}
Let \(n\ge3\), \(d=n^2+2n-4\), and let \(R,R'\) be
regular pencils.  The following are equivalent:
\begin{enumerate}
\item the unmarked schemes \(Z_R^{[d]}\) and \(Z_{R'}^{[d]}\)
are isomorphic;
\item there is an isomorphism \(f:L_R\to L_{R'}\) with
\(\mathcal T_R\simeq f^*\mathcal T_{R'}\);
\item the pencils \(R,R'\) are projectively equivalent.
\end{enumerate}
In particular the finite scheme determines the full root-labelled
spectral torsion sheaf, not only its Fitting divisor or its support.
This is an equivalence of isomorphism classes of objects; it is not
an assertion identifying their automorphism groups or moduli stacks.
\end{theorem}
\begin{proof}
The equivalence of the first and third conditions is
Theorem~\ref{thm:sharp-finite-pencil}.  A congruence of pencils
induces an isomorphism of \eqref{eq:spectral-torsion-sequence},
so the third implies the second.

For the converse choose a common affine coordinate whose point at
infinity avoids the two spectral supports, using \(f\) to identify
the lines.  Write the corresponding symmetric pairs as \((A,B)\)
and \((A',B')\), with \(A,A'\) invertible, and set
\(C=A^{-1}B\), \(C'=A'^{-1}B'\).  Choose the presentation \(zA-B\) for the affine
coordinate on both lines.  Lemma~\ref{lem:spectral-sheaf-similarity}
identifies the two sheaves with the modules of \(C,C'\), so their
isomorphism gives a similarity.  After applying
that change of coordinates, assume \(C'=C\).  Both \(A\) and
\(A'\) are nondegenerate symmetric forms for which \(C\) is
self-adjoint.  Then \(K=A^{-1}A'\) commutes with \(C\)
and is self-adjoint for \(A\).  Lemma~\ref{lem:artinian-polynomial-square-root}
provides a polynomial \(h\) with \(h(K)^2=K\).  Put \(H=h(K)\).  It commutes with \(C\), is
self-adjoint for \(A\), and satisfies
\[
 H^tAH=A',\qquad H^tBH=H^tACH=A'C=B'.
\]
Thus the pairs, and hence the pencils, are congruent.  This also shows
why the spectral sheaf alone suffices here; no additional pairing on
that sheaf is required over \(\C\).
\end{proof}

For \(x\in\Supp\mathcal T_R\), put
\begin{equation}\label{eq:spectral-length-profile}
 \ell_a(x)=\length\bigl(\mathcal T_{R,x}/
                  \mathfrak m_x^a\mathcal T_{R,x}\bigr),
 \qquad a\ge0.
\end{equation}
If the elementary-divisor exponents at \(x\) are
\(e_1,\ldots,e_b\), the structure theorem over the discrete
valuation ring \(\OO_{L_R,x}\) gives
\begin{equation}\label{eq:spectral-partition-inversion}
 \ell_a(x)=\sum_{j=1}^b\min(a,e_j),\qquad
 \ell_a(x)-\ell_{a-1}(x)=\#\{j:e_j\ge a\}.
\end{equation}
Thus these intrinsic local lengths recover every exponent, including
repeated roots.  The corank of the pencil at \(x\) gives only
\(\ell_1(x)\), and the multiplicity of the discriminant gives
only the eventual value \(\ell_a(x)=\sum_j e_j\).  The
intermediate values are genuinely additional information.

\subsection{A specialization with constant discriminant and coranks}

The following family keeps both ends of this length profile fixed.
Fix \(i\in\C\) with \(i^2=-1\), and set
\begin{equation}\label{eq:rank-preserving-four-block}
 N_0=\begin{pmatrix}1&i\\i&-1\end{pmatrix},\qquad
 P_t=\frac t2N_0,\qquad
 A_0=\begin{pmatrix}0&I_2\\I_2&0\end{pmatrix},\qquad
 C_t=\begin{pmatrix}P_t&I_2\\0&P_t\end{pmatrix}.
\end{equation}
Put
\begin{equation}\label{eq:symmetric-specialization-pair}
 B_t=A_0C_t=\begin{pmatrix}0&P_t\\P_t&I_2\end{pmatrix},
 \qquad R_t=\langle A_0,B_t\rangle\subset\Sym^2\C^4.
\end{equation}
Both matrices are symmetric and \(A_0\) is invertible.  They
are independent for every \(t\).  For \(n>4\), append
\(I_{n-4}\) to \(A_0\) and append
\(\diag(\mu_5,\ldots,\mu_n)\) to \(B_t\), where the
\(\mu_j\) are fixed distinct nonzero numbers.  Denote the resulting
matrices by \(A^{(n)},B_t^{(n)}\) and their pencil again by \(R_t\).

\begin{theorem}[Separation inside fixed reduced rank data]
\label{thm:rank-preserving-specialization}
For every \(n\ge4\), the family \(R_t\), \(t\in\A^1\),
has the following properties.
\begin{enumerate}
\item Its full discriminant, with multiplicities, is independent of \(t\):
\begin{equation}\label{eq:fixed-discriminant-specialization}
 \det(sA^{(n)}+uB_t^{(n)})
       =s^4\prod_{j=5}^n(s+\mu_j u).
\end{equation}
For every minor size, the reduced determinantal locus on the marked
parameter line is also independent of \(t\).
\item At the root \(s=0\), the spectral module has exponents
\((3,1)\) for \(t\ne0\) and \((2,2)\) for \(t=0\).
Consequently \(Z_{R_t}^{[d]}\not\simeq Z_{R_0}^{[d]}\)
for every \(t\ne0\), whereas all the punctured fibres are isomorphic.
\item The schemes \(Z_{R_t}^{[d]}\) form a flat projective family
over \(\A^1\), with fixed reduction.  After one identification
of the socle quotient bundles, every truncation of order \(k<d\)
is a constant family.  Thus the first separating order remains
exactly \(d=n^2+2n-4\) even under the constraints in part~(1).
\end{enumerate}
\end{theorem}
\begin{proof}
The identities \(N_0^2=0\) and \(\rank N_0=1\) give
\begin{equation}\label{eq:rank-preserving-power-identities}
 C_t^2=\begin{pmatrix}0&tN_0\\0&0\end{pmatrix},\qquad
 C_t^3=0,\qquad
 C_t=\binom{I_2}{P_t}\begin{pmatrix}P_t&I_2\end{pmatrix}.
\end{equation}
The right factor in the last expression is surjective and the left
factor is injective.  Hence \(\rank C_t=2\) for every \(t\),
whereas \(\rank C_t^2=1\) precisely when \(t\ne0\).
The nilpotent partitions are therefore \((3,1)\) and \((2,2)\).
Also \(\det A_0=1\) and \(\det(sI_4+uC_t)=s^4\),
proving \eqref{eq:fixed-discriminant-specialization}.  At \(s=0\)
the full pencil has corank two for all \(t\); at each remaining
root it has corank one, and elsewhere corank zero.  Equality of these
functions on \(\Pj^1(\C)\) implies equality of the reduced
closed determinantal subschemes for each minor size.

At the unique multiplicity-four root, formula
\eqref{eq:spectral-partition-inversion} yields respectively
\[
 (\ell_1,\ell_2,\ell_3,\ldots)=(2,3,4,\ldots),
 \qquad (2,4,4,\ldots).
\]
No automorphism of the parameter line can exchange this root with a
simple one.  Theorem~\ref{thm:spectral-finite-torelli} separates
the special fibre from the punctured fibres and identifies the latter
with one another.  This is a statement about geometric fibres; no
choice of a simultaneous trivializing congruence on \(\A^1\setminus0\)
is needed.

For flatness we use the relative gluing of
Theorem~\ref{thm:universal-finite-neighbourhood}; the following
coordinates identify its pullback in the present family.  First note that \(t\mapsto R_t\) is a morphism
to \(\Gr(2,\Sym^2V)\), not merely a collection of pencils.
In symmetric matrix coordinates the two fixed linear functionals
\[
 a(Q)=Q_{13}+iQ_{14},\qquad b(Q)=Q_{33}
\]
satisfy \(a(A^{(n)})=b(B_t^{(n)})=1\) and
\(a(B_t^{(n)})=b(A^{(n)})=0\).  Their common kernel is a
fixed complement to every \(R_t\).  It trivializes
\(S_{R_t}=\Sym^2V/R_t\), and hence identifies the relative
socle Grassmannian with \(B\times\A^1\).

On a frame chart of its tautological bundle, the graph description
\eqref{eq:finite-homogeneous-ideal} identifies the truncated algebra
with
\begin{equation}\label{eq:relative-spectral-finite-algebra}
 \bigoplus_{j=0}^{d-1}\Sym^j E^*
       \ \oplus\ \bigl(\Sym^d E^*/K_{R_t}\bigr).
\end{equation}
By Theorem~\ref{thm:finite-parameter-immersion}, \(K_{R_t}\)
is a subbundle of constant rank \(\binom N2\), with locally
free quotient.  This statement commutes with pullback to the present
base and with right changes of frame, so the local descriptions glue.
Thus \eqref{eq:relative-spectral-finite-algebra} is locally free
on \(B\times\A^1\), of rank
\(\binom{n^2+d}{d}-\binom N2\).  The graph projection is a
finite flat morphism to that base.  Since \(B\) is projective,
its composition with projection to \(\A^1\) is flat and
projective.  Below degree \(d\) there is no relation, and the
fixed complement identifies the entire lower truncated algebra,
including multiplication, with a constant family.  These arguments
prove flatness of the finite neighbourhoods; they make no assertion
about flatness of the unrestricted failure schemes.
\end{proof}

The qualifier ``reduced'' in part~(1) is necessary.  At \(s=0\),
in a local parameter \(z=s/u\), the gcd of the three-row minors of
the four-dimensional block is a unit times \(z\) for \(t\ne0\)
and a unit times \(z^2\) for \(t=0\).  Both radicals are
\((z)\).  Thus the family retains exactly the higher Fitting
information discarded by reduction; it does not contradict the
classical classification by the full determinantal ideals.

\subsection{An independently visible change of reciprocal geometry}

For a regular pencil let its reciprocal curve be the projective closure
of the inverses of its nonsingular members.  These curves and their
strata are studied independently of multiplication-failure schemes in
\cite[\S3]{FevolaMandelshtamSturmfels}.

\begin{corollary}[A conic specializes to a line]
\label{cor:reciprocal-conic-line}
For the four-dimensional family \eqref{eq:symmetric-specialization-pair},
the reciprocal curve is a smooth plane conic when \(t\ne0\)
and a projective line when \(t=0\).  The unmarked twentieth
neighbourhood distinguishes these two cases, while all earlier
neighbourhoods, the full discriminant, and all reduced rank loci agree.
\end{corollary}
\begin{proof}
For \(s\ne0\), multiplication on the right by the fixed
invertible matrix \(A_0^{-1}\) identifies the projective inverse
with
\[
 (sI_4+uC_t)^{-1}
   =s^{-1}I_4-us^{-2}C_t+u^2s^{-3}C_t^2.
\]
If \(t\ne0\), the matrices \(I_4,C_t,C_t^2\) are
linearly independent: their span has dimension equal to the degree
three of the minimal polynomial.  The closure is therefore the
quadratic Veronese image \([s:u]\mapsto[s^2:-su:u^2]\),
a smooth conic in its spanning plane.  At \(t=0\) the last term
vanishes, and \(I_4,C_0\) are independent, giving a line.
The remaining assertions follow from
Theorem~\ref{thm:rank-preserving-specialization}, with \(d=20\).
The degree drop does not assert that the reciprocal curves themselves
form a flat family; the flat family proved above consists of the
finite failure neighbourhoods.
\end{proof}

% BEGIN PRESERVED PART 22-relative-spectral-strata.tex
\section{The universal finite neighbourhood and spectral strata}
\label{sec:relative-spectral-strata}

The finite-order construction is a family of actual neighbourhoods in
multiplication-failure schemes.  This fact is needed when passing from
individual pencils to degenerations and infinitesimal parameter spaces.
We establish it before discussing the spectral incidence schemes.

Put \(P=\Gr(2,\Sym^2V)\), and write \(\mathcal R\) and
\(\mathcal S=(\Sym^2V\otimes\OO_P)/\mathcal R\) for its
universal relation and quotient bundles.  Let
\[
 \pi:\mathcal B=\Gr_P(n,\mathcal S)\longrightarrow P,
 \qquad \mathcal E=\Hom(\mathcal U,V\otimes\OO_{\mathcal B}),
\]
where \(\mathcal U\) is tautological of rank \(n\).  As before,
\(N=n(n+1)/2\), \(p=N-2\), \(d=n+2p\), and
\(M=\binom N2\).

\begin{theorem}[Universal neighbourhood and arbitrary base change]
\label{thm:universal-finite-neighbourhood}
The universal multiplication-failure scheme has a canonical order-\(d\)
neighbourhood \(\mathcal Z^{[d]}\) of \(\mathcal B\).  There is a
rank-\(M\) vector subbundle
\(\mathcal K\subset\Sym^d\mathcal E^*\) such that
\begin{equation}\label{eq:universal-truncated-algebra}
 \mathcal Z^{[d]}=
 \Spec_{\mathcal B}\left(
 \Sym\mathcal E^*/
 (\mathcal K,\Sym^{d+1}\mathcal E^*)\right).
\end{equation}
This is finite locally free over \(\mathcal B\), of rank
\begin{equation}\label{eq:universal-finite-rank}
                   \binom{n^2+d}{d}-\binom N2,
\end{equation}
and is projective and flat over \(P\).  Its construction, its
ideal-adic filtration, and its identification with the actual
multiplication-failure neighbourhood commute with every morphism of
complex schemes \(T\to P\), including nonreduced ones.
For \(k<d\), its order-\(k\) truncation is
\(\Spec_{\mathcal B}\bigoplus_{j=0}^k\Sym^j\mathcal E^*\),
with truncated multiplication.  After trivializing \(\mathcal S\)
on the parameter base, these lower-order families are constant.
\end{theorem}
\begin{proof}
The locus in \(\Gr_P(n,V\oplus\mathcal S)\) on which projection
to \(\mathcal S\) is injective is canonically \(\Tot\mathcal E\).
It is an open neighbourhood of \(\mathcal B\), not a choice of
formal coordinates.  On a frame chart of \(\mathcal U\), let
\(T:\mathcal U\to V\) be the universal graph map and let
\(\gamma:\Sym^2V\to\mathcal S\) be the quotient.  The same
block elimination that proves
\eqref{eq:finite-homogeneous-ideal} works over its coordinate ring:
the ideal of the multiplication-failure scheme is
\[
                   (\det T)I_p(\gamma\Sym^2T).
\]
Every generator is homogeneous of degree \(d\).

Here is also a coordinate-free description of its first relation
bundle.  The coefficient map, followed by determinant multiplication,
is a morphism
\begin{equation}\label{eq:universal-relation-bundle-map}
 \det V^*\otimes\det\mathcal U\otimes\det\mathcal S^*
       \otimes\bigwedge^p\Sym^2\mathcal U
 \longrightarrow \Sym^d\mathcal E^*.
\end{equation}
The inclusion used in this expression is
\(\det\mathcal S^*\hookrightarrow(\bigwedge^p\Sym^2V)^*\).
The determinant factors evaluate \(\det T\), and the remaining
factors evaluate \(\bigwedge^p\Sym^2T\) followed by the quotient.
Thus the image is exactly the span of the displayed maximal minors,
not a separately inserted coefficient module.

On a frame chart this map is the fixed injective linear coefficient
map \(\nu\) of Theorem~\ref{thm:finite-parameter-immersion},
restricted to a line subbundle tensored with the fixed space
\(\bigwedge^p\Sym^2\C^n\).  Its image is a subbundle of rank
\(M\): locally split the line in the source, and split the fixed
complex linear inclusion \(\nu\).  Both splittings remain
splittings over any complex algebra.  Denote its image by
\(\mathcal K\).  A right change of frame of \(\mathcal U\)
changes the graph coordinates and this coefficient map
contragrediently.  The determinant character in
\eqref{eq:universal-relation-bundle-map} is exactly the one arising
from \(\det T\).  Hence these local ideals and their degree-\(d\)
subbundles agree on overlaps.  In particular, they glue as ideals
of the given graph neighbourhood.

Let \(\mathfrak b\) be the ideal of the zero section inside the
multiplication-failure scheme.  In graph coordinates its preimage
is the positive-degree ideal.  Quotienting by
\(\mathfrak b^{d+1}\) therefore gives precisely
\eqref{eq:universal-truncated-algebra}.  This identifies the gluing
with the actual ideal-adic neighbourhood, not merely with algebras
having the correct geometric fibres.  Its underlying module is
\[
 \bigoplus_{j=0}^{d-1}\Sym^j\mathcal E^*
       \ \oplus\ (\Sym^d\mathcal E^*/\mathcal K).
\]
All terms are locally free.  Counting monomials in \(n^2\)
variables gives \eqref{eq:universal-finite-rank}.  The morphism
\(\mathcal B\to P\) is smooth and projective, so the asserted
flatness and projectivity follow.  Finally the graph description,
the split coefficient inclusion, and the homogeneous quotient all
commute with base change.  Below degree \(d\) the ideal contributes
no relation.  This proves the final assertion as well.
\end{proof}

\begin{remark}\label{rem:relative-reduction-scope}
Over a nonreduced base, the theorem concerns the specified relative
socle subscheme \(\mathcal B_T\) and its ideal.  It does not identify
that ideal with the absolute nilradical of \(\mathcal Z_T^{[d]}\),
which can also contain nilpotents from the base.  On geometric fibres
\(\mathcal B_t\) is exactly the reduction.  Nor is this theorem an
equivalence of unmarked deformation groupoids: extra automorphisms
of truncated algebras are not being identified with pencil
congruences.
\end{remark}

\subsection{Spectral incidence schemes from the first relation space}

On the regular-pencil locus, pass locally to ordered generators
\((A,B)\) with \(A\) invertible, and put \(C=A^{-1}B\).
For a marked point with affine coordinate \(z\), set \(X=C-zI\).
The next elementary presentation works over an arbitrary coefficient
ring; no diagonalization or division by eigenvalue differences is
involved.

\begin{lemma}[Truncated spectral presentation]
\label{lem:truncated-spectral-presentation}
Let \(D\) be a commutative ring, \(X\in\operatorname{Mat}_n(D)\),
and \(a\ge1\).  There is a natural isomorphism
\begin{equation}\label{eq:truncated-module-power}
 \coker(wI-X)\otimes_{D[w]}D[w]/(w^a)
                  \simeq \coker(X^a:D^n\to D^n).
\end{equation}
Equivalently, the \(an\)-square block presentation with diagonal
blocks \(X\) and subdiagonal blocks \(-I\) is equivalent, by
invertible row and column operations over \(D\), to
\(I_{n(a-1)}\oplus X^a\).
\end{lemma}
\begin{proof}
Map the class of \(\sum_{j=0}^{a-1}w^jv_j\) to the class of
\(\sum_{j=0}^{a-1}X^jv_j\) modulo \(X^aD^n\).
The relation \(wv=Xv\) and the relation \(w^av=0\) show that
this map is well-defined.  Sending a constant vector to its class
is inverse.  Eliminating successively the blocks with coefficient
\(-I\) gives the stated presentation equivalence, using only unit
pivots.  This also proves naturality under ring homomorphisms.
\end{proof}

\begin{theorem}[Scheme-theoretic spectral readout]
\label{thm:relative-spectral-readout}
In the fixed tensor presentation of
Theorem~\ref{thm:finite-parameter-immersion}, the first relation
subspace determines, on the regular-pencil locus, every closed
spectral incidence scheme
\begin{equation}\label{eq:spectral-incidence-minors}
 \mathfrak S_{a,h}
       =V\bigl(I_{n-h+1}((C-zI)^a)\bigr),
       \qquad 1\le a,h\le n.
\end{equation}
This determination commutes with arbitrary complex base change.
At a geometric point, membership is equivalent to
\(\length(\mathcal T_z/\mathfrak m_z^a\mathcal T_z)\ge h\).
Consequently the schemes of specified spectral length profiles, and
not only their reduced supports, can be read from the order-\(d\)
relation subspace.
\end{theorem}
\begin{proof}
The closed immersion \(\Phi\) of
Theorem~\ref{thm:finite-parameter-immersion} has an inverse on its
scheme-theoretic image.  Thus a family of relation subspaces in that
image recovers its morphism to \(P\), over any base, and hence its
universal pencil.  In the displayed chart its spectral sheaf has
presentation \((z+w)I-C=wI-X\).  Lemma~\ref{lem:truncated-spectral-presentation}
gives the presentation \(X^a\) for restriction to the order-\(a\)
fat point centred at \(z\).  Its \((h-1)\)-st Fitting ideal is
exactly the ideal in \eqref{eq:spectral-incidence-minors}; Fitting
ideals of finitely presented modules commute with base change.
On a field the cokernel has dimension at least \(h\) exactly when
these minors vanish.

These descriptions are independent of the chosen frames.  A change
of basis in \(V\) conjugates \(X\); a change of affine coordinate
replaces the parameter ideal at the marked point by a unit multiple
locally.  Alternatively, restriction of the spectral sheaf to the
fat diagonal on its parameter line defines the same finite module
without a coordinate.  Thus the Fitting ideals glue.  Rank equality
is imposed by the corresponding rank inequality together with the
open complement of the next inequality.  Intersecting these locally
closed conditions gives each specified length profile.  Formula
\eqref{eq:spectral-partition-inversion} then identifies its geometric
fibres with the prescribed elementary-divisor partition.
\end{proof}

\subsection{Realizing the complete Segre stratification}

Let \(P_\sigma\) denote the reduced closure in \(P\) of a
regular Segre stratum, and let \(J_\sigma=\Phi(P_\sigma)\).
Because \(\Phi\) is a closed immersion, it gives isomorphisms
\begin{equation}\label{eq:segre-scheme-realization}
 P_\sigma\simeq J_\sigma,
 \qquad
 P_{\sigma_1}\cap_P\cdots\cap_P P_{\sigma_s}
 \simeq
 J_{\sigma_1}\cap_{\Phi(P)}\cdots\cap_{\Phi(P)}J_{\sigma_s},
\end{equation}
where the intersections carry their scheme structures.  In
particular their local rings and formal completions are preserved.
Pulling back Theorem~\ref{thm:universal-finite-neighbourhood} realizes
these parameter schemes by flat projective families of the actual
finite failure neighbourhoods.  The classical inclusions among
Segre-stratum closures \cite[Theorem~5.1]{FevolaMandelshtamSturmfels}
are therefore all realized, rather than just the one degeneration
of Theorem~\ref{thm:rank-preserving-specialization}.

There are two distinct assertions here.  The Segre closure order is
classical, and transport of closed subschemes through a closed
immersion is formal.  The geometric input proved here is the
base-change-compatible realization by actual multiplication-failure
neighbourhoods and the recovery of their spectral Fitting schemes.
The ambient parameter space in \eqref{eq:segre-scheme-realization}
is the image of \(\Phi\), not a moduli space of every possible
unmarked deformation of a finite scheme.

% BEGIN PRESERVED PART 25-fixed-spectral-strata-v145.tex
\section{Fixed-discriminant strata and all admissible specializations}
\label{sec:fixed-spectral-strata}

The universal family of Section~\ref{sec:relative-spectral-strata}
and the example of Section~\ref{sec:spectral-specialization}
lead to a classification of the entire fixed-discriminant,
fixed-reduced-rank locus. The spectral orbit order is classical.
Every allowed specialization is realized here by a flat family of
the original finite multiplication-failure neighbourhoods, and the
same first relation detects every strict transition.

Fix distinct \(\alpha_1,\ldots,\alpha_s\in\C\), positive
\(m_i\) with \(\sum m_i=n\ge3\), and \(1\le c_i\le m_i\).
Fix a nonsingular symmetric \(A\), and let \(\mathfrak p_A\)
be its self-adjoint endomorphisms. Consider the reduced locally closed
variety
\begin{equation}\label{eq:fixed-spectral-locus}
 X_{\boldsymbol m,\boldsymbol c}(A)=
 \left\{C\in\mathfrak p_A:
 \det(zI-C)=\prod_i(z-\alpha_i)^{m_i},\quad
 \dim\ker(C-\alpha_i I)=c_i\right\}.
\end{equation}
Exclude \(s=1,c_1=n\), for which \(C\) is scalar and
\(\langle A,AC\rangle\) is not a pencil. Every other datum
is realized. Put \(R_C=\langle A,AC\rangle\),
\(d=n^2+2n-4\), and \(Z_C=Z_{R_C}^{[d]}\).
The full binary determinant is
\(\det(A)\prod_i(u+\alpha_i v)^{m_i}\). The coranks and
labelled roots specify all reduced determinantal loci on the line.

Write \(\mathcal P(m,c)\) for partitions into exactly \(c\)
positive parts. For equal-size partitions, \(\mu\preceq\lambda\)
means \(\sum_{j\le a}\mu_j\le\sum_{j\le a}\lambda_j\)
for every \(a\); write \(\lambda'\) for the conjugate partition.

\subsection{The classical orthogonal input}
\label{sec:orthogonal-input-v145}

We isolate the orbit theorem used below, with its precise group and
source.  Let \((H,A)\) be a nondegenerate complex symmetric space of
dimension \(m\).  The acting group is the full orthogonal group
\(O(A)\), and the operators are \(A\)-self-adjoint, not elements
of the skew-adjoint Lie algebra \(\mathfrak o(A)\).

\begin{lemma}[Self-adjoint nilpotent orbits]
\label{lem:orthogonal-orbits-v145}
The full \(O(A)\)-orbits of self-adjoint nilpotent operators are
indexed by all partitions \(\lambda\) of \(m\).  If
\(\mathcal O_\lambda\) denotes the orbit, then
\[
 \dim\mathcal O_\lambda=
       \frac12\left(m^2-\sum_a(\lambda'_a)^2\right),\qquad
 \mathcal O_\mu\subset\overline{\mathcal O_\lambda}
       \ \Longleftrightarrow\ \mu\preceq\lambda.
\]
An \(O(A)\)-orbit is one \(SO(A)\)-orbit when \(\lambda\)
has an odd part.  When every part is even, it is the disjoint union
of two \(SO(A)\)-orbits, interchanged by an isometry of determinant
\(-1\).  The closure statement above is for the full orthogonal
orbit and does not identify the closures of its two components.
\end{lemma}
\begin{proof}
The classification and closure assertions are
\cite[Proposition~1, p.~444; Theorem~1, p.~447]{Ohta1986}, respectively,
with \(\varepsilon=1\) and \(G(V)\) as defined in
\cite[\S1.1, p.~443]{Ohta1986}.  Remark~1 on p.~444 specifies that
all partitions occur; the order in \S1.4 on p.~446 is precisely the
row-dominance order used here.  The proof of Theorem~1 occupies
pp.~447--451.  These references concern the full group, not its
identity component.

For the dimension formula there is a short independent calculation.
The trace pairing on \(\End H\) is nondegenerate on both its
self-adjoint and skew-adjoint summands.  For a self-adjoint nilpotent
\(N\), the commutator with \(N\) exchanges the summands, and the
two resulting maps are adjoint up to sign for that pairing.  They
have equal rank.  Consequently
\(\dim\operatorname{GL}(H)N=2\dim O(A)N\), also the identity
in \cite[\S2.4, Remark~8(i), p.~456]{Ohta1986}.
The general linear centralizer has dimension
\(\sum_{i,j}\min(\lambda_i,\lambda_j)
=\sum_a(\lambda'_a)^2\), as follows by writing homomorphisms
between the cyclic \(\C[t]\)-modules of the Jordan decomposition.
This proves the stated dimension.

We give the component argument explicitly.  The good bases of
\cite[Lemma~1, p.~444]{Ohta1986} decompose the symmetric space
orthogonally into cyclic Jordan blocks, each with a reversed-identity
Gram matrix.  If one block has odd size, its negative identity,
extended by the identity on the other blocks, centralizes \(N\)
and has determinant \(-1\).  The centralizer then meets both
components of \(O(A)\), so the orbit is a single \(SO(A)\)-orbit.

For the converse, write
\[
 E_k=\ker N^k/(\ker N^{k-1}+N\ker N^{k+1}).
\]
The form induced by \(A(u,N^{k-1}v)\) on \(E_k\) is
nondegenerate and symmetric; this is immediate on the same cyclic
blocks.  A commuting isometry induces \(g_k\in O(E_k)\).
A filtration of the Jordan module by its cyclic lengths and powers
of \(t\) gives
\[
                        \det g=\prod_k(\det g_k)^k.
\]
Indeed the multiplicity-space action of a length-\(k\) block occurs
\(k\) times on the associated graded; the strictly triangular
terms of a module automorphism do not change its determinant.
If all lengths are even, this product is one.  The centralizer lies
in \(SO(A)\), and the index-two quotient of \(O(A)\) therefore
gives exactly two \(SO(A)\)-orbits.  Since \(SO(A)\) is connected,
these are the connected components of the full orbit.  This does not
preclude their closures from meeting.
\end{proof}


\begin{theorem}[Classification and specialization at the first relation]
\label{thm:fixed-spectral-classification}
The following statements hold.
\begin{enumerate}
\item The \(O(A)\)-orbits in \(X_{\boldsymbol m,\boldsymbol c}(A)\)
are indexed by \(\boldsymbol\lambda\in\prod_i\mathcal P(m_i,c_i)\),
with dimensions
\begin{equation}\label{eq:fixed-spectral-orbit-dimension}
 \dim\mathcal O_{\boldsymbol\lambda}
 =\frac12\left(n^2-\sum_i\sum_a(\lambda'_{i,a})^2\right).
\end{equation}
Unmarked isomorphism classes of \(Z_C\) are these tuples modulo
the permutation group induced on the roots by projective transformations
preserving their multiplicities and coranks.
\item A specialization from \(\boldsymbol\lambda\) to
\(\boldsymbol\mu\) inside \eqref{eq:fixed-spectral-locus}
exists if and only if
\begin{equation}\label{eq:fixed-spectral-dominance}
                 \mu_i\preceq\lambda_i\quad\hbox{for every }i.
\end{equation}
Every such pair is realized over a smooth pointed algebraic curve,
after shrinking at the special point, with \(A\) fixed and with
constant full determinant and reduced rank loci.
\item Every such pencil family induces a flat projective family of
\(Z_C\). After shrinking the curve, every truncation of order
\(k<d\) is constant. If at least one dominance inequality is
strict, the generic and special \(Z_C\) are nonisomorphic;
their first distinguishing order is \(d\).
\item The recovered spectral module at \(\alpha_i\) has length
profile
\begin{equation}\label{eq:fixed-spectral-lengths}
                 L_{i,a}=\sum_j\min(a,\lambda_{i,j}).
\end{equation}
Every \(L_{i,a}\) can only increase in the stated specialization.
All these inequalities together are equivalent to
\eqref{eq:fixed-spectral-dominance}, giving an intrinsic numerical
criterion for every admissible transition.
\end{enumerate}
\end{theorem}
\begin{proof}
We use Lemma~\ref{lem:orthogonal-orbits-v145}.  Its full-orthogonal
classification and dominance order have exactly the group and
self-adjoint convention needed here.  The comparison for Grassmannian
Segre strata is \cite[Theorem~5.1]{FevolaMandelshtamSturmfels}.

The primary spaces \(V_i=\ker(C-\alpha_iI)^{m_i}\) are an
orthogonal direct sum: the Chinese remainder idempotents are
polynomials in the self-adjoint \(C\), hence self-adjoint projections.
Each restricted form is nondegenerate, and
\(C|_{V_i}=\alpha_iI+N_i\), with \(N_i\) self-adjoint
nilpotent. The prescribed corank is its number of Jordan blocks.
The classical theorem gives the indexing. Conversely, a Jordan block
is self-adjoint for the reversed-identity Gram matrix; an isometry
with the chosen nondegenerate form realizes each partition.

The space of orthogonal decompositions has dimension
\((n^2-\sum_i m_i^2)/2\). Adding the nilpotent orbit dimensions
proves \eqref{eq:fixed-spectral-orbit-dimension}. These are dimensions
in the fixed-\(A\) matrix space, not in an unmarked moduli stack.
Full orthogonal orbits need not be connected; no irreducibility
assertion is needed. Theorem~\ref{thm:spectral-finite-torelli}
identifies the unmarked neighbourhood with the pencil, and its spectral
module determines the partition at each root. A projective change of
the line can therefore do exactly the stated root permutations.
For one or two roots the projective stabilizer may be infinite;
only its finite permutation image is used.

For necessity, the same polynomial idempotents decompose a family
with fixed characteristic polynomial into subbundles of ranks \(m_i\).
The rank of \(N_i^a\) cannot increase at specialization. Since
\(m_i-\rank N_i^a=\sum_j\min(a,\lambda_{i,j})\), this gives
the length inequalities. Summing the first \(a\) columns of
conjugate Young diagrams shows that they are equivalent to dominance:
conjugation reverses the dominance order.

For sufficiency, fix an orthogonal decomposition of \((V,A)\).
The closure criterion puts a chosen representative of \(\mu_i\)
in the closure of the orbit for \(\lambda_i\). Thus the special
tuple lies in the closure of the product of the generic orbits.
Choose an irreducible component of that closure containing the point
and meeting the generic orbit product. Successive general hyperplane
sections through the point, avoiding the closed complement, give an
integral algebraic curve with generic point in that open orbit product.
Normalize it and choose a point over the prescribed point. A normal
complex curve is smooth. Remove the other points outside the generic
orbit product. A constant curve suffices in the equality case.
The resulting block operator \(\bigoplus_i(\alpha_iI+N_i(t))\)
is self-adjoint for the same \(A\). Its determinant is unchanged,
and its coranks agree generically and specially, hence throughout the
shrunk curve. This gives the required specialization. It does not
assert constancy of the full higher Fitting ideals.

For part~(3), pull back Theorem~\ref{thm:universal-finite-neighbourhood}
along the curve's morphism to \(P=\Gr(2,\Sym^2V)\).
That theorem identifies the family with the actual ideal-adic
neighbourhoods and proves flatness and projectivity under arbitrary
base change. No inference from fibrewise isomorphisms is needed.
An affine Grassmannian chart around the special pencil supplies a
fixed complement to nearby \(R_C\), identifying \(S_{R_C}\),
\(B\), and \(E\). There is no relation below degree \(d\),
so all lower neighbourhoods become constant. For a strict dominance
transition, \(\sum_{i,a}(\lambda'_{i,a})^2\) strictly increases;
this sum is invariant under permutations of the decorated roots.
The generic and special tuples therefore cannot be related by such
a permutation. Part~(1) separates their unmarked degree-\(d\)
neighbourhoods. Flatness is asserted only for these finite
neighbourhoods, not unrestricted failure schemes.
\end{proof}

\begin{corollary}[Exactly when the coarse data lose a labelled class]
\label{cor:fixed-spectral-ambiguity}
The number of root-labelled congruence classes with the prescribed
data is \(\prod_i\#\mathcal P(m_i,c_i)\). It is one precisely
when, for every \(i\), either \(c_i=1\) or \(m_i-c_i\le1\).
The first possible ambiguity with fixed full determinant and all
reduced rank loci occurs in dimension four, at \((m,c)=(4,2)\),
with partitions \((3,1)\) and \((2,2)\).
\end{corollary}
\begin{proof}
Subtracting one from every part identifies \(\mathcal P(m,c)\)
with partitions of \(m-c\) into at most \(c\) parts.
There is one if \(c=1\) or \(m-c\le1\). Otherwise
\((m-c)\) and \((m-c-1,1)\) are distinct admissible partitions.
The theorem gives the count and the minimal example.
\end{proof}

\subsection{The first relation as a standard graded Betti space}

The finite invariant has a usual Hilbert-scheme interpretation.
The first Betti space means the generator space in the first free
module of a minimal graded resolution of the quotient, not the space
of relations among generators of its defining ideal. The following
statement retains the marked quotient presentation.

\begin{proposition}[A graded Hilbert locus and its first Betti space]
\label{prop:first-syzygy-hilbert}
In the fixed coordinates of Theorem~\ref{thm:finite-parameter-immersion},
let \(S=\Sym H^*\), \(N=\binom{n+1}{2}\), \(M=\binom N2\),
and \(Q_R=S/((K_R)+S_{\ge d+1})\). Then
\begin{equation}\label{eq:first-relation-tor}
                     \Tor_1^S(Q_R,\C)_d=K_R.
\end{equation}
This identity commutes with arbitrary complex base change.
The functor of graded quotients with Hilbert function
\[
 h_j=\begin{cases}\dim S_j,&0\le j<d,\\
 \dim S_d-M,&j=d,\\0,&j>d\end{cases}
\]
is represented by \(\Gr(M,S_d)\). The pencil locus is a closed
copy of \(\Gr(2,\Sym^2V)\). Its degree-\(d\) first Betti
space recovers every scheme-theoretic pencil parameter, including
parameters over nonreduced bases.
\end{proposition}
\begin{proof}
Put \(J=(K_R)+S_{\ge d+1}\) and \(S_+=\bigoplus_{j>0}S_j\).
The augmented sequence \(0\to J\to S\to Q_R\to0\) gives
\(\Tor_1^S(Q_R,\C)=J/S_+J\), since its map to \(S/S_+\)
is zero. Its degree-\(d\) part is \(K_R\). The same argument
over a base algebra applies to the degree-\(d\) subbundle with
locally free quotient and commutes with pullback.
A graded quotient with the indicated locally free components has
zero ideal below degree \(d\), a rank-\(M\) subbundle in degree
\(d\), and the full ideal in higher degrees. Conversely, every
such subbundle defines an ideal, as its positive-degree products lie
above degree \(d\). This is the Grassmannian functor. The closed
immersion and its inverse on the image are
Theorem~\ref{thm:finite-parameter-immersion}. This is not an inference
from closed-point Torelli to unmarked deformation-groupoid equivalence.
\end{proof}

% BEGIN PRESERVED PART 17-assistance.tex
\section*{Acknowledgment of assistance}
AI assistance was used in developing and drafting the arguments and
finite verification scripts.  All structural proofs are supplied for
independent mathematical scrutiny; computational receipts are not
formal proof certificates.  No AI system is listed as an author.

% BEGIN PRESERVED PART 43-assistance-specific-v151.tex
For the present revision, GPT-6 Astra Pro assisted specifically with the
incidence normalization and overlap criterion in
Section~\ref{sec:canonical-boundary-v151}, the universal-property proof
in Section~\ref{sec:rigidification-v151}, the explicit specialization
in Section~\ref{sec:closed-pencil-v151}, and the associated exact tests
and exposition.  These arguments require independent human verification.
No statement here represents a completed journal review or a formal
proof-assistant certification.

% BEGIN PRESERVED PART 47-assistance-specific-v152.tex
\paragraph{Additional disclosure for the present revision.}
GPT-6 Astra Pro assisted specifically with the conductor and collision
models, the scalar-matrix degeneration, the flag-jet calculation,
the reciprocal-polynomial fibre presentation, literature comparisons,
and exact computational checks.  The proofs above state the complete
mathematical arguments; finite computations are consistency checks,
not universal proof certificates.  No journal decision or historical
priority determination is inferred from the build and source audits.

\section*{Assistance in the present revision}
An AI assistant helped formulate and write the binary conductor stratification,
the simultaneous incidence atlas, the evaluated singular Fitting ideal, and
the minimal-relation-module argument for pure-power fibres. It also prepared
the source-preserving reorganization and finite exact consistency checks.
The written proofs, rather than those finite checks, are the mathematical
basis of the statements and are submitted to independent scrutiny.

\paragraph{Revision 154.}
AI assistance was used for the polarized moduli compatibility,
spectral divisor-fibre construction, homological calculations and
proof expansions. The fixed five-variable Betti table uses exact
computer-assisted rational elimination with its complete finite
inputs recorded. The general theorems rely on the written arguments
and the explicitly attributed classical results, not on regression
counts or source-preservation receipts.

\paragraph{Revision 155.}
AI assistance was used to formulate and draft the exact quadratic
coefficient-section theorem and the reciprocal power-ideal description
of spectral Fitting schemes, to check their finite examples, and to
prepare the companion-paper organization.  Classical apolar and
orthogonal-orbit results are explicitly attributed.  The general
claims rely on the written proofs; finite exact tests and source
preservation are not formal proof certificates.

\paragraph{Revision 157.}
AI assistance was used in deriving and drafting the universal power-ideal
formula, its complete-quadric graph interpretation, the contact-divisor
and resolved-pencil statements, and their exact finite checks. Classical
determinantal, apolar and complete-quadric results are attributed at their
points of use. The proofs, rather than the computations, support the
general statements. No formal proof certification or journal endorsement
is claimed.

\paragraph{Scope of the new intrinsic construction.}
The source-normalized envelope is an algebra bundle with an intrinsic
projective power diagram. It is not a claimed preferred subquotient
of the sharp failure algebra. Singular-pencil classification is read
through the classical complex symmetric Kronecker form. The explicit
Hilbert tail is a boundary calculation for the stated transverse
family, not a classification of every fibre of the incidence
compactification. Exact computations accompanying the source check
finite instances; they do not certify the general proofs.

\paragraph{Revision 169.}
AI assistance was used in deriving and drafting the all-order row
factorization, monic chart equations, associated-point and genus-correction
calculations, equivariant chart gluing, Quot comparison, and finite exact
checks. The written arguments are offered for independent mathematical
scrutiny. No external independent audit of the sharp inverse, formal
proof-assistant certification, or journal endorsement is represented as
completed by this revision.


\paragraph{Revision 170.}
AI assistance was used in deriving and drafting the ramified parameter-fibre,
conductor, and normalization-label arguments and in preparing auxiliary
exact checks. The written proofs are submitted for independent scrutiny.
Compilation, finite checks, and source preservation are not represented as
an external proof audit, a formal proof certificate, or a journal endorsement.

\begin{thebibliography}{99}

\bibitem{Ballico93}
E.~Ballico,
On the failure locus of higher order properties of embeddings in
projective spaces,
\emph{Math. Nachr.} \textbf{163} (1993), 5--13.
doi:10.1002/mana.19931630102.

\bibitem{Ballico96}
E.~Ballico,
On the failure cycles for the quadratic normality of a projective
variety,
\emph{Pacific J. Math.} \textbf{172} (1996), 307--313.

\bibitem{ChengWang}
S.-J. Cheng and W. Wang, Howe duality for Lie superalgebras,
\emph{Compositio Math.} \textbf{128} (2001), 55--94;
arXiv:math/0008093v1, \S2 and Theorem 3.5.

\bibitem{FevolaMandelshtamSturmfels}
C.~Fevola, Y.~Mandelshtam, and B.~Sturmfels,
Pencils of quadrics: old and new,
\emph{Le Matematiche} \textbf{76} (2021), no.~2, 319--335.
doi:10.4418/2021.76.2.2; arXiv:2009.04334v2,
Theorem~1.1, Corollary~2.1, Remark~2.3, Conjecture~4.5, and Theorem~5.1.

\bibitem{Ohta1986}
T.~Ohta, The singularities of the closures of nilpotent orbits in certain
symmetric pairs, \emph{Tohoku Math. J. (2)} \textbf{38} (1986), 441--468.
doi:10.2748/tmj/1178228456.

\bibitem{SmithQuadrics}
I.~Smith, Floer cohomology and pencils of quadrics,
\emph{Invent. Math.} \textbf{189} (2012), 149--250.
doi:10.1007/s00222-011-0364-1; arXiv:1006.1099, \S4.2.

\bibitem{EliasRossiShort}
J.~Elias and M.~E.~Rossi,
Isomorphism classes of short Gorenstein local rings via Macaulay's
inverse system,
\emph{Trans. Amer. Math. Soc.} \textbf{364} (2012), no.~9, 4589--4604.
doi:10.1090/S0002-9947-2012-05430-4;
arXiv:0911.3565v1, Theorem~3.3 and Corollaries~3.4--3.5.

\bibitem{EliasRossiCompressed}
J.~Elias and M.~E.~Rossi,
Analytic isomorphisms of compressed local algebras,
\emph{Proc. Amer. Math. Soc.} \textbf{143} (2015), no.~3, 973--987.
doi:10.1090/S0002-9939-2014-12313-6;
arXiv:1207.6919v1, Theorem~3.1 and Examples~3.3--3.5.

\bibitem{MarcusMoyls}
M.~Marcus and B.~N.~Moyls,
Transformations on tensor product spaces,
\emph{Pacific J. Math.} \textbf{9} (1959), no.~4, 1215--1221.
doi:10.2140/pjm.1959.9.1215, Theorem~1.

\bibitem{StacksGroups}
The Stacks Project Authors,
\emph{The Stacks Project}, Section~39.8, Lemma~39.8.2,
\url{https://stacks.math.columbia.edu/tag/0BF6}, accessed September 24, 2026.
\bibitem{Milne2017}
J.~S. Milne, \emph{Algebraic Groups: The Theory of Group Schemes of
Finite Type over a Field}, Cambridge Studies in Advanced Mathematics
\textbf{170}, Cambridge University Press, Cambridge, 2017.

\bibitem{StacksClosedImmersion}
The Stacks Project Authors, \emph{The Stacks Project},
Lemma~41.7.2 (Tag 04XV),
\url{https://stacks.math.columbia.edu/tag/04XV}, accessed September 24, 2026.

\bibitem{StacksQuotient}
The Stacks Project Authors, \emph{The Stacks Project},
Theorem~97.17.2 and Section~97.18,
\url{https://stacks.math.columbia.edu/tag/06PI}, accessed September 24, 2026.

\bibitem{StacksSmoothQuotient}
The Stacks Project Authors, \emph{The Stacks Project},
Lemma~101.33.7 (Tag 0DLS),
\url{https://stacks.math.columbia.edu/tag/0DLS}, accessed September 24, 2026.

\bibitem{StacksFiniteV151}
The Stacks Project Authors, \emph{The Stacks Project},
Lemma~37.44.1 (Tag 02LS),
\url{https://stacks.math.columbia.edu/tag/02LS}, accessed September 24, 2026.

\bibitem{StacksImageV151}
The Stacks Project Authors, \emph{The Stacks Project},
Section~29.6, Scheme theoretic image (Tag 01R5),
\url{https://stacks.math.columbia.edu/tag/01R5}, accessed September 24, 2026.

\bibitem{StacksNormalV151}
The Stacks Project Authors, \emph{The Stacks Project},
Section~29.55, Normalization (Tag 035E),
\url{https://stacks.math.columbia.edu/tag/035E}, accessed September 24, 2026.

\bibitem{StacksRigidV151}
The Stacks Project Authors, \emph{The Stacks Project},
Subsection~112.5.7, Rigidification (Tag 04V2), including the discussion
of noncentral normal inertia subgroups,
\url{https://stacks.math.columbia.edu/tag/04V2}, accessed September 24, 2026.


\bibitem{Chipalkatti2003}
J.~V. Chipalkatti, On equations defining coincident root loci,
\emph{J. Algebra} \textbf{267} (2003), 246--271.
doi:10.1016/S0021-8693(03)00336-3;
arXiv:math/0110224v1, Proposition~5.1, Theorem~5.4,
and Corollary~5.8.

\bibitem{Kurmann2012}
S.~Kurmann, Some remarks on equations defining coincident root loci,
\emph{J. Algebra} \textbf{352} (2012), 223--231.
doi:10.1016/j.jalgebra.2011.10.045;
arXiv:1108.4532v1, Proposition~1.3, Theorem~1.5,
and Proposition~2.1.

\bibitem{Ferrand2003}
D.~Ferrand, Conducteur, descente et pincement,
\emph{Bull. Soc. Math. France} \textbf{131} (2003), no.~4, 553--585.
doi:10.24033/bsmf.2455.

\bibitem{AOV2008}
D.~Abramovich, M.~Olsson, and A.~Vistoli, Tame stacks in positive
characteristic, \emph{Ann. Inst. Fourier (Grenoble)}
\textbf{58} (2008), no.~4, 1057--1091.
doi:10.5802/aif.2378; Appendix~A, Theorem~A.1.

\bibitem{HuLinShao2007}
Y.~Hu, J.~Lin, and Y.~Shao, A compactification of the space of
algebraic maps from $\Pj^1$ to $\Pj^n$,
arXiv:math/0701255v1 (2007), Theorem~1.1,
\S2.2, and Corollary~3.6.


\bibitem{Grinberg2019}
D.~Grinberg, \emph{A basis for a quotient of symmetric polynomials},
arXiv:1910.00207, version of September 24, 2021, Theorem~2.7.

\bibitem{MFK1994}
D.~Mumford, J.~Fogarty, and F.~Kirwan,
\emph{Geometric Invariant Theory}, third edition,
Ergebnisse der Mathematik und ihrer Grenzgebiete (2), vol.~34,
Springer-Verlag, Berlin, 1994.

\bibitem{MillerSturmfels2005}
E.~Miller and B.~Sturmfels, \emph{Combinatorial Commutative Algebra},
Graduate Texts in Mathematics, vol.~227, Springer, New York, 2005.
Chapter~5 (Stanley--Reisner rings and Hochster's formula).

\bibitem{ShafieiSymmetric}
S.~Shafiei, \emph{Apolarity for determinants and permanents of generic
symmetric matrices}, arXiv:1303.1860v3 (2021), Proposition~3.1,
Theorem~3.10, and Corollary~2.7.

\bibitem{Krishna2026}
A.~Krishna, \emph{A universal discriminant formula for pencils of
quadrics}, arXiv:2607.05340v2 (2026), Theorem~1.1 and Lemma~3.1.

\bibitem{StacksLocalV154}
The Stacks Project Authors, \emph{The Stacks Project},
Tags \href{https://stacks.math.columbia.edu/tag/00R4}{00R4}
(miracle flatness),
\href{https://stacks.math.columbia.edu/tag/07QV}{07QV}
(quasi-excellent rings are Nagata), and
\href{https://stacks.math.columbia.edu/tag/0BJ0}{0BJ0}
(geometrically reduced formal fibres), accessed September 25, 2026.

\bibitem{NagaokaWachi2024}
T.~Nagaoka and A.~Wachi,
\emph{The strong Lefschetz property of Gorenstein algebras generated
by relative invariants}, arXiv:2403.05492v1 (2024),
Theorem~4, Example~5, Example~9, Proposition~10, Remark~11,
Lemma~13 and Proposition~16.
Published version: doi:10.1090/proc/16870.

\bibitem{BrunsConca2003}
W.~Bruns and A.~Conca,
\emph{Gr\"obner bases and determinantal ideals},
arXiv:math/0302058v3 (2003), Lemma~2.5 and Theorem~2.4.

\bibitem{HenriquesVarbaro2016}
I.~Henriques and M.~Varbaro,
\emph{Test, multiplier and invariant ideals},
Adv. Math. \textbf{287} (2016), 704--732,
doi:10.1016/j.aim.2015.09.028; arXiv:1407.4324.

\bibitem{Vainsencher1982}
I.~Vainsencher,
\emph{Schubert calculus for complete quadrics},
in Enumerative Geometry and Classical Algebraic Geometry (Nice, 1981),
Progress in Mathematics \textbf{24}, Birkh\"auser (1982), 199--235.

\bibitem{Massarenti2020}
A.~Massarenti,
\emph{On the birational geometry of spaces of complete forms I:
collineations and quadrics},
Proc. Lond. Math. Soc. \textbf{121} (2020), 1579--1618,
doi:10.1112/plms.12377; arXiv:1803.09161v2, \S2.

\bibitem{CasarottiCornianiMassarenti2023}
A. Casarotti, E. Corniani and A. Massarenti,
\emph{Complete singular collineations and quadrics},
International Mathematics Research Notices \textbf{2023} (2023),
no.~18, 15370--15407. DOI: 10.1093/imrn/rnac271.
The cited construction is also in arXiv:2111.02940,
Definition~2.5, Remark~2.6 and Theorem~2.14.

\bibitem{Thompson1991}
R. C. Thompson, \emph{Pencils of complex and real symmetric and skew
matrices}, Linear Algebra and its Applications \textbf{147} (1991),
323--371. DOI: 10.1016/0024-3795(91)90238-R.


\bibitem[Stacks Project]{StacksProject160}
The Stacks Project Authors, \emph{The Stacks Project},
Tags 0806 (universal property of blowing up), 02LS (proper quasi-finite
morphisms), and 0H1G (smooth centres and their blow-ups),
\url{https://stacks.math.columbia.edu}, accessed September 25, 2026.

\bibitem[Li(2009)]{Li2009v161}
L. Li, Wonderful compactification of an arrangement of subvarieties,
\emph{Michigan Math. J.} \textbf{58} (2009), 535--563;
\url{https://arxiv.org/abs/math/0611412}.
\bibitem[De Teran--Dmytryshyn--Dopico(2018)]{DeTeranDmytryshynDopico2018v161}
F. De Ter\'an, A. Dmytryshyn and F. M. Dopico,
\emph{Generic symmetric matrix pencils with bounded rank},
arXiv:1808.03118 (2018), Theorem~2.1;
\url{https://arxiv.org/abs/1808.03118}.
\bibitem[Gathmann(1998)]{Gathmann1998v161}
A. Gathmann, \emph{Gromov--Witten invariants of blow-ups},
arXiv:math/9804043 (1998), Section~7;
\url{https://arxiv.org/abs/math/9804043}.
\bibitem[Ballico(1996)]{Ballico1996v161}
E. Ballico, On the failure cycles for the quadratic normality of a
projective variety, \emph{Pacific J. Math.} \textbf{172} (1996),
307--314.
\bibitem[Stacks Project]{Stacks161}
The Stacks Project Authors, \emph{The Stacks Project},
Section~29.55 (normalization), Tags 080C (strict transform),
0806 (universal property of blowing up), and 02LS (finite morphisms),
\url{https://stacks.math.columbia.edu}, accessed September 25, 2026.

\bibitem[Dmytryshyn(2011/2018)]{Dmytryshyn162}
A. Dmytryshyn, \emph{Miniversal deformations of pairs of symmetric
matrices under congruence}, arXiv:1104.2530, version 2 (2018),
Theorem~2.1, Corollary~2.1 and Lemma~3.1;
\url{https://arxiv.org/abs/1104.2530}.
\bibitem[Hu--Lin--Shao(2011)]{HuLinShao162}
Y. Hu, J. Lin and Y. Shao, A compactification of the space of algebraic
maps from $\Pj^1$ to $\Pj^n$, \emph{Comm. Anal. Geom.}
\textbf{19} (2011), 1--30; arXiv:math/0701255.
\bibitem[Chung--Hong--Kiem]{ChungHongKiem162}
K. Chung, J. Hong and Y.-H. Kiem, \emph{Compactified moduli spaces
of rational curves in projective homogeneous varieties},
arXiv:1010.0068; \url{https://arxiv.org/abs/1010.0068}.
\bibitem[Stacks Project]{Stacks162}
The Stacks Project Authors, \emph{The Stacks Project},
Tags 062Y (relative effective Cartier divisors), 02LS (finite morphisms),
and 0806 (the blow-up universal property),
\url{https://stacks.math.columbia.edu}, accessed September 25, 2026.

\bibitem[Bhatt(2014)]{Bhatt163}
B. Bhatt, \emph{Algebraization and Tannaka duality},
arXiv:1404.7483, 2014, Theorem 1.4 and \S5.
\bibitem[Marian--Oprea--Pandharipande(2011)]{MOP163}
A. Marian, D. Oprea, and R. Pandharipande,
\emph{The moduli space of stable quotients}, Geom. Topol.
\textbf{15} (2011), 1651--1706; arXiv:0904.2992.
\bibitem[Ciocan-Fontanine--Kim--Maulik(2011)]{CFKM163}
I. Ciocan-Fontanine, B. Kim, and D. Maulik,
\emph{Stable quasimaps to GIT quotients}, arXiv:1106.3724, 2011.
\bibitem[Manolache(2013)]{Manolache163}
C. Manolache, \emph{Stable maps and stable quotients},
arXiv:1301.4393, 2013.
\bibitem[Lahoz--Lehn--Macr\`i--Stellari(2016)]{LLMS163}
M. Lahoz, M. Lehn, E. Macr\`i, and P. Stellari,
\emph{Generalized twisted cubics on a cubic fourfold as a moduli
space of stable objects}, arXiv:1609.04573, 2016.
\bibitem[Stacks Project]{Stacks163}
The Stacks Project Authors, \emph{The Stacks Project},
Lemma 37.53.6 (connected fibres of a proper birational morphism)
and Lemma 41.11.3 (completed local criterion for etaleness),
\url{https://stacks.math.columbia.edu}, accessed September 25, 2026.

\bibitem[Stacks Project]{Stacks164}
The Stacks Project Authors, \emph{The Stacks Project},
Tag 0539 (torsion-free modules over valuation rings),
\url{https://stacks.math.columbia.edu/tag/0539}, accessed September 26, 2026.
\bibitem[Chung and Moon(2017)]{ChungMoon164}
K.~Chung and H.-B.~Moon, \emph{Mori's program for the moduli space
of conics in Grassmannian}, Taiwanese J. Math. \textbf{21} (2017),
621--652; arXiv:1608.00181.

\bibitem[Grothendieck(1961)]{Grothendieck166}
A. Grothendieck, \emph{Techniques de construction et th\'eor\`emes
d'existence en g\'eom\'etrie alg\'ebrique IV: les sch\'emas de Hilbert},
S\'eminaire Bourbaki, exp. 221 (1960/61), 249--276; with the published
erratum to exp. 221.
\bibitem[Stacks Project]{Stacks166}
The Stacks Project Authors, \emph{The Stacks Project},
Tags 00MI (fibrewise exactness and flatness), 0BJ0 (geometrically
reduced formal fibres), 05PS (universal flattening), and 08S3
(square-zero extensions and the naive cotangent complex),
\url{https://stacks.math.columbia.edu}, accessed September 26, 2026.
\bibitem[Li(2001)]{JunLi166}
J. Li, \emph{Stable morphisms to singular schemes and relative
stable morphisms}, J. Differential Geom. \textbf{57} (2001),
509--578; arXiv:math/0009097.
\bibitem[Li and Wu(2011)]{LiWu166}
J. Li and B. Wu, \emph{Good degeneration of Quot-schemes and
coherent systems}, arXiv:1110.0390, 2011.


\bibitem[Gotzmann(1978)]{Gotzmann167}
G.~Gotzmann, \emph{Eine Bedingung f\"ur die Flachheit und das
Hilbertpolynom eines graduierten Ringes}, Math. Z. \textbf{158}
(1978), 61--70, doi:10.1007/BF01214566.
\bibitem[Bayer--Morrison(1988)]{BayerMorrison167}
D.~Bayer and I.~Morrison, \emph{Standard bases and geometric invariant
theory. I. Initial ideals and state polytopes}, J. Symbolic Comput.
\textbf{6} (1988), no.~2--3, 209--217,
doi:10.1016/S0747-7171(88)80043-9.


\bibitem[Stacks Project]{Stacks169}
The Stacks Project Authors, \emph{The Stacks Project},
Section~31.33, ``Blowing up'', Tag~01OF,
\url{https://stacks.math.columbia.edu/tag/01OF}.
\bibitem[Speyer--Sturmfels(2004)]{SpeyerSturmfels169}
D.~Speyer and B.~Sturmfels, \emph{The tropical Grassmannian},
Adv. Geom. \textbf{4} (2004), 389--411;
arXiv:math/0304218.
\bibitem[Manubens--Montes(2006)]{ManubensMontes169}
M.~Manubens and A.~Montes, \emph{Improving DISPGB algorithm using
the discriminant ideal}, J. Symbolic Comput. \textbf{41} (2006),
1245--1263; arXiv:math/0601763.
\bibitem[Marian--Oprea--Pandharipande(2011)]{MOP169}
A.~Marian, D.~Oprea, and R.~Pandharipande,
\emph{The moduli space of stable quotients}, Geom. Topol.
\textbf{15} (2011), 1651--1706; arXiv:0904.2992.
\bibitem[Manolache(2014)]{Manolache169}
C.~Manolache, \emph{Stable maps and stable quotients},
Compos. Math. \textbf{150} (2014), 1457--1481; arXiv:1301.4393.


\bibitem{GonzalezTeissier170}
P.~D. Gonz\'alez P\'erez and B. Teissier,
\emph{Toric geometry and the Semple--Nash modification},
Rev. R. Acad. Cienc. Exactas F\'is. Nat. Ser. A Mat. RACSAM,
DOI: \href{https://doi.org/10.1007/s13398-012-0096-0}{10.1007/s13398-012-0096-0}.
\bibitem{AlexeevPardini170}
V. Alexeev and R. Pardini,
\emph{On the existence of ramified abelian covers},
Rend. Sem. Mat. Univ. Politec. Torino \textbf{71} (2013), 307--315.

\end{thebibliography}

\end{document}
